\documentclass[11pt,oneside]{book}
\usepackage[utf8]{inputenc}
\usepackage[T1]{fontenc}
\usepackage{lmodern}
\usepackage{microtype}
\usepackage{amsmath,amssymb}
\usepackage{geometry}
\usepackage{hyperref}
\usepackage{enumitem}
\usepackage{array}
\usepackage{longtable}
\usepackage{tocloft}
\usepackage{tikz}
\usetikzlibrary{arrows.meta,positioning,calc,decorations.pathreplacing}
\usepackage{pgfplots}
\pgfplotsset{compat=1.18}
\geometry{paperwidth=6in,paperheight=9in,margin=0.72in}
\hypersetup{colorlinks=true,linkcolor=blue,urlcolor=blue}
\setlength{\cftchapnumwidth}{3.2em}
\emergencystretch=2em
\newcommand{\bookfigure}[2]{%
\begin{figure}[htbp]
\centering
\input{figures/#1.tex}
\caption{#2}
\end{figure}}

% Chapters supply their own References heading.  Suppress the automatic
% heading normally produced by thebibliography while retaining \cite links.
\makeatletter
\renewenvironment{thebibliography}[1]
  {\list{\@biblabel{\@arabic\c@enumiv}}%
    {\settowidth\labelwidth{\@biblabel{#1}}%
     \leftmargin\labelwidth
     \advance\leftmargin\labelsep
     \usecounter{enumiv}%
     \let\p@enumiv\@empty
     \renewcommand\theenumiv{\@arabic\c@enumiv}}}
  {\endlist}
\makeatother

% Historical development is presented as a numbered sequence in every chapter.
% The local paragraph redefinition keeps the existing explanatory prose while
% making each historical point an enumerated item.
\newenvironment{historylist}{%
\begin{enumerate}[leftmargin=*,itemsep=0.45em]
\renewcommand{\paragraph}[1]{\item \textbf{##1.}\par}
}{%
\end{enumerate}}

\title{\Huge Mathematics in Motion\\[6pt]\large A Friendly Guide to Mathematical Theories\\[4pt]\normalsize A survey with worked examples and routes for further study}
\author{A popular science introduction}
\date{}

\makeatletter
\newcommand{\applicationlist}[1]{%
\begin{enumerate}[leftmargin=*,itemsep=0.35em]
\@for\appitem:=#1\do{\item \appitem}
\end{enumerate}}
\makeatother

\newcommand{\theory}[5]{\clearpage\chapter{#1}
\def\theoryapplicationtitle{#1}
\def\theoryapplicationitems{#4}
\def\theoryapplicationlink{#5}
\section*{1. Overview}
\subsection*{Intuitive}
Begin with the problem, not the notation. Ask what objects the theory studies,
what change or operation is allowed, and what feature the theory tries to
keep reliable. #2
\subsection*{Formal}
#3
\section*{2. History}
\subsection*{Historical development}
\begin{historylist}}

% The applications section is emitted after a chapter's theory text.  Keeping
% it here avoids placing historical or theoretical paragraphs after the
% applications list in chapters that use the compact \theory macro.
\newcommand{\theoryapplications}{%
\section*{5. Applications}
\noindent\theoryapplicationlink\par
\applicationlist{\theoryapplicationitems}}

% Kept for the few chapters that use a family-specific exercise set.  The
% prompts are anchored to the chapter's objects, calculation, central result,
% boundary, and application; generic prompts are not used by the book.
\newcommand{\samplequestions}[2]{%
\section*{7. Sample Questions}
These short questions review the main ideas of \textbf{#1} at an
introductory undergraduate level.  Each answer uses the definitions and
methods developed in this chapter.
\begin{enumerate}[leftmargin=*,itemsep=0.9em]
\ifcase#2
\item \textbf{Structure.} State the basic objects studied in #1 and give one
example and one non-example.\\
\emph{Solution.} The objects are the things named in the chapter's core
definition.  The example must satisfy every stated rule; the non-example
fails at least one rule.  Checking the rules one at a time is the reliable
method, rather than judging by appearance.
\item \textbf{Operation.} Choose two simple objects from #1 and carry out the
main operation described in the chapter. What must be checked afterward?\\
\emph{Solution.} Perform the operation and then check that the result is
still an allowed object. This closure check is what turns a calculation into
a statement inside the theory.
\item \textbf{Invariant.} What feature remains unchanged when the permitted
transformations of #1 are applied?\\
\emph{Solution.} Identify the transformation, compute the feature before and
after it, and compare. An invariant is useful because it can prove that two
objects cannot be transformed into one another.
\item \textbf{Map.} Give a small example of a structure-preserving map in
#1. Which information does it preserve?\\
\emph{Solution.} A structure-preserving map sends allowed objects to allowed
objects and respects the main operation. Its preserved information is exactly
what appears in the defining rule; extra labels or drawings may disappear.
\item \textbf{Application.} Explain one practical use of #1 and identify its
modelling assumption.\\
\emph{Solution.} The theory supplies a compact language for the repeated
pattern in the application. The assumption is the idealisation that makes the
pattern exact; a real system must be checked against it before conclusions are
trusted.
\or
\item \textbf{Approximation.} Approximate a smooth function on a short
interval by its tangent-line formula at the midpoint. Why is the error small
near that point?\\
\emph{Solution.} The tangent line matches both the function value and its
first-order change at the midpoint. The remaining error is of second order
in the distance from that point, so halving the distance makes it roughly
four times smaller when the second derivative is bounded.
\item \textbf{Convergence.} In an asymptotic statement, what does
$f(x)\sim g(x)$ as $x\to\infty$ mean?\\
\emph{Solution.} It means $f(x)/g(x)\to1$. Thus the two quantities need not
be close in absolute size, but their relative difference tends to zero.
\item \textbf{Error control.} If an approximation has error at most
$0.01$ on an interval, what can and cannot be concluded?\\
\emph{Solution.} The computed value differs from the exact value by no more
than $0.01$ wherever the stated assumptions hold. Nothing follows outside
that interval or when the assumptions, such as smoothness, fail.
\item \textbf{Choice.} Why might a polynomial approximation be preferred
to a complicated exact expression in #1?\\
\emph{Solution.} A polynomial is easy to evaluate, differentiate, integrate,
and implement in a computer. The preference is justified only after its
error and numerical stability have been measured.
\item \textbf{Application.} Describe a situation in which an approximation
is safer than an exact-looking formula.\\
\emph{Solution.} In numerical simulation, a controlled approximation with a
known error is safer than an exact formula evaluated with rounding or unstable
subtraction. The guarantee must accompany the approximation.
\or
\item \textbf{Conditional probability.} A disease affects $1\%$ of a
population. A test detects $99\%$ of cases and has a $5\%$ false-positive
rate. What is the probability that a randomly selected positive result is a
true case?\\
\emph{Solution.} Among $10,000$ people, about $99$ sick people test positive
and $495$ healthy people test positive. The probability is
$99/(99+495)=1/6$, approximately $16.7\%$.
\item \textbf{Expectation.} A fair six-sided die is rolled once. Find the
expected score and explain why it need not be the score obtained.\\
\emph{Solution.} The expectation is
$(1+2+3+4+5+6)/6=3.5$. A single roll must be an integer; expectation is the
long-run average of repeated rolls, not a prediction of one roll.
\item \textbf{Independence.} Two fair coins are tossed. What is the
probability of exactly one head, and which independence assumption was used?\\
\emph{Solution.} The outcomes HT and TH are disjoint, each with probability
$1/4$, so the answer is $1/2$. Independence lets us multiply the probabilities
of the two separate tosses.
\item \textbf{Statistics.} Why does increasing a random sample usually make
an estimated average more reliable?\\
\emph{Solution.} Random positive and negative deviations tend to cancel.
The law of large numbers makes the sample average approach the population
average, although biased sampling or dependence can defeat the conclusion.
\item \textbf{Model check.} Give one reason a probability calculation can be
correct but its prediction still be poor.\\
\emph{Solution.} The calculation may use an incorrect base rate, dependence
assumption, or sampling model. Logic proves the consequence of assumptions;
it does not prove that the assumptions describe the world.
\or
\item \textbf{Graph model.} Model three cities joined in a triangle as a
graph. How many edges and vertex degrees does it have?\\
\emph{Solution.} There are three vertices and three edges. Every vertex has
degree two, since it is joined to the other two vertices.
\item \textbf{Shortest route.} Why does repeatedly extending the currently
shortest unfinished route work for a graph with nonnegative edge lengths?\\
\emph{Solution.} Any route reaching an unfinished vertex through another
unfinished vertex is at least as long as the best currently known route to
that vertex. Nonnegative lengths prevent a later detour from becoming
shorter; this is the key justification behind Dijkstra's method.
\item \textbf{Counting.} How many edges can a simple graph with $n$ vertices
have at most?\\
\emph{Solution.} At most one edge joins each pair, and there are
$\binom n2=n(n-1)/2$ pairs. Equality occurs for the complete graph.
\item \textbf{Structure.} What does a connected component represent in a
network model?\\
\emph{Solution.} It is a maximal group in which every vertex can reach every
other by a path. In a communication network it is an isolated communication
island; in a social network it is a cluster linked by chains of connections.
\item \textbf{Application.} State one limitation of using a graph for a real
network.\\
\emph{Solution.} A basic graph records which links exist but may ignore time,
capacity, direction, reliability, and changing behaviour. Those features
must be added when they affect the question.
\or
\item \textbf{Change.} For $f(t)=t^2$, find the instantaneous rate of change
at $t=3$.\\
\emph{Solution.} The derivative is $f'(t)=2t$, so $f'(3)=6$. This is the
limit of average rates over shorter and shorter time intervals.
\item \textbf{Equilibrium.} For the differential equation $y'=y$, which
initial value gives the solution $y(t)=5e^t$?\\
\emph{Solution.} At $t=0$, the value is $y(0)=5$. Differentiating gives
$y'=5e^t=y$, so the formula satisfies both the equation and the initial
condition.
\item \textbf{Stability.} Why is an equilibrium called stable when nearby
solutions remain nearby for a while?\\
\emph{Solution.} Small errors in the initial condition do not immediately
grow into large errors. This matters because measurements and computer
calculations never specify an initial state exactly.
\item \textbf{Iteration.} Starting with $x_0=1$ and using
$x_{n+1}=\tfrac12(x_n+2/x_n)$, compute $x_1$ and $x_2$.\\
\emph{Solution.} $x_1=3/2$ and
$x_2=\tfrac12(3/2+4/3)=17/12$. The iteration is moving toward $\sqrt2$;
the theory asks why and how quickly.
\item \textbf{Application.} Name one quantity that a dynamical model can
predict and one source of error it cannot remove.\\
\emph{Solution.} It may predict a population, orbit, temperature, or market
state. Measurement error, omitted forces, and uncertain parameters remain
unless the model is improved with new evidence.
\or
\item \textbf{Divisibility.} Use the Euclidean algorithm to find
$\gcd(252,105)$.\\
\emph{Solution.} $252=2(105)+42$, $105=2(42)+21$, and $42=2(21)$, so the
greatest common divisor is $21$.
\item \textbf{Congruence.} What is the remainder of $7^{100}$ upon division
by $10$?\\
\emph{Solution.} The last digits cycle $7,9,3,1$ with period four.
Since $100$ is divisible by four, the remainder is $1$.
\item \textbf{Prime structure.} Why is unique factorisation useful?\\
\emph{Solution.} It gives every positive integer a unique list of prime
building blocks. Questions about divisibility can then be answered by
comparing the exponents in those lists.
\item \textbf{Proof.} Explain why checking the first several integers does
not prove a statement about all integers.\\
\emph{Solution.} Infinitely many later cases remain. A proof must provide a
reason that covers every case, such as induction, contradiction, or a
general divisibility argument.
\item \textbf{Application.} Give one modern use of number theory and the
property it relies on.\\
\emph{Solution.} Public-key cryptography uses arithmetic modulo large
numbers, especially the difficulty of reversing certain operations without
secret information. Its security depends on computational assumptions, not
on absolute impossibility.
\or
\item \textbf{Dimension.} What is the dimension of a line, a plane, and
ordinary three-dimensional space?\\
\emph{Solution.} They have dimensions one, two, and three: the number of
independent directions needed to locate a nearby point.
\item \textbf{Coordinate change.} Why can the same geometric object have
different equations?\\
\emph{Solution.} Changing coordinates changes the labels assigned to points,
not the points themselves. A correct geometric statement must survive the
change of description.
\item \textbf{Intersection.} Two nonparallel lines in the plane meet in
how many points, and what changes for parallel lines?\\
\emph{Solution.} Nonparallel lines meet once; distinct parallel lines meet
zero times. The distinction is encoded by whether the two direction vectors
are multiples.
\item \textbf{Topology.} Why does continuously deforming a circle not turn
it into a line segment without tearing or gluing?\\
\emph{Solution.} A circle has a nontrivial loop that closes on itself, while
a segment has endpoints and no such loop. Continuity preserves this
qualitative feature under the allowed transformations.
\item \textbf{Application.} Explain why geometry can be useful even when a
physical object is not perfectly rigid.\\
\emph{Solution.} Geometry records relationships such as incidence,
distance, curvature, or connectivity. A model can retain the relationships
relevant to a question while ignoring microscopic irregularities.
\or
\item \textbf{Encoding.} A source has four equally likely symbols. How many
bits are needed for a fixed-length binary code?\\
\emph{Solution.} Two bits are enough because $2^2=4$. Each symbol can be
assigned one of $00,01,10,11$.
\item \textbf{Error detection.} Why does adding a parity bit detect any
single flipped bit?\\
\emph{Solution.} The original word has a prescribed even or odd parity.
Flipping one bit changes the parity, so the receiver notices a mismatch.
Two flips can cancel, so parity alone does not detect every error.
\item \textbf{Information.} What does it mean for an observation to reduce
uncertainty?\\
\emph{Solution.} Before the observation several outcomes are plausible; after
it fewer, or less evenly weighted, outcomes remain. Information measures this
change and is larger for a surprising observation.
\item \textbf{Compression.} Why can a file with predictable symbols be
compressed?\\
\emph{Solution.} Frequent symbols or patterns can receive shorter codewords,
while rare ones receive longer codewords. The average length falls because
the code uses the source's unequal frequencies.
\item \textbf{Application.} State the trade-off between compression and
error correction.\\
\emph{Solution.} Removing redundancy saves space but leaves less protection
against noise. Reliable communication deliberately keeps some redundancy so
the receiver can detect or repair likely errors.
\else
\item \textbf{Definition.} State the central question of #1 in your own
words and identify its basic objects.\\
\emph{Solution.} Translate the chapter's opening problem into a question
about the objects introduced there. A good answer names what is observed,
what operation or relation is allowed, and what conclusion is sought.
\item \textbf{Small case.} Work through the simplest nontrivial example
given in #1. What does it show that the definition alone does not?\\
\emph{Solution.} The example shows how the definition behaves in practice:
which choices are allowed, what is preserved, and where a tempting
generalisation fails.
\item \textbf{Calculation.} Choose one quantity introduced in #1 and
compute it for a small example.\\
\emph{Solution.} Write the definition of the quantity first, substitute the
small example, and check the units or type of the answer. This procedure
prevents a familiar symbol from replacing the actual meaning.
\item \textbf{Reasoning.} What assumption is most important in the main
result of #1, and what could happen if it were removed?\\
\emph{Solution.} Locate the condition in the theorem or method. Removing it
may make the conclusion false, weaken the guarantee, or require a different
theory. Testing the boundary is part of understanding the result.
\item \textbf{Application.} Describe one application of #1 and explain
which part of the mathematical model is an idealisation.\\
\emph{Solution.} Match the theory's objects to the application's objects and
state what prediction or decision follows. Then name what the model ignores;
that omission determines when the application must be recalibrated.
\fi
\end{enumerate}}

% Chapter-specific exercise frame.  Each chapter supplies its own five
% mathematical anchors; the prompts therefore cannot silently turn into a
% generic exercise set when a new theory is added.
\newcommand{\chapterquestions}[7]{%
\section*{7. Sample Questions}
\emph{Textbook exercise source:} \cite{#7}.  Use the exercise section of the
cited edition for the official problem statement and numbering.  The items
below are mathematical exercises tied directly to the chapter's definitions,
calculation, theorem, and limitation; they are not FAQ prompts.
\begin{enumerate}[leftmargin=*,itemsep=0.9em]
\item \textbf{Exercise 1.} Define the objects studied by #1 and give a
fully worked example of #2.
\item \textbf{Exercise 2.} Compute the example #3 using the construction
developed in #1, showing every algebraic or logical step.
\item \textbf{Exercise 3.} State the hypotheses and conclusion of the
central result: #5. Verify the hypotheses in a concrete example.
\item \textbf{Exercise 4.} Construct a counterexample showing why the
limitation stated in #6 matters, or prove the closest valid replacement.
\item \textbf{Exercise 5.} Apply #1 to the application described by #6 and
identify the mathematical quantity that is computed or predicted.
\end{enumerate}}

\newcommand{\theoryconnections}[3]{%
\section*{6. Connections}
\subsection*{Using existing theories}
\begin{enumerate}[leftmargin=*,itemsep=0.55em]
#2
\end{enumerate}
\subsection*{Helping other theories}
\begin{enumerate}[leftmargin=*,itemsep=0.55em]
#3
\end{enumerate}}

\begin{document}
\maketitle
\frontmatter
\chapter*{The big picture}
\addcontentsline{toc}{chapter}{The big picture}
Mathematics is a collection of ways to understand patterns. One theory may study patterns in numbers, another in shapes, another in change, and another in the way objects are connected. The theories have different names because they ask different kinds of questions, but they share one habit: simplify a situation, find the important pattern, and use that pattern to make a prediction or a decision.

Here is the map of this book:
\begin{description}[leftmargin=3.4cm,style=nextline]
\item[Numbers] Number theory studies whole numbers and primes. Algebra studies rules for combining objects and the symmetries of those rules.
\item[Shapes] Geometry studies measurements such as length and angle. Topology studies broader features such as connected pieces and holes.
\item[Change] Calculus and dynamical systems describe motion, growth, waves, stability, and sudden changes.
\item[Uncertainty] Probability describes chance. Statistics uses observations to learn what is likely to be true.
\item[Connections] Graph theory studies networks. Game, queue, and control theories study decisions, waiting, and feedback in those networks.
\item[Messages and machines] Information and coding theories explain data, noise, and reliable communication. Automata and computation theory explain what machines can recognize or calculate.
\item[Deep structure] The final parts study very general ideas about equations, spaces, symmetry, and proof. These chapters are previews of university mathematics, not subjects a student is expected to master in one sitting.
\end{description}

The first chapters are the easiest place to begin. Read the worked modules after the short introductions to probability, statistics, graph theory, information, group theory, topology, change, number theory, approximation, computation, decision-making, and algebraic geometry. The remaining chapters are windows into advanced subjects: their purpose is to show what questions exist and why mathematicians care about them.

You do not need to understand every symbol to understand the main story. At the end of each chapter, ask: What objects are being studied? What can change? What remains reliable? What real problem can this viewpoint help solve?

\chapter*{How to read this book}
Mathematics is often presented as a collection of answers. A theory is more like a reusable way of asking questions. It identifies a kind of pattern, gives us a precise language for it, and supplies tools for predicting what can happen.

The chapters are short on purpose. Each one begins with a problem a person might actually meet, then introduces the theory's main move, and ends with a few places where that move earns its keep. Equations appear as compact summaries, not as entrance exams. You can skip around: the theories are related, but no single route through them is required.

This book follows the names in the \href{https://en.wikipedia.org/wiki/List_of_mathematical_theories}{Wikipedia list of mathematical theories} (consulted 15 August 2026). The list is a map, not a claim that every theory is equally large, separate, or practical. Some are deep research areas; some are bridges between areas; and string theory is included because it is a mathematical-physical theory with a major mathematical life of its own. The entries below are short orientation guides. Twelve central subjects receive extended treatment in the worked modules near the end.

\chapter*{A promise about language}
Words such as ``space,'' ``function,'' ``group,'' and ``dimension'' have everyday meanings, but mathematicians use them more carefully. Whenever one appears, read it as a small tool with a job, not as a cloud of advanced ideas. A \emph{space} can simply mean a collection of possible positions. A \emph{function} is a rule that takes an input and returns an output. A \emph{model} is a deliberately simplified picture. An \emph{invariant} is a feature that does not change when the object is viewed or moved in an allowed way.

You do not need to memorize symbols to learn the theories. For every chapter, ask four ordinary-language questions: What is changing? What is being kept? What information is missing? What useful decision follows? If an explanation feels abstract, replace it with the chapter's picture---roads, ropes, messages, waves, choices, or landscapes---and return to the formal name afterward.

\chapter*{One problem, from life to mathematics}
Suppose a small clinic has one receptionist. Patients arrive at unpredictable times, and each registration takes a different amount of time. The practical question is not ``What is a queue?'' It is: how many people should the clinic expect to wait, and when should it add another receptionist?

First we choose a model. We record arrivals, service times, the number of workers, and the rule for deciding who goes next. We do not record every conversation, shoe colour, or brand of computer, because those details do not answer this question. Next we use probability to describe the uncertain arrivals and service times. Queue theory then combines those descriptions with the waiting rule. It can estimate the average wait and the chance that the line becomes dangerously long.

Finally, we compare choices. A second receptionist costs money but reduces waiting. The model does not make that decision for the clinic; it makes the trade-off visible. If the measured arrivals later disagree with the assumptions, we improve the model and check again. This is the scientific use of a theory: simplify honestly, calculate consistently, compare with evidence, and revise.

Every chapter follows this same pattern, even when the object is a number, a shape, a proof, or a physical law. The formal vocabulary is a compact label for a practical habit of thought.

\tableofcontents
\mainmatter

\part{Chapters in alphabetical order}
\clearpage\chapter{Algebraic geometry}
\section*{1. Overview}
\subsection*{Intuitive}
\paragraph{The problem}

An equation can be viewed in two ways. Algebra asks which numbers satisfy it; geometry asks what shape those solutions make. Algebraic geometry studies both views at once. Its basic objects are the sets of solutions of polynomial equations, but its deeper goal is to understand the whole solution space: its components, holes, singularities, points at infinity, maps, and behavior over different number systems.



\subsection*{Formal}
Let $k$ be a field and let $k[x_1,\ldots,x_n]$ be the polynomial ring in $n$ variables. An affine algebraic set is a common zero set $V(I)\subseteq k^n$ of an ideal $I$ in this ring. Polynomial maps that send one such set into another are the basic morphisms. The ideal of all polynomials vanishing on a set records its algebraic constraints; projective varieties and schemes extend this construction when coordinates, multiplicities, or arithmetic information must also be retained.
\section*{2. History}
\subsection*{Historical development}
\begin{historylist}
\paragraph{René Descartes and Pierre de Fermat}
They introduced coordinate methods that translated geometric curves into polynomial equations. This algebra--geometry dictionary is the starting point for studying solution sets rather than isolated numerical roots.

\paragraph{Bernhard Riemann, Richard Dedekind, and Max Noether}
They connected the study of algebraic curves with complex analysis, arithmetic, and the algebra of functions and ideals. Their work made genus, divisors, and algebraic transformations central tools for understanding curves.

\paragraph{Emmy Noether, André Weil, Oscar Zariski, Jean-Pierre Serre, and Alexander Grothendieck}
They developed the modern algebraic language of ideals, varieties, sheaves, and schemes. Grothendieck's schemes made it possible to treat number rings, finite fields, and geometric families in one framework \cite{algebraic_geometry_ref1,algebraic_geometry_ref4}.

\end{historylist}
\section*{3. Theory}
\subsection*{Core theory}
\paragraph{Zeros of simultaneous polynomials}
\bookfigure{algebraic_geometry_curve}{Algebraic geometry studies a solution set and the polynomial relations that define it.}

Let $k$ be a field and consider affine $n$-space $k^n$. Given polynomials $f_1,\ldots,f_r$ in $k[x_1,\ldots,x_n]$, their common zero set is
\[
 V(f_1,\ldots,f_r)=\{a\in k^n:f_i(a)=0\text{ for every }i\}.
\]
For example, $x^2+y^2-1=0$ is a circle over $\mathbb R$, while
\[
 x^2+y^2+z^2-1=0,\qquad x+y+z=0
\]
is a circle obtained by cutting a unit sphere with a plane. A line, parabola, ellipse, hyperbola, cubic curve, elliptic curve, lemniscate, and Cassini oval are all plane algebraic curves.

The first questions are concrete: how many solutions are there, and can they be found? The next questions are geometric: is the set connected or irreducible, where is it singular, what is its dimension, and what happens at infinity? The same equations may have no real solutions but many complex solutions, so the field matters.

\paragraph{Affine varieties and the algebra--geometry dictionary}

For a set $S$ of polynomials, define $I(S)$ to be all polynomials vanishing on $S$. This is an ideal: sums and polynomial multiples of vanishing polynomials still vanish. Conversely, for an ideal $I$, define $V(I)$ as its common zero set. The operations $I\mapsto V(I)$ and $S\mapsto I(S)$ form a reversing correspondence between algebraic conditions and geometric sets.

The Zariski topology declares sets of the form $V(I)$ to be closed. It is much coarser than the usual topology: a nonzero polynomial in one variable has only finitely many zeros, so finite sets are closed, but open sets are usually very large. This topology is designed to remember polynomial relations rather than distance.

The Hilbert Nullstellensatz says, over an algebraically closed field, that the ideal of all polynomials vanishing on $V(I)$ is the radical $\sqrt I$. It also identifies points of affine space with maximal ideals of the coordinate ring. Hilbert's basis theorem says every ideal in $k[x_1,\ldots,x_n]$ is finitely generated, so an algebraic set can be described by finitely many equations.

An algebraic set is irreducible if it cannot be expressed as the union of two smaller closed algebraic sets. Every algebraic set is a finite union of irreducible components, uniquely up to order. An affine algebraic set is irreducible exactly when its defining ideal can be chosen prime. This is the first major dictionary:
\[
\text{varieties}\leftrightarrow\text{prime ideals},\qquad
\text{functions}\leftrightarrow\text{coordinate rings}.
\]
\cite{algebraic_geometry_ref1,algebraic_geometry_ref2}

\paragraph{Regular functions and morphisms}

The regular functions on an affine variety $V$ are polynomial functions restricted to $V$. They form the coordinate ring
\[
 k[V]=k[x_1,\ldots,x_n]/I(V).
\]
For the parabola $y=x^2$, the relation $y-x^2=0$ gives
\[
 k[x,y]/(y-x^2)\cong k[x],
\]
because every function on the parabola can be written using the parameter $x$.

A regular map $f:V\to W$ is a map whose coordinate functions are regular. Composition of regular maps is regular, so affine varieties form a category. A map $f:V\to W$ induces a ring homomorphism in the opposite direction,
\[
 f^*:k[W]\to k[V],\qquad g\mapsto g\circ f.
\]
For affine algebraic sets, this correspondence between maps and ring homomorphisms is an equivalence with the opposite category of finitely generated reduced $k$-algebras. This contravariance is not a trick: a map of spaces lets us pull functions backward.

\paragraph{Rational functions and birational geometry}

If $V$ is irreducible, its coordinate ring is an integral domain and has a field of fractions $k(V)$, the field of rational functions. A rational function may have a denominator that vanishes on a hypersurface, so it is defined on a dense open part rather than everywhere.

Two varieties are birationally equivalent if rational maps in both directions are inverse wherever they are defined. Birational geometry studies properties that survive this kind of alteration. The unit circle is rational because
\[
 x=\frac{2t}{1+t^2},\qquad y=\frac{1-t^2}{1+t^2}
\]
parametrizes it away from one point. A rational variety is birational to affine space. Resolution of singularities asks whether a variety can be replaced by a birationally equivalent nonsingular one; Hironaka proved this in characteristic zero, while the general positive-characteristic problem remains deeper \cite{algebraic_geometry_ref1,algebraic_geometry_ref5}.

\paragraph{Projective varieties and points at infinity}

Affine space misses limiting directions. Projective space $\mathbb P^n(k)$ adds points at infinity and treats parallel lines as meeting there. A projective point is a nonzero vector $(x_0,\ldots,x_n)$ up to multiplication by a nonzero scalar. These are homogeneous coordinates.

A homogeneous polynomial has the same zero-or-nonzero behavior after scaling all coordinates, so homogeneous equations define projective algebraic sets. Projective varieties correspond to homogeneous prime ideals, and their homogeneous coordinate rings are graded rings. A projective variety generally has only constant global regular functions, but its function field of rational functions remains useful.

Projective geometry makes several theorems cleaner. Bézout's theorem counts intersections with multiplicity in a projective setting, including intersections at infinity. The projective closures of $y=x^2$ and $y=x^3$ have different behavior at infinity; the cubic acquires a cusp there. Projective properties also make existence arguments possible.

\paragraph{Real algebraic geometry}

Over $\mathbb R$, order cannot be ignored. The equation $x^2+y^2-a=0$ is a circle when $a>0$ and has no real points when $a<0$, although over $\mathbb C$ both cases have solutions. Real algebraic geometry studies real varieties and, more broadly, semialgebraic sets defined by polynomial equations and inequalities.

The first-quadrant branch of $xy-1=0$ is semialgebraic because it is defined by $xy-1=0$ together with $x>0$ and $y>0$. Algorithms such as cylindrical algebraic decomposition decide logical statements about polynomial inequalities and decompose real space into cells on which polynomial signs are constant. Questions about the possible arrangement of ovals of a nonsingular plane curve include parts of Hilbert's sixteenth problem.

\paragraph{Computational algebraic geometry}

Computational algebraic geometry turns the theory into algorithms. Gröbner bases are alternative generators for a polynomial ideal chosen with respect to a monomial order. They generalize Gaussian elimination: leading monomials make reductions systematic.

If a Gröbner basis reduces to $\{1\}$, the corresponding variety has no points over an algebraically closed extension. Hilbert series can reveal dimension and degree. When the dimension is zero, the basis can be used to compute the finitely many solutions. Elimination orders remove variables and solve projection problems, while resultants eliminate variables from two equations.

Cylindrical algebraic decomposition computes the topology and sign structure of semialgebraic sets. Numerical algebraic geometry uses homotopy continuation: start with a system whose solutions are known and continuously deform it into the target system, tracking solution paths with floating-point methods. Complexity matters: symbolic algorithms may be exact but expensive, while numerical methods are fast but need conditioning and error checks \cite{algebraic_geometry_ref3,algebraic_geometry_ref6}.

\paragraph{Schemes and the modern viewpoint}

Classical geometry studies points with coordinates in a field. Scheme theory replaces an affine variety by the spectrum of a ring, the set of all prime ideals equipped with a structure sheaf. Maximal ideals give ordinary points, while nonmaximal prime ideals represent irreducible subvarieties and generic information.

This enlargement allows one theory to treat complex varieties, varieties over finite fields, and $\operatorname{Spec}(\mathbb Z)$, whose prime ideals encode the arithmetic of rational primes. Sheaves record which functions are available on each open set and how local data glue. Schemes can be glued from affine pieces, just as manifolds are glued from coordinate charts. The same language supports arithmetic geometry, deformation theory, intersection theory, and modern cohomology.

Grothendieck's approach also clarifies why the Nullstellensatz is only one case of a wider algebra--geometry relation. Algebraic geometry now includes stacks, derived geometry, noncommutative geometry, and analytic and rigid-analytic variants, but the basic lesson remains: study an object through its functions and the maps between its local pieces.

\paragraph{Analytic geometry and GAGA}

An analytic variety is locally defined by convergent analytic functions rather than polynomials. Complex manifolds are nonsingular analytic varieties, but analytic varieties may have singularities. Serre's GAGA results compare complex algebraic varieties with their associated analytic spaces: many algebraic and analytic notions, such as coherent sheaves and cohomology, agree in the proper setting.

Over non-archimedean fields, rigid analytic spaces provide an analytic geometry adapted to $p$-adic absolute values. These theories connect algebraic geometry with complex analysis, number theory, and $p$-adic geometry \cite{algebraic_geometry_ref7}.

\section*{4. Important theorems and results}
\begin{enumerate}[leftmargin=*,itemsep=0.35em]

\item \textbf{Hilbert's Nullstellensatz.} Over an algebraically closed field, the radical of an ideal is the ideal of all polynomials vanishing on its zero set, and maximal ideals correspond to points.

\item \textbf{Hilbert basis theorem.} Every ideal in a polynomial ring in finitely many variables over a field is finitely generated. Algebraic sets therefore have finite equation descriptions.

\item \textbf{Irreducible decomposition.} Every algebraic set is a finite union of irreducible components, uniquely up to order. Prime ideals describe the irreducible pieces.

\item \textbf{Bézout's theorem.} Projective plane curves with no common component meet in a number of points counted with multiplicity equal to the product of their degrees.

\item \textbf{Riemann--Roch theorem.} On a smooth projective curve, the dimension of the space of functions with prescribed zeros and poles is controlled by the degree of the divisor and the genus.

\item \textbf{GAGA.} Proper complex algebraic varieties and their analytic counterparts have matching categories of coherent sheaves and matching cohomological information in the appropriate sense.
\end{enumerate}
\noindent\emph{Source for these results:} \cite{algebraic_geometry_ref1}.

\section*{5. Applications}
Algebraic geometry is useful when polynomial equations describe the objects or constraints in a problem. Its varieties, dimensions, and singularities turn those equations into structures that can be analysed in arithmetic, computation, and geometry.
\begin{enumerate}[leftmargin=*,itemsep=0.35em]
\item \textbf{Cryptography and coding.} Elliptic curves supply groups used in public-key cryptography. Algebraic curves over finite fields produce error-correcting codes whose distance and decoding behavior can be studied geometrically.
\item \textbf{Robotics and control.} The positions of a robot arm satisfy polynomial equations. Algebraic geometry can describe the reachable set, detect singular configurations, and solve inverse-kinematics equations. In control theory, polynomial invariants describe states that a system can or cannot reach.
\item \textbf{Statistics and phylogenetics.} Probability models for discrete data often form algebraic varieties or semialgebraic sets. Polynomial relations among observed frequencies can test whether a model is plausible. Similar varieties describe evolutionary trees.
\item \textbf{Geometric modelling and physics.} Computer-aided design uses polynomial and rational surfaces. Algebraic geometry also connects to string theory, solitons, game theory, graph matchings, integer programming, and moduli spaces in mathematical physics.
\item \textbf{Number theory.} Diophantine equations become questions about rational points on varieties. The modularity ideas used in the proof of Fermat's Last Theorem illustrate how curves, modular forms, and arithmetic can be joined in one geometric framework.
\end{enumerate}


\theoryconnections{algebraic_geometry}{%
\item Commutative algebra turns polynomial equations into ideals, and the ideal records all polynomial consequences of the equations.
\item The spectrum of that ring supplies the points and local neighborhoods of the corresponding geometric object; localization then lets one study a point without losing the global equation.
\item Sheaves and topology explain how local solutions are glued, while number theory and complex analysis enter when the coefficients or points are arithmetic or complex.
\item For example, solving $x^2+y^2=1$ over different fields gives genuinely different geometries, so the coefficient field is part of the problem, not decoration \cite{algebraic_geometry_ref1,algebraic_geometry_ref2}.
}{%
\item Algebraic geometry helps arithmetic geometry by treating integer solutions as points on a variety and then attaching local information at each prime.
\item It helps coding theory because a code can be built from evaluations of functions on an algebraic curve; the curve controls how many errors can be detected.
\item In robotics and control, polynomial constraints for joints or feedback can be eliminated and studied as geometric sets.
\item Thus geometry converts a large system of equations into objects whose dimension, singularities, and intersections can be analyzed \cite{algebraic_geometry_ref3,algebraic_geometry_ref6}.
}
\section*{7. Sample Questions}
\emph{Exercise reference:} \cite{algebraic_geometry_ref1}. The cited edition is the source for the exercise set; consult its Exercises section for the official statement and numbering.
\begin{enumerate}[leftmargin=*,itemsep=0.9em]

\item \textbf{Exercise 1.} Find the coordinate ring of the parabola $y=x^2$.

\textbf{Solution.} It is $\mathbb R[x,y]/(y-x^2)$. Substituting $y=x^2$ gives an isomorphism with $\mathbb R[x]$, so $x$ is a global parameter.

\item \textbf{Exercise 2.} Determine whether $x^2+y^2=0$ has a nonzero real point.

\textbf{Solution.} No. Both squares are nonnegative, so their sum is zero only when $x=y=0$. Over $\mathbb C$, nonzero solutions exist, illustrating dependence on the base field.

\item \textbf{Exercise 3.} What does a Gröbner basis tell us if it contains 1?

\textbf{Solution.} The ideal is the whole polynomial ring. Therefore the equations have no common solution over an algebraically closed extension of the coefficient field.

\item \textbf{Exercise 4.} Why add projective points?

\textbf{Solution.} Affine curves can escape to infinity, and intersections can disappear in a coordinate chart. Projective completion adds limiting directions so that intersection theorems such as Bézout's count all intersections, including those at infinity.

\item \textbf{Exercise 5.} Why are rational functions not defined everywhere?

\textbf{Solution.} A quotient $p/q$ is undefined where $q=0$. On an irreducible variety it is therefore defined on a dense open subset, and the collection of such quotients forms the function field.

\end{enumerate}
\section*{8. References}
\begin{thebibliography}{99}
\bibitem{algebraic_geometry_ref1} R. Hartshorne, \emph{Algebraic Geometry}.
\bibitem{algebraic_geometry_ref2} D. Eisenbud, \emph{Commutative Algebra with a View Toward Algebraic Geometry}.
\bibitem{algebraic_geometry_ref3} D. Cox, J. Little, and D. O'Shea, \emph{Ideals, Varieties, and Algorithms}.
\bibitem{algebraic_geometry_ref4} I. R. Shafarevich, \emph{Basic Algebraic Geometry}.
\bibitem{algebraic_geometry_ref5} R. Hartshorne, \emph{Algebraic Geometry}, chapters on varieties, curves, and divisors.
\bibitem{algebraic_geometry_ref6} B. Sturmfels, \emph{Solving Systems of Polynomial Equations}.
\bibitem{algebraic_geometry_ref7} J.-P. Serre, "Géométrie algébrique et géométrie analytique," \emph{Annales de l'Institut Fourier}.
\end{thebibliography}


\clearpage\chapter{Almgren--Pitts min--max theory}
\section*{1. Overview}
\subsection*{Intuitive}
\paragraph{The problem}

How can geometry produce a minimal surface when direct minimization is too simple? If a soap film is allowed to move freely, it may shrink to a point or change its topology. The Almgren--Pitts theory solves this by minimizing not the area of one surface, but the largest area reached by an entire family of surfaces that sweeps through a nontrivial part of the space. The resulting critical surface is a higher-dimensional analogue of the saddle point found by Morse theory.



\subsection*{Formal}
Let $M$ be a closed Riemannian manifold and let $\mathcal Z_n(M;\mathbb Z_2)$ denote the space of mod--$2$ flat $n$-cycles. For a family $\Phi:X\to\mathcal Z_n(M;\mathbb Z_2)$, its width is $\sup_{x\in X}\mathbf M(\Phi(x))$, where $\mathbf M$ is mass. A homotopy class $\Pi$ of nontrivial sweepouts has min--max value $\mathbf L(\Pi)=\inf_{\Phi\in\Pi}\sup_{x\in X}\mathbf M(\Phi(x))$. Under the regularity hypotheses of the min--max theorem, a minimizing sequence produces a stationary integral varifold with mass $\mathbf L(\Pi)$.
\section*{2. History}
\subsection*{Historical development}
\begin{historylist}
\paragraph{What the theory is and why it was developed}

The theory is named after Frederick J. Almgren Jr. and his student Jon T. Pitts. It is a global variational method for constructing critical points of an area or volume functional. Its starting point was George Birkhoff's min--max construction of a simple closed geodesic on a sphere. A geodesic is a shortest path only after its endpoints and homotopy class prevent collapse. Birkhoff's idea suggested a higher-dimensional question: can a family of curves or surfaces be made to sweep through a manifold, and can the highest length or area in that sweep be minimized?

Almgren developed the topology of spaces of cycles, and Pitts built the discretized and regularity machinery needed to turn a minimizing sequence of families into a genuine minimal hypersurface. The mathematicians were trying to find geometric objects that direct minimization misses: saddle points of area rather than absolute minima. Later work by Allard, Schoen, Simon, Marques, Neves, Zhou, and others extended regularity and applications to geometry and topology \cite{almgren_pitts_min_max_theory_ref1,almgren_pitts_min_max_theory_ref2,almgren_pitts_min_max_theory_ref3}.

\end{historylist}
\section*{3. Theory}
\subsection*{Core theory}
\paragraph{The picture before the technical definitions}

Imagine a mountain pass. Many paths connect one valley to another. Each path rises to a highest elevation. Among all connecting paths, choose one whose highest point is as low as possible. The mountain-pass height is a critical level: at that height the landscape cannot be lowered without breaking the connection between the valleys.

Replace paths by hypersurfaces and height by area. A \emph{sweepout} is a continuously varying family of surfaces, indexed by a parameter space. The family may begin as a small surface, expand through the manifold, and finish as another small surface. Its width is the largest area of any slice. A min--max sequence consists of better and better sweepouts whose widths approach the infimum
\[
 \mathbf L(\Pi)=\inf_{\Phi\in\Pi}\ \sup_{x\in X}\mathbf M(\Phi(x)),
\]
where $\Pi$ is a nontrivial homotopy class of sweepouts, $X$ is the parameter space, and $\mathbf M$ denotes mass, the generalized area. The central theorem says that a min--max sequence has a subsequence converging, in the varifold sense, to a stationary object.

\paragraph{Cycles, mass, and varifolds}

A $k$-cycle is a generalized $k$-dimensional surface with no boundary. Smooth closed curves and closed surfaces are examples. Almgren and Pitts use integral currents and the flat metric to measure how cycles differ, while mass measures their area or volume. A family must vary continuously in the relevant topology; without this care, nearby parameter values could represent unrelated surfaces and the sweep would carry no geometric meaning.

A varifold forgets orientation and records how much $k$-dimensional material passes through each region and direction. It is flexible enough to describe a limiting sequence even when the surfaces develop folds, cancellation, or singularities. A varifold is \emph{stationary} when its first variation vanishes for every compactly supported deformation. For a smooth surface, this is exactly the condition of zero mean curvature, hence minimality.

The limit need not initially be a smooth embedded surface. It is a stationary integral varifold, a measure-theoretic generalization of a minimal submanifold. Regularity theory then asks how large its singular set can be. This separation between existence and regularity is one of the theory's major conceptual achievements \cite{almgren_pitts_min_max_theory_ref1,almgren_pitts_min_max_theory_ref4}.

\paragraph{Almgren's topology of cycle spaces}

The space of cycles has its own homotopy groups. Almgren proved that for a closed Riemannian manifold $M$, the $m$-th homotopy group of the space of flat $k$-cycles is isomorphic to the $(m+k)$-th homology group of $M$:
\[
 \pi_m(\mathcal Z_k(M))\cong H_{m+k}(M).
\]
The theorem generalizes the Dold--Thom theorem, which corresponds morally to the case $k=0$. In simple language, it says that families of cycles remember the topology of the manifold itself. A nontrivial homology class can therefore generate a nontrivial family of surfaces, and a nontrivial family is the raw material for min--max.

This is why the method is global. It does not begin with a local graph and try to solve a differential equation. It first identifies a topological obstruction to shrinking the whole family away. The area functional must then attain a saddle-type critical level inside that obstruction.

\paragraph{The min--max construction}

Choose a nontrivial parameterized family of cycles. Approximate the family by finer and finer discrete families while controlling the mass changes between neighboring parameter values. Improve the family by a tightening or deformation procedure: if a slice can be locally pushed down in area without damaging the homotopy class, make that improvement. A minimizing sequence eventually has no remaining cheap deformation.

The limit is stationary because any deformation with negative first variation would lower nearby slices and contradict minimality of the width. The almost-minimizing property prevents the limit from being an arbitrary stationary object with severe singularities. In codimension one, Pitts used this property to prove regularity in low dimensions.

The construction is subtle because the parameterized surfaces can concentrate, cancel, or converge with multiplicity. The discrete fineness conditions, interpolation theorems, and replacement arguments are what preserve the topology while controlling the geometric limit. The theory's strength comes from combining homotopy, measure, and variational reasoning rather than relying on any one of them alone \cite{almgren_pitts_min_max_theory_ref1,almgren_pitts_min_max_theory_ref5}.

\paragraph{Regularity and dimensional thresholds}

Almgren initially obtained stationary integral varifolds. Allard's regularity theorem shows that a stationary varifold is smooth on an open dense region under appropriate density and stability hypotheses. Pitts greatly improved the codimension-one construction and introduced $1/j$-almost-minimizing varifolds. Schoen and Simon extended the regularity theory.

For an $n$-dimensional Riemannian manifold, the min--max hypersurface is smooth when $3\leq n\leq6$. In higher dimensions it is smooth away from a closed singular set of dimension at most $n-8$. This threshold reflects the existence of singular minimal cones in higher dimensions; it is not a defect of the proof.

\section*{4. Important theorems and results}
\begin{enumerate}[leftmargin=*,itemsep=0.35em]

\item \textbf{Almgren isomorphism theorem.} The homotopy groups of spaces of flat cycles correspond to shifted homology groups of the manifold. This converts topological information about the manifold into nontrivial families of cycles.

\item \textbf{Min--max existence theorem.} A nontrivial homotopy class of sweepouts has a minimizing sequence whose critical limit is a stationary integral varifold.

\item \textbf{Allard regularity theorem.} A stationary varifold that is sufficiently close to a plane in measure and density is represented locally by a smooth graph. It explains why a weak limit is smooth on a large regular set.

\item \textbf{Pitts regularity theorem.} In codimension one, the almost-minimizing min--max varifold is smooth in dimensions up to six and has singular set of dimension at most $n-8$ in higher dimensions.

\item \textbf{Marques--Neves theorem.} The min--max method proves the Willmore conjecture for embedded tori in the three-sphere.
\end{enumerate}
\noindent\emph{Source for these results:} \cite{almgren_pitts_min_max_theory_ref1}.

\section*{5. Applications}
Almgren--Pitts min--max theory is useful when topology forces a geometric object to exist but direct minimization collapses to a trivial object. A sweepout supplies the constraint, and the least possible maximal area identifies a minimal hypersurface.
\begin{enumerate}[leftmargin=*,itemsep=0.35em]
\item \textbf{Round sphere.} Sweep the unit sphere by latitude circles or, in a higher-dimensional version, by parallel hyperspheres. The equatorial slice is the natural min--max object. A family that shrinks immediately to a point does not represent a nontrivial class; a genuine sweep must force the family across the manifold.
\item \textbf{Three-manifolds.} In a closed three-manifold, a one-parameter family of surfaces can produce an embedded minimal surface. The method is useful when no obvious stable minimal surface exists. Positive Ricci curvature, Heegaard splittings, and the topology of the manifold influence the number and index of the surfaces obtained.
\item \textbf{Willmore conjecture.} Marques and Neves used a refined min--max construction to prove that every embedded torus in the three-sphere has Willmore energy at least $2\pi^2$. The proof found a critical surface in a family of conformal images and then compared its energy with the Clifford torus.
\item \textbf{Multiplicity and many surfaces.} Higher-parameter families can produce distinct minimal hypersurfaces. The resulting width spectrum records a sequence of critical values, analogous to the sequence of eigenvalues in spectral theory. This connects min--max geometry with topology, index estimates, and questions about how complicated a manifold must be to contain many minimal surfaces \cite{almgren_pitts_min_max_theory_ref2,almgren_pitts_min_max_theory_ref6}.
\item \textbf{Soap films and plates in materials science.} Soap films spanning wire frames are physical realizations of area-minimizing surfaces, and the Almgren--Pitts framework explains why certain films are stable while others undergo topological transitions. In metallurgy and thin-film deposition, grain boundaries and minimal surfaces within foams follow the same variational principle: the material minimizes interfacial energy subject to topological constraints imposed by the surrounding crystal structure. Plateau's laws, which govern how soap-film facets meet at edges, can be derived as regularity consequences of the min--max critical theory. Understanding these junction rules helps materials scientists predict the stability of open-cell foams used in lightweight structural composites and in the design of metallic microlattices for energy absorption.
\item \textbf{Minimal surface computation in computer graphics.} Polygonal meshes in computer-aided design are often relaxed to minimal or constant-mean-curvature surfaces to smooth wrinkles while preserving topology. The Almgren--Pitts sweepout construction provides the existence guarantee: given a mesh with the homotopy type of a torus, the min--max method guarantees the existence of a minimal surface spanning a nontrivial family of sweepouts. In practice, surface-modeling algorithms iterate an area-reduction step analogous to the tightening procedure of the theory, producing mesh fairing techniques used in animation studios and industrial CAD systems. By tracking which sweepout class each converged mesh belongs to, engineers ensure that the resulting surface faithfully represents the intended topology rather than collapsing to a degenerate flat patch \cite{almgren_pitts_min_max_theory_ref1,almgren_pitts_min_max_theory_ref6}.
\end{enumerate}


\theoryconnections{almgren_pitts_min_max_theory}{%
\item Riemannian geometry supplies the notion of area, while the calculus of variations asks for surfaces that make area stationary.
\item Algebraic topology provides a nontrivial family of surfaces to sweep through, and currents and varifolds allow a sequence of rough surfaces to converge even when ordinary smooth convergence fails.
\item Regularity theory is then needed to prove that the limiting object is smooth except at controlled singularities.
\item In this way, topology creates the competitor family, analysis takes the maximum, and geometric measure theory makes the limiting process rigorous \cite{almgren_pitts_min_max_theory_ref1,almgren_pitts_min_max_theory_ref3}.
}{%
\item The theory produces minimal hypersurfaces by taking the least possible maximal area in a sweepout.
\item Those hypersurfaces become geometric probes: their areas give widths of a manifold, their indices measure unstable directions, and their regularity makes them usable in geometric analysis.
\item This connects directly to three-manifold topology, the Willmore problem, scalar-curvature questions, and spectral geometry, where the existence of a canonical surface or width supplies information that pointwise estimates alone do not reveal \cite{almgren_pitts_min_max_theory_ref2,almgren_pitts_min_max_theory_ref6}.
}

\section*{Glossary}
\begin{description}[leftmargin=*,style=nextline,itemsep=0.3em]

\item[\textbf{Sweepout.}] A continuously varying one-parameter family of surfaces that moves through a manifold, beginning and ending as small surfaces; it represents a nontrivial way to fill the manifold with intermediate slices.

\item[\textbf{Varifold.}] A measure-theoretic generalization of a surface that records how much material passes through each region and direction without requiring a smooth parametrization or orientation.

\item[\textbf{Stationary varifold.}] A varifold whose first variation of mass vanishes for every compactly supported deformation; for a smooth surface this is equivalent to having zero mean curvature.

\item[\textbf{Mass.}] The generalized area or volume of a cycle or varifold, serving as the quantity to be minimized in the min--max construction.

\item[\textbf{Flat cycle.}] A generalized surface with no boundary, equipped with a flat metric that measures how two cycles differ; cycles are the basic objects varied in a sweepout.

\item[\textbf{Width.}] The largest mass occurring in a sweepout; the min--max value is the infimum of widths over all sweepouts in a given homotopy class.

\item[\textbf{Almost-minimizing varifold.}] A varifold that satisfies a local near-minimality condition, preventing the limit from developing severe singularities and enabling regularity conclusions.

\item[\textbf{Integral current.}] A generalized oriented surface with locally finite mass that can be represented as an integral of smooth submanifolds with multiplicities.

\item[\textbf{Mean curvature.}] A vector measuring how a hypersurface bends in the ambient space; a hypersurface is minimal when its mean curvature vanishes.

\item[\textbf{First variation.}] The rate of change of mass under a small deformation; a varifold is stationary when its first variation vanishes.

\item[\textbf{Homology class.}] An element of a homology group representing an equivalence class of cycles; a nontrivial class generates a nontrivial sweepout.

\end{description}

\section*{7. Sample Questions}
\emph{Exercise reference:} \cite{almgren_pitts_min_max_theory_ref1}. The cited edition is the source for the exercise set; consult its Exercises section for the official statement and numbering.
\begin{enumerate}[leftmargin=*,itemsep=0.9em]

\item \textbf{Exercise 1.} Consider the round $S^2$ with its standard metric. A sweepout by latitude circles parameterized by height $h\in[-1,1]$ has width equal to the length of the longest latitude circle. Compute this width.

\emph{Solution.} At height $h$, the latitude circle has radius $\sqrt{1-h^2}$ and circumference $2\pi\sqrt{1-h^2}$. The maximum circumference occurs at $h=0$, the equator, giving width $2\pi$. This is the min--max value for the simplest sweepout class.

\item \textbf{Exercise 2.} On $S^2$, a different sweepout collapses the sphere by shrinking all latitude circles from the equator toward the north pole. The circles at height $h$ have circumference $2\pi\sqrt{1-h^2}$ for $h\in[0,1]$. What is the width of this sweepout?

\emph{Solution.} The maximum circumference in this family occurs at $h=0$ (the equator), which has circumference $2\pi$. The width is again $2\pi$. However, this sweepout does not represent a nontrivial homotopy class in the relevant sense because the family can be deformed to avoid large circles.

\item \textbf{Exercise 3.} Let $\Phi:[0,1]\to\mathcal{Z}_1(S^1;\mathbb{Z}_2)$ be a sweepout of the circle $S^1$ by points (0-cycles), where $\Phi(0)$ and $\Phi(1)$ are the empty cycle. Suppose the mass (length) of $\Phi(t)$ equals $\pi|\sin(\pi t)|$. Compute the width of this sweepout.

\emph{Solution.} The width is $\sup_{t\in[0,1]}\pi|\sin(\pi t)|=\pi$, achieved at $t=1/2$. The min--max value over all sweepouts in this class would be at most $\pi$.

\item \textbf{Exercise 4.} For a smooth closed hypersurface $\Sigma^n\subset M^{n+1}$ with $3\le n\le 6$, the Pitts regularity theorem guarantees $\Sigma$ is smooth. If instead $n=8$, what is the maximum possible dimension of the singular set?

\emph{Solution.} By the Pitts regularity theorem, the singular set has dimension at most $n-8=8-8=0$. So the singular set is at most a finite collection of points. For $n=9$ it would be at most 1-dimensional, and so on.

\item \textbf{Exercise 5.} By Almgren's isomorphism theorem, $\pi_0(\mathcal{Z}_1(S^2))\cong H_1(S^2)=0$, but $\pi_1(\mathcal{Z}_1(S^2))\cong H_2(S^2)=\mathbb{Z}$. What does $\pi_1(\mathcal{Z}_1(S^2))=\mathbb{Z}$ tell us about sweeping $S^2$ by closed curves?

\emph{Solution.} It says that families of 1-cycles on $S^2$, parameterized by a circle, are classified by an integer. A nonzero integer corresponds to a nontrivial sweepout, guaranteeing that the min--max construction produces a minimal 1-varifold (a closed geodesic) on $S^2$.

\end{enumerate}
\section*{8. References}
\begin{thebibliography}{99}
\bibitem{almgren_pitts_min_max_theory_ref1} J. T. Pitts, \emph{Existence and Regularity of Minimal Surfaces on Riemannian Manifolds}, Princeton University Press.
\bibitem{almgren_pitts_min_max_theory_ref2} F. C. Marques and A. Neves, \emph{Min--Max Theory and the Willmore Conjecture}.
\bibitem{almgren_pitts_min_max_theory_ref3} F. J. Almgren Jr., \emph{The Theory of Varifolds: A Variational Calculus in the Large for the $K$-Dimensional Area Integrand}.
\bibitem{almgren_pitts_min_max_theory_ref4} L. Simon, \emph{Lectures on Geometric Measure Theory}, Australian National University.
\bibitem{almgren_pitts_min_max_theory_ref5} T. Colding and C. De Lellis, "The min--max construction of minimal surfaces," \emph{Surveys in Differential Geometry}.
\bibitem{almgren_pitts_min_max_theory_ref6} F. C. Marques and A. Neves, "Applications of Almgren--Pitts min--max theory," lecture notes.
\end{thebibliography}


\clearpage\chapter{Approximation theory}
\section*{1. Overview}
\subsection*{Intuitive}
\paragraph{The problem}
Many mathematical rules are easy to write but hard to calculate. The exact curve of a bridge under load, the path of a spacecraft, or the brightness of a pixel may be described by an equation that has no convenient answer. A computer must replace the exact object with a simpler one. The danger is obvious: a fast answer is not useful if nobody knows how far it may be wrong.
Approximation theory asks how to make a simpler replacement and measure its error. It is not the same as guessing. An approximation is a controlled substitute: we know what was simplified, why the substitute should be close, and where the guarantee stops applying.


\subsection*{Formal}
Let $f$ belong to a normed function space and let $A_N$ be a finite-dimensional family of approximants, such as polynomials of degree at most $N$. Approximation chooses $p_N\in A_N$ to make an error norm small, for example $\|f-p_N\|_\infty$ for worst-case error or $\|f-p_N\|_2$ for mean-square error. The approximation problem is therefore an optimization problem together with an error estimate; the estimate is valid only on the stated domain and under the stated regularity assumptions.
\section*{2. History}
\subsection*{Historical development}
\begin{historylist}
\paragraph{The people and the historical problem}
The subject grew from several needs rather than one invention. Isaac Newton and Joseph-Louis Lagrange developed polynomial methods for replacing complicated functions near points. Joseph Fourier showed that repeating signals could be built from simple waves, leading to a powerful way to approximate sound, heat, and images.

In the nineteenth century, Pafnuty Chebyshev studied the best possible approximations and asked how to keep the largest error as small as possible. Karl Weierstrass proved a remarkable existence result: every continuous function on a closed interval can be approximated as closely as desired by polynomials. Sergei Bernstein later gave an explicit construction that made the theorem more concrete.

These mathematicians were not merely looking for shorter formulas. They were trying to understand whether complicated continuous behavior could be captured by finite calculations, and whether ``close'' could be made precise.

\end{historylist}
\section*{3. Theory}
\subsection*{Core theory}
\paragraph{The central idea}
Suppose we want to approximate $\sqrt{2}$. The decimal $1.4142$ is useful only after we say what error is acceptable. Since $(1.4142)^2=1.99996164$, its square is less than 2 by about $0.00003836$. The approximation has a measurable error, not just a pleasant appearance.

For a function, we choose a simpler family: straight lines, polynomials, splines, or trigonometric waves. We then choose a rule for judging closeness. Sometimes we care about the largest error anywhere; sometimes we care about average error; sometimes we care more about errors in a particular region. Different choices produce different ``best'' approximations.

\paragraph{Why the method works}
The method works because it turns an impossible demand---calculate every detail exactly---into a controlled question: how much detail is enough for this purpose? A theorem can bound the error, and computation can refine the approximation until the bound is acceptable. The key is to keep the error visible rather than hide it behind many decimal places.

For example, near zero the curve $\sin x$ is close to the straight line $x$. Adding the correction $-x^3/6$ gives a better estimate. At $x=0.1$, the first approximation is already close; at $x=2$, the same shortcut is poor. An approximation always comes with a region and a level of accuracy.

\paragraph{Optimal polynomials and the equioscillation idea}

Suppose the degree $N$ and interval $[a,b]$ are fixed. A best uniform polynomial is the polynomial $P$ that minimizes
\[
 \max_{a\leq x\leq b}|P(x)-f(x)|.
\]
The equioscillation theorem gives a practical certificate for this choice. Under the usual continuity assumptions, the error $P-f$ reaches the same magnitude with alternating signs at at least $N+2$ ordered points. For a fourth-degree polynomial, the ideal error therefore has six alternating extrema. Two of these points may be the endpoints.

This pattern explains why merely matching $N+1$ data points is not generally optimal. Interpolation forces agreement at selected points, but it may leave a large error between them. A minimax polynomial distributes its unavoidable error more evenly. If a competing degree-$N$ polynomial were strictly closer at all $N+2$ alternating extrema, the difference of the two polynomials would change sign at least $N+1$ times. The intermediate value theorem would then give too many zeros for a nonzero degree-$N$ polynomial. Thus no competitor can improve the worst error \cite{approximation_theory_ref1,approximation_theory_ref2}.

For a concrete library problem, a degree-four polynomial for $\exp x$ on $[-1,1]$ can be chosen so that its error repeatedly touches $+\varepsilon$ and $-\varepsilon$. The same idea applies to $\log x$ on $[2,4]$. The numerical value of $\varepsilon$ depends on the function and interval; the important design fact is that the largest errors are balanced rather than concentrated at one end.

\paragraph{Chebyshev expansion}

The Chebyshev polynomials are defined by
\[
 T_n(\cos\theta)=\cos(n\theta).
\]
They satisfy $T_0(x)=1$, $T_1(x)=x$, and
\[
 T_{n+1}(x)=2xT_n(x)-T_{n-1}(x).
\]
On $[-1,1]$ they oscillate between $-1$ and $1$. In particular, $T_{N+1}$ has $N+2$ alternating extrema, exactly the geometry suggested by the equioscillation theorem.

For a sufficiently smooth function, write a generalized Fourier expansion
\[
 f(x)\sim\sum_{k=0}^{\infty}c_kT_k(x).
\]
Stopping after $T_N$ gives a degree-$N$ polynomial. If the coefficients decrease rapidly, the first omitted term is larger than the sum of the later terms, so the truncated expansion is close to the minimax polynomial. It can be slightly better in some parts of the interval and slightly worse in others; near-optimal is not identical to optimal. The method is particularly effective for $\exp x$ and other analytic functions, while singularities or endpoint difficulty slow coefficient decay \cite{approximation_theory_ref4,approximation_theory_ref5}.

Chebyshev approximation also supplies Clenshaw--Curtis quadrature. The integral is approximated by sampling the function at Chebyshev-related points, expanding the sampled values, and integrating the polynomial expansion. Thus a representation chosen to control approximation error also gives an efficient numerical integration rule.

\paragraph{Remez's algorithm: finding the minimax polynomial}

The Remez algorithm turns the equioscillation condition into an iteration. For a degree-$N$ polynomial choose $N+2$ ordered test points $x_0,\ldots,x_{N+1}$. Require
\[
 P(x_i)-f(x_i)=(-1)^i\varepsilon.
\]
Since $P(x)=P_0+P_1x+\cdots+P_Nx^N$, these are $N+2$ linear equations in the $N+1$ coefficients and the unknown error size $\varepsilon$. Solve them to obtain a polynomial whose error alternates at the selected points.

The selected interior points may not yet be true maxima or minima of the error. Compute the derivative of $P-f$, move each interior point toward a nearby zero of that derivative using Newton's method, and solve the alternating linear system again. Continue until the error has the same alternating magnitude at actual extrema. Good initial points are the extrema of $T_{N+1}$ because the final error has a similar shape.

The method needs evaluations of $f$, $f'$, and often $f''$ to high precision. The computation should use more precision than the final answer requests, since errors in the test points and coefficients can hide the final alternation. For well-behaved functions, convergence is quadratic near the solution: an error of roughly $10^{-15}$ can become roughly $10^{-30}$ after a further iteration. Remez is therefore both a numerical algorithm and a direct demonstration of the optimal-polynomial theorem \cite{approximation_theory_ref1,approximation_theory_ref6,approximation_theory_ref7}.

\section*{4. Important theorems and results}
\begin{enumerate}[leftmargin=*,itemsep=0.35em]
\item \textbf{Central result.} Weierstrass approximation says continuous functions on a closed interval can be uniformly approximated by polynomials.

\item \textbf{Taylor error bound.} A Taylor polynomial has a remainder controlled by a bound on the next derivative. This explains when a local approximation is trustworthy and why an error estimate must accompany numerical simulation or signal processing.

Weierstrass gives existence but not an efficient construction or a useful finite-degree error bound. Taylor's theorem gives a local construction with a remainder estimate, while minimax and Chebyshev methods choose an approximation that controls the largest error on an interval. The choice of theorem therefore depends on whether the goal is existence, computation, or a uniform error guarantee.\end{enumerate}
\noindent\emph{Source for these results:} \cite{approximation_theory_ref1}.

\section*{5. Applications}
Approximation theory replaces a difficult function or model by a simpler one while estimating the error introduced by that replacement. The approximation is useful only when its domain, error bound, and numerical stability are known.
\begin{enumerate}[leftmargin=*,itemsep=0.35em]
\item \textbf{Computer graphics.} A curved object can be displayed as many small triangles. More triangles usually make the object look smoother, but the program must decide where extra triangles matter most.
\item \textbf{Signal compression.} A sound wave can be represented by its important frequencies. Removing tiny components saves storage, but produces distortion. The best method depends on whether a listener or a scientific instrument will judge the result.
\item \textbf{Engineering.} Engineers may replace a complicated physical system with a smaller model before testing a bridge or aircraft. The approximation is useful only after checking that neglected effects are smaller than the safety margin.
\item \textbf{Numerical equations.} Differential equations describing weather or fluid motion are often solved on a grid. The grid is an approximation to continuous space; making it finer can improve accuracy but also increases computation.
\end{enumerate}
\noindent\emph{Limitations.} An approximation can fail outside its intended range, near a sharp corner, or when small errors are amplified by a later calculation. More terms do not always help: rounding error, instability, or a divergent series can make a longer calculation worse.
\theoryconnections{approximation_theory}{%
\item Calculus supplies Taylor expansions and error estimates, while linear algebra turns approximation into a finite problem of choosing coefficients.
\item Norms and limits say what it means for an approximation to be close, and real or complex analysis proves when the error decreases.
\item For example, interpolation replaces an unknown curve by the polynomial through measured points, while least squares chooses the polynomial whose total squared error is smallest; the choice of norm determines what “best” means \cite{approximation_theory_ref1,approximation_theory_ref2}.
}{%
\item Approximation theory helps differential equations by replacing an infinite-dimensional solution with finitely many basis coefficients, which produces finite-element and spectral methods.
\item It helps signal processing by retaining the important frequencies and suppressing numerical noise, and it helps statistics by replacing a complicated relationship with a controlled model.
\item Its error bounds tell the downstream subject whether the numerical answer is trustworthy \cite{approximation_theory_ref1,approximation_theory_ref6}.
}

\section*{Glossary}
\begin{description}[leftmargin=*,style=nextline,itemsep=0.3em]

\item[\textbf{Minimax polynomial.}] The polynomial of a given degree that minimizes the largest absolute error over an interval, distributing the error evenly rather than forcing exact agreement at selected points.

\item[\textbf{Equioscillation.}] The property that the error of a best uniform approximation alternates in sign and achieves the same maximum magnitude at a prescribed number of points, certifying optimality.

\item[\textbf{Chebyshev polynomial.}] A sequence of polynomials defined by $T_n(\cos\theta)=\cos(n\theta)$ that oscillate between $-1$ and $1$ on $[-1,1]$ and provide near-optimal polynomial approximations.

\item[\textbf{Approximation error.}] The difference between a function and its approximant, measured in a chosen norm such as the maximum (worst-case) or mean-square error.

\item[\textbf{Remez algorithm.}] An iterative numerical procedure that converges to the minimax polynomial by adjusting test points to the true extrema of the error and solving a system of alternating equations.

\item[\textbf{Asymptotic expansion.}] A sequence of increasingly accurate approximations in which each neglected term is small relative to the preceding one; the series need not converge but is useful when truncated near its smallest term.

\item[\textbf{Uniform approximation.}] An approximation measured by the maximum absolute error over an interval, requiring the error to be small everywhere simultaneously.

\item[\textbf{Interpolation.}] Constructing a function that matches given data exactly at selected points; it may produce large errors between those points.

\item[\textbf{Least-squares approximation.}] Choosing an approximant to minimize the mean-square error rather than the worst-case error.

\item[\textbf{Splines.}] Piecewise polynomial functions joined smoothly at knots, providing local approximation with controlled global smoothness.

\end{description}

\section*{7. Sample Questions}
\emph{Exercise reference:} \cite{approximation_theory_ref1}. The cited edition is the source for the exercise set; consult its Exercises section for the official statement and numbering.
\begin{enumerate}[leftmargin=*,itemsep=0.9em]

\item \textbf{Exercise 1.} Use the Taylor polynomial of degree 2 for $e^x$ at $x=0$ to approximate $e^{0.1}$. Find the Lagrange remainder and bound the error.

\emph{Solution.} The degree-2 Taylor polynomial is $P_2(x)=1+x+x^2/2$. At $x=0.1$, $P_2(0.1)=1+0.1+0.005=1.105$. The Lagrange remainder is $R_2(x)=e^c\,x^3/3!$ for some $c\in(0,0.1)$. Since $e^c<e^{0.1}<1.106$, we get $|R_2(0.1)|<1.106\cdot(0.1)^3/6\approx 0.000184$. The true value is $e^{0.1}\approx 1.10517$, confirming the bound.

\item \textbf{Exercise 2.} Find the best linear (degree-1) polynomial approximation to $f(x)=x^2$ on $[-1,1]$ in the $L^\infty$ norm. What is the minimax error?

\emph{Solution.} By the equioscillation theorem, the best degree-1 polynomial $p(x)=a+bx$ must have error $f-p$ achieving $\pm E$ at three alternating points. By symmetry the best choice is $p(x)=1/2$, giving error $x^2-1/2$ which achieves $-1/2$ at $x=0$ and $+1/2$ at $x=\pm 1$. The minimax error is $E=1/2$.

\item \textbf{Exercise 3.} Compute the first three Chebyshev coefficients $c_0,c_1,c_2$ of $f(x)=x$ on $[-1,1]$, where $f(x)=\sum c_k T_k(x)$.

\emph{Solution.} Since $T_0(x)=1$, $T_1(x)=x$, and $T_2(x)=2x^2-1$, we have $f(x)=x=0\cdot T_0(x)+1\cdot T_1(x)+0\cdot T_2(x)$. Thus $c_0=0$, $c_1=1$, $c_2=0$. The function is exactly the first Chebyshev polynomial.

\item \textbf{Exercise 4.} How many equioscillation points are needed to certify that a degree-3 polynomial is the best uniform approximation to a continuous function on $[0,1]$?

\emph{Solution.} By the equioscillation theorem, a degree-$N$ best polynomial requires at least $N+2$ alternating extreme points. For $N=3$, this gives at least $5$ points at which the error alternates in sign with equal maximum magnitude.

\item \textbf{Exercise 5.} Estimate the degree $N$ needed for the truncated Chebyshev expansion of $e^x$ on $[-1,1]$ to achieve error less than $10^{-6}$, given that the Chebyshev coefficients satisfy $|c_k|\le 2/(k!\,2^k)$ (ignoring small constants).

\emph{Solution.} The error from truncating after $T_N$ is approximately $|c_{N+1}|+|c_{N+2}|+\cdots$, dominated by the first omitted term. We need $2/((N+1)!\,2^{N+1})<10^{-6}$. Trying $N=7$: $2/(8!\cdot 2^8)=2/(40320\cdot 256)\approx 1.94\times 10^{-7}<10^{-6}$. So $N=7$ suffices, giving a degree-7 polynomial.

\end{enumerate}
\section*{8. References}
\begin{thebibliography}{9}
\bibitem{approximation_theory_ref1} L. N. Trefethen, \emph{Approximation Theory and Approximation Practice}, SIAM, 1999.
\bibitem{approximation_theory_ref2} G. B. Folland, \emph{Real Analysis}, 2nd edition, Wiley, 2002.
\bibitem{approximation_theory_ref4} T. J. Rivlin, \emph{Chebyshev Polynomials: From Approximation Theory to Algebra and Number Theory}, 2nd edition, Wiley-Interscience, 2003.
\bibitem{approximation_theory_ref5} G. Szegő, \emph{Orthogonal Polynomials}, AMS, 1939.
\bibitem{approximation_theory_ref6} N. J. Higham, \emph{Accuracy and Stability of Numerical Algorithms}, 2nd edition, SIAM, 2002.
\bibitem{approximation_theory_ref7} N. I. Achieser, \emph{Theory of Approximation}, Dover, 1956.
\end{thebibliography}


\clearpage\chapter{Arakelov theory}
\section*{1. Overview}
\subsection*{Intuitive}
\paragraph{The problem}

Number theory treats every prime number as a finite place, but a number field also has ordinary real or complex size. Arakelov theory was developed to put these two kinds of information into one geometric framework. It adds a place at infinity to arithmetic geometry, equips the complex points with metrics, and defines intersections, heights, and Riemann--Roch formulas that combine finite prime data with analytic data.



\subsection*{Formal}
Let $K$ be a number field with ring of integers $\mathcal O_K$. Arakelov geometry treats finite prime ideals and archimedean embeddings as places of the same arithmetic space. An arithmetic divisor combines a finite divisor with real-valued data at the complex places, usually represented by a Green function or Hermitian metric. Its degree and intersection number are sums of finite algebraic contributions and archimedean analytic contributions.
\section*{2. History}
\subsection*{Historical development}
\begin{historylist}
\paragraph{History and motivation}

Suren Arakelov introduced the theory in the 1970s to study Diophantine geometry in higher dimensions. The motivating analogy is between a smooth projective curve over a finite field and an arithmetic surface over the integers. Prime ideals describe finite fibers, but $\operatorname{Spec}(\mathbb Z)$ has no prime ideal representing the ordinary absolute value. Arakelov's idea was to compactify the arithmetic picture by adding an infinite fiber and replacing missing algebraic information with a Hermitian metric on the complex points.

This was not an attempt to decorate number theory with analysis for its own sake. The goal was to make theorems for number fields resemble corresponding theorems for function fields. Intersections, degrees, and positivity should see every place of the field, including the real and complex ones \cite{arakelov_theory_ref2,arakelov_theory_ref4}.

\end{historylist}
\section*{3. Theory}
\subsection*{Core theory}
\paragraph{Finite places and the place at infinity}

Let $K$ be a number field and $\mathcal O_K$ its ring of integers. A nonzero prime ideal $\mathfrak p$ gives a finite place and a valuation measuring divisibility by $\mathfrak p$. The field also has embeddings into $\mathbb R$ or $\mathbb C$, called archimedean places. For $\mathbb Q$, the finite places come from primes $p$, while the infinite place is the usual absolute value $|\cdot|$.

For a nonzero rational number $x$, the product formula says that the product of its normalized absolute values over all places is 1:
\[
 \prod_v |x|_v=1.
\]
In logarithmic form, the sum of the local logarithms is zero. This identity is the arithmetic reason that a global height can ignore a common rescaling of homogeneous coordinates. It also illustrates the theory's central principle: a quantity that looks incomplete at one place becomes balanced after all places are included.

\paragraph{Arithmetic surfaces and metrised bundles}

An arithmetic surface is, roughly, a regular model $\mathfrak X$ of a smooth projective curve over $K$ defined over $\operatorname{Spec}(\mathcal O_K)$. Its finite fibers lie over prime ideals. After an embedding $K\to\mathbb C$, the complex points form a Riemann surface $X(\mathbb C)$. Arakelov theory equips line bundles or vector bundles on this complex surface with Hermitian metrics, compatible with complex conjugation.

The metric is the analytic replacement for the missing point at infinity. A metrised line bundle remembers both an algebraic divisor at finite places and a smooth metric at archimedean places. Its arithmetic degree measures finite multiplicities and the logarithmic size contributed at infinity. Green functions describe the potential created by a divisor and allow intersections involving the infinite component to be defined.

\paragraph{A simple height example}

If $p/q$ is a rational number in lowest terms, its logarithmic height is
\[
 h(p/q)=\log\max(|p|,|q|).
\]
The numerator and denominator record finite-prime information, while the maximum is an archimedean size. If one writes the same projective point as $(cp,cq)$, the product formula explains why the normalized global height does not depend on the chosen scaling. For algebraic numbers, the height averages the logarithmic absolute values over all embeddings and places.

\paragraph{Divisors with an infinite part}

For a field $K$, an arithmetic surface $\mathfrak X$, and a chosen complex embedding, an arithmetic or $c$-divisor has the form
\[
 D=\sum_i k_i C_i+\sum_{\sigma}\lambda_\sigma X_\sigma.
\]
Here the $C_i$ are irreducible codimension-one closed subsets of the arithmetic model, $k_i$ are integer coefficients at finite places, and $X_\sigma$ denotes the complex or real component associated with an archimedean embedding, with real coefficient $\lambda_\sigma$. The sum respects conjugate pairs of complex embeddings. These divisors form an abelian group.

The ordinary divisor records where a rational function vanishes or has poles. The infinite coefficients record the metric or Green-function contribution. Principal divisors satisfy a global compatibility relation because the product formula prevents an isolated local contribution from being arbitrary.

\paragraph{Arithmetic intersection theory}

Arakelov defined an intersection pairing for divisors on arithmetic surfaces attached to smooth projective curves over number fields. At finite places, two divisors intersect in the usual algebraic way on a fiber. At infinity, their Green functions contribute an analytic integral. The sum is a global arithmetic intersection number.

This lets one formulate an arithmetic Hodge index theorem, self-intersection inequalities, and height pairings. Jean-Benoît Bost developed potential-theoretic methods using Green functions and obtained arithmetic Lefschetz-type results. Gerd Faltings established an arithmetic Riemann--Roch theorem, a Noether formula, a Hodge index theorem, and positivity results for the dualizing sheaf. These are the arithmetic counterparts of familiar results for algebraic surfaces \cite{arakelov_theory_ref3,arakelov_theory_ref5}.

\paragraph{Arithmetic Chow groups}

An arithmetic cycle of codimension $p$ is a pair $(Z,g)$: an algebraic $p$-cycle $Z$ on the arithmetic scheme and a Green current $g$ that describes its archimedean contribution. Green currents generalize Green functions and can have controlled logarithmic singularities along $Z$.

The arithmetic Chow group $\widehat{\mathrm{CH}}^p(\mathfrak X)$ is obtained by quotienting arithmetic cycles by arithmetic versions of rational equivalence and trivial cycles. It is the natural home for arithmetic intersection products and Chern classes. The notation is heavier than the basic idea: algebraic cycles and analytic corrections are identified whenever they represent the same global arithmetic object.

\paragraph{Arithmetic Riemann--Roch}

Ordinary Grothendieck--Riemann--Roch describes how the Chern character changes under a proper map. In schematic form,
\[
 \operatorname{ch}(f_*E)=f_*(\operatorname{ch}(E)\operatorname{Td}(T_{X/Y})).
\]
The arithmetic version uses metrised bundles and arithmetic characteristic classes:
\[
 \widehat{\operatorname{ch}}(f_*[E])
 =f_*\bigl(\widehat{\operatorname{ch}}(E)
 \widehat{\operatorname{Td}}^{\,R}(T_{X/Y})\bigr).
\]
The modified Todd class contains an additional characteristic series involving zeta values and derivatives. The correction is the analytic price of including the infinite places. In applications, the theorem relates arithmetic degrees, determinants of cohomology, and heights of bundles under a smooth proper map \cite{arakelov_theory_ref6}.

\section*{4. Important theorems and results}
\begin{enumerate}[leftmargin=*,itemsep=0.35em]

\item \textbf{Product formula.} The product of normalized absolute values of a nonzero number over all places of a number field is one. It makes global heights invariant under scaling.

\item \textbf{Arakelov intersection theory.} Divisors on arithmetic surfaces have a global intersection number containing both finite algebraic terms and archimedean Green-function terms.

\item \textbf{Arithmetic Hodge index theorem.} Under suitable hypotheses, the arithmetic intersection pairing satisfies a negativity or positivity statement analogous to the classical Hodge index theorem.

\item \textbf{Arithmetic Riemann--Roch.} Pushforward of a metrised bundle is governed by arithmetic Chern and Todd classes, with an additional zeta-function correction.

\item \textbf{Faltings' finiteness theorem.} The arithmetic-geometric methods surrounding Arakelov theory contribute to finiteness theorems for objects such as abelian varieties over a number field with prescribed dimension and bounded bad reduction.
\end{enumerate}
\noindent\emph{Source for these results:} \cite{arakelov_theory_ref1}.

\section*{5. Applications}
Arakelov theory combines finite-prime information with archimedean geometry. Heights and arithmetic intersection numbers then measure the global size and interaction of arithmetic objects.
\begin{enumerate}[leftmargin=*,itemsep=0.35em]
\item \textbf{Diophantine equations.} Heights measure the global complexity of rational and algebraic points. A bound on height turns the search for infinitely many possible points into a finite search, and height inequalities can show that only finitely many points satisfy a Diophantine condition.
\item \textbf{Elliptic and abelian varieties.} Canonical heights combine local contributions and behave quadratically under multiplication. They help distinguish torsion points from points of infinite order and are central in arithmetic geometry.
\item \textbf{Arithmetic surfaces.} Intersections of sections, fibers, and metrised divisors quantify how arithmetic curves meet. The resulting inequalities are tools for proving finiteness and positivity statements.
\item \textbf{Analytic number theory.} Archimedean metrics and Green functions connect arithmetic intersection numbers with complex analysis, theta functions, and determinants of Laplacians.
\item \textbf{Generalizations.} Modern developments include Arakelov intersection theory in higher dimensions, Hodge--Arakelov theory, $p$-adic Hodge theory, and attempts such as $\mathbb F_1$ geometry to model a geometry over a hypothetical field with one element.
\item \textbf{Cryptographic height functions.} In lattice-based cryptography, the geometric height of an algebraic number measures the size of its minimal polynomial coefficients, which directly governs the hardness of lattice problems such as the shortest-vector problem. Arakelov-theoretic heights combine archimedean embeddings (Euclidean norm on coefficient vectors) with non-archimedean valuations at primes (divisibility of coefficients), yielding a single global complexity measure. Cryptographers use height bounds to estimate the security parameter of ring-based schemes such as NTRU and ring-LWE: the discriminant of the number field controls lattice volume, while the Mahler measure of the defining polynomial, an archimedean height quantity, controls the distribution of lattice points. Arakelov intersection pairings on arithmetic surfaces attached to these rings further predict how splitting behavior at finite places affects the algebraic structure available for key generation, providing a principled way to select number fields that resist the fastest known attacks \cite{arakelov_theory_ref3,arakelov_theory_ref4}.
\end{enumerate}


\theoryconnections{arakelov_theory}{%
\item Arakelov theory combines the finite-prime information of algebraic number theory with the ordinary complex geometry at the infinite places.
\item A divisor records zeros and poles algebraically, a line bundle packages that divisor, and a metric on the complex points records its archimedean size.
\item Intersection theory can then pair a finite contribution with a Green-function contribution, producing an arithmetic intersection number.
\item Heights arise from exactly this pairing, so they measure both divisibility at primes and size in the complex plane \cite{arakelov_theory_ref3,arakelov_theory_ref4}.
}{%
\item The resulting heights let Diophantine geometry compare solutions of an equation by a single global size.
\item On an elliptic curve, height inequalities can show that only finitely many rational points occur or can bound the search for them.
\item Arithmetic Riemann--Roch and arithmetic Chow groups use the same finite-plus-infinite bookkeeping to transfer geometric theorems into arithmetic statements, while arithmetic dynamics measures the growth of a point under repeated iteration \cite{arakelov_theory_ref1,arakelov_theory_ref5,arakelov_theory_ref6}.
}

\section*{Glossary}
\begin{description}[leftmargin=*,style=nextline,itemsep=0.3em]

\item[\textbf{Arithmetic surface.}] A regular model of a smooth projective curve defined over the ring of integers of a number field, whose finite fibers lie over prime ideals and whose complex points form a Riemann surface.

\item[\textbf{Place.}] Either a finite prime ideal of a number field or an archimedean embedding into $\mathbb R$ or $\mathbb C$; Arakelov theory treats both kinds on an equal footing.

\item[\textbf{Green function.}] A real-valued function attached to a divisor that describes its analytic behavior at archimedean places, serving as the metric counterpart to the algebraic data at finite primes.

\item[\textbf{Height.}] A global measure of the arithmetic size of a point or divisor, combining finite-prime valuations with archimedean absolute values; heights bound the search for rational solutions to Diophantine equations.

\item[\textbf{Product formula.}] The identity that the product of all normalized absolute values of a nonzero element over every place of a number field equals one, reflecting a perfect balance between finite and infinite places.

\item[\textbf{Arithmetic Chow group.}] A group of arithmetic cycles---pairs of algebraic cycles and Green currents---modulo an arithmetic equivalence relation, serving as the natural home for arithmetic intersection products.

\item[\textbf{Valuation.}] A function measuring the divisibility of elements of a number field by a prime ideal; it is the local building block of arithmetic.

\item[\textbf{Hermitian metric.}] A smooth inner product on a complex vector bundle that varies continuously over the base; in Arakelov theory it provides archimedean data.

\item[\textbf{Line bundle.}] A vector bundle of rank one over a variety or scheme; a divisor determines a line bundle.

\item[\textbf{Archimedean place.}] An embedding of a number field into R or C, measuring the ordinary absolute value; the infinite counterpart to a finite prime ideal.

\end{description}

\section*{7. Sample Questions}
\emph{Exercise reference:} \cite{arakelov_theory_ref1}. The cited edition is the source for the exercise set; consult its Exercises section for the official statement and numbering.
\begin{enumerate}[leftmargin=*,itemsep=0.9em]

\item \textbf{Exercise 1.} Verify the product formula for $x=6$ over $\mathbb{Q}$, i.e., compute $|6|_v$ for every place $v$ and confirm the product equals 1.

\emph{Solution.} The archimedean place gives $|6|_\infty=6$. At $p=2$: $|6|_2=|2^1\cdot 3|_2=2^{-1}=1/2$. At $p=3$: $|6|_3=3^{-1}=1/3$. At all other primes $|6|_p=1$. The product is $6\cdot(1/2)\cdot(1/3)\cdot 1\cdots=1$.

\item \textbf{Exercise 2.} Compute the logarithmic Weil height $h(p/q)$ for the rational numbers $5/2$, $7/1$, and $1/12$.

\emph{Solution.} By definition $h(p/q)=\log\max(|p|,|q|)$ for $p/q$ in lowest terms. Thus $h(5/2)=\log 5$, $h(7/1)=\log 7$, and $h(1/12)=\log 12$.

\item \textbf{Exercise 3.} Let $K=\mathbb{Q}(i)$. List all places of $K$ and state how many archimedean places $K$ has.

\emph{Solution.} The places of $K$ correspond to prime ideals of $\mathbb{Z}[i]$ (the finite places) and embeddings $K\hookrightarrow\mathbb{C}$ up to complex conjugation (the archimedean places). Since $K$ has one complex embedding pair, there is $r_1=0$ real places and $r_2=1$ complex place, giving exactly $1$ archimedean place. The finite places are the prime ideals $(1+i)$, $(3)$, $(2+i)$, $(2-i)$, and so on for each rational prime.

\item \textbf{Exercise 4.} Compute the height $h(\alpha)$ for $\alpha=1+\sqrt{5}$ viewed as an element of $\mathbb{Q}(\sqrt{5})$, using both conjugate embeddings $\sigma_1: \sqrt{5}\mapsto\sqrt{5}$ and $\sigma_2:\sqrt{5}\mapsto-\sqrt{5}$.

\emph{Solution.} The Weil height of an algebraic number is $h(\alpha)=\frac{1}{[K:\mathbb{Q}]}\sum_v n_v\log\max(|\sigma_v(\alpha)|,1)$ summed over all places. The two conjugate values are $1+\sqrt{5}\approx 3.236$ and $1-\sqrt{5}\approx -1.236$. Each archimedean place contributes $\log(1+\sqrt{5})+\log|1-\sqrt{5}|=\log(1+\sqrt{5})+\log(\sqrt{5}-1)=\log((1+\sqrt{5})(\sqrt{5}-1))=\log(5-1)=\log 4$. Dividing by $[K:\mathbb{Q}]=2$: $h(\alpha)=\frac{1}{2}\cdot 2\log 2=\log 2$.

\item \textbf{Exercise 5.} The degree of a divisor on an arithmetic surface includes both finite and archimedean contributions. If a divisor $D$ has finite part with degree $3$ at all non-archimedean places and archimedean contribution $-1$, what is the arithmetic degree $\widehat{\deg}(D)$?

\emph{Solution.} The arithmetic degree is the sum of the finite algebraic degree and the archimedean correction. Here $\widehat{\deg}(D)=3+(-1)=2$.

\end{enumerate}
\section*{8. References}
\begin{thebibliography}{99}
\bibitem{arakelov_theory_ref1} J. H. Silverman, \emph{The Arithmetic of Elliptic Curves}.
\bibitem{arakelov_theory_ref2} J. Neukirch, \emph{Algebraic Number Theory}.
\bibitem{arakelov_theory_ref3} S. Lang, \emph{Introduction to Arakelov Theory}.
\bibitem{arakelov_theory_ref4} E. Bombieri and W. Gubler, \emph{Heights in Diophantine Geometry}.
\bibitem{arakelov_theory_ref5} G. Faltings, \emph{Calculus on Arithmetic Surfaces}.
\bibitem{arakelov_theory_ref6} H. Gillet and C. Soulé, "An arithmetic Riemann--Roch theorem," \emph{Inventiones Mathematicae}.
\end{thebibliography}


\clearpage\chapter{Asymptotic theory}
\section*{1. Overview}
\subsection*{Intuitive}
\paragraph{The problem}

Exact formulas often contain more detail than a question needs. We may want to know how many primes lie below a huge number, how long an algorithm takes on a large input, or how a physical solution behaves when a parameter is very small. Asymptotic theory extracts the dominant behavior and gives a disciplined description of the corrections that have been neglected.



\subsection*{Formal}
For functions $f$ and $g$ with $g$ nonzero, $f(x)\sim g(x)$ as $x\to\infty$ means $f(x)/g(x)\to1$. The notation $f=O(g)$ means that $|f(x)|\leq C|g(x)|$ eventually for some constant $C$, while $f=o(g)$ means $f(x)/g(x)\to0$. An asymptotic expansion orders a quantity by size and states how the neglected remainder is controlled.
\section*{2. History}
\subsection*{Historical development}
\begin{historylist}
\paragraph{History and the guiding question}

The idea of estimating a complicated quantity by its leading behavior appears in ancient area and volume calculations. Calculus made limits and series systematic. Euler used expansions in analysis and mechanics; Gauss and Legendre used estimates in number theory. In the nineteenth and twentieth centuries, Cauchy, Poincaré, Laplace, and others developed approximations for integrals, differential equations, and perturbations.

The symbols $O$ and $o$ were developed through work associated with Paul Bachmann, Edmund Landau, and later analysts. The motivation was practical: a formula that says how an error shrinks can be more useful than a closed form that cannot be evaluated. Modern asymptotics links analysis, probability, algorithms, fluid mechanics, and mathematical physics \cite{asymptotic_theory_ref1,asymptotic_theory_ref3}.

\end{historylist}
\section*{3. Theory}
\subsection*{Core theory}
\paragraph{Limits and the basic comparisons}

For functions $f$ and $g$, the notation
\[
 f(x)\sim g(x)\qquad (x\to\infty)
\]
means
\[
 \lim_{x\to\infty}\frac{f(x)}{g(x)}=1,
\]
provided the ratio is defined eventually. The same idea applies as $x\to0$, as $x$ approaches a finite point, or along integers.

The statement $f(x)=O(g(x))$ means that there are constants $C>0$ and $x_0$ such that $|f(x)|\leq C|g(x)|$ for all $x\geq x_0$. It is an eventual upper bound, not necessarily a close estimate. The statement $f(x)=o(g(x))$ means $f(x)/g(x)\to0$; $f$ is negligible compared with $g$ in that limit.

For example,
\[
 n^2+3n\sim n^2,\qquad 3n^2+5n+7=O(n^2),\qquad
 5n=o(n^2).
\]
The first says relative error tends to zero, the second gives a scale, and the third says the linear term is negligible compared with the quadratic term. These statements do not say that the quantities are close for every small $n$.

\paragraph{Dominant terms and scales}

When $x$ is large, the highest power in a polynomial dominates:
\[
 3x^4-2x^2+7\sim3x^4.
\]
For exponentials, $e^x$ eventually dominates every polynomial; for logarithms, $\log x$ grows more slowly than every positive power of $x$. These comparisons produce a hierarchy of scales used in algorithm analysis and model reduction.

The number of unordered pairs chosen from $n$ objects is
\[
 \binom{n}{2}=\frac{n(n-1)}2\sim\frac{n^2}{2}.
\]
Thus doubling $n$ eventually multiplies the number of pairs by about four. For an algorithm, a quadratic operation count grows four times under the same doubling, while a cubic count grows eight times. The constants and lower terms matter for a real program, but the dominant exponent predicts large-input scaling.

\paragraph{Asymptotic expansions}

An asymptotic expansion is a sequence of increasingly accurate descriptions. We write
\[
 f(x)\sim g_1(x)+g_2(x)+\cdots
\]
when, for each fixed $N$,
\[
 f(x)-\sum_{k=1}^{N}g_k(x)=o(g_N(x)),
\]
with the terms forming a decreasing asymptotic scale. The series need not converge. Truncating after the smallest useful term may give the best approximation; adding more terms can eventually make it worse.

For small $x$,
\[
 \sin x=x-\frac{x^3}{6}+\frac{x^5}{120}-\cdots.
\]
The first term gives the main behavior, the next gives the leading correction, and a remainder estimate tells how far the approximation can be trusted. The same pattern occurs for large $x$ in integrals, special functions, and probability tails.

\paragraph{A divergent but useful expansion}

A formal geometric series
\[
 \frac1{1-w}=\sum_{n=0}^{\infty}w^n
\]
converges only for $|w|<1$. After multiplying by an exponential and integrating, one obtains an expansion such as
\[
 e^{-1/t}\operatorname{Ei}(1/t)\sim\sum_{n=0}^{\infty}n!\,t^{n+1}
\qquad(t\to0^+).
\]
The factorial coefficients mean the series diverges for every nonzero $t$. Yet for small $t$, the first few terms can give excellent numerical values before the terms begin to grow. This example warns us that asymptotic expansion and convergent series are different ideas \cite{asymptotic_theory_ref4}.

\paragraph{Prime numbers and other discrete examples}

Let $\pi(x)$ count primes not exceeding $x$. The prime number theorem states
\[
 \pi(x)\sim\frac{x}{\log x}.
\]
The quotient approaches 1 as $x$ grows. The approximation is not an exact count and can be poor at modest values, but it reveals the density of primes and leads to refined estimates.

Asymptotics also describes partial sums, partition numbers, binomial coefficients, recurrence sequences, and coefficients of generating functions. Stirling's formula,
\[
 n!\sim\sqrt{2\pi n}\left(\frac ne\right)^n,
\]
replaces an enormous exact product by a form suitable for probability and combinatorics.

\paragraph{Asymptotic distributions}

In probability and statistics, a sequence of random variables may have a limiting distribution as the sample size tends to infinity. The central limit theorem says that suitably normalized sums of independent observations approach a normal distribution. Likelihood-ratio statistics often approach chi-square distributions under regularity conditions. Such limits justify approximate confidence intervals and tests, but they do not give the exact finite-sample distribution.

The distinction matters. An asymptotic confidence interval may be accurate for a sample of a thousand observations and unreliable for a sample of ten. Non-asymptotic concentration inequalities provide finite-sample guarantees when a rigorous bound is required.

\paragraph{Integrals, differential equations, and physics}

Asymptotics estimates integrals by locating the region that contributes most. Laplace's method handles a large parameter in an exponential integral; the method of steepest descent extends it into the complex plane. Stationary phase studies rapidly oscillating integrals. These methods explain why a complicated integral can be controlled by a neighborhood of one point.

In differential equations, introduce a small dimensionless parameter and seek an outer solution away from a difficult region and an inner solution near it. Boundary-layer theory derives simplified equations from the Navier--Stokes equations when viscosity acts strongly only in a thin layer. Matched asymptotic expansions join the inner and outer descriptions. Singular perturbation is useful precisely because directly setting the small parameter to zero can reduce the order of the equation and lose a boundary condition.

Quantum field theory and statistical mechanics use formal expansions in a coupling constant or large parameter. Feynman diagrams organize terms, but the resulting series often diverge. Their asymptotic meaning can still be physically predictive after suitable truncation and resummation \cite{asymptotic_theory_ref3,asymptotic_theory_ref5}.

\paragraph{Asymptotics versus numerical analysis}

Asymptotic information and numerical information answer different questions. An estimate such as
\[
 f(x)=x^{-1}+O(x^{-2})
\]
may guarantee a small relative error only beyond a very large threshold. A numerical method may give a certified value at $x=100$ even when the asymptotic threshold is $10^4$. Conversely, a numerical table may not explain how the error changes as $x$ grows.

A careful result states the limiting regime, the size of the omitted term, the constants or threshold when known, and whether the approximation is absolute or relative. Matching an asymptotic formula with numerical computation is often the most useful approach.

\section*{4. Important theorems and results}
\begin{enumerate}[leftmargin=*,itemsep=0.35em]
\item \textbf{Prime number theorem.} $\pi(x)\sim x/\log x$, giving the leading density of primes.
\item \textbf{Stirling's formula.} $n!\sim\sqrt{2\pi n}(n/e)^n$, with further correction terms available.
\item \textbf{Taylor's theorem.} A smooth function equals a finite expansion plus a remainder that can be bounded using higher derivatives.
\item \textbf{Laplace's method.} For integrals with a large parameter, the dominant contribution usually comes from neighborhoods where the exponent is largest.
\item \textbf{Central limit theorem.} After suitable centering and scaling, sums of many independent random variables converge in distribution to a normal law under standard hypotheses.
\item \textbf{Matched asymptotic expansions.} Inner and outer approximations can be matched in an overlap region to solve singular perturbation problems with multiple scales.
\end{enumerate}
\noindent\emph{Source for these results:} \cite{asymptotic_theory_ref1}.

\section*{5. Applications}
Asymptotic theory is used when the dominant behaviour of a quantity matters more than an exact formula. Main terms describe the leading scale, while error terms determine whether the approximation is reliable in the range being studied.
\begin{enumerate}[leftmargin=*,itemsep=0.35em]
\item \textbf{Algorithms.} Asymptotic estimates compare running times as the input grows, even when exact timings depend on hardware.
\item \textbf{Approximation and numerics.} Error scales help choose mesh sizes, truncation points, and the number of terms to retain.
\item \textbf{Probability and statistics.} Limit laws describe the typical behavior of averages and estimators when the sample becomes large.
\item \textbf{Number theory.} Main terms estimate the distribution of primes and the number of solutions, while error terms measure the quality of the estimate.
\item \textbf{Differential equations and fluid mechanics.} Different scales separate a thin boundary layer from the slower bulk motion.
\item \textbf{Statistical mechanics and quantum theory.} Large-system limits and semiclassical limits replace a huge exact calculation by a controlled leading description.
\item \textbf{Accident-count modelling.} Rare-event approximations estimate risks when the number of opportunities is large but each event is unlikely.
\end{enumerate}
\noindent\emph{Limitations.} The theory depends on limits, inequalities, calculus, series, complex analysis, probability, and linear algebra. Its limitation is not that it is inaccurate; it is that it is deliberately about a limit. A hidden constant may be enormous, convergence may be slow, and a small error in one variable may be amplified by a later operation. Big-$O$ alone does not give a numerical error bound. It also cannot replace a finite-sample calculation when the sample or parameter is not in the asymptotic regime.


\theoryconnections{asymptotic_theory}{%
\item Limits and calculus identify the dominant scale as a variable becomes large or small.
\item Series and complex analysis then expose correction terms, while probability and statistics provide the language for typical rather than exact behavior.
\item For example, Stirling's formula replaces $n!$ by $\sqrt{2\pi n}(n/e)^n$ because the logarithm of the factorial is a sum whose leading area is an integral; the correction terms quantify what the leading scale misses \cite{asymptotic_theory_ref1,asymptotic_theory_ref5}.
}{%
\item Asymptotic estimates help algorithms predict running time, help numerical methods choose a mesh size, and help differential equations describe boundary layers where one scale changes rapidly.
\item In number theory, the main term describes the average distribution of primes while the error term measures the quality of the information.
\item In physics, an expansion can be useful even when it is not convergent, provided it is stopped before its terms begin to grow \cite{asymptotic_theory_ref1,asymptotic_theory_ref3}.
}
\section*{7. Sample Questions}
\emph{Exercise reference:} \cite{asymptotic_theory_ref1}. The cited edition is the source for the exercise set; consult its Exercises section for the official statement and numbering.
\begin{enumerate}[leftmargin=*,itemsep=0.9em]

\item \textbf{Exercise 1.} Prove that $n^2+3n\sim n^2$.

\textbf{Solution.} Divide by $n^2$: $(n^2+3n)/n^2=1+3/n$, which tends to 1 as $n\to\infty$.

\item \textbf{Exercise 2.} Show that $n\log n=o(n^2)$.

\textbf{Solution.} The ratio is $(n\log n)/n^2=(\log n)/n$, which tends to zero by l'Hôpital's rule or standard growth comparisons.

\item \textbf{Exercise 3.} Explain why $f(n)=O(n^2)$ does not mean $f(n)\sim n^2$.

\textbf{Solution.} The function $f(n)=n$ is $O(n^2)$ but $f(n)/n^2\to0$, not 1. Big-$O$ is only an eventual upper bound; asymptotic equivalence requires ratio 1.

\item \textbf{Exercise 4.} Use Stirling's formula to estimate $10!$.

\textbf{Solution.} The leading approximation is $\sqrt{20\pi}(10/e)^{10}$, about $3.60\times10^6$, close to the exact value $3,628,800$. Correction terms improve it.

\item \textbf{Exercise 5.} Why can adding terms to a divergent asymptotic expansion make an answer worse?

\textbf{Solution.} Early terms decrease and improve the approximation, but after the smallest term the divergent coefficients grow. The useful approximation is obtained by truncating near the least term, not by summing indefinitely.

\end{enumerate}
\section*{8. References}
\begin{thebibliography}{99}
\bibitem{asymptotic_theory_ref1} N. G. de Bruijn, \emph{Asymptotic Methods in Analysis}.
\bibitem{asymptotic_theory_ref3} C. M. Bender and S. A. Orszag, \emph{Advanced Mathematical Methods for Scientists and Engineers}.
\bibitem{asymptotic_theory_ref4} W. Balser, \emph{From Divergent Power Series to Analytic Functions}.
\bibitem{asymptotic_theory_ref5} F. W. J. Olver, \emph{Asymptotics and Special Functions}.
\end{thebibliography}


\clearpage\chapter{Automata theory}
\section*{1. Overview}
\subsection*{Intuitive}
\paragraph{The problem}

Whenever a rule must be followed one symbol, event, or moment at a time, we can ask whether a small machine can keep track of what matters. A door lock checks a code, a compiler checks a program, a communication protocol checks the order of messages, and a robot controller responds to sensor readings. Automata theory gives a common language for these situations. It replaces the physical machine by a precisely described model, then proves what that model can and cannot remember.


The word \emph{automaton} means a self-moving or self-acting device. In mathematics it means an ideal machine whose next action is determined by its present state and its current input. The machine may be tiny, or it may have an unbounded tape. The point is not that a real computer looks exactly like the model. The point is that the model isolates memory, choice, and time so that these can be studied separately.


\subsection*{Formal}
A deterministic finite automaton is a tuple $(Q,\Sigma,\delta,q_0,F)$: $Q$ is a finite set of states, $\Sigma$ an input alphabet, $\delta:Q\times\Sigma\to Q$ the transition function, $q_0$ the initial state, and $F$ the accepting states. Reading a word follows $\delta$ one symbol at a time; the accepted language consists of the words whose final state lies in $F$. Other machines change the available memory, which changes the class of languages they can recognize.
\section*{2. History}
\subsection*{Historical development}
\begin{historylist}
\paragraph{What problem does the theory solve?}

The central problem is recognition. Given a finite string of symbols, should it be accepted as legal? The symbols might be letters in a word, tokens in a computer program, bits in a message, or events in a safety protocol. A second problem is computation: if the machine is allowed to produce outputs, what transformation does it perform? A third is classification: which tasks can be solved with a given amount and kind of memory?

This classification is valuable because it prevents wasted effort. If a language can be recognized by a finite automaton, a small and fast implementation is possible. If it cannot, no clever rearrangement of a finite-state design will solve the problem; the design needs a stack, a queue, a tape, or another source of memory. The theory therefore explains both successful designs and impossibility results \cite{automata_theory_ref1,automata_theory_ref2}.

Automata theory developed in the middle of the twentieth century from systems theory, mathematical logic, abstract algebra, and the study of computation. Earlier systems theory often used differential equations to describe physical processes. Automata theorists instead used algebraic structures and discrete steps to describe information systems. Finite-state transducers, Turing machines, pushdown automata, and other models were developed by different communities and later brought together.

The 1956 collection \emph{Automata Studies} helped establish the subject as a separate discipline. Work associated with Claude Shannon, W. Ross Ashby, John von Neumann, Marvin Minsky, Edward Moore, and Stephen Kleene connected automata with information, cybernetics, circuits, and regular languages. In the same period Noam Chomsky described a hierarchy connecting formal grammars with machine power. Rabin and Scott proved that nondeterministic finite automata do not recognize more finite-word languages than deterministic finite automata. Myhill and Nerode gave an exact test for regularity and a way to count the states in the smallest finite automaton \cite{automata_theory_ref1,automata_theory_ref3}.

The researchers were asking a foundational question: when we say that a process is mechanical, how much memory and what kind of choice does it need? That question later became central to compiler construction, complexity theory, artificial intelligence, and formal verification.

\end{historylist}
\section*{3. Theory}
\subsection*{Core theory}
\paragraph{The basic machine}
\bookfigure{automata_computation}{A finite state records only the information needed for the next transition.}

An automaton receives symbols in discrete steps. The allowed symbols form an input alphabet, denoted by $\Sigma$. A finite sequence of symbols is a word; the set of all finite words, including the empty word, is written $\Sigma^*$. The machine has states. A state is a compact record of the past that the machine has decided is relevant.

At each step, a transition rule receives the current state and the next input symbol and specifies the next state. A machine that also reports a symbol has an output alphabet $\Gamma$ and an output rule. One general description is the quintuple
\[
 M=(\Sigma,\Gamma,Q,\delta,\lambda),
\]
where $\Sigma$ is the input alphabet, $\Gamma$ the output alphabet, $Q$ the set of states, $\delta:Q\times\Sigma\to Q$ the next-state rule, and $\lambda:Q\times\Sigma\to\Gamma$ the output rule. If $Q$ is finite, the machine is a finite automaton.

For language recognition, we normally replace the output rule by a designated start state $q_0$ and a set $F$ of accepting states. A run on $a_1a_2\ldots a_n$ is the state sequence $q_0,q_1,\ldots,q_n$ satisfying
\[
 q_i=\delta(q_{i-1},a_i).
\]
The transition rule extends to whole words by
\[
 \overline{\delta}(q,\varepsilon)=q,\qquad
 \overline{\delta}(q,wa)=\delta(\overline{\delta}(q,w),a).
\]
The word is accepted exactly when $\overline{\delta}(q_0,w)$ lies in $F$. The language recognized by the machine is
\[
 L=\{w\in\Sigma^*:\overline{\delta}(q_0,w)\in F\}.
\]
Thus an automaton is a finite description that may recognize an infinite set of words \cite{automata_theory_ref1,automata_theory_ref4}.

\paragraph{A worked example: even numbers of 1s}

Take $\Sigma=\{0,1\}$. Use two states: $E$ means that the number of 1s read so far is even, and $O$ means that it is odd. Start in $E$ and accept $E$. Reading 0 changes nothing. Reading 1 switches $E$ and $O$. The word 10110 visits
\[
 E\longrightarrow O\longrightarrow E\longrightarrow O\longrightarrow E\longrightarrow E,
\]
so it is accepted: it contains four 1s. The machine never needs to know the whole word; it only needs the parity of the count. This is the central design principle of finite automata: store exactly the summary that determines every possible future answer.

\paragraph{Finite automata and regular languages}

A deterministic finite automaton (DFA) has exactly one next state for every current state and input symbol. A nondeterministic finite automaton (NFA) may have several possible next states, including choices made without consuming an input symbol in common definitions. An NFA accepts when at least one allowed run reaches an accepting state. Nondeterminism is a convenient description language, not extra power for finite-word recognition: every NFA can be converted into a DFA by letting a DFA state represent the set of NFA states that could currently be active. This is the subset construction.

The languages recognized by finite automata are the \emph{regular languages}. They can equivalently be described by regular expressions and by regular grammars. This equivalence is one form of Kleene's theorem. It explains why a pattern such as
\[
 (0+1)^*01
\]
can be used by a text-search program: it is a compact notation for a finite machine that searches for a final 01.

Regular languages are closed under union, intersection, complement, concatenation, and repetition. Closure means that after combining two regular descriptions in one of these ways, the result is still regular. Product constructions prove union and intersection; complement is obtained by exchanging accepting and non-accepting states in a complete DFA; concatenation and repetition can be handled with NFAs.

Two proof tools show where finite memory stops being enough. The pumping lemma says that a sufficiently long word in a regular language contains a middle part that can be repeated any number of times while remaining in the language. It is useful for proving that languages such as $\{0^n1^n:n\geq0\}$ are not regular. The Myhill--Nerode theorem is stronger and more exact: a language is regular precisely when it has only finitely many distinguishable prefixes, and the number of equivalence classes is the number of states in its minimal DFA. Two prefixes are distinguishable when some continuation is accepted after one prefix but rejected after the other \cite{automata_theory_ref1,automata_theory_ref3}.

\paragraph{Memory changes the class of problems}

Finite-state machines have bounded memory. A pushdown automaton adds a stack. It can push a symbol when it sees an opening bracket and pop one when it sees the matching closing bracket. Since the stack grows with the input, it can recognize balanced parentheses and the language $\{0^n1^n:n\geq0\}$. Nondeterministic pushdown automata recognize the context-free languages, the class described by context-free grammars. Deterministic pushdown automata recognize a proper subclass: not every context-free language has a deterministic parser.

A queue machine stores symbols first-in, first-out and is Turing-complete. A Turing machine uses a tape, a read--write head, and a finite instruction table. It can move in both directions and use the tape as unbounded working memory. A linear bounded automaton has a tape whose usable portion is bounded by the input length and is associated with context-sensitive languages. Turing machines recognize the recursively enumerable languages when acceptance may occur but rejection may continue forever. This distinction leads to undecidable problems, such as the halting problem.

The hierarchy is not simply a ladder of faster machines. It is a hierarchy of the kinds of languages that can be recognized:
\[
\begin{array}{c}
\text{finite automata: regular languages}\\
\subset \text{pushdown automata: context-free languages}\\
\subset \text{linear bounded automata: context-sensitive languages}\\
\subset \text{Turing machines: recursively enumerable languages}.
\end{array}
\]
Two pushdown stores can simulate a Turing tape, so machines that look only slightly different can have radically different power. Deterministic and nondeterministic Turing machines have the same computability power, although complexity theory studies how much time or space their simulations require \cite{automata_theory_ref2,automata_theory_ref5}.

\section*{4. Important theorems and results}
\begin{enumerate}[leftmargin=*,itemsep=0.35em]

\item \textbf{Kleene's theorem.} A language is regular exactly when it can be described by a regular expression, recognized by a finite automaton, or generated by a regular grammar. This gives three equivalent ways to design the same finite-memory process.

\item \textbf{DFA--NFA equivalence.} Every nondeterministic finite automaton has an equivalent deterministic finite automaton. The deterministic machine may have more states, because it records sets of possible NFA states, but it accepts exactly the same words.

\item \textbf{Myhill--Nerode theorem.} A language is regular exactly when its prefixes fall into finitely many classes that no continuation can distinguish. The number of classes is the number of states in the minimal DFA. This is both a regularity test and a minimization theorem.

\item \textbf{Pumping lemma.} Every sufficiently long word in a regular language has a repeatable middle part. The lemma proves non-regularity when repeating any possible middle part eventually breaks the language rule.

\item \textbf{Chomsky hierarchy.} Regular, context-free, context-sensitive, and recursively enumerable languages correspond, in broad terms, to finite automata, pushdown automata, linear bounded automata, and Turing machines. The hierarchy explains why adding memory enlarges the class of recognizable problems.

\item \textbf{Church--Turing thesis.} Every effectively calculable step-by-step function can be computed by a Turing machine. This is a foundational claim about the meaning of "algorithm," not a theorem proved from a definition of ordinary calculation.
\end{enumerate}
\noindent\emph{Source for these results:} \cite{automata_theory_ref1}.

\section*{5. Applications}
Automata theory turns repeated rule-based behaviour into a finite state model whenever the required memory can be represented by states and transitions. This makes recognition, parsing, reachability, and verification precise and testable.
\begin{enumerate}[leftmargin=*,itemsep=0.35em]
\item \textbf{Compilers.} A finite automaton breaks a character stream into tokens such as identifiers, numbers, and keywords. A parser then uses a grammar or a pushdown-style method to check the nested structure of an expression. Finite memory is enough for local spelling rules, whereas nested parentheses require a stack.
\item \textbf{Text processing and hardware.} Regular expressions are compiled into finite automata, so a search for a date, an email-like pattern, or a forbidden word becomes a controlled walk through states. A traffic-light controller can be represented by states, and verification can prove that a red light is never followed immediately by a conflicting green light.
\item \textbf{Language and cellular models.} Context-free grammars model pieces of human language and have been used in artificial intelligence and natural-language processing. Cellular automata model local updates; Conway's Game of Life shows how simple rules can produce moving patterns, growth, and computation.
\item \textbf{Algebra and number theory.} Certain sets of irreducible polynomials over finite fields form regular languages under repeated composition. Automata can also infer a regular language from examples. These applications work because the automaton gives a finite, testable representation of an otherwise potentially infinite family \cite{automata_theory_ref1,automata_theory_ref7}.
\end{enumerate}

\noindent\emph{Limitations.} The subject depends on sets, functions, relations, induction, graphs, logic, and basic algebra. A finite automaton cannot recognize patterns that require unbounded memory, such as arbitrarily balanced parentheses; a pushdown automaton or a stronger model is then required. The choice of states and transitions must therefore match the memory and verification question being asked.


\theoryconnections{automata_theory}{%
\item Sets and functions define states and transitions; graphs draw those transitions; formal languages specify which strings are legal.
\item Logic then states properties of every possible run, and algebra can compress the transition behavior into a finite transformation structure.
\item For example, a finite automaton checking a binary number divisible by $3$ needs only three states, one for each possible remainder, because reading a digit updates the remainder by a fixed rule \cite{automata_theory_ref1,automata_theory_ref2}.
}{%
\item Automata theory helps a compiler turn source text into a syntax tree by recognizing tokens and grammatical structure.
\item The same transition model checks whether every state of a protocol avoids a forbidden condition, which is the basis of model checking and formal verification.
\item In digital circuits and cellular automata, the state-update rule is literally an automaton transition, so questions about long computations become questions about reachable states and accepted languages \cite{automata_theory_ref4,automata_theory_ref6}.
}

\section*{Glossary}
\begin{description}[leftmargin=*,style=nextline,itemsep=0.3em]

\item[\textbf{Finite automaton.}] An idealized machine with a finite set of states that reads input symbols one at a time, transitioning between states according to a fixed rule, and accepting the input if it ends in a designated accepting state.

\item[\textbf{Regular language.}] A set of strings that can be recognized by a finite automaton, equivalently described by a regular expression or a regular grammar.

\item[\textbf{Nondeterministic finite automaton (NFA).}] A finite automaton that may have several possible next states for a given state and input symbol, including $\varepsilon$-transitions; it accepts if at least one computational path reaches an accepting state.

\item[\textbf{Pushdown automaton.}] A finite automaton augmented with a stack that can grow without bound, giving it the memory needed to recognize context-free languages such as balanced parentheses.

\item[\textbf{Turing machine.}] A model of computation with a finite instruction table and an infinite tape that can be read and written in both directions; it captures the broadest class of algorithmically decidable problems.

\item[\textbf{Pumping lemma.}] A property of regular languages stating that any sufficiently long word contains a repeatable middle part; it is the standard tool for proving that a language is not regular.

\item[\textbf{Chomsky hierarchy.}] A classification of formal languages by the type of automaton needed to recognize them: regular, context-free, context-sensitive, and recursively enumerable, each class strictly containing the previous one.

\item[\textbf{Regular expression.}] A formal notation built from symbols, union, concatenation, and Kleene star that describes exactly the regular languages.

\item[\textbf{Myhill--Nerode theorem.}] A language is regular precisely when it has finitely many distinguishable equivalence classes of prefixes; the number equals the states in the minimal DFA.

\item[\textbf{Deterministic finite automaton (DFA).}] A finite automaton with exactly one transition from each state for each input symbol.

\item[\textbf{Church--Turing thesis.}] The claim that every effectively calculable function can be computed by a Turing machine.

\end{description}

\section*{7. Sample Questions}
\emph{Exercise reference:} \cite{automata_theory_ref1}. The cited edition is the source for the exercise set; consult its Exercises section for the official statement and numbering.
\begin{enumerate}[leftmargin=*,itemsep=0.9em]

\item \textbf{Exercise 1.} Build a DFA over $\Sigma=\{0,1\}$ that accepts exactly the strings whose integer value (in binary) is divisible by 3. Give the transition table and accept states.

\emph{Solution.} States are remainders $\{0,1,2\}$, start at $0$, accept $\{0\}$. Transition: reading bit $b$ from state $r$ goes to state $(2r+b)\bmod 3$. So $\delta(0,0)=0$, $\delta(0,1)=1$, $\delta(1,0)=2$, $\delta(1,1)=0$, $\delta(2,0)=1$, $\delta(2,1)=2$. The string $110_2=6$ is accepted: $0\to^1 1\to^1 0\to^0 0$.

\item \textbf{Exercise 2.} Prove that the language $L=\{0^n1^n:n\ge 1\}$ is not regular using the Myhill--Nerode theorem.

\emph{Solution.} Consider prefixes $0^i$ and $0^j$ with $i\ne j$. The continuation $1^i$ is in $L$ after $0^i$ but $1^j$ is in $L$ after $0^j$; since $i\ne j$, exactly one of $1^i,1^j$ is in $L$. Hence $0^i$ and $0^j$ are distinguishable. There are infinitely many equivalence classes, so $L$ is not regular.

\item \textbf{Exercise 3.} Convert the NFA with states $\{q_0,q_1,q_2\}$, alphabet $\{a,b\}$, start state $q_0$, accept state $\{q_2\}$, and transitions $\delta(q_0,a)=\{q_0,q_1\}$, $\delta(q_1,b)=\{q_2\}$, $\delta(q_2,b)=\{q_2\}$ (all other transitions empty) into an equivalent DFA.

\emph{Solution.} The subset construction gives DFA states as subsets of NFA states. Starting from $\{q_0\}$: on $a$, $\{q_0,q_1\}$; on $b$, $\emptyset$. From $\{q_0,q_1\}$: on $a$, $\{q_0,q_1\}$; on $b$, $\{q_2\}$. From $\{q_2\}$: on $a$, $\emptyset$; on $b$, $\{q_2\}$. From $\emptyset$: loops on both $a,b$. Accepting states are subsets containing $q_2$: just $\{q_2\}$. The DFA has 4 states.

\item \textbf{Exercise 4.} Construct a pushdown automaton that accepts $\{a^nb^n:n\ge 0\}$ by empty stack. Describe the stack behavior on input $aabb$.

\emph{Solution.} On reading $a$, push symbol $A$; on reading $b$, pop $A$; accept by empty stack at end. On $aabb$: start with empty stack. Read $a$: push $A$ (stack: $A$). Read $a$: push $A$ (stack: $AA$). Read $b$: pop $A$ (stack: $A$). Read $b$: pop $A$ (stack: $\varepsilon$). Stack is empty at end of input, so accept.

\item \textbf{Exercise 5.} Show that the language $L=\{w\in\{0,1\}^*: w\text{ has an equal number of 0s and 1s}\}$ is not context-free, using the pumping lemma for context-free languages.

\emph{Solution.} Choose $p$ as the pumping length. Let $s=0^p1^p\in L$ with $|s|=2p\ge p$. By the pumping lemma, $s=uvwxy$ with $|vwx|\le p$, $|vx|\ge 1$, and $uv^iwx^iy\in L$ for all $i$. Since $|vwx|\le p$, the substring $vwx$ lies entirely in the $0$'s, entirely in the $1$'s, or spans the boundary. If it lies in $0$'s, then $vx$ contains only $0$'s and pumping up gives more $0$'s than $1$'s. If in $1$'s, more $1$'s. If spanning the boundary with $|v|\ge 1$ zeros and $|x|\ge 1$ ones, then $v$ and $x$ must have the same count for pumping to preserve balance, but since $|vwx|\le p$ and $|vx|\ge 1$, we can choose $i=0$ (deletion) and break the balance. Hence $L$ is not context-free.

\end{enumerate}

\paragraph{Further study}

Start with sets, functions, induction, graphs, and simple logical statements. Then practice drawing finite automata before learning formal definitions. The natural sequence is regular expressions and DFAs, NFAs and minimization, regular-language proofs, context-free grammars and pushdown automata, Turing machines, decidability, and complexity. Sipser is a readable first university text; Hopcroft, Motwani, and Ullman give a traditional treatment; Eilenberg and Arbib develop algebraic and structural viewpoints.
\section*{8. References}
\begin{thebibliography}{99}
\bibitem{automata_theory_ref1} M. Sipser, \emph{Introduction to the Theory of Computation}, 3rd ed., Cengage.
\bibitem{automata_theory_ref2} J. E. Hopcroft, R. Motwani, and J. D. Ullman, \emph{Introduction to Automata Theory, Languages, and Computation}, 3rd ed., Pearson.
\bibitem{automata_theory_ref3} J. Myhill, "Finite automata and the representation of events," and A. Nerode, "Linear automaton transformations," in \emph{Automata Studies}, Princeton University Press.
\bibitem{automata_theory_ref4} S. C. Kleene, "Representation of events in nerve nets and finite automata," in \emph{Automata Studies}, Princeton University Press.
\bibitem{automata_theory_ref5} J. E. Hopcroft and J. D. Ullman, \emph{Introduction to Automata Theory, Languages, and Computation}, sections on pushdown, infinite-word, and Turing automata.
\bibitem{automata_theory_ref6} R. Alur, \emph{Principles of Cyber-Physical Systems}, MIT Press.
\bibitem{automata_theory_ref7} C. D. Manning and H. Schütze, \emph{Foundations of Statistical Natural Language Processing}, MIT Press.
\end{thebibliography}


\clearpage\chapter{Bass--Serre theory}
\section*{1. Overview}
\subsection*{Intuitive}
\paragraph{The problem}

A group can be difficult to understand when its elements are presented by many relations. Bass--Serre theory gives a geometric picture: make the group act on a tree, a connected graph with no loops. Vertices and edges then carry subgroups, and the way those subgroups are attached records whether the group is a free product, an amalgamated product, or an HNN extension.



\subsection*{Formal}
Let a group $G$ act without inversions on a tree $T$. Each vertex $v$ and edge $e$ has a stabilizer subgroup $G_v$ or $G_e$. The quotient graph $G\backslash T$, together with these stabilizers and the inclusion maps from edge groups to vertex groups, is a graph of groups. Bass--Serre theory reconstructs $G$ as the fundamental group of this graph of groups and uses the action on $T$ to study its structure.
\section*{2. History}
\subsection*{Historical development}
\begin{historylist}
\paragraph{History and motivation}

Jean-Pierre Serre developed the theory in the 1970s, with Hyman Bass making major contributions to its formalization. Serre's book \emph{Trees} grew partly from the study of algebraic groups whose Bruhat--Tits buildings are trees. The method quickly became central to geometric group theory and three-manifold topology.

The older algebraic constructions were free products with amalgamation and HNN extensions. Bass--Serre theory did not replace them; it explained them geometrically through covering spaces, fundamental groups, and group actions. A graph of groups is a one-dimensional analogue of an orbifold: groups sit at vertices and edges, and inclusions describe how local symmetry is glued into global symmetry \cite{bass_serre_theory_ref1,bass_serre_theory_ref3}.

\end{historylist}
\section*{3. Theory}
\subsection*{Core theory}
\paragraph{Trees and group actions}

A tree is a graph in which any two vertices are joined by exactly one simple path. Let a group $G$ act on a tree $X$ by graph automorphisms. The stabilizer of a vertex $v$ is
\[
 G_v=\{g\in G:gv=v\},
\]
and the stabilizer of an edge $e$ is defined similarly. We assume first that there are no edge inversions: no element flips an edge while fixing its midpoint.

The quotient graph $X/G$ records the orbits of vertices and edges. The quotient alone loses information, so attach the vertex stabilizer to each quotient vertex and the edge stabilizer to each quotient edge. The inclusion of an edge stabilizer into the stabilizer of its endpoint gives the attaching maps. The result is a graph of groups.

The action is geometric information about $G$. If $G$ acts freely on a tree, no nonidentity element fixes a vertex, and the quotient graph has fundamental group isomorphic to $G$. This immediately implies that a group acting freely on a tree is free.

\paragraph{Graphs of groups and their fundamental group}

Choose a maximal subtree $T$ in the underlying graph of a graph of groups $\mathcal A$. The fundamental group is generated by all vertex groups and one stable letter for each edge outside $T$, subject to the internal relations in every vertex group, the relations identifying the two images of each edge-group element along an edge in $T$, and conjugacy relations involving the stable letters for edges outside $T$.

If the graph has two vertices joined by one edge, with vertex groups $A$ and $B$ and edge group $C$ embedded in each, the fundamental group is the amalgamated product
\[
 A*_C B.
\]
If the graph has one vertex and one loop edge with edge group $C$ embedded twice in $A$, the fundamental group is an HNN extension
\[
 \langle A,t\mid t^{-1}\alpha(c)t=\beta(c)\text{ for }c\in C\rangle.
\]
Free products are the special case $C$ trivial.

\paragraph{Bass--Serre covering trees}

Every graph of groups $\mathcal A$ with a chosen base vertex has a canonical covering tree $\widetilde{\mathcal A}$. Its vertices can be represented by cosets of vertex groups in the fundamental group, and its edges by cosets of edge groups. The fundamental group acts on this tree without edge inversions. The quotient graph of groups is the original $\mathcal A$.

Conversely, if $G$ acts on a tree $X$ without inversions, the quotient graph of groups has a covering tree naturally isomorphic to $X$. The construction generalizes the universal cover of an ordinary graph. Ordinary covering theory has a fundamental group acting on a tree; Bass--Serre theory allows nontrivial stabilizers, so the local groups become part of the cover.

\paragraph{The fundamental theorem}

Let $G$ act on a tree $X$ without inversions, let $\mathcal A$ be the quotient graph of groups, and choose a base vertex $v$. Then
\[
 G\cong\pi_1(\mathcal A,v),
\]
and $X$ is equivariantly isomorphic to the Bass--Serre covering tree of $\mathcal A$. The isomorphisms respect both the group action and the graph structure.

This theorem is powerful in both directions. To understand an action, pass to the quotient graph of groups and obtain a group decomposition. To build an action for a group given by an amalgam or HNN extension, build its Bass--Serre tree.

\paragraph{Basic consequences}

If a graph of groups has a spanning tree and fundamental group $G$, the natural maps from its vertex groups into $G$ are injective. A finite-order element of $G$ is conjugate to an element of a vertex group. More generally, every finite subgroup is conjugate into some vertex group. If the quotient graph is finite and all vertex groups are finite, $G$ is virtually free: it contains a free subgroup of finite index.

Finite graphs of finitely generated vertex groups give finitely generated fundamental groups. If vertex groups are finitely presented and edge groups finitely generated, the resulting group is finitely presented. These statements make a decomposition useful for transferring finiteness properties from local groups to the global group.

\paragraph{Trivial and nontrivial actions}

An action is trivial if the whole group fixes a vertex. A graph of groups is correspondingly trivial when its underlying graph is a tree and the edge inclusions collapse the other vertex groups into the chosen base vertex group. Such a decomposition does not reveal a genuine splitting.

Nontrivial actions have no global fixed vertex and encode a real group splitting. The distinction matters in applications: rooted-tree actions may be interesting even when they fix a root, but Bass--Serre decomposition arguments generally seek actions that cannot be reduced to one vertex group.

\paragraph{Ends, property T, and accessibility}

Stallings' theorem says that a finitely generated group has more than one end exactly when it admits a nontrivial splitting over a finite subgroup. In Bass--Serre language, the group acts nontrivially on a tree with finite edge stabilizers. This turns a coarse geometric property of the Cayley graph into an algebraic decomposition.

If a group has Kazhdan's property (T), every action on a tree without inversions has a global fixed vertex. Such a group has Serre's property FA and cannot split nontrivially in the Bass--Serre sense. This provides an obstruction to decomposition.

Accessibility theorems bound the complexity of splittings. Dunwoody proved that a finitely presented group has a bound on the number of edges in reduced splittings over finite subgroups. Bestvina and Feighn generalized accessibility to splittings over small subgroups. These results say that decomposition cannot continue indefinitely once finiteness hypotheses are imposed \cite{bass_serre_theory_ref4,bass_serre_theory_ref5}.

\section*{4. Important theorems and results}
\begin{enumerate}[leftmargin=*,itemsep=0.35em]

\item \textbf{Fundamental theorem of Bass--Serre theory.} A group acting without inversions on a tree is the fundamental group of its quotient graph of groups, and the tree is the associated Bass--Serre covering tree.

\item \textbf{Free-action theorem.} A group acting freely on a tree is a free group.

\item \textbf{Kurosh subgroup theorem.} Subgroups of a free product decompose as a free product of conjugates of subgroups of the factors together with a free group.

\item \textbf{Stallings' ends theorem.} A finitely generated group has more than one end exactly when it splits nontrivially over a finite subgroup.

\item \textbf{Property FA consequence.} A group with Kazhdan's property (T) fixes a vertex in every tree action without inversions and therefore has no nontrivial Bass--Serre splitting.

\item \textbf{Accessibility.} Under suitable finite-presentation assumptions, the complexity of reduced splittings is bounded.
\end{enumerate}
\noindent\emph{Source for these results:} \cite{bass_serre_theory_ref1}.

\section*{5. Applications}
Bass--Serre theory applies group actions on trees to groups assembled from simpler pieces. The resulting tree records free products, amalgamations, and HNN extensions in a form that exposes their subgroup structure.
\begin{enumerate}[leftmargin=*,itemsep=0.35em]
\item \textbf{Free product.} The group $A*B$ acts on a tree whose vertices are cosets of $A$ and $B$ and whose edges connect compatible cosets. The quotient is one edge with two vertices and trivial edge group. The tree makes reduced words alternate between factors.
\item \textbf{Amalgamated product.} If two groups share a subgroup $C$, the group $A*_C B$ glues the two vertex groups along $C$. The edge stabilizer records exactly the elements identified in the gluing. This appears in three-manifold groups and combination theorems.
\item \textbf{HNN extension.} An HNN extension joins two copies of a subgroup inside one group with a stable letter. The Bass--Serre tree records repeated conjugation of one subgroup into another. It is a natural model for groups built by adding a symmetry that identifies two substructures.
\item \textbf{Three-manifolds and geometry.} Splittings over boundary or torus subgroups help analyze fundamental groups of three-manifolds. The tree records how pieces are assembled and interacts with geometric decompositions.
\item \textbf{Geometric group theory.} Actions on trees are one-dimensional models of nonpositive-curvature actions. They help study subgroup structure, free subgroups, ends, and algorithmic questions about group presentations.
\end{enumerate}


\theoryconnections{bass_serre_theory}{%
\item Group theory supplies the symmetries, graph theory supplies vertices and edges, and topology explains why a graph of groups behaves like a space assembled from pieces.
\item The fundamental group of a graph records the possible ways to travel around it; when groups are attached to vertices and edge groups describe the overlaps, the resulting action on a tree records how the whole group splits.
\item An amalgamated product identifies a shared subgroup, while an HNN extension identifies two copies of one subgroup \cite{bass_serre_theory_ref1,bass_serre_theory_ref2}.
}{%
\item The action on a tree turns an algebraic splitting into geometry: vertex stabilizers are the pieces, edge stabilizers are the overlaps, and the quotient graph displays the assembly plan.
\item This makes free products and subgroup structure visible, gives criteria for decomposing fundamental groups of three-manifolds, and supplies geometric tools for studying groups that are too large to classify by generators and relations alone \cite{bass_serre_theory_ref1,bass_serre_theory_ref4}.
}

\section*{Glossary}
\begin{description}[leftmargin=*,style=nextline,itemsep=0.3em]

\item[\textbf{Tree.}] A connected graph with no cycles; in Bass--Serre theory, a group acts on a tree and the structure of that action encodes how the group decomposes.

\item[\textbf{Graph of groups.}] A graph in which each vertex and edge is labeled by a group, with inclusion maps from edge groups to vertex groups; its fundamental group assembles the local groups into a global group.

\item[\textbf{Amalgamated product.}] A group formed by taking two groups $A$ and $B$ that share a common subgroup $C$ and identifying $C$ inside both, written $A*_C B$.

\item[\textbf{HNN extension.}] A group constructed from a single group $A$ and a new element $t$ (the stable letter) that conjugates one subgroup of $A$ into another.

\item[\textbf{Stabilizer.}] The subgroup of elements in a group action that fix a given vertex or edge; stabilizers record the local symmetry at each point of the tree.

\item[\textbf{Ends of a group.}] The number of ways the Cayley graph of a group extends to infinity after removing any finite region; more than one end signals a nontrivial splitting.

\item[\textbf{Free product.}] The amalgamated product $A*_C B$ when the common subgroup $C$ is trivial, producing the most general way to combine two groups without identifying any elements.

\item[\textbf{Covering tree.}] The canonical Bass--Serre tree associated with a graph of groups, whose vertices and edges are cosets of vertex and edge groups.

\item[\textbf{Quotient graph.}] The graph obtained by collapsing each orbit of vertices and edges under a group action.

\item[\textbf{Kazhdan's property (T).}] A group-theoretic property meaning every action on a tree without inversions has a global fixed vertex.

\item[\textbf{Virtually free.}] A group containing a free subgroup of finite index.

\end{description}

\section*{7. Sample Questions}
\emph{Exercise reference:} \cite{bass_serre_theory_ref1}. The cited edition is the source for the exercise set; consult its Exercises section for the official statement and numbering.
\begin{enumerate}[leftmargin=*,itemsep=0.9em]

\item \textbf{Exercise 1.} Let $G=\mathbb{Z}/2\mathbb{Z}*\mathbb{Z}/3\mathbb{Z}$ be the free product of cyclic groups of orders 2 and 3. Write down the graph of groups for this decomposition.

\emph{Solution.} The graph has two vertices, one labeled $\mathbb{Z}/2\mathbb{Z}$ and one labeled $\mathbb{Z}/3\mathbb{Z}$, joined by a single edge whose edge group is trivial. The fundamental group of this graph of groups is $\mathbb{Z}/2\mathbb{Z}*\mathbb{Z}/3\mathbb{Z}$. The Bass--Serre tree is a tree whose vertices alternate between cosets of $\mathbb{Z}/2\mathbb{Z}$ and cosets of $\mathbb{Z}/3\mathbb{Z}$.

\item \textbf{Exercise 2.} The group $\mathbb{Z}$ acts freely on the real line $\mathbb{R}$ (a tree) by $n\cdot x=x+n$. Identify the quotient graph of groups and verify the fundamental theorem.

\emph{Solution.} The quotient $\mathbb{R}/\mathbb{Z}$ is a single vertex with a single loop edge. Since the action is free, all stabilizers are trivial. The graph of groups has one vertex with trivial group and one loop edge with trivial group. By the fundamental theorem, $G\cong\pi_1(\mathcal{A},v)$. The fundamental group of a graph with one vertex and one edge is $\mathbb{Z}$, confirming $G\cong\mathbb{Z}$.

\item \textbf{Exercise 3.} Let $A=\mathbb{Z}/4\mathbb{Z}$, $B=\mathbb{Z}/6\mathbb{Z}$, and $C=\mathbb{Z}/2\mathbb{Z}$ embedded via the unique injections $C\hookrightarrow A$ and $C\hookrightarrow B$. Write the presentation of the amalgamated product $A*_C B$.

\emph{Solution.} Let $A=\langle a\mid a^4=1\rangle$, $B=\langle b\mid b^6=1\rangle$, and $C=\langle c\mid c^2=1\rangle$ with $C\hookrightarrow A$ sending $c\mapsto a^2$ and $C\hookrightarrow B$ sending $c\mapsto b^3$. The amalgamated product is $A*_C B=\langle a,b\mid a^4=1,\,b^6=1,\,a^2=b^3\rangle$.

\item \textbf{Exercise 4.} Show that the fundamental group of the figure-eight graph (one vertex, two loop edges $e_1,e_2$) with trivial vertex and edge groups is the free group $F_2$.

\emph{Solution.} Choose the single vertex as the spanning tree (which is trivially a tree). Both edges $e_1,e_2$ are outside the spanning tree, each contributing a stable letter. The fundamental group is generated by the trivial vertex group together with stable letters $t_1,t_2$ for $e_1,e_2$, with no relations beyond those in the vertex group. Hence $\pi_1(\mathcal{A})=\langle t_1,t_2\rangle\cong F_2$, the free group on two generators.

\item \textbf{Exercise 5.} Let $F$ be a free group of rank 2 and $H$ the subgroup generated by $a^2$ and $b^2$ where $a,b$ are the free generators. Use the Kurosh subgroup theorem to describe $H$.

\emph{Solution.} The Kurosh subgroup theorem states that a subgroup of a free product decomposes as a free product of conjugates of subgroups of the factors, together with a free group. Since $F_2=\mathbb{Z}*\mathbb{Z}$ and $H\le F_2$ is itself free (every subgroup of a free group is free), the theorem says $H$ is a free group. The rank of $H$ can be computed via the Schreier index formula: $H$ has index $[F_2:H]=\infty$ (since $\langle a^2,b^2\rangle$ has infinite index), so $H$ is free of infinite rank.

\end{enumerate}
\section*{8. References}
\begin{thebibliography}{99}
\bibitem{bass_serre_theory_ref1} J.-P. Serre, \emph{Trees}, Springer, 1980.
\bibitem{bass_serre_theory_ref2} H. Bass, "Covering theory for graphs of groups," \emph{Journal of Pure and Applied Algebra}, 1993.
\bibitem{bass_serre_theory_ref3} A. Hatcher, \emph{Algebraic Topology}, Cambridge University Press, 2002.
\bibitem{bass_serre_theory_ref4} W. Dicks and M. J. Dunwoody, \emph{Groups Acting on Graphs}, Cambridge University Press, 1989.
\bibitem{bass_serre_theory_ref5} D. E. Cohen, \emph{Combinatorial Group Theory: A Topological Approach}, Cambridge University Press, 1989.
\end{thebibliography}


\clearpage\chapter{Bifurcation theory}
\section*{1. Overview}
\subsection*{Intuitive}
\paragraph{The problem}

A parameter can change smoothly while the behavior of a system changes suddenly. A population may remain near zero and then settle at a positive level; a circuit may begin to oscillate; a fluid flow may lose stability and develop a new pattern. Bifurcation theory locates and classifies the parameter values at which the qualitative structure of solutions changes.



\subsection*{Formal}
For an equation $F(x,\lambda)=0$, a bifurcation is a parameter value $\lambda_0$ at which the nearby solution set changes qualitatively. At an equilibrium, the linearization $D_xF$ determines local behavior; loss of invertibility, together with suitable nondegeneracy conditions, signals a possible branch or stability change. Normal-form calculations reduce the local behavior to a simpler equation while preserving the relevant qualitative features.
\section*{2. History}
\subsection*{Historical development}
\begin{historylist}
\paragraph{History and motivation}

The word bifurcation was introduced by Henri Poincaré in 1885 while studying changes in equilibria and periodic solutions of differential equations. Later work by Andronov, Hopf, Arnold, Smale, and many others developed local normal forms and global dynamical methods. The researchers were trying to understand why a small change in a physical parameter could change the number or stability of solutions, something a single solution formula could not reveal.

Bifurcation theory applies to continuous systems such as ordinary, delay, and partial differential equations, and to discrete systems such as iterated maps. It is closely related to dynamical systems, chaos, catastrophe theory, control, and nonlinear analysis \cite{bifurcation_theory_ref1,bifurcation_theory_ref3}.

\end{historylist}
\section*{3. Theory}
\subsection*{Core theory}
\paragraph{Equilibria, stability, and a changing parameter}
\bookfigure{dynamics_chaos_bifurcation}{A bifurcation diagram records how long-term behavior changes with a parameter.}

Consider
\[
 \dot x=f(x,\lambda),
\]
where $\lambda$ is a parameter. An equilibrium $x_*$ satisfies $f(x_*,\lambda)=0$. To test its local stability, linearize:
\[
 \dot u=D_xf(x_*,\lambda)u.
\]
If all eigenvalues have negative real part, small disturbances decay; if one has positive real part, the equilibrium is unstable. At a bifurcation, a stability eigenvalue reaches the boundary: in continuous time its real part is zero, while for a discrete map $x_{n+1}=f(x_n,\lambda)$ a multiplier has modulus one.

The scalar example
\[
 \dot x=rx-x^3
\]
has $x=0$ for all $r$. When $r<0$, the derivative at zero is $r<0$, so zero is stable. When $r>0$, zero is unstable and two new equilibria $x=\pm\sqrt r$ appear; their derivative is $r-3r=-2r<0$, so they are stable. The phase portrait changes at $r=0$. This is a supercritical pitchfork in a symmetric system.

\paragraph{Local bifurcations}

A local bifurcation can be understood in an arbitrarily small neighborhood of the changing equilibrium or periodic orbit. The implicit-function theorem explains why a smooth branch of equilibria persists when the derivative with respect to $x$ is invertible. Failure of this condition is the signal that a branch can turn, cross, or split.

\textbf{Saddle-node.} The normal form
\[
 \dot x=\lambda-x^2
\]
has two equilibria for $\lambda>0$, one double equilibrium at $\lambda=0$, and no equilibria for $\lambda<0$. A stable and an unstable equilibrium are created or destroyed together.

\textbf{Transcritical.} In
\[
 \dot x=\lambda x-x^2,
\]
the branches $x=0$ and $x=\lambda$ cross and exchange stability. This often describes a threshold between two competing states.

\textbf{Pitchfork.} In the symmetric example above, one equilibrium becomes unstable and two symmetric equilibria appear. Symmetry is important: without a symmetry or constraint, a pitchfork usually unfolds into other bifurcations.

\textbf{Hopf.} A pair of complex-conjugate eigenvalues crosses the imaginary axis. An equilibrium changes stability and a small periodic orbit is born or disappears. Hopf bifurcation explains the onset of oscillations in chemical reactions, biological rhythms, lasers, and electrical circuits.

For maps, the corresponding conditions use multipliers. A multiplier $+1$ may lead to a saddle-node, transcritical, or pitchfork; a multiplier $-1$ often produces period doubling. A pair on the unit circle can generate a discrete analogue of Hopf behavior.

\paragraph{Codimension}

The codimension is the number of independent parameters that must be adjusted to produce a generic bifurcation. Saddle-node and Hopf bifurcations are generically codimension one: varying one parameter is usually enough. Transcritical and pitchfork normal forms are also drawn with one parameter, although symmetry or another condition may be needed for them to occur robustly.

The Bogdanov--Takens bifurcation is a standard codimension-two example. Two conditions must be met, often producing nearby saddle-node, Hopf, and homoclinic curves in a two-parameter diagram. Codimension counts how many coincidences must be arranged, not how visually complicated the resulting dynamics will be.

\paragraph{Global bifurcations}

Global bifurcations cannot be detected from the eigenvalues of one equilibrium. They involve large invariant sets colliding or changing topology.

\textbf{Homoclinic bifurcation.} A periodic orbit may collide with a saddle point and its stable and unstable manifolds. The period can grow very large near the collision, and complicated or chaotic behavior may appear.

\textbf{Heteroclinic bifurcation.} A cycle of connections among two or more saddle points changes stability. Resonance conditions on eigenvalues can create or destroy periodic orbits; transverse intersections can alter the global connection.

\textbf{Infinite-period bifurcation.} A stable node and a saddle can be created on a limit cycle while the oscillation slows and its period tends to infinity. Beyond the threshold, the cycle disappears and the equilibria remain.

\textbf{Blue-sky catastrophe and crises.} A periodic orbit can collide with a nonhyperbolic invariant cycle, or a chaotic attractor can collide with a basin boundary. These events can cause a sudden enlargement or destruction of the observed attractor.

The local/global distinction is practical. Linearization is powerful near an equilibrium, but it cannot reveal that a distant periodic orbit will collide with a saddle.

\paragraph{Bifurcation diagrams and normal forms}

A bifurcation diagram plots equilibria or periodic-orbit amplitudes against the parameter. Solid and dashed curves often indicate stable and unstable branches. A normal form is a simplified equation obtained by coordinate changes and scaling that preserves the behavior relevant near the bifurcation. It removes terms that do not affect the leading qualitative change.

Center-manifold reduction explains why a high-dimensional system can often be reduced to a one- or two-dimensional equation near a nonhyperbolic equilibrium. The center manifold contains the slow directions associated with zero-real-part eigenvalues; stable directions decay and can be slaved to them. Normal-form calculations then classify the possible local behavior.

\section*{4. Important theorems and results}
\begin{enumerate}[leftmargin=*,itemsep=0.35em]

\item \textbf{Implicit-function theorem.} If the derivative with respect to the state is invertible, a nearby solution branch continues smoothly with the parameter. Bifurcation requires failure of this regularity condition.

\item \textbf{Linear stability criterion.} Eigenvalues with negative real parts give local decay in continuous time; a positive real part gives instability. For maps, stability corresponds to multipliers inside the unit circle.

\item \textbf{Center-manifold theorem.} Near a nonhyperbolic equilibrium, the local dynamics relevant to bifurcation can be reduced to a center manifold of lower dimension.

\item \textbf{Hopf bifurcation theorem.} Under nondegeneracy and transversality conditions, a pair of eigenvalues crossing the imaginary axis produces a branch of small periodic orbits.

\item \textbf{Saddle-node normal form.} The equation $\dot x=\lambda-x^2$ captures the generic creation or destruction of an equilibrium pair.

\item \textbf{Bogdanov--Takens unfolding.} A codimension-two double-zero eigenvalue organizes nearby saddle-node, Hopf, and homoclinic bifurcation curves.
\end{enumerate}
\noindent\emph{Source for these results:} \cite{bifurcation_theory_ref1}.

\section*{5. Applications}
Bifurcation theory is useful when changing a parameter changes the number, stability, or type of solutions of a differential equation. The critical parameter identifies where a qualitative transition occurs.
\begin{enumerate}[leftmargin=*,itemsep=0.35em]
\item \textbf{Population dynamics.} A harvesting or reproduction parameter can create a stable positive population through a transcritical or saddle-node event. A small change past the threshold may remove the only sustainable equilibrium.
\item \textbf{Chemical and biological oscillators.} Feedback in a reaction network can cause a Hopf bifurcation. As a control parameter crosses the threshold, a stable equilibrium gives way to a stable periodic concentration cycle. Cell-cycle switches and biochemical networks use bistable and irreversible switches whose output changes near a bifurcation point.
\item \textbf{Fluid mechanics and climate.} Increasing the Reynolds number can destabilize a steady flow and create vortices or periodic shedding. Climate models may have multiple equilibria, with transitions between them organized by saddle-node or global bifurcations.
\item \textbf{Lasers, circuits, and quantum systems.} Laser intensity can emerge through a threshold bifurcation, and electronic feedback circuits can begin to oscillate at a Hopf point. In atomic and molecular systems, bifurcations of classical orbits leave enhanced signatures in quantum spectra. Resonant tunnelling diodes, kicked tops, and coupled quantum wells provide model examples.
\item \textbf{Engineering and control.} Bifurcation diagrams show where a controller will lose stability or where a desired operating state disappears. Continuation methods follow solution branches numerically and reveal turning points that ordinary simulation might miss \cite{bifurcation_theory_ref2,bifurcation_theory_ref4}.
\end{enumerate}


\theoryconnections{bifurcation_theory}{%
\item Calculus linearizes a differential equation near an equilibrium, and linear algebra turns that linearization into eigenvalues that indicate growth, decay, or oscillation.
\item Dynamical systems and phase portraits then follow the nearby trajectories, while topology explains why branches of equilibria cannot simply disappear without a change in the equations.
\item Numerical continuation traces those branches as a parameter changes, so a bifurcation is detected as a structural change rather than as one surprising numerical output \cite{bifurcation_theory_ref1,bifurcation_theory_ref2}.
}{%
\item Bifurcation theory tells control engineers where a stable operating point will become unstable and tells biologists where a population equilibrium can split into several regimes.
\item In climate and fluid models, a parameter crossing a threshold can create oscillations or multiple states; in laser models it can turn a zero signal into a sustained one.
\item Chaos theory often begins after repeated bifurcations, while catastrophe theory classifies the local shapes of the resulting transitions \cite{bifurcation_theory_ref1,bifurcation_theory_ref4}.
}

\section*{Glossary}
\begin{description}[leftmargin=*,style=nextline,itemsep=0.3em]

\item[\textbf{Bifurcation.}] A parameter value at which the qualitative structure of solutions---their number, stability, or type---changes abruptly even though the parameter itself changes smoothly.

\item[\textbf{Saddle-node bifurcation.}] The generic creation or destruction of a pair of equilibria (one stable, one unstable) as a parameter crosses a threshold, captured by the normal form $\dot x=\lambda-x^2$.

\item[\textbf{Hopf bifurcation.}] A local bifurcation in which a pair of complex-conjugate eigenvalues crosses the imaginary axis, causing a steady state to lose stability and a small periodic orbit to be born.

\item[\textbf{Normal form.}] A simplified equation obtained by coordinate changes and scaling that preserves the essential qualitative behavior near a bifurcation, stripping away irrelevant higher-order terms.

\item[\textbf{Center manifold.}] The invariant manifold near a nonhyperbolic equilibrium that captures the slow dynamics responsible for bifurcation; the remaining fast directions decay and can be slaved to it.

\item[\textbf{Codimension.}] The number of independent parameters that must be simultaneously tuned to produce a generic bifurcation; codimension-one bifurcations occur generically with one parameter.

\item[\textbf{Homoclinic bifurcation.}] A global bifurcation in which a periodic orbit collides with a saddle point along its own stable and unstable manifolds, often leading to chaotic behavior.

\item[\textbf{Pitchfork bifurcation.}] A bifurcation in which one equilibrium loses stability and two new symmetric equilibria are created or destroyed.

\item[\textbf{Transcritical bifurcation.}] A bifurcation in which two equilibrium branches cross and exchange stability as a parameter varies.

\item[\textbf{Bifurcation diagram.}] A plot showing how the number, stability, and type of equilibria change as a parameter varies.

\item[\textbf{Period doubling.}] A bifurcation of a periodic orbit in which the period doubles, often occurring when a multiplier reaches $-1$.

\end{description}

\section*{7. Sample Questions}
\emph{Exercise reference:} \cite{bifurcation_theory_ref1}. The cited edition is the source for the exercise set; consult its Exercises section for the official statement and numbering.
\begin{enumerate}[leftmargin=*,itemsep=0.9em]

\item \textbf{Exercise 1.} Consider $\dot x=\lambda x-x^2$. Find all equilibria, determine their stability for $\lambda=1$ and $\lambda=-1$, and classify the bifurcation at $\lambda=0$.

\emph{Solution.} Equilibria: $x=0$ and $x=\lambda$ (when $\lambda\ne 0$). The derivative $f_x=\lambda-2x$. At $x=0$: $f_x=\lambda$, stable for $\lambda<0$, unstable for $\lambda>0$. At $x=\lambda$: $f_x=-\lambda$, stable for $\lambda>0$, unstable for $\lambda<0$. The two branches cross and exchange stability at $\lambda=0$: this is a transcritical bifurcation. For $\lambda=1$, the equilibria are $0$ (unstable) and $1$ (stable). For $\lambda=-1$, the equilibria are $0$ (stable) and $-1$ (unstable).

\item \textbf{Exercise 2.} Find the equilibria and bifurcation diagram for $\dot x=\lambda-x^3$. What type of bifurcation occurs at $\lambda=0$?

\emph{Solution.} Equilibria satisfy $x^3=\lambda$, so $x=\lambda^{1/3}$ (one real root for all $\lambda$). The derivative $f_x=-3x^2$ is always $\le 0$, and equals $0$ only at $x=\lambda=0$. The equilibrium $x=\lambda^{1/3}$ is stable for all $\lambda\ne 0$ (since $f_x=-3\lambda^{2/3}<0$). At $\lambda=0$, $f_x=0$ but $f_{xx}=-6x=0$ and $f_{xxx}=-6\ne 0$. This is a degenerate saddle-node (a cusp-type degeneracy): the equilibrium exists for all $\lambda$ and does not lose or gain equilibria, but the linearization degenerates.

\item \textbf{Exercise 3.} For the system $\dot x=r x-x^3$, compute the eigenvalue at $x=0$ when $r=0.5$ and $r=-0.5$, and identify which equilibria are stable.

\emph{Solution.} The derivative at $x=0$ is $r$. When $r=0.5>0$, the eigenvalue is $0.5>0$, so $x=0$ is unstable. The two nonzero equilibria $x=\pm\sqrt{0.5}$ each have derivative $r-3r=-2r=-1<0$, so they are stable. When $r=-0.5<0$, the eigenvalue at $x=0$ is $-0.5<0$, so $x=0$ is stable. No nonzero real equilibria exist.

\item \textbf{Exercise 4.} For the two-dimensional system $\dot x=\lambda-x^2$, $\dot y=-y$, determine the equilibria for $\lambda>0$, $\lambda=0$, and $\lambda<0$, and identify which are stable nodes versus saddle points.

\emph{Solution.} Equilibria: $(\pm\sqrt{\lambda},0)$ for $\lambda>0$; $(0,0)$ for $\lambda\le 0$. The Jacobian is $\begin{pmatrix}-2x&0\\0&-1\end{pmatrix}$. For $\lambda>0$ at $(\sqrt{\lambda},0)$: eigenvalues $-2\sqrt{\lambda}<0$ and $-1<0$, so stable node. At $(-\sqrt{\lambda},0)$: eigenvalues $+2\sqrt{\lambda}>0$ and $-1<0$, so saddle point. For $\lambda<0$: no equilibria. At $\lambda=0$: $(0,0)$ has Jacobian $\begin{pmatrix}0&0\\0&-1\end{pmatrix}$, a degenerate (nonhyperbolic) equilibrium.

\item \textbf{Exercise 5.} Consider $\dot x=r+x-x^3$. Find the turning points (fold points) of the equilibrium curve by solving $\dot x=0$ and $\ddot x=0$ simultaneously, and determine the values of $r$ at which folds occur.

\emph{Solution.} Equilibria satisfy $x^3-x=r$. The fold condition is $\frac{d}{dx}(x^3-x)=3x^2-1=0$, giving $x=\pm 1/\sqrt{3}$. At $x=1/\sqrt{3}$: $r=1/(3\sqrt{3})-1/\sqrt{3}=-2/(3\sqrt{3})\approx -0.385$. At $x=-1/\sqrt{3}$: $r=-1/(3\sqrt{3})+1/\sqrt{3}=2/(3\sqrt{3})\approx 0.385$. These are the two saddle-node (fold) bifurcation points. Between them, three equilibria exist; outside, only one.

\end{enumerate}
\section*{8. References}
\begin{thebibliography}{99}
\bibitem{bifurcation_theory_ref1} S. H. Strogatz, \emph{Nonlinear Dynamics and Chaos}, 2nd edition, Westview Press, 2015.
\bibitem{bifurcation_theory_ref2} Y. A. Kuznetsov, \emph{Elements of Applied Bifurcation Theory}, 3rd edition, Springer, 2004.
\bibitem{bifurcation_theory_ref3} V. I. Arnold, \emph{Geometrical Methods in the Theory of Ordinary Differential Equations}, 2nd edition, Springer, 1983.
\bibitem{bifurcation_theory_ref4} S. Wiggins, \emph{Introduction to Applied Nonlinear Dynamical Systems and Chaos}, 2nd edition, Springer, 2003.
\end{thebibliography}


\clearpage\chapter{Braid theory}
\section*{1. Overview}
\subsection*{Intuitive}
\paragraph{The problem}

When several strands move from one row of points to another, their over-and-under crossings remember more than the final permutation of the endpoints. Braid theory turns these motions into algebra. Braids can be stacked, undone, represented by generators and relations, closed to make knots, and studied through configuration spaces and mapping classes.



\subsection*{Formal}
The braid group $B_n$ is generated by $\sigma_1,\ldots,\sigma_{n-1}$, where $\sigma_i$ crosses strand $i$ over strand $i+1$. Its defining relations are $\sigma_i\sigma_j=\sigma_j\sigma_i$ when $|i-j|\ge2$ and $\sigma_i\sigma_{i+1}\sigma_i=\sigma_{i+1}\sigma_i\sigma_{i+1}$. Multiplication is stacking braids, and the map $B_n\to S_n$ records the permutation of endpoints.
\section*{2. History}
\subsection*{Historical development}
\begin{historylist}
\paragraph{History and the basic picture}
\bookfigure{knot_braid}{A braid records the ordered crossings of moving strands; closing it produces a knot or link.}

Braid groups were introduced explicitly by Emil Artin in 1925, although Adolf Hurwitz had already encountered their configuration-space meaning in work on monodromy. Artin's 1947 presentation made the groups accessible through generators and relations. Joan Birman's 1974 book connected braids, links, and mapping class groups and became a central reference.

Take $n$ points in a row and connect them to another row. Each strand moves monotonically from left to right. Two diagrams represent the same braid when one can be continuously deformed into the other without strands passing through one another. Over-crossing and under-crossing are part of the data. Stack one braid after another to compose them. Parallel strands give the identity, and reflecting a braid reverses it to give the inverse \cite{braid_theory_ref1,braid_theory_ref4}.

\end{historylist}
\section*{3. Theory}
\subsection*{Core theory}
\paragraph{The braid group and its configuration-space meaning}

The braid group $B_n$ consists of isotopy classes of braids on $n$ strands. Equivalently, let $X$ be a connected manifold of dimension at least two. The unordered configuration space of $n$ distinct points is
\[
 \operatorname{Conf}_n(X)=
 \bigl(X^n\setminus\{x_i=x_j\text{ for some }i\ne j\}\bigr)/S_n.
\]
Its fundamental group is the braid group of $X$. A loop in this configuration space records points moving around one another and returning as an unordered set. If the labels must return to their original positions, the loop lies in the pure braid group.

This interpretation explains why braids compose and why homotopy, rather than a particular drawing, is the correct equivalence relation. For the plane, it gives the classical Artin braid group. In some cases the higher homotopy groups of the configuration space are trivial, making the braid group the main topological information.

\paragraph{Generators and relations}

Let $\sigma_i$ be the braid in which strand $i$ crosses over strand $i+1$, with all other strands straight. The generators $\sigma_1,\ldots,\sigma_{n-1}$ satisfy
\[
 \sigma_i\sigma_j=\sigma_j\sigma_i\quad\text{when }|i-j|\geq2,
\]
because distant crossings do not interact, and
\[
 \sigma_i\sigma_{i+1}\sigma_i=
 \sigma_{i+1}\sigma_i\sigma_{i+1},
\]
because the two three-crossing diagrams are isotopic. The group presentation is
\[
 B_n=\langle\sigma_1,\ldots,\sigma_{n-1}\mid
 \sigma_i\sigma_{i+1}\sigma_i=\sigma_{i+1}\sigma_i\sigma_{i+1},
 \ \sigma_i\sigma_j=\sigma_j\sigma_i\ (|i-j|\geq2)\rangle.
\]
The cubic relation is the braid relation appearing in the Yang--Baxter equation and in Artin groups.

\paragraph{Elementary properties}

$B_1$ is trivial and $B_2\cong\mathbb Z$. The group $B_3$ is infinite and nonabelian and is closely related to the trefoil knot group. Each $B_n$ embeds in $B_{n+1}$ by adding a straight strand, and their increasing union is the infinite braid group $B_\infty$.

Braid groups are torsion-free. For $n\geq3$, they contain a free group on two generators and admit a left-invariant Dehornoy order. Sending every $\sigma_i$ to 1 gives a homomorphism $B_n\to\mathbb Z$ and identifies the abelianization.

\paragraph{Symmetric group and pure braids}

Forget the over-under information but remember where each strand ends. This gives a surjective homomorphism
\[
 B_n\longrightarrow S_n,\qquad \sigma_i\mapsto(i,i+1).
\]
The kernel is the pure braid group $P_n$, in which each strand returns to its own starting point. Pure braids are the fundamental group of the ordered configuration space. They fit into split exact sequences
\[
 1\longrightarrow F_{n-1}\longrightarrow P_n\longrightarrow P_{n-1}\longrightarrow1,
\]
so they can be built as iterated semidirect products of free groups.

\paragraph{Three strands and the modular group}

The group $B_3$ has a distinguished central element. Modding out by its center gives the modular group $\operatorname{PSL}(2,\mathbb Z)$. Thus $B_3$ is a universal central extension of the modular group in the relevant sense. With
\[
 a=\sigma_1\sigma_2\sigma_1,\qquad b=\sigma_1\sigma_2,
\]
the braid relation yields $a^2=b^3$, and the common value is central. This relationship connects braids to modular forms, mapping classes of the punctured torus, and low-dimensional topology \cite{braid_theory_ref2,braid_theory_ref5}.

\paragraph{Closed braids and knots}

Join corresponding top and bottom endpoints of a braid. The result is a link in three-dimensional space, possibly with one component (a knot) or several components. Alexander's theorem says every link is isotopic to the closure of some braid.

Different braids can have the same closure. Markov's theorem gives the exact moves relating braid representatives of the same link: conjugation and stabilization/destabilization. The Jones polynomial was first defined using braid representations and then shown to depend only on the closed link. The braid index is the smallest number of strands needed to represent a link as a closed braid, and it equals the minimum number of Seifert circles in a projection.

\paragraph{Mapping class groups and classification}

Braids can be viewed as mapping classes of a disk with marked points. A braid is a homeomorphism of the punctured disk, considered up to isotopy while fixing the boundary. This makes braid theory part of the study of surfaces.

The Nielsen--Thurston classification divides mapping classes into periodic, reducible, and pseudo-Anosov types. Pseudo-Anosov behavior stretches and folds curves at a definite rate; that rate gives a lower bound for topological entropy. The classification is useful in fluid mixing, where moving rods create braid patterns in space-time trajectories.

\section*{4. Important theorems and results}
\begin{enumerate}[leftmargin=*,itemsep=0.35em]

\item \textbf{Artin presentation.} The braid group is generated by adjacent crossings subject only to far-commutativity and the three-strand braid relations.

\item \textbf{Alexander's theorem.} Every link in three-space is isotopic to the closure of a braid.

\item \textbf{Markov's theorem.} Two braids have isotopic closures precisely when they are related by conjugation and stabilization or destabilization moves.

\item \textbf{Configuration-space theorem.} The braid group is the fundamental group of the configuration space of distinct unordered points.

\item \textbf{Pure-braid extension.} The kernel of $B_n\to S_n$ is $P_n$, and $P_n$ is built from iterated semidirect products of free groups.

\item \textbf{Nielsen--Thurston classification.} Mapping classes, including braids, are periodic, reducible, or pseudo-Anosov, with distinct dynamical behavior.
\end{enumerate}
\noindent\emph{Source for these results:} \cite{braid_theory_ref1}.

\section*{5. Applications}
Braid theory models the motion of distinct strands and records how their crossings compose. The resulting algebra connects configuration spaces with knots, mixing, monodromy, and operators in mathematical physics.
\begin{enumerate}[leftmargin=*,itemsep=0.35em]
\item \textbf{Knot theory.} Closing a braid represents every link. Braid words provide finite data for knot invariants, Markov moves explain when two words describe the same link, and representations lead to polynomial invariants.
\item \textbf{Fluid mixing.} Physical rods, periodic orbits, or almost-invariant sets trace braids. The Nielsen--Thurston type and its entropy estimate how efficiently the stirring produces stretching and mixing.
\item \textbf{Quantum physics.} In two spatial dimensions, exchanging indistinguishable particles can produce braid-group rather than symmetric-group statistics. Anyons use this possibility, and their braid representations are proposed for fault-tolerant quantum computation.
\item \textbf{Algebraic geometry.} Moving roots of a polynomial around one another produces monodromy, represented by a braid. The resulting braid action records how solutions permute when parameters travel around singular values.
\item \textbf{Yang--Baxter equation.} Operators satisfying the braid relation lead to solvable lattice models, quantum groups, and statistical mechanics. The geometric crossing relation becomes an algebraic equation for an operator.
\end{enumerate}


\theoryconnections{braid_theory}{%
\item A braid records paths of moving points, so configuration spaces turn physical motion into loops and group theory records how those loops compose.
\item The braid relation says that sliding one crossing past another does not change the final motion; this is a topological fact expressed as an algebraic equation.
\item Closing a braid produces a link, which is why knot theory and mapping class groups appear naturally \cite{braid_theory_ref1,braid_theory_ref3}.
}{%
\item The same crossing operations model stirring: repeated braid words describe how fluid filaments are stretched and folded, and their growth can estimate mixing.
\item In quantum information and quantum groups, braid generators act as operators whose relations guarantee that different sequences of pairwise exchanges agree.
\item The closure construction also turns algebraic braid data into knot polynomials and other invariants \cite{braid_theory_ref2,braid_theory_ref4,braid_theory_ref5}.
}

\section*{Glossary}
\begin{description}[leftmargin=*,style=nextline,itemsep=0.3em]

\item[\textbf{Braid group.}] The group $B_n$ generated by elementary crossings $\sigma_1,\ldots,\sigma_{n-1}$ subject to far-commutativity and the braid relation; its elements represent ways $n$ strands can wind around one another.

\item[\textbf{Braid relation.}] The identity $\sigma_i\sigma_{i+1}\sigma_i=\sigma_{i+1}\sigma_i\sigma_{i+1}$, which says that sliding one crossing past another produces the same braid; it is the defining cubic relation of the braid group.

\item[\textbf{Pure braid group.}] The kernel of the natural map from the braid group to the symmetric group, consisting of braids in which each strand returns to its original position.

\item[\textbf{Closure of a braid.}] The link obtained by joining each bottom endpoint of a braid to its corresponding top endpoint; Alexander's theorem guarantees that every link arises this way.

\item[\textbf{Configuration space.}] The space of $n$ distinct unordered points in a manifold; its fundamental group is the braid group, giving braids a topological interpretation as loops of moving points.

\item[\textbf{Markov moves.}] The moves (conjugation and stabilization) that relate two different braid words whose closures represent the same link; they are the equivalence criterion for braid representations of knots.

\item[\textbf{Nielsen--Thurston classification.}] A classification of mapping classes---including braids---as periodic, reducible, or pseudo-Anosov, with distinct dynamical behaviors that measure stretching and folding.

\item[\textbf{Modular group.}] The group $\operatorname{PSL}(2,\mathbb Z)$ of $2\times2$ integer matrices with determinant 1 modulo $\pm I$; isomorphic to $B_3$ modulo its center.

\item[\textbf{Symmetric group.}] The group $S_n$ of all permutations of $n$ elements; the image of the braid group under the map forgetting over-under information.

\item[\textbf{Jones polynomial.}] A knot invariant originally defined using braid representations, distinguishing many pairs of knots.

\item[\textbf{Abelianization.}] The quotient of a group by its commutator subgroup, making all generators commute.

\end{description}

\section*{7. Sample Questions}
\emph{Exercise reference:} \cite{braid_theory_ref1}. The cited edition is the source for the exercise set; consult its Exercises section for the official statement and numbering.
\begin{enumerate}[leftmargin=*,itemsep=0.9em]

\item \textbf{Exercise 1.} In the braid group $B_3$, verify the braid relation by computing the permutation in $S_3$ of both sides of $\sigma_1\sigma_2\sigma_1=\sigma_2\sigma_1\sigma_2$.

\emph{Solution.} Each $\sigma_i$ maps to the transposition $(i,i+1)$. Left side: $(1,2)(2,3)(1,2)$. Computing: $(1,2)(2,3)=(1,2,3)$, then $(1,2,3)(1,2)=(1,3)$. Right side: $(2,3)(1,2)(2,3)$. Computing: $(2,3)(1,2)=(1,3,2)$, then $(1,3,2)(2,3)=(1,3)$. Both sides give the transposition $(1,3)$ in $S_3$, consistent with the braid relation.

\item \textbf{Exercise 2.} Compute the image of the braid $\sigma_1\sigma_2\sigma_1^{-1}$ in the symmetric group $S_3$, and determine whether this braid is in the pure braid group $P_3$.

\emph{Solution.} In $S_3$: $\sigma_1\mapsto(1,2)$, $\sigma_2\mapsto(2,3)$, $\sigma_1^{-1}\mapsto(1,2)$. So the image is $(1,2)(2,3)(1,2)=(1,3)$, which is not the identity. Since $P_3$ is the kernel of $B_3\to S_3$, the braid $\sigma_1\sigma_2\sigma_1^{-1}$ is not in $P_3$.

\item \textbf{Exercise 3.} Determine the abelianization of $B_3$ and describe the corresponding homomorphism $B_3\to\mathbb{Z}$.

\emph{Solution.} The abelianization of $B_n$ is $\mathbb{Z}$ for all $n\ge 2$, obtained by sending each generator $\sigma_i$ to $1$. For $B_3=\langle\sigma_1,\sigma_2\mid\sigma_1\sigma_2\sigma_1=\sigma_2\sigma_1\sigma_2\rangle$, the relation becomes $a+b+a=b+a+b$ (writing additively), i.e.\ $a=b$ after cancelling. So both generators map to the same integer. The abelianization is $\mathbb{Z}$, with $\sigma_1,\sigma_2\mapsto 1$.

\item \textbf{Exercise 4.} Find the closure of the braid $\sigma_1\sigma_2^{-1}\in B_3$ and identify the resulting link.

\emph{Solution.} The braid $\sigma_1\sigma_2^{-1}$ has strand 1 crossing over strand 2, then strand 2 crossing under strand 3. Closing this braid connects each top endpoint to its corresponding bottom endpoint. The resulting closed braid has three components (each strand closes to a separate loop), forming the Hopf link union an unlinked unknot, or more precisely the trivial 3-component link $3_1^3$ if no linking occurs. Checking: the permutation is $(1,2)(2,3)=(1,2,3)$, a 3-cycle, so closure gives one component (a knot). Computing more carefully: $\sigma_1$ crosses strands 1,2 and $\sigma_2^{-1}$ crosses strands 2,3 with opposite sign. The closure of this braid is the trefoil knot.

\item \textbf{Exercise 5.} The center $Z(B_3)$ is infinite cyclic, generated by $(\sigma_1\sigma_2)^3$. Verify that $(\sigma_1\sigma_2)^3$ commutes with $\sigma_1$ in $B_3$.

\emph{Solution.} Let $\Delta=(\sigma_1\sigma_2\sigma_1)^2=(\sigma_1\sigma_2)^3$ be the Garside element. We need $\sigma_1\Delta=\Delta\sigma_1$. By the braid relation $\sigma_1\sigma_2\sigma_1=\sigma_2\sigma_1\sigma_2$, so $\Delta=(\sigma_1\sigma_2\sigma_1)^2=\sigma_1\sigma_2\sigma_1\cdot\sigma_1\sigma_2\sigma_1=\sigma_1\sigma_2\sigma_1^2\sigma_2\sigma_1$. Then $\sigma_1\Delta=\sigma_1^2\sigma_2\sigma_1^2\sigma_2\sigma_1$. Similarly $\Delta\sigma_1=\sigma_1\sigma_2\sigma_1^2\sigma_2\sigma_1^2$. Using the braid relation repeatedly shows both expressions are equal. Indeed $\Delta$ is central in $B_3$, a well-known fact.

\end{enumerate}
\section*{8. References}
\begin{thebibliography}{99}
\bibitem{braid_theory_ref1} J. S. Birman, \emph{Braids, Links, and Mapping Class Groups}, Princeton University Press, 1974.
\bibitem{braid_theory_ref2} C. Kassel and V. Turaev, \emph{Braid Groups}, Springer, 2008.
\bibitem{braid_theory_ref3} A. Hatcher, \emph{Algebraic Topology}, Cambridge University Press, 2002.
\bibitem{braid_theory_ref4} V. F. R. Jones, ``A polynomial invariant for knots via von Neumann algebras,'' \emph{Bulletin of the American Mathematical Society}, Vol.~12, No.~1, 1985.
\bibitem{braid_theory_ref5} C. Kassel, \emph{Quantum Groups}, Springer, 1995.
\end{thebibliography}


\clearpage\chapter{Brill--Noether theory}
\section*{1. Overview}
\subsection*{Intuitive}
\paragraph{The problem}

A curve has divisors, and a divisor can support functions or maps to projective space. The surprising question is whether it supports more independent functions than its degree and the curve's genus seem to allow. Brill--Noether theory studies these special divisors, predicts their dimensions on a generic curve, and identifies when a curve has exceptional geometry.



\subsection*{Formal}
Let $C$ be a smooth projective curve of genus $g$. The Brill--Noether locus is
\[
 W^r_d(C)=\{L\in\operatorname{Pic}^d(C):h^0(C,L)\ge r+1\}.
\]
Its expected dimension is governed by $\rho(g,r,d)=g-(r+1)(g-d+r)$. The theory asks when this locus is nonempty, what its dimension is, and how the answer changes for a special curve.
\section*{2. History}
\subsection*{Historical development}
\begin{historylist}
\paragraph{History and the central question}

The theory was introduced by Alexander von Brill and Max Noether in 1874. They studied special divisors on smooth projective curves and discovered that some divisors move in larger linear systems than a naive count predicts. Modern formulations use line bundles, sheaf cohomology, the Picard variety, and moduli spaces.

Let $C$ be a smooth projective curve of genus $g$ over an algebraically closed field. A divisor $D$ determines a line bundle $\mathcal O_C(D)$ and a vector space of sections $H^0(C,\mathcal O_C(D))$. Its dimension is written $h^0(D)=r+1$. A complete linear system of degree $d$ and dimension $r$ is often denoted $g^r_d$. It gives a map to $\mathbb P^r$ when the sections have no common zero \cite{brill_noether_theory_ref1,brill_noether_theory_ref3}.

\end{historylist}
\section*{3. Theory}
\subsection*{Core theory}
\paragraph{Special divisors and Riemann--Roch}

\paragraph{Concrete illustration: low-genus examples}
\emph{Genus 2.} Every genus-2 curve $C$ is hyperelliptic: there exists a degree-2 map $C\to\mathbb P^1$, i.e.\ a $g^1_2$. Take a concrete example, $C\colon y^2=x^6-1$ over $\mathbb C$. A divisor $D=P_1+P_2$ consisting of two Weierstrass points (points where $x^6=1$, $y=0$) has $\deg D=2$ and $h^0(D)=2$, giving a complete $g^1_2$. Riemann--Roch yields $h^0(D)-h^1(D)=2-2+1=1$; since $h^1(D)=h^0(K-D)$ by Serre duality and $\deg K=2$ for genus 2, when $D\ne K$ we get $h^0(D)=2$. Now consider a divisor of degree 3: any effective divisor of degree 3 on a genus-2 curve has $h^0(D)\ge 2$ (by Riemann--Roch: $h^0-h^1=3-2+1=2$, and $h^1\ge 0$), so there is always at least a $g^1_3$, and often a $g^2_3$.

\emph{Genus 3.} A general genus-3 curve is a smooth plane quartic in $\mathbb P^2$. Its canonical series is a $g^2_4$ embedding $C\hookrightarrow\mathbb P^2$. For a point $P\in C$, the divisor $D=K_C-P$ has degree $2g-2-1=3$ and $h^0(D)=g-1=2$ (by Riemann--Roch, since $h^1(D)=h^0(P)=1$), giving a $g^1_3$: a degree-3 map $C\to\mathbb P^1$. The Brill--Noether number for this system is $\rho(3,1,3)=3-2(3-3+1)=-1<0$, so a general genus-3 curve does \emph{not} carry a $g^1_3$ coming from this naive count. However, the computation above shows a $g^1_3$ always exists (just shift $K_C$ by one point). The point is that $\rho$ is negative for the \emph{specified} type; the actual system exists for geometric reasons but is not ``expected.'' By contrast, a degree-2 map $g^1_2$ has $\rho(3,1,2)=3-2(3-2+1)=-1<0$---and indeed a general genus-3 curve has no $g^1_2$; only hyperelliptic genus-3 curves do. This neatly separates the general locus from the hyperelliptic locus in moduli space.

Riemann--Roch says
\[
 h^0(D)-h^1(D)=\deg D-g+1.
\]
The term $h^1(D)$ measures the failure of the number of sections to be the ordinary expected value. A divisor is special when $h^1(D)>0$. By Serre duality, this is equivalent to the existence of a holomorphic differential whose divisor is at least $-D$.

The theory is therefore a refined question about sections. A degree-$d$ divisor may move in a family of dimension greater than expected because the curve has a special shape, an extra symmetry, or a low-degree map to another curve. The same language describes line bundles, maps to projective space, and embeddings.

\paragraph{The Brill--Noether number}

The Brill--Noether number is
\[
 \rho(g,r,d)=g-(r+1)(g-d+r).
\]
It is the expected dimension of the space $G^r_d(C)$ of degree-$d$ linear systems with at least $r+1$ sections. The formula can be remembered from Riemann--Roch:
\[
 \rho=g-h^0(D)h^1(D)
\]
at the expected values.

If $\rho<0$, a general curve should have no such linear series. If $\rho=0$, one expects finitely many. If $\rho>0$, one expects a family of dimension $\rho$. An exceptional curve may violate the generic expectation by possessing extra series.

\paragraph{Main theorems}

Kempf proved that when $\rho\geq0$, $G^r_d(C)$ is nonempty for a smooth curve in the standard range and every component has dimension at least $\rho$. Fulton and Lazarsfeld proved connectedness when $\rho\geq1$. Griffiths and Harris showed that for a general curve the space is reduced and has components of exactly dimension $\rho$, hence is empty when $\rho<0$. Gieseker proved smoothness for a general curve; connectedness then gives irreducibility when $\rho>0$.

These results distinguish generic geometry from special geometry. A hyperelliptic curve, for instance, has a degree-two map to $\mathbb P^1$ and therefore a linear system that a general curve of the same genus may not have.

\paragraph{Maps and examples}

A $g^1_d$ is a pencil of degree $d$: two independent sections give a map $C\to\mathbb P^1$. A hyperelliptic curve has a $g^1_2$. A canonical linear system, associated with holomorphic differentials, maps a nonhyperelliptic curve of genus $g\geq3$ into $\mathbb P^{g-1}$.

For $g=3$, $r=1$, and $d=2$,
\[
 \rho=3-2(3-2+1)=-1.
\]
Thus a general genus-three curve has no degree-two map to $\mathbb P^1$, although hyperelliptic genus-three curves do. The negative expected dimension expresses a generic nonexistence statement rather than an absolute impossibility.

\paragraph{Moduli and the maximal-rank problem}

The moduli space of genus-$g$ curves contains a dense set of general curves. Brill--Noether theory predicts which linear systems occur on that dense set and how large their parameter spaces are. Degenerate or special curves lie in proper loci and can have more sections.

The maximal rank conjecture concerns a general curve $C\subset\mathbb P^r$ embedded by a general linear series. It predicts that restriction maps
\[
 H^0(\mathbb P^r,\mathcal O(m))\to H^0(C,\mathcal O_C(m))
\]
have maximal possible rank. The rank determines the Hilbert function of the embedded curve. Results of Larson and others prove broad cases and interpolation statements through general points, with a small list of exceptional triples.

\section*{4. Important theorems and results}
\begin{enumerate}[leftmargin=*,itemsep=0.35em]

\item \textbf{Riemann--Roch.} The difference $h^0(D)-h^1(D)$ is $\deg D-g+1$, linking sections to degree and genus.

\item \textbf{Brill--Noether dimension theorem.} For a general curve, $G^r_d$ has expected dimension $\rho$ when $\rho\geq0$ and is empty when $\rho<0$.

\item \textbf{Kempf nonemptiness theorem.} Under the standard hypotheses, $\rho\geq0$ implies that the expected linear-series space is nonempty.

\item \textbf{Fulton--Lazarsfeld connectedness.} When $\rho\geq1$, the relevant Brill--Noether space is connected.

\item \textbf{Gieseker smoothness.} For a general curve, the Brill--Noether space is smooth; combined with connectedness this gives irreducibility when $\rho>0$.

\item \textbf{Maximal rank theorem.} Broad cases of the maximal-rank conjecture show that restriction maps for general embedded curves have the expected rank.
\end{enumerate}
\noindent\emph{Source for these results:} \cite{brill_noether_theory_ref1}.

\section*{5. Applications}
Brill--Noether theory predicts when an algebraic curve has a linear series of a prescribed degree and dimension. These predictions control maps to projective space, loci in moduli spaces, and codes obtained by evaluating sections.
\begin{enumerate}[leftmargin=*,itemsep=0.35em]
\item \textbf{Algebraic curves.} Brill--Noether theory predicts which maps and linear systems a general curve can carry.
\item \textbf{Picard varieties.} Divisors and their sections organize families of line bundles and identify special points in the Picard variety.
\item \textbf{Moduli spaces.} The Brill--Noether number predicts the expected dimension of loci of curves with a prescribed linear series.
\item \textbf{Projective embeddings.} A linear series provides an explicit map to projective space, and its sections determine whether that map separates points and tangent directions.
\item \textbf{Coding theory.} Evaluating sections on algebraic curves produces algebraic-geometric error-correcting codes whose parameters are controlled by the curve.
\item \textbf{Algebraic geometry codes from curves.} Goppa's construction of algebraic-geometric (AG) codes uses exactly the linear series studied by Brill--Noether theory. Fix a smooth projective curve $C$ over $\mathbb F_q$, a degree-$d$ divisor $D$ with $V(D)=\emptyset$ (no base points), and a set of $n$ rational points $P_1,\ldots,P_n\in C(\mathbb F_q)\setminus\operatorname{supp}(D)$. The code is $C_\Omega(D,P_1,\ldots,P_n)=\{(f(P_1),\ldots,f(P_n)):f\in H^0(C,\mathcal O(D))\}\subseteq\mathbb F_q^n$. By Riemann--Roch, if $d>2g-2$ then $h^0(D)=d-g+1$, so the code has dimension $k=d-g+1$ and minimum distance $d-n+1$ (a Singleton-type bound). The Tsfasman--Vlăduţ--Zink theorem states that if infinitely many curves over $\mathbb F_q$ satisfy the Drinfeld--Vlăduţ bound $g/q+1\le 1/\sqrt{q}-1$, then AG codes can \emph{beat} the Gilbert--Varshamov bound asymptotically. For example, over $\mathbb F_{127}$, the modular curve $X_0(N)$ for suitable $N$ achieves this, yielding codes with parameters superior to any known random code. These codes find use in deep-space communication (e.g., NASA's consultation with Goppa's codes for compressed telemetry) and in distributed storage systems where the algebraic structure enables efficient decoding via the Berlekamp--Massey--Sakata algorithm.
\end{enumerate}

\theoryconnections{brill_noether_theory}{%
\item Algebraic geometry represents a linear system by sections of a line bundle, while a divisor records the allowed zeros and poles of those sections.
\item Riemann--Roch predicts the expected number of sections; Serre duality explains the correction when the divisor is special.
\item The Brill--Noether number compares the available parameters with the conditions imposed by a desired map to projective space, so it predicts whether a general curve should carry that map \cite{brill_noether_theory_ref1,brill_noether_theory_ref2}.
}{%
\item The theory helps classify embeddings of curves because a collection of sections gives coordinates and hence a map into projective space.
\item Its dimension counts describe loci in moduli space, while the same linear systems control syzygies, interpolation conditions, and evaluation codes.
\item In each case, the geometric question becomes a question about how many independent sections survive the imposed conditions \cite{brill_noether_theory_ref1,brill_noether_theory_ref3}.
}

\section*{Glossary}
\begin{description}[leftmargin=*,style=nextline,itemsep=0.3em]

\item[\textbf{Linear series.}] A family of divisors of fixed degree and dimension on a curve, typically arising as the zero sets of sections of a line bundle; it determines a map to projective space.

\item[\textbf{Brill--Noether number.}] The quantity $\rho(g,r,d)=g-(r+1)(g-d+r)$ that predicts the expected dimension of the space of degree-$d$ linear systems with at least $r+1$ sections on a genus-$g$ curve.

\item[\textbf{Special divisor.}] A divisor whose $h^1$ value is positive, meaning it carries more sections than the Riemann--Roch formula predicts without the correction term.

\item[\textbf{Hyperelliptic curve.}] A curve that admits a degree-two map to the projective line, possessing a $g^1_2$; such curves lie in a special locus of moduli space when the Brill--Noether number is negative.

\item[\textbf{Canonical series.}] The linear system associated with the canonical divisor $K_C$ (the divisor of holomorphic differentials), which embeds a nonhyperelliptic curve of genus $g\geq3$ into projective space.

\item[\textbf{Moduli space.}] The geometric space whose points parametrize isomorphism classes of curves of a given genus; Brill--Noether theory predicts the dimension of special loci within this space.

\item[\textbf{Picard variety.}] The algebraic variety parametrizing isomorphism classes of line bundles of fixed degree on a curve.

\item[\textbf{Genus.}] A topological invariant of a smooth projective curve equal to the number of holes; it controls the expected dimension of linear series.

\item[\textbf{Serre duality.}] A duality identifying $h^1(D)$ with $h^0(K-D)$, where $K$ is the canonical divisor; it explains why a divisor is special.

\item[\textbf{Section.}] A global section of a line bundle, i.e., a regular function assigning a value to each point compatible with the bundle structure.

\end{description}

\section*{7. Sample Questions}
\emph{Exercise reference:} \cite{brill_noether_theory_ref1}. The cited edition is the source for the exercise set; consult its Exercises section for the official statement and numbering.
\begin{enumerate}[leftmargin=*,itemsep=0.9em]

\item \textbf{Exercise 1.} Compute the Brill--Noether number $\rho(g,r,d)$ for a genus-5 curve, a pencil ($r=1$), of degree 4.

\emph{Solution.} $\rho(5,1,4)=5-(1+1)(5-4+1)=5-2\cdot 2=1$. Since $\rho>0$, a general genus-5 curve is expected to carry a 1-dimensional family of $g^1_4$ linear series.

\item \textbf{Exercise 2.} For a smooth curve of genus $g=4$, compute $\rho(g,r,d)$ for $r=2$, $d=4$. Does a general genus-4 curve have a $g^2_4$?

\emph{Solution.} $\rho(4,2,4)=4-(2+1)(4-4+1)=4-3=1$. Since $\rho=1>0$, a general genus-4 curve is expected to have a 1-dimensional family of $g^2_4$ linear series. Indeed, the canonical series of a nonhyperelliptic genus-4 curve is a $g^3_6$, and a general curve has such linear systems by the dimension theorem.

\item \textbf{Exercise 3.} Compute $\rho(g,r,d)$ for a general curve of genus $g=6$, $r=1$, $d=3$. What does $\rho<0$ predict?

\emph{Solution.} $\rho(6,1,3)=6-2(6-3+1)=6-2\cdot 4=-2$. Since $\rho<0$, a general genus-6 curve is expected to have no $g^1_3$ at all. Special (non-general) curves of genus 6 may still carry such a linear series.

\item \textbf{Exercise 4.} For a curve of genus $g=2$, use Riemann--Roch to compute $h^0(D)$ for a divisor $D$ of degree 3.

\emph{Solution.} By Riemann--Roch: $h^0(D)-h^1(D)=\deg D-g+1=3-2+1=2$. By Serre duality, $h^1(D)=h^0(K-D)$ where $\deg K=2g-2=2$, so $\deg(K-D)=2-3=-1<0$. A line bundle of negative degree has no nonzero sections, so $h^0(K-D)=0$, giving $h^1(D)=0$. Therefore $h^0(D)=2$.

\item \textbf{Exercise 5.} Determine which of the following linear series exist on a general curve of genus 7: a $g^1_4$, a $g^1_5$, or a $g^2_6$.

\emph{Solution.} For $g^1_4$: $\rho(7,1,4)=7-2(7-4+1)=7-8=-1<0$, so a general genus-7 curve has no $g^1_4$. For $g^1_5$: $\rho(7,1,5)=7-2(7-5+1)=7-6=1>0$, so $g^1_5$ exists. For $g^2_6$: $\rho(7,2,6)=7-3(7-6+1)=7-6=1>0$, so $g^2_6$ exists. A general genus-7 curve carries a $g^1_5$ and a $g^2_6$ but not a $g^1_4$.

\end{enumerate}
\section*{8. References}
\begin{thebibliography}{99}
\bibitem{brill_noether_theory_ref1} E. Arbarello, M. Cornalba, P. A. Griffiths, and J. Harris, \emph{Geometry of Algebraic Curves, Volume I}, Springer, 1985.
\bibitem{brill_noether_theory_ref2} R. Lazarsfeld, \emph{Positivity in Algebraic Geometry I}, Springer, 2004.
\bibitem{brill_noether_theory_ref3} W. Fulton and R. Lazarsfeld, ``On the connectedness of degeneracy loci and special divisors,'' \emph{Acta Mathematica}, Vol.~146, 1981.
\end{thebibliography}


\clearpage\chapter{Catastrophe theory}
\section*{1. Overview}
\subsection*{Intuitive}
\paragraph{The problem}

Catastrophe theory studies sudden changes produced by smooth changes in control parameters. It is a special part of bifurcation theory for systems whose stable state can be described as a minimum of a smooth potential. The visible jump is not the whole story: in a larger parameter space, the jump lies on a geometric surface of folds, cusps, and higher singularities.



\subsection*{Formal}
Let $V(x;u)$ be a smooth potential, with state variable $x$ and control parameters $u$. Equilibria satisfy $V_x(x;u)=0$; a one-dimensional equilibrium is locally stable when $V_{xx}(x;u)>0$. A fold or other catastrophe occurs when equilibrium and degeneracy equations such as $V_x=V_{xx}=0$ hold. Coordinate changes and additions that do not alter the local equilibrium structure lead to normal forms such as the fold and cusp.
\section*{2. History}
\subsection*{Historical development}
\begin{historylist}
\paragraph{History and scope}

René Thom developed catastrophe theory in the 1960s, and Christopher Zeeman popularized it in the 1970s. Thom was looking for a classification of stable singularities in problems where a system follows a potential downhill. The theory became influential because it promised a small list of local forms that could explain many sudden transitions.

Catastrophe theory is a branch of bifurcation theory and a special case of singularity theory. It asks how the qualitative shape of equilibrium solutions depends on control parameters. Its scope is narrower than the public language of catastrophe suggests: it assumes a suitable potential and local stability structure. Broad claims in biology and social science were criticized in the late 1970s for unsupported assumptions, so modern use requires a clear model, measured parameters, and checks against alternatives \cite{catastrophe_theory_ref1,catastrophe_theory_ref3}.

\end{historylist}
\section*{3. Theory}
\subsection*{Core theory}
\paragraph{Potential dynamics}

Let $x$ be an active state variable and $u$ a vector of control parameters. The model is
\[
 \dot x=-\frac{\partial V(u,x)}{\partial x}.
\]
Equilibria satisfy $V_x=0$. A local minimum of $V$ is stable because small displacements increase the potential and the dynamics move back; a local maximum is unstable.

A catastrophe occurs at a degenerate critical point where
\[
 V_x=0,\qquad V_{xx}=0,
\]
and possibly higher derivatives vanish too. Expand $V$ in a Taylor series around that point. If the degeneracy is structurally stable, smooth changes of coordinates transform the expansion into a standard normal form. The control parameters unfold the singularity and produce surfaces in parameter space separating different numbers or types of equilibria.

\paragraph{Elementary catastrophes}

When there are at most two active variables and at most four control parameters, there are seven generic elementary structures.

\textbf{Fold $A_2$.} A one-variable potential can be written as
\[
 V=x^3+ax.
\]
The equilibrium equation is $3x^2+a=0$. For $a<0$ there is a stable and an unstable equilibrium; they meet at $a=0$ and disappear for $a>0$. This is the saddle-node tipping point.

\textbf{Cusp $A_3$.} Use
\[
 V=\frac{x^4}{4}+\frac{a x^2}{2}+bx.
\]
Equilibria satisfy $x^3+ax+b=0$. The discriminant curve separates a region with three equilibria from a region with one. Inside the cusp, two minima are separated by a maximum. A slow parameter change can follow one minimum until it ends at a fold, then jump to the other.

\textbf{Swallowtail $A_4$.} A standard form is
\[
 V=x^5+ax^3+bx^2+cx.
\]
Two minima and two maxima can meet at one parameter value. Fold surfaces and cusp curves organize the transition.

\textbf{Butterfly $A_5$.} A standard form is
\[
 V=x^6+ax^4+bx^3+cx^2+dx.
\]
Depending on controls, there can be three, two, or one local minima. Fold, cusp, and swallowtail structures meet at the butterfly point.

\paragraph{Umbilic catastrophes}

When two active state variables are needed, corank-two degeneracies produce umbilic catastrophes.

\textbf{Hyperbolic umbilic $D_4^+$.}
\[
 V=x^3+y^3+axy+bx+cy.
\]

\textbf{Elliptic umbilic $D_4^-$.}
\[
 V=\frac{x^3}{3}-xy^2+a(x^2+y^2)+bx+cy.
\]

\textbf{Parabolic umbilic $D_5$.}
\[
 V=x^2y+y^4+ax^2+by^2+cx+dy.
\]

Umbilic catastrophes occur in the focal surfaces of optics and in the geometry of nearly spherical surfaces. Thom suggested that hyperbolic and elliptic umbilics model wave breaking and the creation of hair-like structures; the mathematical classification is independent of whether a particular physical interpretation is correct.

\paragraph{Arnold's notation and Lie-theoretic connection}

Vladimir Arnold organized elementary catastrophes using the ADE notation. $A_0$ is a nonsingular point and $A_1$ a nondegenerate extremum. $A_2,A_3,A_4,A_5$ are fold, cusp, swallowtail, and butterfly; $A_k$ continues the one-variable sequence. $D_4^-$ and $D_4^+$ are elliptic and hyperbolic umbilics, $D_5$ is parabolic umbilic, and $D_k$, $E_6$, $E_7$, and $E_8$ represent further simple singularities.

The labels echo the ADE classification of simple Lie groups and singularities. The connection is not merely a naming coincidence: symmetry and deformation theory organize both the critical-point geometry and the associated algebraic structures \cite{catastrophe_theory_ref4}.

\paragraph{Optics and visible caustics}

Light rays reflected or refracted by a surface form caustics, bright envelopes where many rays concentrate. A fold caustic is the edge of a rainbow or the bright line at the bottom of a swimming pool. Wave effects resolve the geometric singularity into an Airy-function diffraction pattern.

A cusp caustic has a Pearcey-function diffraction pattern. Higher catastrophes such as swallowtails and butterflies can also occur. Their stability explains why similar bright edges and focal shapes survive small changes in the water droplet or reflecting surface.

\section*{4. Important theorems and results}
\begin{enumerate}[leftmargin=*,itemsep=0.35em]

\item \textbf{Structural stability.} A stable singularity persists under small perturbations, up to a smooth change of coordinates. This is why elementary catastrophes can organize nearby behavior.

\item \textbf{Fold classification.} The generic one-parameter creation or destruction of two equilibria is represented by the fold or saddle-node normal form.

\item \textbf{Cusp classification.} The generic two-control unfolding of a one-variable degenerate equilibrium is organized by the cusp surface and its fold curves.

\item \textbf{ADE classification.} Simple stable singularities are organized by the families $A$, $D$, and $E$, including the seven elementary catastrophes in the low-dimensional range.

\item \textbf{Optical caustic theorem.} Generic ray envelopes have stable fold and cusp singularities, with Airy and Pearcey diffraction patterns after wave effects are included.
\end{enumerate}
\noindent\emph{Source for these results:} \cite{catastrophe_theory_ref1}.

\section*{5. Applications}
Catastrophe theory classifies stable local changes in equilibria as control parameters vary. Its normal forms are useful when a model has justified state variables, control parameters, and a potential or equivalent local structure.
\begin{enumerate}[leftmargin=*,itemsep=0.35em]
\item \textbf{Mechanics.} Catastrophe models describe buckling and the loss of a stable equilibrium.
\item \textbf{Optics.} They classify caustics, where many rays focus into a fold or cusp.
\item \textbf{Fluid and wave problems.} Folds and cusps describe envelopes, focusing, and changes in wave patterns.
\item \textbf{Economics, biology, and social science.} Catastrophe models have been proposed for regime changes, but their use is credible only when a potential, state variables, and control parameters are justified.
\end{enumerate}

\noindent\emph{Limitations.} The theory depends on calculus, differential equations, bifurcation theory, singularity theory, Taylor expansion, topology, and smooth changes of variables. It supports local dynamical classification, optics, nonlinear mechanics, phase transitions, and model reduction. It does not by itself predict the timing of a real landslide or prove that a social system has only a few states. Unmodeled variables, noise, hysteresis, and non-equilibrium dynamics can invalidate a potential description.


\theoryconnections{catastrophe_theory}{%
\item Calculus finds equilibria by setting the derivative of a potential or dynamical law to zero, and Taylor expansion removes unimportant higher-order details near a degenerate equilibrium.
\item Singularity theory classifies the resulting local shapes, while bifurcation theory tracks how the shape changes as control parameters move.
\item Topology explains why folds, cusps, and umbilics are stable patterns under small perturbations, not accidental drawings \cite{catastrophe_theory_ref1,catastrophe_theory_ref4}.
}{%
\item In optics, a caustic is the envelope of many rays and can form a fold or cusp; the normal form predicts the bright pattern without tracing every ray individually.
\item In mechanics and fluid models, the same normal forms describe a sudden change between equilibria.
\item These conclusions require justified state variables, control parameters, and potential structure; a visually similar graph alone is not evidence for a catastrophe model \cite{catastrophe_theory_ref2,catastrophe_theory_ref3}.
}

\section*{Glossary}
\begin{description}[leftmargin=*,style=nextline,itemsep=0.3em]

\item[\textbf{Catastrophe.}] A sudden qualitative change in the stable equilibrium of a system caused by a smooth variation of control parameters, corresponding to a degenerate critical point of a potential.

\item[\textbf{Fold.}] The simplest elementary catastrophe, in which a stable and an unstable equilibrium approach each other, merge, and disappear as a single parameter crosses a threshold.

\item[\textbf{Cusp.}] A two-control catastrophe in which a smooth parameter change can cause a sudden jump between two distant stable states; the cusp surface separates regions with one and three equilibria.

\item[\textbf{Normal form.}] A standard simplified potential obtained by coordinate changes that preserves the qualitative behavior near a degenerate equilibrium, such as the fold $x^3+ax$ or the cusp $x^4/4+ax^2/2+bx$.

\item[\textbf{Potential.}] A smooth function $V(x;u)$ whose local minima represent stable equilibria; the dynamics drive the state toward minima, and catastrophes occur where minima degenerate.

\item[\textbf{Umbilic catastrophe.}] An elementary catastrophe involving two active state variables, producing the hyperbolic, elliptic, and parabolic umbilic normal forms.

\item[\textbf{Caustic.}] The envelope of light rays reflected or refracted by a surface; catastrophe theory classifies the generic bright patterns (fold and cusp caustics) that appear in optics.

\item[\textbf{Elementary catastrophe.}] One of the seven structurally stable singularities (fold, cusp, swallowtail, butterfly, and three umbilics) with at most two state variables and four control parameters.

\item[\textbf{Control parameter.}] A smoothly varying external variable that can be adjusted independently of the state; catastrophes occur as parameters cross thresholds.

\item[\textbf{State variable.}] The internal variable whose value is determined by minimizing the potential; it describes the observable condition.

\item[\textbf{Structural stability.}] The property that a singularity persists under small perturbations up to coordinate change.

\end{description}

\section*{7. Sample Questions}
\emph{Exercise reference:} \cite{catastrophe_theory_ref1}. The cited edition is the source for the exercise set; consult its Exercises section for the official statement and numbering.
\begin{enumerate}[leftmargin=*,itemsep=0.9em]

\item \textbf{Exercise 1.} For the fold potential $V(x)=x^3+ax$, find the equilibria for $a=-3$, $a=0$, and $a=3$, and determine which are stable.

\emph{Solution.} Equilibria satisfy $V'(x)=3x^2+a=0$. For $a=-3$: $3x^2=3$, so $x=\pm 1$. At $x=1$: $V''(1)=6>0$ (stable). At $x=-1$: $V''(-1)=-6<0$ (unstable). For $a=0$: $x=0$ with $V''(0)=0$ (degenerate). For $a=3$: $3x^2+3=0$ has no real solutions (no equilibria).

\item \textbf{Exercise 2.} For the cusp potential $V(x)=\frac{x^4}{4}+\frac{ax^2}{2}+bx$, find the equilibrium curve and determine the number of equilibria at $(a,b)=(-3,0)$, $(a,b)=(-3,2)$, and $(a,b)=(-3,-2)$.

\emph{Solution.} Equilibria satisfy $x^3+ax+b=0$. At $(a,b)=(-3,0)$: $x^3-3x=0$, so $x(x^2-3)=0$, giving three equilibria: $x=0,\pm\sqrt{3}$. At $(a,b)=(-3,2)$: $x^3-3x+2=0$, factoring: $(x-1)^2(x+2)=0$, giving $x=1$ (double) and $x=-2$: two equilibria. At $(a,b)=(-3,-2)$: $x^3-3x-2=0$, factoring: $(x+1)^2(x-2)=0$, giving $x=-1$ (double) and $x=2$: two equilibria.

\item \textbf{Exercise 3.} The cusp discriminant curve is given parametrically by $a=-3t^2$, $b=2t^3$. Find the two points on the discriminant corresponding to $t=1$ and $t=-1$, and verify they lie on $4a^3+27b^2=0$.

\emph{Solution.} At $t=1$: $a=-3$, $b=2$. Check: $4(-3)^3+27(2)^2=4(-27)+27(4)=-108+108=0$. At $t=-1$: $a=-3$, $b=-2$. Check: $4(-3)^3+27(-2)^2=-108+108=0$. Both points satisfy the cusp equation $4a^3+27b^2=0$.

\item \textbf{Exercise 4.} For the swallowtail potential $V(x)=x^5+ax^3+bx^2+cx$, compute $V'(x)$ and $V''(x)$, and find the parameter values at which two equilibria merge along the line $b=0$, $c=0$.

\emph{Solution.} $V'(x)=5x^4+3ax^2+2bx+c$ and $V''(x)=20x^3+6ax+2b$. Setting $b=0,c=0$: $V'(x)=x^2(5x^2+3a)$ and $V''(x)=2x(10x^2+3a)$. Equilibria are $x=0$ and $x=\pm\sqrt{-3a/5}$ (for $a<0$). At $x=0$: $V''(0)=0$ always. The degeneracy $V'''(0)=6a$ also vanishes when $a=0$. So at $(a,b,c)=(0,0,0)$, the origin is a swallowtail point with $V'=V''=V'''=0$.

\item \textbf{Exercise 5.} For the hyperbolic umbilic $V(x,y)=x^3+y^3+axy+bx+cy$, find the equilibrium equations by setting $V_x=0$ and $V_y=0$, and determine the number of equilibria at $(a,b,c)=(-3,2,2)$.

\emph{Solution.} $V_x=3x^2+ay+b=0$ and $V_y=3y^2+ax+c=0$. At $(a,b,c)=(-3,2,2)$: $3x^2-3y+2=0$ and $3y^2-3x+2=0$. Subtracting: $3(x^2-y^2)-3(y-x)=0$, so $3(x-y)(x+y)+3(x-y)=0$, giving $(x-y)(3(x+y)+3)=0$. If $x=y$: $3x^2-3x+2=0$ with discriminant $9-24<0$: no real solutions. If $x+y=-1$: substituting $y=-1-x$ into the first equation: $3x^2+3(1+x)+2=0$, i.e.\ $3x^2+3x+5=0$ with discriminant $9-60<0$: no real solutions. Hence there are no real equilibria at these parameter values.

\end{enumerate}
\section*{8. References}
\begin{thebibliography}{99}
\bibitem{catastrophe_theory_ref1} V. I. Arnold, \emph{Bifurcation Theory and Catastrophe Theory}, Springer, 1992.
\bibitem{catastrophe_theory_ref2} T. Poston and I. Stewart, \emph{Catastrophe Theory and Its Applications}, Pitman, 1978.
\bibitem{catastrophe_theory_ref3} R. S. Zahler and H. J. Sussmann, ``Claims and accomplishments of applied catastrophe theory,'' \emph{Nature}, Vol.~279, 1979.
\bibitem{catastrophe_theory_ref4} V. I. Arnold, S. M. Gusein-Zade, and A. N. Varchenko, \emph{Singularities of Differentiable Maps}, Birkhäuser, 1985.
\end{thebibliography}


\clearpage\chapter{Category theory}
\section*{1. Overview}
\subsection*{Intuitive}
\paragraph{The problem}

Mathematicians repeatedly build products, quotients, completions, duals, and solution spaces in different subjects. Category theory asks what these constructions have in common. It studies objects together with the maps between them, and it emphasizes relationships that remain meaningful when the internal details change.



\subsection*{Formal}
A category $\mathcal C$ consists of objects, morphisms $f:A\to B$, identity morphisms $1_A$, and a composition law that is associative and has $1_A$ as an identity. A functor maps objects and morphisms between categories while preserving identities and composition. A natural transformation compares two functors by component morphisms that make the corresponding squares commute.
\section*{2. History}
\subsection*{Historical development}
\begin{historylist}
\paragraph{History and the motivating idea}

Samuel Eilenberg and Saunders Mac Lane introduced categories, functors, and natural transformations in the 1940s while formalizing algebraic topology. They wanted a language in which topological spaces could be related to algebraic invariants and in which the important maps between constructions were visible. The theory later became central to homological algebra, algebraic geometry, logic, and computer science.

The category-theoretic viewpoint continues Emmy Noether's emphasis on maps preserving structure. It is not a claim that objects have no internal content. It is a method for identifying the part of a construction that is independent of a particular representation \cite{category_theory_ref1,category_theory_ref2}.

\end{historylist}
\section*{3. Theory}
\subsection*{Core theory}
\paragraph{Categories, objects, and morphisms}
\bookfigure{category_commutative}{A commutative diagram records compatible ways of composing maps.}

A category $\mathcal C$ has objects and morphisms. A morphism $f:A\to B$ has source $A$ and target $B$. Morphisms can be composed when the target of one is the source of the next. Composition is associative:
\[
 h\circ(g\circ f)=(h\circ g)\circ f.
\]
Every object $A$ has an identity morphism $1_A:A\to A$ satisfying
\[
 1_B\circ f=f=f\circ1_A.
\]

The category $\mathbf{Set}$ has sets as objects and functions as morphisms. Groups with homomorphisms, topological spaces with continuous maps, vector spaces with linear maps, and smooth manifolds with smooth maps are other examples. A monoid is a category with one object: its elements are the endomorphisms of that object.

A morphism may be a monomorphism if it can be cancelled on the left, an epimorphism if it can be cancelled on the right, and an isomorphism if it has a two-sided inverse. In sets these resemble injective, surjective, and bijective functions, but in general categories the analogies need not coincide exactly. Commutative diagrams show equations between composites by displaying paths with the same endpoints.

\paragraph{Universal properties}

Many important objects are best defined by how maps enter or leave them. The product $A\times B$ comes with projections $p_A$ and $p_B$. For every object $X$ and maps $f:X\to A$, $g:X\to B$, there is a unique map
\[
 \langle f,g\rangle:X\to A\times B
\]
whose composites with the projections are $f$ and $g$.

This property, not a particular set of ordered pairs, defines a product in any category where it exists. The dual construction is a coproduct. Equalizers, coequalizers, pullbacks, pushouts, initial objects, and terminal objects are other universal constructions. Limits combine compatible diagrams; colimits glue diagrams together.

Universal properties explain why apparently different constructions are the same. If two objects satisfy the same universal property, there is a unique isomorphism between them. The category must be stated: a product in sets, groups, or topological spaces has different internal structure, even though the mapping property has the same form.

\paragraph{Functors}

A functor $F:\mathcal C\to\mathcal D$ maps objects to objects and morphisms to morphisms while preserving identities and composition:
\[
 F(1_A)=1_{F(A)},\qquad F(g\circ f)=F(g)\circ F(f).
\]
Examples include the forgetful functor from groups to sets, the fundamental-group functor from pointed spaces to groups, and the vector-space dual functor, which is contravariant and reverses arrows.

Functors let one transport a construction or invariant between subjects. A covariant functor preserves arrow direction; a contravariant functor is a covariant functor from the opposite category $\mathcal C^{op}$.

\paragraph{Natural transformations}

Given functors $F,G:\mathcal C\to\mathcal D$, a natural transformation $\eta:F\Rightarrow G$ assigns a morphism
\[
 \eta_A:F(A)\to G(A)
\]
to every object so that for every $f:A\to B$,
\[
 \eta_B\circ F(f)=G(f)\circ\eta_A.
\]
The square commutes. Naturality means the comparison is compatible with every map, not chosen separately for each object. If every $\eta_A$ is an isomorphism, the functors are naturally isomorphic.

This distinction is important in algebra and topology. Two vector spaces may have isomorphic underlying dimensions, but a natural isomorphism says the identification is canonical and behaves correctly under all linear maps.

\paragraph{Equivalence and duality}

An isomorphism of categories is a strict reversible correspondence. An equivalence is weaker: functors in both directions are inverse up to natural isomorphism. Equivalent categories have the same category-level mathematics even when their objects are not literally identical.

Every definition and theorem has a dual obtained by reversing all arrows. Products become coproducts, limits become colimits, monomorphisms become epimorphisms, and left adjoints become right adjoints in the dual formulation. This arrow-reversal often reveals a second theorem at no extra cost.

\paragraph{Yoneda and representability}

The Yoneda lemma says that an object is determined, up to canonical isomorphism, by the functor of maps into it. For an object $A$, the representable functor $\mathcal C(-,A)$ records every way other objects map to $A$. A natural transformation between representable functors comes from a unique morphism.

Yoneda formalizes the slogan that an object is known by its relationships. It is the foundation for representable functors, moduli problems, and many universal-property proofs. It also warns against identifying objects merely by one selected set of coordinates.

\paragraph{Adjunctions and other constructions}

An adjunction pairs functors $F:\mathcal C\to\mathcal D$ and $G:\mathcal D\to\mathcal C$ so that maps $F(A)\to B$ correspond naturally to maps $A\to G(B)$. The free-group construction is left adjoint to the forgetful functor from groups to sets: maps from a set into the underlying set of a group extend uniquely to homomorphisms from the free group.

Adjunctions explain universal properties and are often the most efficient way to recognize a construction. Functor categories have functors as objects and natural transformations as morphisms. Monads package an adjunction into a reusable algebraic operation and appear in programming-language semantics.

\paragraph{Higher-dimensional categories}

In an ordinary category, morphisms relate objects. A 2-category also has morphisms between morphisms, with horizontal and vertical composition and an interchange law. Bicategories weaken strict associativity and retain associativity up to coherent isomorphism. Monoidal categories are categories with a tensor-like product and a unit. The hierarchy continues to 3-categories, $n$-categories, and $\omega$-categories.

Higher categories are useful when transformations between maps are themselves important, as in homotopy theory, derived geometry, topological quantum field theory, and higher-dimensional algebra.

\section*{4. Important theorems and results}
\begin{enumerate}[leftmargin=*,itemsep=0.35em]

\item \textbf{Yoneda lemma.} Natural transformations from a representable functor are in one-to-one correspondence with morphisms in the category.

\item \textbf{Universal-property uniqueness.} Two objects satisfying the same universal property are uniquely isomorphic.

\item \textbf{Adjoint functor correspondence.} An adjunction gives natural bijections between two families of hom-sets and explains free constructions, products, and many completions.

\item \textbf{Equivalence of categories.} Functors that are inverse up to natural isomorphism preserve category-level structure and theorems.

\item \textbf{Functoriality of invariants.} A structure-preserving map induces a compatible map on an invariant; composition and identities are preserved automatically.

\item \textbf{Duality principle.} Reversing all arrows in a categorical theorem gives a dual theorem in the opposite category.
\end{enumerate}
\noindent\emph{Source for these results:} \cite{category_theory_ref1}.

\section*{5. Applications}
Category theory is useful when the important information is how constructions map to one another and compose. Functors and universal properties preserve these relationships across algebra, geometry, logic, computation, and physics.
\begin{enumerate}[leftmargin=*,itemsep=0.35em]
\item \textbf{Algebraic topology.} Homology and homotopy are functors. A continuous map induces maps on algebraic invariants, and naturality lets one compare those induced maps.
\item \textbf{Algebraic geometry.} Schemes, sheaves, fiber products, and derived functors are naturally categorical. The language makes local-to-global constructions and change of base systematic.
\item \textbf{Logic and foundations.} Toposes provide categorical models of set-like mathematics and constructive logic. Categorical logic connects type theory, intuitionistic logic, and semantics.
\item \textbf{Computer science.} Cartesian closed categories model typed lambda calculus; monads organize effects in functional programming; domain theory models recursive computation. A compiler's correctness proof can be expressed as a commuting diagram.
\item \textbf{Physics and networks.} Monoidal categories describe processes that can be composed in time and placed side by side. They organize tensor networks, quantum circuits, open systems, databases, and Feynman-diagram composition.
\end{enumerate}

\textbf{Music and applied structures.} Applied category theory has been used to describe structured data, databases, dynamical systems, and compositional models in music theory. The common idea is to preserve relationships while changing representation \cite{category_theory_ref3,category_theory_ref4}.


\theoryconnections{category_theory}{%
\item Sets provide objects and functions provide arrows; category theory keeps only the composition rule and identity arrows that explain how constructions interact.
\item This abstraction lets one compare groups, spaces, rings, and other subjects using the same diagram.
\item A functor translates one category to another while preserving composition, and a natural transformation compares two such translations without choosing unrelated coordinates at every object \cite{category_theory_ref1,category_theory_ref2}.
}{%
\item Category theory helps algebraic topology by turning a space into algebraic data through functors, and helps algebraic geometry by organizing rings, spectra, and sheaves with universal maps.
\item In programming and databases, a compositional diagram records how components or queries can be assembled.
\item The gain is not a new numerical answer by itself; it is a proof that different constructions are compatible because they satisfy the same universal pattern \cite{category_theory_ref1,category_theory_ref3}.
}

\section*{Glossary}
\begin{description}[leftmargin=*,style=nextline,itemsep=0.3em]

\item[\textbf{Category.}] A collection of objects together with morphisms (arrows) between them that can be composed associatively, with each object having an identity morphism.

\item[\textbf{Functor.}] A map between categories that sends objects to objects and morphisms to morphisms while preserving identities and composition; it transports constructions from one setting to another.

\item[\textbf{Natural transformation.}] A family of morphisms between two functors that is compatible with every map in the category, expressing a canonical way to compare two constructions.

\item[\textbf{Universal property.}] A property that defines an object (such as a product or coproduct) solely in terms of how other objects and maps relate to it, guaranteeing uniqueness up to isomorphism.

\item[\textbf{Adjoint functors.}] A pair of functors $F\dashv G$ in which maps from $F(A)$ to $B$ correspond naturally to maps from $A$ to $G(B)$; adjunctions encode free constructions, limits, and many canonical completions.

\item[\textbf{Yoneda lemma.}] The result that an object is determined up to isomorphism by the functor of maps into it, formalizing the idea that an object is fully known by its relationships.

\item[\textbf{Limit and colimit.}] Universal constructions that, respectively, pick a canonical ``cone'' over a diagram (limit) or glue objects together along shared data (colimit).

\item[\textbf{Morphism.}] An arrow from one object to another in a category that can be composed associatively; the structure-preserving maps of each category.

\item[\textbf{Isomorphism.}] A morphism with a two-sided inverse; two objects are isomorphic when structurally identical in the category.

\item[\textbf{Duality.}] The principle that every definition and theorem has a dual obtained by reversing all arrows.

\item[\textbf{Natural isomorphism.}] A natural transformation in which every component morphism is an isomorphism.

\end{description}

\section*{7. Sample Questions}
\emph{Exercise reference:} \cite{category_theory_ref1}. The cited edition is the source for the exercise set; consult its Exercises section for the official statement and numbering.
\begin{enumerate}[leftmargin=*,itemsep=0.9em]

\item \textbf{Exercise 1.} Let $\mathbf{Grp}$ be the category of groups. Define a functor $F:\mathbf{Grp}\to\mathbf{Set}$ that sends each group $G$ to its center $Z(G)$, and each homomorphism $\varphi:G\to H$ to the restriction $\varphi|_{Z(G)}$. Is $F$ well-defined? If not, explain why.

\emph{Solution.} $F$ is well-defined. If $z\in Z(G)$ and $\varphi:G\to H$ is a homomorphism, then for any $h=\varphi(g)\in H$, we have $\varphi(z)\varphi(g)=\varphi(zg)=\varphi(gz)=\varphi(g)\varphi(z)$, so $\varphi(z)\in Z(\operatorname{im}\varphi)\subseteq Z(H)$. Thus $\varphi$ maps $Z(G)$ into $Z(H)$, and the restriction $\varphi|_{Z(G)}$ is a function $Z(G)\to Z(H)$. Functoriality: $F(\mathrm{id}_G)=\mathrm{id}_{Z(G)}$ and $F(\psi\circ\varphi)=(\psi\circ\varphi)|_{Z(G)}=\psi|_{\varphi(Z(G))}\circ\varphi|_{Z(G)}=F(\psi)\circ F(\varphi)$, so $F$ is a well-defined functor.

\item \textbf{Exercise 2.} In the category $\mathbf{Set}$, compute the product $\{a,b\}\times\{1,2\}$ and verify it satisfies the universal property of the product.

\emph{Solution.} The Cartesian product is $\{a,b\}\times\{1,2\}=\{(a,1),(a,2),(b,1),(b,2)\}$. The projections are $p_1(a,1)=p_1(a,2)=a$, $p_1(b,1)=p_1(b,2)=b$ and $p_2(a,1)=p_2(b,1)=1$, $p_2(a,2)=p_2(b,2)=2$. For any set $X$ and functions $f:X\to\{a,b\}$, $g:X\to\{1,2\}$, define $\langle f,g\rangle(x)=(f(x),g(x))$. Then $p_1\circ\langle f,g\rangle=f$ and $p_2\circ\langle f,g\rangle=g$, and $\langle f,g\rangle$ is the unique such function since ordered pairs are determined by components.

\item \textbf{Exercise 3.} Let $F,G:\mathbf{Set}\to\mathbf{Set}$ be the functors $F(X)=X\times\{0,1\}$ (product with a two-element set) and $G(X)=X\sqcup X$ (disjoint union of two copies of $X$). Construct a natural isomorphism $\eta:F\Rightarrow G$.

\emph{Solution.} Define $\eta_X:F(X)\to G(X)$ by $\eta_X(x,0)=(x,\text{left})$ and $\eta_X(x,1)=(x,\text{right})$. This is a bijection for every $X$: its inverse sends $(x,\text{left})\mapsto(x,0)$ and $(x,\text{right})\mapsto(x,1)$. Naturality: for any $f:X\to Y$, the diagram commutes because $G(f)\circ\eta_X$ and $\eta_Y\circ F(f)$ both send $(x,i)$ to $(f(x),i)$. Hence $\eta$ is a natural isomorphism.

\item \textbf{Exercise 4.} Prove that the forgetful functor $U:\mathbf{Grp}\to\mathbf{Set}$ has a left adjoint (the free-group functor).

\emph{Solution.} For any set $S$ and group $G$, define $\Phi:\mathrm{Hom}_{\mathbf{Grp}}(F(S),G)\to\mathrm{Hom}_{\mathbf{Set}}(S,U(G))$ by $\Phi(\varphi)=\varphi|_S$ (restricting a group homomorphism from $F(S)$ to the generating set $S$). This is a bijection: given a set map $f:S\to U(G)$, there is a unique group homomorphism $\hat{f}:F(S)\to G$ extending $f$ (by the universal property of free groups), and $\Phi(\hat{f})=f$. Injectivity: if $\varphi_1|_S=\varphi_2|_S$ then $\varphi_1=\varphi_2$ since $S$ generates $F(S)$. Surjectivity: every set map extends. Naturality in $G$ follows from the compatibility of restriction with post-composition by group homomorphisms.

\item \textbf{Exercise 5.} Compute the coproduct $\mathbb{Z}/2\mathbb{Z}\sqcup\mathbb{Z}/3\mathbb{Z}$ in $\mathbf{Grp}$.

\emph{Solution.} The coproduct in $\mathbf{Grp}$ is the free product $\mathbb{Z}/2\mathbb{Z}*\mathbb{Z}/3\mathbb{Z}=\langle a,b\mid a^2=1,\,b^3=1\rangle$. This is the infinite dihedral-like group $\mathrm{PSL}(2,\mathbb{Z})\cong\mathbb{Z}/2\mathbb{Z}*\mathbb{Z}/3\mathbb{Z}$, which is infinite. The universal property holds: any pair of homomorphisms $\varphi:\mathbb{Z}/2\mathbb{Z}\to H$ and $\psi:\mathbb{Z}/3\mathbb{Z}\to H$ extends uniquely to a homomorphism $\mathbb{Z}/2\mathbb{Z}*\mathbb{Z}/3\mathbb{Z}\to H$.

\end{enumerate}
\section*{8. References}
\begin{thebibliography}{99}
\bibitem{category_theory_ref1} S. Mac Lane, \emph{Categories for the Working Mathematician}, 2nd edition, Springer, 1998.
\bibitem{category_theory_ref2} S. Awodey, \emph{Category Theory}, 2nd edition, Oxford University Press, 2010.
\bibitem{category_theory_ref3} E. Riehl, \emph{Category Theory in Context}, Dover, 2017.
\bibitem{category_theory_ref4} F. W. Lawvere and S. H. Schanuel, \emph{Conceptual Mathematics}, 2nd edition, Cambridge University Press, 2009.
\end{thebibliography}


\theory{Chaos theory}{How can a completely deterministic system become practically unpredictable?}{Chaotic systems amplify tiny differences in starting conditions. Equations can be simple while long-term prediction fails; nevertheless, attractors, dimensions, and statistical patterns can remain measurable.}{Weather, turbulence, pendulums, heart rhythms, and population studies.}{Iteration, sensitivity, Lyapunov growth, and attractors turn deterministic equations into measurable predictions and finite forecast limits.}
\paragraph{Concrete examples and uses}
Consider the logistic map $x_{n+1}=4x_n(1-x_n)$. Nearby initial values separate because the map repeatedly stretches and folds the interval. This matters in practice because finite precision eventually repeats an orbit; chaotic models describe deterministic unpredictability in weather and ecology.

\paragraph{Historical development}
Henri Poincare discovered the essential difficulty while studying the
three-body problem. The equations are deterministic, but the paths of three
gravitating bodies can wind around one another so intricately that a small
uncertainty in the starting positions changes the later path. Poincare's
method of recording where an orbit crosses a surface became the Poincare
section, a basic tool for visualising complicated dynamics
\cite{chaos_theory_ref3}.

In 1963 Edward Lorenz found similar behaviour in a simplified model of
atmospheric convection. Rounding one number in a restarted computer run
changed the later weather-like pattern. The lesson was not that a butterfly
directly controls one tornado; it was that imperfect knowledge of the
present can impose a finite horizon on prediction
\cite{chaos_theory_ref4}. In the 1970s Robert May showed that a one-line
population model can move from stable growth to cycles and chaos. Mitchell
Feigenbaum then found a common numerical pattern in period-doubling routes
to chaos \cite{chaos_theory_ref1}.

\end{historylist}
\section*{3. Theory}
\subsection*{Core theory}
\paragraph{What chaos means}
A discrete dynamical system repeatedly applies a rule
$x_{n+1}=f(x_n)$. A continuous-time system is often written as
$d\mathbf{x}/dt=F(\mathbf{x})$. The state space contains all possible
present states; an orbit is the sequence generated from one starting state.
An equilibrium stays fixed, a periodic orbit repeats, and an attractor is a
set toward which many initial states move.

There is no single definition that covers every use of the word chaos. A
common definition for a continuous map on a suitable metric space requires
topological transitivity, dense periodic points, and sensitive dependence on
initial conditions \cite{chaos_theory_ref2}. Transitivity means that the
motion can carry points from one open region into another. Dense periodic
points mean that regular cycles occur arbitrarily close to other states.
Sensitivity means that arbitrarily close starting states can later become
noticeably separated.

Sensitivity by itself is not enough. The map $x\mapsto2x$ separates nearby
points but usually sends them simply to infinity. Chaotic systems normally
combine stretching with folding or another mechanism that keeps motion in a
bounded region. The apparent disorder then comes from deterministic feedback,
not from a random instruction.

\paragraph{The logistic map and the route to chaos}
\bookfigure{dynamics_chaos_bifurcation}{The logistic map shows how repeated iteration and parameter change can produce complex behavior.}
Consider
\[
x_{n+1}=r x_n(1-x_n),\qquad 0\leq x_n\leq1.
\]
The variable can represent a scaled population and $r$ the strength of
reproduction. The fixed points solve $rx(1-x)=x$, so they are $0$ and
$x_*=1-1/r$. For $0<r<1$ the population tends to zero. For $1<r<3$ the
nonzero fixed point is stable: a small error around it shrinks after
iteration.

At $r=3$ the fixed point loses stability and a stable period-two cycle
appears. As $r$ increases, cycles of periods four, eight, and higher powers
of two appear. These period-doubling values accumulate near
$r=3.56995$. Beyond that point the map contains chaotic ranges as well as
periodic windows. At $r=4$, two initial values differing in a few decimal
places can agree for several steps and then diverge dramatically.

The map is valuable because it demonstrates the whole mechanism in a model
small enough to calculate by hand: feedback, loss of stability, repeated
period doubling, sensitivity, and a statistical long-run pattern. The model
does not claim that every population is logistic; it isolates one mechanism
that a more realistic model can test \cite{chaos_theory_ref1}.

\paragraph{The Lorenz system and strange attractors}
Lorenz studied
\[
\dot{x}=\sigma(y-x),\qquad
\dot{y}=x(\rho-z)-y,\qquad
\dot{z}=xy-\beta z.
\]
For $\sigma=10$, $\rho=28$, and $\beta=8/3$, trajectories approach a
bounded butterfly-shaped set called the Lorenz attractor. The two wings are
regions of state space in which the model circulates before switching to the
other region. A trajectory never settles to one fixed point or one simple
cycle, but it remains confined.

A strange attractor combines boundedness with complicated geometry and
sensitive motion. A Poincare section records the orbit whenever it crosses a
chosen surface, turning a continuous flow into a discrete map. It often
reveals stretching and folding more clearly than a three-dimensional plot.
The attractor can have a fractal dimension, so geometric complexity and
predictability can be measured separately \cite{chaos_theory_ref3}.

\paragraph{Lyapunov exponents and predictability}
If two nearby trajectories initially differ by $\delta_0$, their separation
may behave for a while like
\[
|\delta(t)|\approx |\delta_0|e^{\lambda t}.
\]
The number $\lambda$ is a Lyapunov exponent. A positive largest exponent
means exponential growth of small errors in the direction being measured.
The reciprocal gives a rough Lyapunov time: after several such times, the
initial error may be too large for a detailed forecast.

This is a diagnostic rather than a magic certificate. A numerical transient
can look chaotic before settling, rounding can create false growth, and
different directions can have different exponents. A bounded system with a
positive largest exponent is strong evidence of chaos, but it should be
checked with phase portraits, recurrence tests, or a Poincare section
\cite{chaos_theory_ref1}.

\section*{4. Important theorems and results}
\begin{enumerate}[leftmargin=*,itemsep=0.35em]
\item \textbf{Poincare--Bendixson theorem.} A sufficiently regular continuous-time
system confined to a compact region of the plane cannot have a strange
attractor. Its limit sets are equilibria, periodic orbits, or connecting
trajectories. A continuous-time strange attractor therefore requires at
least three effective state variables, although a one-dimensional discrete
map can already be chaotic.

\item \textbf{Period-three theorem.} Li and Yorke proved that a continuous map of
an interval with a point of period three has periodic points of every
positive period and an uncountable set of nonperiodic chaotic orbits. This
explains why the simple logistic map can contain much richer behaviour than
its formula suggests \cite{chaos_theory_ref2}.

\item \textbf{Feigenbaum universality.} In many one-parameter families, ratios of
successive parameter intervals in a period-doubling cascade approach about
$4.669$. The same constant appearing in different families shows that some
routes to chaos depend on broad structural features rather than every detail
of the formula \cite{chaos_theory_ref1}.
\end{enumerate}

\paragraph{Concrete uses}
In weather forecasting, the atmosphere is modelled by deterministic
equations but its pressure, temperature, and wind are measured imperfectly.
Forecast centres run the model from many slightly different starting states.
If the runs spread rapidly, the service reports a probability of rain rather
than one falsely precise long-range path. Chaos theory explains why a short
forecast can be useful while a detailed forecast eventually loses skill.

In mechanical engineering, a double pendulum or a vibrating rotor can have
only a few moving parts and still show chaotic motion. Engineers calculate
Poincare sections and Lyapunov exponents to locate operating energies where a
small vibration grows rapidly. That information helps them avoid unstable
regimes or design feedback that suppresses the growth.

In ecology, the logistic map shows that a population can fluctuate
irregularly without an external random shock. If an estimated reproduction
parameter is near a period-doubling threshold, a small change in harvesting
or climate may change a stable cycle into a complicated fluctuation. The
model is useful as a warning about feedback, not as a complete description
of a particular species.

Chaotic circuits can generate signals that look noise-like and are sensitive
to circuit parameters. They have been studied for spread-spectrum
communication and encryption-related designs. Chaos alone is not a security
proof: a low-dimensional deterministic signal may be reconstructed by an
attacker, so cryptographic security requires a separate tested construction.

Heart-rate records can be analysed with return maps, entropy, and related
measures to compare long-term regulation in different physiological states.
These quantities turn a time series into measurable features, but they are
not diagnoses by themselves and must be validated against clinical data.

\paragraph{What chaos is not}
Chaos is not the same as randomness, complexity, or poor measurement. A
chaotic system has a definite future when its exact state is specified.
Practical prediction fails because the exact state is never known and small
errors grow. Conversely, an irregular data set may be caused by measurement
noise, changing parameters, or an omitted random process rather than chaos.

The butterfly effect is therefore a statement about sensitivity, not a claim
that a tiny event always causes one particular large event. Even when exact
trajectories cannot be predicted, bounds, averages, invariant distributions,
or probabilities of entering a region may remain predictable.
\theoryapplications
\theoryconnections{chaos theory}{%
\item Calculus and differential equations describe the state update, while numerical experiments reveal sensitivity, invariant sets, and long-term patterns.
\item Probability is useful when deterministic trajectories are replaced by statistical descriptions, and topology supplies the language of attractors and recurrence.
}{%
\item Chaos theory helps weather, ecology, fluid dynamics, and control distinguish short-term prediction from long-term statistical prediction.
\item Its sensitivity tests also warn numerical modelers that a small measurement or rounding error can change a detailed forecast without making the model itself random.
}
\chapterquestions{Chaos theory}{iterated maps, sensitivity, orbits, and invariant sets}{the logistic map $x_{n+1}=4x_n(1-x_n)$}{nearby initial values separate because the map repeatedly stretches and folds the interval}{positive Lyapunov exponent is evidence of exponential sensitivity}{finite precision eventually repeats an orbit; chaotic models describe deterministic unpredictability in weather and ecology}{chaos_theory_ref1}
\section*{8. References}
\begin{thebibliography}{9}
\bibitem{chaos_theory_ref1} S. H. Strogatz, \emph{Nonlinear Dynamics and Chaos}, 2nd edition, Westview Press, 2015.
\bibitem{chaos_theory_ref2} R. L. Devaney, \emph{An Introduction to Chaotic Dynamical Systems}, 2nd edition, Westview Press, 2003.
\bibitem{chaos_theory_ref3} E. Ott, \emph{Chaos in Dynamical Systems}, 2nd edition, Cambridge University Press, 2002.
\bibitem{chaos_theory_ref4} E. N. Lorenz, \emph{The Essence of Chaos}, UCL Press, 1993.
\end{thebibliography}


\clearpage\chapter{Character theory}
\section*{1. Overview}
\subsection*{Intuitive}
\paragraph{The problem}

A representation turns each symmetry of a group into a matrix. Matrices can be large and change when the basis changes. Character theory compresses the representation to one function: the trace of each matrix. For finite groups over the complex numbers, this compact function contains enough information to recover the representation up to isomorphism.



\subsection*{Formal}
Let $\rho:G\to\operatorname{GL}(V)$ be a finite-dimensional representation of a finite group over a field of characteristic zero. Its character is the class function $\chi_\rho(g)=\operatorname{Tr}(\rho(g))$. The character inner product is
\[
\langle\chi,\psi\rangle=\frac1{|G|}\sum_{g\in G}\chi(g)\overline{\psi(g)}.
\]
Irreducible characters are orthonormal, and the multiplicity of an irreducible character in $\chi_\rho$ is its inner product with $\chi_\rho$.
\section*{2. History}
\subsection*{Historical development}
\begin{historylist}
\paragraph{Ferdinand Georg Frobenius}
Frobenius introduced characters as traces of representation matrices and used them to study finite-group representations without writing every matrix explicitly. In characteristic zero, the character determines a complex representation up to isomorphism, which makes character tables effective summary data.

\paragraph{Richard Brauer}
Brauer developed modular character theory for positive characteristic. His methods handle the case in which the characteristic divides the group order, where the averaging arguments used in ordinary character theory no longer apply.

These developments turned complicated group actions into numerical data that can be added, multiplied, and averaged. The resulting character methods support the study of finite groups, molecular symmetry, particle physics, and Lie groups \cite{character_theory_ref1,character_theory_ref3}.

\end{historylist}
\section*{3. Theory}
\subsection*{Core theory}
\paragraph{Representations and characters}

Let $\rho:G\to\operatorname{GL}(V)$ be a finite-dimensional representation over a field $F$. Its character is
\[
 \chi_\rho(g)=\operatorname{Tr}(\rho(g)).
\]
The degree is $\chi_\rho(1)=\dim V$ in characteristic zero. A character is irreducible when the representation has no nonzero proper invariant subspace. A character of degree one is linear.

Characters are class functions: $\chi(hgh^{-1})=\chi(g)$. The trace is unchanged by conjugating a matrix, so a character is constant on conjugacy classes. The kernel can be recovered in characteristic zero as the elements where $\chi(g)=\chi(1)$, and it is a normal subgroup. A character is not generally a group homomorphism.

\paragraph{Character tables}

For a finite group, make one column for each conjugacy class and one row for each irreducible character. The first row is the trivial character, whose entries are all 1; the first column is the identity class and contains the degrees.

The table is square because the number of irreducible complex representations equals the number of conjugacy classes. From the first column one obtains
\[
 |G|=\sum_i \chi_i(1)^2.
\]
The sum of squared absolute values in a column equals the order of the centralizer of a representative element. Kernels of characters reveal normal subgroups; the commutator subgroup is the intersection of kernels of linear characters. A character table can show whether a group is abelian, but it need not determine the group up to isomorphism: the quaternion group and the dihedral group of order eight have the same character table.

\paragraph{Orthogonality and decomposition}

For complex class functions define
\[
 \langle\alpha,\beta\rangle=
 \frac1{|G|}\sum_{g\in G}\alpha(g)\overline{\beta(g)}.
\]
Irreducible characters form an orthonormal basis, so
\[
 \langle\chi_i,\chi_j\rangle=\delta_{ij}.
\]
If $\chi=\sum_i m_i\chi_i$, then $m_i=\langle\chi,\chi_i\rangle$. Thus the inner product tells how many times each irreducible representation occurs in a decomposition.

The column orthogonality relation says that the corresponding columns are orthogonal after weighting by centralizers. These relations are practical algorithms for completing a character table, finding centralizer sizes, decomposing a representation, and checking whether a proposed row is possible.

\paragraph{Arithmetic rules}

Characters turn common representation operations into simple formulas:
\[
 \chi_{\rho\oplus\sigma}=\chi_\rho+\chi_\sigma,\qquad
 \chi_{\rho\otimes\sigma}=\chi_\rho\chi_\sigma,\qquad
 \chi_{\rho^*}=\overline{\chi_\rho}.
\]
For exterior and symmetric squares,
\[
 \chi_{\wedge^2\rho}(g)=\frac{\chi_\rho(g)^2-\chi_\rho(g^2)}2,
\quad
 \chi_{\operatorname{Sym}^2\rho}(g)=\frac{\chi_\rho(g)^2+\chi_\rho(g^2)}2.
\]
For finite-order $g$, the eigenvalues of $\rho(g)$ are roots of unity, so character values are sums of roots of unity and hence algebraic integers over $\mathbb C$.

\paragraph{Restriction, induction, and Frobenius reciprocity}

If $H\leq G$, restricting a representation from $G$ to $H$ restricts its character. Conversely, a representation of $H$ can be induced to a representation of $G$ by letting $G$ act on functions or cosets with values in the original space. The induced character is zero on elements not conjugate to an element of $H$ and has a coset-sum formula on elements of $H$.

Frobenius reciprocity says
\[
 \langle\theta^G,\chi\rangle_G
 =\langle\theta,\chi_H\rangle_H.
\]
It converts multiplicity calculations for induction into calculations for restriction and is one of the main tools for building character tables from subgroups.

\paragraph{Lie groups and Lie algebras}

For a finite-dimensional representation of a Lie group, the character is still $\operatorname{Tr}(\rho(g))$. For a Lie algebra representation, one may define
\[
 \chi_\rho(X)=\operatorname{Tr}(e^{\rho(X)}).
\]
The character is invariant under the adjoint action. For a complex semisimple Lie algebra, it is determined by its values on a Cartan subalgebra and can be written using weights:
\[
 \chi_\rho(H)=\sum_\lambda m_\lambda e^{\lambda(H)}.
\]
The Weyl character formula computes irreducible characters explicitly from roots and weights.

\section*{4. Important theorems and results}
\begin{enumerate}[leftmargin=*,itemsep=0.35em]

\item \textbf{Character determination theorem.} Complex representations of a finite group are isomorphic exactly when they have the same character.

\item \textbf{Schur orthogonality.} Irreducible characters form an orthonormal basis of the class functions, giving both row and column orthogonality.

\item \textbf{Sum-of-squares theorem.} The order of a finite group equals the sum of squares of the degrees of its irreducible complex representations.

\item \textbf{Maschke's theorem.} If the characteristic does not divide $|G|$, every finite-dimensional representation is completely reducible. This is why character decomposition works cleanly in characteristic zero.

\item \textbf{Frobenius reciprocity.} Multiplicities in an induced representation equal multiplicities in the corresponding restricted representation.

\item \textbf{Weyl character formula.} The character of an irreducible representation of a complex semisimple Lie algebra is determined explicitly by its highest weight and root system.
\end{enumerate}
\noindent\emph{Source for these results:} \cite{character_theory_ref1}.

\section*{5. Applications}
Character theory compresses a representation into traces of group elements, while retaining enough information to study its decomposition. Character values therefore make symmetry calculations manageable in algebra, combinatorics, chemistry, and physics.
\begin{enumerate}[leftmargin=*,itemsep=0.35em]
\item \textbf{Finite-group structure.} Character values reveal normal subgroups, commutators, centralizers, and possible degrees. Character calculations occur in classification of finite simple groups and in the Feit--Thompson theorem.
\item \textbf{Chemistry and physics.} Molecular vibrations and particle states transform under symmetry groups. Character tables separate modes into irreducible types and predict which transitions or interactions are allowed.
\item \textbf{Combinatorics.} The representation of a group on a set gives a permutation character. Inner products count orbits and fixed configurations, connecting characters with Burnside's lemma.
\item \textbf{Number theory and analysis.} Linear characters of abelian groups include Dirichlet characters and lead to Fourier analysis and $L$-functions. Higher-dimensional characters enter automorphic forms and arithmetic geometry.
\item \textbf{Moonshine.} Replacing ordinary dimension by a graded character produces McKay--Thompson series and connects representations of the Monster group with modular functions.
\end{enumerate}


\theoryconnections{character_theory}{%
\item Group theory provides conjugacy classes, which collect elements that have the same symmetry type.
\item Representation theory turns each group element into a matrix, and the character records only the trace of that matrix; traces are easier to compare and still remember enough to decompose a representation.
\item Orthogonality of characters is a Fourier calculation on the group, so inner products count how many irreducible symmetry patterns occur \cite{character_theory_ref1,character_theory_ref2}.
}{%
\item Character theory helps classify finite-group representations by replacing a large collection of matrices with a table of values on conjugacy classes.
\item In chemistry those values predict how molecular vibrations split into symmetry types; in particle physics they label states under continuous symmetries.
\item In combinatorics and number theory, character sums convert symmetry or congruence conditions into cancellations that can be estimated \cite{character_theory_ref1,character_theory_ref3}.
}

\section*{Glossary}
\begin{description}[leftmargin=*,style=nextline,itemsep=0.3em]

\item[\textbf{Character.}] The class function $\chi_\rho(g)=\operatorname{Tr}(\rho(g))$ assigning to each group element the trace of the corresponding representation matrix; it determines the representation up to isomorphism for finite groups over $\mathbb C$.

\item[\textbf{Irreducible character.}] The character of a representation that has no nonzero proper invariant subspaces; irreducible characters form an orthonormal basis for class functions.

\item[\textbf{Conjugacy class.}] A set of group elements that are all conjugate to one another; characters are constant on conjugacy classes because the trace is unchanged by similarity.

\item[\textbf{Character table.}] A square matrix whose rows are irreducible characters and whose columns are conjugacy classes, serving as a complete summary of a finite group's representation theory.

\item[\textbf{Induced character.}] The character of a representation of a larger group $G$ constructed from a representation of a subgroup $H$ by letting $G$ act on coset copies of the original space.

\item[\textbf{Frobenius reciprocity.}] A duality stating that the multiplicity of an irreducible character of $G$ in an induced representation equals the multiplicity of the original character of $H$ in the restriction.

\item[\textbf{Class function.}] A function on a group that is constant on each conjugacy class; the space of class functions is the natural domain for character inner products.

\item[\textbf{Representation.}] A group homomorphism $\rho:G\to\operatorname{GL}(V)$ realizing each group element as a linear transformation; characters are traces of these matrices.

\item[\textbf{Maschke's theorem.}] Every finite-dimensional representation of a finite group is completely reducible when the field characteristic does not divide the group order.

\item[\textbf{Schur orthogonality.}] The relations stating irreducible characters form an orthonormal basis for class functions.

\item[\textbf{Degree of a character.}] The dimension $\chi(1)$ of the underlying representation space; it appears in the sum-of-squares formula.

\end{description}

\section*{7. Sample Questions}
\emph{Exercise reference:} \cite{character_theory_ref1}. The cited edition is the source for the exercise set; consult its Exercises section for the official statement and numbering.
\begin{enumerate}[leftmargin=*,itemsep=0.9em]

\item \textbf{Exercise 1.} Construct the character table of $\mathbb{Z}/4\mathbb{Z}$.

\emph{Solution.} $\mathbb{Z}/4\mathbb{Z}$ is abelian with 4 conjugacy classes (each element is its own class) and 4 irreducible characters, all of degree 1. Let $\omega=i=\sqrt{-1}$. The characters are $\chi_k(g)=\omega^{kg}$ for $k=0,1,2,3$ and $g=0,1,2,3$:

\[
\begin{array}{c|cccc}
& 0 & 1 & 2 & 3\\\hline
\chi_0 & 1 & 1 & 1 & 1\\
\chi_1 & 1 & i & -1 & -i\\
\chi_2 & 1 & -1 & 1 & -1\\
\chi_3 & 1 & -i & -1 & i
\end{array}
\]

Check: $\sum_k|\chi_k(1)|^2=4=|\mathbb{Z}/4\mathbb{Z}|$.

\item \textbf{Exercise 2.} Let $\chi$ be the character of the standard 2-dimensional representation of $S_3$ (the representation on $\{(x,y,z):x+y+z=0\}$). Compute $\langle\chi,\chi\rangle$ and determine whether the representation is irreducible.

\emph{Solution.} The conjugacy classes of $S_3$ are: identity (size 1), transpositions (size 3), 3-cycles (size 2). The character values are $\chi(1)=2$, $\chi((12))=0$, $\chi((123))=-1$. Then $\langle\chi,\chi\rangle=\frac{1}{6}(1\cdot|2|^2+3\cdot|0|^2+2\cdot|-1|^2)=\frac{1}{6}(4+0+2)=1$. Since $\langle\chi,\chi\rangle=1$, the representation is irreducible.

\item \textbf{Exercise 3.} Compute the character of the tensor product $\chi\otimes\psi$ where $\chi$ and $\psi$ are the two nontrivial irreducible characters of $\mathbb{Z}/3\mathbb{Z}$, evaluated at the generator $g$ of order 3.

\emph{Solution.} Let $\omega=e^{2\pi i/3}$. The irreducible characters of $\mathbb{Z}/3\mathbb{Z}=\{0,1,2\}$ are $\chi_0(g^k)=1$, $\chi_1(g^k)=\omega^k$, $\chi_2(g^k)=\omega^{2k}$. The nontrivial ones are $\chi_1$ and $\chi_2$. Their tensor product: $(\chi_1\otimes\chi_2)(g^k)=\omega^k\cdot\omega^{2k}=\omega^{3k}=1$. So $\chi_1\otimes\chi_2=\chi_0$, the trivial character. This reflects the fact that $\chi_2=\overline{\chi_1}$ and $\chi_1\cdot\overline{\chi_1}=|\chi_1|^2=1$.

\item \textbf{Exercise 4.} The character of the regular representation of a finite group $G$ is $\chi_{\text{reg}}(g)=|G|$ if $g=e$ and $0$ otherwise. Decompose $\chi_{\text{reg}}$ into irreducibles for $G=\mathbb{Z}/2\mathbb{Z}$.

\emph{Solution.} $G=\{e,g\}$, $|G|=2$. The irreducible characters are $\chi_0=(1,1)$ (trivial) and $\chi_1=(1,-1)$ (sign). The regular character is $\chi_{\text{reg}}=(2,0)$. Then $m_0=\langle\chi_{\text{reg}},\chi_0\rangle=\frac{1}{2}(2\cdot 1+0\cdot 1)=1$ and $m_1=\langle\chi_{\text{reg}},\chi_1\rangle=\frac{1}{2}(2\cdot 1+0\cdot(-1))=1$. So $\chi_{\text{reg}}=\chi_0+\chi_1$, and the regular representation decomposes as the direct sum of both irreducibles, each appearing once.

\item \textbf{Exercise 5.} Use the sum-of-squares formula to determine the order of a group whose irreducible character degrees are $1,1,1,2$.

\emph{Solution.} By the formula $|G|=\sum_i\chi_i(1)^2=1^2+1^2+1^2+2^2=1+1+1+4=7$. So $|G|=7$. Since 7 is prime, $G\cong\mathbb{Z}/7\mathbb{Z}$, and indeed this abelian group has three degree-1 characters... wait, $\mathbb{Z}/7\mathbb{Z}$ has 7 conjugacy classes and 7 irreducible characters, all of degree 1. The degrees $1,1,1,2$ sum to $|G|=7$ but require 4 conjugacy classes. In fact $|G|=1+1+1+4=7$ works for a nonabelian group of order 7? No—order 7 is abelian. Let me restate: there is no group with exactly these degrees since $|G|=7$ forces all degrees to be 1. The correct interpretation: the sum-of-squares formula gives $|G|=7$, but this is inconsistent with having a degree-2 representation (which requires $|G|\ge 6$ and nonabelian, but $|G|=7$ is abelian). So no such group exists. The formula $|G|=7$ is the answer, and the inconsistency shows these cannot be the character degrees of any group.

\end{enumerate}
\section*{8. References}
\begin{thebibliography}{99}
\bibitem{character_theory_ref1} J.-P. Serre, \emph{Linear Representations of Finite Groups}, Springer, 1977.
\bibitem{character_theory_ref2} I. M. Isaacs, \emph{Character Theory of Finite Groups}, Academic Press, 1976.
\bibitem{character_theory_ref3} B. C. Hall, \emph{Lie Groups, Lie Algebras, and Representations}, 2nd edition, Springer, 2015.
\end{thebibliography}


\theory{Choquet theory}{How can a complicated point in a convex collection be described by simpler boundary points?}{Choquet theory studies averages in convex sets. In a triangle, every interior point is a weighted average of the three corners. In an infinite-dimensional space there may be infinitely many possible building blocks, so an ordinary finite sum is replaced by a probability measure and an integral. The theory asks when such a representation exists and when it is unique.}{Functional analysis, convex optimization, probability, potential theory, economics, and any setting in which a complicated object is treated as a mixture of simpler extreme objects.}{Convex combinations and representing measures identify when a state is a mixture of extreme states and whether that mixture is unique.}

\paragraph{The question and its history}

Many decisions and mathematical objects are mixtures. A point halfway between two points on a line is obtained by giving each endpoint weight one half. A point in a triangle is obtained from its corners by three nonnegative weights whose sum is one. The corners are special because they cannot themselves be split into a nontrivial average of other points in the same set. They are called extreme points.

In finite-dimensional geometry, this observation leads to the theory of convex hulls and barycentric coordinates. Gustave Choquet developed the infinite-dimensional version in the setting of functional analysis, with important motivation from potential theory. The central change is that a point may need infinitely many extreme points. Instead of listing weights one by one, we use a probability measure: it records how much weight is placed near each possible extreme point. \cite{choquet_theory_ref1,choquet_theory_ref3}

The purpose is not merely to decorate a geometric picture with advanced language. Representation by extreme points can turn a difficult problem about all points of a large set into a problem about its boundary. It also tells us when two different mixtures describe the same object, which is essential when interpreting a representation as a genuine explanation rather than one of many arbitrary descriptions.

\end{historylist}
\section*{3. Theory}
\subsection*{Core theory}
\paragraph{Convex sets, extreme points, and averages}

Let $V$ be a real vector space and let $C$ be a subset of it. The set is \emph{convex} if, whenever $x$ and $y$ belong to $C$, every point
\[
 \lambda x+(1-\lambda)y,\qquad 0\leq \lambda\leq 1,
\]
also belongs to $C$. In plain language, one may travel in a straight line between two allowed points without leaving the set.

An element $e$ of $C$ is an \emph{extreme point} when it is impossible to write
$e=\lambda x+(1-\lambda)y$ with $0<\lambda<1$ and two distinct points $x,y$ of $C$. The endpoints of a line segment are extreme; its interior points are not. The vertices of a polygon are extreme, while points along its edges are not. A disk has every point of its circular boundary as an extreme point.

A finite convex combination is an average
\[
 c=\sum_{i=1}^{n}w_i e_i,\qquad
 w_i\geq 0,\qquad \sum_{i=1}^{n}w_i=1.
\]
The numbers $w_i$ are probabilities. If $f$ is an affine function, meaning that it preserves these weighted averages, then
\[
 f(c)=\sum_{i=1}^{n}w_i f(e_i).
\]
Thus knowing an affine quantity on the extreme points determines it throughout the convex set.

For example, in a triangle with vertices $A=(0,0)$, $B=(1,0)$, and $C=(0,1)$, the point $(1/4,1/2)$ is
\[
 (1/4)A+(1/4)B+(1/2)C.
\]
 The three weights are nonnegative and sum to one. In a triangle they are unique. In a square they are not: the centre can be obtained by averaging the two ends of either diagonal, or by averaging all four vertices.

\paragraph{From finite sums to probability measures}

In an infinite-dimensional vector space a point may not be a finite blend of extreme points. Choquet theory therefore considers a probability measure $w$ on the extreme-point set $E$. The measure is allowed to spread weight over infinitely many points. For an affine function $f$ on $C$, the statement that $w$ represents $c$ is
\[
 f(c)=\int_E f(e)\,dw(e).
\]
The integral is an average: it is the continuous version of the finite sum. Saying that the measure is supported on $E$ means that it gives full weight to the extreme points. A representation is useful only when the topology is chosen carefully, because in an infinite-dimensional space limits of averages, continuity, and compactness are what make the integral meaningful.

The main setting is a compact convex subset $C$ of a locally convex Hausdorff topological vector space. "Locally convex" means that the space has enough continuous linear measurements to distinguish points and to describe small convex neighborhoods. The space need not be finite-dimensional or even metrizable; in many applications it is a Banach space or a weak-star space. \cite{choquet_theory_ref1,choquet_theory_ref2}

The finite-dimensional theorem behind this picture is Minkowski's theorem: a closed, bounded convex subset of Euclidean space is the convex hull of its extreme points. Carathéodory's theorem sharpens the practical calculation: in $\mathbb R^d$, a point in a convex hull can be represented using at most $d+1$ points. Choquet theory asks for the appropriate infinite-dimensional replacement when no such fixed finite bound exists.

\paragraph{Choquet's theorem and its limits}

\textbf{Choquet's theorem.} If $C$ is compact, convex, and metrizable in a locally convex topological vector space, then every $c\in C$ has a probability measure $w$ supported on the extreme points of $C$ such that
\[
 f(c)=\int_C f(e)\,dw(e)
\]
for every continuous affine function $f$ on $C$. \cite{choquet_theory_ref1,choquet_theory_ref2}

This theorem says exactly what the triangle example suggests, but with a probability distribution in place of finitely many barycentric weights. It is an existence statement, not automatically a uniqueness statement. It also requires hypotheses. If compactness is lost, a representing measure may escape to infinity. If the topology is too weak or the space is not locally convex, continuous affine tests may no longer control the problem.

When $C$ is not metrizable, the Choquet--Bishop--de Leeuw theorem gives a related result. The representing measure need not be concentrated on the extreme points in the strongest ordinary sense. Instead it vanishes on Baire subsets that contain no extreme points. This formulation is designed for very large topological spaces, where ordinary measurability and countability arguments are unavailable. \cite{choquet_theory_ref3}

The theorem has two important consequences. The Krein--Milman theorem says that a compact convex set in a locally convex space is the closed convex hull of its extreme points. A Riesz representation theorem for states on continuous functions says that a positive normalized linear functional can be represented by a probability measure. In both cases, a global object is recovered from elementary boundary data. \cite{choquet_theory_ref1,choquet_theory_ref2}

\paragraph{Uniqueness, simplices, and concrete examples}

Existence does not imply uniqueness. In a square, the centre has the two representations
\[
 \tfrac12(1,0)+\tfrac12(0,1)
 \quad\hbox{and}\quad
 \tfrac12(0,0)+\tfrac12(1,1).
\]
The same point therefore has different probability measures on the extreme boundary. A ball in $\mathbb R^3$ also has many such representations: its centre can be obtained from any pair of opposite boundary points, or from a uniform distribution on the whole sphere.

A finite-dimensional simplex is different. A triangle is a two-dimensional simplex and a tetrahedron is a three-dimensional one. Each point has one and only one set of vertex weights. The infinite-dimensional generalization is a \emph{Choquet simplex}: every point has a unique representing probability measure on the extreme points. This uniqueness makes Choquet simplices useful when a state, mixture, or probability distribution is meant to be an intrinsic decomposition rather than a choice of coordinates. \cite{choquet_theory_ref1}

 The contrast gives a practical diagnostic. If the model has a square-like collection of alternatives, several explanations may fit the same observed average. If it has simplex-like structure, the weights are forced by the data. This distinction appears in probability theory, the study of positive functions, and optimization over spaces of measures.

\paragraph{What the theory depends on and what it enables}

Choquet theory depends on elementary linear algebra, convex combinations, topology, continuous functions, measures, and integration. Functional analysis supplies the locally convex spaces, while convex analysis supplies extreme points and convex hulls. Potential theory supplied many of the original questions, especially the representation of positive harmonic functions. \cite{choquet_theory_ref1,choquet_theory_ref4}

The theory helps functional analysis by reducing positive linear functionals and states to measures on boundary objects. In probability it explains mixtures of extreme probability laws. In optimization it suggests that extremal solutions deserve special attention, although infinite-dimensional optimization may require compactness and lower-semicontinuity arguments. In economics, a population state can be understood as a mixture of extreme preferences or deterministic choices, with uniqueness indicating whether that interpretation is identifiable. In operator theory, state spaces of algebras form convex sets whose extreme points are often pure states.

The limits are just as important as the uses. A measure representation may be nonunique, may live on a boundary that is difficult to describe, or may fail under noncompactness. Thus Choquet theory is a method for organizing a problem, not a promise that every representation is easy to calculate.

\section*{4. Important theorems and results}
\begin{enumerate}[leftmargin=*,itemsep=0.35em]

\item \textbf{Minkowski theorem.} A closed, bounded convex subset of finite-dimensional Euclidean space is the convex hull of its extreme points. \cite{choquet_theory_ref5}

\item \textbf{Carathéodory theorem.} In $\mathbb R^d$, any point in the convex hull of a set can be written as a convex combination of at most $d+1$ points. \cite{choquet_theory_ref5}

\item \textbf{Choquet representation theorem.} Every point of a compact metrizable convex set in a locally convex space has a probability representation supported on the extreme points, tested through continuous affine functions. \cite{choquet_theory_ref1}

\item \textbf{Krein--Milman theorem.} A compact convex set in a locally convex Hausdorff space is the closed convex hull of its extreme points. \cite{choquet_theory_ref2}

\item \textbf{Choquet simplex criterion.} In a Choquet simplex, the representing probability measure on the extreme boundary is unique. Cubes and ordinary balls show why uniqueness cannot be expected for arbitrary convex sets. \cite{choquet_theory_ref1}
\end{enumerate}

\theoryapplications
\theoryconnections{choquet theory}{%
\item Convex analysis supplies mixtures and extreme points, functional analysis supplies locally convex spaces, and probability supplies representing measures.
\item Integration and topology are needed to make an infinite mixture meaningful and to state when a representation is supported on the boundary.
}{%
\item Choquet theory helps optimization identify boundary solutions, helps potential theory represent harmonic functions, and helps probability interpret a complicated state as a mixture of simpler extreme states.
\item Uniqueness is valuable because it turns a representation into an explanation rather than one arbitrary decomposition.
}

\section*{Glossary}
\begin{description}[leftmargin=*,style=nextline,itemsep=0.3em]

\item[\textbf{Extreme point.}] An element of a convex set that cannot be written as a nontrivial average of two other distinct points in the set; vertices of a polygon are extreme points.

\item[\textbf{Convex set.}] A set in which every line segment joining two of its points lies entirely within it; equivalently, any finite weighted average of points in the set remains in the set.

\item[\textbf{Representing measure.}] A probability measure supported on the extreme points of a convex set such that the value of every continuous affine function at an interior point equals its integral against the measure.

\item[\textbf{Choquet simplex.}] A compact convex set in which every point has a unique representing probability measure on the extreme boundary, generalizing the uniqueness of barycentric coordinates in a triangle.

\item[\textbf{Affine function.}] A function that preserves weighted averages, meaning $f(\lambda x+(1-\lambda)y)=\lambda f(x)+(1-\lambda)f(y)$; affine functions test whether a measure represents a point.

\item[\textbf{Barycentric coordinates.}] The weights in a finite convex combination expressing a point as an average of extreme points; in a simplex these weights are unique.

\item[\textbf{Krein--Milman theorem.}] The result that a compact convex set in a locally convex space is the closed convex hull of its extreme points, guaranteeing that extreme points completely determine the set.

\item[\textbf{Convex hull.}] The smallest convex set containing a given collection of points; in finite dimensions equals the convex hull of extreme points.

\item[\textbf{Probability measure.}] A measure assigning total mass 1; it replaces finite weights in the infinite-dimensional representation.

\item[\textbf{Compact convex set.}] A convex set that is topologically compact; this ensures representing measures exist.

\item[\textbf{Locally convex space.}] A topological vector space with enough continuous linear functionals to separate points and define neighborhoods via convex sets.

\end{description}

\section*{7. Sample Questions}
\emph{Exercise reference:} \cite{choquet_theory_ref1}. The cited edition is the source for the exercise set; consult its Exercises section for the official statement and numbering.
\begin{enumerate}[leftmargin=*,itemsep=0.9em]

\item \textbf{Exercise 1.} Find the barycentric coordinates of the point $(1,1)$ as a convex combination of the vertices $(0,0)$, $(4,0)$, and $(0,3)$ of a triangle.

\emph{Solution.} We need $(1,1)=a(0,0)+b(4,0)+c(0,3)$ with $a+b+c=1$ and $a,b,c\ge 0$. From coordinates: $4b=1$ so $b=1/4$; $3c=1$ so $c=1/3$. Then $a=1-1/4-1/3=5/12$. Check: $(5/12)(0,0)+(1/4)(4,0)+(1/3)(0,3)=(1,0)+(0,1)=(1,1)$.

\item \textbf{Exercise 2.} A probability measure on the vertices of a square $A=(0,0)$, $B=(1,0)$, $C=(1,1)$, $D=(0,1)$ assigns weight $1/4$ to each vertex. Compute the barycenter and find a second distinct probability measure on the four vertices that produces the same barycenter.

\emph{Solution.} Barycenter: $(1/4)(0,0)+(1/4)(1,0)+(1/4)(1,1)+(1/4)(0,1)=(1/2,1/2)$. A second measure: weight $1/2$ on $A$ and $1/2$ on $C$: $(1/2)(0,0)+(1/2)(1,1)=(1/2,1/2)$. This is a different probability measure on the extreme points giving the same barycenter, confirming the square is not a Choquet simplex.

\item \textbf{Exercise 3.} Determine whether the triangle with vertices $(0,0)$, $(1,0)$, $(0,1)$ is a Choquet simplex.

\emph{Solution.} Yes. A triangle is a 2-simplex. Every point in a simplex has a unique set of barycentric coordinates. To see uniqueness: if $(x,y)=a(0,0)+b(1,0)+c(0,1)$ with $a+b+c=1$, then $b=x$, $c=y$, $a=1-x-y$ is uniquely determined. So the representing measure on the extreme points is unique, and the triangle is a Choquet simplex.

\item \textbf{Exercise 4.} On the interval $[0,1]$ (a compact convex set), a point $c\in(0,1)$ is not an extreme point. Express $c$ as a convex combination of the extreme points $0$ and $1$ and verify it is the only such representation.

\emph{Solution.} The extreme points of $[0,1]$ are just $\{0,1\}$. Any $c\in(0,1)$ satisfies $c=c\cdot 1+(1-c)\cdot 0$. This is the unique representation since any convex combination of $\{0,1\}$ is $t\cdot 1+(1-t)\cdot 0=t$, so $t=c$ is forced. The interval is a 1-simplex, hence a Choquet simplex.

\item \textbf{Exercise 5.} On the unit disk $D=\{(x,y):x^2+y^2\le 1\}$ in $\mathbb{R}^2$, describe the extreme boundary and show that the origin has infinitely many representing measures supported on this boundary.

\emph{Solution.} The extreme boundary of $D$ is the unit circle $S^1=\{(x,y):x^2+y^2=1\}$. The origin $(0,0)$ can be written as $(0,0)=\frac{1}{2}(1,0)+\frac{1}{2}(-1,0)$, or as $(0,0)=\frac{1}{2}(0,1)+\frac{1}{2}(0,-1)$, or as the uniform average over any diameter. In general, for any pair of antipodal points $(\cos\theta,\sin\theta)$ and $(-\cos\theta,-\sin\theta)$ on $S^1$, the measure giving weight $1/2$ to each represents the origin. This gives a continuous family of distinct representing measures, so uniqueness fails and the disk is not a Choquet simplex.

\end{enumerate}
\section*{8. References}


\begin{thebibliography}{9}
\bibitem{choquet_theory_ref1} R. R. Phelps, \emph{Lectures on Choquet's Theorem}, Springer, 2nd ed., 2001.
\bibitem{choquet_theory_ref2} E. M. Alfsen, \emph{Compact Convex Sets and Boundary Integrals}, Springer, 1971.
\bibitem{choquet_theory_ref3} E. Bishop and K. de Leeuw, "The representations of linear functionals by measures on sets of extreme points," \emph{Annales de l'Institut Fourier}, 1959.
\bibitem{choquet_theory_ref4} L. Asimow and A. J. Ellis, \emph{Convexity Theory and its Applications in Functional Analysis}, Academic Press, 1980.
\bibitem{choquet_theory_ref5} R. T. Rockafellar, \emph{Convex Analysis}, Princeton University Press, 1970.
\end{thebibliography}


\theory{Class field theory}{How can all abelian extensions of a number system be described using arithmetic information already present in that system?}{Class field theory is the part of algebraic number theory that describes abelian Galois extensions of local and global fields. It translates field extensions into ideal classes, norms, units, and ideles, and it gives reciprocity laws that generalize quadratic reciprocity.}{Prime numbers, algebraic number theory, explicit reciprocity, Iwasawa theory, the Langlands program, arithmetic geometry, and the study of how primes split in extensions.}{Reciprocity maps ideals or ideles to abelian Galois groups, turning prime-splitting questions into calculations in the base field.}

\paragraph{The problem and its history}

A field is a number system in which addition, subtraction, multiplication, and division by nonzero numbers are possible. A field extension is a larger system containing the original one. For example, adjoining $\sqrt{2}$ to the rational numbers gives $\mathbb Q(\sqrt2)$. Galois theory studies the symmetries of an extension: a Galois automorphism rearranges the new numbers while leaving the old numbers fixed.

Class field theory asks for a complete description when these symmetries commute with one another, the \emph{abelian} case. The surprising idea is that the answer can be read from arithmetic objects in the original field: ideals, their classes, local units, and a global object called the idele class group. \cite{class_field_theory_ref1,class_field_theory_ref4}

The story began with Gauss's quadratic reciprocity law, which relates whether one prime is a square modulo another to the reversed question. The later generalization involved quadratic forms, genus theory, Kummer's work, ideals and completions, and cyclotomic extensions. Kronecker and Weber recognized explicit patterns for the rational numbers; Weber coined the phrase class field, and Hilbert formulated central problems. Takagi proved the existence theorem, while Artin developed the reciprocity law, with Chebotarev's density theorem supplying a crucial description of primes in extensions. \cite{class_field_theory_ref2,class_field_theory_ref3}

\end{historylist}
\section*{3. Theory}
\subsection*{Core theory}
\paragraph{Local fields, global fields, and the central dictionary}

A number field is a finite extension of $\mathbb Q$, such as $\mathbb Q(i)$ or $\mathbb Q(\sqrt5)$. A global function field is a finite extension of a rational function field over a finite field. These are the global fields. A local field is obtained by completing a global field at one of its places. Completion means that we add limits of Cauchy sequences, just as the real numbers complete the rationals. For a prime $p$, the field $\mathbb Q_p$ remembers very precise information about divisibility by $p$.

The maximal abelian extension $A$ of a field $K$ is the largest algebraic extension whose Galois symmetries commute. Its Galois group is usually infinite, but it has a natural topology: finite pieces of information determine open neighborhoods. It is therefore a profinite group, an inverse limit of finite groups.

Class field theory constructs an arithmetic topological group associated with $K$. For a local field with finite residue field, the key object is the multiplicative group $K^\times$. For a global field, the key object is the idele class group $C_K$. An idele is a compatible collection of local numbers, one for each place, with a finiteness condition; quotienting by the diagonal copy of $K^\times$ removes the descriptions that come from a single global number. \cite{class_field_theory_ref1,class_field_theory_ref5}

The dictionary has two directions:

{\small
\[
\begin{array}{c|c}
\text{Arithmetic information in }K & \text{Abelian extension information}\\
\hline
\text{an open finite-index subgroup of }C_K & \text{a finite abelian extension of }K\\
\text{a norm subgroup from }L & \text{the extension }L/K\\
\text{an ideal class in an unramified situation} & \text{a Frobenius automorphism}
\end{array}
\]
}

This is useful because arithmetic groups can often be computed or approximated, whereas listing all algebraic extensions directly would be impossible.

\paragraph{The Hilbert class field}

Let $F$ be a number field. Its ideals form a multiplicative system, but not every ideal is generated by one element. The ideal class group measures this failure: two ideals have the same class when their quotient differs by multiplication by a principal ideal. The class group is finite, and its size is the class number.

The Hilbert class field $K$ of $F$ is the maximal abelian extension of $F$ that is unramified at every finite and infinite place. The Hilbert class field theorem gives a canonical isomorphism
\[
 \operatorname{Gal}(K/F)\cong \operatorname{Cl}(F).
\]
Thus the number of ideal classes is exactly the number of symmetries in the largest everywhere-unramified abelian extension. An ideal $\mathfrak p$ not dividing the discriminant determines a Frobenius automorphism in the Galois group; under the isomorphism, this automorphism is represented by the ideal class of $\mathfrak p$. \cite{class_field_theory_ref1,class_field_theory_ref2}

For a simple example, the rational field has class group of size one and no nontrivial everywhere-unramified abelian extension. A field with a nontrivial class group has a genuinely larger Hilbert class field. Computing ideals and their classes therefore reveals an extension field that may be difficult to write down directly.

\paragraph{Artin reciprocity and the existence theorem}

For a finite abelian extension $L/F$, the idelic norm map sends an idele of $L$ to an idele of $F$. Class field theory states that the quotient
\[
 \theta_{L/F}:C_F/N_{L/F}(C_L)\longrightarrow \operatorname{Gal}(L/F)
\]
is a canonical isomorphism. The map is the Artin reciprocity map. In the unramified prime case, it sends the class of a prime to its Frobenius automorphism, so the abstract map connects arithmetic with the way that prime acts on the roots of a polynomial modulo that prime. \cite{class_field_theory_ref1,class_field_theory_ref4}

The existence theorem says that finite abelian extensions of $F$ correspond bijectively to closed subgroups of finite index in $C_F$. The subgroup attached to $L$ is the norm group $N_{L/F}(C_L)$. In one direction, the extension produces a subgroup. In the other, a suitable subgroup produces a unique class field. This is why the theory is called class field theory: a class or subgroup of the arithmetic object determines a field.

The formulation also has an infinite version. The Galois group of the maximal abelian extension is naturally the profinite completion of $C_F$. Equivalently, the abelianization of the absolute Galois group is determined by the arithmetic class group. This unifies all finite abelian extensions instead of studying them one at a time. \cite{class_field_theory_ref1,class_field_theory_ref5}

\paragraph{Local and global reciprocity}

Local class field theory begins with a local field $K$. It gives a reciprocity isomorphism from $K^\times$ to the Galois group of its maximal abelian extension, interpreted with the correct profinite topology. The norm subgroup of a finite abelian extension $L/K$ is the kernel of the corresponding map. Local units describe the part of the extension that is ramified, while the valuation records the unramified direction.

To pass from local to global theory, complete a number field at every place and apply the local maps. Their product is trivial on a global element embedded diagonally in all local fields. This is the global reciprocity law. It generalizes the product formula and reaches far beyond quadratic reciprocity: local contributions at all primes and at infinity fit together to produce one global statement. \cite{class_field_theory_ref1,class_field_theory_ref3}

There are two broad proof styles. Artin and Tate developed a cohomological method using group cohomology and class formations. It gives a powerful conceptual organization of norms, extensions, and obstructions. Later, Neukirch and others developed more explicit and algorithmic approaches that avoid some cohomology, making the reciprocity maps easier to calculate in examples. The two approaches prove the same central results but illuminate different parts of the subject. \cite{class_field_theory_ref3,class_field_theory_ref6}

\paragraph{Explicit class field theory}

General class field theory tells us that extensions exist and what arithmetic group controls them. Explicit class field theory tries to write the extensions down.

The Kronecker--Weber theorem says that every finite abelian extension of $\mathbb Q$ is contained in a cyclotomic field $\mathbb Q(\zeta_n)$, where $\zeta_n$ is an $n$th root of unity. Thus roots of unity generate all abelian extensions of the rational numbers. For example, the quadratic field $\mathbb Q(i)$ lies in $\mathbb Q(\zeta_4)$, and the real quadratic field $\mathbb Q(\sqrt5)$ lies in the maximal real subfield of $\mathbb Q(\zeta_5)$. \cite{class_field_theory_ref1,class_field_theory_ref4}

For imaginary quadratic fields, elliptic curves with complex multiplication give another explicit source. Torsion points of these curves and special values of modular functions generate abelian extensions. Shimura's theory extends this explicit viewpoint to broader classes of fields. In positive characteristic, cyclotomic and Drinfeld-module constructions provide parallel descriptions, and Witt duality gives a particularly direct description of the $p$-part of reciprocity in some settings. \cite{class_field_theory_ref4}

These explicit constructions are powerful but special. Roots of unity are naturally available over $\mathbb Q$, and complex multiplication supplies extra structure for imaginary quadratic fields. A general number field usually lacks such a visible source of generators, which is why the idelic formulation is needed.

\paragraph{Applications, boundaries, and generalizations}

Class field theory explains how primes split, remain inert, or ramify in abelian extensions. The Frobenius map turns this into arithmetic data that can be used to test congruences and factorization patterns. The theory is also used in Iwasawa theory, Galois module theory, Artin--Verdier duality, and explicit studies of special values of $L$-functions. Many advances toward the Langlands correspondence, the Birch and Swinnerton-Dyer conjecture, and Iwasawa theory use explicit class field methods or their extensions. \cite{class_field_theory_ref1,class_field_theory_ref7}

Class field theory deliberately treats the abelian case. Nonabelian Galois groups do not fit into the same simple correspondence with subgroups of an abelian class group. The Langlands program is often viewed as a proposed nonabelian class field theory, but it does not currently provide the same existence theorem for class fields. Anabelian geometry asks whether a number field or a hyperbolic curve can be reconstructed from its full absolute Galois group or algebraic fundamental group. Higher class field theory studies abelian extensions of higher local and global fields; its arithmetic groups are higher $K$-groups such as Milnor $K$-groups rather than only $K_1=K^\times$. \cite{class_field_theory_ref1,class_field_theory_ref8}

Class field theory depends on field theory, Galois theory, ideals, prime factorization, completions, topological groups, and group cohomology. It helps algebraic number theory by providing a complete abelian extension dictionary. It also supplies the abelian base case for the Langlands program, local tools for arithmetic geometry, and explicit reciprocity laws used in the study of elliptic curves and $L$-functions.

\section*{4. Important theorems and results}
\begin{enumerate}[leftmargin=*,itemsep=0.35em]

\item \textbf{Hilbert class field theorem.} For a number field $F$, the Galois group of its maximal everywhere-unramified abelian extension is canonically isomorphic to the ideal class group of $F$. \cite{class_field_theory_ref1}

\item \textbf{Artin reciprocity law.} For every finite abelian extension $L/F$, the quotient $C_F/N_{L/F}(C_L)$ is canonically isomorphic to $\operatorname{Gal}(L/F)$. \cite{class_field_theory_ref1,class_field_theory_ref3}

\item \textbf{Existence theorem.} Finite abelian extensions of $F$ correspond to closed finite-index subgroups of the idele class group $C_F$. \cite{class_field_theory_ref1}

\item \textbf{Kronecker--Weber theorem.} Every finite abelian extension of $\mathbb Q$ is contained in a cyclotomic field. \cite{class_field_theory_ref4}

\item \textbf{Chebotarev density theorem.} In a finite Galois extension, prime ideals are distributed among Frobenius conjugacy classes with densities determined by the sizes of those classes. In the abelian case, this turns reciprocity classes into predictable families of primes. \cite{class_field_theory_ref2}
\end{enumerate}

\paragraph{Applications}
\begin{enumerate}[leftmargin=*,itemsep=0.35em]
\item \textbf{Cryptographic protocols using class groups.} Class groups of imaginary quadratic fields provide abelian groups of unknown structure, which is a desirable foundation for cryptographic constructions. The group $\operatorname{Cl}(\mathbb Q(\sqrt{-d}))$ is hard to compute explicitly for large square-free $d$, yet its order and structure can be verified using class field theory's prediction that the Hilbert class field has Galois group isomorphic to $\operatorname{Cl}(\mathbb Q(\sqrt{-d}))$. Cryptosystems based on this hardness include commitment schemes, verifiable delay functions, and group-action-based protocols analogous to Diffie--Hellman: two parties publicly exchange ideal classes and combine them with their secret ideals, arriving at a shared class without revealing the underlying ideals. The security rests on the assumption that computing the class number or finding a principal ideal representation is computationally intractable for appropriately chosen discriminants \cite{class_field_theory_ref1,class_field_theory_ref4}.
\item \textbf{Primality proving via class field theory.} The theory of complex multiplication gives an explicit construction of abelian extensions of imaginary quadratic fields, and the resulting class polynomials can be used to prove that an integer is prime. Specifically, one computes the Hilbert class polynomial $H_d(X)$ for discriminant $-d$ and verifies that a given elliptic curve $E$ over $\mathbb F_p$ has $j$-invariant equal to a root of $H_d(X)$ modulo $p$. Class field theory guarantees that $p$ splits completely in the Hilbert class field $\mathbb Q(\sqrt{-d},E)$ only when $p$ is prime and satisfies the splitting conditions predicted by Artin reciprocity. This method, implemented in software such as APRCL and CM-based primality tests, proves primality of numbers with thousands of digits and is part of the infrastructure used by cryptographic libraries to generate large prime moduli for RSA and Diffie--Hellman key exchange \cite{class_field_theory_ref2,class_field_theory_ref3}.
\item \textbf{Ideal class group and arithmetic dynamics.} The growth of denominators in orbits of rational points under polynomial iteration is governed by height functions whose local components are determined by splitting behavior in abelian extensions. Class field theory predicts how primes ramify or remain inert in these extensions, which in turn controls the canonical height that measures orbit growth. This connection has been used to prove that certain polynomial maps have only preperiodic rational points---a dynamical analogue of the Mordell conjecture---by reducing the height computation to Artin reciprocity and Chebotarev density arguments \cite{class_field_theory_ref1,class_field_theory_ref7}.
\end{enumerate}
\theoryconnections{class field theory}{%
\item Galois theory supplies extension groups, algebraic number theory supplies ideals and units, and local fields supply completions at primes and infinite places.
\item Topological groups and harmonic analysis are needed because the relevant global Galois groups are profinite and the idele class group is topological.
}{%
\item Class field theory helps arithmetic geometry and Diophantine analysis by predicting how primes split in extensions.
\item Its reciprocity map also provides the abelian part of the wider Langlands viewpoint, where arithmetic information is related to representations and analytic functions.
}

\section*{Glossary}
\begin{description}[leftmargin=*,style=nextline,itemsep=0.3em]

\item[\textbf{Ideal class group.}] The group that measures the failure of unique factorization into principal ideals in a number field; its size (the class number) counts how many distinct classes of ideals exist.

\item[\textbf{Abelian extension.}] A field extension whose Galois symmetries commute with one another; class field theory provides a complete description of all such extensions from arithmetic data in the base field.

\item[\textbf{Reciprocity map.}] A canonical isomorphism from an arithmetic group (such as the idele class group) to the Galois group of an abelian extension, translating prime-splitting behavior into abstract symmetry.

\item[\textbf{Idele class group.}] The quotient of the idele group of a number field by its diagonal copy of the field's nonzero elements; it is the arithmetic object that class field theory pairs with the abelian Galois group.

\item[\textbf{Frobenius automorphism.}] The Galois automorphism associated to an unramified prime in an extension, which raises residue field elements to the power of the prime's cardinality; it records how that prime acts on the extension.

\item[\textbf{Hilbert class field.}] The maximal abelian extension of a number field that is unramified at every place; its Galois group is isomorphic to the ideal class group of the base field.

\item[\textbf{Ramification.}] The phenomenon in which a prime ideal of the base field fails to remain prime or to have the expected number of factors in an extension; ramified primes are the ones where the extension ``wraps around.''

\item[\textbf{Galois group.}] The group of field automorphisms of an extension that fix the base field; class field theory studies its abelian quotient.

\item[\textbf{Number field.}] A finite extension of $\mathbb Q$, such as $\mathbb Q(i)$; class field theory describes abelian extensions of number fields.

\item[\textbf{Completion.}] Adding limits of Cauchy sequences at a given place, producing a local field; bridges global and local class field theory.

\item[\textbf{Profinite group.}] A topological group that is an inverse limit of finite groups, hence compact, Hausdorff, and totally disconnected.

\item[\textbf{Chebotarev density theorem.}] In a finite Galois extension, Frobenius elements are distributed among conjugacy classes with densities proportional to class sizes.

\end{description}

\section*{7. Sample Questions}
\emph{Exercise reference:} \cite{class_field_theory_ref1}. The cited edition is the source for the exercise set; consult its Exercises section for the official statement and numbering.
\begin{enumerate}[leftmargin=*,itemsep=0.9em]

\item \textbf{Exercise 1.} Compute the ideal class group of $F=\mathbb{Q}(\sqrt{-5})$. Specifically, show that the ideal $(2,1+\sqrt{-5})$ is non-principal.

\emph{Solution.} The ring of integers is $\mathbb{Z}[\sqrt{-5}]$. The ideal $I=(2,1+\sqrt{-5})$ has norm $N(I)=2$ (since $2$ ramifies: $(2)=(2,1+\sqrt{-5})^2$). If $I$ were principal, say $I=(\alpha)$, then $N(\alpha)=N(I)=2$, so $a^2+5b^2=2$ for $\alpha=a+b\sqrt{-5}$. This has no integer solutions (since $a^2+5b^2\ge 5$ for $|b|\ge 1$ and $a^2\le 2$ gives $a=0,\pm 1$, none yielding $a^2+5b^2=2$). So $I$ is non-principal, and the class number is at least 2. In fact $h(\mathbb{Q}(\sqrt{-5}))=2$.

\item \textbf{Exercise 2.} Use the Kronecker--Weber theorem to show that $\mathbb{Q}(\sqrt{5})$ is contained in a cyclotomic field. Find the smallest such cyclotomic field.

\emph{Solution.} By Kronecker--Weber, every abelian extension of $\mathbb{Q}$ is contained in $\mathbb{Q}(\zeta_n)$ for some $n$. We need $\sqrt{5}\in\mathbb{Q}(\zeta_n)$. The discriminant of $\mathbb{Q}(\sqrt{5})$ is $5$, and by class field theory, $\mathbb{Q}(\sqrt{5})\subset\mathbb{Q}(\zeta_5)$ since the quadratic subfield of $\mathbb{Q}(\zeta_5)$ is $\mathbb{Q}(\sqrt{5})$ (the unique quadratic subfield of $\mathbb{Q}(\zeta_p)$ for an odd prime $p$ is $\mathbb{Q}(\sqrt{(-1)^{(p-1)/2}p})$, which for $p=5$ gives $\mathbb{Q}(\sqrt{5})$). The smallest cyclotomic field is $\mathbb{Q}(\zeta_5)$, of degree 4 over $\mathbb{Q}$.

\item \textbf{Exercise 3.} Determine the Frobenius element $\left(\frac{\mathbb{Q}(\sqrt{5})/\mathbb{Q}}{3}\right)$ for the prime $p=3$ in the extension $\mathbb{Q}(\sqrt{5})/\mathbb{Q}$.

\emph{Solution.} We determine how $3$ splits in $\mathbb{Q}(\sqrt{5})$ by computing the Legendre symbol $\left(\frac{5}{3}\right)$. Since $5\equiv 2\pmod{3}$ and $\left(\frac{2}{3}\right)=-1$, the prime $3$ is inert in $\mathbb{Q}(\sqrt{5})$. The Frobenius element is therefore the nontrivial element of $\operatorname{Gal}(\mathbb{Q}(\sqrt{5})/\mathbb{Q})\cong\mathbb{Z}/2\mathbb{Z}$, i.e., the conjugation $\sqrt{5}\mapsto -\sqrt{5}$.

\item \textbf{Exercise 4.} For the extension $\mathbb{Q}(\zeta_7)/\mathbb{Q}$, find the decomposition of the prime $p=2$ by computing the order of $2$ modulo $7$.

\emph{Solution.} The prime $p=2$ is unramified in $\mathbb{Q}(\zeta_7)/\mathbb{Q}$. Its Frobenius element is determined by $2\bmod 7$, and the order of the Frobenius in $\operatorname{Gal}(\mathbb{Q}(\zeta_7)/\mathbb{Q})\cong(\mathbb{Z}/7\mathbb{Z})^\times$ equals the order of $2$ in $(\mathbb{Z}/7\mathbb{Z})^\times$. Compute: $2^1=2$, $2^2=4$, $2^3=8\equiv 1\pmod 7$. So the order is 3. The prime $2$ splits into $[L:\mathbb{Q}]/3=6/3=2$ primes, each with inertia degree $3$.

\item \textbf{Exercise 5.} Compute the class number of $\mathbb{Q}(\sqrt{-1})=\mathbb{Q}(i)$ by showing every ideal in $\mathbb{Z}[i]$ is principal.

\emph{Solution.} The ring $\mathbb{Z}[i]$ is a Euclidean domain with norm $N(a+bi)=a^2+b^2$: for any $\alpha,\beta\in\mathbb{Z}[i]$ with $\beta\ne 0$, there exist $q,r\in\mathbb{Z}[i]$ with $\alpha=q\beta+r$ and $N(r)<N(\beta)$. This follows from rounding $\alpha/\beta$ to the nearest Gaussian integer. Every Euclidean domain is a PID, and every PID has trivial class group. Therefore $h(\mathbb{Q}(i))=1$.

\end{enumerate}
\section*{8. References}


\begin{thebibliography}{9}
\bibitem{class_field_theory_ref1} J. Neukirch, \emph{Class Field Theory}, Springer, 1986.
\bibitem{class_field_theory_ref2} J. W. S. Cassels and A. Fröhlich, eds., \emph{Algebraic Number Theory}, Academic Press, 1967.
\bibitem{class_field_theory_ref3} J. Neukirch, A. Schmidt, and K. Wingberg, \emph{Cohomology of Number Fields}, Springer, 2nd ed., 2008.
\bibitem{class_field_theory_ref4} S. Lang, \emph{Algebraic Number Theory}, Springer, 2nd ed., 1994.
\bibitem{class_field_theory_ref5} J. W. S. Cassels and A. Fröhlich, \emph{Algebraic Number Theory}, Springer reprint.
\bibitem{class_field_theory_ref6} J. Neukirch, \emph{Algebraic Number Theory}, Springer, 1999.
\bibitem{class_field_theory_ref7} R. Greenberg, "Iwasawa theory for elliptic curves," in \emph{Arithmetic Theory of Elliptic Curves}, Springer.
\bibitem{class_field_theory_ref8} A. Grothendieck, "Brief an G. Faltings," in discussions of anabelian geometry and arithmetic fundamental groups.
\end{thebibliography}


\theory{Cobordism theory}{When should two spaces be regarded as the two ends of one larger space?}{Cobordism compares manifolds by asking whether they jointly form the boundary of a compact manifold one dimension higher. It is a deliberately coarse comparison: two spaces that are not the same may still be connected by a higher-dimensional bridge. The resulting groups and rings can be calculated using surgery, Morse functions, characteristic numbers, and homotopy theory.}{Geometric topology, algebraic topology, manifold classification, surgery theory, generalized cohomology, topological quantum field theory, and mathematical physics.}{Cobordism replaces direct manifold classification by boundary equivalence, characteristic numbers, and controlled higher-dimensional bridges.}

\paragraph{The problem and the historical idea}

An $n$-dimensional manifold is a space that looks like ordinary $\mathbb R^n$ when viewed close to any one point. A surface is a two-dimensional manifold; a solid region is three-dimensional. A manifold with boundary is allowed to look like a half-space near its edge. A compact manifold is one that is bounded and contains all its limiting points. A closed manifold is compact and has no boundary.

Classifying manifolds up to exact sameness is difficult. Two manifolds can have the same dimension and similar local appearance while differing globally in ways that are hard to detect. Cobordism asks a simpler question: can the two manifolds be placed together as the boundary of a compact manifold of one dimension higher? The higher-dimensional object is the cobordism. \cite{cobordism_theory_ref1,cobordism_theory_ref2}

The roots go back to Poincare's attempts to understand homology through manifolds. Lev Pontryagin used the idea in geometric topology, and Rene Thom made the subject systematic. Thom showed that cobordism groups could be studied through homotopy theory and a construction now called the Thom space or Thom complex. This connected geometric questions with algebraic topology and helped establish cobordism as an extraordinary cohomology theory. \cite{cobordism_theory_ref3,cobordism_theory_ref4}

\end{historylist}
\section*{3. Theory}
\subsection*{Core theory}
\paragraph{Definition and elementary examples}

Two closed $n$-manifolds $M$ and $N$ are cobordant if there is a compact $(n+1)$-manifold $W$ whose boundary is their disjoint union:
\[
 \partial W=M\sqcup N.
\]
More precisely, a cobordism is a compact manifold $W$ together with embeddings of $M$ and $N$ into its boundary whose images are disjoint and cover $\partial W$. The notation $(W;M,N)$ records the bridge and its two ends. The cobordism is not required to be connected.

The interval $[0,1]$ is a one-dimensional cobordism between two points. For any closed manifold $M$, the cylinder $M\times[0,1]$ is a cobordism between its two copies at the ends. A circle bounds a disk, so it is cobordant to the empty manifold and is called null-cobordant. Every sphere bounds a disk of one dimension higher. Every orientable closed surface bounds a handlebody, so it is null-cobordant in the unoriented setting.

The pair of pants is a two-dimensional cobordism: one circular boundary component at one end and two circles at the other. It models splitting one loop into two. Reversing the direction models joining two loops. A disjoint union is cobordant to a connected sum; the bridge records the surgery that joins the pieces. These examples show why cobordisms are useful: they describe a controlled change rather than claiming that the endpoints are identical. \cite{cobordism_theory_ref1,cobordism_theory_ref5}

Cobordism is an equivalence relation on closed manifolds. Reflexivity comes from the cylinder, symmetry from turning $W$ around, and transitivity from gluing two bridges along their common boundary. It is much coarser than homeomorphism or diffeomorphism. That coarseness is a strength: exact classification becomes impossible in many dimensions, while cobordism classes can often be calculated.

\paragraph{Unoriented, oriented, and structured versions}

In unoriented cobordism no direction is assigned to the manifolds. One may instead require orientations. The boundary orientation convention makes the two ends appear with opposite induced orientations, so an oriented cobordism from $M$ to $N$ has boundary $M\sqcup(-N)$. The oriented cobordism groups are written $\Omega_n^{SO}$.

More generally, a manifold can carry a $G$-structure: extra data such as an orientation, a complex structure after stabilization, or a framing. The structure on the bridge must restrict to the given structures at its ends. Important choices include $G=O$ for unoriented cobordism, $G=SO$ for oriented cobordism, and $G=U$ for complex cobordism. Piecewise-linear and topological manifolds produce PL and TOP versions, which are generally harder to calculate than smooth versions. \cite{cobordism_theory_ref3,cobordism_theory_ref6}

Disjoint union gives addition. Cartesian product gives multiplication, because an $m$-manifold times an $n$-manifold has dimension $m+n$. Under suitable conditions the groups form a graded ring
\[
 \Omega_*^G=\bigoplus_{n\geq 0}\Omega_n^G.
\]
The grading remembers dimension, and the ring records how geometric operations interact. This is already a useful summary: two manifolds may be compared by adding and multiplying their classes instead of manipulating every point.

\paragraph{Surgery: changing a manifold in a controlled way}

Surgery removes a standard piece and glues in a complementary piece. Suppose an $n$-manifold contains an embedded
\[
 S^p\times D^q,\qquad p+q=n,
\]
where $S^p$ is a $p$-sphere and $D^q$ is a $q$-dimensional disk. Remove its interior and glue in
\[
 D^{p+1}\times S^{q-1}.
\]
Their boundaries are both $S^p\times S^{q-1}$, so the gluing can be made along the common boundary. The result is another $n$-manifold, and the trace of the operation is an $(n+1)$-dimensional cobordism.

For a circle, cutting out two small intervals and joining the exposed endpoints can either leave one circle or produce two circles. For a sphere, removing a cylinder and inserting two disks gives two spheres; removing two disks and inserting a cylinder can give a torus or a Klein bottle depending on the gluing orientation. In higher dimensions surgery changes handles, holes, and sometimes the fundamental group while preserving the cobordism class. \cite{cobordism_theory_ref1,cobordism_theory_ref7}

Surgery is valuable because it turns a global classification problem into a sequence of local modifications. In high-dimensional topology, one first builds a manifold with desired algebraic invariants and then uses surgery to remove unwanted homotopy or homology. The h-cobordism theorem and the broader surgery program use this strategy to classify manifolds under strong hypotheses.

\paragraph{Morse functions and handle decompositions}

A Morse function is a smooth height function whose critical points are nondegenerate: near each critical point it looks like a sum of positive and negative squares. As the level of the function rises, the level surface changes only when it passes a critical point. A critical point of index $p+1$ produces a $p$-surgery on the level manifold.

If $f:W\to[0,1]$ is a Morse function on a cobordism with $f^{-1}(0)=M$ and $f^{-1}(1)=N$, then $W$ is obtained from $M\times[0,1]$ by attaching one handle for each critical point. A handle of index $p+1$ has the form $D^{p+1}\times D^{n-p}$ and realizes the corresponding surgery. Conversely, a handle decomposition produces a suitable Morse function. This Morse--handle correspondence is one of the central practical tools of cobordism theory. \cite{cobordism_theory_ref1,cobordism_theory_ref7}

For example, a three-dimensional cobordism obtained by attaching a one-handle to a sphere can have a torus as its other boundary. The height function sees the attachment as a critical point of index one. Instead of staring at the whole three-dimensional shape, one records the order and indices of the handles.

\paragraph{Cobordism as algebraic topology}

Thom's work identifies cobordism classes with homotopy classes of maps into Thom spaces. The details depend on the chosen structure, but the message is simple: a geometric object can be encoded by a map between standard spaces. This made calculation possible using homotopy theory, characteristic classes, and cohomology operations.

Cobordism is an extraordinary cohomology theory. Ordinary cohomology assigns groups to spaces using singular chains; extraordinary theories have different coefficient objects and detect geometric information that ordinary cohomology may miss. Complex $K$-theory and cobordism are important examples. Cobordism groups appear as coefficient groups of the corresponding generalized homology theory. \cite{cobordism_theory_ref4,cobordism_theory_ref6}

Characteristic numbers provide computable obstructions. A characteristic class is a cohomology class built from the tangent bundle of a manifold. Evaluating products of these classes on the fundamental class gives numbers. If two manifolds are cobordant, appropriate characteristic numbers agree. Conversely, in important settings, a complete list of characteristic numbers determines the cobordism class. This turns a geometric boundary question into arithmetic with integers.

\paragraph{Categories and quantum field theory}

Cobordisms themselves can be treated as morphisms in a category. The objects are closed manifolds; a morphism from $M$ to $N$ is a cobordism $(W;M,N)$. Composition is gluing: if $W$ joins $M$ to $N$ and $W'$ joins $N$ to $P$, glue along $N$ to obtain a cobordism from $M$ to $P$. The common boundary is the interface at which the processes are joined.

This category is the geometric domain of a topological quantum field theory. An Atiyah--Segal type field theory assigns a vector space or other algebraic object to each boundary manifold and a linear map to each cobordism. The pair of pants becomes multiplication or comultiplication, while a cylinder acts as an identity. Because the assignment depends only on the topological cobordism, the resulting physical model ignores distances and angles but retains global topology. \cite{cobordism_theory_ref8,cobordism_theory_ref9}

\paragraph{What the theory depends on and where it is useful}

Cobordism theory depends on manifolds, boundaries, smooth maps, homotopy, homology, characteristic classes, and category theory. It helps Morse theory by translating critical points into handles, helps surgery theory by organizing changes of manifolds, and helps generalized cohomology by supplying geometric representatives for algebraic classes.

In geometric topology it gives a manageable equivalence relation for manifolds that cannot be classified exactly. In algebraic topology it supplies generalized homology and cohomology theories. In mathematical physics it gives the language for topological quantum field theory. In the Atiyah--Singer index theorem and Hirzebruch--Riemann--Roch, cobordism methods helped organize characteristic numbers and index calculations. \cite{cobordism_theory_ref4,cobordism_theory_ref8}

The limitations are instructive. Cobordism forgets much information: two manifolds may be cobordant while having different fundamental groups or geometries. A cobordism also need not be a product, so it need not give a deformation through identical-looking intermediate spaces. More refined theories, including h-cobordism, framed cobordism, and surgery theory, restore some of the information that ordinary cobordism discards.

\section*{4. Important theorems and results}
\begin{enumerate}[leftmargin=*,itemsep=0.35em]

\item \textbf{Equivalence theorem.} Cobordism is an equivalence relation on closed manifolds of a fixed dimension. The cylinder, reversal, and gluing give reflexivity, symmetry, and transitivity. \cite{cobordism_theory_ref1}

\item \textbf{Morse--handle correspondence.} A Morse function on a cobordism gives a handle decomposition, and a handle decomposition comes from a suitable Morse function. Passing a critical point of index $p+1$ performs a $p$-surgery. \cite{cobordism_theory_ref7}

\item \textbf{Thom's theorem.} Cobordism groups can be related to homotopy groups of Thom spaces, allowing geometric classification to be attacked with homotopy theory. \cite{cobordism_theory_ref4}

\item \textbf{Characteristic-number principle.} In major smooth cobordism theories, characteristic numbers provide algebraic invariants that detect and, in appropriate settings, classify cobordism classes. \cite{cobordism_theory_ref3}

\item \textbf{Cobordism-category principle.} Closed manifolds and cobordisms form a category under gluing, which is the input category for topological quantum field theories. \cite{cobordism_theory_ref8}
\end{enumerate}

\theoryapplications
\theoryconnections{cobordism theory}{%
\item Differential topology supplies manifolds and smooth maps, algebraic topology supplies homotopy and homology, and Morse theory turns a manifold into handles attached at critical points.
\item Surgery then changes a manifold in controlled local steps while tracking the algebraic obstruction to success.
}{%
\item Cobordism helps classify manifolds and organize generalized cohomology theories.
\item It also supplies the geometric language behind topological quantum field theory, where a manifold is viewed as a process between boundary manifolds and composition is given by gluing.
}
\section*{Glossary}
\begin{description}[leftmargin=*,style=nextline,itemsep=0.3em]
\item[\textbf{Manifold.}] A space that locally resembles ordinary Euclidean space at every point, such as a surface or a higher-dimensional analogue.
\item[\textbf{Cobordism.}] An equivalence relation on closed manifolds that asks whether two manifolds can together form the boundary of a compact manifold one dimension higher.
\item[\textbf{Surgery.}] A controlled operation that removes a standard piece of a manifold and glues in a complementary piece, producing a new manifold with different topology.
\item[\textbf{Morse function.}] A smooth function on a manifold whose critical points are nondegenerate, so the topology changes in a predictable way as level sets cross each critical point.
\item[\textbf{Characteristic class.}] A cohomology class assigned to a vector bundle that detects geometric properties such as orientability or twisting, computed from the tangent bundle of a manifold.
\item[\textbf{Handle decomposition.}] A description of a manifold built by starting with a simple piece and attaching handles---standard building blocks indexed by their dimension---one for each critical point of a Morse function.
\item[\textbf{Topological quantum field theory.}] A mathematical framework that assigns algebraic data to manifolds and linear maps to cobordisms, encoding physical processes that depend only on topology.
\item[\textbf{Null-cobordant.}] A manifold that is cobordant to the empty manifold, i.e., it bounds a compact manifold of one dimension higher.
\item[\textbf{Graded ring.}] A ring whose elements are organized into additive components indexed by degree, with multiplication respecting the grading.
\item[\textbf{Closed manifold.}] A compact manifold with no boundary.
\end{description}

\section*{7. Sample Questions}
\emph{Exercise reference:} \cite{cobordism_theory_ref1}. The cited edition is the source for the exercise set; consult its Exercises section for the official statement and numbering.
\begin{enumerate}[leftmargin=*,itemsep=0.9em]

\item \textbf{Exercise 1.} Why are two copies of a closed manifold $M$ cobordant?

\emph{Solution.} The cylinder $M\times[0,1]$ is compact and has boundary $M\times\{0\}\sqcup M\times\{1\}$. It is therefore a cobordism between the two copies.

\item \textbf{Exercise 2.} Show that the circle is null-cobordant.

\emph{Solution.} The boundary of the disk $D^2$ is the circle $S^1$. Taking one boundary component to be $S^1$ and the other to be empty gives a cobordism from the circle to the empty manifold.

\item \textbf{Exercise 3.} What does a critical point of Morse index $p+1$ do to a level manifold?

\emph{Solution.} It attaches a handle $D^{p+1}\times D^{n-p}$ to the cobordism. On the level set this removes $S^p\times D^{n-p}$ and inserts $D^{p+1}\times S^{n-p-1}$, which is $p$-surgery.

\item \textbf{Exercise 4.} Why is cobordism easier to classify than diffeomorphism in many dimensions?

\emph{Solution.} Diffeomorphism requires matching the complete smooth structure. Cobordism only asks whether the two manifolds are the boundary of a higher-dimensional bridge, so many differences are ignored. The reduced information permits algebraic invariants and surgery calculations.

\item \textbf{Exercise 5.} How are two cobordisms composed?

\emph{Solution.} Glue the outgoing boundary of the first cobordism to the incoming boundary of the second. The common manifold becomes an internal seam, and the result has the original incoming and outgoing boundaries.

\end{enumerate}
\section*{8. References}


\begin{thebibliography}{9}
\bibitem{cobordism_theory_ref1} J. Milnor, \emph{Topology from the Differentiable Viewpoint}, Princeton University Press, 1997.
\bibitem{cobordism_theory_ref2} A. Hatcher, \emph{Algebraic Topology}, Cambridge University Press, 2002.
\bibitem{cobordism_theory_ref3} J. Milnor and J. Stasheff, \emph{Characteristic Classes}, Princeton University Press, 1974.
\bibitem{cobordism_theory_ref4} R. E. Stong, \emph{Notes on Cobordism Theory}, Princeton University Press, 1968.
\bibitem{cobordism_theory_ref5} W. S. Massey, \emph{A Basic Course in Algebraic Topology}, Springer, 1991.
\bibitem{cobordism_theory_ref6} J. P. May, \emph{A Concise Course in Algebraic Topology}, University of Chicago Press, 1999.
\bibitem{cobordism_theory_ref7} M. W. Hirsch, \emph{Differential Topology}, Springer, 1976.
\bibitem{cobordism_theory_ref8} M. Atiyah, "Topological quantum field theories," \emph{Publications Mathematiques de l'IHES}, 1988.
\bibitem{cobordism_theory_ref9} J. C. Baez and J. Dolan, "Higher-dimensional algebra and topological quantum field theory," \emph{Journal of Mathematical Physics}, 1995.
\end{thebibliography}


\theory{Coding theory}{How can a receiver recover an intended message after symbols are changed, lost, or observed by an adversary?}{Coding theory designs representations of information for a particular purpose. Source coding removes unnecessary repetition to make data smaller; channel coding adds controlled repetition so errors can be detected and corrected; cryptographic coding protects information from adversaries; and line coding turns symbols into physical signals. A good code balances reliability, rate, delay, power, and computational effort.}{Data compression, mobile phones, QR codes, compact discs, hard drives, satellite and deep-space communication, computer networks, cryptography, medical testing, and neural information processing.}{Hamming distance and entropy quantify redundancy, error correction, and the rate at which information can be transmitted reliably.}

\paragraph{The communication problem and its history}
\bookfigure{information_coding_channel}{The same communication diagram explains compression, noise, redundancy, and recovery.}

A message begins as information at a source. An encoder changes it into a string over an alphabet, such as the binary symbols $0$ and $1$. The string travels through a channel, where noise may flip bits, erase symbols, add disturbances, or hide the signal. The decoder receives the damaged string and tries to recover the original message.

Coding theory asks how to design the encoder and decoder so that the result is efficient and reliable. It is studied in mathematics, information theory, electrical engineering, computer science, and linguistics. Its central tension is easy to state: removing redundancy saves space, while adding redundancy makes errors easier to detect. \cite{coding_theory_ref1,coding_theory_ref2}

Claude Shannon's 1948 paper, A Mathematical Theory of Communication, supplied the statistical model and introduced entropy, mutual information, channel capacity, and the noisy-channel coding theorem. Shannon's work explained the best possible rates even before practical codes were known. Richard Hamming developed error-detecting and error-correcting codes at Bell Labs. The binary Golay code followed in 1949. Later developments included the discrete cosine transform for JPEG, MPEG, and MP3, and modern turbo and low-density parity-check codes. \cite{coding_theory_ref2,coding_theory_ref3}

\end{historylist}
\section*{3. Theory}
\subsection*{Core theory}
\paragraph{Four kinds of coding}

Source coding, also called data compression, removes redundancy that is predictable from the source. DEFLATE, used in common compressed files, represents repeated patterns compactly. Lossless compression permits perfect reconstruction; lossy compression, used in images and sound, keeps the important perceptual information while discarding details.

Channel coding adds useful redundancy before transmission or storage. A compact disc uses Reed--Solomon coding and interleaving so that a scratch affects different parts of the encoded word rather than destroying one contiguous block. Mobile phones combat radio noise and fading with channel codes. Modems and the NASA Deep Space Network use codes adapted to their different kinds of disturbances.

Cryptographic coding transforms information so that an unauthorized observer cannot read or alter it. Modern cryptography aims at confidentiality, integrity, authentication, and non-repudiation. It protects passwords, bank cards, electronic commerce, and network connections, and may be secure because of a computational assumption or, as with a one-time pad, secure even against unlimited computing power.

Line coding maps digital symbols to a physical waveform for baseband transmission. Unipolar, polar, bipolar, and Manchester encodings choose voltage changes so that a receiver can distinguish symbols, recover timing, and operate within the limits of the wire or radio channel. \cite{coding_theory_ref1,coding_theory_ref4}

\paragraph{Codes, alphabets, and distance}

Let $\mathcal X$ be the set of possible source messages and let $\Sigma$ be a finite alphabet. A code is a function
\[
 C:\mathcal X\longrightarrow\Sigma^*,
\]
where $\Sigma^*$ denotes all finite strings over $\Sigma$. The string $C(x)$ is the codeword for message $x$, and its length is $l(C(x))$. If a random source symbol $X$ has probability $\mathbb P[X=x]$, the expected code length is
\[
 l(C)=\sum_{x\in\mathcal X}l(C(x))\mathbb P[X=x].
\]

A code is non-singular when different source symbols receive different codewords. It is uniquely decodable when every concatenated sequence of codewords has only one possible splitting into source symbols. It is instantaneous, or prefix-free, when no codeword is a proper prefix of another; then the decoder can decide as soon as a codeword ends. These distinctions matter because a code that is individually unambiguous may still become ambiguous when codewords are placed next to one another. \cite{coding_theory_ref1,coding_theory_ref5}

For error correction, the Hamming distance between two equal-length words is the number of positions in which they differ. The minimum distance $d_{\min}$ of a code is the smallest distance between two distinct valid codewords. If a received word is closer to one valid codeword than to all others, nearest-codeword decoding chooses that word.

Geometrically, codewords are points in a high-dimensional space. Large minimum distance puts them far apart, leaving room for noise to move a received word without crossing into a neighbor's region. The same picture connects coding with sphere packing: each codeword has a ball of correctable errors around it, and the balls must fit without overlapping.

\paragraph{Source coding and entropy}

Source coding tries to represent likely messages with short words and rare messages with longer words. A source's entropy is its average uncertainty. For a discrete source,
\[
 H(X)=-\sum_x \mathbb P[X=x]\log_2\mathbb P[X=x].
\]
If a source is highly predictable, its entropy is small and strong compression is possible. If all symbols are equally likely and independent, there is less redundancy to remove.

The source-coding theorem says that no lossless code can, in the long run, use fewer than approximately $H(X)$ bits per source symbol, while suitable codes can approach this limit. Huffman and arithmetic coding approach the limit for different source models. Run-length coding replaces a long sequence such as 000000 with a symbol saying “zero repeated six times.” The same idea appears in facsimile transmission and file compression.

The entropy bound is not a claim that every individual file has exactly its entropy as length. It is an average statement for long sequences under a chosen probability model. If the model is wrong, the promised savings may disappear. \cite{coding_theory_ref2,coding_theory_ref6}

\paragraph{Channel coding and error correction}

Suppose the message bit $0$ is sent as 000 and $1$ as 111. The two codewords have distance three. If at most one bit flips, the received word is closer to the correct codeword; for example, 101 is decoded as 111. This simple repetition code demonstrates the principle but wastes many bits.

In general, a code with minimum distance $d_{\min}$ detects up to $d_{\min}-1$ errors and corrects up to
\[
 t=\left\lfloor\frac{d_{\min}-1}{2}\right\rfloor
\]
errors by nearest-neighbor decoding. Detection and correction are different goals: detecting an error may be enough if the sender can retransmit, while correction is necessary for one-way deep-space communication.

Interleaving spreads neighboring symbols across different codewords. A scratch on a CD then becomes many small errors, which a Reed--Solomon code can repair. Deep-space links face continuous thermal noise, telephone modems face network noise, and mobile radios face rapid fading. The best code depends on which error pattern is likely. \cite{coding_theory_ref1,coding_theory_ref7}

\paragraph{Algebraic and linear codes}

Algebraic coding theory expresses code properties using algebra. A linear block code is a set of fixed-length words over a finite field such that the sum of two codewords is again a codeword. It is often summarized by parameters $(n,k,d_{\min})$: $n$ is the word length, $k$ is the number of source symbols encoded at once, and $d_{\min}$ is the minimum Hamming distance. Its rate is $k/n$.

Examples include repetition, parity-check, cyclic, BCH, Hamming, Reed--Solomon, Reed--Muller, algebraic-geometric, perfect, and locally recoverable codes. A parity check can detect an odd number of bit errors. Hamming codes add several parity equations so that the pattern of failed checks, called the syndrome, identifies a single bad position.

Perfect codes fill the available error spheres exactly. The useful nontrivial examples include distance-three Hamming codes with parameters $(2^r-1,2^r-1-r,3)$ and the binary $(23,12,7)$ and ternary $(11,6,5)$ Golay codes. The sphere-packing picture also reveals a limitation: as dimension grows, each codeword has many nearby neighbors, so even a small probability of choosing the wrong neighbor can accumulate. \cite{coding_theory_ref1,coding_theory_ref8}

Convolutional codes do not work on isolated blocks only. Each output symbol is a weighted sum of current and earlier input symbols stored in encoder registers. The encoder is often a small circuit of memory and exclusive-or gates. The Viterbi algorithm finds the most likely path through the encoder's state diagram. Convolutional codes have been used in voiceband modems, GSM phones, satellite systems, and military communication. \cite{coding_theory_ref7}

\paragraph{Cryptographic coding, synchronization, and shared channels}

Cryptographic coding assumes that an adversary may listen to or modify a communication. Encryption hides the content; authentication codes reveal whether the content was altered and identify the sender. Public-key systems allow secure exchange without first sharing a secret, while symmetric systems are usually faster once a key is available. Computationally secure systems depend on problems believed to be hard, such as factoring or discrete logarithms. A one-time pad gives information-theoretic security but requires a truly random key as long as the message and must never reuse it. \cite{coding_theory_ref4}

Codes can also solve synchronization problems. A receiver must know where one message ends and the next begins, and must detect a phase shift or timing error. Special patterns allow the receiver to find alignment. Code-division multiple access assigns different, nearly uncorrelated sequences to users; after demodulation, the other users appear approximately as low-level noise, allowing several phones to share one radio channel.

Automatic repeat-request protocols add check bits and ask for retransmission when the checks fail. Internet protocols such as TCP use sequence numbers to distinguish a new packet from a retransmission. Thus coding is not only about correcting a corrupted symbol; it can coordinate the entire conversation between sender and receiver.

\paragraph{Group testing and other applications}

Group testing uses a code to find a small number of unusual items among a large population with as few tests as possible. A test pools several items and returns positive if at least one member has the property. The problem arose during the Second World War when pooled tests were considered for detecting syphilis among soldiers. Modern applications include defective-product screening, medical testing, compressed sensing, and identifying infected samples.

Analog coding studies error correction, compression, and encryption when signals are continuous rather than strings of bits. Neural coding asks how networks of neurons represent sensory information and how the nervous system may compress, detect, and correct signals. These fields show that coding ideas apply to physical and biological systems as well as digital computers. \cite{coding_theory_ref9,coding_theory_ref10}

\paragraph{How coding theory is used in real systems}

When a QR code is photographed at an angle or with a small stain, its added patterns let the reader recover missing modules. A compact disc survives dust because interleaving and Reed--Solomon checks turn a long scratch into a collection of repairable errors. A mobile phone chooses codes suited to noise and fading, then changes its rate as the radio conditions change. A deep-space probe uses strong coding because a retransmission may take hours or years.

Compression and error correction are sometimes combined. A compressed file has little redundancy, so a few errors may make it unusable; a protected packet therefore adds redundancy after compression. The order matters: compressing after adding correction may destroy the patterns that the decoder needs. In storage, the design also balances capacity against lifetime, decoding energy, and the chance that multiple nearby bits fail together.

Coding theory depends on probability, logarithms, linear algebra, finite fields, algorithms, and signal models. It helps information theory by turning entropy and capacity bounds into practical constructions. It helps computer engineering by supplying decoders and storage protocols, cryptography by formalizing secure encodings, and statistics by providing designs for pooled tests and noisy observations. \cite{coding_theory_ref2,coding_theory_ref6}

\section*{4. Important theorems and results}
\begin{enumerate}[leftmargin=*,itemsep=0.35em]

\item \textbf{Source coding theorem.} For a memoryless source, the entropy is the fundamental lower limit on the average number of bits per source symbol for reliable lossless compression, and suitable codes approach this limit for long messages. \cite{coding_theory_ref2}

\item \textbf{Noisy-channel coding theorem.} A noisy channel has a capacity; communication with arbitrarily small error probability is possible below that capacity, while reliable communication is impossible above it. \cite{coding_theory_ref2}

\item \textbf{Distance bound.} A code with minimum distance $d_{\min}$ detects at most $d_{\min}-1$ errors and corrects at most $\lfloor(d_{\min}-1)/2\rfloor$ errors by nearest-codeword decoding. \cite{coding_theory_ref1}

\item \textbf{Hamming sphere-packing bound.} If a binary code corrects $t$ errors, the radius-$t$ error spheres around its codewords must be disjoint; this limits how many messages can be packed into a fixed word length. \cite{coding_theory_ref1}

\item \textbf{Singleton bound.} A code with length $n$, size described by dimension $k$, and minimum distance $d$ satisfies $d\leq n-k+1$ in the standard linear setting. Codes meeting the bound, such as many Reed--Solomon codes, are called maximum-distance separable. \cite{coding_theory_ref8}
\end{enumerate}

\theoryapplications
\theoryconnections{coding theory}{%
\item Linear algebra supplies vector-space codes, algebra supplies finite fields and polynomial constructions, and probability models errors in a communication channel.
\item Information theory supplies entropy and capacity, which say how much compression or reliability is possible before a code is designed.
}{%
\item Coding theory helps digital communication, storage, satellite links, and deep-space transmission by adding structured redundancy.
\item It also supports distributed storage and data integrity: a receiver can detect or repair errors because the valid messages occupy separated positions in a larger space.
}
\section*{Glossary}
\begin{description}[leftmargin=*,style=nextline,itemsep=0.3em]
\item[\textbf{Hamming distance.}] The number of positions in which two equal-length strings differ, used to measure how far apart codewords are.
\item[\textbf{Entropy.}] A measure of the average uncertainty in a random source, expressed in bits; it sets the fundamental limit on lossless compression.
\item[\textbf{Codeword.}] The encoded string produced by applying a code to a source message, typically longer than the original to add redundancy for error protection.
\item[\textbf{Minimum distance.}] The smallest Hamming distance between any two distinct valid codewords in a code; it determines how many errors the code can detect or correct.
\item[\textbf{Linear block code.}] A set of fixed-length words over a finite field in which the sum of any two codewords is again a codeword, described by parameters $(n,k,d)$.
\item[\textbf{Syndrome.}] The pattern of failed parity-check equations that identifies which errors occurred, allowing the decoder to locate and correct faulty positions.
\item[\textbf{Interleaving.}] A technique that redistributes symbols from different codewords across a physical medium so that a burst of adjacent errors is spread out among many codewords.
\item[\textbf{Channel capacity.}] The maximum rate at which information can be transmitted through a noisy channel with arbitrarily small error probability.
\item[\textbf{Prefix-free code.}] A code in which no codeword is a proper prefix of another, enabling instantaneous uniquely decodable decoding.
\item[\textbf{Reed-Solomon code.}] A non-binary cyclic error-correcting code especially effective at correcting burst errors.
\item[\textbf{Source coding.}] The process of compressing a source's output by removing redundancy, also called data compression.
\item[\textbf{Finite field.}] A field with finitely many elements, such as $\mathbb{F}_q$; the structure over which linear block codes are defined.
\end{description}

\section*{7. Sample Questions}
\emph{Exercise reference:} \cite{coding_theory_ref1}. The cited edition is the source for the exercise set; consult its Exercises section for the official statement and numbering.
\begin{enumerate}[leftmargin=*,itemsep=0.9em]

\item \textbf{Exercise 1.} What is the Hamming distance between 101101 and 100001?

\emph{Solution.} Compare positions: they differ in the third and fourth positions, and agree elsewhere. The distance is $2$.

\item \textbf{Exercise 2.} How many errors can a code with minimum distance $7$ detect and correct?

\emph{Solution.} It can detect up to $6$ errors. It can correct up to $\lfloor(7-1)/2\rfloor=3$ errors by nearest-codeword decoding.

\item \textbf{Exercise 3.} Decode 101 using the repetition code $\{000,111\}$.

\emph{Solution.} The distance from 101 to 000 is two, while its distance from 111 is one. Nearest-codeword decoding therefore returns 111, representing the message bit one.

\item \textbf{Exercise 4.} What do the parameters $(n,k,d)$ of a linear block code mean?

\emph{Solution.} The codeword has length $n$, $k$ source symbols are encoded at once, and distinct codewords are at least distance $d$ apart. The rate is $k/n$, while $d$ controls error detection and correction.

\item \textbf{Exercise 5.} Why can interleaving help a Reed--Solomon code correct a scratch on a disc?

\emph{Solution.} A scratch creates a burst of adjacent errors. Interleaving places adjacent stored symbols in different codewords, turning one burst into many smaller errors. Each codeword then stays within the correction ability of the Reed--Solomon decoder.

\end{enumerate}
\section*{8. References}


\begin{thebibliography}{9}
\bibitem{coding_theory_ref1} F. J. MacWilliams and N. J. A. Sloane, \emph{The Theory of Error-Correcting Codes}, North-Holland, 1977.
\bibitem{coding_theory_ref2} C. E. Shannon, "A mathematical theory of communication," \emph{Bell System Technical Journal}, 1948.
\bibitem{coding_theory_ref3} T. M. Cover and J. A. Thomas, \emph{Elements of Information Theory}, Wiley, 2nd ed., 2006.
\bibitem{coding_theory_ref4} J. Katz and Y. Lindell, \emph{Introduction to Modern Cryptography}, CRC Press, 3rd ed., 2020.
\bibitem{coding_theory_ref5} D. E. Knuth, \emph{The Art of Computer Programming, Volume 3: Sorting and Searching}, Addison-Wesley, 1998.
\bibitem{coding_theory_ref6} T. M. Cover, \emph{Elements of Information Theory}, Wiley.
\bibitem{coding_theory_ref7} J. G. Proakis and M. Salehi, \emph{Digital Communications}, McGraw-Hill, 5th ed., 2008.
\bibitem{coding_theory_ref8} R. Hill, \emph{A First Course in Coding Theory}, Oxford University Press, 1986.
\bibitem{coding_theory_ref9} H. Nyquist, "Certain topics in telegraph transmission theory," \emph{Transactions of the AIEE}, 1928.
\bibitem{coding_theory_ref10} W. Bialek, \emph{Biophysics: Searching for Principles}, Princeton University Press, 2012.
\end{thebibliography}


\theory{Cohomology theory}{How can we turn local measurements into algebraic information about a whole space?}{Cohomology assigns groups, rings, and other algebraic objects to spaces. A cohomology class records a pattern that is locally consistent but may fail to come from a simpler global choice. This makes cohomology a language for holes, twists, gluing obstructions, fields, and characteristic numbers.}{Algebraic topology, differential geometry, algebraic geometry, topology of manifolds, gauge theory, physics, and the classification of bundles and global structures.}{Cocycles, coboundaries, and cup products convert local compatibility data into global algebraic invariants.}

\paragraph{The central idea and history}

Suppose a surface looks ordinary in every small neighborhood but has a global twist. A strip made by joining the ends of a rectangle with one half-turn is a Mobius strip. Locally it looks like a rectangle, but globally it is not a cylinder. Cohomology records the obstruction to choosing one consistent description all the way around.

The word comes from homology, which counts holes using chains. Cohomology reverses the arrows: instead of building geometric pieces, it studies functions or measurements on those pieces. The theory grew from topology in the nineteenth and early twentieth centuries and became a central tool in geometry and algebra. Singular cohomology, de Rham cohomology, sheaf cohomology, and generalized cohomology theories now share a common pattern. \cite{cohomology_theory_ref1,cohomology_theory_ref4}

The advantage is computability. A difficult shape is replaced by groups and a product operation. A continuous map between spaces produces a map in the opposite direction on cohomology, so geometric relationships become algebraic equations that can be calculated.

\end{historylist}
\section*{3. Theory}
\subsection*{Core theory}
\paragraph{Cochains, coboundaries, and classes}

A cochain complex is a sequence of groups or vector spaces
\[
 C^0\xrightarrow{d^0}C^1\xrightarrow{d^1}C^2\xrightarrow{d^2}\cdots
\]
in which two consecutive arrows compose to zero:
\[
 d^{n+1}d^n=0.
\]
The map $d$ is called the coboundary operator. An element in $\ker d^n$ is a cocycle; it satisfies the consistency condition. An element in $\operatorname{im}d^{n-1}$ is a coboundary; it came from a simpler object one degree earlier. The $n$th cohomology group is
\[
 H^n(C^\bullet)=\ker(d^n)/\operatorname{im}(d^{n-1}).
\]
Thus cohomology keeps consistent data and identifies two data sets when their difference is an elementary change.

For differential forms, the operator is ordinary differentiation. A form is closed when its derivative is zero and exact when it is the derivative of another form. Every exact form is closed, but a closed form need not be exact if the space has a hole. The quotient of closed forms by exact forms is de Rham cohomology.

On a circle, a periodic function has derivative with integral zero over one complete turn. The form $d\theta$ is closed but has integral $2\pi$, so it is not exact on the circle. Every one-form $g(\theta)d\theta$ differs from an exact form by a constant multiple of $d\theta$. Hence
\[
 H^0(S^1,\mathbb R)\cong\mathbb R,\qquad
 H^1(S^1,\mathbb R)\cong\mathbb R,
\]
and higher groups vanish. The first group records the one independent circular direction. \cite{cohomology_theory_ref2,cohomology_theory_ref3}

\paragraph{Singular cohomology and functoriality}

For a topological space $X$, singular homology uses continuous maps from standard simplices into $X$. The groups of formal sums of these simplices form a chain complex with boundary maps. Singular cochains are homomorphisms from the chain groups to a coefficient group $A$. Dualizing the boundary maps gives coboundary maps, and their kernels modulo images define
\[
 H^n(X;A).
\]
The coefficient group matters. Integer coefficients can detect torsion, while real or rational coefficients may simplify calculations and ignore some finite information.

A continuous map $f:X\to Y$ induces a pullback
\[
 f^*:H^n(Y;A)\longrightarrow H^n(X;A).
\]
The direction reverses because a measurement on $Y$ can be evaluated after sending a point of $X$ to $Y$. Homotopic maps induce the same pullback, so cohomology is a topological invariant rather than an accidental feature of coordinates.

The cohomology groups usually form a graded-commutative ring under the cup product. Classes in degrees $p$ and $q$ multiply to a class in degree $p+q$. The product stores how different holes and measurements interact, which is why the cohomology ring can distinguish spaces with the same groups. \cite{cohomology_theory_ref1,cohomology_theory_ref5}

\paragraph{Computing cohomology}

The Mayer--Vietoris sequence breaks a space into overlapping parts. If $X=U\cup V$, the cohomology of $U$, $V$, and $U\cap V$ fit into a long exact sequence containing the cohomology of $X$. This is the algebraic version of solving a global problem by comparing local pieces.

Relative cohomology $H^n(X,Y;A)$ measures what remains in $X$ after treating a subspace $Y$ as trivial. It fits into a long exact sequence with $H^n(X)$ and $H^n(Y)$. The universal coefficient theorem relates cohomology to homology:
\[
0\to \operatorname{Ext}^1_{\mathbb Z}(H_{n-1}(X;\mathbb Z),A)
\to H^n(X;A)
\to \operatorname{Hom}_{\mathbb Z}(H_n(X;\mathbb Z),A)
\to0.
\]
Over a field, this reduces to the dual of homology as a vector space.

The diagonal in $X\times X$ is another important object. Its cohomology class encodes the identity map and is used in intersection theory and fixed-point formulas. The cap product pairs cohomology classes with homology classes, allowing a measurement to be evaluated on a geometric cycle. \cite{cohomology_theory_ref1,cohomology_theory_ref2}

\paragraph{Poincare duality and characteristic classes}

For a closed oriented $n$-manifold $M$, Poincare duality pairs degree $k$ cohomology with degree $n-k$ homology. Geometrically, a $k$-dimensional measurement corresponds to an $(n-k)$-dimensional complementary cycle. For an oriented surface, the class of a loop is paired with a class recording how another loop intersects it.

Characteristic classes are cohomology classes attached to vector bundles. A vector bundle assigns a vector space to every point in a continuous way, as the tangent planes of a surface do. Stiefel--Whitney, Chern, and Pontryagin classes detect whether a bundle can be oriented, trivialized, or given additional structure. Their products and evaluations give characteristic numbers that obstruct two manifolds or bundles from being equivalent.

These classes connect cohomology to geometry. The Euler class detects zeros of vector fields; Chern classes measure twisting of complex bundles; Pontryagin classes enter the signature theorem. In physics, gauge fields are connections on bundles and their global charges are expressed through characteristic classes. \cite{cohomology_theory_ref3,cohomology_theory_ref6}

\paragraph{Sheaf cohomology and algebraic geometry}

A sheaf assigns local data to each open region and specifies how data restricts to smaller regions. Sections may exist on every small patch but fail to glue consistently on the whole space. Sheaf cohomology measures exactly this failure. Degree zero records global sections; higher groups record increasingly subtle gluing obstructions.

This idea is fundamental in algebraic geometry. A line bundle may be locally described by one formula on each chart, while transition functions on overlaps twist the charts together. Sheaf cohomology determines whether the bundle has global sections and counts them through Euler characteristics. It also appears in complex analysis, where the existence of a global solution to a local analytic problem becomes a cohomology question.

Cohomology of varieties includes singular, de Rham, etale, crystalline, and other theories. They use different kinds of coefficients and different geometric information, but they aim at the same kind of invariant: a computable record of global structure. \cite{cohomology_theory_ref4,cohomology_theory_ref7}

\paragraph{Axioms and generalized theories}

Ordinary cohomology satisfies the Eilenberg--Steenrod axioms: homotopy invariance, exactness, excision, additivity, and a dimension condition. If the dimension condition is removed, generalized or extraordinary cohomology theories appear. K-theory, cobordism, and stable cohomotopy are examples. They detect structure that ordinary cohomology cannot.

Eilenberg--Mac Lane spaces provide a bridge between cohomology and homotopy. A degree-$n$ cohomology class can be represented by a homotopy class of maps into a suitable space $K(A,n)$. This gives a topological meaning to what otherwise looks like an algebraic quotient.

Cohomology theories are also used in group cohomology, Lie algebra cohomology, cyclic and Hochschild cohomology, Galois cohomology, and mathematical physics. The same pattern persists: consistent objects are divided by trivial changes, and the remainder records an obstruction or invariant. \cite{cohomology_theory_ref4,cohomology_theory_ref8}

\paragraph{What the theory depends on and where it is useful}

Cohomology depends on algebraic groups, linear maps, chain complexes, topology, differential forms, and integration. It helps homology by supplying dual measurements, helps homotopy theory by converting maps into algebraic invariants, and helps geometry by expressing twisting and intersection in equations.

In topology, the cohomology ring distinguishes spaces that have identical numbers of holes but different ways for those holes to interact. In differential geometry, de Rham cohomology turns global questions about differential equations into integrals. In algebraic geometry, sheaf cohomology controls global sections and the geometry of bundles. In physics, characteristic classes express quantized charges and obstructions to global gauge choices. The theory is useful precisely because local formulas are easy to write while global consistency is hard.

\section*{4. Important theorems and results}
\begin{enumerate}[leftmargin=*,itemsep=0.35em]

\item \textbf{Cochain-complex definition.} For $d^{n+1}d^n=0$, the quotient $\ker d^n/\operatorname{im}d^{n-1}$ is the $n$th cohomology group. It measures consistent data modulo elementary changes. \cite{cohomology_theory_ref1}

\item \textbf{Homotopy invariance.} Homotopic continuous maps induce the same map on cohomology. Therefore cohomology depends on the homotopy type of a space. \cite{cohomology_theory_ref2}

\item \textbf{Mayer--Vietoris sequence.} The cohomology of a union can be calculated from the cohomology of two pieces and their intersection through a long exact sequence. \cite{cohomology_theory_ref2}

\item \textbf{Poincare duality.} For a closed oriented $n$-manifold, degree $k$ cohomology and degree $n-k$ homology are paired by evaluation on the fundamental class. \cite{cohomology_theory_ref3}

\item \textbf{de Rham theorem.} The cohomology of smooth differential forms on a manifold is naturally isomorphic to its singular cohomology with real coefficients. \cite{cohomology_theory_ref6}
\end{enumerate}

\theoryapplications
\theoryconnections{cohomology theory}{%
\item Algebraic topology supplies chains and boundaries, linear algebra turns them into kernels and quotients, and category theory explains functoriality and long exact sequences.
\item Differential geometry supplies de Rham forms, while algebraic geometry supplies sheaves and derived constructions.
}{%
\item Cohomology helps topology detect holes, helps geometry solve global extension problems, and helps number theory organize arithmetic information.
\item Its functorial maps and duality statements let different theories transfer local calculations into global conclusions.
}
\section*{Glossary}
\begin{description}[leftmargin=*,style=nextline,itemsep=0.3em]
\item[\textbf{Cocycle.}] An element of a cochain complex that is mapped to zero by the coboundary operator, representing locally consistent data that may fail to be globally trivial.
\item[\textbf{Coboundary.}] An element that arises as the image of a simpler object under the coboundary operator, representing a trivial or elementary change that should be quotiented out.
\item[\textbf{Cup product.}] An operation that multiplies cohomology classes in different degrees to produce a class in a higher degree, giving cohomology a ring structure.
\item[\textbf{Sheaf.}] A structure that assigns compatible local data to every open region of a space and specifies how the data restricts to smaller regions.
\item[\textbf{Poincaré duality.}] A pairing between degree-$k$ cohomology and degree-$(n-k)$ homology for a closed oriented $n$-manifold, reflecting the symmetry between a measurement and its complementary geometric cycle.
\item[\textbf{Pullback.}] The map induced on cohomology by a continuous map between spaces, sending a measurement on the target to its composition with the map on the source.
\item[\textbf{de Rham cohomology.}] The cohomology theory built from smooth differential forms, where the coboundary operator is the exterior derivative; it is isomorphic to singular cohomology with real coefficients.
\item[\textbf{Exact form.}] A differential form that is the exterior derivative of another form; every exact form is closed but not conversely.
\item[\textbf{Closed form.}] A differential form whose exterior derivative is zero; closed forms that are not exact represent nontrivial cohomology classes.
\item[\textbf{Long exact sequence.}] An exact sequence of cohomology groups linking the cohomology of a space, a subspace, and the pair.
\item[\textbf{Mayer-Vietoris sequence.}] A long exact sequence computing the cohomology of a union from the cohomology of two overlapping pieces.
\end{description}

\section*{7. Sample Questions}
\emph{Exercise reference:} \cite{cohomology_theory_ref1}. The cited edition is the source for the exercise set; consult its Exercises section for the official statement and numbering.
\begin{enumerate}[leftmargin=*,itemsep=0.9em]

\item \textbf{Exercise 1.} Why is every exact differential form closed?

\emph{Solution.} If $\omega=df$, then $d\omega=d(df)=0$ because applying the exterior derivative twice gives zero. Hence exact forms lie inside the closed forms.

\item \textbf{Exercise 2.} Why is $d\theta$ on a circle not exact?

\emph{Solution.} If it were $df$ for a globally defined periodic function, its integral around the circle would be $f(2\pi)-f(0)=0$. But $\int_0^{2\pi}d\theta=2\pi$. Therefore it is closed but not exact and represents a nonzero class.

\item \textbf{Exercise 3.} What does a pullback do to a cohomology class?

\emph{Solution.} A map $f:X\to Y$ takes a measurement or differential form defined on $Y$ and evaluates it after applying $f$. This produces a class on $X$, giving $f^*:H^n(Y)\to H^n(X)$.

\item \textbf{Exercise 4.} What information does the cup product add?

\emph{Solution.} It records how classes in different degrees combine. Two spaces may have the same cohomology groups but different products, so the ring structure can distinguish their global geometry.

\item \textbf{Exercise 5.} What is a gluing obstruction?

\emph{Solution.} It is local data that is compatible on overlaps but cannot be produced by one global choice. Sheaf cohomology records the obstruction and can show whether a global section, solution, or bundle trivialization exists.

\end{enumerate}
\section*{8. References}


\begin{thebibliography}{9}
\bibitem{cohomology_theory_ref1} A. Hatcher, \emph{Algebraic Topology}, Cambridge University Press, 2002.
\bibitem{cohomology_theory_ref2} J. R. Munkres, \emph{Elements of Algebraic Topology}, Addison-Wesley, 1984.
\bibitem{cohomology_theory_ref3} J. Milnor and J. Stasheff, \emph{Characteristic Classes}, Princeton University Press, 1974.
\bibitem{cohomology_theory_ref4} R. Bott and L. W. Tu, \emph{Differential Forms in Algebraic Topology}, Springer, 1982.
\bibitem{cohomology_theory_ref5} E. H. Spanier, \emph{Algebraic Topology}, Springer, 1966.
\bibitem{cohomology_theory_ref6} F. Warner, \emph{Foundations of Differentiable Manifolds and Lie Groups}, Springer, 1983.
\bibitem{cohomology_theory_ref7} R. Hartshorne, \emph{Algebraic Geometry}, Springer, 1977.
\bibitem{cohomology_theory_ref8} K. S. Brown, \emph{Cohomology of Groups}, Springer, 1982.
\end{thebibliography}


\theory{Continuum theory}{How can we study connected shapes that are too irregular to be treated as smooth curves or surfaces?}{A continuum is a nonempty compact connected metric space, or in a broader version a compact connected Hausdorff space. Continuum theory studies the subcontinua, decompositions, inverse limits, and symmetries of such spaces. It provides a language for wild and fractal-like shapes without requiring coordinates or differentiability.}{Point-set topology, geometric topology, topological dynamics, fractal geometry, inverse limits, and the study of complicated connected shapes.}{Compactness and connectedness, refined by inverse limits and decompositions, describe spaces whose geometry is too wild for coordinates alone.}

\paragraph{Hans Hahn and Stefan Mazurkiewicz}
They established foundational results about connected compact metric spaces and about continuous images of intervals. These results supplied basic models and structural tests for continua.

\paragraph{Gordon T. Whyburn and later continuum theorists}
Whyburn developed decomposition and separation methods for continua. Later researchers used inverse limits and decompositions to construct and analyze connected spaces with wild local behavior.
\end{historylist}
\section*{3. Theory}
\subsection*{Core theory}
\paragraph{Why continua matter}

The interval $[0,1]$, a circle, a disk, a tree-shaped network, and the Sierpinski carpet are all examples of connected shapes. A topologist may care about whether a shape stays connected after a piece is removed, whether every pair of points can be joined by a path, and whether the shape can be decomposed into simpler connected pieces.

Smooth geometry is not enough for many examples. A fractal may have infinitely many branches or no tangent line at many points. Continuum theory studies the shape using only stable topological properties: compactness, connectedness, continuous maps, and the arrangement of subspaces. The subject grew in point-set topology and geometric topology and is also used in topological dynamics, where an interval-like state space evolves under repeated maps. \cite{continuum_theory_ref1,continuum_theory_ref2}

\paragraph{Definition and basic vocabulary}

A continuum is a nonempty compact connected metric space. Compactness means that every sequence has a convergent subsequence with its limit still in the space, while connectedness means that the space cannot be split into two disjoint nonempty open pieces. Some authors use compact connected Hausdorff spaces instead of metric spaces.

A continuum containing more than one point is nondegenerate. A subset $A$ of a continuum $X$ is a subcontinuum when $A$ is itself a continuum with the subspace topology. A planar continuum is a space homeomorphic to a subcontinuum of the Euclidean plane.

The interval is an arc; a circle is a simple closed curve. A continuum is locally connected if each point has arbitrarily small connected neighborhoods. A locally connected continuum is often called a Peano continuum. Local connectedness is stronger than connectedness: the topologist's sine curve is connected but has a complicated failure of local connectedness.

A continuum is homogeneous if for every two points $x,y$ there is a homeomorphism of the whole continuum sending $x$ to $y$. A circle is homogeneous; a usual interval is not because its endpoints have a different local role from interior points. An indecomposable continuum cannot be written as the union of two proper subcontinua. Such spaces demonstrate how far connected sets can be from ordinary curves. \cite{continuum_theory_ref1,continuum_theory_ref3}

\paragraph{Connectedness is not path connectedness}

A path is a continuous map from $[0,1]$ to the space. Path connectedness requires a path between every pair of points. Every path-connected space is connected, but the reverse implication fails. The topologist's sine curve and related examples are connected although some pairs of points cannot be joined by a path.

This distinction matters in applications. A model may have no sudden separation into two regions while still lacking a usable route between two states. Continuum theory keeps these questions separate rather than replacing all kinds of connectedness by one word.

Continuous images preserve the defining properties: the continuous image of a compact space is compact, and the continuous image of a connected space is connected. Therefore a continuous image of a continuum is a continuum. A projection, deformation, or limit map can simplify a shape without destroying its continuum status. \cite{continuum_theory_ref2}

\paragraph{Constructing continua by nested intersections}

One fundamental construction uses a nested sequence
\[
 X_1\supseteq X_2\supseteq X_3\supseteq\cdots
\]
of continua. If every $X_n$ is nonempty, compact, and connected, then
\[
 X=\bigcap_{n=1}^{\infty}X_n
\]
is nonempty, compact, and connected. The intersection may be much more complicated than any one stage.

This construction explains how a fractal can be built. Begin with a connected compact set, remove selected pieces while preserving connectedness or a controlled collection of components, and repeat. The limiting set can have infinitely many holes, branches, or boundary details. The theorem guarantees that the limit remains a continuum when the nested pieces obey the hypotheses.

Products supply another construction. A finite or countable product of continua is a continuum in the standard product topology. The product of two intervals is a square, while infinite products create compact connected spaces with infinitely many coordinates. \cite{continuum_theory_ref1,continuum_theory_ref4}

\paragraph{Inverse limits}

An inverse sequence consists of continua $X_1,X_2,\ldots$ and continuous bonding maps
\[
 f_n:X_{n+1}\longrightarrow X_n.
\]
The inverse limit is the set of compatible choices
\[
 \varprojlim(X_n,f_n)
 =\{(x_1,x_2,\ldots):f_n(x_{n+1})=x_n\text{ for every }n\}.
\]
It is a subset of the product of all coordinate spaces. Compatibility says that the description at a later, more detailed stage agrees with the description seen at an earlier stage.

The inverse limit of a sequence of continua is again a continuum under standard compactness assumptions. This is a powerful way to construct spaces that are not simple subsets of the plane. For example, inverse limits of intervals under a map can produce indecomposable continua and spaces that arise naturally as invariant sets in dynamical systems.

The method also explains approximation. A complicated limit object is represented by a sequence of simpler coordinate spaces and maps. Instead of drawing the limit exactly, one studies the finite stages and checks that the compatibility maps preserve the needed properties. \cite{continuum_theory_ref1,continuum_theory_ref5}

\paragraph{Decompositions, planar continua, and dynamics}

A decomposition divides a continuum into subcontinua. The quotient space records which points have been identified as belonging to one piece. Questions include whether the quotient is again a well-behaved continuum, whether the pieces vary continuously, and whether the original space can be reconstructed from them.

Planar continua include boundaries of regions, dendrites, lakes of Wada, and many fractal sets. A dendrite is a locally connected continuum containing no simple closed curve; it behaves like a tree that may branch infinitely often. The Menger sponge is a three-dimensional universal compact space for many embedding questions, while planar continua expose restrictions caused by the plane.

In topological dynamics, a continuous map $T:X\to X$ moves points of a continuum through time. An invariant continuum satisfies $T(X)=X$ or contains a dynamically preserved subcontinuum. Inverse limits and continua are used to describe attractors, interval maps, and complicated invariant sets. A continuum may remain connected even while the individual trajectories are chaotic. \cite{continuum_theory_ref5,continuum_theory_ref6}

\paragraph{What the theory depends on and where it is useful}

Continuum theory depends on metric spaces, compactness, connectedness, continuous maps, homeomorphisms, product topology, and inverse limits. It helps topological dynamics by describing state spaces and invariant sets, helps shape theory by retaining information when a space is too wild for ordinary homotopy methods, and helps geometric topology by classifying embeddings and decompositions.

The theory is useful whenever the shape is more important than a coordinate formula. It can model a branching physical object, a limiting construction, or an invariant set from a dynamical system. Its limitations are clear: continua need not be smooth, locally connected, path connected, or easy to visualize, and many classification problems remain open. These difficulties are not defects in the definition; they are the reason a general theory is needed.

\section*{4. Important theorems and results}
\begin{enumerate}[leftmargin=*,itemsep=0.35em]

\item \textbf{Continuous-image theorem.} The continuous image of a continuum is a continuum. Compactness and connectedness are both preserved by continuous maps. \cite{continuum_theory_ref2}

\item \textbf{Nested-intersection theorem.} The intersection of a nested nonempty family of compact connected sets is a continuum when the standard compactness hypotheses hold. \cite{continuum_theory_ref1}

\item \textbf{Inverse-limit theorem.} The inverse limit of a sequence of continua with continuous bonding maps is a continuum under the usual nonemptiness and compactness assumptions. \cite{continuum_theory_ref1}

\item \textbf{Product theorem.} A finite or countable product of continua is a continuum in the product topology. \cite{continuum_theory_ref4}

\item \textbf{Peano continuum principle.} Local connectedness gives much stronger control over arcs and neighborhoods than connectedness alone; its failure explains important wild continua. \cite{continuum_theory_ref3}
\end{enumerate}

\theoryapplications
\theoryconnections{continuum theory}{%
\item Point-set topology supplies compactness, connectedness, and continuous maps; metric spaces provide concrete examples and inverse limits provide a way to build complicated continua from simpler ones.
\item Dimension theory and planar topology refine the basic question of how a connected compact set can be shaped.
}{%
\item Continuum theory helps geometric topology and dynamical systems describe invariant compact sets.
\item It also supplies counterexamples that test whether a proposed theorem really needs local connectedness, path connectedness, or a dimension assumption.
}
\section*{Glossary}
\begin{description}[leftmargin=*,style=nextline,itemsep=0.3em]
\item[\textbf{Continuum.}] A nonempty, compact, connected metric space; it is the central object of study in continuum theory.
\item[\textbf{Subcontinuum.}] A subset of a continuum that is itself a continuum with the subspace topology.
\item[\textbf{Peano continuum.}] A locally connected continuum, meaning every point has arbitrarily small connected neighborhoods; these are the well-behaved continua.
\item[\textbf{Indecomposable continuum.}] A continuum that cannot be written as the union of two proper subcontinua, illustrating how complicated a connected compact set can be.
\item[\textbf{Inverse limit.}] A space constructed from a sequence of simpler spaces and bonding maps by collecting all compatible sequences of points, one from each space.
\item[\textbf{Homogeneous continuum.}] A continuum in which any point can be mapped to any other by a global homeomorphism; a circle is homogeneous but an interval is not.
\item[\textbf{Dendrite.}] A locally connected continuum that contains no simple closed curve, behaving like a tree that may branch infinitely often.
\item[\textbf{Locally connected.}] A space in which every point has arbitrarily small connected neighborhoods.
\item[\textbf{Arc.}] A continuum homeomorphic to the closed interval $[0,1]$.
\item[\textbf{Bonding map.}] A continuous map in an inverse sequence connecting successive coordinate spaces.
\end{description}

\section*{7. Sample Questions}
\emph{Exercise reference:} \cite{continuum_theory_ref1}. The cited edition is the source for the exercise set; consult its Exercises section for the official statement and numbering.
\begin{enumerate}[leftmargin=*,itemsep=0.9em]

\item \textbf{Exercise 1.} Why is $[0,1]$ a continuum?

\emph{Solution.} It is nonempty, compact in the usual metric, and connected because it cannot be separated into two disjoint nonempty open subsets.

\item \textbf{Exercise 2.} Why is the rational interval $\mathbb Q\cap[0,1]$ not a continuum?

\emph{Solution.} It is not compact in the subspace topology: a sequence of rationals can converge to an irrational number in $[0,1]$, whose limit is missing. It is also totally disconnected.

\item \textbf{Exercise 3.} What does homogeneity mean for a continuum?

\emph{Solution.} Every point can be moved to every other point by a homeomorphism of the entire space. The circle is homogeneous, but an interval is not because an endpoint cannot be mapped to an interior point while preserving local structure.

\item \textbf{Exercise 4.} What is an inverse-limit point?

\emph{Solution.} It is a sequence $(x_1,x_2,\ldots)$ with $x_n\in X_n$ such that $f_n(x_{n+1})=x_n$ for every $n$. Each coordinate is a consistent view of the same limiting object.

\item \textbf{Exercise 5.} Why can connectedness fail to provide a path?

\emph{Solution.} Connectedness only forbids a separation into two open pieces. A path requires a continuous parametrized route between every pair of points. Wild continua can satisfy the first condition but not the second.

\end{enumerate}
\section*{8. References}


\begin{thebibliography}{9}
\bibitem{continuum_theory_ref1} S. B. Nadler, \emph{Continuum Theory: An Introduction}, Marcel Dekker, 1992.
\bibitem{continuum_theory_ref2} J. R. Munkres, \emph{Topology}, 2nd ed., Prentice Hall, 2000.
\bibitem{continuum_theory_ref3} R. H. Bing, \emph{Concerning Hereditarily Indecomposable Continua}, Pacific Journal of Mathematics, 1951.
\bibitem{continuum_theory_ref4} J. L. Kelley, \emph{General Topology}, Springer, 1975.
\bibitem{continuum_theory_ref5} M. Barge and J. Kennedy, "Continuum theory and topological dynamics," in \emph{Open Problems in Topology}, 1990.
\bibitem{continuum_theory_ref6} W. T. Ingram and W. S. Mahavier, \emph{Inverse Limits}, Springer, 2012.
\end{thebibliography}


\theory{Computational learning theory}{How much data is needed to learn a concept, and when can a concept be learned efficiently?}{Computational learning theory studies what can be inferred from examples and how much evidence is needed. It asks whether a learner can find a hypothesis that is probably approximately correct, how many samples suffice, and what properties of a concept class determine its learnability.}{Machine learning, statistics, data mining, artificial intelligence, computer vision, natural language processing, bioinformatics, and economics.}{Sample complexity, VC dimension, and Rademacher complexity measure how much data a concept class needs before a learner can generalise reliably.}

\paragraph{The learning problem}
A learner receives examples of an unknown function $f$ drawn from a fixed but unknown distribution $\mathcal D$ over an instance space $\mathcal X$. The learner chooses a hypothesis $h$ from a hypothesis class $\mathcal H$. The error of $h$ is the probability that $h(x)\neq f(x)$ when $x$ is drawn from $\mathcal D$. The learner does not know $\mathcal D$ or $f$; it sees only a finite sample. The central question is how large that sample must be so that the chosen $h$ has small error with high probability.

\paragraph{Historical development}
\begin{historylist}
\paragraph{Origins in psychology and philosophy}
The problem of learning from examples is older than its mathematical formulation. William of Ockham suggested in the fourteenth century that the simplest explanation consistent with observations is preferred. In the twentieth century, questionnaires about concept acquisition led to early artificial intelligence systems that tried to infer rules from positive and negative examples.

\paragraph{Valiant's PAC framework}
Leslie Valiant proposed the Probably Approximately Correct model in 1984. He defined learnability as a complexity question: a concept class is learnable when a learner can, from a polynomial number of examples, produce a hypothesis whose error is below a chosen tolerance with high probability. This made learning a topic of computational complexity theory \cite{colt_ref1}.

\paragraph{VC theory}
Vladimir Vapnik and Alexey Chervonenkis introduced the dimension that bears their name in 1971. The VC dimension measures the largest set of points that a hypothesis class can shatter, meaning it can realize every possible labeling of those points. It controls sample complexity: a class with finite VC dimension needs only finitely many examples to learn, regardless of the distribution.

\paragraph{Online learning and boosting}
In the 1990s, online learning models received new attention. The weighted majority algorithm and its variants combined weak learners into strong ones. Robert Schapire and Yoav Freund developed AdaBoost, which provably reduces error by combining classifiers that are only slightly better than random guessing. These results connected computational learning theory with practical machine learning \cite{colt_ref3}.

\paragraph{Contemporary directions}
Modern work extends the PAC framework to distribution-free settings, to classes with infinite VC dimension measured by Rademacher complexity, to agnostic learning where no target concept in the class is assumed, and to statistical learning theory where the emphasis is on generalisation bounds rather than computational efficiency \cite{colt_ref2}.

\end{historylist}
\section*{3. Theory}
\subsection*{Core theory}
\paragraph{The PAC model}
In the PAC model, the learner receives $m$ independent examples $(x_i,f(x_i))$ drawn from $\mathcal D$. It outputs a hypothesis $h\in\mathcal H$. The hypothesis has error $\operatorname{err}(h)=\Pr_{x\sim\mathcal D}[h(x)\neq f(x)]$. The learner is probably approximately correct when for every $\varepsilon,\delta>0$, it produces $h$ with $\operatorname{err}(h)\leq\varepsilon$ with probability at least $1-\delta$, provided the sample size $m$ is large enough.

The sample complexity is the minimum $m$ that guarantees this. For a finite hypothesis class $\mathcal H$ and a realisable setting where the target $f$ belongs to $\mathcal H$, the bound
\[
m\geq\frac{1}{\varepsilon}\left(\ln|\mathcal H|+\ln\frac{1}{\delta}\right)
\]
suffices. It follows from a union bound: each bad hypothesis can be eliminated by seeing a counterexample, and the number of hypotheses is $|\mathcal H|$.

\paragraph{VC dimension}
When $\mathcal H$ is infinite, VC dimension replaces $|\mathcal H|$. A set of points $S=\{x_1,\ldots,x_n\}$ is shattered by $\mathcal H$ when for every subset $T\subseteq S$, there exists $h\in\mathcal H$ such that $h(x_i)=1$ if and only if $x_i\in T$. The VC dimension $\operatorname{VC}(\mathcal H)$ is the largest $n$ for which some set of size $n$ is shattered.

A half-plane in $\mathbb R^2$ has VC dimension three: it can shatter any three non-collinear points but cannot shatter four. The growth function $\Pi_{\mathcal H}(m)$ counts the number of distinct labelings that $\mathcal H$ produces on $m$ points. Sauer's lemma bounds it:
\[
\Pi_{\mathcal H}(m)\leq\sum_{i=0}^{\operatorname{VC}(\mathcal H)}\binom{m}{i}\leq\left(\frac{em}{\operatorname{VC}(\mathcal H)}\right)^{\operatorname{VC}(\mathcal H)}.
\]
This implies that $\mathcal H$ shatters at most a polynomial number of labelings even when it is infinite, so the sample bound generalises:
\[
m\geq C\frac{\operatorname{VC}(\mathcal H)\ln(1/\varepsilon)+\ln(1/\delta)}{\varepsilon^2}
\]
examples suffice for PAC learning \cite{colt_ref2}.

\paragraph{Rademacher complexity}
Rademacher complexity measures the richness of a hypothesis class by how well it can fit random noise. For a sample $S=(x_1,\ldots,x_m)$ and a class $\mathcal H$ of real-valued functions, the empirical Rademacher complexity is
\[
\hat{\mathfrak R}_S(\mathcal H)=\mathbb E_\sigma\left[\sup_{h\in\mathcal H}\frac{1}{m}\sum_{i=1}^m\sigma_i h(x_i)\right],
\]
where $\sigma_i$ are independent Rademacher random variables taking values $\pm1$ with equal probability. A generalisation bound of the form
\[
\operatorname{err}(h)\leq\hat{\operatorname{err}}(h)+2\hat{\mathfrak R}_S(\mathcal H)+\sqrt{\frac{\ln(1/\delta)}{2m}}
\]
holds simultaneously for all $h\in\mathcal H$ with probability at least $1-\delta$. When $\mathcal H$ consists of linear classifiers in $\mathbb R^d$ with bounded norm, $\hat{\mathfrak R}_S(\mathcal H)$ is of order $\sqrt{d/m}$, giving a parametric rate even when the class is infinite \cite{colt_ref4}.

\paragraph{Agnostic learning}
In agnostic learning, the target $f$ may not belong to $\mathcal H$. The best achievable error is
\[
\operatorname{err}^*=\inf_{h\in\mathcal H}\operatorname{err}(h).
\]
A learner is agnostically PAC when it finds $h$ with $\operatorname{err}(h)\leq\operatorname{err}^*+\varepsilon$ with high probability. The sample complexity depends on the pseudo-dimension of $\mathcal H$, a generalisation of VC dimension for real-valued function classes. Agnostic learning is harder than realisable learning: without the assumption that $f\in\mathcal H$, more examples are needed and the hypothesis class must satisfy additional structural conditions such as convexity or finite metric entropy \cite{colt_ref2}.

\paragraph{Boosting}
Boosting combines many weak learners, each slightly better than random, into a single strong learner. AdaBoost maintains a distribution over training examples, focusing on those that the current ensemble misclassifies. After $T$ rounds, the combined hypothesis has error
\[
\operatorname{err}\leq\prod_{t=1}^T2\sqrt{\varepsilon_t(1-\varepsilon_t)},
\]
where $\varepsilon_t$ is the weighted error of the $t$-th weak learner. If each weak learner has error below $1/2-\gamma$, the product decreases exponentially with $T$. Schapire showed that boosting is possible in principle; Freund and Schapire gave a practical algorithm. The connection to margin theory explains why boosting often improves generalisation even when training error reaches zero \cite{colt_ref3}.

\paragraph{Online learning and regret minimization}
In online learning, examples arrive one at a time and the learner must predict before seeing the label. Regret compares the learner's cumulative loss with the best fixed hypothesis in hindsight. The weighted majority algorithm achieves regret $O(\sqrt{T\ln|\mathcal H|})$ over $T$ rounds. The hedge algorithm and follow-the-regularised-leader extend this to infinite strategy sets. Online-to-batch conversion turns an online algorithm with low regret into a batch learner with good generalisation \cite{colt_ref5}.

\paragraph{Computational efficiency}
PAC learning requires the hypothesis to have low error, but it does not require the learner to run in polynomial time. Proper PAC learning requires outputting a hypothesis from $\mathcal H$; improper learning allows outputting any computable function. Some classes are efficiently improperly learnable but not efficiently properly learnable. The relationship between computational and statistical complexity remains a central open question: many classes with finite VC dimension are statistically learnable but not known to be computationally learnable in polynomial time \cite{colt_ref1}.

\paragraph{Applications}
In machine learning, PAC bounds guide model selection and tell practitioners how much data is needed. In natural language processing, VC bounds on recurrent networks help estimate sample requirements. In computational biology, learning theory explains why gene-expression classifiers need many samples when the number of features is large. In economics and game theory, online learning models describe how agents improve strategies through repeated interaction. In computer vision, sample complexity bounds explain the data hunger of deep networks and motivate data augmentation \cite{colt_ref4}.

\paragraph{What the theory depends on and where it is useful}
Computational learning theory depends on probability, combinatorics, linear algebra, optimisation, and computational complexity. It helps statistics quantify generalisation, machine learning guide model choice, and computer science separate tractable from intractable learning problems. Its bounds are often conservative: they hold for the worst case over distributions, so practical learning may need fewer examples than the theory guarantees.

\section*{4. Important theorems and results}
\begin{enumerate}[leftmargin=*,itemsep=0.35em]
\item \textbf{Fundamental theorem of PAC learning.} A finite hypothesis class is PAC learnable if and only if it is finite, and the sample complexity is $\Theta((\ln|\mathcal H|+\ln(1/\delta))/\varepsilon)$ in the realisable case \cite{colt_ref1}.

\item \textbf{VC dimension theorem.} A hypothesis class with VC dimension $d$ is PAC learnable with sample complexity $\Theta((d+\ln(1/\delta))/\varepsilon^2)$ in the realisable case and $\Theta((d+\ln(1/\delta))/\varepsilon^2)$ in the agnostic case \cite{colt_ref2}.

\item \textbf{Sauer's lemma.} The growth function of a class with VC dimension $d$ is bounded by $(em/d)^d$, ensuring that finite VC dimension implies finite sample complexity \cite{colt_ref2}.

\item \textbf{No-free-lunch theorem.} Without assumptions on the target distribution, no learner can guarantee error below $1/2$ for all possible target functions. Every PAC learning result requires either a restricted hypothesis class or a restriction on the data-generating process \cite{colt_ref1}.

\item \textbf{Boosting theorem.} Any weak learner that produces hypotheses with edge $\gamma$ over random guessing can be boosted into a strong learner with arbitrarily small error by running AdaBoost for $O(\ln(1/\varepsilon)/\gamma^2)$ rounds \cite{colt_ref3}.
\end{enumerate}

\theoryapplications
\theoryconnections{computational learning theory}{%
\item Probability supplies concentration inequalities and Rademacher complexity, while statistics provides hypothesis testing and confidence bounds that compare with PAC guarantees.
\item Computational complexity separates tractable from intractable learning problems, and information theory bounds the rate at which a learner can reduce uncertainty.
}{%
\item Computational learning theory helps machine learning estimate how much data is needed, guides model selection, and explains why boosting, regularisation, and margin maximisation work.
\item It also connects to online optimisation, game theory, and control by analysing how agents improve through repeated interaction and feedback.
}
\section*{Glossary}
\begin{description}[leftmargin=*,style=nextline,itemsep=0.3em]
\item[\textbf{PAC learning.}] A framework in which a learner, given enough examples, produces a hypothesis whose error is below a chosen tolerance with high probability---probably approximately correct.
\item[\textbf{VC dimension.}] The largest number of points that a hypothesis class can shatter, meaning it can realize every possible binary labeling of those points; it controls sample complexity.
\item[\textbf{Rademacher complexity.}] A measure of how well a hypothesis class can fit random noise on a given sample, used to derive distribution-dependent generalization bounds.
\item[\textbf{Hypothesis class.}] The set of candidate functions from which a learner selects a hypothesis during training.
\item[\textbf{Sample complexity.}] The minimum number of training examples needed to guarantee that the learner's error is below a specified threshold with a specified confidence.
\item[\textbf{Agnostic learning.}] A setting in which the target function need not belong to the hypothesis class, so the learner can only hope to approximate the best hypothesis in the class.
\item[\textbf{Boosting.}] A technique that combines many weak learners, each only slightly better than random guessing, into a single strong learner with much lower error.
\item[\textbf{Shatter.}] A hypothesis class shatters a set of points if it can realize every possible binary labeling of those points.
\item[\textbf{Generalization.}] The ability of a learner to perform well on unseen data drawn from the same distribution as training data.
\item[\textbf{Weak learner.}] A classifier that performs only slightly better than random guessing; boosting combines many weak learners into a strong one.
\item[\textbf{Growth function.}] A function counting the number of distinct labelings a hypothesis class produces on $m$ points.
\end{description}

\section*{7. Sample Questions}
\emph{Exercise reference:} \cite{colt_ref2}. The cited edition is the source for the exercise set; consult its Exercises section for the official statement and numbering.
\begin{enumerate}[leftmargin=*,itemsep=0.9em]
\item \textbf{Exercise 1.} Compute the VC dimension of the hypothesis class of intervals $[a,b]$ on $\mathbb{R}$.

\emph{Solution.} The VC dimension is $2$. Given two points $x_1<x_2$, every subset is realized: $\varnothing$ by $[3,4]$ (disjoint from both), $\{x_1\}$ by $[x_1,x_1]$, $\{x_2\}$ by $[x_2,x_2]$, and $\{x_1,x_2\}$ by $[x_1,x_2]$. For three points $x_1<x_2<x_3$, the labeling $(+, -,+)$ cannot be realized: any interval containing both $x_1$ and $x_3$ must contain $x_2$ as well.

\item \textbf{Exercise 2.} A finite hypothesis class $|\mathcal H|=128$ is PAC-learnable in the realisable setting. How many examples suffice for $\varepsilon=0.05$ and $\delta=0.01$?

\emph{Solution.} The sample bound is $m\geq\frac{1}{\varepsilon}(\ln|\mathcal H|+\ln(1/\delta))=\frac{1}{0.05}(\ln 128+\ln 100)=20\,(6.772+4.605)=20\times11.377=227.5$. So $m=228$ examples suffice.

\item \textbf{Exercise 3.} Apply Sauer's lemma to bound the growth function $\Pi_{\mathcal H}(10)$ for a hypothesis class with $\operatorname{VC}(\mathcal H)=3$.

\emph{Solution.} Sauer's lemma gives $\Pi_{\mathcal H}(m)\leq\sum_{i=0}^{d}\binom{m}{i}$ where $d=\operatorname{VC}(\mathcal H)=3$. For $m=10$: $\Pi_{\mathcal H}(10)\leq\binom{10}{0}+\binom{10}{1}+\binom{10}{2}+\binom{10}{3}=1+10+45+120=176$. So at most $176$ of the $2^{10}=1024$ possible labelings can be realized.

\item \textbf{Exercise 4.} In AdaBoost, five weak learners each have edge $\gamma=0.1$ over random guessing. What is the upper bound on the combined error after five rounds?

\emph{Solution.} The error bound is $\prod_{t=1}^{5}2\sqrt{\varepsilon_t(1-\varepsilon_t)}$. With $\varepsilon_t=0.5-0.1=0.4$: each factor is $2\sqrt{0.4\times0.6}=2\sqrt{0.24}\approx0.9798$. After five rounds: $(0.9798)^5\approx0.904$. The combined error is at most $\approx 0.904$, which decreases further with more rounds or larger edge.

\item \textbf{Exercise 5.} Compute the empirical Rademacher complexity of $\mathcal H=\{h_1,h_2\}$ on the sample $S=(1,2,3,4,5)$ where $h_1(x)=x$ and $h_2(x)=-x$.

\emph{Solution.} For each Rademacher vector $\sigma=(\sigma_1,\ldots,\sigma_5)$, compute $\max\bigl|\sum_{i=1}^5\sigma_i h(x_i)\bigr|=\max\bigl|\sum\sigma_i x_i\bigr|,\bigl|-\sum\sigma_i x_i\bigr|\bigr|=\bigl|\sum_{i=1}^5\sigma_i x_i\bigr|$. Since $x_i=i$, the maximum over $\sigma$ of $\bigl|\sum\sigma_i i\bigr|$ is $1+2+3+4+5=15$ (achieved when all $\sigma_i=1$). Averaging over all $2^5=32$ choices of $\sigma$: $\hat{\mathfrak R}_S=\frac{1}{5}\,\mathbb{E}\bigl|\sum\sigma_i i\bigr|$. By symmetry $\mathbb{E}\sum\sigma_i i=0$, but $\mathbb{E}\bigl|\sum\sigma_i i\bigr|>0$; a direct calculation gives $\hat{\mathfrak R}_S=\frac{1}{5}\cdot\frac{15\cdot\binom{4}{2}}{2^4}=\frac{1}{5}\cdot\frac{90}{16}\approx1.125$. (The exact value depends on the distribution of $\sum\sigma_i i$.)

\end{enumerate}
\section*{8. References}

\begin{thebibliography}{9}
\bibitem{colt_ref1} L. G. Valiant, "A theory of the learnable," \emph{Communications of the ACM}, 27 (1984), 1134--1142.
\bibitem{colt_ref2} S. Shalev-Shwartz and S. Ben-David, \emph{Understanding Machine Learning: From Theory to Algorithms}, Cambridge University Press, 2014.
\bibitem{colt_ref3} R. E. Schapire, "The strength of weak learnability," \emph{Machine Learning}, 5 (1990), 197--227.
\bibitem{colt_ref4} M. Mohri, A. Rostamizadeh, and A. Talwalkar, \emph{Foundations of Machine Learning}, MIT Press, 2018.
\bibitem{colt_ref5} N. Cesa-Bianchi and G. Lugosi, \emph{Prediction, Learning, and Games}, Cambridge University Press, 2006.
\end{thebibliography}


\theory{Control theory}{How can a system reach a desired state despite disturbances, delay, and imperfect measurements?}{Control theory designs inputs that make a dynamical system behave as desired. A controller compares a measured process variable with a reference value, uses the error as feedback, and chooses a corrective action. The design must balance speed, accuracy, stability, energy, constraints, and uncertainty.}{Aircraft and spacecraft, robots, industrial plants, vehicles, power grids, communication systems, economics, biology, artificial intelligence, and any process in which feedback changes behavior.}{State-space models, feedback, and stability criteria connect measured system behavior with inputs designed to regulate it.}

\paragraph{The problem and its history}

A dynamical system changes with time. A room has a temperature, a vehicle has a position and speed, and a spacecraft has an orientation and fuel level. An input can change the system, while disturbances such as wind, friction, noise, or an unexpected load push it away from the desired state.

The controller measures a process variable, compares it with a reference or set point, and forms the error
\[
 e(t)=\text{set point}-\text{measured value}.
\]
It uses this error to generate a control action. A thermostat turns heating on when the measured temperature is below the target. Cruise control changes engine power when the speed differs from the set speed. The aim is not simply to act strongly; excessive correction causes overshoot or oscillation. \cite{control_theory_ref1,control_theory_ref2}

Control systems existed in mechanical form long before the modern theory. James Clerk Maxwell analyzed the stability of centrifugal governors in 1868 and explained how delay can produce self-oscillation. Edward Routh, Charles Sturm, and Adolf Hurwitz developed stability criteria. Nicolas Minorsky developed proportional-integral-derivative control for ship steering in the 1920s. Aircraft control, World War II guidance, the Space Race, robotics, and industrial automation then drove the subject forward. \cite{control_theory_ref3,control_theory_ref4}

\end{historylist}
\section*{3. Theory}
\subsection*{Core theory}
\paragraph{Open-loop and closed-loop control}
\bookfigure{control_feedback}{Feedback compares a desired value with a measured output and corrects the input.}

An open-loop controller gives an input without using the measured result. A washing machine may run a fixed program for a fixed time. Open-loop control is simple, but it cannot correct an unplanned disturbance.

Closed-loop control uses feedback. A sensor measures the output; the controller compares it with the target; and the actuator changes the input. Negative feedback means that the controller acts to reduce the error. A block diagram commonly shows the reference, summing point, controller, plant, sensor, and feedback path.

Feedback can stabilize an unstable process, reject disturbances, and make performance less sensitive to uncertain parameters. It can also cause trouble. A delayed or excessively strong feedback signal can make the system oscillate. Every design therefore asks whether the closed loop is stable and whether it remains stable when the model is imperfect.

\paragraph{Mathematical models and transfer functions}

A simple linear time-invariant model is
\[
 \dot x=Ax+Bu,\qquad y=Cx+Du,
\]
where $x$ is the state, $u$ is the input, and $y$ is the measured output. With state feedback $u=-Kx$, the closed-loop equation becomes
\[
 \dot x=(A-BK)x.
\]
Choosing $K$ changes the eigenvalues of $A-BK$, which determine whether small disturbances decay or grow.

In classical control, the relation between input and output is often described by a transfer function
\[
 G(s)=\frac{Y(s)}{U(s)},
\]
obtained from differential equations using the Laplace transform. The variable $s$ is complex. Poles of $G(s)$ describe natural modes; zeros describe cancellations or frequencies at which the response is suppressed. A block diagram lets engineers combine components by multiplication and feedback formulas.

The model is an approximation. A motor may be linearized near one operating point, while friction, saturation, backlash, and changing loads make the actual system nonlinear. Control theory studies how much performance survives such errors. \cite{control_theory_ref1,control_theory_ref5}

\paragraph{Classical control: SISO and PID}

Classical control usually treats one input and one output, called SISO control. The design can use a step response, a root-locus plot, a Bode plot, a Nyquist plot, or a frequency-response measurement.

The proportional-integral-derivative controller is
\[
 u(t)=K_Pe(t)+K_I\int_0^t e(\tau)\,d\tau+K_D\frac{de}{dt}.
\]
The proportional term reacts to the present error. The integral term accumulates past error and removes steady-state offset. The derivative term anticipates the trend and can reduce overshoot, although it amplifies measurement noise. Real controllers use filters and limits because actuators cannot supply infinite force.

Design specifications include settling time, rise time, percent overshoot, steady-state error, bandwidth, and gain and phase margins. Classical controllers remain common because they are inexpensive, interpretable, and easy to tune on site, even when modern state-space methods could describe the same plant more fully. \cite{control_theory_ref2,control_theory_ref6}

\paragraph{Modern control: state space and MIMO systems}

Modern control represents all important internal variables as a state vector. This is essential when several inputs and outputs interact. A large telescope may have many mirror segments, actuators, and sensors; a nuclear reactor and a biological cell may also require multiple interacting variables.

A system is controllable if inputs can move it from the relevant initial state to a desired state in finite time. It is observable if the internal state can be reconstructed from the measured outputs. For the linear system $(A,B)$, controllability is tested by the rank of
\[
 [B\;AB\;A^2B\;\cdots\;A^{n-1}B].
\]
For $(A,C)$, observability is tested by the rank of
\[
 \begin{bmatrix}C\\CA\\CA^2\\\vdots\\CA^{n-1}\end{bmatrix}.
\]
Without controllability, some modes cannot be moved by the actuator. Without observability, a controller cannot know which internal mode is active.

Observers estimate hidden states from measurements. The Kalman filter combines a model with noisy data and is optimal for a standard linear-Gaussian estimation problem. Pole placement chooses feedback so that the closed-loop eigenvalues have desired locations, provided the system is controllable. \cite{control_theory_ref7,control_theory_ref8}

\paragraph{Frequency-domain and time-domain analysis}

Frequency-domain analysis replaces differential equations by algebraic relations in frequency. Laplace, Fourier, and Z transforms convert continuous or discrete time signals into functions of a complex frequency. For linear systems this simplifies convolution into multiplication and makes resonance, bandwidth, gain, phase, and stability visible in plots.

Time-domain state-space analysis keeps the state as a function of time and works naturally with multiple inputs, nonzero initial conditions, nonlinear terms, and simulations. Frequency methods are especially convenient for linear SISO design; state-space methods are more flexible for modern MIMO and nonlinear systems. Engineers often use both: a simulation checks the transient response, while a frequency plot checks robustness and margins.

\paragraph{Stability and robustness}

Stability means that disturbances do not make the system diverge. Lyapunov stability asks whether trajectories that start close remain close; asymptotic stability additionally requires them to approach the desired state. Bounded-input bounded-output stability requires every bounded input to produce a bounded output for a causal system.

For a continuous-time linear system, poles must lie in the open left half of the complex plane. For a discrete-time system, poles must lie inside the unit circle. A pole on the boundary can produce a permanent oscillation or marginal behavior. For example,
\[
 x[n]=0.5^n u[n]
\]
has a Z-transform with a pole at $0.5$, so it decays. Replacing $0.5$ by $1.5$ gives a pole outside the unit circle and a growing response.

Robust control designs for uncertainty in the model. A controller should tolerate small changes in mass, delay, friction, sensor calibration, and unmodeled dynamics. H-infinity loop shaping, sliding-mode control, and structured uncertainty methods explicitly trade nominal performance for guaranteed margins. \cite{control_theory_ref2,control_theory_ref9}

\paragraph{Linear, nonlinear, decentralized, and stochastic systems}

Linear methods are powerful because superposition holds: the response to two inputs is the sum of the separate responses. Nonlinear systems violate this rule. Robots and aircraft contain saturation, trigonometric geometry, contact, and aerodynamic effects. Nonlinear control uses Lyapunov functions, feedback linearization, backstepping, sliding modes, trajectory linearization, and differential geometry.

In decentralized control, several controllers coordinate through communication rather than one central computer. This is useful for power networks, fleets of robots, and geographically distributed plants. The communication itself may be delayed or unreliable, creating another feedback loop.

Deterministic control assumes that the model and disturbances are known functions. Stochastic control includes random shocks. It seeks policies that optimize an expected cost or a risk-sensitive criterion. Filtering and stochastic control are central in finance, robotics, navigation, and adaptive decision-making.

\paragraph{Main control strategies}

Optimal control chooses an input that minimizes a cost, such as fuel, time, tracking error, or energy. Pontryagin's maximum principle gives necessary conditions, while dynamic programming leads to Bellman equations. Linear-quadratic-Gaussian control combines a quadratic cost, linear dynamics, Gaussian noise, and Kalman filtering.

Model predictive control repeatedly solves a finite-horizon optimization problem, applies the first action, measures the new state, and solves again. It can include actuator and safety constraints, which explains its importance in chemical plants, engines, and energy systems. Robust control protects against model uncertainty; adaptive control identifies parameters online and changes controller gains.

Hierarchical and networked control arrange devices and software in levels. Intelligent control combines neural networks, Bayesian methods, fuzzy logic, evolutionary algorithms, or machine learning with a dynamic model. These methods can help with poorly known systems, but safety and stability cannot be assumed merely because a prediction is accurate on training data. \cite{control_theory_ref6,control_theory_ref10}

\paragraph{Applications and connections}

Control theory depends on differential equations, linear algebra, transforms, probability, optimization, numerical simulation, and dynamical systems. It helps robotics by turning desired motion into actuator commands, aerospace by stabilizing aircraft and spacecraft, manufacturing by keeping temperature and pressure within limits, and medicine by regulating drug delivery and artificial organs.

Feedback also appears in economics, operations research, biology, and artificial intelligence. A market model responds to prices; a population responds to resources; a learning system changes its parameters from errors. The same mathematical questions recur: what is measured, what can be changed, how fast does the correction act, and what happens when the model is wrong?

\section*{4. Important theorems and results}
\begin{enumerate}[leftmargin=*,itemsep=0.35em]

\item \textbf{Routh--Hurwitz criterion.} The coefficients of a real polynomial determine, through a finite array calculation, whether all roots lie in the open left half-plane. This gives a stability test without solving for every root. \cite{control_theory_ref3}

\item \textbf{Pole-location criterion.} A continuous-time linear transfer function is asymptotically stable when all poles have negative real part; a discrete-time system is stable when all poles lie inside the unit circle. \cite{control_theory_ref2}

\item \textbf{Controllability and observability tests.} The Kalman rank tests determine whether a finite-dimensional linear system can be driven and reconstructed from its inputs and outputs. \cite{control_theory_ref7}

\item \textbf{Separation principle.} Under standard linear assumptions, a controller designed from an estimated state and a Kalman observer can be designed in separate pieces: the feedback and observer poles combine in the closed loop. \cite{control_theory_ref8}

\item \textbf{Bellman principle.} An optimal control policy has optimal remaining decisions after every intermediate state, leading to dynamic programming and the Hamilton--Jacobi--Bellman equation. \cite{control_theory_ref10}
\end{enumerate}

\theoryapplications
\theoryconnections{control theory}{%
\item Differential equations model the plant, linear algebra describes state and output maps, and complex analysis locates poles and zeros.
\item Optimization and probability enter when control effort, noise, uncertainty, and competing performance goals must be balanced.
}{%
\item Control theory helps robotics, aircraft, industrial processes, power systems, and autonomous vehicles keep a desired state despite disturbances.
\item Its stability and observability tests also guide system identification, signal filtering, and the design of feedback in biological and economic models.
}
\section*{Glossary}
\begin{description}[leftmargin=*,style=nextline,itemsep=0.3em]
\item[\textbf{Feedback.}] A control scheme in which the measured output is compared with a reference value and the resulting error drives a corrective input to the system.
\item[\textbf{Transfer function.}] The ratio of the output to the input in the Laplace domain for a linear time-invariant system, encoding poles and zeros that determine stability and response.
\item[\textbf{PID controller.}] A controller that combines proportional, integral, and derivative actions on the error signal to balance speed, accuracy, and steady-state error.
\item[\textbf{Controllability.}] The property that inputs can move a system from any initial state to any desired state in finite time.
\item[\textbf{Observability.}] The property that the internal state of a system can be fully determined from its measured outputs over time.
\item[\textbf{Pole.}] A value of the complex frequency variable at which the transfer function becomes unbounded, determining the natural modes and stability of the system.
\item[\textbf{Lyapunov stability.}] A notion in which trajectories that start close to an equilibrium remain close for all future time, without necessarily approaching it.
\item[\textbf{State-space model.}] A representation of a dynamical system using a vector of internal state variables and first-order differential equations.
\item[\textbf{Steady-state error.}] The residual difference between the desired set point and actual output after transients have decayed.
\item[\textbf{Kalman filter.}] An optimal state estimator for linear-Gaussian systems combining a dynamic model with noisy measurements.
\item[\textbf{Pole placement.}] A design technique choosing feedback so closed-loop eigenvalues have prescribed locations.
\end{description}

\section*{7. Sample Questions}
\emph{Exercise reference:} \cite{control_theory_ref1}. The cited edition is the source for the exercise set; consult its Exercises section for the official statement and numbering.
\begin{enumerate}[leftmargin=*,itemsep=0.9em]

\item \textbf{Exercise 1.} For $\dot x=ax$ with feedback $u=-kx$ and plant equation $\dot x=ax+u$, when is the closed loop asymptotically stable?

\emph{Solution.} The closed loop is $\dot x=(a-k)x$. It decays exactly when $a-k<0$, or $k>a$. Feedback can stabilize an unstable plant with $a>0$ if the actuator affects the state.

\item \textbf{Exercise 2.} What is the purpose of the integral term in a PID controller?

\emph{Solution.} It accumulates the error over time. If a constant disturbance leaves a nonzero steady-state error, the integral term keeps increasing the control action until that offset is removed, subject to actuator saturation.

\item \textbf{Exercise 3.} What does controllability mean in plain language?

\emph{Solution.} It means that the available inputs can move the system from its current state to the desired state, at least within the model, in finite time. An unactuated mode makes the system uncontrollable.

\item \textbf{Exercise 4.} Why might a system be observable but not controllable?

\emph{Solution.} Sensors may reveal an internal mode while actuators cannot influence it. Observation concerns what can be inferred from outputs; controllability concerns what can be changed by inputs. They are logically different properties.

\item \textbf{Exercise 5.} Why can feedback destabilize a system?

\emph{Solution.} If the correction is delayed, too strong, or based on a noisy measurement, the controller may overcorrect. The resulting closed-loop poles can move outside the stable region, creating growing oscillations rather than reducing the error.

\end{enumerate}
\section*{8. References}


\begin{thebibliography}{9}
\bibitem{control_theory_ref1} K. J. Åström and R. M. Murray, \emph{Feedback Systems: An Introduction for Scientists and Engineers}, Princeton University Press, 2008.
\bibitem{control_theory_ref2} G. F. Franklin, J. D. Powell, and A. Emami-Naeini, \emph{Feedback Control of Dynamic Systems}, Pearson, 8th ed., 2019.
\bibitem{control_theory_ref3} J. C. Maxwell, "On governors," \emph{Proceedings of the Royal Society}, 1868.
\bibitem{control_theory_ref4} E. J. Routh, \emph{A Treatise on the Stability of a Given State of Motion}, Macmillan, 1877.
\bibitem{control_theory_ref5} R. H. Middleton and G. C. Goodwin, \emph{Digital Control and Estimation}, Prentice Hall, 1990.
\bibitem{control_theory_ref6} S. Skogestad and I. Postlethwaite, \emph{Multivariable Feedback Control}, Wiley, 2nd ed., 2005.
\bibitem{control_theory_ref7} R. E. Kalman, "On the general theory of control systems," in \emph{Proceedings of the First International Congress of Automatic Control}, 1960.
\bibitem{control_theory_ref8} R. E. Kalman, "A new approach to linear filtering and prediction problems," \emph{Journal of Basic Engineering}, 1960.
\bibitem{control_theory_ref9} K. Zhou, J. C. Doyle, and K. Glover, \emph{Robust and Optimal Control}, Prentice Hall, 1996.
\bibitem{control_theory_ref10} D. P. Bertsekas, \emph{Dynamic Programming and Optimal Control}, Athena Scientific, 4th ed., 2017.
\end{thebibliography}


\theory{Deformation theory}{If a mathematical object is nudged slightly, which nearby objects are genuinely different, and which changes are merely a change of coordinates?}{Deformation theory studies families of solutions that vary with a small parameter. It first finds infinitesimal changes, then identifies obstructions to extending them to actual families, and finally organizes all families in moduli spaces. The same pattern applies to shapes, equations, algebras, representations, bundles, and arithmetic objects.}{Algebraic geometry, complex geometry, singularity theory, moduli spaces, mathematical physics, representation theory, arithmetic geometry, and perturbation problems.}{Infinitesimal deformations and obstruction classes determine which structures can change and which changes are only coordinate effects.}

\paragraph{The question and the historical development}

Imagine a rigid triangle, a flexible surface, or a polynomial curve whose coefficients are changed slightly. Some changes produce a genuinely new object; others are only a relabeling of the old one. Deformation theory separates these possibilities.

Write a family as $P_\varepsilon$, where $\varepsilon$ is small. Keeping only first-order terms means that $\varepsilon^2$ and higher powers are ignored. The resulting linear equations describe infinitesimal deformations. But a first-order direction may fail to extend to second order, or may extend in many different ways. The obstruction to continuing the deformation is often a cohomology class.

The ideas have roots in perturbation theory, geometry, and mechanics. Kodaira and Spencer developed a systematic theory for complex manifolds. Grothendieck recast deformations in algebraic geometry using flat maps and universal families. Gerstenhaber studied deformations of algebras, and later work connected deformation theory with Hochschild cohomology, quantum groups, and string theory. \cite{deformation_theory_ref1,deformation_theory_ref2}

\end{historylist}
\section*{3. Theory}
\subsection*{Core theory}
\paragraph{Infinitesimal changes and trivial changes}

Suppose an object is defined by equations $F(x)=0$. A deformation changes the equations to
\[
 F(x)+\varepsilon G(x)+O(\varepsilon^2)=0.
\]
Substituting a moving solution and comparing coefficients of $\varepsilon$ gives a linearized condition on $G$ and the first-order motion of $x$. The solutions of this linear condition form the tangent space to the deformation problem.

Some solutions are trivial. If a coordinate change moves the original object without changing its isomorphism class, its derivative represents a direction that should be quotiented out. The true tangent space is therefore “allowed first-order changes modulo changes of coordinates.” This is why deformation theory naturally produces quotient spaces and cohomology groups.

For a plane curve, changing coefficients can smooth a singular point, move it, or leave the isomorphism class unchanged. A deformation of $y^2=x^n$ such as
\[
 y^2-x^n+\varepsilon=0
\]
replaces the singular central fiber by smoother nearby fibers. The central object and the nearby objects are different fibers of one family.

\paragraph{Flat families and universal deformations}

The most general geometric formulation is a flat map
\[
 f:\mathcal X\longrightarrow S.
\]
The fiber over a point $s\in S$ is the object being studied at parameter $s$. Flatness is the algebraic condition that prevents components from suddenly appearing or disappearing in an uncontrolled way; it is the correct replacement for “the family changes continuously” in algebraic geometry.

A universal family $\mathfrak X\to B$ is a family from which every nearby deformation is obtained by pullback along a unique map from the parameter space to $B$. If uniqueness holds only up to an allowed equivalence, the family is versal rather than universal. Hilbert schemes and Quot schemes often supply parameter spaces, while moduli spaces identify points corresponding to isomorphic objects. \cite{deformation_theory_ref1,deformation_theory_ref3}

The universal picture explains why a deformation problem is not just one calculation. The base $B$ records all nearby choices, singular points of $B$ record obstructed or singular moduli behavior, and quotienting by symmetries prevents counting the same object repeatedly.

\paragraph{Complex manifolds and curves}

For a complex manifold $X$, the Kodaira--Spencer theory identifies first-order deformations of its complex structure with a sheaf cohomology group
\[
 H^1(X,\Theta_X),
\]
where $\Theta_X$ is the holomorphic tangent bundle. The obstruction lies in a related degree-two group, often $H^2(X,\Theta_X)$.

For a Riemann surface, the obstruction group vanishes for dimension reasons. The Riemann sphere has no nontrivial deformations: its complex structure is isolated. A genus-one curve has a one-parameter family of complex structures. Every elliptic curve can be written as
\[
 y^2=x^3+ax+b,
\]
but different pairs $(a,b)$ can describe isomorphic curves. The invariant $j$ records the true isomorphism class, while coefficient changes may be mere coordinate changes.

For a curve of genus $g\geq2$, Serre duality identifies the deformation tangent space with the dual of the space of holomorphic quadratic differentials. Riemann--Roch gives dimension $3g-3$ for the Teichmuller space of smooth curves. This is a concrete example of cohomology measuring the number of genuine shape changes. \cite{deformation_theory_ref2,deformation_theory_ref4}

\paragraph{Singularities and analytic algebras}

A germ of an analytic space near a point can be represented by a quotient
\[
 A=\mathbb C\{x_1,\ldots,x_n\}/I,
\]
where the braces denote convergent power series. A deformation is a flat family over a base germ whose special fiber is $A$. For the plane-curve singularity $y^2-x^n=0$, adding a parameter $s$ gives
\[
 y^2-x^n+s=0.
\]
The fiber at $s=0$ is singular; nonzero fibers may be smooth. Such a nearby smooth fiber is a Milnor fiber.

The number and shape of deformations can be organized by a deformation functor. Instead of immediately constructing a geometric moduli space, one assigns to each local Artin algebra the set or groupoid of deformations over that algebra. Tangent vectors correspond to first-order extensions by dual numbers, and obstruction classes determine whether extensions to larger nilpotent bases exist. Schlessinger's criteria and cotangent complexes provide general tests and controlling objects. \cite{deformation_theory_ref1,deformation_theory_ref5}

\paragraph{Applications in geometry, arithmetic, and physics}

Mori's bend-and-break uses deformations of curves on algebraic varieties. A curve is moved through a family until it breaks into components, reducing degree or genus and helping prove the existence of rational curves on Fano varieties.

Arithmetic deformation theory asks how a variety over $\mathbb F_p$ can lift to $\mathbb Z_p$ or to higher powers $\mathbb Z/p^n\mathbb Z$. For smooth curves, vanishing obstruction groups can allow compatible lifts through all orders, producing a formal scheme. The Serre--Tate theorem says that deformations of an abelian scheme are controlled by its $p$-divisible group.

Galois deformation theory starts with a representation
\[
 \rho:G\longrightarrow\operatorname{GL}_n(\mathbb F_p)
\]
and asks whether it can be lifted to a representation over $\mathbb Z_p$ or another coefficient ring. This is central in modern number theory and in the study of modularity.

In deformation quantization, a commutative algebra of classical observables is deformed into a noncommutative algebra of quantum observables. In string theory and mathematical physics, deformations of algebras and geometric structures are organized by Hochschild cohomology and related higher operations. \cite{deformation_theory_ref6,deformation_theory_ref7}

\paragraph{What the theory depends on and where it is useful}

Deformation theory depends on calculus, linearization, algebraic geometry, sheaf cohomology, category theory, and moduli. It helps moduli theory by describing local parameter spaces, singularity theory by smoothing or classifying singular objects, and arithmetic geometry by lifting structures across rings of increasing precision.

Its fundamental warning is that an infinitesimal solution is not automatically an actual family. The first-order equations are necessary, but higher-order compatibility can fail. This is the reason obstruction groups and universal properties are as important as tangent spaces.

\section*{4. Important theorems and results}
\begin{enumerate}[leftmargin=*,itemsep=0.35em]

\item \textbf{Kodaira--Spencer principle.} First-order deformations of a complex manifold are controlled by $H^1(X,\Theta_X)$, with obstructions in a related degree-two group. \cite{deformation_theory_ref2}

\item \textbf{Riemann-surface dimension result.} The deformation space of a smooth genus-$g$ curve has dimension $0$ for $g=0$, dimension $1$ for $g=1$, and $3g-3$ for $g\geq2$. \cite{deformation_theory_ref4}

\item \textbf{Flat-family principle.} Flat maps provide algebraic families whose fibers vary without uncontrolled changes in dimension or multiplicity. \cite{deformation_theory_ref1}

\item \textbf{Schlessinger principle.} Under suitable finiteness and lifting conditions, a deformation functor has a hull or pro-representing object that records all infinitesimal deformations. \cite{deformation_theory_ref5}

\item \textbf{Serre--Tate theorem.} Deformations of an abelian scheme are controlled by deformations of its associated $p$-divisible group. \cite{deformation_theory_ref6}
\end{enumerate}

\theoryapplications
\theoryconnections{deformation theory}{%
\item Algebraic geometry supplies schemes and flat families, differential geometry supplies tangent spaces and vector fields, and cohomology records infinitesimal changes and obstructions.
\item Category theory packages the deformation problem as a functor and makes universal objects precise.
}{%
\item Deformation theory helps moduli problems, algebraic geometry, representation theory, and mathematical physics distinguish genuine parameters from changes caused by a coordinate choice.
\item Obstruction classes tell when a first-order construction can be extended to an actual family.
}
\section*{Glossary}
\begin{description}[leftmargin=*,style=nextline,itemsep=0.3em]
\item[\textbf{Infinitesimal deformation.}] A first-order change to a mathematical object, described by a small parameter while ignoring higher-order terms, representing a tangent direction in the space of possible deformations.
\item[\textbf{Flat family.}] A family of objects varying over a parameter space in which the algebraic structure of each fiber is controlled and no components appear or disappear unexpectedly.
\item[\textbf{Obstruction class.}] A cohomology class that measures whether a first-order deformation can be extended to the next order; it vanishes exactly when the extension is possible.
\item[\textbf{Kodaira--Spencer map.}] A map that identifies first-order deformations of a complex manifold's structure with a sheaf cohomology group, with obstructions lying in a higher-degree group.
\item[\textbf{Universal deformation.}] A family of deformations from which every nearby deformation can be obtained by a unique pullback along a map of parameter spaces.
\item[\textbf{Versal deformation.}] Like a universal deformation but with uniqueness only up to allowed equivalences; it still captures all infinitesimal deformations.
\item[\textbf{Moduli space.}] A geometric space whose points parametrize isomorphic classes of a given type of mathematical object, recording all genuinely different configurations.
\item[\textbf{Tangent space.}] The vector space of first-order deformations, whose elements represent allowable leading-order changes.
\item[\textbf{Milnor fiber.}] The nearby smooth fiber of an isolated singularity in a flat family, capturing local topology.
\item[\textbf{Teichmuller space.}] The space parameterizing complex structures on a surface up to diffeomorphisms isotopic to the identity.
\end{description}

\section*{7. Sample Questions}
\emph{Exercise reference:} \cite{deformation_theory_ref1}. The cited edition is the source for the exercise set; consult its Exercises section for the official statement and numbering.
\begin{enumerate}[leftmargin=*,itemsep=0.9em]

\item \textbf{Exercise 1.} What is an infinitesimal deformation?

\emph{Solution.} It is a first-order change described by a parameter $\varepsilon$ while ignoring $\varepsilon^2$ and higher powers. It is a tangent direction to the space of solutions, before checking whether the direction extends to a genuine family.

\item \textbf{Exercise 2.} Why are coordinate changes removed?

\emph{Solution.} A coordinate change may move the equation or picture without changing the underlying object. Deformation theory studies new isomorphism classes, so it quotients out these trivial directions.

\item \textbf{Exercise 3.} What does flatness prevent in a family?

\emph{Solution.} It prevents pathological jumps in the algebraic behavior of fibers, such as components suddenly appearing or disappearing. It formalizes a controlled variation of the object over the parameter space.

\item \textbf{Exercise 4.} What is the role of an obstruction class?

\emph{Solution.} It measures why a first-order or lower-order deformation cannot be extended to the next order. If the obstruction vanishes, an extension may exist; if not, that infinitesimal direction does not integrate.

\item \textbf{Exercise 5.} Why can different equations describe the same elliptic curve?

\emph{Solution.} A change of coordinates can alter the coefficients $a$ and $b$ without changing the isomorphism class. The invariant $j$ removes this redundancy and records the true complex structure.

\end{enumerate}
\section*{8. References}


\begin{thebibliography}{9}
\bibitem{deformation_theory_ref1} R. Hartshorne, \emph{Deformation Theory}, Springer, 2010.
\bibitem{deformation_theory_ref2} K. Kodaira and D. C. Spencer, "On deformations of complex analytic structures," \emph{Annals of Mathematics}, 1958.
\bibitem{deformation_theory_ref3} R. Hartshorne, \emph{Algebraic Geometry}, Springer, 1977.
\bibitem{deformation_theory_ref4} J. Harris and I. Morrison, \emph{Moduli of Curves}, Springer, 1998.
\bibitem{deformation_theory_ref5} M. Schlessinger, "Functors of Artin rings," \emph{Transactions of the American Mathematical Society}, 1968.
\bibitem{deformation_theory_ref6} J. Tate, "Serre--Tate local moduli," in \emph{Algebraic Number Theory}, Academic Press, 1967.
\bibitem{deformation_theory_ref7} M. Gerstenhaber, "On the deformation of rings and algebras," \emph{Annals of Mathematics}, 1964.
\end{thebibliography}


\theory{Dempster--Shafer theory}{How can several sources of incomplete or conflicting evidence be combined without pretending that ignorance is a precise probability?}{Dempster--Shafer theory, also called evidence theory or the theory of belief functions, assigns support to sets of possibilities rather than only to single outcomes. It reports a lower value called belief and an upper value called plausibility. The gap between them records what the evidence has not decided.}{Sensor fusion, diagnosis, decision support, intelligence analysis, missing data, reliability engineering, possibility theory, and imprecise probability.}{Mass assignments, belief, plausibility, and combination rules preserve both supporting evidence and unresolved ambiguity.}

\paragraph{The problem and its history}

Suppose a detector says that a signal is either red or yellow but cannot distinguish which. Ordinary probability forces us to split the evidence between red and yellow, even if no reason for the split exists. Evidence theory keeps the statement as a set-valued piece of information. It distinguishes support for a claim from room left for the claim to be true.

Arthur Dempster introduced upper and lower probabilities in statistical inference. Glenn Shafer developed the framework into a general mathematical theory of evidence. Later work introduced the transferable belief model, theory of hints, alternative fusion operators, random-set interpretations, geometric models, and relational measures. \cite{dempster_shafer_theory_ref1,dempster_shafer_theory_ref2}

The subject is not simply probability with different notation. A probability assigns mass to individual outcomes or events with additivity. Dempster--Shafer theory assigns basic mass to subsets, allowing a source to say “one of these possibilities” without saying how the mass is divided among them.

\end{historylist}
\section*{3. Theory}
\subsection*{Core theory}
\paragraph{Frame of discernment and basic belief assignments}

Let $X$ be the set of possible states. Its power set $2^X$ contains every subset, including the empty set and the whole set. A basic belief assignment is a function
\[
 m:2^X\longrightarrow[0,1]
\]
with
\[
 m(\varnothing)=0,\qquad \sum_{A\subseteq X}m(A)=1.
\]
The sets with positive mass are called focal elements. The mass $m(A)$ supports the statement that the true state lies somewhere in $A$, but does not claim which member of $A$ is true.

For $X=\{\text{alive},\text{dead}\}$, a detector might assign mass $0.2$ to alive, $0.5$ to dead, and $0.3$ to the whole set $X$. The final $0.3$ represents ignorance, not an equal probability of $0.15$ for each state. If later evidence distinguishes the states, that ignorance can be redistributed.

The theory also allows a random-set interpretation. Begin with an ordinary probability space of observations. A set-valued map sends each observation to the subset of states that remain possible. A hidden die whose faces 1 and 2 are indistinguishable produces the set $\{1,2\}$ when that hidden face appears. Missing data in science, weather records, and computer vision naturally produce this type of set-valued observation. \cite{dempster_shafer_theory_ref1,dempster_shafer_theory_ref3}

\paragraph{Belief and plausibility}

For a hypothesis $A\subseteq X$, the belief is
\[
 \operatorname{Bel}(A)=\sum_{B\subseteq A}m(B).
\]
It counts all evidence that directly supports $A$ or an even more specific subset. The plausibility is
\[
 \operatorname{Pl}(A)=\sum_{B\cap A\neq\varnothing}m(B)
 =1-\operatorname{Bel}(X\setminus A).
\]
It counts evidence that does not contradict $A$. Therefore
\[
 \operatorname{Bel}(A)\leq P(A)\leq\operatorname{Pl}(A)
\]
for every ordinary probability $P$ compatible with the evidence.

Return to the cat example. With masses $m(\text{alive})=0.2$, $m(\text{dead})=0.5$, and $m(X)=0.3$,
\[
 \operatorname{Bel}(\text{alive})=0.2,\qquad
 \operatorname{Pl}(\text{alive})=0.2+0.3=0.5.
\]
The evidence supports alive by $0.2$ and leaves at most $0.5$ room for alive. The interval $[0.2,0.5]$ is more honest than a forced number when the detector cannot resolve the situation.

Belief is a lower probability and plausibility is an upper probability. For finite $X$, knowing all masses, all beliefs, or all plausibilities determines the other two through finite sums and inversion. For infinite $X$, a belief and plausibility function may exist even when no ordinary mass function on every subset is convenient. \cite{dempster_shafer_theory_ref2,dempster_shafer_theory_ref4}

\paragraph{Combining evidence}

Suppose two independent sources assign masses $m_1$ and $m_2$. Dempster's rule forms intersections of focal sets. If $B$ is supported by the first source and $C$ by the second, their common claim is $B\cap C$. The conflicting mass is
\[
 K=\sum_{B\cap C=\varnothing}m_1(B)m_2(C).
\]
For a nonempty $A$, the normalized combination is
\[
 (m_1\oplus m_2)(A)
 =\frac{1}{1-K}
 \sum_{B\cap C=A}m_1(B)m_2(C).
\]
The normalization discards conflict and rescales the nonempty intersections.

This rule is appropriate when the masses express compatible belief constraints from independent sources. Other situations need cumulative, cautious, proportional, or disjunctive fusion rules. The choice of rule is a modeling assumption, not a theorem that one formula is always correct.

For two friends choosing among films $X,Y,Z$, Alice may assign $0.99$ to $X$ and $0.01$ to $Y$, while Bob assigns $0.99$ to $Z$ and $0.01$ to $Y$. If the numbers mean preferences and the goal is a film both can agree to watch, the only shared preference is $Y$, so the normalized combination assigns it all the mass. If the sources make completely opposed claims, $K=1$ and the rule is undefined; that is a warning that no common supported conclusion exists. \cite{dempster_shafer_theory_ref2,dempster_shafer_theory_ref5}

\paragraph{Conflict and counterintuitive examples}

Zadeh's example shows why the meaning of a mass assignment must be stated. One doctor assigns mass $0.99$ to brain tumor and $0.01$ to meningitis. Another assigns $0.99$ to concussion and $0.01$ to meningitis. Dempster's rule normalizes the small agreement on meningitis and can assign it mass one. This is counterintuitive if the numbers mean each doctor's confidence in diagnoses: both doctors considered meningitis unlikely. It may be sensible if the numbers describe constraints and the only common conclusion is meningitis, but the interpretation must be explicit.

Even low conflict can surprise. If both doctors assign large mass to brain tumor but disagree slightly about whether the alternative is meningitis or concussion, normalization can give brain tumor complete support. A model that treats belief as proof from evidence, rather than probability that a diagnosis is true, can explain some of these outcomes. Pearl and other authors criticized interpretations that confuse probability of truth with probability of provability, while later work studied which fusion rules are justified for which tasks. \cite{dempster_shafer_theory_ref6,dempster_shafer_theory_ref7}

\paragraph{Relation to Bayesian probability}

Bayesian probability is recovered as a special case when every focal element is a singleton. Then there is no unassigned set-level ignorance, and belief and plausibility agree:
\[
 \operatorname{Bel}(A)=\operatorname{Pl}(A)=P(A).
\]
If masses are assigned to sets such as red-or-green, the theory expresses information that is weaker than a probability distribution on the individual colors.

The advantage is explicit ignorance. The cost is that belief is not generally additive on disjoint sets, and some combinations can violate intuitions or decision-theoretic principles. A Bayesian approximation redistributes the mass of each focal set among its members to obtain a probability distribution. This is useful when a downstream method accepts only singleton probabilities, but it necessarily loses information about ignorance and dependence on the chosen redistribution rule. \cite{dempster_shafer_theory_ref2,dempster_shafer_theory_ref8}

\paragraph{Other uncertainty measures and geometry}

A belief function determines a credal set: the convex set of all ordinary probability distributions $P$ satisfying
\[
 P(A)\geq\operatorname{Bel}(A)\quad\text{for every }A\subseteq X.
\]
The corresponding plausibility is the upper probability over this set. Thus evidence theory is a structured form of imprecise probability. Possibility measures correspond to consonant belief functions, in which the focal sets are nested. Probability boxes represent lower and upper cumulative distributions.

For a finite frame, mass and belief values can be treated as coordinates of a point in a convex belief space. The valid points form a simplex-like geometric region, with ordinary probability distributions appearing as a special face. This picture makes mixtures, constraints, and distance between evidence sources visible.

Relational measures replace the ordinary real-number order by a partial order in a bounded lattice. A measure $\mu$ assigns an ordered value to each subset of criteria, is monotone under inclusion, sends the empty set to the least element, and sends the whole set to the greatest element. This generalizes evidence combination to preferences and multi-criteria reasoning. \cite{dempster_shafer_theory_ref9}

\paragraph{Applications and dependencies}

Dempster--Shafer theory depends on sets, probability, combinatorics, convex geometry, and decision theory. It helps sensor fusion by retaining ambiguity when cameras, radar, and microphones disagree. In diagnosis, it can distinguish evidence directly supporting a disease from evidence that merely fails to rule it out. In intelligence analysis it represents reports that identify a group of possible sources but not one individual.

The theory is useful when data are incomplete, sources report constraints rather than calibrated frequencies, and preserving ignorance matters. It is dangerous when users interpret belief intervals as ordinary confidence intervals or apply Dempster's rule to evidence with the wrong independence semantics. A responsible application states the frame, mass assignment meaning, fusion rule, conflict handling, and final decision criterion.

\section*{4. Important theorems and results}
\begin{enumerate}[leftmargin=*,itemsep=0.35em]

\item \textbf{Belief--plausibility bounds.} For every hypothesis $A$, $\operatorname{Bel}(A)$ is the sum of masses contained in $A$, $\operatorname{Pl}(A)$ is the sum of masses intersecting $A$, and every compatible probability lies between them. \cite{dempster_shafer_theory_ref2}

\item \textbf{Duality.} Plausibility and belief satisfy $\operatorname{Pl}(A)=1-\operatorname{Bel}(X\setminus A)$. \cite{dempster_shafer_theory_ref2}

\item \textbf{Bayesian special case.} If all focal elements are singletons, belief and plausibility coincide with an ordinary probability distribution. \cite{dempster_shafer_theory_ref2}

\item \textbf{Dempster combination rule.} Independent mass assignments combine through nonempty intersections and normalization by $1-K$, where $K$ is the conflicting mass. The rule is undefined when $K=1$. \cite{dempster_shafer_theory_ref1,dempster_shafer_theory_ref5}

\item \textbf{Credal-set representation.} A belief function determines a convex set of compatible probability distributions, so evidence theory can be viewed as a structured lower-and-upper probability model. \cite{dempster_shafer_theory_ref9}
\end{enumerate}

\theoryapplications
\theoryconnections{dempster shafer theory}{%
\item Probability supplies events and conditioning, while set theory allows a mass to be assigned to a group of possible states rather than to one state.
\item Combinatorics computes intersections of focal sets, and decision theory is needed when belief and plausibility lead to an interval rather than one probability.
}{%
\item Dempster--Shafer theory helps sensor fusion and diagnosis combine evidence from sources that are incomplete or partly conflicting.
\item Its interval output preserves uncertainty that an ordinary probability model might hide, although the combination rule must be checked when conflict is large.
}
\section*{7. Sample Questions}
\emph{Exercise reference:} \cite{dempster_shafer_theory_ref1}. The cited edition is the source for the exercise set; consult its Exercises section for the official statement and numbering.
\begin{enumerate}[leftmargin=*,itemsep=0.9em]

\item \textbf{Exercise 1.} Let $X=\{a,b\}$ with $m(\{a\})=0.3$, $m(\{b\})=0.2$, and $m(X)=0.5$. Find belief and plausibility for $\{a\}$.

\emph{Solution.} Belief sums masses contained in $\{a\}$, so $\operatorname{Bel}(\{a\})=0.3$. Plausibility sums masses intersecting $\{a\}$, so it is $0.3+0.5=0.8$.

\item \textbf{Exercise 2.} Why is $m(\{a,b\})=0.5$ not the same as assigning probability $0.25$ to each state?

\emph{Solution.} The set mass says only that the state is one of $a$ or $b$. It does not choose a division. Assigning $0.25$ to each adds an equal-likelihood assumption that the evidence did not provide.

\item \textbf{Exercise 3.} What happens when two sources have conflict $K=1$?

\emph{Solution.} Every pair of focal sets has empty intersection, so the normalizing factor $1-K$ is zero. Dempster's normalized rule is undefined. The correct response is to inspect the model or choose a fusion rule that retains conflict.

\item \textbf{Exercise 4.} When does Dempster--Shafer theory reduce to Bayesian probability?

\emph{Solution.} When all nonzero mass is assigned to singleton states. Then the belief and plausibility of every event are equal, and the mass function is an ordinary probability distribution.

\item \textbf{Exercise 5.} Why must an application specify what the mass means?

\emph{Solution.} The same numerical mass can represent a belief constraint, a probability of provability, a preference, or a sensor reliability. Dempster's normalization can be sensible for one meaning and misleading for another, especially under conflict.

\end{enumerate}
\section*{8. References}


\begin{thebibliography}{9}
\bibitem{dempster_shafer_theory_ref1} A. P. Dempster, "Upper and lower probabilities induced by a multivalued mapping," \emph{The Annals of Mathematical Statistics}, 1967.
\bibitem{dempster_shafer_theory_ref2} G. Shafer, \emph{A Mathematical Theory of Evidence}, Princeton University Press, 1976.
\bibitem{dempster_shafer_theory_ref3} I. Molchanov, \emph{Theory of Random Sets}, Springer, 2017.
\bibitem{dempster_shafer_theory_ref4} K. Sentz and S. Ferson, \emph{Combination of Evidence in Dempster--Shafer Theory}, Sandia National Laboratories, 2002.
\bibitem{dempster_shafer_theory_ref5} A. Jøsang and P. Simon, "Dempster's rule as seen by little colored balls," \emph{Computational Intelligence}, 2012.
\bibitem{dempster_shafer_theory_ref6} J. Pearl, \emph{Probabilistic Reasoning in Intelligent Systems}, Morgan Kaufmann, 1988.
\bibitem{dempster_shafer_theory_ref7} N. Wilson, "The assumptions behind Dempster's rule," in \emph{Proceedings of UAI}, 1993.
\bibitem{dempster_shafer_theory_ref8} F. Voorbraak, "A computationally efficient approximation of Dempster--Shafer theory," \emph{International Journal of Man--Machine Studies}, 1989.
\bibitem{dempster_shafer_theory_ref9} F. Cuzzolin, \emph{The Geometry of Uncertainty}, Springer, 2021.
\end{thebibliography}


\theory{Dimension theory}{How can the size or number of independent directions of an algebraic object be measured when the object is not a simple line, plane, or vector space?}{Dimension theory studies several notions of dimension and the conditions under which they agree. In commutative algebra and algebraic geometry, the main notion is Krull dimension: the maximum length of a chain of prime ideals. It measures how many layers of algebraic subspaces a ring or scheme can contain.}{Algebraic geometry, commutative algebra, singularities, algebraic number theory, module theory, homological algebra, and the comparison of geometric and algebraic complexity.}{Prime chains, local rings, and homological dimensions turn the informal size of an object into computable algebraic quantities.}

\paragraph{Wolfgang Krull}
Krull defined the dimension of a commutative ring through chains of prime ideals. This algebraic definition measures how many successive layers of algebraic subspaces can occur.

\paragraph{Henri Lebesgue and topological dimension theorists}
Lebesgue and later covering-dimension researchers developed ways to assign dimension to general topological spaces. These notions allow algebraic, geometric, and topological dimensions to be compared when their hypotheses overlap.
\end{historylist}
\section*{3. Theory}
\subsection*{Core theory}
\paragraph{Why dimension needs a theory}

For a vector space, dimension is the number of vectors in a basis. For a line, plane, or manifold, it is the number of independent local coordinates. Algebraic sets can be singular, reducible, or infinitely complicated, so these familiar definitions may no longer agree or may not even apply.

Dimension theory studies several questions. How long can a chain of proper algebraic subspaces be? How many independent parameters describe a family? How many equations are needed locally? How complicated is a module's resolution? For regular geometric objects these answers often coincide; for singular or non-Noetherian objects they may differ. \cite{dimension_theory_ref1,dimension_theory_ref2}

\paragraph{Krull dimension and prime chains}

Let $R$ be a commutative ring. A prime ideal is a set of elements closed under addition and multiplication by arbitrary ring elements such that if a product $ab$ lies in the ideal, then $a$ or $b$ lies in it. Prime ideals represent irreducible algebraic conditions.

The Krull dimension of $R$ is the supremum of the lengths of chains
\[
 \mathfrak p_0\subsetneq\mathfrak p_1\subsetneq\cdots\subsetneq\mathfrak p_n
\]
of prime ideals. The height of a prime $\mathfrak p$ is the dimension of the localization $R_{\mathfrak p}$, equivalently the maximum length of a chain ending at $\mathfrak p$.

For a field $k$, the ring $k$ has dimension zero because its only prime ideal is $(0)$. The polynomial ring $k[x]$ has the chain $(0)\subsetneq(x)$ and dimension one. The ring $k[x,y]$ has
\[
 (0)\subsetneq(x)\subsetneq(x,y),
\]
so its dimension is two. Algebraically, the plane has two independent coordinates and two layers of prime conditions.

\paragraph{Polynomial rings, varieties, and chains}

For a Noetherian ring $R$, adjoining one polynomial variable raises the dimension:
\[
 \dim R[x]=\dim R+1.
\]
Thus a polynomial ring in $n$ variables over a field has dimension $n$. Quotients lower dimension according to the equations imposed. A hypersurface defined by one nonzero equation in affine $n$-space usually has dimension $n-1$, although singularities and components require care.

The prime spectrum $\operatorname{Spec}R$ turns the set of prime ideals into a geometric space. Inclusion of primes becomes specialization: a larger prime represents a more constrained point. Krull dimension is then the dimension of the longest chain of geometric strata. Localization focuses on a neighborhood, while quotienting by a prime restricts to a closed irreducible subspace.

For finitely generated algebras over a field, algebraic, geometric, and transcendence dimensions agree in the expected cases. The dimension of an irreducible variety is the transcendence degree of its function field over the ground field. \cite{dimension_theory_ref1,dimension_theory_ref3}

\paragraph{Noetherian rings and principal ideal bounds}

A ring is Noetherian when every ascending chain of ideals eventually stops, or equivalently when every ideal has a finite generating set. This condition makes dimension theory tractable. Polynomial rings over Noetherian rings remain Noetherian by the Hilbert basis theorem.

Krull's principal ideal theorem says that a minimal prime over a principal ideal generated by a nonzero nonunit has height at most one. More generally, a minimal prime over an ideal generated by $r$ elements has height at most $r$ under standard Noetherian hypotheses. This formalizes the geometric expectation that one equation cuts dimension by at most one.

A ring is catenary when saturated prime chains between two fixed primes have the same length. Catenarity ensures that dimension behaves predictably under passage between strata. Many finitely generated algebras over fields are catenary, while arbitrary non-Noetherian rings may violate familiar dimension formulas. \cite{dimension_theory_ref2,dimension_theory_ref4}

\paragraph{Regularity, singularities, and local dimension}

At a point of an algebraic variety, the local ring records functions defined near that point. A regular local ring has dimension equal to the dimension of its cotangent or tangent space. If the tangent space is larger than the local Krull dimension, the point is singular: too many first-order directions meet there.

For example, a crossing such as $xy=0$ has two branches meeting at the origin. The local ring has dimension one, but the tangent space has two independent directions, signaling a singularity. A smooth curve has matching dimensions and tangent directions.

Regular rings are important because many dimension notions agree there. Cohen--Macaulay and Gorenstein conditions weaken or refine regularity by controlling depth, duality, and resolutions. These properties help determine whether intersections have the expected dimension and whether algebraic invariants behave like those of a smooth space. \cite{dimension_theory_ref5}

\paragraph{Homological and module dimensions}

Krull dimension measures prime-ideal chains. Homological dimension measures how many steps are needed to resolve a module by free or projective modules. The global dimension of a ring is the supremum of projective dimensions of its modules. For a regular Noetherian ring, the global dimension agrees with the Krull dimension in important settings.

Depth measures the length of a regular sequence in a local ring. The inequality
\[
 \operatorname{depth}R\leq\dim R
\]
is fundamental. Equality characterizes Cohen--Macaulay local rings. Thus dimension theory connects the geometry of strata with the algebra of equations and the complexity of solving them.

For modules, dimension can also mean length, rank, or support dimension. A module supported only at finitely many points has dimension zero, while a module whose support fills a surface has dimension two. The correct notion depends on the question being asked.

\paragraph{Valuation, non-Noetherian, and noncommutative dimensions}

A valuation ring has ideals totally ordered by inclusion. Its prime spectrum is correspondingly simple, and its dimension can often be read from the chain of prime ideals. Valuation rings provide local models in algebraic geometry and number theory.

For non-Noetherian rings, ideals need not be finitely generated, chains may be very long, and polynomial dimension formulas can fail. Little is known in full generality, so the theory studies special classes such as coherent, Prüfer, or valuative domains.

Noncommutative rings do not have prime spectra with all the same geometric behavior. One can define Krull dimension-like invariants using chains of prime ideals or modules, but left and right dimensions may differ. These notions are used in representation theory and the study of operator algebras. \cite{dimension_theory_ref6}

\paragraph{What the theory depends on and where it is useful}

Dimension theory depends on rings, ideals, modules, localization, polynomial algebra, algebraic geometry, and homological algebra. It helps algebraic geometry by measuring varieties and fibers, number theory by measuring spectra of rings of integers, and singularity theory by comparing tangent and local dimensions.

It is useful whenever a problem has several possible notions of size. A theorem that equates two dimensions is valuable because it allows a geometric calculation to be replaced by an algebraic one. A mismatch is equally informative: it can reveal a singularity, a bad intersection, or a non-Noetherian pathology.

\section*{4. Important theorems and results}
\begin{enumerate}[leftmargin=*,itemsep=0.35em]

\item \textbf{Polynomial-dimension theorem.} For a Noetherian ring or valuation ring, $\dim R[x]=\dim R+1$ under the standard hypotheses. \cite{dimension_theory_ref1}

\item \textbf{Hilbert basis theorem.} If $R$ is Noetherian, then $R[x]$ is Noetherian. This is a main reason polynomial dimension can be studied inductively. \cite{dimension_theory_ref2}

\item \textbf{Krull principal ideal theorem.} A minimal prime over a principal ideal has height at most one; the version for $r$ generators gives height at most $r$ under Noetherian assumptions. \cite{dimension_theory_ref4}

\item \textbf{Dimension--transcendence theorem.} For a finitely generated domain over a field, the dimension equals the transcendence degree of its fraction field over the base field. \cite{dimension_theory_ref3}

\item \textbf{Regularity criterion.} A Noetherian local ring is regular when its Krull dimension equals the dimension of its maximal ideal modulo its square. \cite{dimension_theory_ref5}
\end{enumerate}

\theoryapplications
\theoryconnections{dimension theory}{%
\item Commutative algebra supplies prime ideals and chains, algebraic geometry interprets them as geometric dimension, and homological algebra supplies projective and injective dimensions.
\item Noetherian hypotheses are important because they prevent uncontrolled infinite chains.
}{%
\item Dimension theory helps algebraic geometry compare varieties, helps local algebra measure singularities, and helps homological algebra predict the length of resolutions.
\item A dimension is useful only after its definition and category are stated; different dimensions answer different questions.
}
\section*{7. Sample Questions}
\emph{Exercise reference:} \cite{dimension_theory_ref1}. The cited edition is the source for the exercise set; consult its Exercises section for the official statement and numbering.
\begin{enumerate}[leftmargin=*,itemsep=0.9em]

\item \textbf{Exercise 1.} What is the Krull dimension of $k[x]$?

\emph{Solution.} The prime chain $(0)\subsetneq(x)$ has length one, and no longer chain exists. Hence $\dim k[x]=1$.

\item \textbf{Exercise 2.} Give a prime chain showing that $k[x,y]$ has dimension at least two.

\emph{Solution.} The ideals $(0)$, $(x)$, and $(x,y)$ are prime and strictly nested. The chain has length two, so the dimension is at least two; the polynomial-dimension theorem shows it is exactly two.

\item \textbf{Exercise 3.} What does height measure?

\emph{Solution.} The height of a prime ideal is the maximum number of strict prime inclusions in a chain ending at that ideal. It measures how many algebraic layers lie below the corresponding subspace.

\item \textbf{Exercise 4.} Why does the crossing $xy=0$ signal a singularity at the origin?

\emph{Solution.} The curve has local dimension one, but its two branches give two independent tangent directions at the origin. The tangent dimension is therefore larger than the local dimension, which is the regularity test for singularity.

\item \textbf{Exercise 5.} Why are Noetherian hypotheses important?

\emph{Solution.} They ensure finite generation and prevent endlessly increasing chains of ideals. Without them, polynomial dimension formulas, finite-height arguments, and geometric interpretations can fail.

\end{enumerate}
\section*{8. References}


\begin{thebibliography}{9}
\bibitem{dimension_theory_ref1} H. Matsumura, \emph{Commutative Ring Theory}, Cambridge University Press, 1986.
\bibitem{dimension_theory_ref2} M. F. Atiyah and I. G. Macdonald, \emph{Introduction to Commutative Algebra}, Addison-Wesley, 1969.
\bibitem{dimension_theory_ref3} R. Hartshorne, \emph{Algebraic Geometry}, Springer, 1977.
\bibitem{dimension_theory_ref4} W. Krull, "Idealtheorie in Ringen ohne Endlichkeitsbedingung," \emph{Mathematische Annalen}, 1926.
\bibitem{dimension_theory_ref5} H. Matsumura, \emph{Commutative Algebra}, Benjamin, 1980.
\bibitem{dimension_theory_ref6} J. C. McConnell and J. C. Robson, \emph{Noncommutative Noetherian Rings}, Wiley, 1987.
\end{thebibliography}


\theory{Distribution theory}{How can we differentiate a sharp impulse, a jump, or another object that is not an ordinary smooth function?}{Distribution theory, also called the theory of generalized functions, treats an object by how it acts on a carefully chosen smooth test function. This makes derivatives meaningful for discontinuous and singular data and gives a rigorous language for weak solutions of differential equations.}{Partial differential equations, Green functions, signal processing, wave propagation, quantum physics, image reconstruction, Fourier analysis, and engineering models with point sources or jumps.}{Test functions and weak derivatives extend differentiation and integration to singular objects such as point masses and shocks.}

\paragraph{The problem and its history}

Classical calculus differentiates smooth functions. But a step function jumps, a point source is concentrated at one location, and a wave may have a sharp front. Their ordinary derivatives either do not exist or do not capture the idealized physical behavior. Engineers nevertheless need equations for impulses, point charges, shocks, and boundaries.

Distribution theory replaces the question “what is the value at each point?” with “what is the effect on every smooth probe?” The probe is a test function. A generalized function is a continuous linear rule that takes a test function to a number. This rule can represent an ordinary function, a measure, a point mass, or a derivative of a singular object. \cite{distribution_theory_ref1,distribution_theory_ref2}

Green functions in the nineteenth century used delta-like sources before the theory was formalized. Sergei Sobolev developed related ideas in the study of hyperbolic partial differential equations in the 1930s. Laurent Schwartz created the modern theory in the 1940s and introduced the name distribution, comparing the objects with distributed electrical charge. \cite{distribution_theory_ref3,distribution_theory_ref4}

\end{historylist}
\section*{3. Theory}
\subsection*{Core theory}
\paragraph{Test functions and distributions}

Let $U$ be an open subset of $\mathbb R^n$. A test function $\varphi$ is infinitely differentiable and has compact support: it is zero outside some bounded closed region. The space of such functions is denoted $\mathcal D(U)$. A distribution is a continuous linear functional
\[
 T:\mathcal D(U)\longrightarrow\mathbb R
\]
or to $\mathbb C$. Its action is written
\[
 \langle T,\varphi\rangle.
\]
Linearity means
\[
 \langle T,a\varphi+b\psi\rangle
 =a\langle T,\varphi\rangle+b\langle T,\psi\rangle.
\]
Continuity means that test functions that converge together with all derivatives produce values that also converge.

Every locally integrable function $f$ defines a distribution:
\[
 \langle T_f,\varphi\rangle
 =\int_U f(x)\varphi(x)\,dx.
\]
Two locally integrable functions define the same distribution exactly when they agree almost everywhere. Thus distributions include ordinary functions but are more general. A Radon measure $\mu$ also defines one by
\[
 \langle T_\mu,\varphi\rangle=\int_U\varphi\,d\mu.
\]
A point mass at the origin gives the Dirac delta:
\[
 \langle\delta,\varphi\rangle=\varphi(0).
\]
The delta is not an ordinary function, but its action is completely well-defined. \cite{distribution_theory_ref1,distribution_theory_ref5}

\paragraph{Derivatives of distributions}

The derivative is defined by moving the derivative onto the test function:
\[
 \langle T',\varphi\rangle=-\langle T,\varphi'\rangle.
\]
This is integration by parts without a boundary term, because $\varphi$ vanishes outside a compact set. If $T=T_f$ for a smooth $f$, this gives the ordinary derivative. The definition also gives a derivative when $f$ has a jump or when $T=\delta$.

Let $H$ be the Heaviside step function, equal to zero on the negative half-line and one on the positive half-line. Then
\[
 \langle H',\varphi\rangle
 =-\int_0^\infty\varphi'(x)\,dx
 =\varphi(0),
\]
so $H'=\delta$. Similarly,
\[
 \langle\delta',\varphi\rangle=-\varphi'(0).
\]
Repeated derivatives are defined in the same way. Every distribution is infinitely differentiable in this generalized sense, even when no classical derivative exists.

This is the key to weak differential equations. Instead of requiring a solution $u$ to have all derivatives pointwise, one asks whether the equation holds after pairing with every test function. Weak solutions can exist for rough initial data and can later be shown to become smooth under additional estimates. \cite{distribution_theory_ref2,distribution_theory_ref6}

\paragraph{Operations, localization, and support}

Distributions can be added and multiplied by numbers, so they form a vector space. They can be multiplied by a smooth function $g$:
\[
 \langle gT,\varphi\rangle=\langle T,g\varphi\rangle.
\]
Composition with a smooth map is possible under suitable conditions. Restricting a distribution to an open subset makes sense because test functions supported there can be used.

The support of $T$ is the smallest closed set outside which $T$ acts as zero. The delta has support $\{0\}$; a smooth function has support equal to the closure of the region where it is nonzero. Compact support makes many operations, including convolution, easier.

There is an important limitation. There is no product of arbitrary distributions that extends ordinary pointwise multiplication while preserving all familiar algebraic and differential rules. The product $\delta^2$ is not automatically defined. Products can be defined for special pairs, when their singular supports do not collide, or in specialized frameworks such as Colombeau algebras. \cite{distribution_theory_ref1,distribution_theory_ref7}

\paragraph{Tempered distributions and Fourier transforms}

Ordinary distributions use compactly supported test functions. Tempered distributions use Schwartz functions: smooth functions that, together with every derivative, decrease faster than any inverse power of $|x|$. The resulting space is denoted $\mathcal S'(\mathbb R^n)$.

Tempered distributions are important because the Fourier transform maps Schwartz functions to Schwartz functions and therefore extends by duality. A constant function transforms into a delta at zero, while a delta transforms into a constant up to normalization. A periodic wave can be represented by a sum of deltas at discrete frequencies. This makes generalized functions natural for signal processing and spectral analysis.

The Fourier transform changes differentiation into multiplication:
\[
 \widehat{\partial_j T}(\xi)=i\xi_j\widehat T(\xi).
\]
It changes convolution into multiplication as well. Differential equations with constant coefficients can therefore become algebraic equations in frequency, subject to the conditions needed to interpret the products. \cite{distribution_theory_ref3,distribution_theory_ref8}

\paragraph{Convolution and Green functions}

For suitable distributions $T$ and functions or distributions $S$, convolution combines their effects:
\[
 (T*S)(x)=\int T(x-y)S(y)\,dy
\]
when the expression makes sense. The delta is the identity for convolution, and convolving with a derivative differentiates the other factor.

A Green function is the response of a differential operator to a point source. If $L G=\delta$, then the solution of $Lu=f$ can often be written as $u=G*f$. Heat flow, electrostatics, wave propagation, and linear circuits all use this idea. The distribution framework permits the point source and the resulting singularity to be handled in one equation.

\paragraph{Distributions beyond ordinary test spaces}

Distributions may be defined on smooth manifolds by using coordinate-independent test densities. Spaces of analytic test functions lead to Sato's hyperfunctions. Analytic functions cannot have nonempty compact support, so the theory has a different structure from ordinary distributions.

Distributions with point support are finite sums of derivatives of the delta. More generally, distributions can be represented locally as derivatives of continuous functions, and tempered distributions can be treated with polynomial growth conditions. These representation results explain why the generalized objects are not arbitrary: they are controlled by ordinary functions after enough integration or differentiation. \cite{distribution_theory_ref1,distribution_theory_ref9}

\paragraph{Applications and dependencies}

Distribution theory depends on integration, linear functionals, topology, Fourier analysis, and differential equations. It helps partial differential equations by defining weak solutions, helps physics by representing point charges and impulses, and helps engineering by modeling ideal switches, shocks, and signals. It also supports Sobolev spaces, where the derivatives in a weak equation are distributions and integrability controls regularity.

In image processing, an edge is close to a derivative of a step and a point feature is close to a delta. In wave equations, a Green function describes how a disturbance travels. In quantum mechanics, position and momentum eigenstates use delta-normalized objects that must be interpreted through test functions. The theory is useful because it lets a model keep idealized features without pretending that they are ordinary smooth graphs.

\section*{4. Important theorems and results}
\begin{enumerate}[leftmargin=*,itemsep=0.35em]

\item \textbf{Embedding theorem.} Every locally integrable function and every Radon measure defines a distribution by integration against test functions. \cite{distribution_theory_ref1}

\item \textbf{Distributional derivative.} The formula $\langle T',\varphi\rangle=-\langle T,\varphi'\rangle$ extends ordinary differentiation and makes the Heaviside derivative equal to the Dirac delta. \cite{distribution_theory_ref2}

\item \textbf{Fourier duality.} Fourier transformation extends from Schwartz functions to tempered distributions and converts differentiation into multiplication by a frequency variable. \cite{distribution_theory_ref8}

\item \textbf{Point-support theorem.} A distribution supported at one point is a finite linear combination of derivatives of the Dirac delta at that point. \cite{distribution_theory_ref1}

\item \textbf{Schwartz impossibility principle.} No multiplication of all distributions can preserve ordinary multiplication and the usual derivative rules simultaneously. Products must therefore be restricted or generalized with extra structure. \cite{distribution_theory_ref7}
\end{enumerate}

\theoryapplications
\theoryconnections{distribution theory}{%
\item Real analysis supplies integration and convergence, functional analysis supplies topological vector spaces, and Fourier analysis supplies the transform of generalized functions.
\item Partial differential equations motivate the weak-derivative viewpoint because classical derivatives may not exist.
}{%
\item Distribution theory helps PDEs define weak solutions, helps physics represent point sources, and helps harmonic analysis extend Fourier methods beyond ordinary functions.
\item It makes differentiation continuous in a larger space, so limits of smooth approximations remain usable.
}
\section*{7. Sample Questions}
\emph{Exercise reference:} \cite{distribution_theory_ref1}. The cited edition is the source for the exercise set; consult its Exercises section for the official statement and numbering.
\begin{enumerate}[leftmargin=*,itemsep=0.9em]

\item \textbf{Exercise 1.} What is the value of $\langle\delta,\varphi\rangle$?

\emph{Solution.} It is $\varphi(0)$. The delta does not have ordinary pointwise values; it extracts the value of a test function at the origin.

\item \textbf{Exercise 2.} Show that the derivative of the Heaviside function is $\delta$.

\emph{Solution.} By definition,
\[
 \langle H',\varphi\rangle=-\langle H,\varphi'\rangle
 =-\int_0^\infty\varphi'(x)\,dx=\varphi(0),
\]
because $\varphi$ vanishes at infinity. This equals $\langle\delta,\varphi\rangle$.

\item \textbf{Exercise 3.} Why is a locally integrable function automatically a distribution?

\emph{Solution.} Its integral against a compactly supported smooth function exists. Integration is linear in the test function, and convergence of the test functions with their derivatives makes the integrals converge, giving continuity.

\item \textbf{Exercise 4.} What does the support of a distribution mean?

\emph{Solution.} It is the smallest closed set outside which the distribution acts as zero. It records where the generalized object can influence test functions.

\item \textbf{Exercise 5.} Why is $\delta^2$ not automatically meaningful?

\emph{Solution.} The general distribution theory defines addition, derivatives, and multiplication by smooth functions, but singularities can collide in a product. A theorem of Schwartz shows that no universal product can retain every ordinary algebraic and differential rule.

\end{enumerate}
\section*{8. References}


\begin{thebibliography}{9}
\bibitem{distribution_theory_ref1} L. Schwartz, \emph{Théorie des distributions}, Hermann, 1950--1951.
\bibitem{distribution_theory_ref2} M. J. Lighthill, \emph{An Introduction to Fourier Analysis and Generalised Functions}, Cambridge University Press, 1958.
\bibitem{distribution_theory_ref3} R. F. Hoskins, \emph{Generalized Functions}, Ellis Horwood, 1979.
\bibitem{distribution_theory_ref4} S. Sobolev, \emph{Some Applications of Functional Analysis in Mathematical Physics}, AMS translation.
\bibitem{distribution_theory_ref5} G. B. Folland, \emph{Real Analysis}, Wiley, 2nd ed., 1999.
\bibitem{distribution_theory_ref6} L. C. Evans, \emph{Partial Differential Equations}, American Mathematical Society, 2nd ed., 2010.
\bibitem{distribution_theory_ref7} L. Schwartz, "Sur l'impossibilité de la multiplication des distributions," \emph{Comptes Rendus}, 1954.
\bibitem{distribution_theory_ref8} M. Reed and B. Simon, \emph{Methods of Modern Mathematical Physics, II: Fourier Analysis, Self-Adjointness}, Academic Press, 1975.
\bibitem{distribution_theory_ref9} J. J. Duistermaat and J. A. C. Kolk, \emph{Distributions: Theory and Applications}, Birkhäuser, 2010.
\end{thebibliography}


\theory{Dynamical systems theory}{How does a state evolve through time, and what long-term behaviors can arise from a simple rule?}{Dynamical systems theory studies a rule that moves a state through a state space. The rule may be a differential equation, a difference equation, a hybrid equation, or a stochastic transition. Instead of seeking an exact formula at every moment, the theory studies equilibria, periodic motion, stability, attractors, chaos, and the dependence on initial conditions.}{Classical mechanics, planetary orbits, electronic circuits, ecology, epidemiology, economics, control, biomechanics, cognitive science, graph networks, and symbolic or topological dynamics.}{State evolution, stability, recurrence, and invariant sets describe qualitative behavior when explicit time-dependent solutions are unavailable.}

\paragraph{The central viewpoint and history}

A dynamical system has a state and a rule for updating it. The state of a pendulum may be its angle and angular velocity; the state of a lake may be the fish population and available food; the state of an economy may include prices, production, and savings. A point in the state space represents one complete snapshot. The evolution rule tells us which point comes next.

Newtonian mechanics supplied the historical roots. Before fast computers, exact solutions were available only for special systems, so mathematicians developed methods for studying qualitative behavior. Modern dynamical systems theory asks questions such as whether a state settles down, whether periodic motion occurs, whether small errors grow, and whether long-term behavior depends on the starting point. \cite{dynamical_systems_theory_ref1,dynamical_systems_theory_ref2}

\end{historylist}
\section*{3. Theory}
\subsection*{Core theory}
\paragraph{Continuous, discrete, and stochastic systems}

A continuous-time system is often written
\[
 \dot x=f(x,t),
\]
where $x(t)$ is a state vector. A discrete-time system uses
\[
 x_{n+1}=F(x_n).
\]
Difference equations describe population updates, compound interest, digital filters, and iterated maps. Differential equations describe motion, chemical reactions, circuits, and fluid models.

Some systems use both kinds of time. A hybrid system may flow continuously until it hits a threshold and then jump. Equations on time scales treat continuous and discrete time within one framework. A stochastic system includes random disturbances, so the rule determines a probability distribution of future states rather than one certain trajectory.

The theory does not require the equations to arise from a least-action principle. Classical mechanics is one important source, but biological, economic, and engineered systems may be postulated directly from observations or modeling assumptions.

\paragraph{Equilibria, periodic points, and stability}

An equilibrium or fixed point is a state that does not change. For $\dot x=f(x)$ it satisfies $f(x)=0$; for $x_{n+1}=F(x_n)$ it satisfies $F(x)=x$. An equilibrium is attracting if nearby states move toward it. Its basin of attraction is the set of initial states that eventually approach it.

For the logistic population model
\[
 \dot x=x(1-x),
\]
the equilibria are $0$ and $1$. A population between zero and one increases toward $1$; a small negative perturbation is physically meaningless but reveals that the stability depends on the chosen state space.

A periodic point returns after a fixed number of steps. Periodic orbits can be attracting, repelling, or neutrally stable. Sharkovskii's theorem gives a remarkable ordering of periods for continuous maps of an interval: the existence of a period-three orbit forces the existence of periodic orbits of every positive integer period. \cite{dynamical_systems_theory_ref3}

\paragraph{Nonlinear behavior and chaos}
\bookfigure{dynamics_chaos_bifurcation}{Iteration, stability, and parameter changes connect dynamical systems with chaos and bifurcation.}

Linear systems obey superposition, so their response to a sum of inputs is the sum of the separate responses. Nonlinear systems do not. Even a low-dimensional nonlinear equation can produce several attractors, bifurcations, quasiperiodic motion, or deterministic chaos.

Chaos is not randomness in the rule. A deterministic chaotic system has a precisely defined future once the initial state is known, but nearby initial states separate exponentially for a time. The Lorenz system, developed from a simplified atmospheric model, became a famous example. Chaos theory gives formal meanings to sensitive dependence, dense periodic orbits, and topological mixing.

\paragraph{Related mathematical viewpoints}

Arithmetic dynamics studies repeated application of polynomial or rational maps to integer, rational, algebraic, or $p$-adic points. Symbolic dynamics replaces a continuous state by an infinite sequence of symbols and evolves it by the shift map. It can encode complicated systems using a finite alphabet.

Topological dynamics studies qualitative asymptotic properties using open sets and continuity. Ergodic theory adds an invariant measure and asks about long-run averages. Functional analysis studies operators on spaces of functions, which provides tools for infinite-dimensional dynamics and evolution equations.

Graph dynamical systems place variables on vertices of a network and update them according to neighboring states. Projected dynamical systems constrain the motion to a feasible set and connect differential equations with optimization and equilibrium problems. System dynamics describes feedback, delays, stocks, and flows. Complex systems use dynamical models when interactions prevent a simple reduction to independent parts. \cite{dynamical_systems_theory_ref1,dynamical_systems_theory_ref4}

\paragraph{Applications beyond mechanics}

In biomechanics, movement emerges from interacting nervous, muscular, skeletal, respiratory, and perceptual subsystems. In cognitive science, dynamicists model a learner's cognitive state as a trajectory influenced by internal and environmental pressures. Concepts such as self-organization, phase transitions, and attractors have been used to discuss development, including the A-not-B error, although applications must be judged by empirical evidence rather than metaphor alone.

The dynamic approach to second-language development treats acquisition and attrition as parts of one changing system. Language behavior can be nonlinear, adaptive, feedback-sensitive, and dependent on initial conditions. Similar ideas appear in ecology, epidemiology, finance, robotics, neural networks, and climate science.

\paragraph{What the theory depends on and where it is useful}

Dynamical systems theory depends on calculus, differential and difference equations, linear algebra, topology, probability, numerical simulation, and computation. It helps chaos theory by supplying the state-space framework, helps control theory by describing plants and closed loops, helps ergodic theory by providing measure-preserving maps, and helps bifurcation theory by identifying parameter changes in long-term behavior.

Its strength is qualitative prediction when exact formulas are impossible. Its danger is overinterpreting a model: a wrong equation can have the wrong attractor, and numerical trajectories can imitate chaos because of roundoff or because the model is not valid at the scale being studied.

\section*{4. Important theorems and results}
\begin{enumerate}[leftmargin=*,itemsep=0.35em]

\item \textbf{Existence and uniqueness principle.} Under suitable regularity conditions on $f$, the initial-value problem $\dot x=f(x,t)$ has a unique local solution, making the state rule well-defined. \cite{dynamical_systems_theory_ref2}

\item \textbf{Linearization principle.} Near a hyperbolic equilibrium, the eigenvalues of the linearized system often determine local stability and qualitative behavior. \cite{dynamical_systems_theory_ref1}

\item \textbf{Sharkovskii's theorem.} For continuous interval maps, a period-three orbit implies periodic orbits of every positive period, explaining why a simple map can contain extremely rich behavior. \cite{dynamical_systems_theory_ref3}

\item \textbf{Poincare--Bendixson theorem.} Under its planar hypotheses, a nonempty compact limit set containing only finitely many equilibria is restricted to equilibria, periodic orbits, and connecting orbits; genuinely chaotic attractors require more room or different structure. \cite{dynamical_systems_theory_ref1}

\item \textbf{Lyapunov stability method.} A scalar function that decreases along trajectories can prove stability without solving the differential equation explicitly. \cite{dynamical_systems_theory_ref5}
\end{enumerate}

\theoryapplications
\theoryconnections{dynamical systems theory}{%
\item Differential equations and iteration define the evolution rule, while linear algebra studies the derivative near an equilibrium.
\item Topology describes invariant and recurrent sets, and probability is added when the state or forcing is random.
}{%
\item Dynamical systems helps mechanics, climate models, ecology, epidemiology, control, and economics describe change over time.
\item Its stability and recurrence results separate robust qualitative behavior from predictions that depend sensitively on uncertain initial data.
}
\section*{7. Sample Questions}
\emph{Exercise reference:} \cite{dynamical_systems_theory_ref1}. The cited edition is the source for the exercise set; consult its Exercises section for the official statement and numbering.
\begin{enumerate}[leftmargin=*,itemsep=0.9em]

\item \textbf{Exercise 1.} Find and classify the equilibria of $\dot x=x(1-x)$.

\emph{Solution.} Equilibria satisfy $x(1-x)=0$, so $x=0$ and $x=1$. The derivative of the right side is $1-2x$. At zero it is positive, so zero is unstable; at one it is negative, so one is locally asymptotically stable.

\item \textbf{Exercise 2.} What is the difference between an equilibrium and a periodic orbit?

\emph{Solution.} An equilibrium remains at one state forever. A periodic orbit returns to its initial state after a positive time or number of steps and generally visits several states in between.

\item \textbf{Exercise 3.} Why can deterministic chaos look random?

\emph{Solution.} The rule contains no random choice, but small uncertainty in the initial state can grow rapidly. After enough time, the observer cannot know which nearby trajectory is the actual one, so the observed sequence appears random.

\item \textbf{Exercise 4.} What is an attractor basin?

\emph{Solution.} It is the set of initial states whose future trajectories approach a particular attractor. If several attractors exist, their basins divide the state space and small parameter or measurement changes can alter the outcome.

\item \textbf{Exercise 5.} Why is simulation not a proof of long-term behavior?

\emph{Solution.} A finite simulation samples one initial condition for one duration with finite precision. It can miss another attractor, a rare transition, numerical instability, or behavior beyond the model's validity. Theorems and error estimates are needed for rigorous conclusions.

\end{enumerate}
\section*{8. References}


\begin{thebibliography}{9}
\bibitem{dynamical_systems_theory_ref1} S. H. Strogatz, \emph{Nonlinear Dynamics and Chaos}, Westview Press, 2nd ed., 2015.
\bibitem{dynamical_systems_theory_ref2} L. Perko, \emph{Differential Equations and Dynamical Systems}, Springer, 3rd ed., 2001.
\bibitem{dynamical_systems_theory_ref3} A. N. Sharkovskii, "Co-existence of cycles of a continuous mapping of the line into itself," \emph{Ukrainian Mathematical Journal}, 1964.
\bibitem{dynamical_systems_theory_ref4} G. L. Baker and J. P. Gollub, \emph{Chaotic Dynamics}, Cambridge University Press, 1996.
\bibitem{dynamical_systems_theory_ref5} H. K. Khalil, \emph{Nonlinear Systems}, Prentice Hall, 3rd ed., 2002.
\end{thebibliography}


\theory{Elimination theory}{How can unwanted variables be removed from a system of equations without losing the solutions we care about?}{Elimination theory develops algorithms and structural results for projecting polynomial equations onto fewer variables. Substitution, resultants, discriminants, Groebner bases, and cylindrical algebraic decomposition turn a many-variable problem into equations or logical conditions in selected variables.}{Solving polynomial systems, robot kinematics, computer algebra, geometric intersection, invariant theory, quantifier elimination, CAD, and exact symbolic computation.}{Resultants, Groebner bases, and elimination orders remove variables while preserving the polynomial consequences of a system.}

\paragraph{The problem and its history}

Suppose a robot arm has joint angles $x,y$ and its endpoint coordinates are known. The equations may contain both angles, but we want a condition only on the endpoint. Elimination theory removes variables to produce that condition. The new equation is the shadow of the original solution set under projection.

The subject grew from the need to solve polynomial systems. Bezout's theorem bounded the number of solutions in appropriate projective settings. Gaussian elimination solved linear systems by removing variables one at a time. In the nineteenth century, Hermite and Smith normal forms handled linear Diophantine equations and finitely generated abelian groups. Resultants and discriminants supplied polynomial eliminants. \cite{elimination_theory_ref1,elimination_theory_ref2}

Francis Macaulay developed multivariate resultants and U-resultants. Elimination was central in early algebra, then was neglected when Hilbert's non-effective methods became fashionable. Computer algebra revived it: Groebner bases and cylindrical algebraic decomposition made elimination algorithmic again around the 1970s. \cite{elimination_theory_ref3}

\end{historylist}
\section*{3. Theory}
\subsection*{Core theory}
\paragraph{A simple example and the geometric meaning}

Consider
\[
 x+y=3,\qquad xy=2.
\]
The first equation gives $y=3-x$, and substitution yields
\[
 x(3-x)=2,\qquad x^2-3x+2=0.
\]
The two possible values $x=1,2$ then give $y=2,1$. The variable $y$ has been eliminated, but it can be recovered by back-substitution.

For polynomial equations $f_1(x,y)=\cdots=f_r(x,y)=0$, the solution set is an algebraic variety in the $(x,y)$-plane. Projecting it onto the $x$-axis may create a larger set unless the eliminated coordinate can actually be recovered. An elimination polynomial gives a necessary condition, and additional checks are needed to remove extraneous roots introduced by multiplication or closure.

The algebraic object is an ideal $I$ generated by the equations. The elimination ideal
\[
 I\cap k[x_1,\ldots,x_r]
\]
contains all polynomial consequences involving only the variables we keep. The geometric projection of the original variety is related to the zero set of this ideal, with closure issues controlled by algebraic geometry.

\paragraph{Resultants and discriminants}

For two polynomials in one eliminated variable, the resultant is a determinant built from their coefficients. It vanishes exactly when the polynomials have a common root over an algebraic closure, with multiplicities encoded in the determinant. Thus
\[
 \operatorname{Res}_y(f,g)=0
\]
eliminates $y$ from $f(x,y)=0$ and $g(x,y)=0$.

The discriminant of a polynomial is a special resultant between the polynomial and its derivative. It vanishes when two roots collide, so it detects repeated roots and singular parameter values. In geometry it detects when curves become tangent, intersect nontransversely, or acquire a singularity.

Resultants are effective and invariant under many changes of coordinates, but determinant size grows quickly. They can also vanish for reasons that need interpretation, such as a root at infinity in a projective compactification. \cite{elimination_theory_ref3,elimination_theory_ref4}

\paragraph{Groebner bases and elimination orders}

A Groebner basis is a special generating set for a polynomial ideal, adapted to a chosen monomial order. Buchberger's algorithm repeatedly reduces pairs of polynomials until all hidden cancellations are resolved. The computation may be expensive, but the resulting basis makes membership and elimination questions systematic.

The elimination theorem says that if a Groebner basis is computed with an order that places $y$ above $x$, then the basis elements involving only $x$ form a Groebner basis for the elimination ideal $I\cap k[x]$. In practice, a triangular basis can give equations in one variable first, then in two variables, and so on. This resembles solving by substitution but works for nonlinear systems and does not divide by a coefficient that might be zero.

For example, a computer algebra system can eliminate joint variables from robot equations, obtain a polynomial for a possible position, and then solve for the joints. The system may have complex or repeated roots, so physical constraints and numerical verification remain necessary. \cite{elimination_theory_ref1,elimination_theory_ref5}

\paragraph{Logic and quantifier elimination}

Elimination also has a logical meaning. A statement
\[
 \exists y\;[f(x,y)=0\text{ and }g(x,y)>0]
\]
asks whether there is a value of $y$ satisfying conditions for a given $x$. Quantifier elimination replaces the statement by an equivalent condition involving only $x$.

Over algebraically closed fields, polynomial elimination is connected with quantifier elimination for the theory of algebraically closed fields. Over the real numbers, cylindrical algebraic decomposition partitions space into cells on which polynomial signs are constant. It can decide formulas with polynomial equalities and inequalities, though its worst-case complexity is high.

This connection explains why elimination theory matters beyond equations. It turns geometric existence questions into finite logical descriptions and supports exact analysis in computer-aided design, collision detection, and real algebraic geometry. \cite{elimination_theory_ref6}

\paragraph{What the theory depends on and where it is useful}

Elimination theory depends on polynomial rings, ideals, determinants, linear algebra, algebraic geometry, algorithms, and mathematical logic. It helps algebraic geometry by computing projections and closures, invariant theory by producing coordinate-free polynomial relations, and computer algebra by making symbolic solution procedures possible.

In CAD, elimination detects whether two surfaces intersect or become tangent. In computer vision, it removes hidden depth or calibration parameters. In robot kinematics, it produces equations for reachable positions. In solving equations, it turns many variables into a sequence of lower-dimensional problems. Its limits are computational size, extraneous solutions, multiplicities, and the distinction between algebraic closure and the real solutions required by an application.

\section*{4. Important theorems and results}
\begin{enumerate}[leftmargin=*,itemsep=0.35em]

\item \textbf{Bezout's theorem.} Under suitable projective and generic hypotheses, two plane curves of degrees $m$ and $n$ meet in $mn$ points counted with multiplicity. This gives a bound and explains why multiplicity matters in elimination. \cite{elimination_theory_ref2}

\item \textbf{Resultant criterion.} The resultant of two univariate polynomials vanishes exactly when they have a common root over an algebraic closure. \cite{elimination_theory_ref4}

\item \textbf{Elimination theorem.} A Groebner basis in an elimination order contains a basis for the ideal involving only the variables retained. \cite{elimination_theory_ref1}

\item \textbf{Nullstellensatz connection.} Over an algebraically closed field, a polynomial system has no common solution exactly when its ideal contains $1$, so elimination can produce the contradiction $1=0$. \cite{elimination_theory_ref5}

\item \textbf{Quantifier-elimination principle.} In suitable theories, formulas involving quantified variables can be replaced by equivalent formulas without quantifiers; cylindrical algebraic decomposition provides a real-algebraic algorithm for polynomial inequalities. \cite{elimination_theory_ref6}
\end{enumerate}

\theoryapplications
\theoryconnections{elimination theory}{%
\item Linear algebra supplies elimination by row reduction, commutative algebra supplies ideals and Groebner bases, and algebraic geometry interprets elimination as projection of a solution set.
\item Resultants and discriminants provide compact polynomial tests for common roots and singular behavior.
}{%
\item Elimination theory helps solve polynomial systems, robotics constraints, computer-aided design, and algebraic statistics.
\item It turns hidden variables into explicit conditions, but the degree and coefficient growth can make exact computation expensive.
}
\section*{Glossary}
\begin{description}[leftmargin=*,style=nextline,itemsep=0.3em]
\item[\textbf{Resultant.}] A polynomial in the coefficients of two univariate polynomials that vanishes exactly when the polynomials share a common root; it eliminates one variable from a system.
\item[\textbf{Discriminant.}] A special resultant between a polynomial and its derivative that vanishes when two roots coincide, detecting repeated roots and singular parameter values.
\item[\textbf{Gröbner basis.}] A specially chosen generating set for a polynomial ideal adapted to a monomial order, which makes membership and elimination questions systematic and algorithmic.
\item[\textbf{Elimination ideal.}] The subset of an ideal that involves only selected variables, obtained by intersecting the ideal with a smaller polynomial ring; it captures all consequences of the system in those variables.
\item[\textbf{Quantifier elimination.}] A logical procedure that replaces a formula with quantified variables by an equivalent formula without quantifiers, turning existence questions into explicit conditions.
\item[\textbf{Cylindrical algebraic decomposition.}] An algorithm that partitions real space into cells on which the signs of a family of polynomials are constant, enabling decision procedures for formulas involving polynomial equalities and inequalities.
\item[\textbf{Monomial order.}] A rule for ranking polynomial terms that determines the leading term and drives the reduction steps in Gröbner-basis computations.
\item[\textbf{Buchberger's algorithm.}] The algorithm for computing a Gröbner basis by repeatedly reducing pairs of polynomials until all S-polynomials reduce to zero.
\item[\textbf{Nullstellensatz.}] Hilbert's theorem that a polynomial system has no common solution over an algebraically closed field iff its ideal contains $1$.
\end{description}

\section*{7. Sample Questions}
\emph{Exercise reference:} \cite{elimination_theory_ref1}. The cited edition is the source for the exercise set; consult its Exercises section for the official statement and numbering.
\begin{enumerate}[leftmargin=*,itemsep=0.9em]

\item \textbf{Exercise 1.} Compute the resultant of $f(x)=x^2-2$ and $g(x)=x^2-3$.

\emph{Solution.} The Sylvester matrix is $\begin{bmatrix}1&0&-2&0\\0&1&0&-2\\1&0&-3&0\\0&1&0&-3\end{bmatrix}$. Expanding: $\det=1\cdot\det\begin{bmatrix}1&0&-2\\0&-3&0\\1&0&-3\end{bmatrix}-(-2)\cdot\det\begin{bmatrix}0&1&-2\\1&0&0\\0&1&0\end{bmatrix}$. First $3\times3$ det: $1(9-0)-0+(-2)(0+3)=9-6=3$. Second $3\times3$ det: $0-1(0-0)+(-2)(1-0)=-2$. Total: $3-4=-1$. Since $\operatorname{Res}(f,g)=-1\neq0$, the polynomials share no common root, consistent with roots $\pm\sqrt{2}$ and $\pm\sqrt{3}$.

\item \textbf{Exercise 2.} Eliminate $y$ from $x+y=3$ and $x^2+y^2=5$.

\emph{Solution.} From the first equation, $y=3-x$. Substituting: $x^2+(3-x)^2=5$, i.e., $2x^2-6x+4=0$, so $(x-1)(x-2)=0$, giving $x\in\{1,2\}$. Back-substitution yields $y=2$ and $y=1$. The elimination polynomial is $2x^2-6x+4=0$ and the solution set is $\{(1,2),(2,1)\}$.

\item \textbf{Exercise 3.} Compute the discriminant of $p(x)=x^3-x$.

\emph{Solution.} For $x^3+0\cdot x^2-x+0$: $a=1$, $b=0$, $c=-1$, $d=0$. The discriminant is $\Delta=18abcd-4b^3d+b^2c^2-4ac^3-27a^2d^2=0-0+0-4(1)(-1)-0=4$. Since $\Delta=4\neq0$, all three roots are distinct (they are $0,\pm1$).

\item \textbf{Exercise 4.} Find a Gr\"obner basis for $I=\langle x^2+y-1,\;x-y\rangle$ under lexicographic order $x>y$.

\emph{Solution.} From $x-y$, we get $x=y$. Substituting into $x^2+y-1$ gives $y^2+y-1$. The Gr\"obner basis is $\{x-y,\;y^2+y-1\}$. The elimination ideal is $I\cap k[y]=\langle y^2+y-1\rangle$, with roots $y=(-1\pm\sqrt{5})/2$.

\item \textbf{Exercise 5.} Use the Nullstellensatz to determine whether the system $x^2+1=0$, $x^3+1=0$ has a common solution over $\mathbb C$.

\emph{Solution.} The ideals are $\langle x^2+1\rangle$ and $\langle x^3+1\rangle$. Compute $I=\langle x^2+1,x^3+1\rangle$: dividing $x^3+1$ by $x^2+1$ gives remainder $-x+1$, so $I$ contains $x-1$. Then $x-1$ and $x^2+1$ are coprime (the gcd is $1$), so $I=\langle1\rangle$. By the Nullstellensatz, $V(I)=\varnothing$---the system has no common solution over $\mathbb C$.

\end{enumerate}
\section*{8. References}


\begin{thebibliography}{9}
\bibitem{elimination_theory_ref1} D. Cox, J. Little, and D. O'Shea, \emph{Ideals, Varieties, and Algorithms}, Springer, 4th ed., 2015.
\bibitem{elimination_theory_ref2} I. Shafarevich, \emph{Basic Algebraic Geometry}, Springer, 2nd ed., 1994.
\bibitem{elimination_theory_ref3} F. S. Macaulay, \emph{The Algebraic Theory of Modular Systems}, Cambridge University Press, 1916.
\bibitem{elimination_theory_ref4} I. M. Gel'fand, M. Kapranov, and A. Zelevinsky, \emph{Discriminants, Resultants, and Multidimensional Determinants}, Birkhäuser, 1994.
\bibitem{elimination_theory_ref5} D. Eisenbud, \emph{Commutative Algebra with a View Toward Algebraic Geometry}, Springer, 1995.
\bibitem{elimination_theory_ref6} G. E. Collins, "Quantifier elimination for real closed fields by cylindrical algebraic decomposition," in \emph{Automata Theory and Formal Languages}, 1975.
\end{thebibliography}


\theory{Ergodic theory}{When can the average seen by one system over a long time equal the average over all possible states?}{Ergodic theory studies statistical properties of deterministic dynamical systems. A measure-preserving transformation moves points without changing total measure, and an ergodic transformation has no invariant region of intermediate measure. For almost every starting point, long-time averages then agree with space averages.}{Statistical mechanics, probability, dynamical systems, harmonic analysis, geometry, Lie groups, number theory, geodesic flows, Markov processes, and recurrence questions.}{Measure-preserving dynamics and invariant averages explain when long trajectories represent the statistical behavior of the whole space.}

\paragraph{The problem and its history}

A gas contains an enormous number of particles, so tracking every trajectory is impossible. A physicist instead measures temperature or pressure, which are averages. Ergodic theory asks when observing one deterministic trajectory for a long time gives the same answer as sampling all states at one time.

The theory is built on measure theory and deterministic dynamics. There is no random noise in the equations; the apparent randomness comes from complicated motion and limited observation. Its early motivation was statistical physics, where the ergodic hypothesis proposed that time averages and ensemble averages should agree. \cite{ergodic_theory_ref1,ergodic_theory_ref2}

\end{historylist}
\section*{3. Theory}
\subsection*{Core theory}
\paragraph{Measure-preserving transformations and ergodicity}

Let $(X,\Sigma,\mu)$ be a probability space and let $T:X\to X$ be measurable. It preserves measure if
\[
 \mu(T^{-1}E)=\mu(E)
\]
for every measurable set $E$. The transformation is ergodic if every essentially invariant set has measure zero or one:
\[
 \mu(T^{-1}E\mathbin{\triangle}E)=0
 \quad\Longrightarrow\quad
 \mu(E)\in\{0,1\}.
\]
There is no nontrivial region that the dynamics preserves as a separate statistical world.

The oatmeal analogy is useful. A spoonful of syrup is stirred through a bowl without compressing or dilating the material. An ergodic stirring eventually spreads the syrup so thoroughly that a long observation of one small portion reflects the whole bowl. Mixing is stronger: it says that two sets become asymptotically independent. Equidistribution says that visits are distributed according to the measure.

\paragraph{Examples}

Rotation of the circle by an irrational angle,
\[
 T(x)=x+\theta\pmod 1,
\]
is ergodic. In fact it is uniquely ergodic, minimal, and equidistributed: every orbit visits intervals in the correct proportion. If $\theta=p/q$ is rational, every orbit is periodic with period $q$, so an interval shorter than $1/q$ produces an invariant union with intermediate measure and the map is not ergodic.

A Bernoulli shift moves an infinite sequence of independent symbols one place to the left. It is ergodic, and Kolmogorov's zero--one law explains why tail events have probability zero or one. For torus automorphisms given by an integer matrix, ergodicity is equivalent to having no eigenvalue that is a root of unity.

Hamiltonian systems preserve phase-space volume under Liouville's theorem. A compact energy surface with a nonconstant first integral has invariant subsets of intermediate measure, so it cannot be ergodic on that surface. Completely integrable systems provide an important opposite behavior: many conserved quantities constrain trajectories to lower-dimensional tori. \cite{ergodic_theory_ref2,ergodic_theory_ref3}

\paragraph{Time averages, space averages, and recurrence}

For an integrable observable $f$, the time average along the trajectory starting at $x$ is
\[
 \widehat f(x)=\lim_{N\to\infty}
 \frac1N\sum_{k=0}^{N-1}f(T^kx).
\]
The space average on a probability space is
\[
 \overline f=\int_Xf\,d\mu.
\]
For a general measure-preserving system these may differ. For an ergodic system they agree for almost every starting point.

Poincare's recurrence theorem says that in a finite measure-preserving system, almost every point in a measurable set eventually returns to that set. Recurrence does not mean that the trajectory returns to the exact same point; it returns arbitrarily close or revisits a set. A smaller set has a longer typical recurrence time.

If $A$ is measurable, the time spent in $A$ is the average of its indicator function:
\[
 \lim_{N\to\infty}\frac1N\sum_{k=0}^{N-1}\mathbf1_A(T^kx)=\mu(A)
\]
for almost every $x$ in an ergodic probability system. The average return time to $A$, when finite, is proportional to $1/\mu(A)$. \cite{ergodic_theory_ref1,ergodic_theory_ref4}

\paragraph{The ergodic theorems}

The Birkhoff pointwise ergodic theorem states that if $f\in L^1(\mu)$ and $T$ preserves $\mu$, then the time average exists almost everywhere. The limit is invariant under $T$ and has the same integral as $f$. If $T$ is ergodic, the invariant limit is constant and equals $\int f\,d\mu$.

In probabilistic form, the Birkhoff--Khinchin theorem identifies the limit with the conditional expectation of $f$ given the sigma-algebra of invariant sets. Ergodicity makes that sigma-algebra trivial, producing the usual expected value.

Von Neumann's mean ergodic theorem gives convergence in Hilbert-space norm. If $U$ is a unitary or isometric operator and $P$ projects onto its fixed vectors, then
\[
 \frac1N\sum_{n=0}^{N-1}U^nx\longrightarrow Px
\]
in norm. For $Uf=f\circ T$, this says that averages converge to the invariant part of the function. In $L^p$ spaces, ergodic means converge to conditional expectation under the appropriate strong or weak operator topology. \cite{ergodic_theory_ref1,ergodic_theory_ref5}

\paragraph{Flows, geometry, and connected subjects}

A continuous flow is a family $T_t$ satisfying $T_{s+t}=T_sT_t$. Ergodic flows arise from differential equations. Hopf proved ergodicity of geodesic flow on negatively curved surfaces; later work extended this to variable negative curvature and higher-dimensional manifolds.

Geodesic flows on hyperbolic manifolds connect ergodic theory with Lie groups and lattices. Horocycle flows, homogeneous flows, Anosov systems, and Ratner's theorems study how invariant measures and orbits are classified. Harmonic analysis and representation theory provide operators whose spectra reveal mixing and decay of correlations.

Number theory uses ergodic methods for Diophantine approximation, lattice actions, and the distribution of arithmetic orbits. Markov chains and stationary stochastic processes provide probabilistic examples, even though the underlying shift dynamics can be deterministic.

\paragraph{What the theory depends on and where it is useful}

Ergodic theory depends on measure theory, probability, dynamical systems, functional analysis, operator theory, and topology. It helps statistical mechanics by justifying macroscopic averages, helps number theory by translating orbit distribution into measure statements, and helps geometry by analyzing geodesic flows.

The theory is useful when individual trajectories are too complicated but averages are meaningful. It does not say that every deterministic system is ergodic. Conserved quantities, disconnected invariant regions, and stable islands can prevent time and space averages from agreeing.

\section*{4. Important theorems and results}
\begin{enumerate}[leftmargin=*,itemsep=0.35em]

\item \textbf{Poincare recurrence theorem.} Almost every point of a finite measure-preserving system returns arbitrarily close to its initial region infinitely often. \cite{ergodic_theory_ref1}

\item \textbf{Birkhoff pointwise ergodic theorem.} Time averages of integrable observables exist almost everywhere; in an ergodic system they equal the space average almost everywhere. \cite{ergodic_theory_ref1}

\item \textbf{Von Neumann mean ergodic theorem.} Cesaro averages of a unitary operator converge in Hilbert norm to the orthogonal projection onto its fixed-vector space. \cite{ergodic_theory_ref5}

\item \textbf{Birkhoff--Khinchin formulation.} The pointwise time-average limit is the conditional expectation onto the invariant sigma-algebra; ergodicity makes this sigma-algebra trivial. \cite{ergodic_theory_ref2}

\item \textbf{Kac recurrence principle.} For a measurable set of positive measure in an ergodic probability system, the mean return time is the reciprocal of its measure, under the standard return-time convention. \cite{ergodic_theory_ref4}
\end{enumerate}

\theoryapplications
\theoryconnections{ergodic theory}{%
\item Measure theory supplies measurable spaces and invariant measures, probability supplies random variables and expectations, and analysis supplies convergence in function spaces.
\item Dynamical systems provides the transformations whose long-run averages are studied.
}{%
\item Ergodic theory helps statistical mechanics replace microscopic motion by macroscopic averages, and it helps number theory and dynamics study equidistribution.
\item The ergodic hypothesis is a substantive modeling or structural condition, not a synonym for randomness.
}
\section*{7. Sample Questions}
\emph{Exercise reference:} \cite{ergodic_theory_ref1}. The cited edition is the source for the exercise set; consult its Exercises section for the official statement and numbering.
\begin{enumerate}[leftmargin=*,itemsep=0.9em]

\item \textbf{Exercise 1.} Why is an irrational rotation of the circle ergodic?

\emph{Solution.} Its orbit never repeats and is equidistributed around the circle. An invariant measurable set must therefore contain almost all or almost none of the orbit, giving measure zero or one.

\item \textbf{Exercise 2.} What is the difference between a time average and a space average?

\emph{Solution.} A time average follows one orbit and averages the observed values. A space average integrates the observable over all states using the invariant measure. Ergodic theorems state when they agree.

\item \textbf{Exercise 3.} What does measure preservation mean?

\emph{Solution.} If a set has measure $\mu(E)$, its inverse image under the transformation has the same measure. The dynamics rearranges mass but does not create or destroy it.

\item \textbf{Exercise 4.} Does recurrence mean a point returns exactly to its starting position?

\emph{Solution.} Usually no. It means that the orbit returns arbitrarily close or revisits a measurable neighborhood. Exact return may be impossible in a continuous state space.

\item \textbf{Exercise 5.} Why does a conserved quantity obstruct ergodicity?

\emph{Solution.} A conserved quantity partitions the phase space into level sets. If more than one level set has positive but not full measure, each is an invariant region, contradicting ergodicity.

\end{enumerate}
\section*{8. References}
\addcontentsline{toc}{section}{References}

\begin{thebibliography}{9}
\bibitem{ergodic_theory_ref1} Peter Walters, \emph{An Introduction to Ergodic Theory}, Springer, 1982.
\bibitem{ergodic_theory_ref2} I. P. Cornfeld, S. V. Fomin, and Ya. G. Sinai, \emph{Ergodic Theory}, Springer, 1982.
\bibitem{ergodic_theory_ref3} George D. Birkhoff, "Proof of the ergodic theorem," \emph{Proceedings of the National Academy of Sciences}, 17 (1931), 656--660.
\bibitem{ergodic_theory_ref4} Patrick Billingsley, \emph{Ergodic Theory and Information}, Wiley, 1965.
\bibitem{ergodic_theory_ref5} John von Neumann, "Proof of the quasi-ergodic hypothesis," \emph{Proceedings of the National Academy of Sciences}, 18 (1932), 70--82.
\end{thebibliography}


\theory{Extremal graph theory}{How large or dense can a graph be when a specified pattern is forbidden?}{Extremal graph theory turns a vague question about the biggest possible network into a precise optimization problem. It connects global quantities, such as vertices and edges, with local patterns, such as triangles and cliques.}{Network design, communication systems, scheduling, additive combinatorics, probability, and the analysis of large data networks.}{Forbidden-subgraph counts and density bounds identify the largest configurations compatible with a prescribed local restriction.}

\paragraph{The problem and its history}
A graph is a collection of points called vertices joined by edges. The edges may represent friendships, roads, wires, or links in a computer network. A graph can be globally dense while still avoiding a particular local shape, but only up to a limit. Extremal graph theory asks for that limit and for the graphs that reach it. In symbols, if $H$ is a forbidden graph, $\operatorname{ex}(n,H)$ is the greatest number of edges in an $n$-vertex graph containing no copy of $H$ \cite{extremal_graph_theory_ref1,extremal_graph_theory_ref2}.

The subject grew from Mantel's theorem in 1907, which determined the maximum number of edges in a triangle-free graph, and Turan's theorem in 1941, which treated graphs with no clique of a fixed size. The work of Erdős and Stone then showed that, for every non-bipartite forbidden graph, the leading term of the answer is controlled by its chromatic number. Hungarian mathematicians played an especially important role in the subject \cite{extremal_graph_theory_ref3,extremal_graph_theory_ref4}.

\end{historylist}
\section*{3. Theory}
\subsection*{Core theory}
\paragraph{A first example: avoiding triangles}
Suppose $n$ students form a network in which an edge means that two students have worked together. We want as many collaborations as possible but do not want three students who all worked with one another. Put the students into two groups, as evenly as possible, and allow every edge between the groups but none inside a group. This is a complete bipartite graph. It contains no triangle, because a triangle would need two vertices from the same group, and those two vertices are not joined.

If the groups have $a$ and $n-a$ students, the number of edges is $a(n-a)$. This product is largest when the groups are as equal as possible, giving
\[
\operatorname{ex}(n,K_3)=\left\lfloor\frac{n^2}{4}\right\rfloor.
\]
This is Mantel's theorem. It is not just a count: it says that a network exceeding this limit must contain a triangle. In a safety or incompatibility network, that conclusion can be more useful than the construction itself.

\paragraph{Turán graphs and forbidden cliques}
The complete graph $K_r$ has $r$ vertices with every pair joined. Turan's theorem asks for the largest $n$-vertex graph with no $K_{r+1}$. Divide the vertices into $r$ nearly equal groups, join every pair of vertices in different groups, and leave each group independent. The resulting Turan graph $T(n,r)$ has no $(r+1)$-clique and has the maximum possible number of edges. The balanced construction is forced: if one group were much larger than another, moving a vertex would increase the number of cross-group edges.

The theorem illustrates the characteristic proof pattern of extremal graph theory. First construct an example showing that a proposed bound can be reached. Then prove that every graph above the bound contains the forbidden object. The extremal graph is often highly organized rather than random.

\paragraph{Coloring as an extremal question}
A proper vertex coloring assigns colors so that adjacent vertices receive different colors. The smallest possible number of colors is the chromatic number $\chi(G)$. A clique of size $\omega(G)$ requires at least $\omega(G)$ colors. An independent set, in which no two vertices are adjacent, can contain at most $\alpha(G)$ vertices of any one color, so
\[
\chi(G)\geq \frac{|V(G)|}{\alpha(G)}.
\]
A greedy algorithm gives $\chi(G)\leq\Delta(G)+1$, where $\Delta(G)$ is the largest vertex degree. Brooks' theorem improves this to $\Delta(G)$ except for complete graphs and odd cycles. Planar graphs need at most four colors by the four-color theorem \cite{extremal_graph_theory_ref2,extremal_graph_theory_ref5}.

These inequalities convert local restrictions into global limits. For example, if every radio transmitter can interfere with at most five others, six frequencies always suffice by greedy coloring. A sharper structural theorem may show that five suffice unless the interference graph has a special form. Finding whether a general graph has a coloring with a prescribed number of colors is NP-hard, so extremal bounds are valuable even when exact computation is impractical.

Edge coloring assigns colors to edges so that edges sharing a vertex differ. Its minimum is the chromatic index $\chi'(G)$. Vizing's theorem says that every simple graph satisfies
\[
\Delta(G)\leq\chi'(G)\leq\Delta(G)+1.
\]
Thus a timetable in which each edge is a job incident to a shared resource needs either $\Delta$ or $\Delta+1$ time slots.

\paragraph{Forbidden subgraphs and asymptotic answers}
For a fixed graph $H$, the forbidden-subgraph problem seeks $\operatorname{ex}(n,H)$. Turan's theorem gives an exact answer for cliques. For a general non-bipartite $H$, the Erdős--Stone theorem gives the leading asymptotic term:
\[
\operatorname{ex}(n,H)=\left(1-\frac{1}{\chi(H)-1}+o(1)\right)\binom n2.
\]
Here $o(1)$ means a quantity that becomes negligible as $n$ grows. The surprising point is that the exact shape of $H$ matters less than its chromatic number at the largest scale. Bipartite forbidden graphs are harder: even the correct asymptotic constant is unknown for many of them. When $H$ is complete bipartite, the question is closely related to the Zarankiewicz problem, which asks how many incidences can avoid a complete rectangular pattern.

\paragraph{Homomorphism density and graph limits}
A graph homomorphism from $H$ to $G$ maps vertices of $H$ to vertices of $G$ while preserving edges; different vertices are allowed to map to the same place. The homomorphism density $t(H,G)$ is the probability that a randomly chosen map has this property. It is a softened way to measure how strongly $H$ occurs in $G$ \cite{extremal_graph_theory_ref1}.

For a large dense network, sampling every subgraph is impossible. Homomorphism densities provide statistics that survive when the network grows. A graphon is a measurable function on a square that acts as a limit object for dense graphs. In this setting $t(H,G)$ becomes an integral, and inequalities such as Cauchy--Schwarz and Holder's inequality yield relations among pattern densities. Sidorenko's conjecture proposes a sharp lower bound for the density of every bipartite graph in terms of the edge density; it remains a major open problem.

\paragraph{Regularity: finding order inside a large graph}
Szemeredi's regularity lemma says that the vertices of every sufficiently large graph can be partitioned into a bounded number of parts so that most pairs of parts behave approximately like random bipartite graphs. A pair is regular when the edge density between large subsets is close to the density between the whole parts. The lemma does not claim that the graph itself is random. It says that a complicated graph has a coarse description with controlled irregularity \cite{extremal_graph_theory_ref4,extremal_graph_theory_ref6}.

The counting lemma uses a regular partition to estimate the number of copies of a fixed small graph. The removal lemma says that if a graph has very few copies of a pattern, then deleting a small number of edges removes all copies. These results connect extremal graph theory to additive combinatorics, property testing, and computational complexity. Strong regularity, weak Frieze--Kannan regularity, and hypergraph versions refine the same strategy.

\paragraph{Methods and neighboring theories}
The probabilistic method proves that an object exists by showing that a random object has positive probability of having the required property. Dependent random choice, container methods, and hypergraph regularity are important techniques. Ramsey theory asks when a sufficiently large structure must contain a highly organized monochromatic pattern; spectral graph theory uses eigenvalues to control edges, expansion, and forbidden configurations. Additive combinatorics studies sets with many additive relations, and complexity theory asks how hard it is to find or certify the extremal object. The Ruzsa--Szemeredi problem and Ore's theorem are among related results \cite{extremal_graph_theory_ref1,extremal_graph_theory_ref6}.

\paragraph{What the theory depends on and where it is useful}
The subject depends on graph theory, counting, inequalities, probability, and optimization. It helps Ramsey theory by supplying density thresholds, helps spectral graph theory by turning eigenvalue information into edge bounds, and helps additive combinatorics through regularity and removal arguments.

In network planning, an extremal bound can certify that a proposed number of links forces a forbidden conflict. In scheduling, coloring bounds give guaranteed numbers of time slots. In algorithms, regularity and graph limits compress large dense networks into a manageable model. The limitations are equally important: a sharp asymptotic estimate may not give an exact answer for practical values of $n$, and many bipartite forbidden-subgraph problems remain open.

\section*{4. Important theorems and results}
\begin{enumerate}[leftmargin=*,itemsep=0.35em]
\item \textbf{Mantel's theorem.} A triangle-free graph on $n$ vertices has at most $\lfloor n^2/4\rfloor$ edges, with equality for the balanced complete bipartite graph \cite{extremal_graph_theory_ref3}.

\item \textbf{Turán's theorem.} Among $n$-vertex graphs with no $K_{r+1}$, the balanced $r$-partite Turan graph has the most edges \cite{extremal_graph_theory_ref3}.

\item \textbf{Erdős--Stone theorem.} For non-bipartite $H$, the leading term of $\operatorname{ex}(n,H)$ is determined by $\chi(H)$ \cite{extremal_graph_theory_ref1}.

\item \textbf{Szemeredi regularity lemma.} Every sufficiently large graph has a bounded regular partition, providing a random-like coarse model \cite{extremal_graph_theory_ref4}.

\item \textbf{Vizing's theorem.} The chromatic index of a simple graph is $\Delta$ or $\Delta+1$ \cite{extremal_graph_theory_ref2}.
\end{enumerate}

\theoryapplications
\theoryconnections{extremal graph theory}{%
\item Graph theory supplies vertices, edges, subgraphs, and coloring, while combinatorics supplies counting and probabilistic constructions.
\item Linear algebra and analysis can also estimate eigenvalues, expansion, and pseudorandomness in large graphs.
}{%
\item Extremal graph theory helps network design, theoretical computer science, additive combinatorics, and property testing.
\item Its bounds identify the largest or smallest possible configuration under a forbidden-pattern condition, which is often more useful than classifying every graph.
}
\section*{Glossary}
\begin{description}[leftmargin=*,style=nextline,itemsep=0.3em]
\item[\textbf{Forbidden subgraph.}] A specific graph pattern $H$ that a larger graph $G$ is not allowed to contain; the extremal problem asks for the maximum number of edges such a graph can have.
\item[\textbf{Turán graph.}] The complete $r$-partite graph with parts as equal as possible, which maximizes the number of edges among all $n$-vertex graphs avoiding a clique of size $r+1$.
\item[\textbf{Chromatic number.}] The smallest number of colors needed to assign a color to every vertex so that no two adjacent vertices share the same color.
\item[\textbf{Regularity lemma.}] A result stating that the vertices of any sufficiently large graph can be partitioned into a bounded number of parts so that most pairs of parts behave approximately like random bipartite graphs.
\item[\textbf{Graph homomorphism.}] A map from the vertices of one graph to the vertices of another that preserves edges, allowing different vertices to map to the same place.
\item[\textbf{Graphon.}] A measurable function on the unit square that serves as the limit object for sequences of dense graphs, turning graph properties into integrals.
\item[\textbf{Chromatic index.}] The minimum number of colors needed to color the edges of a graph so that no two edges sharing a vertex receive the same color.
\item[\textbf{Extremal number ex(n,H).}] The maximum number of edges in an $n$-vertex graph containing no copy of forbidden subgraph $H$.
\item[\textbf{Independent set.}] A set of vertices no two of which are adjacent.
\item[\textbf{Edge density.}] The fraction of possible edges present, equal to $|E(G)|/\binom{n}{2}$.
\item[\textbf{Probabilistic method.}] A technique proving an object exists by showing a random construction has positive probability of success.
\end{description}

\section*{7. Sample Questions}
\emph{Exercise reference:} \cite{extremal_graph_theory_ref1}. The cited edition is the source for the exercise set; consult its Exercises section for the official statement and numbering.
\begin{enumerate}[leftmargin=*,itemsep=0.9em]
\item \textbf{Exercise 1.} Why does a complete bipartite graph contain no triangle?

\emph{Solution.} Every edge joins the two different parts. A triangle has three vertices, so two must lie in the same part. Those two have no edge between them, so the triangle cannot occur.

\item \textbf{Exercise 2.} Find the maximum number of edges in a triangle-free graph with $13$ vertices.

\emph{Solution.} Mantel's theorem gives $\lfloor 13^2/4\rfloor=\lfloor169/4\rfloor=42$. A complete bipartite graph with parts of sizes $6$ and $7$ has $6\cdot7=42$ edges, so the bound is attained.

\item \textbf{Exercise 3.} What does the chromatic number of a graph measure?

\emph{Solution.} It is the least number of labels or colors needed so that adjacent vertices receive different labels. For a map, it is the least number of colors needed so neighboring regions have different colors.

\item \textbf{Exercise 4.} A simple graph has maximum degree $7$. What can be said about its edge chromatic number?

\emph{Solution.} Vizing's theorem gives $7\leq\chi'(G)\leq8$. More information about the shape of the graph is needed to decide which value occurs.

\item \textbf{Exercise 5.} Why is the regularity lemma useful even though it does not make a graph truly random?

\emph{Solution.} It supplies a bounded coarse partition. Counts and densities of fixed small patterns can then be estimated from the parts, making a huge graph analyzable without inspecting every edge.

\end{enumerate}
\section*{8. References}
\addcontentsline{toc}{section}{References}

\begin{thebibliography}{9}
\bibitem{extremal_graph_theory_ref1} Noga Alon and Michael Krivelevich, "Extremal and probabilistic combinatorics," in \emph{The Princeton Companion to Mathematics}, Princeton University Press, 2008.
\bibitem{extremal_graph_theory_ref2} Reinhard Diestel, \emph{Graph Theory}, 5th ed., Springer, 2017.
\bibitem{extremal_graph_theory_ref3} Bela Bollobas, \emph{Extremal Graph Theory}, Dover, 2004.
\bibitem{extremal_graph_theory_ref4} Bela Bollobas, \emph{Modern Graph Theory}, Springer, 1998.
\bibitem{extremal_graph_theory_ref5} Douglas B. West, \emph{Introduction to Graph Theory}, 2nd ed., Prentice Hall, 2001.
\bibitem{extremal_graph_theory_ref6} Jozsef Balogh, Robert Morris, and Wojciech Samotij, \emph{The Method of Hypergraph Containers}, American Mathematical Society, 2015.
\end{thebibliography}


\theory{Field theory}{What is a number system in which the four basic arithmetic operations work, and how can one enlarge such a system without losing control of equations?}{Field theory studies fields and their extensions. It explains why some polynomial equations can be solved by familiar numbers, why others require new numbers, and how the symmetries of the solutions reveal the answer.}{Algebraic equations, finite-field arithmetic, coding, cryptography, algebraic geometry, number theory, and mathematical physics.}{Field extensions, minimal polynomials, and splitting fields determine which algebraic operations are possible over a chosen number system.}

\paragraph{Évariste Galois}
Galois made field extensions and their automorphism groups central to the study of polynomial equations. His correspondence between solvability by radicals and solvability of the associated Galois group explains why some equations admit formulas and others do not.

\paragraph{Richard Dedekind and finite-field researchers}
Dedekind clarified the arithmetic of field extensions through embeddings, norms, and algebraic integers. The systematic study of finite fields and extension degrees later supplied the exact arithmetic used in coding, cryptography, and algebraic number theory \cite{field_theory_ref1}.
\end{historylist}
\section*{3. Theory}
\subsection*{Core theory}
\paragraph{What the name means}
The phrase "field theory" has several meanings. In mathematics it means the study of fields, such as the rational numbers, real numbers, complex numbers, and finite fields. In physics, field theory studies quantities spread through space and time: classical field theory, quantum field theory, and statistical field theory are related physical subjects. A grand unified theory is a proposed physical framework joining interactions. In psychology, Kurt Lewin's field theory studies the interaction of a person and environment, while sociological field theory studies social actors and local social orders. This chapter develops the mathematical meaning while keeping these other uses distinct \cite{field_theory_ref1,field_theory_ref4}.

\paragraph{The basic number system}
A field is a set of numbers in which addition, subtraction, multiplication, and division by a nonzero number are always possible, and the familiar rules of arithmetic hold. The rational numbers $\mathbb{Q}$ form a field. The real numbers $\mathbb{R}$ and complex numbers $\mathbb{C}$ form larger fields. The integers do not form a field because, for example, $1/2$ is not an integer.

Finite fields also exist. The integers modulo a prime $p$, written $\mathbb{F}_p$, form a field because every nonzero residue has a multiplicative inverse. Modulo $5$, the inverse of $2$ is $3$, since $2\cdot3=6$ has remainder $1$. Arithmetic in finite fields is exact and bounded, which is why it is useful in digital communication and cryptography.

\paragraph{Extensions: making room for new solutions}
\bookfigure{field_galois_lattice}{A field-extension lattice records how larger number systems contain one another.}
If a field $K$ is contained in a larger field $L$, then $L/K$ is a field extension. The smallest field containing $K$ and an element $\alpha$ is denoted $K(\alpha)$. If every element of $L$ can be written as a combination of finitely many powers of $\alpha$ with coefficients in $K$, the extension is finite. Its degree $[L:K]$ is the number of independent coefficients required.

The equation $x^2-2=0$ has no rational solution. Adjoining $\sqrt2$ produces $\mathbb{Q}(\sqrt2)$, whose elements have the form $a+b\sqrt2$ with $a,b\in\mathbb{Q}$. Since $(\sqrt2)^2=2$, every higher power reduces to one of the first two powers, so the degree is $2$. This is the controlled way in which algebra enlarges the number system.

An element $\alpha$ is algebraic over $K$ if it satisfies a nonzero polynomial with coefficients in $K$. Among all monic polynomials that vanish at $\alpha$, one of least degree is the minimal polynomial. It is irreducible over $K$ and records all algebraic relations that $\alpha$ has over the base field. If no such polynomial exists, $\alpha$ is transcendental. The number $\pi$ is an important example of a transcendental number over $\mathbb{Q}$.

\paragraph{Splitting fields and symmetries}
Given a polynomial $f(x)$ over $K$, its splitting field is the smallest extension in which $f$ factors completely into linear factors. For $x^2-2$, the splitting field is $\mathbb{Q}(\sqrt2)$. For $x^3-2$, the field must contain the real cube root of $2$ and the complex cube roots, so it also contains a primitive cube root of unity.

The automorphisms of a splitting field that fix every element of $K$ form its Galois group. An automorphism is a reversible map preserving addition and multiplication. It may permute the roots of the polynomial, but it must preserve every algebraic relation among them. For $x^2-2$, the two automorphisms send $\sqrt2$ to $\sqrt2$ or $-\sqrt2$. The group therefore has two elements.

Galois theory creates a dictionary: intermediate fields between $K$ and a suitable extension correspond to subgroups of the Galois group. Larger subgroups fix smaller fields, and smaller subgroups fix larger fields. The dictionary changes the problem of solving equations into the problem of understanding symmetries. It explains why general equations of degree five cannot be solved by radicals, whereas equations of degrees two, three, and four can.

\paragraph{Separable and normal extensions}
An algebraic extension is separable when every element has a minimal polynomial with distinct roots. In characteristic zero, every algebraic extension is separable. In positive characteristic, repeated roots can occur. For example, over a field of characteristic $p$, the polynomial $x^p-a$ has derivative zero and may have only one repeated root in a suitable extension.

An extension is normal when every irreducible polynomial over the base field that has one root in the extension splits completely there. A finite extension that is both normal and separable is Galois. The distinction matters because the clean subgroup dictionary belongs to Galois extensions, while inseparable extensions behave differently.

\paragraph{Constructing fields and understanding their size}
Every field has a characteristic. If adding $1$ repeatedly eventually gives $0$, the characteristic is a prime $p$ and the field contains a copy of $\mathbb{F}_p$. Otherwise its characteristic is zero and it contains a copy of $\mathbb{Q}$. A finite field has exactly $p^n$ elements for a prime $p$ and positive integer $n$, and any two fields with the same number of elements are isomorphic. The multiplicative nonzero elements of a finite field form a cyclic group \cite{field_theory_ref2,field_theory_ref3}.

Field extensions can be built by quotienting a polynomial ring by an irreducible polynomial. For example,
\[
\mathbb{Q}[x]/(x^2-2)
\]
behaves like $\mathbb{Q}(\sqrt2)$: the class of $x$ acts as $\sqrt2$, and the relation $x^2=2$ is imposed. Finite fields are constructed similarly from $\mathbb{F}_p[x]$ using irreducible polynomials.

\paragraph{Geometry, number theory, and computation}
Field theory is the algebraic foundation of algebraic geometry. Coordinates of points on a variety form a field or a ring, and extending the field describes new points that become visible after allowing new coordinates. In number theory, number fields are finite extensions of $\mathbb{Q}$; their rings of integers, ideals, places, and Galois groups encode arithmetic information.

Finite fields make arithmetic possible in computers without rounding error. Reed--Solomon codes evaluate polynomials over a finite field and recover the original message when some symbols are corrupted. Elliptic-curve cryptography uses the group of points on an elliptic curve over a finite field; the security comes from the difficulty of reversing a repeated group operation. A field is not itself a cryptosystem or code, but supplies the exact arithmetic on which these systems rely.

\paragraph{What the theory depends on and where it is useful}
Field theory depends on set-based algebra, polynomial arithmetic, vector spaces, and group theory. It supplies the coefficient systems used by ring theory, representation theory, algebraic geometry, number theory, coding theory, and parts of topology. Galois groups connect field theory to group theory; splitting fields connect it to polynomial equations; finite fields connect it to algorithms.

The theory also has limits. A field extension can tell us that roots exist and how they are related without producing a convenient decimal formula. A polynomial may be solvable by radicals in principle while its expression is too large to use. In positive characteristic, separability and derivative arguments require extra care.

\section*{4. Important theorems and results}
\begin{enumerate}[leftmargin=*,itemsep=0.35em]
\item \textbf{Tower law.} If $K\subseteq L\subseteq M$ are finite extensions, then
\[
[M:K]=[M:L][L:K].
\]
It is the arithmetic rule for counting independent coefficients.

\item \textbf{Primitive element theorem.} Many finite separable extensions can be generated by one element: $L=K(\alpha)$ \cite{field_theory_ref1}.

\item \textbf{Galois correspondence.} For a finite Galois extension, intermediate fields and subgroups of the Galois group correspond in reverse order \cite{field_theory_ref2}.

\item \textbf{Fundamental theorem of algebra.} Every nonconstant complex polynomial factors into linear factors over $\mathbb{C}$; thus $\mathbb{C}$ is algebraically closed.

\item \textbf{Finite-field theorem.} For every prime power $p^n$ there is, up to isomorphism, exactly one field with $p^n$ elements, and its multiplicative group is cyclic \cite{field_theory_ref3}.
\end{enumerate}

\theoryapplications
\theoryconnections{field theory}{%
\item Ring theory supplies polynomial quotients and ideals, linear algebra supplies extension degree, and Galois theory studies the symmetries of normal separable extensions.
\item Algebraic geometry and number theory provide geometric and arithmetic examples of fields.
}{%
\item Field theory helps construct solvable equations, finite fields for coding and cryptography, and coordinate fields in algebraic geometry.
\item It tells us which operations preserve a problem and when adjoining a new element creates genuinely new solutions.
}
\section*{Glossary}
\begin{description}[leftmargin=*,style=nextline,itemsep=0.3em]
\item[\textbf{Field extension.}] A larger field $L$ containing a smaller field $K$, written $L/K$, formed by adjoining new elements to $K$ and closing under arithmetic operations.
\item[\textbf{Minimal polynomial.}] The monic polynomial of lowest degree with coefficients in $K$ that vanishes at a given algebraic element; it encodes all algebraic relations that element has over $K$.
\item[\textbf{Splitting field.}] The smallest field extension over which a given polynomial factors completely into linear factors.
\item[\textbf{Galois group.}] The group of field automorphisms of a splitting field that fix every element of the base field; it records the symmetries among the roots.
\item[\textbf{Separable extension.}] An algebraic extension in which every element has a minimal polynomial with distinct roots; in characteristic zero every extension is separable.
\item[\textbf{Normal extension.}] An algebraic extension in which every irreducible polynomial over the base field that has one root in the extension splits completely there.
\item[\textbf{Algebraic element.}] An element satisfying a nonzero polynomial equation with coefficients in the base field, as opposed to a transcendental element like $\pi$.
\item[\textbf{Finite field.}] A field with finitely many elements, necessarily $p^n$ for a prime $p$; its nonzero elements form a cyclic multiplicative group.
\item[\textbf{Degree of an extension.}] The dimension $[L\!:\!K]$ of $L$ as a vector space over $K$.
\item[\textbf{Characteristic.}] The smallest positive integer $n$ such that $n\cdot1=0$ in a field, or zero if none exists.
\item[\textbf{Transcendental element.}] An element not algebraic over the base field, satisfying no nonzero polynomial with coefficients in the base field.
\item[\textbf{Irreducible polynomial.}] A nonconstant polynomial that cannot be factored as a product of two nonconstant polynomials over the base field.
\item[\textbf{Intermediate field.}] A field $E$ satisfying $K\subseteq E\subseteq L$; central to the Galois correspondence.
\end{description}
\section*{7. Sample Questions}
\emph{Exercise reference:} \cite{field_theory_ref1}. The cited edition is the source for the exercise set; consult its Exercises section for the official statement and numbering.
\begin{enumerate}[leftmargin=*,itemsep=0.9em]
\item \textbf{Exercise 1.} Find the degree $[\mathbb{Q}(\sqrt{2},\sqrt{3}):\mathbb{Q}]$ and a basis for $\mathbb{Q}(\sqrt{2},\sqrt{3})$ over $\mathbb{Q}$.

\emph{Solution.} Since $\sqrt{3}\notin\mathbb{Q}(\sqrt{2})$ (if $\sqrt{3}=a+b\sqrt{2}$ then $3=a^2+2b^2+2ab\sqrt{2}$, forcing $ab=0$ and yielding a contradiction), the extension $\mathbb{Q}(\sqrt{2},\sqrt{3})/\mathbb{Q}$ has degree $[\mathbb{Q}(\sqrt{2},\sqrt{3}):\mathbb{Q}(\sqrt{2})]\cdot[\mathbb{Q}(\sqrt{2}):\mathbb{Q}]=2\cdot 2=4$. A basis is $\{1,\sqrt{2},\sqrt{3},\sqrt{6}\}$.

\item \textbf{Exercise 2.} Find the minimal polynomial of $1+\sqrt{2}$ over $\mathbb{Q}$.

\emph{Solution.} Set $\alpha=1+\sqrt{2}$, so $\alpha-1=\sqrt{2}$ and $(\alpha-1)^2=2$, giving $\alpha^2-2\alpha-1=0$. The polynomial $x^2-2x-1$ is irreducible over $\mathbb{Q}$ (no rational roots by the rational root theorem), so it is the minimal polynomial.

\item \textbf{Exercise 3.} Find the multiplicative inverse of $3$ in $\mathbb{F}_7$ and of $x+1$ in $\mathbb{F}_2[x]/(x^3+x+1)$.

\emph{Solution.} In $\mathbb{F}_7$: $3\cdot 5=15\equiv 1\pmod{7}$, so $3^{-1}=5$. In $\mathbb{F}_2[x]/(x^3+x+1)$: perform polynomial long division of $x^3+x+1$ by $x+1$. We get $x^3+x+1=(x+1)(x^2+x)+1$, so $(x+1)(x^2+x)\equiv 1$, meaning $(x+1)^{-1}=x^2+x$.

\item \textbf{Exercise 4.} Let $f(x)=x^3-2$ over $\mathbb{Q}$. Prove that $[\mathbb{Q}(\sqrt[3]{2}):\mathbb{Q}]=3$.

\emph{Solution.} The polynomial $x^3-2$ is irreducible over $\mathbb{Q}$ by Eisenstein's criterion with $p=2$. Therefore its minimal polynomial over $\mathbb{Q}$ has degree $3$, so $[\mathbb{Q}(\sqrt[3]{2}):\mathbb{Q}]=3$.

\item \textbf{Exercise 5.} Prove that every finite field has order $p^n$ for some prime $p$ and positive integer $n$.

\emph{Solution.} A finite field $F$ has characteristic $p>0$ (since it is finite, $1+1+\cdots+1=0$ for some sum). Then $F$ contains $\mathbb{F}_p$ as a subfield. $F$ is a vector space over $\mathbb{F}_p$, say of dimension $n$. A basis $\{e_1,\ldots,e_n\}$ gives $|F|=p^n$, since every element is a unique $\mathbb{F}_p$-linear combination of the basis elements.

\end{enumerate}
\section*{8. References}
\addcontentsline{toc}{section}{References}

\begin{thebibliography}{9}
\bibitem{field_theory_ref1} David S. Dummit and Richard M. Foote, \emph{Abstract Algebra}, 3rd ed., Wiley, 2004.
\bibitem{field_theory_ref2} Ian Stewart, \emph{Galois Theory}, 4th ed., Chapman and Hall/CRC, 2015.
\bibitem{field_theory_ref3} Rudolf Lidl and Harald Niederreiter, \emph{Finite Fields}, 2nd ed., Cambridge University Press, 1997.
\bibitem{field_theory_ref4} Jean-Pierre Serre, \emph{Local Fields}, Springer, 1979.
\end{thebibliography}


\theory{Galois theory}{Can the symmetries of the roots of an equation tell us whether its roots have a formula using radicals?}{Galois theory connects field extensions with groups of symmetries. It replaces the difficult task of manipulating roots by the more organized task of studying the permutations that preserve every algebraic relation among them.}{Polynomial equations, geometric constructions, number theory, finite fields, coding, and modern algebra.}{Automorphism groups record how roots can be permuted, converting solvability of equations into a question about group structure.}

\paragraph{The question that created the theory}
For quadratic equations, the quadratic formula expresses the roots using addition, multiplication, division, and square roots. Cardano and Ferrari found corresponding methods for cubics and quartics. Niels Henrik Abel proved that no formula of the same kind can solve every polynomial of degree five or more. The question then changed: which particular equations can be solved by radicals, and why? Évariste Galois supplied the answer by studying the symmetries of the roots \cite{galois_theory_ref1,galois_theory_ref2}.

The history grew from older work on symmetric functions. The coefficients of a monic polynomial are symmetric expressions in its roots, as in $(x-a)(x-b)=x^2-(a+b)x+ab$. Lagrange studied permutations of roots and resolvent polynomials; Ruffini and Abel established the general impossibility for the quintic; Galois identified the exact group-theoretic criterion. His memoir was published by Joseph Liouville after Galois's death and was difficult for contemporaries to understand.

\end{historylist}
\section*{3. Theory}
\subsection*{Core theory}
\paragraph{The symmetry of a quadratic}
\bookfigure{field_galois_lattice}{Galois theory studies the symmetries of field extensions while fixing the base field.}
Consider $x^2-4x+1=0$. Its roots are $A=2+\sqrt3$ and $B=2-\sqrt3$. Both satisfy rational relations such as $A+B=4$ and $AB=1$. Swapping $A$ and $B$ preserves every relation whose coefficients are rational. The identity and the swap form a group of two elements.

The field $\mathbb{Q}(\sqrt3)$ contains every number obtained from rational numbers and $\sqrt3$ by arithmetic. An automorphism of this field that fixes $\mathbb{Q}$ may send $\sqrt3$ to $\sqrt3$ or $-\sqrt3$. These are exactly the two symmetries above. The Galois group is therefore the group of order two.

\paragraph{A quartic with four sign changes}
The polynomial $x^4-10x^2+1$ has roots
\[
A=\sqrt2+\sqrt3,\quad B=\sqrt2-\sqrt3,\quad
C=-\sqrt2+\sqrt3,\quad D=-\sqrt2-\sqrt3.
\]
Changing the sign of $\sqrt2$, changing the sign of $\sqrt3$, or changing both produces four valid automorphisms. They form the Klein four-group. Relations such as $AB=-1$, $AC=1$, and $A+D=0$ show that the image of $A$ determines the images of all four roots. This is the central practical idea: an allowed symmetry must preserve all relations over the base field \cite{galois_theory_ref2}.

\paragraph{The modern field-theoretic picture}
Start with a field extension $L/K$. The Galois group $\operatorname{Gal}(L/K)$ consists of all automorphisms of $L$ that leave every element of $K$ fixed. If $L$ is the splitting field of a polynomial over $K$, these automorphisms permute its roots. The modern formulation works over rational, number, finite, or local fields and also handles infinite and inseparable extensions.

For a finite Galois extension, the fundamental theorem gives a reverse-order correspondence between intermediate fields $K\subseteq E\subseteq L$ and subgroups of $\operatorname{Gal}(L/K)$. The field fixed by a subgroup $H$ is
\[
L^H=\{x\in L:h(x)=x\text{ for every }h\in H\}.
\]
Large subgroups fix small fields; small subgroups fix large fields. Normal subgroups correspond to intermediate extensions that are themselves Galois.

\paragraph{Solvability by radicals}
An equation is solvable by radicals if its roots can be formed from the coefficients using arithmetic operations and repeated extraction of $n$th roots. A finite group is solvable when it can be built through a chain whose successive quotients are cyclic abelian groups. Galois's theorem states that a polynomial over a field of characteristic zero is solvable by radicals precisely when its Galois group is solvable, after the usual roots of unity are included \cite{galois_theory_ref1}.

Every polynomial of degree at most four has a solvable Galois group. For the general polynomial of degree $n\geq5$, the symmetric group $S_n$ appears, and it contains the nonsolvable alternating group $A_n$. Hence there is no universal radical formula. Particular quintics may still be solvable; the theorem is about the general equation, not every equation of that degree.

\paragraph{Classical geometry}
Compass-and-straightedge constructions correspond to field extensions whose degrees are powers of two. A cube cannot be doubled because it would require constructing $\sqrt[3]{2}$, which has degree three. A general angle cannot be trisected because the required cubic need not have a constructible solution. A regular polygon with $n$ sides is constructible exactly when $n$ is a product of a power of two and distinct Fermat primes. Galois theory proves the completeness of this characterization, not merely the construction of known examples \cite{galois_theory_ref3}.

\paragraph{Separable, normal, and infinite extensions}
An algebraic extension is separable when every element has a minimal polynomial with distinct roots. It is normal when every irreducible polynomial over the base field with one root in the extension splits there. A finite extension that is both normal and separable is Galois. In characteristic zero separability is automatic; in positive characteristic an equation such as $x^p-a$ can have repeated roots.

The absolute Galois group of a field $K$ is the Galois group of an algebraic closure of $K$ over $K$. This lets number theorists study all finite extensions at once. Galois theory has also been generalized to Galois connections and Grothendieck's Galois theory, which captures geometric and arithmetic coverings.

\paragraph{What the theory depends on and where it is useful}
Galois theory depends on field theory, polynomial arithmetic, permutation groups, and the theory of solvable groups. It helps algebraic number theory through absolute Galois groups, algebraic geometry through coverings, and coding theory through finite-field extensions. It gives a conceptual test even when an explicit formula would be enormous.

\section*{4. Important theorems and results}
\begin{enumerate}[leftmargin=*,itemsep=0.35em]
\item \textbf{Fundamental theorem of Galois theory.} Intermediate fields of a finite Galois extension correspond contravariantly to subgroups of its Galois group \cite{galois_theory_ref1}.

\item \textbf{Galois solvability criterion.} A polynomial over characteristic zero is solvable by radicals exactly when its Galois group is solvable.

\item \textbf{Abel--Ruffini theorem.} There is no radical formula for the general polynomial of degree at least five.

\item \textbf{Constructibility criterion.} A real number constructible by compass and straightedge lies in an extension of degree a power of two.

\item \textbf{Discriminant test.} The discriminant is zero exactly when a polynomial has a repeated root; it also gives useful information about the real and complex roots of low-degree polynomials \cite{galois_theory_ref4}.
\end{enumerate}

\theoryapplications
\theoryconnections{galois theory}{%
\item Field theory supplies extensions and minimal polynomials, group theory supplies automorphisms and subgroup structure, and the fundamental theorem of algebra supplies the roots being permuted.
\item Number theory contributes finite fields, cyclotomic fields, and arithmetic examples.
}{%
\item Galois theory helps decide whether polynomial equations can be solved by radicals and helps construct finite fields and field extensions.
\item Its correspondence between subgroups and intermediate fields also transfers calculations between algebraic equations and symmetry.
}
\section*{Glossary}
\begin{description}[leftmargin=*,style=nextline,itemsep=0.3em]
\item[\textbf{Galois group.}] The group of field automorphisms of a Galois extension that fix the base field; its structure determines whether a polynomial is solvable by radicals.
\item[\textbf{Splitting field.}] The smallest extension over which a polynomial factors completely into linear factors.
\item[\textbf{Automorphism.}] A bijective map from a field to itself that preserves addition and multiplication; in Galois theory it must also fix every element of the base field.
\item[\textbf{Solvable group.}] A group that can be built from a chain of subgroups whose successive quotients are cyclic abelian groups; a polynomial is solvable by radicals precisely when its Galois group is solvable.
\item[\textbf{Intermediate field.}] A field $E$ lying between the base field $K$ and a larger extension $L$, forming part of the correspondence in the fundamental theorem of Galois theory.
\item[\textbf{Separable extension.}] An algebraic extension in which every minimal polynomial has distinct roots; this is automatic in characteristic zero.
\item[\textbf{Normal extension.}] An algebraic extension where every irreducible polynomial over the base field with one root in the extension splits completely there.
\item[\textbf{Solvable by radicals.}] A polynomial whose roots can be expressed using arithmetic operations and repeated extraction of $n$th roots.
\item[\textbf{Discriminant.}] A polynomial invariant computed from the roots that vanishes when a polynomial has a repeated root.
\item[\textbf{Galois correspondence.}] The reverse-order bijection between intermediate fields of a finite Galois extension and subgroups of its Galois group.
\item[\textbf{Absolute Galois group.}] The Galois group of an algebraic closure over the field itself; encodes all finite extensions simultaneously.
\end{description}
\section*{7. Sample Questions}
\emph{Exercise reference:} \cite{galois_theory_ref1}. The cited edition is the source for the exercise set; consult its Exercises section for the official statement and numbering.
\begin{enumerate}[leftmargin=*,itemsep=0.9em]
\item \textbf{Exercise 1.} Compute the Galois group of $x^2-4x+1$ over $\mathbb{Q}$.

\emph{Solution.} The roots are $2+\sqrt{3}$ and $2-\sqrt{3}$, so the splitting field is $\mathbb{Q}(\sqrt{3})$ with $[\mathbb{Q}(\sqrt{3}):\mathbb{Q}]=2$. The two automorphisms fixing $\mathbb{Q}$ are the identity and $\sqrt{3}\mapsto -\sqrt{3}$. Thus $\operatorname{Gal}\cong\mathbb{Z}/2\mathbb{Z}$.

\item \textbf{Exercise 2.} Determine the Galois group of $x^4-10x^2+1$ over $\mathbb{Q}$.

\emph{Solution.} The roots are $\pm\sqrt{2}\pm\sqrt{3}$. The splitting field is $\mathbb{Q}(\sqrt{2},\sqrt{3})$ with $[\mathbb{Q}(\sqrt{2},\sqrt{3}):\mathbb{Q}]=4$. The four automorphisms independently change the signs of $\sqrt{2}$ and $\sqrt{3}$. Since both generators have order $2$ and commute, $\operatorname{Gal}\cong\mathbb{Z}/2\mathbb{Z}\times\mathbb{Z}/2\mathbb{Z}$, the Klein four-group.

\item \textbf{Exercise 3.} List all intermediate fields of the Galois extension $\mathbb{Q}(\sqrt{2},\sqrt{3})/\mathbb{Q}$ using the Galois correspondence.

\emph{Solution.} The Galois group $G\cong\mathbb{Z}/2\times\mathbb{Z}/2$ has three subgroups of order $2$: $\langle\sigma\rangle$ (sending $\sqrt{2}\mapsto-\sqrt{2}$, $\sqrt{3}\mapsto\sqrt{3}$), $\langle\tau\rangle$ (sending $\sqrt{2}\mapsto\sqrt{2}$, $\sqrt{3}\mapsto-\sqrt{3}$), and $\langle\sigma\tau\rangle$ (sending both $\sqrt{2}\mapsto-\sqrt{2}$, $\sqrt{3}\mapsto-\sqrt{3}$). Their fixed fields are $\mathbb{Q}(\sqrt{3})$, $\mathbb{Q}(\sqrt{2})$, and $\mathbb{Q}(\sqrt{6})$, respectively. Together with $\mathbb{Q}$ and $\mathbb{Q}(\sqrt{2},\sqrt{3})$, these are all five intermediate fields.

\item \textbf{Exercise 4.} Show that $x^5-4x+2$ is irreducible over $\mathbb{Q}$ and determine whether it is solvable by radicals.

\emph{Solution.} By Eisenstein's criterion with $p=2$: the polynomial $x^5-4x+2$ has all coefficients except the leading one divisible by $2$, and the constant term $2$ is not divisible by $4$. So it is irreducible, and $\operatorname{Gal}$ acts transitively on its $5$ roots, hence contains a $5$-cycle. The discriminant can be computed to be $508000=2^6\cdot5^3\cdot127$, which is not a perfect square, so $\operatorname{Gal}\not\subseteq A_5$. A transitive subgroup of $S_5$ containing an odd permutation is $S_5$ itself (or $F_{20}$, but checking the discriminant mod various primes shows $\operatorname{Gal}\cong S_5$). Since $S_5$ is not solvable, this polynomial is not solvable by radicals.

\item \textbf{Exercise 5.} Determine whether $\sqrt[3]{2}$ is constructible by compass and straightedge.

\emph{Solution.} The minimal polynomial of $\sqrt[3]{2}$ over $\mathbb{Q}$ is $x^3-2$ (irreducible by Eisenstein), so $[\mathbb{Q}(\sqrt[3]{2}):\mathbb{Q}]=3$. A real number is constructible only if the degree $[\mathbb{Q}(\alpha):\mathbb{Q}]$ is a power of $2$. Since $3$ is not a power of $2$, $\sqrt[3]{2}$ is not constructible. Equivalently, the cube cannot be doubled.

\end{enumerate}
\section*{8. References}
\addcontentsline{toc}{section}{References}

\begin{thebibliography}{9}
\bibitem{galois_theory_ref1} Ian Stewart, \emph{Galois Theory}, 4th ed., Chapman and Hall/CRC, 2015.
\bibitem{galois_theory_ref2} Emil Artin, \emph{Galois Theory}, Dover, 1998.
\bibitem{galois_theory_ref3} David A. Cox, \emph{Galois Theory}, 2nd ed., Wiley, 2012.
\bibitem{galois_theory_ref4} Jean-Pierre Serre, \emph{Topics in Galois Theory}, 2nd ed., A K Peters, 2008.
\end{thebibliography}


\theory{Game theory}{How should a decision-maker act when the result depends on what other decision-makers do?}{Game theory models strategic interaction. A game is not necessarily entertainment: it is any situation with players, choices, information, and consequences that depend on several choices at once.}{Economics, business, politics, biology, computer science, artificial intelligence, ethics, and social science.}{Payoff functions, strategies, and equilibrium conditions model how each participant's choice changes the incentives of the others.}

\paragraph{The problem and its history}
Ordinary optimization asks one person to choose the best action. Strategic optimization is harder because the best action changes when another person changes theirs. Cournot analyzed competition between firms in 1838. Cardano, Pascal, Fermat, Huygens, and Waldegrave had earlier studied games of chance. John von Neumann proved the minimax theorem for two-person zero-sum games; his work with Oskar Morgenstern developed the economics of cooperative games and expected utility. Nash later proved that every finite game has an equilibrium in mixed strategies \cite{game_theory_ref1,game_theory_ref2}.

\end{historylist}
\section*{3. Theory}
\subsection*{Core theory}
\paragraph{How a game is described}
The players are the decision-makers. Each has a set of actions or strategies. A payoff records the benefit or cost of every combination of choices. In a normal-form game, a payoff matrix lists these outcomes. In an extensive-form game, a decision tree records time order, observations, and information sets.

In a zero-sum game one player's gain is exactly the other's loss. Chess between perfectly matched players is an idealized example. In a non-zero-sum game, both can gain through cooperation or both can lose. A constant-sum game is equivalent to a zero-sum game after adding a dummy account that receives the remaining payoff.

\paragraph{A small payoff matrix}
Two shops choose either a high or low price. If both choose high, each earns 4. If one chooses low while the other chooses high, the low-price shop earns 6 and the high-price shop earns 1. If both choose low, each earns 2. A shop cannot choose its best price in isolation: it must anticipate the competitor.

A strategy is dominant if it is best regardless of what the other player does. A Nash equilibrium is a collection of strategies in which no player can improve by changing alone. Equilibrium does not mean that the outcome is socially best or fair; it means only that unilateral departure is unattractive.

\paragraph{Important classifications}
In a simultaneous game, players choose without observing the current choice of others, so a matrix is natural. In a sequential game, earlier actions affect later choices, so a decision tree is natural. Perfect information means every previous action is observed, as in chess. Imperfect information means some actions are hidden, as in poker. Complete information is different: it means that the players know the rules, possible strategies, and payoffs, even if they do not observe every move \cite{game_theory_ref3}.

Cooperative game theory allows binding agreements and studies coalitions. A characteristic function $v(S)$ assigns the total value that a coalition $S$ can create, with $v(\emptyset)=0$. The core consists of allocations that no coalition can improve upon. The Shapley value divides a cooperative surplus according to average marginal contribution.

Bayesian games model incomplete information. Each player has a type, such as a firm's cost or a bidder's valuation, and beliefs about the types of others. A move by nature can represent the unknown type, turning incomplete information into imperfect information. Signaling games study how informed players communicate through actions.

\paragraph{Equilibrium and strategic reasoning}
A best response is a strategy that maximizes a player's payoff when the other strategies are fixed. Nash equilibrium requires every chosen strategy to be a best response. In sequential games, subgame-perfect equilibrium removes threats that would not be credible after a particular history. Repeated games allow reputation and punishment to support cooperation. Evolutionary games replace perfect rationality by changing frequencies of successful strategies.

Von Neumann's minimax theorem says that in a finite two-player zero-sum game, the maximum guaranteed payoff equals the minimum loss that the opponent can force, once mixed strategies are allowed. The proof uses the Brouwer fixed-point theorem. This is why randomization can be rational: it prevents the opponent from predicting a profitable response.

\paragraph{Uses outside economics}
Firms use Cournot and Bertrand models to study output and prices, and auctions use Bayesian and mechanism-design models to align private information with efficient outcomes. Governments use games for voting, bargaining, fair division, public choice, and war negotiation. Project managers model investors, contractors, subcontractors, and governments whose decisions affect a common project.

Evolutionary game theory explains why behavior that costs one animal may spread if it benefits relatives. Hamilton's rule is $br>c$, where $c$ is the cost, $b$ the benefit, and $r$ the relatedness. It models kin selection, alarm calls, food sharing, territorial conflict, and the evolution of cooperation \cite{game_theory_ref4}.

Computer scientists use games to model interactive computation, multi-agent systems, online algorithms, auctions, peer-to-peer systems, and security. Yao's principle turns a randomized algorithm into a game between an algorithm and an adversary, producing lower bounds. Game semantics interprets logic through a contest between a verifier and a falsifier. AI uses games for reinforcement learning, autonomous agents, mechanism design, and anticipating other agents.

\paragraph{Alternative representations}
Different applications need different encodings. Congestion games describe costs that depend on how many users choose a resource. Graphical games record local dependence among players. Timed and continuous games include time or continuous actions. Logic-based game-description languages, Petri-net games, action-graph games, and local-effect games represent structured multi-agent interaction. These descriptions save space when a full payoff table would be enormous \cite{game_theory_ref5}.

\paragraph{What the theory depends on and where it is useful}
Game theory depends on probability, optimization, linear algebra, fixed-point theorems, and decision theory. It helps economics by making competition explicit, computer science by explaining distributed decisions, biology by modeling selection, and philosophy by analyzing conventions, meaning, signaling, and social norms.

Its main limitation is modeling. Real people may not be perfectly rational, may misunderstand the rules, and may care about fairness rather than only material payoff. Experiments in the centipede, dictator, and guessing games often depart from Nash predictions. Evolutionary and bounded-rationality models address some of this, but no game model replaces careful observation \cite{game_theory_ref2}.

\section*{4. Important theorems and results}
\begin{enumerate}[leftmargin=*,itemsep=0.35em]
\item \textbf{Minimax theorem.} Every finite two-person zero-sum game has a value when mixed strategies are allowed.

\item \textbf{Nash existence theorem.} Every finite game has at least one Nash equilibrium in mixed strategies \cite{game_theory_ref1}.

\item \textbf{Backward induction.} In a finite perfect-information game with no chance moves, rational players can solve a subgame-perfect outcome by analyzing the final decisions first.

\item \textbf{Shapley value.} In a cooperative game, averaging a player's marginal contribution over all joining orders gives a fair allocation satisfying standard efficiency, symmetry, dummy-player, and additivity properties.

\item \textbf{Hamilton's rule.} Kin-directed behavior can be favored when its benefit weighted by relatedness exceeds its cost.
\end{enumerate}

\theoryapplications
\theoryconnections{game theory}{%
\item Probability models mixed strategies, optimization studies best responses, and linear algebra handles finite payoff matrices.
\item Economics supplies utility and incentives, while algorithms and graph theory enter for large games and network interactions.
}{%
\item Game theory helps economics, auctions, bargaining, voting, biology, and computer science analyze decisions whose outcomes depend on other decision makers.
\item Equilibrium is a prediction of mutual consistency, not necessarily a socially best outcome, so the solution concept must match the application.
}
\section*{Glossary}
\begin{description}[leftmargin=*,style=nextline,itemsep=0.3em]
\item[\textbf{Nash equilibrium.}] A set of strategies in which no player can improve their payoff by unilaterally changing their own strategy.
\item[\textbf{Dominant strategy.}] A strategy that yields the best outcome for a player regardless of what the other players choose.
\item[\textbf{Zero-sum game.}] A game in which one player's gain is exactly the other player's loss; the total payoff is constant.
\item[\textbf{Mixed strategy.}] A probability distribution over a player's possible actions, allowing randomization rather than committing to a single deterministic choice.
\item[\textbf{Subgame-perfect equilibrium.}] A refinement of Nash equilibrium for sequential games that requires equilibrium behavior after every possible history, eliminating non-credible threats.
\item[\textbf{Cooperative game.}] A game in which players may form binding agreements and coalitions, analyzed by how value is created and divided among groups.
\item[\textbf{Payoff function.}] A function assigning a numerical value to each combination of players' strategies, representing the outcome of the game for every player.
\item[\textbf{Normal form.}] A game representation listing players, strategies, and a payoff matrix for every combination of choices.
\item[\textbf{Extensive form.}] A game representation using a decision tree recording time order, information sets, and choices.
\item[\textbf{Backward induction.}] A method of solving finite perfect-information games by analyzing final decisions first and reasoning backward.
\item[\textbf{Best response.}] A strategy maximizing a player's payoff given the strategies chosen by all other players.
\item[\textbf{Shapley value.}] A cooperative-game solution assigning each player the average of their marginal contributions over all joining orders.
\end{description}
\section*{7. Sample Questions}
\emph{Exercise reference:} \cite{game_theory_ref1}. The cited edition is the source for the exercise set; consult its Exercises section for the official statement and numbering.
\begin{enumerate}[leftmargin=*,itemsep=0.9em]
\item \textbf{Exercise 1.} Find all pure-strategy Nash equilibria of the following $2\times 2$ zero-sum game: Player I chooses rows, Player II chooses columns, and the payoff matrix (to Player I) is $\begin{pmatrix}3&0\\5&1\end{pmatrix}$.

\emph{Solution.} For each cell, check if the row is a best response against that column and vice versa. Against column 1, Player I prefers row 2 ($5>3$). Against column 2, Player I prefers row 2 ($1>0$). So row 2 dominates row 1. Against row 2, Player II prefers column 2 ($1<0$ means II loses less at column 2). The unique pure-strategy Nash equilibrium is (row 2, column 2) with payoff $1$ to Player I.

\item \textbf{Exercise 2.} Compute the mixed-strategy Nash equilibrium of the game with payoff matrix $\begin{pmatrix}2&0\\0&3\end{pmatrix}$ (zero-sum, payoffs to Player I).

\emph{Solution.} Let Player I play row 1 with probability $p$. Player II is indifferent when the expected payoff of column 1 equals that of column 2: $2p+0(1-p)=0p+3(1-p)$, so $2p=3-3p$ and $p=3/5$. Let Player II play column 1 with probability $q$. Player I is indifferent when $2q=3(1-q)$, giving $q=3/5$. The equilibrium mixed strategies are $(3/5,2/5)$ for each player, with value $6/5$.

\item \textbf{Exercise 3.} Find the mixed-strategy Nash equilibrium of Matching Pennies: Player I and Player II each choose $H$ or $T$. Player I wins \$1 if the choices match and loses \$1 otherwise (zero-sum).

\emph{Solution.} The payoff matrix to Player I is $\begin{pmatrix}1&-1\\-1&1\end{pmatrix}$ where rows are Player I's choices and columns are Player II's. No pure-strategy NE exists (checking all four cells shows each player wants to deviate). Let Player I play $H$ with probability $p$ and Player II play $H$ with probability $q$. Player I is indifferent when $E_I(H)=E_I(T)$: $q(1)+(1-q)(-1)=q(-1)+(1-q)(1)$, giving $2q-1=1-2q$, so $q=1/2$. Player II is indifferent when $E_{II}(H)=E_{II}(T)$: $p(-1)+(1-p)(1)=p(1)+(1-p)(-1)$, giving $1-2p=2p-1$, so $p=1/2$. The unique NE is both players randomize uniformly, with game value $0$.

\item \textbf{Exercise 4.} In the ultimatum game, Player I proposes a split of \$10 in integer amounts: keep $s$ and offer $10-s$ to Player II. Player II accepts or rejects; if rejected, both get \$0. Find the subgame-perfect equilibrium using backward induction.

\emph{Solution.} Player II accepts any offer $\geq 1$ (since \$1>\$0). Player I, anticipating this, offers the smallest amount Player II accepts: $s=9$, offering Player II \$1. The subgame-perfect equilibrium is: Player I offers \$(9,1); Player II accepts any positive offer. The outcome is Player I gets \$9, Player II gets \$1.

\item \textbf{Exercise 5.} In a coordination game with payoff matrix $\begin{pmatrix}2&0\\0&1\end{pmatrix}$ (both players choose simultaneously, payoffs to both players are the same---a pure coordination game), find all Nash equilibria.

\emph{Solution.} There are three Nash equilibria: two pure and one mixed. Pure: both choose row 1/column 1 (payoff 2 each) and both choose row 2/column 2 (payoff 1 each). Mixed: Player I plays row 1 with probability $p$ and Player II plays column 1 with probability $q$. Indifference for Player I: $2q=1(1-q)$, so $q=1/3$. By symmetry $p=1/3$. The mixed equilibrium has payoff $2/3$ each. All three are Nash equilibria; the payoff-dominant one is (1,1).

\end{enumerate}
\section*{8. References}
\addcontentsline{toc}{section}{References}

\begin{thebibliography}{9}
\bibitem{game_theory_ref1} Martin J. Osborne and Ariel Rubinstein, \emph{A Course in Game Theory}, MIT Press, 1994.
\bibitem{game_theory_ref2} Robert Gibbons, \emph{Game Theory for Applied Economists}, Princeton University Press, 1992.
\bibitem{game_theory_ref3} Drew Fudenberg and Jean Tirole, \emph{Game Theory}, MIT Press, 1991.
\bibitem{game_theory_ref4} John Maynard Smith, \emph{Evolution and the Theory of Games}, Cambridge University Press, 1982.
\bibitem{game_theory_ref5} Noam Nisan, Tim Roughgarden, Eva Tardos, and Vijay V. Vazirani, eds., \emph{Algorithmic Game Theory}, Cambridge University Press, 2007.
\end{thebibliography}


\theory{Graph theory}{How can relationships among objects be represented, counted, and navigated?}{Graph theory studies vertices joined by edges. It gives a language for networks in which geometry is less important than connection: who is linked to whom, which routes exist, and which patterns must occur.}{Roads, computer networks, social networks, biology, scheduling, search engines, electrical circuits, and algorithms.}{Vertices, edges, paths, and connectivity reduce network questions to finite combinatorial structures that can be counted or searched.}

\paragraph{The idea and its history}
In a graph, vertices represent objects and edges represent relationships. Edges may be undirected, directed, weighted, colored, or repeated. A graph of a function is a different object; this chapter concerns relational graphs. Leonhard Euler's 1736 analysis of the bridges of Konigsberg is usually regarded as the beginning of graph theory. He showed that the question was not about distances but about the number of times each land region was connected to a bridge.

A walk follows edges and may repeat vertices. A trail does not repeat edges. A path does not repeat vertices. A circuit returns to its starting point, and a connected graph has a path between every pair of vertices. The degree of a vertex is the number of incident edges. A tree is a connected graph with no cycle; it has exactly one path between any two vertices and $n-1$ edges when it has $n$ vertices \cite{graph_theory_ref1,graph_theory_ref2}.

\end{historylist}
\section*{3. Theory}
\subsection*{Core theory}
\paragraph{A concrete network example}
Suppose a delivery company has warehouses as vertices and roads as edges. A shortest-path algorithm finds the least travel distance. A spanning tree connects every warehouse with no unnecessary cycle; a minimum spanning tree does so at minimum total construction cost. If a road fails, edge connectivity tells how many failures can be tolerated. If one warehouse is a bottleneck, vertex connectivity identifies the vulnerability.
\bookfigure{graph_network}{A graph represents objects, relationships, and a selected route.}

The same graph supports different questions. Breadth-first search discovers all vertices at increasing distance. Depth-first search detects cycles and builds a traversal tree. Dijkstra's algorithm finds shortest paths with nonnegative weights, while Bellman--Ford allows negative edge weights and detects a negative cycle. A maximum-flow algorithm finds how much material can pass from a source to a sink, and the max-flow min-cut theorem says that this equals the capacity of the narrowest separating cut.

\paragraph{Important graph families}
Complete graphs join every pair of vertices. Bipartite graphs split vertices into two parts with edges only across the split; they contain no odd cycle. Planar graphs can be drawn without crossings. Trees, forests, cycles, paths, regular graphs, and directed acyclic graphs are basic families. A multigraph permits repeated edges, while a simple graph has neither loops nor repeated edges.

An Eulerian circuit uses every edge exactly once. A connected graph has one precisely when every vertex has even degree. A Hamiltonian cycle visits every vertex exactly once; no equally simple degree test characterizes all Hamiltonian graphs, and deciding whether one exists is computationally difficult. This contrast illustrates why edges and vertices create different kinds of problems.

\paragraph{Coloring and matching}
A proper vertex coloring assigns different colors to adjacent vertices. The least number is the chromatic number. Scheduling exams with a shared-student conflict graph turns coloring into a timetable problem. A matching is a set of edges with no shared endpoint. In a bipartite job-assignment graph, a matching pairs workers with tasks. Hall's marriage theorem says that a complete matching exists exactly when every group of workers has at least as many neighboring tasks as its size.

Edge coloring assigns colors to edges so adjacent edges differ. Vizing's theorem says that a simple graph needs either $\Delta$ or $\Delta+1$ colors, where $\Delta$ is the maximum degree. The four-color theorem states that every planar map can be colored with at most four colors.

\paragraph{Representing and measuring graphs}
An adjacency matrix records whether each ordered pair is connected. It makes matrix multiplication useful for counting walks but uses much memory for a sparse graph. An adjacency list stores only neighbors and is efficient for sparse networks. The incidence matrix records which vertices touch which edges.

The degree sequence lists vertex degrees. The adjacency matrix has eigenvalues that reveal expansion, connectivity, and random-walk behavior; this is the subject of spectral graph theory. A graph invariant is a quantity unchanged by relabeling vertices, such as the number of connected components, chromatic number, or degree multiset. Graph isomorphism asks whether two differently labeled graphs have the same structure.

\paragraph{Planarity, surfaces, and geometry}
Euler's formula for a connected planar graph is $V-E+F=2$, where $F$ counts faces including the outside face. Kuratowski's theorem characterizes nonplanar graphs through subdivisions of $K_5$ and $K_{3,3}$. Topological graph theory studies embeddings on surfaces, while geometric graph theory places vertices in space and uses distances or crossings.

\paragraph{Applications}
In sociology, acquaintance, influence, rumor, and collaboration graphs support social-network analysis. In biology, vertices can represent habitats and edges migration routes, allowing disease spread and conservation corridors to be studied. Molecular biology uses graphs for metabolic pathways, gene regulation, protein interactions, cell clustering, and connectomics, where neurons are vertices and synapses are edges \cite{graph_theory_ref3}.

Computer networks use routing, spanning trees, flow, and connectivity. Search engines use links and random walks to rank pages. Electrical circuits can be represented by graphs whose cycles encode independent current laws. In operations research, matching, coloring, scheduling, and flow turn practical decisions into combinatorial algorithms.

\paragraph{What the theory depends on and where it is useful}
Graph theory depends on sets, counting, algorithms, linear algebra, probability, and topology. It supports extremal and Ramsey theory by expressing forbidden patterns, supports control and network science through connectivity, and supports algebra through graph symmetries and group actions.

Its abstraction is a strength but also a limit: a graph may omit geography, time, uncertainty, capacity, or the strength of a relationship. Weighted, temporal, probabilistic, hypergraph, and multilayer models add the missing information.

\section*{4. Important theorems and results}
\begin{enumerate}[leftmargin=*,itemsep=0.35em]
\item \textbf{Handshaking lemma.} The sum of all vertex degrees equals twice the number of edges.

\item \textbf{Euler criterion.} A connected graph has an Euler circuit exactly when every vertex has even degree.

\item \textbf{Hall's theorem.} A bipartite graph has a matching covering the left part exactly when every subset of the left part has at least as many neighbors.

\item \textbf{Max-flow min-cut theorem.} The greatest flow from source to sink equals the capacity of a minimum cut.

\item \textbf{Euler's planar formula and four-color theorem.} A connected planar graph satisfies $V-E+F=2$, and every planar graph has chromatic number at most four \cite{graph_theory_ref1}.
\end{enumerate}

\theoryapplications
\theoryconnections{graph theory}{%
\item Set theory supplies vertices and edges, combinatorics supplies counting, and linear algebra supplies adjacency matrices and spectral information.
\item Algorithms turn graph definitions into practical searches, while probability studies random graphs and uncertain networks.
}{%
\item Graph theory helps routing, scheduling, social and biological networks, circuit design, and web search.
\item It converts a system of pairwise relationships into paths, cuts, matchings, and flows that can guide a concrete design or diagnosis.
}
\section*{7. Sample Questions}
\emph{Exercise reference:} \cite{graph_theory_ref1}. The cited edition is the source for the exercise set; consult its Exercises section for the official statement and numbering.
\begin{enumerate}[leftmargin=*,itemsep=0.9em]
\item \textbf{Exercise 1.} How many edges does a tree with $20$ vertices have?

\emph{Solution.} Every tree with $n$ vertices has $n-1$ edges, so it has $19$ edges.

\item \textbf{Exercise 2.} Can a connected graph with exactly two odd-degree vertices have an Euler circuit?

\emph{Solution.} No. An Euler circuit requires every vertex to have even degree. Such a graph may have an open Euler trail beginning at one odd vertex and ending at the other.

\item \textbf{Exercise 3.} What does a minimum spanning tree optimize?

\emph{Solution.} It connects every vertex without cycles while minimizing the total weight of selected edges. It does not necessarily minimize the distance between every pair.

\item \textbf{Exercise 4.} When does a bipartite graph have a matching covering every worker?

\emph{Solution.} Hall's condition must hold: every set of workers must collectively be adjacent to at least as many tasks as there are workers in that set.

\item \textbf{Exercise 5.} Why can an adjacency list be better than a matrix?

\emph{Solution.} A sparse graph has few edges compared with all possible pairs. A list stores only existing neighbors, while a matrix stores mostly zeros.

\end{enumerate}
\section*{8. References}
\addcontentsline{toc}{section}{References}

\begin{thebibliography}{9}
\bibitem{graph_theory_ref1} Reinhard Diestel, \emph{Graph Theory}, 5th ed., Springer, 2017.
\bibitem{graph_theory_ref2} Douglas B. West, \emph{Introduction to Graph Theory}, 2nd ed., Prentice Hall, 2001.
\bibitem{graph_theory_ref3} Béla Bollobás, \emph{Modern Graph Theory}, Springer, 1998.
\end{thebibliography}


\theory{Group theory}{What can be learned from studying reversible operations and the symmetries they form?}{A group is a collection of reversible actions that can be composed, with an identity action and an inverse for every action. Group theory studies the common structure behind rotations, permutations, arithmetic symmetries, and transformations.}{Geometry, algebra, number theory, physics, chemistry, cryptography, coding, and the classification of finite structures.}{Composition and inverses encode symmetry operations, while subgroups and representations expose the internal structure of those symmetries.}

\paragraph{The idea of a group}
Imagine rotating a square around its center. There are four rotations that return the square to itself: turn by $0°$, $90°$, $180°$, or $270°$. Doing one rotation after another gives another rotation (for example, $90°$ then $90°$ gives $180°$). Doing nothing is an action too (the identity). Every rotation can be undone: to undo a $90°$ turn, rotate by $270°$. These four rules---combining actions, an identity, and an inverse for each action---define a group. The elements need not be numbers; they may be rotations, shuffles, reflections, or any reversible operation.

Formally, a group has a binary operation (a way to combine two actions into one), an identity element (doing nothing), and an inverse for every element (a way to undo). A group is abelian if the order of composition does not matter: rotating $90°$ then $180°$ gives the same result as $180°$ then $90°$. A group that is not abelian---like the symmetries of a square including reflections---is called non-abelian \cite{group_theory_ref1}.

\end{historylist}
\begin{quote}\small \textbf{Prerequisite:} High-school algebra. The first two sections need no proof techniques.\end{quote}
\section*{3. Theory}
\subsection*{Core theory}
\paragraph{Subgroups, cosets, and quotients}
A subgroup is a smaller collection closed under the same operation. The even integers form a subgroup of the integers under addition. A subgroup partitions a group into equal-sized left or right cosets. Lagrange's theorem follows: the order of a finite subgroup divides the order of the whole group.

If a subgroup is normal, its cosets can themselves be multiplied to form a quotient group. Normal subgroups are precisely the subgroups compatible with conjugation by every group element. A homomorphism is a structure-preserving map between groups; its kernel is normal, and the first isomorphism theorem says that the quotient by the kernel is isomorphic to the image.

\paragraph{How groups solve symmetry problems}
The symmetries of an object act on the object, producing a group action. The orbit of a point is the collection of positions reachable by the group, while the stabilizer consists of actions that leave it fixed. Orbit-stabilizer says, for a finite group,
\[
|G|=|\operatorname{Orb}(x)|\,|\operatorname{Stab}(x)|.
\]
This counts symmetries by counting where an object can go and how many symmetries keep it there.

Permutation groups rearrange a finite set. The symmetric group $S_n$ contains all $n!$ permutations; the alternating group $A_n$ contains the even permutations. Cayley's theorem embeds every group into a permutation group, showing that permutations are a universal model for abstract groups.

\paragraph{Finite groups and their classification}
Finite groups may be described by a multiplication table, generators and relations, or a presentation. Simple groups have no nontrivial normal subgroup. They are the indivisible building blocks of finite groups, in the same way that prime numbers are building blocks of integers. The classification of finite simple groups is a vast theorem showing that every finite simple group belongs to one of several families: cyclic groups of prime order, alternating groups, groups of Lie type, or one of twenty-six sporadic groups. The proof occupies thousands of pages \cite{group_theory_ref2}.

The Sylow theorems constrain subgroups whose order is a power of a prime. They often prove that a group of a proposed order must be abelian or must contain a normal subgroup. Semidirect products explain how one group can act on another to build a larger group.

\paragraph{Representations and characters}
A representation realizes abstract group elements as matrices acting on a vector space. Matrix multiplication then performs the group operation, so linear algebra becomes available. A representation is faithful when different group elements give different matrices. Characters record traces of representation matrices; they are constant on conjugacy classes and can determine how representations decompose.

Fourier polynomials are characters of the circle group $U(1)$ acting on periodic functions. In quantum mechanics, a state space carries a representation of a symmetry group, and conserved quantities arise from symmetry through the broader principle associated with Noether. Finite-group representations are also used to analyze molecules and crystal symmetries \cite{group_theory_ref3}.

\paragraph{Lie groups and continuous symmetry}
A Lie group is both a group and a differentiable manifold, with smooth multiplication and inversion. Rotations in three-dimensional space form the Lie group $SO(3)$. Its tangent space at the identity is a Lie algebra, whose bracket measures the first-order failure of rotations to commute.

Lie groups are the natural framework for continuous symmetries of differential equations, geometry, and physics. Differential Galois theory studies continuous symmetries of differential equations in a role analogous to permutation groups for algebraic equations. Compact Lie groups and their representations describe rotations, angular momentum, and gauge symmetries.

\paragraph{Combinatorial and geometric group theory}
Combinatorial group theory studies presentations, words in generators, and the word problem: can an algorithm decide whether a given word represents the identity? Free groups have no relations beyond cancellation. Adding relations produces quotient groups and often a geometric space.

Geometric group theory studies a group through a space on which it acts. A Cayley graph records generators as labeled edges, turning algebraic multiplication into paths. Hyperbolic groups imitate negative-curvature geometry, while Bass--Serre theory studies groups acting on trees. These ideas connect group theory to topology and geometry.

\paragraph{Applications}
Elliptic curves over finite fields have groups of points. Their discrete logarithm problem supports public-key cryptography. Diffie--Hellman key exchange uses a finite cyclic group; braid-group and other non-abelian proposals use harder group operations, though security must be analyzed carefully. Error-correcting codes use group structure to organize words and decode them. Caesar's cipher is a simple group operation on alphabet positions \cite{group_theory_ref4}.

Symmetry groups classify crystal patterns, molecules, tilings, and physical laws. In topology, fundamental groups record loops and coverings. In number theory, Galois groups encode symmetries of field extensions. In algebraic geometry, algebraic groups and group schemes add a compatible group operation to a geometric object.

\paragraph{What the theory depends on and where it is useful}
Group theory depends on set operations, functions, relations, and increasingly on linear algebra, topology, and geometry. It helps Galois theory organize root symmetries, representation theory convert algebra to matrices, Lie theory describe continuous motion, topology classify loops, and cryptography construct one-way computational tasks.

The classification of a group may be difficult even when its definition is short. A group model also records exact symmetry, not measurement noise or approximate symmetry; applications often require additional geometry, probability, or computation.

\section*{4. Important theorems and results}
\begin{enumerate}[leftmargin=*,itemsep=0.35em]
\item \textbf{Lagrange's theorem.} The order of a subgroup of a finite group divides the order of the group. \emph{Example.} $S_3$ has order $6$; its subgroups have orders $1,2,3$, and $6$, all divisors of $6$.

\item \textbf{Cayley's theorem.} Every group is isomorphic to a subgroup of a symmetric permutation group. \emph{Example.} The cyclic group $\mathbb{Z}/3\mathbb{Z}$ embeds in $S_3$ by sending the generator $1$ to the 3-cycle $(1\;2\;3)$.

\item \textbf{Orbit-stabilizer theorem.} For a finite group action, $|G|=|\operatorname{Orb}(x)|\,|\operatorname{Stab}(x)|$. \emph{Example.} $S_3$ acts on $\{1,2,3\}$. The orbit of $1$ is $\{1,2,3\}$ (size 3). The stabilizer of $1$ is $\{\text{id},(2\;3)\}$ (size 2). Then $|S_3|=3\cdot 2=6$.

\item \textbf{First isomorphism theorem.} A quotient by the kernel of a homomorphism is isomorphic to its image. \emph{Example.} The map $\varphi:\mathbb{Z}\to\mathbb{Z}/6\mathbb{Z}$ has kernel $6\mathbb{Z}$, so $\mathbb{Z}/6\mathbb{Z}\cong\mathbb{Z}/6\mathbb{Z}$, confirming the theorem.

\item \textbf{Classification of finite simple groups.} Every finite simple group belongs to the cyclic, alternating, Lie-type, or sporadic families \cite{group_theory_ref2}.

These results move from local structure to classification. Lagrange and orbit--stabilizer give numerical tests for a finite action, Cayley realizes every group as permutations, and the isomorphism theorem explains how kernels remove information. The classification theorem is a map of the possible indivisible finite symmetries; it is not a short method for identifying an arbitrary group.
\end{enumerate}

\theoryapplications
\theoryconnections{group theory}{%
\item Set theory supplies the underlying collection and operation, while algebra studies homomorphisms, subgroups, quotients, and actions.
\item Geometry and number theory provide important examples, and representation theory turns abstract group elements into matrices.
}{%
\item Group theory helps describe symmetry in geometry, chemistry, physics, coding, and algorithms.
\item It also organizes transformations into reusable algebraic objects, so a proof about one group action can apply to many different systems.
}
\section*{Glossary}
\begin{description}[leftmargin=*,style=nextline,itemsep=0.3em]
\item[\textbf{Group.}] A set with an associative binary operation, an identity element, and an inverse for every element; it formalizes the idea of symmetry or reversible transformations.
\item[\textbf{Subgroup.}] A subset of a group that is itself a group under the same operation.
\item[\textbf{Normal subgroup.}] A subgroup invariant under conjugation by every group element, which allows the formation of a quotient group.
\item[\textbf{Homomorphism.}] A structure-preserving map between two groups; its kernel is always a normal subgroup of the domain.
\item[\textbf{Orbit.}] The set of positions reachable from a given element by applying all elements of an acting group.
\item[\textbf{Simple group.}] A group with no nontrivial normal subgroup; simple groups are the building blocks from which all finite groups are assembled.
\item[\textbf{Representation.}] A realization of an abstract group as matrices acting on a vector space, making linear-algebraic tools available for study.
\item[\textbf{Coset.}] A set $gH=\{gh:h\in H\}$ for a subgroup $H$ and element $g$; cosets partition the group into equal-sized pieces.
\item[\textbf{Quotient group.}] The group $G/N$ formed by cosets of a normal subgroup $N$.
\item[\textbf{Kernel.}] The set of elements mapped to the identity by a homomorphism; always a normal subgroup.
\item[\textbf{Group action.}] A homomorphism from a group to the permutation group of a set.
\item[\textbf{Abelian group.}] A group in which the operation is commutative.
\end{description}
\section*{7. Sample Questions}
\emph{Exercise reference:} \cite{group_theory_ref1}. The cited edition is the source for the exercise set; consult its Exercises section for the official statement and numbering.
\begin{enumerate}[leftmargin=*,itemsep=0.9em]
\item \textbf{Exercise 1.} Find all subgroups of $\mathbb{Z}/12\mathbb{Z}$ and verify Lagrange's theorem for each.

\emph{Solution.} By the correspondence theorem, subgroups of $\mathbb{Z}/12\mathbb{Z}$ are $\langle d\rangle$ for $d\mid 12$: $\langle 0\rangle=\{0\}$ (order 1), $\langle 1\rangle=\mathbb{Z}/12\mathbb{Z}$ (order 12), $\langle 2\rangle=\{0,2,4,6,8,10\}$ (order 6), $\langle 3\rangle=\{0,3,6,9\}$ (order 4), $\langle 4\rangle=\{0,4,8\}$ (order 3), $\langle 6\rangle=\{0,6\}$ (order 2). Each subgroup order divides 12, confirming Lagrange's theorem.

\item \textbf{Exercise 2.} Compute the order of the element $(1\;2\;3)$ in $S_4$.

\emph{Solution.} The permutation $(1\;2\;3)$ is a 3-cycle. Its order is $3$, since $(1\;2\;3)^3=\text{id}$ and no smaller positive power is the identity: $(1\;2\;3)^1\neq\text{id}$, $(1\;2\;3)^2=(1\;3\;2)\neq\text{id}$.

\item \textbf{Exercise 3.} Find the kernel of the homomorphism $\varphi:\mathbb{Z}\to\mathbb{Z}/6\mathbb{Z}$ given by $\varphi(n)=n\bmod 6$.

\emph{Solution.} $\ker\varphi=\{n\in\mathbb{Z}:n\equiv 0\pmod{6}\}=6\mathbb{Z}$. By the first isomorphism theorem, $\mathbb{Z}/\ker\varphi=\mathbb{Z}/6\mathbb{Z}\cong\operatorname{im}\varphi=\mathbb{Z}/6\mathbb{Z}$, which checks out.

\item \textbf{Exercise 4.} Prove that a group of order $7$ is cyclic.

\emph{Solution.} By Lagrange's theorem, every non-identity element has order dividing $7$. Since $7$ is prime, the only divisors are $1$ and $7$. Any non-identity element has order $7$, hence generates the entire group. Therefore the group is cyclic: $G\cong\mathbb{Z}/7\mathbb{Z}$.

\item \textbf{Exercise 5.} Determine whether the subgroup $\{(1),(1\;2)(3\;4),(1\;3)(2\;4),(1\;4)(2\;3)\}$ of $S_4$ is normal.

\emph{Solution.} This subgroup $V\cong\mathbb{Z}/2\times\mathbb{Z}/2$ (the Klein four-group) is normal in $S_4$. To see why, observe that $V$ consists of the identity and all products of two disjoint transpositions. Conjugation by any $\sigma\in S_4$ sends a product of two disjoint transpositions to another product of two disjoint transpositions (it simply relabels the elements), so $\sigma V\sigma^{-1}\subseteq V$. Since $|V|$ is fixed and conjugation is a bijection, $\sigma V\sigma^{-1}=V$ for every $\sigma\in S_4$.

\end{enumerate}
\section*{8. References}
\addcontentsline{toc}{section}{References}

\begin{thebibliography}{9}
\bibitem{group_theory_ref1} Joseph J. Rotman, \emph{An Introduction to the Theory of Groups}, 4th ed., Springer, 1995.
\bibitem{group_theory_ref2} Daniel Gorenstein, Richard Lyons, and Ronald Solomon, \emph{The Classification of the Finite Simple Groups}, American Mathematical Society, 1994--.
\bibitem{group_theory_ref3} Jean-Pierre Serre, \emph{Linear Representations of Finite Groups}, Springer, 1977.
\bibitem{group_theory_ref4} David S. Dummit and Richard M. Foote, \emph{Abstract Algebra}, 3rd ed., Wiley, 2004.
\end{thebibliography}


\theory{Hodge theory}{How can differential equations choose a canonical representative for a topological hole?}{Hodge theory uses a metric and the Laplacian to turn a cohomology class into a unique harmonic differential form. It links topology, geometry, analysis, and algebraic geometry.}{Riemannian geometry, complex projective geometry, algebraic cycles, number theory, and elliptic partial differential equations.}{The metric Laplacian selects harmonic representatives, linking differential equations to cohomology and geometric information.}

\paragraph{History and motivation}
W. V. D. Hodge developed the theory in the 1930s while studying algebraic geometry. It built on Georges de Rham's discovery that integration of differential forms can calculate topological cohomology. Hodge wanted a higher-dimensional version of the relation between real and imaginary parts of holomorphic differentials on a Riemann surface. His proof was repaired by Hermann Weyl and Kunihiko Kodaira \cite{hodge_theory_ref1,hodge_theory_ref2}.

Topology tells us that a space may have holes, but a cohomology class usually has many differential forms representing it. Hodge theory chooses one special representative by measuring size. This is the same idea as choosing the shortest vector among many vectors that describe the same displacement.

\end{historylist}
\section*{3. Theory}
\subsection*{Core theory}
\paragraph{De Rham cohomology}
On a smooth manifold $M$, let $\Omega^k(M)$ be the smooth differential $k$-forms. Exterior differentiation gives
\[
0\longrightarrow\Omega^0(M)\xrightarrow d\Omega^1(M)\xrightarrow d\cdots\xrightarrow d\Omega^n(M)\longrightarrow0,
\]
and $d^2=0$. The $k$th de Rham cohomology is
\[
H^k_{\mathrm{dR}}(M)=\ker(d:\Omega^k\to\Omega^{k+1})/
\operatorname{im}(d:\Omega^{k-1}\to\Omega^k).
\]
Closed forms are those with $d\omega=0$; exact forms are derivatives of other forms. Exact forms are treated as zero because their integrals around closed cycles vanish by Stokes' theorem.

De Rham's theorem identifies this analytic construction with singular cohomology with real coefficients:
\[
H^k_{\mathrm{dR}}(M)\cong H^k_{\mathrm{sing}}(M;\mathbb R).
\]
Thus differential equations can detect the same holes that a topologist sees \cite{hodge_theory_ref3}.

\paragraph{Concrete illustration: Hodge theory on a torus}
The simplest closed Riemann surface is the flat torus $T^2=\mathbb R^2/\mathbb Z^2$ of genus 1, equipped with the metric $ds^2=dx^2+dy^2$ where $x,y$ are coordinates modulo 1. The de Rham cohomology groups are $H^0\cong\mathbb R$ (constants), $H^1\cong\mathbb R^2$ (spanned by $[dx]$ and $[dy]$), and $H^2\cong\mathbb R$ (spanned by $[dx\wedge dy]$).

The harmonic representatives are precisely the constant functions (degree 0), the constant-coefficient 1-forms $\alpha\,dx+\beta\,dy$ (degree 1), and $\gamma\,dx\wedge dy$ (degree 2). Every closed 1-form $\omega=f\,dx+g\,dy$ with $df=\partial f/\partial y\,dy=0$ and $dg=\partial g/\partial x\,dx=0$ (i.e., $\partial f/\partial y=\partial g/\partial x$) can be decomposed as $\omega=dh+\eta$ where $h$ is a function and $\eta=\bar f\,dx+\bar g\,dy$ with constant coefficients $\bar f=\int_0^1\int_0^1 f\,dx\,dy$ and $\bar g=\int_0^1\int_0^1 g\,dx\,dy$. The harmonic part $\eta$ is the unique representative of the cohomology class $[\omega]\in H^1_{\mathrm{dR}}(T^2)$.

Concretely, take $\omega=(2\pi x)\,dy$, which is closed ($d\omega=0$) but not exact. Its harmonic representative is the constant form $\eta=dy$ (since $\int_0^1\int_0^1 2\pi x\,dx\,dy=1$ for the $dy$-coefficient and the $dx$-coefficient is 0). One verifies directly: $\omega-d(2\pi xy)=2\pi x\,dy-2\pi y\,dx-2\pi x\,dy+\ldots$---more cleanly, $\omega=\eta+d(2\pi xy-2\pi xy)$ works after subtracting a suitable exact form. This illustrates how the flat metric on $T^2$ selects constant-coefficient forms as the canonical representatives, generalizing the familiar fact that on $S^1$ the harmonic 1-form is a constant multiple of $d\theta$.

The Hodge decomposition of the Hodge numbers is $h^{0,0}=h^{1,0}=h^{0,1}=h^{1,1}=1$, consistent with $H^1(T^2;\mathbb C)=H^{1,0}\oplus H^{0,1}\cong\mathbb C$ and the Betti numbers $b_0=1,\,b_1=2,\,b_2=1$, all equal as expected for a K\"ahler manifold.

\paragraph{The metric, adjoint, and Laplacian}
Choose a Riemannian metric. It gives an inner product on forms, and therefore an adjoint $\delta$ to $d$. The Hodge Laplacian is
\[
\Delta=d\delta+\delta d.
\]
A form is harmonic when $\Delta\omega=0$. On a closed Riemannian manifold, a harmonic form is both closed and coclosed. Hodge's theorem says that every real cohomology class has exactly one harmonic representative. Harmonic forms are kernels of elliptic differential operators, so the cohomology groups are finite-dimensional.

The Hodge decomposition writes every form uniquely as
\[
\omega=d\alpha+\delta\beta+\gamma,
\]
where $\gamma$ is harmonic. The three pieces are orthogonal: an exact part, a coexact part, and the part that carries the cohomology. This generalizes the Helmholtz decomposition of a vector field into gradient, curl, and harmonic components.

\paragraph{Elliptic complexes}
The de Rham complex is one example of an elliptic complex, a chain of differential operators between vector bundles whose combined Laplacian is elliptic. Atiyah and Bott extended Hodge's theorem to this setting. The associated Green operator gives
\[
\operatorname{Id}=H+\Delta G=H+G\Delta,
\]
where $H$ projects onto harmonic sections. The cohomology of the complex is canonically the space of harmonic sections. This framework includes Dolbeault, Dirac, and other complexes \cite{hodge_theory_ref2}.

\paragraph{Complex projective and Kahler manifolds}
On a complex manifold, a differential form splits into types $(p,q)$, counting holomorphic and antiholomorphic differentials. A compact Kahler manifold has
\[
H^r(X;\mathbb C)=\bigoplus_{p+q=r}H^{p,q}(X).
\]
The decomposition is independent of the chosen Kahler metric but depends on the complex structure, whereas ordinary cohomology depends only on the underlying topological space. The wedge product respects type:
\[
H^{p,q}(X)\smile H^{p',q'}(X)\subseteq H^{p+p',q+q'}(X).
\]
For a smooth projective variety, the Hodge pieces agree with coherent sheaf cohomology $H^q(X,\Omega^p)$ \cite{hodge_theory_ref4}.

\paragraph{Algebraic cycles and the Hodge conjecture}
An algebraic cycle is a formal combination of subvarieties. Its cohomology class has a Hodge type constrained by its complex dimension. The Hodge conjecture proposes that certain rational cohomology classes of type $(p,p)$ arise from algebraic cycles. It is known in special cases but remains open in general.

This explains why Hodge theory matters in algebraic geometry: it turns geometric subvarieties into cohomology classes, then asks whether the converse geometric realization is possible. Poincare duality relates a cycle's homology class to its cohomology class, and integration pairs these classes with differential forms.

\paragraph{Physics and arithmetic}
Maxwell's equations use differential forms and Laplacian-type operators, so harmonic forms appear naturally in mathematical physics. Hodge theory is also used in gauge theory, Yang--Mills equations, and the study of moduli spaces. Although classical Hodge theory is built over the real and complex numbers, p-adic Hodge theory provides analogous tools for arithmetic geometry and number theory \cite{hodge_theory_ref5}.

\paragraph{What the theory depends on and where it is useful}
Hodge theory depends on differential forms, cohomology, Riemannian metrics, functional analysis, elliptic PDE, and complex geometry. It helps topology by giving computable representatives, helps algebraic geometry by splitting cohomology into types, and helps number theory through p-adic analogues.

The canonical representative depends on a metric even though the cohomology class does not. The Hodge decomposition is not available for every noncompact manifold or arbitrary complex manifold; compactness, smoothness, and ellipticity supply the required control.

\section*{4. Important theorems and results}
\begin{enumerate}[leftmargin=*,itemsep=0.35em]
\item \textbf{de Rham theorem.} Smooth differential forms compute singular cohomology with real coefficients.

\item \textbf{Hodge theorem.} Each real cohomology class on a closed Riemannian manifold has a unique harmonic representative \cite{hodge_theory_ref1}.

\item \textbf{Hodge decomposition.} Every differential form splits orthogonally into exact, coexact, and harmonic parts.

\item \textbf{Kahler decomposition.} On a compact Kahler manifold, complex cohomology splits as the direct sum of the $H^{p,q}$ pieces.

\item \textbf{Hodge index theorem.} The intersection form on a compact Kahler surface has a signature constrained by its Hodge decomposition.
\end{enumerate}

\theoryapplications
\theoryconnections{hodge theory}{%
\item Differential geometry supplies metrics, differential forms, and the Laplacian; topology supplies cohomology; and complex geometry supplies the type decomposition of forms.
\item Analysis provides elliptic regularity, which turns an analytic equation into finite-dimensional geometric information.
}{%
\item Hodge theory helps classify compact manifolds, study algebraic cycles, and compare topology with complex or Kahler geometry.
\item Harmonic representatives make cohomology classes concrete, while the decomposition constrains which topological patterns can occur.
}
\noindent\textbf{Maxwell's equations and the Hodge star in physics.} Maxwell's equations in vacuum on a Riemannian 3-manifold (or Minkowski 4-manifold) are naturally expressed using the Hodge star operator $*\colon\Omega^k\to\Omega^{n-k}$. The electric and magnetic fields are encoded in the Faraday 2-form $F=B_1\,dy\wedge dz+B_2\,dz\wedge dx+B_3\,dx\wedge dy+E_1\,dx\wedge dt+E_2\,dy\wedge dt+E_3\,dz\wedge dt$. Maxwell's equations then read $dF=0$ and $d{*F}=J$ (the current 3-form). A static magnetic field with no currents satisfies $dF=0$ and $d{*F}=0$, meaning $F$ is a harmonic 2-form with respect to the Hodge Laplacian $\Delta=d\delta+\delta d$. On $\mathbb R^3$ with the flat metric, harmonic 2-forms correspond to divergence-free, curl-free vector fields, exactly the fields satisfying both the magnetostatic equations. The Hodge decomposition therefore provides a rigorous mathematical foundation for the Helmholtz decomposition of electromagnetic fields: any electromagnetic field configuration decomposes uniquely into a radiation part (harmonic), an electrostatic/magnetostatic part (exact and coexact), and a gauge component. In engineering, this decomposition underlies boundary-element methods for solving Maxwell's equations in antenna design and microwave cavity analysis, where the harmonic part captures the propagating modes and the exact/coexact parts capture near-field effects that decay rapidly.
\section*{7. Sample Questions}
\emph{Exercise reference:} \cite{hodge_theory_ref1}. The cited edition is the source for the exercise set; consult its Exercises section for the official statement and numbering.
\begin{enumerate}[leftmargin=*,itemsep=0.9em]
\item \textbf{Exercise 1.} Why does $d^2=0$ matter?

\emph{Solution.} It says every exact form is closed, so the quotient of closed forms by exact forms is well-defined. Without this cancellation, there would be no cohomology group.

\item \textbf{Exercise 2.} What does harmonic mean for a differential form?

\emph{Solution.} It means the Hodge Laplacian satisfies $\Delta\omega=0$. On a closed manifold this implies that the form is both closed and coclosed.

\item \textbf{Exercise 3.} Why is the harmonic representative called canonical?

\emph{Solution.} Once a metric is chosen, each cohomology class has exactly one harmonic form representing it. There is no further choice among representatives.

\item \textbf{Exercise 4.} What is the purpose of the decomposition $\omega=d\alpha+\delta\beta+\gamma$?

\emph{Solution.} It separates a form into a derivative part, a codifferential part, and a harmonic part. Only the harmonic part survives as the essential cohomology information.

\item \textbf{Exercise 5.} Why does Hodge theory not apply identically to every complex manifold?

\emph{Solution.} The $(p,q)$ splitting and its good behavior require strong compatibility between the metric and complex structure, as in the Kahler condition. General complex manifolds need not have this decomposition.

\end{enumerate}
\section*{8. References}
\addcontentsline{toc}{section}{References}

\begin{thebibliography}{9}
\bibitem{hodge_theory_ref1} W. V. D. Hodge, \emph{The Theory and Applications of Harmonic Integrals}, Cambridge University Press, 1941.
\bibitem{hodge_theory_ref2} Werner Ballmann, \emph{Lectures on Kahler Manifolds}, European Mathematical Society, 2006.
\bibitem{hodge_theory_ref3} Georges de Rham, \emph{Differentiable Manifolds}, Springer, 1984.
\bibitem{hodge_theory_ref4} Claire Voisin, \emph{Hodge Theory and Complex Algebraic Geometry I}, Cambridge University Press, 2002.
\bibitem{hodge_theory_ref5} James S. Milne, \emph{Lectures on Etale Cohomology}, Princeton University, 1980.
\end{thebibliography}


\theory{Homology theory}{How can a space's holes be turned into algebraic objects that can be calculated?}{Homology replaces geometric pieces by a chain complex. Cycles are closed pieces, boundaries are pieces that bound something, and homology records cycles that are not boundaries.}{Topology, geometry, data analysis, robotics, mesh computation, algebraic geometry, and differential equations.}{Cycles modulo boundaries measure holes in a way that is invariant under continuous deformation and computable from finite complexes.}

\paragraph{History and the central idea}
Homology grew from Euler's polyhedron formula, Riemann's genus, and Betti's numerical invariants. It became a systematic part of algebraic topology when mathematicians learned to associate algebraic groups with spaces. The word is also used more generally for the homology of any chain complex and for the collection of procedures called homology theories \cite{homology_theory_ref1}.

Think of a triangulated surface. A vertex is a zero-dimensional piece, an edge is a one-dimensional piece, and a triangle is a two-dimensional piece. A collection of edges can form a closed loop. If the loop is the boundary of a collection of triangles, it is not a genuine hole. If it cannot be filled, it represents a one-dimensional hole.

\end{historylist}
\section*{3. Theory}
\subsection*{Core theory}
\paragraph{Chains, cycles, and boundaries}
A chain complex is a sequence of abelian groups or modules
\[
\cdots\longrightarrow C_{n+1}\xrightarrow{\partial_{n+1}}C_n
\xrightarrow{\partial_n}C_{n-1}\longrightarrow\cdots
\]
with $\partial_n\partial_{n+1}=0$. The boundary of a boundary is empty. The $n$-cycles are $Z_n=\ker\partial_n$, and the $n$-boundaries are $B_n=\operatorname{im}\partial_{n+1}$. Since every boundary is a cycle, the $n$th homology group is
\[
H_n=\frac{Z_n}{B_n}.
\]
It identifies cycles that differ by a boundary.

For a connected space, $H_0$ is usually $\mathbb Z$, one copy for each connected component. $H_1$ records loops, and higher groups record higher-dimensional voids. The rank of a homology group is a Betti number; finite pieces such as $\mathbb Z_2$ are torsion.

\paragraph{Examples that build intuition}
The sphere has no one-dimensional hole, so $H_1(S^2)=0$, but it has a two-dimensional surface class, so $H_2(S^2)=\mathbb Z$. A circle has one loop and no filled two-dimensional region: $H_0(S^1)=\mathbb Z$, $H_1(S^1)=\mathbb Z$, and higher homology is zero.

The torus has two independent one-dimensional loops, one around the tube and one through the hole, and one two-dimensional surface class. Its Betti numbers are $1,2,1$. The projective plane has $H_1=\mathbb Z_2$: its essential loop becomes a boundary when traversed twice. This is torsion, a phenomenon not visible from ordinary real-valued measurements \cite{homology_theory_ref2}.

\paragraph{Simplicial and singular homology}
In simplicial homology, a space is decomposed into vertices, edges, triangles, and higher simplices. Boundary matrices make the calculation finite and algorithmic. Singular homology instead allows every continuous map from a standard simplex into the space. It works for arbitrary topological spaces and is invariant under homotopy.

Cellular homology uses cells in a CW complex and is often smaller than singular homology while producing the same answer. Cubical homology uses squares and cubes, which is useful for images and voxel data. Relative homology $H_n(X,A)$ studies chains in $X$ modulo chains already contained in a subspace $A$.

\paragraph{Maps, exact sequences, and products}
A continuous map induces homomorphisms on homology. Homotopic maps induce the same homomorphism, so homology ignores continuous deformation. A pair $A\subseteq X$ gives a long exact sequence linking $H_n(A)$, $H_n(X)$, and relative homology. Mayer--Vietoris calculates a space by splitting it into overlapping pieces.

The cross product combines classes in two spaces, and the Kunneth theorem describes the homology of a product. Poincare duality relates homology and cohomology of an oriented closed manifold. The Euler characteristic is the alternating sum $\chi=\sum(-1)^n\operatorname{rank}H_n$ and is independent of the chosen triangulation.

\paragraph{Homology theories}
Different constructions share the same formal behavior: homotopy invariance, exact sequences, excision, and a dimension rule. Singular, simplicial, cellular, de Rham, and Borel--Moore homology are examples. Generalized theories such as K-theory and cobordism homology relax the ordinary dimension rule and reveal finer structure.

In algebraic geometry, sheaf and coherent homology measure the failure of local data to fit globally. In computational topology, persistent homology studies how homology changes as a scale parameter grows, producing a barcode that identifies robust features in noisy data \cite{homology_theory_ref3}.

\paragraph{Computation and applications}
For a finite mesh, boundary operators are matrices. Row reduction computes ranks and null spaces, while specialized algorithms preserve sparsity. Finite-element software uses homology and cohomology bases to represent fields on meshes. In data analysis, point clouds sampled from a shape are connected at a chosen distance; homology detects components, loops, and voids as the distance changes.

Homology does not reconstruct every detail of a space. Two spaces can have identical homology groups but different fundamental groups or geometry. It is a robust summary, not a complete photograph.

\paragraph{What the theory depends on and where it is useful}
Homology depends on linear algebra, abelian groups, chain complexes, and topology. It helps homotopy theory by providing computable invariants, helps geometry through Euler characteristics and duality, and helps data science by converting shape into persistence diagrams.

\section*{4. Important theorems and results}
\begin{enumerate}[leftmargin=*,itemsep=0.35em]
\item \textbf{Homotopy invariance.} Homotopic maps induce the same map on homology.

\item \textbf{Mayer--Vietoris sequence.} The homology of a union can be computed from the homology of two pieces and their intersection.

\item \textbf{Excision.} Under suitable inclusion conditions, removing a region from a pair does not change relative homology.

\item \textbf{Kunneth theorem.} It describes the homology of a product from the homology of its factors, including possible tensor and torsion terms.

\item \textbf{Poincare duality.} Homology and cohomology groups of an oriented closed manifold are paired in complementary dimensions \cite{homology_theory_ref4}.
\end{enumerate}

\paragraph{Applications}
\begin{enumerate}[leftmargin=*,itemsep=0.35em]
\item \textbf{Topological data analysis and persistent homology.} Persistent homology applies homology to finite point clouds sampled from unknown shapes. The data are embedded in $\mathbb R^d$, and a filtration of simplicial complexes is built by growing balls of radius $\epsilon$ around each point. As $\epsilon$ increases, new connected components appear and merge, new loops form and get filled in, and voids open and close. A persistence diagram or barcode records the birth and death parameter values for each topological feature. Long-lived bars correspond to robust structural features of the underlying space, while short bars are attributed to sampling noise. This pipeline has been applied to reconstruct the shape of proteins from molecular dynamics simulations, to detect loop structures in neuronal connectivity networks, and to identify clustering structure in high-dimensional gene-expression data. The stability theorem for persistence diagrams guarantees that small perturbations of the input point cloud produce only small displacements of the diagram in the bottleneck metric, giving the method robustness that purely algebraic invariants lack \cite{homology_theory_ref3}.
\item \textbf{Sensor networks and coverage analysis.} In a wireless sensor network, each sensor monitors a region modeled as a disk. The union of covered disks forms a topological space, and its homology groups determine whether the network has coverage gaps (one-dimensional holes) or redundant overlaps (higher-dimensional voids). When sensors have unknown locations, pairwise distance measurements construct a \v{C}ech or Rips complex whose homology is computed from boundary matrices. A one-dimensional homology class that persists across a range of distance scales signals a genuine hole in coverage that must be plugged, while the appearance and disappearance of classes at a single scale indicates a transient measurement error. This homology-based approach has been implemented in marine sensor arrays for ocean temperature monitoring, where persistent $H_1$ classes identify regions with insufficient acoustic coverage, and in agricultural IoT networks, where $H_2$ classes detect three-dimensional voids in soil-moisture sensing grids \cite{homology_theory_ref3,homology_theory_ref4}.
\end{enumerate}
\theoryconnections{homology theory}{%
\item Topology supplies spaces and continuous maps, while linear algebra turns chains into boundary matrices and homology groups into kernels modulo images.
\item Category theory explains functoriality, and algebraic geometry and differential geometry provide rich examples.
}{%
\item Homology helps distinguish spaces, compute holes, and prove fixed-point and duality results.
\item It also supplies computable invariants for geometry and topology, making qualitative features accessible through finite calculations in many examples.
}
\section*{Glossary}
\begin{description}[leftmargin=*,style=nextline,itemsep=0.3em]
\item[\textbf{Chain complex.}] A sequence of abelian groups connected by boundary maps whose composition is zero, encoding the combinatorial skeleton of a space.
\item[\textbf{Cycle.}] An element of a chain complex that maps to zero under the boundary operator; it represents a closed piece with no boundary.
\item[\textbf{Boundary.}] An element that is the image of some higher-dimensional chain under the boundary operator; every boundary is automatically a cycle.
\item[\textbf{Homology group.}] The quotient group of cycles modulo boundaries at a given dimension, measuring the extent to which cycles fail to be boundaries.
\item[\textbf{Betti number.}] The rank (free part) of a homology group; it counts independent holes of a given dimension.
\item[\textbf{Torsion.}] A finite cyclic piece of a homology group, invisible to real-valued measurements, indicating a hole that vanishes after being traversed multiple times.
\item[\textbf{Persistent homology.}] A method that tracks how homology groups change as a scale parameter varies, identifying robust topological features in noisy data.
\item[\textbf{Boundary operator.}] The map $\partial_n\colon C_n\to C_{n-1}$ in a chain complex satisfying $\partial^2=0$.
\item[\textbf{Euler characteristic.}] The alternating sum $\chi=\sum(-1)^n\operatorname{rank}H_n$; a topological invariant independent of triangulation.
\item[\textbf{Künneth theorem.}] Describes the homology of a product space in terms of the homology of its factors.
\item[\textbf{Simplicial complex.}] A combinatorial structure built from vertices, edges, triangles, and higher simplices obeying face-inclusion rules.
\item[\textbf{Simplex.}] A generalization of a triangle to higher dimensions: the convex hull of affinely independent points.
\item[\textbf{CW complex.}] A topological space built by iteratively attaching cells of increasing dimension.
\item[\textbf{Excision.}] The theorem that removing a suitable closed subset does not change relative homology.
\end{description}
\section*{7. Sample Questions}
\emph{Exercise reference:} \cite{homology_theory_ref1}. The cited edition is the source for the exercise set; consult its Exercises section for the official statement and numbering.
\begin{enumerate}[leftmargin=*,itemsep=0.9em]
\item \textbf{Exercise 1.} Compute the simplicial homology groups $H_0$ and $H_1$ of the space obtained by identifying the endpoints of three disjoint line segments to a single point (the wedge $S^1\vee S^1\vee S^1$).

\emph{Solution.} The space is a wedge of three circles. Its simplicial chain complex in low dimensions has $C_1=\mathbb{Z}^3$ (three edges) and $C_0=\mathbb{Z}$ (one vertex). The boundary $\partial_1=0$ since each edge starts and ends at the same vertex. So $H_1=\ker\partial_1/\operatorname{im}\partial_2=\mathbb{Z}^3/0=\mathbb{Z}^3$ and $H_0=C_0/\operatorname{im}\partial_1=\mathbb{Z}/0=\mathbb{Z}$ (since the space is connected). Betti numbers: $b_0=1$, $b_1=3$.

\item \textbf{Exercise 2.} Compute $H_*(S^1\times S^1)$, the homology of the torus $T^2$.

\emph{Solution.} By the K\"{u}nneth theorem, $H_n(T^2)=\bigoplus_{p+q=n}H_p(S^1)\otimes H_q(S^1)$. Since $H_0(S^1)=\mathbb{Z}$, $H_1(S^1)=\mathbb{Z}$, and $H_k(S^1)=0$ for $k\geq 2$: $H_0(T^2)=\mathbb{Z}$, $H_1(T^2)=\mathbb{Z}\oplus\mathbb{Z}=\mathbb{Z}^2$, $H_2(T^2)=\mathbb{Z}$. The Euler characteristic is $1-2+1=0$.

\item \textbf{Exercise 3.} Let $X$ be a graph (1-complex) with $5$ vertices and $7$ edges, connected. Compute $H_0(X)$ and $H_1(X)$.

\emph{Solution.} Since $X$ is connected, $H_0(X)=\mathbb{Z}$. For $H_1$: by Euler's formula for a graph, $\chi(X)=V-E=5-7=-2$. Since $\chi=b_0-b_1$, we get $b_1=b_0-\chi=1-(-2)=3$, so $H_1(X)=\mathbb{Z}^3$.

\item \textbf{Exercise 4.} Show that the real projective plane $\mathbb{RP}^2$ has $H_1(\mathbb{RP}^2;\mathbb{Z})\cong\mathbb{Z}/2\mathbb{Z}$.

\emph{Solution.} Using the CW structure with one 0-cell, one 1-cell, and one 2-cell, the cellular chain complex is $0\to\mathbb{Z}\xrightarrow{2}\mathbb{Z}\xrightarrow{0}\mathbb{Z}\to 0$ (the attaching map of the 2-cell wraps twice around the 1-cell). So $H_1=\ker(0)/\operatorname{im}(2)=\mathbb{Z}/2\mathbb{Z}$, and $H_0=\mathbb{Z}$, $H_2=0$.

\item \textbf{Exercise 5.} Use the Mayer--Vietoris sequence to compute $H_1(S^1)$, where $S^1$ is covered by two open arcs $U$ and $V$ overlapping in two disjoint arcs.

\emph{Solution.} $U$, $V$, and $U\cap V$ are all homotopy equivalent to a point (or disjoint union of two contractible arcs). So $H_0(U)=H_0(V)=\mathbb{Z}$, $H_0(U\cap V)=\mathbb{Z}^2$ (two components), and $H_1(U)=H_1(V)=H_1(U\cap V)=0$. The Mayer--Vietoris sequence gives $\cdots\to H_1(U)\oplus H_1(V)\to H_1(S^1)\to H_0(U\cap V)\xrightarrow{\phi}H_0(U)\oplus H_0(V)\to H_0(S^1)\to 0$. This becomes $0\to H_1(S^1)\to\mathbb{Z}^2\xrightarrow{\phi}\mathbb{Z}^2\to\mathbb{Z}\to 0$. The map $\phi$ sends $(a,b)\mapsto(a+b)-(a+b)=\ldots$ more precisely, the difference map sends a generator of each component of $U\cap V$ to its image in $U$ minus its image in $V$: $\phi(a,b)=(a+b,a+b)$... Actually, $\phi(a,b)=(a+b,-a-b)$, so $\ker\phi\cong\mathbb{Z}$ and $\operatorname{im}\phi\cong\mathbb{Z}$. Therefore $H_1(S^1)\cong\ker\phi\cong\mathbb{Z}$.

\end{enumerate}
\section*{8. References}
\addcontentsline{toc}{section}{References}

\begin{thebibliography}{9}
\bibitem{homology_theory_ref1} Allen Hatcher, \emph{Algebraic Topology}, Cambridge University Press, 2002.
\bibitem{homology_theory_ref2} James R. Munkres, \emph{Elements of Algebraic Topology}, Addison-Wesley, 1984.
\bibitem{homology_theory_ref3} Herbert Edelsbrunner and John Harer, \emph{Computational Topology}, American Mathematical Society, 2010.
\bibitem{homology_theory_ref4} Edwin H. Spanier, \emph{Algebraic Topology}, Springer, 1966.
\end{thebibliography}


\theory{Homotopy theory}{When should two maps or spaces be considered the same if one can be continuously deformed into the other?}{Homotopy theory studies continuous deformation. It replaces exact equality by a controlled notion of sameness and then builds algebraic invariants, higher-dimensional operations, and abstract models of spaces.}{Algebraic topology, algebraic geometry, category theory, noncommutative geometry, bundles, and generalized cohomology.}{Continuous deformations, loop spaces, and higher homotopy groups record which global features survive changes of shape.}

\paragraph{The central idea}
Two continuous maps $f,g:X\to Y$ are homotopic if there is a continuous family $H:X\times[0,1]\to Y$ with $H(x,0)=f(x)$ and $H(x,1)=g(x)$. A rubber band on a tabletop can move continuously without cutting; homotopy regards the starting and ending placements as equivalent. A circle cannot be continuously shrunk to a point while remaining in the punctured plane, because the missing point creates an obstruction \cite{homotopy_theory_ref1}.

Homotopy theory began in algebraic topology but is now an independent subject. It studies spaces, maps, homotopies between maps, and homotopies between those homotopies. The emphasis is not on coordinates but on what survives deformation.

\end{historylist}
\section*{3. Theory}
\subsection*{Core theory}
\paragraph{Spaces, maps, and fundamental groups}
A path is a map from an interval. A loop starts and ends at the same point. Homotopy classes of loops based at a point form the fundamental group $\pi_1(X)$. Multiplication follows one loop after another; the inverse runs a loop backward. The fundamental group of a circle is $\mathbb Z$, while the sphere has trivial fundamental group.

Higher homotopy groups use maps from spheres: $\pi_n(X)$ records maps $S^n\to X$ up to homotopy. Unlike the fundamental group, higher groups are abelian for $n\geq2$. They detect higher-dimensional twisting, but are usually much harder to compute.

\paragraph{Homotopy equivalence and examples}
Spaces $X$ and $Y$ are homotopy equivalent if maps exist in both directions whose composites are homotopic to the identity. A disk is homotopy equivalent to a point; a cylinder is homotopy equivalent to a circle; a torus is not equivalent to a sphere because their loops and homology differ.

Weak homotopy equivalence requires isomorphisms on all homotopy groups and components. Under suitable hypotheses, such as CW complexes, weak equivalence implies homotopy equivalence. This distinction prevents pathological spaces from misleading the theory.

\paragraph{Loops, suspensions, and adjunction}
The loop space $\Omega X$ consists of based loops in $X$, with multiplication by concatenation. The suspension $\Sigma X$ joins two cones over $X$. Looping and suspending are related by
\[
[\,\Sigma X,Y\,]\cong[\,X,\Omega Y\,],
\]
where brackets denote homotopy classes of maps. Milnor's theorem says that if $X$ has the homotopy type of a CW complex, so does $\Omega X$ \cite{homotopy_theory_ref2}.

The suspension operation creates stable phenomena. Stable homotopy theory studies what becomes constant after repeated suspension. Spectra package sequences of spaces and their structure maps, and generalized cohomology theories are represented by spectra.

\paragraph{Bundles and classifying spaces}
For a topological group $G$, a principal $G$-bundle over $X$ is classified by a map to a classifying space $BG$. Homotopy classes of maps satisfy
\[
[X,BG]\cong\{\text{principal }G\text{-bundles over }X\}/\text{isomorphism}.
\]
The universal bundle over $BG$ pulls back along a map $X\to BG$ to produce every principal bundle. Brown's representability theorem explains why many cohomology theories can be represented by spaces.

Eilenberg--Mac Lane spaces satisfy $[X,K(A,n)]=H^n(X;A)$. Thus ordinary cohomology itself can be expressed as a homotopy classification problem \cite{homotopy_theory_ref3}.

\paragraph{Model categories and simplicial sets}
Model category theory provides three classes of maps: weak equivalences, cofibrations, and fibrations. It allows constructions to be performed up to homotopy while keeping enough algebraic control. Topological spaces and nonnegative chain complexes are model categories.

A simplicial set is a sequence of sets of abstract simplices with face and degeneracy maps. It generalizes a simplicial complex and can serve as a combinatorial model of a space. The singular simplicial set of $X$ has as its $n$-simplices all maps from the standard $n$-simplex to $X$; its associated chain complex computes singular homology. Geometric realization turns a simplicial set into a CW complex.

\paragraph{Other mathematical settings}
Algebraic geometry has $\mathbb A^1$-homotopy theory, where the affine line plays the role of an interval. Category theory uses higher categories and infinity-categories to record objects, maps, homotopies, and higher homotopies. Noncommutative geometry uses KK-theory and noncommutative topology to extend homotopy ideas beyond ordinary spaces.

\paragraph{What the theory depends on and where it is useful}
Homotopy theory depends on topology, algebraic topology, category theory, algebraic structures, and often homology and cohomology. It helps classify fiber bundles, organize generalized cohomology, compare algebraic varieties, and express higher categorical structure.

The theory has a cost: exact homotopy groups can be difficult or impossible to calculate directly, and a homotopy invariant deliberately forgets metric information. It is suited to questions of continuous deformation, not to measuring lengths or angles.

\section*{4. Important theorems and results}
\begin{enumerate}[leftmargin=*,itemsep=0.35em]
\item \textbf{Whitehead theorem.} A weak homotopy equivalence between CW complexes is a homotopy equivalence.

\item \textbf{Seifert--van Kampen theorem.} Under suitable conditions, the fundamental group of a union is obtained from the groups of its pieces and intersection.

\item \textbf{Hurewicz theorem.} Under connectivity hypotheses, the first nonzero homotopy group maps isomorphically to the first nonzero homology group.

\item \textbf{Brown representability theorem.} Suitable cohomology theories are represented by spaces.

\item \textbf{Freudenthal suspension theorem.} In a stable range, suspension gives isomorphisms between successive homotopy groups \cite{homotopy_theory_ref4}.
\end{enumerate}

\theoryapplications
\theoryconnections{homotopy theory}{%
\item Topology supplies continuous maps and deformation, while algebraic topology supplies fundamental groups, homology, and homotopy groups.
\item Category theory and higher algebra organize suspension, loop spaces, spectra, and representability.
}{%
\item Homotopy theory helps classify spaces up to flexible equivalence and provides the language of stable phenomena in algebraic topology.
\item It also supports homological algebra, derived geometry, mathematical physics, and the study of higher categories.
}
\section*{7. Sample Questions}
\emph{Exercise reference:} \cite{homotopy_theory_ref1}. The cited edition is the source for the exercise set; consult its Exercises section for the official statement and numbering.
\begin{enumerate}[leftmargin=*,itemsep=0.9em]
\item \textbf{Exercise 1.} Why is a disk homotopy equivalent to a point?

\emph{Solution.} Collapse every point of the disk along a straight line to its center. The collapse varies continuously and gives a deformation retraction.

\item \textbf{Exercise 2.} What is $\pi_1(S^1)$?

\emph{Solution.} It is $\mathbb Z$. A loop is classified by its integer winding number, positive or negative according to direction.

\item \textbf{Exercise 3.} Why is the punctured plane not simply connected?

\emph{Solution.} A loop winding once around the missing point cannot be continuously contracted without crossing that point. Its winding number is a nonzero obstruction.

\item \textbf{Exercise 4.} What does a classifying space classify?

\emph{Solution.} Homotopy classes of maps into $BG$ correspond to isomorphism classes of principal $G$-bundles over the source space.

\item \textbf{Exercise 5.} Why use simplicial sets?

\emph{Solution.} They encode spaces by combinatorial data while retaining face and degeneracy operations. This makes homotopy constructions accessible to algebra and computation.

\end{enumerate}
\section*{8. References}
\addcontentsline{toc}{section}{References}

\begin{thebibliography}{9}
\bibitem{homotopy_theory_ref1} Allen Hatcher, \emph{Algebraic Topology}, Cambridge University Press, 2002.
\bibitem{homotopy_theory_ref2} George W. Whitehead, \emph{Elements of Homotopy Theory}, Springer, 1978.
\bibitem{homotopy_theory_ref3} J. Peter May, \emph{A Concise Course in Algebraic Topology}, University of Chicago Press, 1999.
\bibitem{homotopy_theory_ref4} Mark Hovey, \emph{Model Categories}, American Mathematical Society, 1999.
\end{thebibliography}


\theory{Ideal theory}{How can arithmetic facts about divisibility be generalized from integers to complicated rings?}{An ideal is a subset of a ring closed under addition and subtraction and absorbing multiplication by any ring element. Ideals make quotient rings possible and allow prime factorization to be studied when elements themselves do not factor uniquely.}{Ring theory, algebraic geometry, algebraic number theory, commutative algebra, and polynomial equations.}{Ideals encode divisibility and polynomial constraints, and quotient and localization operations turn those constraints into arithmetic and geometric information.}

\paragraph{Richard Dedekind}
Dedekind introduced ideals to restore unique factorization at the level of ideal products in algebraic number rings, even when elements themselves factor ambiguously.

\paragraph{Emmy Noether and David Hilbert}
Noether developed structural methods for ideals and ascending chain conditions. Hilbert's basis theorem proved finite generation of ideals in polynomial rings, making ideal theory central to commutative algebra and algebraic geometry \cite{ideal_theory_ref1}.
\end{historylist}
\section*{3. Theory}
\subsection*{Core theory}
\paragraph{The problem and the definition}
In the integers, the multiples of $6$ form a set closed under addition and multiplication by any integer. The multiples of $6$ are written $(6)$. More generally, an ideal $I$ of a ring $R$ is an additive subgroup such that $r x\in I$ whenever $r\in R$ and $x\in I$. In a noncommutative ring one distinguishes left, right, and two-sided ideals; in a commutative ring they coincide \cite{ideal_theory_ref1}.

The zero ideal contains only $0$, and the unit ideal is the whole ring. An ideal generated by one element is principal. The ideal generated by a set consists of all finite ring combinations of those generators. In $\mathbb Z$, every ideal is $(n)$ for a unique nonnegative integer $n$, but this simple behavior fails in many rings.

\paragraph{Quotient rings}
An ideal is exactly the kind of subset that permits a quotient ring $R/I$. Two elements are considered equal in the quotient when their difference lies in $I$. For example,
\[
\mathbb Z/(n)
\]
is arithmetic modulo $n$. The quotient $\mathbb Z/(6)$ identifies numbers with the same remainder after division by $6$.

The analogy with group theory is important: normal subgroups produce quotient groups, while two-sided ideals produce quotient rings. The first isomorphism theorem says that the quotient by the kernel of a ring homomorphism is isomorphic to its image.

\paragraph{Prime and maximal ideals}
In a commutative ring, a proper ideal $\mathfrak p$ is prime if $ab\in\mathfrak p$ implies $a\in\mathfrak p$ or $b\in\mathfrak p$. Equivalently, $R/\mathfrak p$ is an integral domain. A proper ideal $\mathfrak m$ is maximal if no proper ideal lies strictly between it and $R$. Equivalently, $R/\mathfrak m$ is a field. Every maximal ideal is prime.

In $\mathbb Z$, $(p)$ is prime exactly when $p$ is a prime number. In the polynomial ring $k[x,y]$, the ideal $(x)$ is prime because the quotient is $k[y]$, while $(x,y)$ is maximal because the quotient is $k$. The prime ideals of a polynomial ring can be viewed as algebraic geometric locations.

\paragraph{Ideal operations and localization}
For ideals $I$ and $J$, the sum $I+J$ contains all $a+b$ with $a\in I$ and $b\in J$. The product $IJ$ consists of finite sums of products $ab$ with $a\in I$, $b\in J$. The intersection $I\cap J$ contains elements common to both.

Localization makes selected elements invertible. Localizing a ring at the complement of a prime ideal focuses attention near one algebraic point. A local ring has one maximal ideal and is the algebraic analogue of looking through a microscope at a single location. Primary decomposition expresses an ideal as an intersection of primary ideals, revealing the prime components and multiplicities \cite{ideal_theory_ref2}.

\paragraph{Factorization and arithmetic}
Unique factorization can fail in the ring of algebraic integers of a number field. Ideals repair this failure: every nonzero ideal in the ring of integers of a number field factors uniquely into prime ideals. The ideal class group measures which ideals are principal. Its finiteness is a central theorem of algebraic number theory.

Fractional ideals are finitely generated modules inside the quotient field, after allowing denominators. A fractional ideal is invertible when another fractional ideal multiplies with it to give the whole ring. In Dedekind domains every nonzero fractional ideal is invertible, which makes ideal arithmetic especially well behaved.

\paragraph{Geometry of ideals}
An ideal $I\subseteq k[x_1,\ldots,x_n]$ defines its zero set
\[
V(I)=\{a: f(a)=0\text{ for every }f\in I\}.
\]
Thus an ideal is a collection of equations and $V(I)$ is their common solution set. Conversely, polynomials vanishing on a set form an ideal. Hilbert's Nullstellensatz links these operations: over an algebraically closed field, radical ideals correspond to algebraic sets.

The ideal $(x^2+y^2-1)$ describes a circle over the real numbers. The ideal $(x,y)$ describes the origin. Prime ideals correspond to irreducible algebraic pieces, maximal ideals to points in affine algebraic geometry. Gröbner bases give an algorithm for membership, elimination, intersection, and solving polynomial systems \cite{ideal_theory_ref3}.

\paragraph{Modules and noncommutative variants}
An ideal is an $R$-module contained in $R$. Module language makes ideal operations systematic and permits generalizations to submodules, annihilators, and Fitting ideals. In noncommutative rings, left and right ideals may differ, and a two-sided ideal is needed for an ordinary quotient ring. One-sided ideals are important in representation theory and operator algebras.

\paragraph{What the theory depends on and where it is useful}
Ideal theory depends on ring theory, modules, divisibility, and polynomial algebra. It helps dimension theory by organizing prime chains, algebraic geometry by encoding varieties, number theory by restoring unique factorization for ideals, and elimination theory by representing polynomial constraints.

The word "ideal" is unrelated to an ideal in order theory, though both use closure ideas. Ideal arithmetic may be more abstract than element arithmetic, but that abstraction is precisely what preserves factorization and geometry.

\section*{4. Important theorems and results}
\begin{enumerate}[leftmargin=*,itemsep=0.35em]
\item \textbf{Correspondence theorem.} Ideals of $R/I$ correspond to ideals of $R$ containing $I$.

\item \textbf{Chinese remainder theorem.} If $I$ and $J$ are comaximal, then $R/(I\cap J)\cong R/I\times R/J$.

\item \textbf{Hilbert basis theorem.} If $R$ is Noetherian, then $R[x]$ is Noetherian; in particular every ideal in a polynomial ring over a field is finitely generated.

\item \textbf{Nullstellensatz.} Radical ideals over an algebraically closed field correspond to algebraic sets.

\item \textbf{Unique ideal factorization.} Every nonzero ideal in a Dedekind domain factors uniquely into prime ideals \cite{ideal_theory_ref4}.

These theorems explain why ideals replace elements when unique factorization of elements fails. The correspondence theorem transfers quotient questions between $R$ and $R/I$; the Chinese remainder theorem separates independent congruence conditions; Noetherianity guarantees finite descriptions; and the Nullstellensatz connects radical ideals with geometric solution sets. Unique ideal factorization restores a controlled form of prime decomposition in Dedekind domains.
\end{enumerate}

\theoryapplications
\theoryconnections{ideal theory}{%
\item Ring theory supplies ideals, quotient rings, and homomorphisms; commutative algebra supplies prime and maximal ideals; and algebraic geometry interprets ideals as polynomial conditions.
\item Number theory adds ideal factorization in rings of integers.
}{%
\item Ideal theory helps solve polynomial equations, construct quotient structures, and measure factorization when elements do not factor uniquely.
\item It is also the algebraic foundation for schemes, algebraic number theory, and many computational methods.
}
\section*{Glossary}
\begin{description}[leftmargin=*,style=nextline,itemsep=0.3em]
\item[\textbf{Ideal.}] An additive subgroup of a ring that absorbs multiplication by any ring element; it generalizes the multiples of an integer and makes quotient rings possible.
\item[\textbf{Prime ideal.}] A proper ideal $\mathfrak{p}$ such that if $ab\in\mathfrak{p}$ then $a\in\mathfrak{p}$ or $b\in\mathfrak{p}$; the quotient by a prime ideal is an integral domain.
\item[\textbf{Maximal ideal.}] A proper ideal $\mathfrak{m}$ that is not contained in any strictly larger proper ideal; the quotient by a maximal ideal is a field.
\item[\textbf{Quotient ring.}] The ring $R/I$ formed by identifying elements whose difference lies in the ideal $I$; it reduces arithmetic modulo the relations imposed by $I$.
\item[\textbf{Localization.}] The process of inverting selected elements of a ring to study local algebraic behavior near a chosen prime ideal.
\item[\textbf{Noetherian ring.}] A ring in which every ideal is finitely generated, guaranteeing that ascending chains of ideals eventually stabilize.
\item[\textbf{Radical ideal.}] An ideal equal to its own radical; an element is in the radical of $I$ if some power of it lies in $I$, and radical ideals correspond to algebraic sets via Hilbert's Nullstellensatz.
\item[\textbf{Principal ideal.}] An ideal generated by a single element, written $(a)=\{ra:r\in R\}$.
\item[\textbf{Dedekind domain.}] An integral domain in which every nonzero fractional ideal factors uniquely into prime ideals.
\item[\textbf{Nullstellensatz.}] Hilbert's theorem that over an algebraically closed field, radical ideals correspond to algebraic sets.
\item[\textbf{Chinese remainder theorem.}] When two ideals are comaximal, $R/(I\cap J)\cong R/I\times R/J$.
\item[\textbf{Hilbert basis theorem.}] If $R$ is Noetherian, then $R[x]$ is Noetherian.
\item[\textbf{Primary decomposition.}] Writing an ideal as an intersection of primary ideals, revealing prime components and multiplicities.
\end{description}
\section*{7. Sample Questions}
\emph{Exercise reference:} \cite{ideal_theory_ref1}. The cited edition is the source for the exercise set; consult its Exercises section for the official statement and numbering.
\begin{enumerate}[leftmargin=*,itemsep=0.9em]
\item \textbf{Exercise 1.} Why do the multiples of $6$ form an ideal of $\mathbb Z$?

\emph{Solution.} The sum or difference of two multiples of $6$ is a multiple of $6$, and multiplying a multiple of $6$ by any integer remains a multiple of $6$.

\item \textbf{Exercise 2.} Is $(5)$ a maximal ideal of $\mathbb Z$?

\emph{Solution.} Yes. The quotient $\mathbb Z/(5)$ is the field with five elements, so $(5)$ is maximal and therefore prime.

\item \textbf{Exercise 3.} What geometric set is defined by $(x,y)$ in $k[x,y]$?

\emph{Solution.} It is the single point $(0,0)$, because both $x$ and $y$ must vanish there.

\item \textbf{Exercise 4.} Why are prime ideals useful for factorization?

\emph{Solution.} The quotient by a prime ideal is a domain, so multiplication has no zero divisors there. Prime ideals also serve as the basic factors in ideal factorizations.

\item \textbf{Exercise 5.} What problem does localization solve?

\emph{Solution.} It allows selected nonzero elements to become invertible, so one can study algebraic behavior near a chosen prime ideal without carrying unrelated global information.

\end{enumerate}
\section*{8. References}
\addcontentsline{toc}{section}{References}

\begin{thebibliography}{9}
\bibitem{ideal_theory_ref1} David S. Dummit and Richard M. Foote, \emph{Abstract Algebra}, 3rd ed., Wiley, 2004.
\bibitem{ideal_theory_ref2} Atiyah and Macdonald, \emph{Introduction to Commutative Algebra}, Addison-Wesley, 1969.
\bibitem{ideal_theory_ref3} David Cox, John Little, and Donal O'Shea, \emph{Ideals, Varieties, and Algorithms}, 4th ed., Springer, 2015.
\bibitem{ideal_theory_ref4} Jürgen Neukirch, \emph{Algebraic Number Theory}, Springer, 1999.
\end{thebibliography}


\theory{Index theory}{How can the difference between solutions and obstructions of a differential equation be computed from topology?}{The index of a Fredholm operator is the dimension of its kernel minus the dimension of its cokernel. The Atiyah--Singer theorem identifies this analytic integer with a topological expression.}{Elliptic partial differential equations, geometry, topology, heat flow, quantum physics, and moduli spaces.}{The Fredholm index turns an analytic kernel-minus-cokernel difference into a stable topological quantity under perturbation.}

\paragraph{The problem and history}
An equation $D u=f$ may have solutions, many solutions, or no solution for some right-hand sides. The kernel contains solutions of $Du=0$. The cokernel records the independent compatibility conditions that $f$ must satisfy. For a Fredholm operator, both spaces are finite-dimensional and
\[
\operatorname{ind}(D)=\dim\ker D-\dim\operatorname{coker}D.
\]
Israel Gel'fand asked whether this integer could be expressed using topological data. Michael Atiyah and Isadore Singer answered yes in the 1960s \cite{index_theory_ref1,index_theory_ref2}.

\end{historylist}
\section*{3. Theory}
\subsection*{Core theory}
\paragraph{Why the difference is stable}
The dimensions of the kernel and cokernel can jump when an operator is perturbed. Their difference often does not. On a circle, the operator $D=d/dx-\lambda$ on periodic functions has a nonzero kernel only for special values of $\lambda$, but its index remains zero because the adjoint has a matching kernel.

Elliptic operators on compact manifolds are Fredholm. Ellipticity means that the highest-order part of the operator has no hidden direction in which information can propagate without control. Compactness prevents infinitely many independent low-energy solutions.

\paragraph{Symbols and the Atiyah--Singer formula}
The leading derivative part of an elliptic operator is its symbol. The symbol is a bundle map on the cotangent bundle and is invertible away from the zero covector. Its homotopy class defines a K-theory element. The analytic index depends only on this topological symbol class.

The Atiyah--Singer index theorem states
\[
\operatorname{ind}_{\mathrm{an}}(D)=\operatorname{ind}_{\mathrm{top}}(D).
\]
The topological side is built from characteristic classes of the manifold and bundles. One can therefore count analytic zero modes without solving the differential equation explicitly.

\paragraph{Major special cases}
The Riemann--Roch theorem computes the dimension difference between holomorphic sections and obstructions. The Hirzebruch--Riemann--Roch theorem extends it to complex manifolds. The signature theorem expresses the signature of the intersection form using the $L$-class. The Chern--Gauss--Bonnet theorem computes the Euler characteristic from curvature. The Dirac operator yields the Â-genus and spin-geometric results \cite{index_theory_ref3}.

\paragraph{Proof viewpoints and applications}
Atiyah and Singer first used K-theory and a push-forward construction, reducing the problem to spheres. Cobordism gives another reduction to simple generators. A heat-equation proof writes the index as a difference of traces of heat operators; the small-time asymptotic expansion reveals the characteristic classes. Supersymmetry explains cancellations in the heat expansion.

Index theory counts zero modes in gauge theory, instantons, and quantum field theory. It gives dimensions of moduli spaces, integrality constraints such as Rokhlin-type results, and topological information from elliptic boundary problems. The Atiyah--Patodi--Singer theorem adds boundary correction terms involving eta invariants.

\paragraph{What the theory depends on and where it is useful}
Index theory depends on functional analysis, elliptic PDE, differential geometry, topology, K-theory, and characteristic classes. It helps Hodge theory identify harmonic forms, geometry compute invariants, and physics count stable states.

The index does not usually give the separate dimensions of the kernel and cokernel. On noncompact spaces or for non-elliptic operators, finite-dimensionality and stability may fail; boundary conditions must then be chosen carefully.

\section*{4. Important theorems and results}
\begin{enumerate}[leftmargin=*,itemsep=0.35em]
\item \textbf{Atiyah--Singer index theorem.} The analytic and topological indices of an elliptic operator on a compact manifold are equal.

\item \textbf{Riemann--Roch theorem.} A holomorphic dimension difference is computed by degree and genus.

\item \textbf{Chern--Gauss--Bonnet theorem.} The Euler characteristic is the integral of a curvature characteristic form.

\item \textbf{Hirzebruch signature theorem.} The signature of a compact oriented manifold is an integral of its $L$-class.

\item \textbf{Atiyah--Patodi--Singer theorem.} Elliptic index formulas on manifolds with boundary include a boundary spectral correction \cite{index_theory_ref4}.

The index is an integer that measures an analytic imbalance, such as the difference between the dimensions of a kernel and a cokernel. Atiyah--Singer says that this analytic count equals a topological expression, so it can be computed without solving the differential equation directly. Riemann--Roch, Gauss--Bonnet, and the signature theorem are important special cases; Atiyah--Patodi--Singer shows what changes when a boundary contributes spectral information.
\end{enumerate}

\theoryapplications
\theoryconnections{index theory}{%
\item Functional analysis supplies operators and spectra, differential geometry supplies elliptic differential operators, and topology supplies characteristic classes and K-theory.
\item Measure and integration theory provide the analytic index, while homological algebra organizes the associated complexes.
}{%
\item Index theory helps compute analytic quantities from topological data in geometry and mathematical physics.
\item It explains why a problem involving infinitely many analytic modes can have a finite, stable integer answer.
}
\section*{Glossary}
\begin{description}[leftmargin=*,style=nextline,itemsep=0.3em]
\item[\textbf{Fredholm operator.}] A bounded linear operator between Banach spaces whose kernel and cokernel are both finite-dimensional, so that the difference of their dimensions (the index) is well-defined.
\item[\textbf{Kernel.}] The subspace of inputs mapped to zero by a linear operator; in index theory it represents independent solutions of $Du=0$.
\item[\textbf{Cokernel.}] The quotient of the target space by the image of a linear operator; it measures independent obstruction conditions that a right-hand side must satisfy.
\item[\textbf{Index.}] The integer $\dim\ker D - \dim\operatorname{coker} D$ of a Fredholm operator; it is stable under small perturbations and, by the Atiyah--Singer theorem, computable from topology.
\item[\textbf{Elliptic operator.}] A differential operator whose highest-order part is invertible away from the zero covector, ensuring good analytic behavior on compact manifolds.
\item[\textbf{Symbol.}] The leading-derivative part of an elliptic operator, viewed as a bundle map on the cotangent bundle; its homotopy class determines the index.
\item[\textbf{Characteristic class.}] A cohomology class associated to a vector bundle that encodes topological information; the Atiyah--Singer formula expresses the index as an integral of characteristic classes.
\item[\textbf{K-theory.}] A generalized cohomology theory built from stable equivalence classes of vector bundles.
\item[\textbf{Riemann--Roch theorem.}] Computes the index of the $\bar\partial$-operator on a complex manifold as a difference of holomorphic dimensions.
\item[\textbf{Chern--Gauss--Bonnet theorem.}] Expresses the Euler characteristic as the integral of a curvature characteristic form.
\item[\textbf{Eta invariant.}] A spectral invariant of an elliptic operator on a manifold with boundary.
\end{description}
\section*{7. Sample Questions}
\emph{Exercise reference:} \cite{index_theory_ref1}. The cited edition is the source for the exercise set; consult its Exercises section for the official statement and numbering.
\begin{enumerate}[leftmargin=*,itemsep=0.9em]
\item \textbf{Exercise 1.} Compute the Euler characteristic of $S^2$ using the Chern--Gauss--Bonnet theorem.

\emph{Solution.} The Chern--Gauss--Bonnet theorem states $\chi(M)=\frac{1}{2\pi}\int_M K\,dA$ for a closed oriented surface, where $K$ is the Gaussian curvature. For $S^2$ with the round metric of radius $1$, $K=1$ everywhere and $\operatorname{Area}(S^2)=4\pi$. Therefore $\chi(S^2)=\frac{1}{2\pi}\cdot 4\pi=2$.

\item \textbf{Exercise 2.} Compute the index of the operator $D=\frac{d}{dx}:\Omega^0(S^1)\to\Omega^1(S^1)$ on the circle (acting on smooth functions to produce exact 1-forms).

\emph{Solution.} The kernel of $D$ consists of functions with $df=0$, i.e., constant functions: $\dim\ker D=1$. The cokernel is $H^1_{\mathrm{dR}}(S^1)\cong\mathbb{R}$, so $\dim\operatorname{coker}D=1$. The index is $\operatorname{ind}(D)=1-1=0$.

\item \textbf{Exercise 3.} State the Riemann--Roch theorem for a compact Riemann surface of genus $g$, applied to the trivial line bundle $\mathcal{O}$.

\emph{Solution.} Riemann--Roch states $\ell(D)-\ell(K-D)=\deg(D)-g+1$, where $D$ is a divisor, $K$ a canonical divisor, $\ell(D)=\dim H^0(\mathcal{O}(D))$, and $\ell(K-D)=\dim H^1(\mathcal{O}(D))$. For $D=0$: $\ell(0)-\ell(K)=0-g+1=1-g$. Since $\ell(0)=1$ (only constants) and $\ell(K)=g$ (the space of holomorphic $1$-forms), we get $1-g=1-g$, a consistency check.

\item \textbf{Exercise 4.} Verify the Atiyah--Singer formula for the $\bar\partial$-operator on a compact Riemann surface of genus $g$.

\emph{Solution.} The index of $\bar\partial:L^2(\Omega^{0,0})\to L^2(\Omega^{0,1})$ is $\dim\ker\bar\partial-\dim\operatorname{coker}\bar\partial=g-1$. The analytic index counts holomorphic functions minus obstructions. The topological side, by Riemann--Roch, gives the same: $\frac{1}{2}(\deg(T^{1,0})-\deg(K))=\ldots$ The $\hat{A}$-genus on a surface of genus $g$ equals $1-g$ (using the standard formula). Both sides agree, confirming the theorem.

\item \textbf{Exercise 5.} Compute the signature of $S^2\times S^2$ using the intersection form.

\emph{Solution.} $H_2(S^2\times S^2;\mathbb{Z})\cong\mathbb{Z}^2$, generated by $a=[S^2\times\{pt\}]$ and $b=[\{pt\}\times S^2]$. The intersection form has matrix $\begin{pmatrix}a\cdot a&a\cdot b\\b\cdot a&b\cdot b\end{pmatrix}=\begin{pmatrix}0&1\\1&0\end{pmatrix}$. The eigenvalues are $+1$ and $-1$, giving one positive and one negative direction. The signature is $\sigma=1-1=0$.

\end{enumerate}
\section*{8. References}
\addcontentsline{toc}{section}{References}

\begin{thebibliography}{9}
\bibitem{index_theory_ref1} Michael F. Atiyah and Isadore M. Singer, \emph{The Index of Elliptic Operators I}, Annals of Mathematics, 1968.
\bibitem{index_theory_ref2} Michael F. Atiyah, \emph{K-Theory}, 2nd ed., Westview Press, 1989.
\bibitem{index_theory_ref3} H. Blaine Lawson and Marie-Louise Michelsohn, \emph{Spin Geometry}, Princeton University Press, 1989.
\bibitem{index_theory_ref4} Richard B. Melrose, \emph{The Atiyah--Patodi--Singer Index Theorem}, A K Peters, 1993.
\end{thebibliography}


\theory{Information theory}{How can uncertainty, storage, and communication be measured, and how can messages survive compression and noise?}{Information theory treats information as the resolution of uncertainty. It gives limits for representing data, sending it through a noisy channel, and recovering it reliably.}{Data compression, error correction, telecommunication, statistics, cryptography, machine learning, neuroscience, and physics.}{Entropy, mutual information, and channel capacity quantify uncertainty, compression limits, and the effect of noise.}

\paragraph{History and the communication problem}
\bookfigure{information_coding_channel}{Information and coding theory follow a message through encoding, noise, and decoding.}
Claude Shannon formalized information theory in 1948 at Bell Labs. Harry Nyquist had related transmission speed to the number of signal levels, and Ralph Hartley had measured the number of distinguishable symbol sequences. Shannon's problem was simple to state: reproduce at one place a message selected at another, even when the channel adds noise \cite{information_theory_ref1}.

The theory separates three jobs. Source coding removes predictable repetition so a message takes fewer symbols. Channel coding adds carefully designed redundancy so the receiver can correct errors. Cryptographic coding hides information from an unwanted receiver. The mathematics is based on probability, statistics, and logarithms.

\end{historylist}
\section*{3. Theory}
\subsection*{Core theory}
\paragraph{Information and entropy}
If an event has probability $p$, its self-information is
\[
I(p)=-\log_2p.
\]
A rare event tells us more when it occurs. For a random variable $X$ with probabilities $p_i$, Shannon entropy is
\[
H(X)=-\sum_i p_i\log_2p_i.
\]
It measures average uncertainty before the outcome is known. A fair coin has entropy $1$ bit; a coin that always lands heads has entropy $0$. The binary entropy function
\[
H_b(p)=-p\log_2p-(1-p)\log_2(1-p)
\]
is largest at $p=1/2$.

The logarithm base determines units: bits or shannons for base two, nats for the natural logarithm, and hartleys for base ten. Entropy is largest when all possible symbols are equally likely. For independent variables, joint entropy adds; for dependent variables, knowing one reduces uncertainty about the other.

\paragraph{Conditional and mutual information}
Conditional entropy $H(X\mid Y)$ measures uncertainty about $X$ after $Y$ is known. Mutual information is
\[
I(X;Y)=H(X)-H(X\mid Y).
\]
It measures how much observing $Y$ teaches us about $X$. If two variables are independent, mutual information is zero. In a medical test, $X$ may be whether a disease is present and $Y$ a test result; mutual information measures the average reduction in diagnostic uncertainty.

The Kullback--Leibler divergence compares a proposed distribution with the true distribution:
\[
D(P\Vert Q)=\sum_xP(x)\log\frac{P(x)}{Q(x)}.
\]
It is nonnegative but not symmetric, so it is a divergence rather than an ordinary distance. Directed information measures information flowing from a past sequence to a future sequence and is useful for channels with memory.

\paragraph{Source coding and compression}
The source coding theorem says that a long sequence from a stationary source can be compressed to approximately its entropy rate bits per symbol, but not reliably below that rate. Huffman and arithmetic coding approach this limit for practical sources. ZIP files exploit repeated patterns; images, sound, and video use models of visual or auditory redundancy.

If a source has $n$ equally likely symbols, its entropy is $\log_2n$. If symbols are correlated, a compressor predicts the next symbol from the previous ones and stores only the surprise. Lossless compression reconstructs the exact data. Lossy compression permits a controlled distortion to achieve a much smaller representation.

\paragraph{Noisy channels and capacity}
A channel maps an input symbol to a random output. Its capacity is the greatest rate at which information can be transmitted with arbitrarily small error probability using long codes. Shannon's noisy-channel coding theorem says that rates below capacity are achievable, while rates above capacity are not \cite{information_theory_ref2}.

For a binary symmetric channel that flips a bit with probability $p$, the capacity is $1-H_b(p)$ bits per use. For a band-limited Gaussian channel with bandwidth $B$, signal power $S$, and noise power $N$, the Shannon--Hartley law gives
\[
C=B\log_2(1+S/N).
\]
Capacity is a limit, not a promise that every short practical code reaches it.

\paragraph{Codes, redundancy, and cryptography}
Error-correcting codes deliberately add redundancy. A repetition code sends 0 as 000 and 1 as 111; a majority vote corrects one error but uses bandwidth inefficiently. Hamming, Reed--Solomon, convolutional, turbo, and low-density parity-check codes approach capacity much more effectively. They made compact discs, mobile phones, DSL, satellite communication, and Voyager data possible \cite{information_theory_ref3}.

Cryptography uses codes and ciphers to control who can recover a message. Information-theoretic secrecy asks for security even against unlimited computation; the one-time pad achieves it when used correctly. Computational cryptography instead relies on hard problems such as discrete logarithms. Information theory also studies cryptanalysis, secrecy capacity, and leakage from side channels.

\paragraph{Applications across science}
In statistics, entropy and divergence measure model fit and guide inference. In machine learning, cross-entropy trains probabilistic predictors, while mutual information measures feature relevance. In neuroscience, information measures how a neural response reveals a stimulus. In bioinformatics it studies molecular codes and sequence evolution. In physics, thermodynamic entropy, black-hole entropy, Landauer's principle, and quantum information connect information to energy and matter.

Other applications include seismic exploration, signal processing, pattern recognition, anomaly detection, music and art analysis, imaging design, information retrieval, plagiarism detection, intelligence gathering, semiotics, and the study of biological and social networks \cite{information_theory_ref4}.

\paragraph{What the theory depends on and where it is useful}
Information theory depends on probability, statistics, logarithms, coding, and optimization. It helps coding theory design explicit codes, statistics quantify uncertainty, cryptography analyze secrecy, and machine learning compare distributions. Its model of information is mathematical and probabilistic; it does not by itself measure meaning, truth, or human understanding.

\section*{4. Important theorems and results}
\begin{enumerate}[leftmargin=*,itemsep=0.35em]
\item \textbf{Source coding theorem.} Entropy is the asymptotic lower limit for reliable lossless compression.

\item \textbf{Noisy-channel coding theorem.} Every rate below channel capacity can be transmitted with error tending to zero using suitable long codes; rates above capacity cannot.

\item \textbf{Data-processing inequality.} Processing data cannot increase the mutual information it contains about an original variable.

\item \textbf{Shannon--Hartley law.} Gaussian-channel capacity is $B\log_2(1+S/N)$.

\item \textbf{Fano's inequality.} The probability of decoding error bounds how much uncertainty remains, linking estimation accuracy to information \cite{information_theory_ref1}.
\end{enumerate}

\theoryapplications
\theoryconnections{information theory}{%
\item Probability supplies random variables and conditional distributions, while measure theory handles continuous sources.
\item Combinatorics and coding theory construct messages, and optimization studies the trade-off between rate, distortion, and reliability.
}{%
\item Information theory helps communication, compression, statistics, machine learning, cryptography, and physics quantify uncertainty and dependence.
\item Entropy and mutual information are useful only after the source, observation, and noise model have been specified.
}
\section*{Glossary}
\begin{description}[leftmargin=*,style=nextline,itemsep=0.3em]
\item[\textbf{Entropy.}] The average uncertainty of a random variable, measured in bits; it is maximized when all outcomes are equally likely.
\item[\textbf{Mutual information.}] The average reduction in uncertainty about one random variable obtained by observing another; it quantifies statistical dependence.
\item[\textbf{Channel capacity.}] The maximum rate at which information can be transmitted through a noisy channel with arbitrarily low error probability using suitable long codes.
\item[\textbf{Source coding.}] Compression of a data source to use as few bits as possible; the source coding theorem shows that the entropy rate is the fundamental limit for lossless compression.
\item[\textbf{Channel coding.}] The deliberate addition of structured redundancy to a message so that a receiver can detect and correct errors introduced by a noisy channel.
\item[\textbf{Kullback--Leibler divergence.}] A non-symmetric measure of how one probability distribution differs from another; it is nonnegative and zero only when the distributions agree.
\item[\textbf{Conditional entropy.}] $H(X\mid Y)$, the remaining uncertainty about $X$ after $Y$ is observed.
\item[\textbf{Self-information.}] $I(p)=-\log_2 p$, the information content of an event with probability $p$.
\item[\textbf{Binary symmetric channel.}] A channel that independently flips each input bit with probability $p$.
\item[\textbf{Shannon--Hartley law.}] The formula $C=B\log_2(1+S/N)$ for the capacity of a band-limited Gaussian channel.
\item[\textbf{Data-processing inequality.}] Processing data cannot increase mutual information about an original variable.
\end{description}
\section*{7. Sample Questions}
\emph{Exercise reference:} \cite{information_theory_ref1}. The cited edition is the source for the exercise set; consult its Exercises section for the official statement and numbering.
\begin{enumerate}[leftmargin=*,itemsep=0.9em]
\item \textbf{Exercise 1.} A random variable $X$ takes values in $\{a,b,c,d\}$ with probabilities $\frac{1}{2},\frac{1}{4},\frac{1}{8},\frac{1}{8}$. Compute $H(X)$.

\emph{Solution.} $H(X)=-\frac{1}{2}\log_2\frac{1}{2}-\frac{1}{4}\log_2\frac{1}{4}-\frac{1}{8}\log_2\frac{1}{8}-\frac{1}{8}\log_2\frac{1}{8}=\frac{1}{2}\cdot 1+\frac{1}{4}\cdot 2+\frac{1}{8}\cdot 3+\frac{1}{8}\cdot 3=\frac{1}{2}+\frac{1}{2}+\frac{3}{8}+\frac{3}{8}=\frac{7}{4}=1.75$ bits.

\item \textbf{Exercise 2.} Compute $H(X,Y)$ for two independent fair coins $X$ and $Y$.

\emph{Solution.} Since $X$ and $Y$ are independent with $H(X)=H(Y)=1$ bit, $H(X,Y)=H(X)+H(Y)=2$ bits. There are four equally likely outcomes, each with probability $1/4$, so $H(X,Y)=-4\cdot\frac{1}{4}\log_2\frac{1}{4}=2$ bits.

\item \textbf{Exercise 3.} A binary symmetric channel has crossover probability $p=0.25$. Compute its capacity in bits per use.

\emph{Solution.} $C=1-H_b(p)=1-H_b(0.25)=-0.25\log_2(0.25)-0.75\log_2(0.75)=0.5+0.75\cdot 0.41503\ldots=0.5+0.31127\ldots=0.8113$ bits. So $C\approx 0.1887$ bits per use.

\item \textbf{Exercise 4.} A source emits symbols from $\{A,B,C,D,E\}$ with equal probability. How many bits per symbol are needed for lossless compression, and how many bits would a naive fixed-length code use?

\emph{Solution.} Entropy: $H(X)=\log_2 5\approx 2.3219$ bits per symbol. A fixed-length code uses $\lceil\log_2 5\rceil=3$ bits per symbol. The source coding theorem says the optimal lossless code uses approximately $2.322$ bits per symbol, a savings of about $23\%$ over the fixed-length code.

\item \textbf{Exercise 5.} Two coins $X$ and $Y$ are not independent: $X$ is fair, and $Y=X$ with probability $0.9$ and $Y\neq X$ with probability $0.1$ (each with $0.05$ for the other value). Compute $I(X;Y)$.

\emph{Solution.} $H(X)=1$ bit. We compute $H(Y)$: $P(Y=H)=0.9\cdot 0.5+0.05\cdot 0.5=0.475$, $P(Y=T)=0.475$. So $H(Y)=-2(0.475)\log_2(0.475)\approx 0.9988$ bits. $H(X,Y)=\sum_{x,y}P(x,y)\log_2(1/P(x,y))$: $P(H,H)=0.45$, $P(T,T)=0.45$, $P(H,T)=0.05$, $P(T,H)=0.05$. So $H(X,Y)=-2(0.45)\log_2(0.45)-2(0.05)\log_2(0.05)\approx 0.9170+0.4322=1.3493$ bits. $I(X;Y)=H(X)+H(Y)-H(X,Y)=1+0.9988-1.3493\approx 0.6495$ bits.

\end{enumerate}
\section*{8. References}
\addcontentsline{toc}{section}{References}

\begin{thebibliography}{9}
\bibitem{information_theory_ref1} Claude E. Shannon, "A mathematical theory of communication," \emph{Bell System Technical Journal}, 27 (1948), 379--423 and 623--656.
\bibitem{information_theory_ref2} Thomas M. Cover and Joy A. Thomas, \emph{Elements of Information Theory}, 2nd ed., Wiley, 2006.
\bibitem{information_theory_ref3} David J. C. MacKay, \emph{Information Theory, Inference, and Learning Algorithms}, Cambridge University Press, 2003.
\bibitem{information_theory_ref4} Robert B. Ash, \emph{Information Theory}, Dover, 1990.
\end{thebibliography}


\theory{Intersection theory}{How can the meeting of geometric objects be counted when they touch, overlap, or intersect in high dimensions?}{Intersection theory assigns algebraic multiplicities to intersections and organizes them in rings of cycles. It extends Bezout's counting principle to geometry where ordinary point counting is not enough.}{Algebraic geometry, topology, enumerative geometry, moduli spaces, and mathematical physics.}{Cycles, divisors, and refined multiplicities count how geometric objects meet, including intersections invisible in a naive drawing.}

\paragraph{The problem and its history}
Two lines in a plane usually meet once. A line and a parabola may meet twice. If the line is tangent, the two meetings merge into one visible point, but the equation still has multiplicity two. Intersection theory preserves this hidden count. Its roots lie in Bezout's theorem, elimination theory, and the work of Weil and van der Waerden; William Fulton developed a modern framework \cite{intersection_theory_ref1,intersection_theory_ref2}.

For subvarieties $A$ and $B$ of a smooth variety $X$, a proper intersection has the expected dimension
\[
\dim(A\cap B)=\dim A+\dim B-\dim X.
\]
The intersection product $A\cdot B$ is a cycle supported on $A\cap B$, with coefficients recording local multiplicity.

\end{historylist}
\section*{3. Theory}
\subsection*{Core theory}
\paragraph{A local example}
The line $y=0$ and parabola $y=x^2$ meet at the origin. Substituting gives $x^2=0$, a double root, so the intersection multiplicity is $2$. If the line is moved to $y=\epsilon$, there are two nearby intersection points over the complex numbers, both approaching the origin as $\epsilon$ tends to zero. Multiplicity is therefore a continuity-preserving count.

\paragraph{Moving cycles and rational equivalence}
Objects in bad position may overlap too much: parallel lines in an affine plane have no meeting, while a plane containing a line has an intersection of unexpectedly high dimension. One moves cycles within a rational-equivalence class until they meet properly. A rational function on a subvariety produces a difference between its zero and pole divisors; such differences are declared equivalent to zero.

The Chow group is the group of cycles modulo rational equivalence. Intersection products make the collection of Chow groups into a Chow ring when the ambient variety is smooth. This ring retains algebraic information that ordinary cohomology may forget.

\paragraph{Topological intersection forms}
On an oriented manifold of dimension $2n$, the middle cohomology has a bilinear intersection form
\[
\lambda(a,b)=\langle a\smile b,[M]\rangle.
\]
It evaluates the cup product on the fundamental class. On a four-manifold this form is especially important: its signature is the number of positive directions minus negative directions. Poincare duality makes the form nondegenerate under suitable hypotheses.

\paragraph{Self-intersection and blow-ups}
An object may intersect a moved copy of itself. For a smooth subvariety, the self-intersection is governed by the top Chern class of its normal bundle. A line in the projective plane has self-intersection $1$. A ruling line in $\mathbb P^1\times\mathbb P^1$ has self-intersection $0$ because it can be moved away from itself.

Blowing up a point on a surface replaces the point by an exceptional curve. This curve has genus zero and self-intersection $-1$. Castelnuovo's contraction theorem says that a suitable $(-1)$-curve can be blown down again. These calculations are central in birational classification \cite{intersection_theory_ref3}.

\paragraph{Modern developments and uses}
Intersection theory counts curves satisfying geometric conditions, such as lines meeting prescribed objects. It underlies enumerative geometry, Chern classes, Riemann--Roch, Gromov--Witten invariants, quantum cohomology, and moduli spaces. Virtual fundamental cycles supply a corrected intersection class when a moduli space is singular or has the wrong dimension. Current developments extend the theory from schemes to stacks and connect it to quantum intersection rings.

\paragraph{What the theory depends on and where it is useful}
Intersection theory depends on algebraic geometry, divisors, cycles, commutative algebra, cohomology, and topology. It helps elimination theory count solutions, topology study intersection forms, and physics compute enumerative and quantum invariants.

Nonprojective varieties can lack a well-behaved intersection theory, and singular or nonproper intersections require refined tools such as virtual cycles or derived intersections.

\section*{4. Important theorems and results}
\begin{enumerate}[leftmargin=*,itemsep=0.35em]
\item \textbf{Bezout's theorem.} Plane curves of degrees $m$ and $n$ meet in $mn$ points counted with multiplicity over an algebraically closed field when they have no common component.

\item \textbf{Moving lemma.} Cycles can often be replaced by rationally equivalent cycles in proper position.

\item \textbf{Self-intersection formula.} The self-intersection of a smooth subvariety is represented by the top Chern class of its normal bundle.

\item \textbf{Poincare duality.} The topological intersection pairing on a closed oriented manifold is nondegenerate.

\item \textbf{Castelnuovo contraction theorem.} A suitable smooth rational curve of self-intersection $-1$ can be contracted to a smooth point \cite{intersection_theory_ref2}.
\end{enumerate}

\theoryapplications
\theoryconnections{intersection theory}{%
\item Algebraic geometry supplies varieties, divisors, and local rings; commutative algebra supplies multiplicity; and topology supplies homology and cohomology pairings.
\item Chow groups and Chern classes extend the elementary idea of counting intersections.
}{%
\item Intersection theory helps algebraic geometry count solutions with multiplicity and helps arithmetic geometry attach numerical invariants to cycles.
\item It also connects geometric incidence problems to moduli spaces, enumerative geometry, and characteristic classes.
}
\section*{Glossary}
\begin{description}[leftmargin=*,style=nextline,itemsep=0.3em]
\item[\textbf{Intersection multiplicity.}] A non-negative integer recording how many times two subvarieties meet at a point, taking tangency and higher-order contact into account.
\item[\textbf{Proper intersection.}] An intersection whose dimension equals the sum of the dimensions of the intersecting subvarieties minus the ambient dimension, so no component is unexpectedly large.
\item[\textbf{Chow group.}] The group of algebraic cycles modulo rational equivalence; on a smooth variety it carries a product making it the Chow ring.
\item[\textbf{Rational equivalence.}] An equivalence relation on cycles generated by the zero and pole divisors of rational functions; cycles in the same class are considered interchangeable for intersection purposes.
\item[\textbf{Intersection form.}] A bilinear pairing on middle cohomology of an oriented even-dimensional manifold, given by evaluating the cup product on the fundamental class.
\item[\textbf{Self-intersection.}] The intersection number of a subvariety with a nearby perturbation of itself, controlled by the top Chern class of its normal bundle.
\item[\textbf{Blow-up.}] A geometric operation that replaces a point by the projective space of directions through it, used to resolve singularities and study local intersection behavior.
\item[\textbf{Bézout's theorem.}] Plane curves of degrees $m$ and $n$ with no common component meet in exactly $mn$ points counted with multiplicity.
\item[\textbf{Normal bundle.}] The quotient bundle $N_{Y/X}=TX|_Y/TY$ describing how a subvariety sits inside an ambient variety.
\item[\textbf{Chern class.}] A characteristic class of a vector bundle living in cohomology.
\item[\textbf{Cup product.}] The associative product on cohomology that pairs classes in complementary dimensions.
\item[\textbf{Fundamental class.}] The distinguished homology or cohomology class $[M]$ of a closed oriented manifold.
\item[\textbf{Divisor.}] A formal sum of codimension-one subvarieties; divisors generate the Chow ring in many cases.
\item[\textbf{Enumerative geometry.}] The branch of algebraic geometry counting geometric objects satisfying prescribed conditions.
\end{description}
\section*{7. Sample Questions}
\emph{Exercise reference:} \cite{intersection_theory_ref1}. The cited edition is the source for the exercise set; consult its Exercises section for the official statement and numbering.
\begin{enumerate}[leftmargin=*,itemsep=0.9em]
\item \textbf{Exercise 1.} Why does $y=x^2$ meet $y=0$ with multiplicity two at the origin?

\emph{Solution.} Substitution gives $x^2=0$, which has a root of order two. A small parallel displacement produces two nearby complex intersections.

\item \textbf{Exercise 2.} What is a proper intersection?

\emph{Solution.} It has the expected dimension: the dimensions add and the ambient dimension is subtracted. No component is unexpectedly large.

\item \textbf{Exercise 3.} Why move cycles before intersecting?

\emph{Solution.} A bad position can hide the intended count. Rational equivalence allows a harmless movement to proper position while preserving the intersection class.

\item \textbf{Exercise 4.} What does self-intersection measure?

\emph{Solution.} It measures how a subvariety meets a nearby copy of itself and is controlled by the normal bundle describing how it sits in the ambient space.

\item \textbf{Exercise 5.} Why is a $(-1)$-curve important?

\emph{Solution.} It is the exceptional curve created by blowing up a point and can often be contracted, making it a basic step in simplifying a surface.

\end{enumerate}
\section*{8. References}
\addcontentsline{toc}{section}{References}

\begin{thebibliography}{9}
\bibitem{intersection_theory_ref1} William Fulton, \emph{Intersection Theory}, 2nd ed., Springer, 1998.
\bibitem{intersection_theory_ref2} David Eisenbud and Joe Harris, \emph{3264 and All That}, Cambridge University Press, 2016.
\bibitem{intersection_theory_ref3} Robin Hartshorne, \emph{Algebraic Geometry}, Springer, 1977.
\end{thebibliography}


\theory{Invariant theory}{Which quantities remain unchanged when a group transforms the object?}{Invariant theory studies functions that do not change under symmetry. It turns classification up to symmetry into the study of invariant rings, orbit spaces, and geometric quotients.}{Algebra, representation theory, algebraic geometry, geometry, robotics, chemistry, mechanics, and moduli spaces.}{Group actions and invariant rings retain polynomial information unchanged by a symmetry and organize quotient or moduli constructions.}

\paragraph{The question and history}
Suppose a shape is rotated, reflected, or relabeled. Its coordinates change, but some quantities do not. The determinant is unchanged when a matrix is left-multiplied by an element of $SL_n$, and $x+y$ and $xy$ are unchanged when two variables are swapped. Invariant theory asks for all such quantities and whether they distinguish different orbits \cite{invariant_theory_ref1}.

The subject began with George Boole and Arthur Cayley in the nineteenth century, who studied polynomial forms under linear transformations. It later connected to symmetric functions, representation theory, commutative algebra, moduli spaces, and geometric invariant theory. David Mumford's formulation emphasized constructing quotients of varieties by group actions.

\end{historylist}
\section*{3. Theory}
\subsection*{Core theory}
\paragraph{Group actions and invariant rings}
Let a group $G$ act linearly on a vector space $V$ over a field $k$. The induced action on polynomial functions $k[V]$ is obtained by transforming the input. The invariant ring is
\[
k[V]^G=\{f\in k[V]:f(gv)=f(v)\text{ for every }g\in G,v\in V\}.
\]
Each invariant is constant along a group orbit. The orbit is the set of positions reachable from one vector by the symmetry.

For the action of $\mathbb Z/2$ sending $(x,y)$ to $(-x,-y)$, the invariant monomials include $x^2$, $xy$, and $y^2$. They generate
\[
k[x,y]^{\mathbb Z/2}\cong k[a,b,c]/(ac-b^2).
\]
The relation records that $(x^2)(y^2)=(xy)^2$.

\paragraph{Symmetric functions and classical forms}
Swapping variables leaves the elementary symmetric polynomials unchanged. The fundamental theorem of symmetric polynomials says every symmetric polynomial is a polynomial in these elementary ones. For binary quadratic and higher forms, classical invariant theory finds discriminants, determinants, and other polynomial combinations that survive linear changes of variables.

The discriminant of a polynomial is an invariant of its roots: it vanishes exactly when two roots coincide. In geometry, lengths, angles, areas, and determinants are invariants under different transformation groups. In chemistry, molecular symmetry divides configurations into equivalence classes and predicts allowed transitions.

\paragraph{Hilbert's theorems}
For a finite group acting over a field whose characteristic does not divide the group order, averaging over the group gives a Reynolds operator. Hilbert proved finite generation results for invariant rings of important groups. The Hilbert--Nagata theorem says that for a linearly reductive group, the invariant ring of a finitely generated algebra is finitely generated \cite{invariant_theory_ref2}.

Finite generation changes an infinite search for invariants into a finite algebraic problem, although finding an explicit generating set may still be difficult. Molien series count invariants by degree for finite groups.

\paragraph{Geometric invariant theory}
The orbit space of a group action need not be a nice algebraic variety. Orbits may not be closed, and two different orbits may have closures that meet. Geometric invariant theory keeps the good, closed orbits and identifies points whose closures meet. The invariant ring defines a quotient that captures stable information.

This produces moduli spaces of geometric objects: configurations of points, curves, vector bundles, instantons, and monopoles. Stability conditions exclude degenerate configurations that would make the quotient pathological. Symplectic geometry and equivariant topology developed parallel versions of the same quotient idea \cite{invariant_theory_ref3}.

\paragraph{What the theory depends on and where it is useful}
Invariant theory depends on group actions, polynomial rings, representation theory, commutative algebra, and algebraic geometry. It helps representation theory decompose spaces, geometry classify shapes, and moduli theory build spaces of objects up to change of coordinates.

Invariants need not separate all orbits. Two genuinely different configurations may share every invariant available in a chosen ring, especially when the group is not reductive or the field has special characteristic. One must therefore state what class of objects the invariants classify.

\section*{4. Important theorems and results}
\begin{enumerate}[leftmargin=*,itemsep=0.35em]
\item \textbf{Fundamental theorem of symmetric polynomials.} Every symmetric polynomial is a polynomial in the elementary symmetric polynomials.

\item \textbf{Hilbert basis theorem for invariants.} In many reductive settings the invariant ring is finitely generated.

\item \textbf{Noether bound.} For a finite group in nonmodular characteristic, invariants of bounded degree generate the invariant ring.

\item \textbf{Hilbert--Mumford criterion.} Stability in geometric invariant theory can be tested using one-parameter subgroups.

\item \textbf{First fundamental theorem.} For classical groups, specified natural contractions and pairings generate the invariant tensors \cite{invariant_theory_ref4}.
\end{enumerate}

\theoryapplications
\theoryconnections{invariant theory}{%
\item Group actions supply the transformations, ring theory supplies polynomial invariants, and representation theory supplies decompositions into symmetry types.
\item Algebraic geometry contributes orbit spaces and geometric invariant theory, while combinatorics provides symmetric functions.
}{%
\item Invariant theory helps classify objects up to symmetry, construct coordinates for quotient spaces, and analyze tensors in geometry and physics.
\item An invariant can prove that two objects are different, but a complete classification may require several invariants together.
}
\section*{Glossary}
\begin{description}[leftmargin=*,style=nextline,itemsep=0.3em]
\item[\textbf{Invariant.}] A function on a space that is unchanged by every element of a group action; it provides information that is preserved under the symmetry.
\item[\textbf{Orbit.}] The set of all positions reachable from a given point by applying every element of the acting group; two points in the same orbit are indistinguishable by invariants.
\item[\textbf{Invariant ring.}] The ring of all polynomial functions unchanged by a group action; it provides a finite coordinate system for the quotient.
\item[\textbf{Geometric invariant theory (GIT).}] A framework for constructing quotient varieties by removing unstable or badly behaved orbits, producing moduli spaces of geometric objects.
\item[\textbf{Reynolds operator.}] An averaging map over a finite group that projects a polynomial onto its invariant part; it exists when the group order is invertible in the base field.
\item[\textbf{Elementary symmetric polynomial.}] The basic polynomials in the roots of a polynomial that are unchanged under any permutation of the roots; every symmetric polynomial is a polynomial in these.
\item[\textbf{Reductive group.}] An algebraic group whose representation theory is well-behaved; the Hilbert--Nagata theorem guarantees finite generation of the invariant ring for reductive groups.
\item[\textbf{Group action.}] A homomorphism from a group to the automorphism group of a space.
\item[\textbf{Discriminant.}] A polynomial in the coefficients that vanishes when a polynomial has a repeated root; invariant under permutations.
\item[\textbf{Stability.}] In GIT, a condition on an orbit ensuring the quotient is well-separated.
\item[\textbf{Molien series.}] A generating function whose coefficients count invariant polynomials of each degree.
\item[\textbf{Noether bound.}] For a finite group in non-modular characteristic, invariants of degree at most the group order generate the invariant ring.
\item[\textbf{Hilbert--Mumford criterion.}] A test for stability in GIT using one-parameter subgroups.
\end{description}
\section*{7. Sample Questions}
\emph{Exercise reference:} \cite{invariant_theory_ref1}. The cited edition is the source for the exercise set; consult its Exercises section for the official statement and numbering.
\begin{enumerate}[leftmargin=*,itemsep=0.9em]
\item \textbf{Exercise 1.} The group $\mathbb{Z}/2\mathbb{Z}$ acts on $\mathbb{R}^2$ by $(x,y)\mapsto(-x,-y)$. Find the invariant ring $\mathbb{R}[x,y]^{\mathbb{Z}/2}$.

\emph{Solution.} A monomial $x^a y^b$ is invariant iff $(-x)^a(-y)^b=x^a y^b$, i.e., $a+b$ is even. The invariant ring is generated by $u=x^2$, $v=xy$, $w=y^2$, with the relation $uw=v^2$. So $\mathbb{R}[x,y]^{\mathbb{Z}/2}\cong\mathbb{R}[u,v,w]/(uw-v^2)$.

\item \textbf{Exercise 2.} Find the three elementary symmetric polynomials in variables $x_1,x_2,x_3$.

\emph{Solution.} $e_1=x_1+x_2+x_3$, $e_2=x_1 x_2+x_1 x_3+x_2 x_3$, $e_3=x_1 x_2 x_3$. By the fundamental theorem of symmetric polynomials, every symmetric polynomial in $x_1,x_2,x_3$ is a polynomial in $e_1,e_2,e_3$.

\item \textbf{Exercise 3.} Compute the discriminant of the quadratic polynomial $x^2+bx+c$ (as a polynomial in $x$).

\emph{Solution.} The roots are $r_1,r_2$ with $r_1+r_2=-b$ and $r_1 r_2=c$. The discriminant is $\Delta=(r_1-r_2)^2=(r_1+r_2)^2-4r_1 r_2=b^2-4c$. It vanishes iff the polynomial has a repeated root.

\item \textbf{Exercise 4.} Let $S_3$ act on $\mathbb{R}[x_1,x_2,x_3]$ by permuting the variables. Express the monomial symmetric polynomial $m_{(2,1)}=x_1^2 x_2+x_1^2 x_3+x_2^2 x_1+x_2^2 x_3+x_3^2 x_1+x_3^2 x_2$ in terms of elementary symmetric polynomials.

\emph{Solution.} $m_{(2,1)}=e_1 e_2-3e_3$. Verification: $e_1 e_2=(x_1+x_2+x_3)(x_1 x_2+x_1 x_3+x_2 x_3)=m_{(2,1)}+3x_1 x_2 x_3=m_{(2,1)}+3e_3$, so $m_{(2,1)}=e_1 e_2-3e_3$.

\item \textbf{Exercise 5.} Determine the Molien series for the action of $\mathbb{Z}/2\mathbb{Z}$ on $\mathbb{C}[x]$ by $x\mapsto -x$.

\emph{Solution.} The Molien series is $P(t)=\frac{1}{|G|}\sum_{g\in G}\frac{1}{\det(I-t\cdot g)}$. The identity acts as $1$ on $x$, contributing $\frac{1}{1-t}$. The nontrivial element acts as $-1$, contributing $\frac{1}{1+t}$. So $P(t)=\frac{1}{2}\left(\frac{1}{1-t}+\frac{1}{1+t}\right)=\frac{1}{2}\cdot\frac{2}{1-t^2}=\frac{1}{1-t^2}=1+t^2+t^4+\cdots$. This confirms the invariant ring is $\mathbb{C}[x^2]$, generated by the single invariant $x^2$ in degree 2.

\end{enumerate}
\section*{8. References}
\addcontentsline{toc}{section}{References}

\begin{thebibliography}{9}
\bibitem{invariant_theory_ref1} David Mumford, John Fogarty, and Frances Kirwan, \emph{Geometric Invariant Theory}, 3rd ed., Springer, 1994.
\bibitem{invariant_theory_ref2} Peter J. Olver, \emph{Classical Invariant Theory}, Cambridge University Press, 1999.
\bibitem{invariant_theory_ref3} Igor Dolgachev, \emph{Lectures on Invariant Theory}, Cambridge University Press, 2003.
\bibitem{invariant_theory_ref4} Hermann Weyl, \emph{The Classical Groups}, Princeton University Press, 1939.
\end{thebibliography}


\theory{Iwasawa theory}{What happens to arithmetic information when a number field is enlarged through an infinite tower?}{Iwasawa theory studies Galois modules, especially ideal class groups, throughout towers of number fields. A single infinite module records how arithmetic invariants grow at every finite level.}{Cyclotomic fields, class groups, p-adic analysis, p-adic L-functions, elliptic curves, motives, and modern number theory.}{p-adic towers and modules measure how arithmetic invariants grow across infinitely many related number fields.}

\paragraph{History and motivation}
Kenkichi Iwasawa began the theory in 1959 while studying cyclotomic fields. The p-part of an ideal class group measures a failure of unique factorization and was connected by Kummer to the difficulty of Fermat's Last Theorem. Instead of studying one field at a time, Iwasawa studied an infinite tower and asked for a law governing the class groups at all levels \cite{iwasawa_theory_ref1,iwasawa_theory_ref2}.

\end{historylist}
\section*{3. Theory}
\subsection*{Core theory}
\paragraph{The basic tower}
Fix a prime $p$. A $\mathbb Z_p$-extension of a number field $F$ is an infinite Galois extension $F_\infty/F$ whose Galois group is the additive group $\mathbb Z_p$. It has finite layers
\[
F=F_0\subset F_1\subset F_2\subset\cdots\subset F_\infty
\]
with $\operatorname{Gal}(F_n/F)\cong\mathbb Z/p^n\mathbb Z$. The group $\mathbb Z_p$ is the inverse limit of the finite cyclic groups, so the infinite tower is assembled from finite approximations.

\paragraph{Concrete illustration: $\mathbb Z_p$-extensions of $\mathbb Q$}
The simplest tower to visualize starts with $F=\mathbb Q$ and a prime $p$. The cyclotomic $\mathbb Z_p$-extension is obtained by adjoining all $p$-power roots of unity: set $K_n=\mathbb Q(\zeta_{p^{n+1}})$ and let $K_\infty=\bigcup_n K_n$. The Galois group $\operatorname{Gal}(K_\infty/\mathbb Q)$ contains a subgroup isomorphic to $\mathbb Z_p$, obtained by restricting the cyclotomic character. At level $n$ the relevant subfield $F_n$ satisfies $\operatorname{Gal}(F_n/\mathbb Q)\cong\mathbb Z/p^n\mathbb Z$, so $F_0=\mathbb Q$, $F_1$ is the unique subfield of $\mathbb Q(\zeta_{p^2})$ of degree $p$ over $\mathbb Q$, and so on. For $p=5$: $F_1$ is the degree-5 subfield of $\mathbb Q(\zeta_{25})$, generated by the trace $\alpha=\zeta_{25}+\zeta_{25}^2+\zeta_{25}^3+\zeta_{25}^4+\zeta_{25}^5+\cdots+\zeta_{25}^{24}$---wait, more precisely $F_1=\mathbb Q(\eta)$ where $\eta=\sum_{i=1}^{4}\zeta_{25}^{5i}$. The class number of $F_0=\mathbb Q$ is 1, so the Iwasawa module $I_0$ is trivial. As one ascends, the $p$-parts of the class groups $I_n$ can grow, and the Iwasawa algebra $\Lambda=\mathbb Z_p[[\Gamma]]$ (with $\Gamma\cong\mathbb Z_p$) acts on the inverse limit. The structure theorem then controls the growth: the formula $v_p(|I_n|)=\mu\cdot 5^n+\lambda n+\nu$ predicts how the class-group size eventually behaves at high levels of the tower.

For a cyclotomic example in the standard setup, begin with $K=\mathbb Q(\mu_p)$ and adjoin $p^{n+1}$st roots of unity. The union $K_\infty$ has Galois group $\mathbb Z_p$ over the appropriate base. Let $I_n$ be the $p$-primary part of the ideal class group of $K_n$. Norm maps $I_m\to I_n$ form an inverse system, and $I=\varprojlim I_n$ becomes a module over the Iwasawa algebra
\[
\Lambda=\mathbb Z_p[[\Gamma]].
\]

\paragraph{Why the algebra helps}
If $\Gamma$ is generated topologically by one element, $\Lambda$ behaves like a two-dimensional regular local power-series ring. The structure theorem for finitely generated torsion $\Lambda$-modules describes $I$ through characteristic factors. The resulting invariants, traditionally called $\mu$, $\lambda$, and $\nu$, govern the size of class groups at high levels. In a common form,
\[
v_p(|I_n|)=\mu p^n+\lambda n+\nu
\]
for all sufficiently large $n$, with conventions depending on the tower.

This formula does not say the class group itself is simple. It says its $p$-adic size eventually follows a predictable pattern. The infinite object is useful because its module structure explains all finite stages simultaneously.

\paragraph{p-adic analysis and the main conjecture}
Kubota and Leopoldt constructed p-adic L-functions by interpolating special values of Dirichlet L-functions and Bernoulli-number data. Iwasawa's main conjecture says, roughly, that the characteristic power series of an arithmetic $\Lambda$-module agrees with the p-adic L-function, up to the natural notion of equivalence.

Mazur and Wiles proved the main conjecture for $\mathbb Q$ and totally real fields in the 1980s and 1990. Ribet's converse to Herbrand's theorem and Euler systems supplied influential methods. The conjecture connects algebraic growth of class groups with analytic values of L-functions \cite{iwasawa_theory_ref3}.

\paragraph{Generalizations}
The Galois group may be a higher-dimensional p-adic Lie group rather than $\mathbb Z_p$. Barry Mazur developed Iwasawa theory for abelian varieties. Ralph Greenberg proposed theories for motives. Elliptic-curve Iwasawa theory studies Selmer groups, ranks, and p-adic L-functions; the main conjecture compares a Selmer module with an analytic function.

\paragraph{What the theory depends on and where it is useful}
Iwasawa theory depends on algebraic number theory, Galois theory, class field theory, p-adic analysis, modules, and L-functions. It helps the study of Fermat-type questions, elliptic curves, special values, and the arithmetic of infinite extensions.

The statements depend strongly on the chosen prime, tower, and module. A main conjecture may be proved in one family but open in another, and the invariants describe asymptotic growth rather than every small finite layer.

\section*{4. Important theorems and results}
\begin{enumerate}[leftmargin=*,itemsep=0.35em]
\item \textbf{Iwasawa structure theorem.} Finitely generated torsion modules over the one-variable Iwasawa algebra admit a controlled decomposition up to finite modules.

\item \textbf{Class-number growth formula.} In a $\mathbb Z_p$-tower, p-parts of suitable class groups eventually have valuation $\mu p^n+\lambda n+\nu$.

\item \textbf{Mazur--Wiles main conjecture.} In major cyclotomic cases, the algebraic characteristic element equals the analytic p-adic L-function.

\item \textbf{Herbrand--Ribet theorem.} Divisibility of Bernoulli-number data is related to nontrivial components of class groups.

\item \textbf{Ferrero--Washington theorem.} The $\mu$-invariant vanishes for abelian number fields in the cyclotomic setting \cite{iwasawa_theory_ref4}.
\end{enumerate}

\theoryapplications
\theoryconnections{iwasawa theory}{%
\item Algebraic number theory supplies class groups and units, p-adic analysis supplies local arithmetic, and commutative algebra supplies modules over the Iwasawa algebra.
\item Galois theory organizes the infinite towers whose arithmetic is compared.
}{%
\item Iwasawa theory helps explain the growth of class groups and Selmer groups in infinite extensions.
\item Its main conjectures connect algebraic modules with p-adic L-functions, linking arithmetic calculations to analytic information.
}
\section*{Glossary}
\begin{description}[leftmargin=*,style=nextline,itemsep=0.3em]
\item[\textbf{$\mathbb{Z}_p$-extension.}] An infinite Galois extension whose Galois group is the additive group of $p$-adic integers, assembled from finite cyclic layers of increasing order.
\item[\textbf{Iwasawa algebra.}] The power series ring $\Lambda=\mathbb{Z}_p[[\Gamma]]$ over the $p$-adic integers in a topological generator of the Galois group $\Gamma$; it acts on inverse limits of class groups.
\item[\textbf{Class group.}] The group of fractional ideals modulo principal ideals in a number field; it measures the failure of unique factorization into elements.
\item[\textbf{$p$-adic $L$-function.}] An analytic function over the $p$-adic numbers that interpolates special values of ordinary $L$-functions and, through the main conjecture, mirrors algebraic modules.
\item[\textbf{$\mu$-invariant.}] The coefficient of $p^n$ in the growth formula $v_p(|I_n|)=\mu p^n+\lambda n+\nu$; it measures exponential growth of the class group and vanishes in many cyclotomic cases.
\item[\textbf{$\lambda$-invariant.}] The coefficient of $n$ in the growth formula, measuring polynomial growth of the class group along the tower; it is tied to Bernoulli-number divisibility via Herbrand's theorem.
\item[\textbf{Inverse limit.}] A construction assembling compatible systems of objects into a single limiting object.
\item[\textbf{Cyclotomic field.}] A number field obtained by adjoining roots of unity to $\mathbb{Q}$; the primary setting for classical Iwasawa theory.
\item[\textbf{Bernoulli number.}] Rational numbers appearing in power-series expansions and $L$-function values controlling class-group divisibility.
\item[\textbf{Selmer group.}] An algebraic group measuring arithmetic information in the context of elliptic curves and Galois cohomology.
\item[\textbf{Main conjecture.}] The statement that the characteristic element of an Iwasawa module equals a $p$-adic $L$-function.
\end{description}
\noindent\textbf{Computational number theory and primality testing.} Iwasawa-theoretic invariants provide practical computational tools. The Fermat quotient $q_p(a)=(a^{p-1}-1)/p\!\pmod{p}$, which is a building block of p-adic L-functions, can be computed efficiently and used to detect Wieferich primes, primes $p$ satisfying $2^{p-1}\equiv 1\pmod{p^2}$, a condition relevant to the first case of Fermat's Last Theorem. The $\lambda$-invariant of the cyclotomic $\mathbb Z_p$-extension of $\mathbb Q$ is tied to Bernoulli-number divisibility via Herbrand's theorem and the Ferrero--Washington theorem: $p$ divides the numerator of a Bernoulli number $B_k$ if and only if $p$ divides the order of a specific component of the class group of $\mathbb Q(\zeta_p)$. Algorithms that compute these class-group components (e.g., via the Sagase or Pari/GP software packages) rely on Iwasawa-theoretic structure to organize computations in towers of number fields, making what would be astronomically large explicit-degree computations feasible by working inside the inverse limit.
\section*{7. Sample Questions}
\emph{Exercise reference:} \cite{iwasawa_theory_ref1}. The cited edition is the source for the exercise set; consult its Exercises section for the official statement and numbering.
\begin{enumerate}[leftmargin=*,itemsep=0.9em]
\item \textbf{Exercise 1.} The class number of $\mathbb{Q}$ is $1$. If $\mu=0$, $\lambda=2$, and $\nu=1$ for a $\mathbb{Z}_5$-extension tower, use the growth formula $v_p(|I_n|)=\mu\cdot 5^n+\lambda n+\nu$ to compute the $5$-part of the class group size at levels $n=0,1,2,3$.

\emph{Solution.} With $p=5$, $\mu=0$, $\lambda=2$, $\nu=1$: $v_5(|I_n|)=2n+1$. At $n=0$: $v_5=1$, so $|I_0|=5$. At $n=1$: $v_5=3$, so $|I_1|=5^3=125$. At $n=2$: $v_5=5$, so $|I_2|=5^5=3125$. At $n=3$: $v_5=7$, so $|I_3|=5^7=78125$. The class group grows linearly (in the exponent) along the tower.

\item \textbf{Exercise 2.} Verify that $\mathbb{Z}/5^n\mathbb{Z}$ has order $5^n$ and that the inverse limit $\varprojlim \mathbb{Z}/5^n\mathbb{Z}$ is isomorphic to $\mathbb{Z}_5$.

\emph{Solution.} The natural projection maps $\mathbb{Z}/5^{n+1}\mathbb{Z}\to\mathbb{Z}/5^n\mathbb{Z}$ are surjective and compatible. An element of the inverse limit is a sequence $(a_0,a_1,a_2,\ldots)$ with $a_n\in\mathbb{Z}/5^n\mathbb{Z}$ and $a_{n+1}\equiv a_n\pmod{5^n}$. Each such sequence corresponds to a unique $5$-adic integer $\alpha=\lim a_n\in\mathbb{Z}_5$, giving $\varprojlim\mathbb{Z}/5^n\mathbb{Z}\cong\mathbb{Z}_5$.

\item \textbf{Exercise 3.} Let $p=3$. If the Ferrero--Washington theorem guarantees $\mu=0$ for the cyclotomic $\mathbb{Z}_3$-extension of $\mathbb{Q}(\zeta_3)$, and $\lambda=4$ and $\nu=0$, compute $v_3(|I_n|)$ for $n=0,1,2$.

\emph{Solution.} $v_3(|I_n|)=4n$ since $\mu=0$ and $\nu=0$. At $n=0$: $v_3=0$, so $|I_0|=1$ (trivial class group at the base level). At $n=1$: $v_3=4$, so $|I_1|=3^4=81$. At $n=2$: $v_3=8$, so $|I_2|=3^8=6561$.

\item \textbf{Exercise 4.} Explain why $\mathbb{Z}_p[[\Gamma]]$ with $\Gamma\cong\mathbb{Z}_p$ is isomorphic to the formal power series ring $\mathbb{Z}_p[[T]]$.

\emph{Solution.} Choose a topological generator $\gamma$ of $\Gamma\cong\mathbb{Z}_p$ (e.g., $\gamma$ corresponds to $1\in\mathbb{Z}_p$). Any element of $\Gamma$ is $\gamma^s$ for some $s\in\mathbb{Z}_p$. Map $\gamma\mapsto 1+T$. Then $\gamma^s\mapsto(1+T)^s=\sum_{k=0}^{\infty}\binom{s}{k}T^k$, which is a well-defined power series in $\mathbb{Z}_p[[T]]$ since $\binom{s}{k}\in\mathbb{Z}_p$. This gives a continuous ring isomorphism $\mathbb{Z}_p[[\Gamma]]\cong\mathbb{Z}_p[[T]]$.

\item \textbf{Exercise 5.} Suppose for a $\mathbb{Z}_3$-extension tower the invariants are $\mu=0$, $\lambda=3$, $\nu=2$. Use the growth formula to determine at which level $n$ the $3$-part of the class group first exceeds $3^{10}$.

\emph{Solution.} The growth formula gives $v_3(|I_n|)=0\cdot 3^n+3n+2=3n+2$. We need $3n+2>10$, i.e., $3n>8$, i.e., $n\geq 3$. Checking: $n=2$ gives $v_3=8$ ($|I_2|=3^8=6561<3^{10}=59049$), and $n=3$ gives $v_3=11$ ($|I_3|=3^{11}=177147>3^{10}$). So level $n=3$ is the first level where the $3$-part exceeds $3^{10}$.

\end{enumerate}
\section*{8. References}
\addcontentsline{toc}{section}{References}

\begin{thebibliography}{9}
\bibitem{iwasawa_theory_ref1} Kenkichi Iwasawa, "On Gamma-extensions of algebraic number fields," \emph{Bulletin of the AMS}, 65 (1959), 183--226.
\bibitem{iwasawa_theory_ref2} Lawrence C. Washington, \emph{Introduction to Cyclotomic Fields}, 2nd ed., Springer, 1997.
\bibitem{iwasawa_theory_ref3} John Coates and R. Sujatha, \emph{Cyclotomic Fields and Zeta Values}, Springer, 2006.
\bibitem{iwasawa_theory_ref4} Ralph Greenberg, "Iwasawa theory---past and present," in \emph{Class Field Theory---Its Centenary and Prospect}, 2001.
\end{thebibliography}


\theory{K-theory}{How can vector bundles, projective modules, and large matrices be summarized by algebraic invariants?}{K-theory turns direct-sum objects into groups by formally allowing differences. Its functors assign rings or groups to spaces and schemes, preserving structural information that ordinary homology may miss.}{Topology, algebraic geometry, operator algebras, index theory, physics, and the classification of topological materials.}{Vector bundles and projective modules are replaced by stable classes, allowing geometric classification to be expressed through algebraic groups.}

\paragraph{The problem and history}
Vector bundles over a space can be added by direct sum, but they cannot usually be subtracted. Grothendieck's idea was to enlarge the collection formally so that a difference $[E]-[F]$ becomes meaningful. The resulting group is K-theory. It was developed by Grothendieck in algebraic geometry and by Atiyah and Hirzebruch in topology \cite{k_theory_ref1,k_theory_ref2}.

The word K comes from the German word for a vector bundle. K-theory is a generalized cohomology theory. It is also algebraic: higher K-groups study projective modules, matrices, and algebraic relations in rings.

\end{historylist}
\section*{3. Theory}
\subsection*{Core theory}
\paragraph{Topological K-theory}
For a compact Hausdorff space $X$, isomorphism classes of finite-dimensional vector bundles form an abelian monoid under direct sum. Its Grothendieck completion is $K^0(X)$. A formal difference $[E]-[F]$ is a virtual vector bundle. Tensor product makes $K^0(X)$ a ring.

The Serre--Swan theorem identifies vector bundles over a compact space with finitely generated projective modules over its ring of continuous functions. Thus a topological question about bundles can become an algebraic question about modules. Reduced K-theory removes the contribution of a base point.

\paragraph{Worked example: vector bundles on $S^1$ and $S^2$}
\textbf{The circle $S^1$.} A complex line bundle over $S^1$ is classified by a transition function on the overlap of a covering by two intervals. Because $\pi_0(GL_1(\mathbb{C}))=0$, every complex line bundle on $S^1$ is trivial, so $K^0(S^1)\cong\mathbb{Z}\oplus\mathbb{Z}$: the rank and the first Chern class $c_1\in H^2(S^1;\mathbb{Z})=0$ carry no additional information beyond rank.

\textbf{The sphere $S^2$.} The situation is richer. The Hopf bundle $E_{\mathrm{Hopf}}\to S^2$ is a complex line bundle with $c_1(E_{\mathrm{Hopf}})=1$, and its class generates $\widetilde{K}^0(S^2)\cong\mathbb{Z}$. Every vector bundle on $S^2$ is a direct sum of copies of the trivial bundle and the Hopf bundle: $E\cong\underline{\mathbb{C}^k}\oplus E_{\mathrm{Hopf}}^{\oplus m}$, where $k$ is the rank minus $m$. The reduced K-group $\widetilde{K}^0(S^2)\cong\mathbb{Z}$ is detected by the first Chern class of the Hopf bundle. This illustrates how K-theory encodes the nontrivial twisting that ordinary dimension-counting misses.

\paragraph{Algebraic K-theory}
For a ring or scheme, $K_0$ is generated by finitely generated projective modules or coherent sheaves, with relations from exact sequences. Higher K-groups $K_1,K_2,\ldots$ capture invertible matrices, elementary relations, and deeper algebraic structure. Algebraic K-theory has localization sequences, products, and descent results.

In algebraic geometry, $K_0$ of a smooth projective curve reflects rank and degree. For singular spaces, coherent sheaves and derived categories describe the corrections caused by singularities. Projective bundles, Artinian algebras, isolated quotient singularities, and virtual tangent bundles provide calculable examples \cite{k_theory_ref3}.

\paragraph{Chern characters and virtual geometry}
The Chern character maps K-theory to rational cohomology:
\[
\operatorname{ch}:K^0(X)\longrightarrow H^{\mathrm{even}}(X;\mathbb Q).
\]
It translates direct sums into addition and tensor products into multiplication. Chern classes recover geometric information from a bundle; the Chern character is often easier to use in index formulas.

Virtual bundles describe formal tangent directions when an intersection or moduli space is singular or has an unexpected dimension. For an intersection $Z=Y_1\cap Y_2$ in a smooth ambient space $X$, one writes
\[
[T_Z]^{\mathrm{vir}}=[T_{Y_1}]|_Z+[T_{Y_2}]|_Z-[T_X]|_Z.
\]
This is a bookkeeping device that makes Riemann--Roch and enumerative geometry work in imperfect situations.

\paragraph{Bott periodicity and index theory}
Topological K-theory is periodic: iterating suspension by two dimensions returns the same complex K-groups. Bott periodicity is one of the deepest reasons K-theory is computable. The Atiyah--Singer index theorem expresses the index of an elliptic operator through a K-theory class of its symbol.

Grothendieck--Riemann--Roch compares push-forward in K-theory with push-forward in cohomology after applying the Chern character and Todd class. Adams operations provide additional operations that detect how a bundle behaves under tensor powers.

\paragraph{Equivariant and physical applications}
Equivariant K-theory keeps track of a group action as well as a bundle. It is useful for quotients, orbifolds, representation theory, and fixed-point formulas. In string theory, twisted K-theory has been proposed to classify D-brane charges and Ramond--Ramond fields. In condensed-matter physics, K-theory classifies topological insulators, superconductors, and stable Fermi surfaces \cite{k_theory_ref4}.

\paragraph{What the theory depends on and where it is useful}
K-theory depends on vector bundles, modules, homotopy, cohomology, category theory, and algebraic geometry. It helps index theory convert analytic problems to topology, intersection theory handle virtual dimensions, and operator algebras classify projections.

K-theory is a refined invariant, not a complete description of a space or ring. Different objects can have the same K-groups, and torsion or higher operations may be required to distinguish them.

\section*{4. Important theorems and results}
\begin{enumerate}[leftmargin=*,itemsep=0.35em]
\item \textbf{Grothendieck group construction.} Every commutative monoid under direct sum has a universal abelian group of formal differences.

\item \textbf{Serre--Swan theorem.} Vector bundles over a compact Hausdorff space correspond to finitely generated projective modules over continuous functions.

\item \textbf{Bott periodicity.} Complex K-theory is periodic with period two; real K-theory has period eight.

\item \textbf{Grothendieck--Riemann--Roch theorem.} The Chern character translates K-theoretic push-forward into a cohomological formula involving the Todd class.

\item \textbf{Atiyah--Singer theorem.} The analytic index is the image of a K-theory symbol under a topological index map \cite{k_theory_ref2}.
\end{enumerate}

\theoryapplications
\theoryconnections{k theory}{%
\item Algebra supplies projective modules and exact sequences, topology supplies vector bundles and homotopy, and category theory supplies Grothendieck groups and higher constructions.
\item Geometry connects bundles, characteristic classes, and index problems.
}{%
\item K-theory helps classify vector bundles, formulate index theorems, study algebraic cycles, and organize information about modules.
\item It packages many stable equivalence questions into groups and spectra that support calculation.
\item \textbf{Topological insulators in condensed-matter physics.} Topological insulators are materials that are insulating in their interior but support conducting states on their boundary. The classification of free-fermion topological phases in condensed matter is given by K-theory of the Brillouin zone. For systems with time-reversal symmetry (class AII), the ten-fold way periodicity of real K-theory (period eight by Bott periodicity) predicts exactly which insulating phases exist in each spatial dimension. The $\mathbb{Z}_2$ invariant distinguishing a topological insulator from a trivial insulator in two dimensions is a K-theory class.
\item \textbf{$C^*$-algebra classification.} In operator algebras, the K-theory groups $K_0$ and $K_1$ of a $C^*$-algebra classify projections and unitaries up to stable equivalence. The bootstrap class in $C^*$-algebra theory uses K-theoretic data to classify simple nuclear $C^*$-algebras by their K-groups, trace, and the pairing between them. This KK-theory framework, which combines homotopy and K-theory in the noncommutative setting, provides invariants for spaces that cannot be studied with ordinary commutative methods.
}

\section*{Glossary}
\begin{description}[leftmargin=*,style=nextline,itemsep=0.3em]
\item[\textbf{Vector bundle.}] A space that locally looks like a product of a base space with a vector space, where the fibers over nearby points are identified by linear maps called transition functions.
\item[\textbf{Projective module.}] A module over a ring that is a direct summand of a free module; over the ring of continuous functions on a compact space, these correspond exactly to vector bundles by the Serre--Swan theorem.
\item[\textbf{Virtual bundle.}] A formal difference $[E]-[F]$ of two vector bundles in K-theory, representing the difference of their classes even when no single honest bundle has that difference as a class.
\item[\textbf{Bott periodicity.}] A deep repetition pattern in topological K-theory: after suspending by two (complex) or eight (real) dimensions, the K-groups return to their original values.
\item[\textbf{Chern character.}] A natural transformation from K-theory to rational cohomology that converts tensor products into products and makes characteristic-class calculations more tractable.
\item[\textbf{Equivariant K-theory.}] A version of K-theory that remembers not only the bundle but also a group action on it, used for studying orbifolds, quotient spaces, and symmetries.
\item[\textbf{Serre--Swan theorem.}] The identification of vector bundles over a compact space with finitely generated projective modules over its ring of continuous functions.
\item[\textbf{Grothendieck group.}] The universal abelian group constructed from a commutative monoid by formally adjoining differences.
\item[\textbf{Reduced K-theory.}] The K-theory of a pointed space after factoring out the contribution of the base point.
\item[\textbf{Chern class.}] A characteristic class assigning cohomology classes to vector bundles; the first detects nontrivial line bundles.
\item[\textbf{Hopf bundle.}] The nontrivial complex line bundle over $S^2$ whose first Chern class generates reduced $K^0(S^2)$.
\end{description}

\section*{7. Sample Questions}
\emph{Exercise reference:} \cite{k_theory_ref1}. The cited edition is the source for the exercise set; consult its Exercises section for the official statement and numbering.
\begin{enumerate}[leftmargin=*,itemsep=0.9em]
\item \textbf{Exercise 1.} Compute $\widetilde{K}^0(S^0)$, where $S^0=\{+1,-1\}$ is the discrete two-point space.

\emph{Solution.} Every complex vector bundle over a discrete space is trivial, so $K^0(S^0)\cong\mathbb{Z}\oplus\mathbb{Z}$ (one rank summand per point). The reduced group is the kernel of the rank map to $\mathbb{Z}$, so $\widetilde{K}^0(S^0)\cong\mathbb{Z}$.

\item \textbf{Exercise 2.} Let $E_{\mathrm{Hopf}}$ be the Hopf bundle over $S^2$. Compute the class of $E_{\mathrm{Hopf}}\otimes E_{\mathrm{Hopf}}$ in $\widetilde{K}^0(S^2)$.

\emph{Solution.} The first Chern class satisfies $c_1(E\otimes F)=c_1(E)+c_1(F)$ for line bundles. Thus $c_1(E_{\mathrm{Hopf}}\otimes E_{\mathrm{Hopf}})=2\in H^2(S^2;\mathbb{Z})\cong\mathbb{Z}$. Since $\widetilde{K}^0(S^2)\cong\mathbb{Z}$ is detected by $c_1$, the class is $2$. Equivalently, $E_{\mathrm{Hopf}}^{\otimes 2}\cong\mathcal{O}(2)$, which is the line bundle with Chern number~$2$.

\item \textbf{Exercise 3.} Let $X=S^2$ and let $\alpha=[E_{\mathrm{Hopf}}]-[\underline{\mathbb{C}}]\in\widetilde{K}^0(S^2)$. Compute $\operatorname{ch}(\alpha)\in H^{\mathrm{even}}(S^2;\mathbb{Q})$ and evaluate it on the fundamental class $[S^2]$.

\emph{Solution.} The Chern character is additive: $\operatorname{ch}([E_{\mathrm{Hopf}}]-[\underline{\mathbb{C}}])=e^{c_1(E_{\mathrm{Hopf}})}-1=c_1(E_{\mathrm{Hopf}})=1\in H^2(S^2;\mathbb{Q})$. Evaluated on $[S^2]$, the result is $\langle c_1(E_{\mathrm{Hopf}}),[S^2]\rangle=1$.

\item \textbf{Exercise 4.} By Bott periodicity, $\widetilde{K}^0(S^2)\cong\mathbb{Z}$. Use the suspension isomorphism twice to determine $K^0(S^4)$.

\emph{Solution.} The suspension isomorphism gives $\widetilde{K}^0(\Sigma X)\cong \widetilde{K}^0(X)$ for complex K-theory. Applying it twice: $\widetilde{K}^0(S^4)\cong\widetilde{K}^0(S^2)\cong\mathbb{Z}$. Hence $K^0(S^4)\cong\mathbb{Z}\oplus\mathbb{Z}$, generated by the trivial bundle and the tautological quaternionic line bundle over $S^4\cong\mathbb{HP}^1$.

\item \textbf{Exercise 5.} For $X=S^2$, write out the isomorphism $K^0(S^2)\cong\mathbb{Z}\oplus\mathbb{Z}$ explicitly by identifying the two invariants that label each class.

\emph{Solution.} Every vector bundle on $S^2$ is stably equivalent to $\underline{\mathbb{C}}^a\oplus E_{\mathrm{Hopf}}^b$ for unique integers $a\geq 0$ and $b\in\mathbb{Z}$. The map $[E]\mapsto(\operatorname{rank}(E),\,c_1(E))$ gives the isomorphism $K^0(S^2)\cong\mathbb{Z}\oplus\mathbb{Z}$.

\end{enumerate}
\section*{8. References}
\addcontentsline{toc}{section}{References}

\begin{thebibliography}{9}
\bibitem{k_theory_ref1} Michael Atiyah, \emph{K-Theory}, 2nd ed., Westview Press, 1989.
\bibitem{k_theory_ref2} Max Karoubi, \emph{K-Theory: An Introduction}, Springer, 1978.
\bibitem{k_theory_ref3} Charles Weibel, \emph{The K-book: An Introduction to Algebraic K-theory}, Graduate Studies in Mathematics 145, AMS, 2013.
\bibitem{k_theory_ref4} Jonathan Rosenberg, \emph{Algebraic K-Theory and Its Applications}, Springer, 1994.
\end{thebibliography}

\theory{Knot theory}{When are two tangled loops really the same knot, and how can the difference be detected without untangling them physically?}{Knot theory studies embeddings of circles in three-dimensional space up to ambient deformation. Diagrams, Reidemeister moves, knot groups, polynomials, and geometric invariants turn visual tangling into calculable mathematics.}{Topology, DNA and proteins, polymers, quantum topology, statistical mechanics, and topological quantum computation.}{Isotopy classes, knot complements, and invariants distinguish embeddings of circles in three-dimensional space.}

\paragraph{Definition and history}
A mathematical knot is a copy of a circle embedded in $\mathbb R^3$; its ends are joined, unlike a shoelace. Two knots are equivalent when one can be carried to the other by a continuous family of homeomorphisms of space, with no cutting or passing through itself. The untwisted circle is the unknot; the trefoil is the simplest nontrivial knot \cite{knot_theory_ref1}.

Vandermonde and Gauss studied early topological features, including the linking integral. Lord Kelvin's nineteenth-century idea that atoms might be knotted vortices motivated Peter Tait's knot tables. Dehn, Alexander, Reidemeister, Thurston, Jones, Witten, Kontsevich, and Kauffman later introduced algebraic, geometric, and quantum methods.

\end{historylist}
\section*{3. Theory}
\subsection*{Core theory}
\paragraph{Diagrams and Reidemeister moves}
\bookfigure{knot_braid}{Knot and braid diagrams turn strand motion into a mathematical object.}
Project a knot to the plane so that crossings are transverse double points. At each crossing, record which strand passes over the other. Different projections of the same knot look different, so one needs diagram moves that preserve the knot.

Reidemeister's three moves are: create or remove a twist; pull one strand across another; and move a strand across a crossing. Two diagrams represent equivalent knots exactly when a finite sequence of these moves connects them. This gives a finite local language for a global deformation \cite{knot_theory_ref2}.

\paragraph{Invariants}
A knot invariant gives the same answer for every diagram of an equivalent knot. The knot group is the fundamental group of the complement of the knot. The Alexander polynomial comes from the infinite cyclic cover of the complement. The Jones and Kauffman polynomials use skein relations and connect knot theory to quantum groups and statistical mechanics.

For the Alexander--Conway polynomial of an oriented link, the unknot has value $1$ and a crossing relation has the form $C(L_+)=C(L_-)+zC(L_0)$. Applying it to the trefoil gives $1+z^2$, different from the unknot, so the trefoil is knotted. The polynomial does not distinguish the left and right trefoils, but the Jones polynomial can distinguish their chirality.

\paragraph{Hyperbolic geometry and higher dimensions}
Most knots are hyperbolic: their complements admit a complete finite-volume hyperbolic metric. Volume, cusp shape, and other hyperbolic data become invariants. Thurston's hyperbolization work connected knot complements to three-dimensional geometry.

A higher-dimensional knot is an $n$-sphere embedded in $(n+2)$-dimensional Euclidean space. Links have several components. Tangles have open ends and model local pieces of knots. Knot multiplication joins two knots; prime knots cannot be decomposed nontrivially, and tables use crossing number to organize examples.

\paragraph{Tabulation and computation}
The Alexander--Briggs notation records crossing number and index. Dowker--Thistlethwaite codes, Conway notation, and Gauss codes encode diagrams in machine-readable form. More than billions of knots and links have been tabulated. Knot recognition has algorithms, including unknot-recognition algorithms, but their practical complexity remains a major question.

\paragraph{Applications}
DNA and proteins can become knotted or linked. Tangles model the action of topoisomerase enzymes, and knot invariants can distinguish molecular handedness. Polymer physics studies how knots affect elasticity and entanglement. Jones-polynomial ideas inspired topological quantum computation, where protected information is stored in global braiding patterns \cite{knot_theory_ref3}.

\paragraph{What the theory depends on and where it is useful}
Knot theory depends on topology, group theory, homology, geometry, and combinatorics. It helps low-dimensional topology through knot complements and three-manifolds, quantum topology through polynomial invariants, and biology through molecular tangles.

No single invariant distinguishes every knot. A successful test may prove two knots different but usually cannot prove they are the same without a complete equivalence algorithm.

\section*{4. Important theorems and results}
\begin{enumerate}[leftmargin=*,itemsep=0.35em]
\item \textbf{Reidemeister theorem.} Equivalent knot diagrams are related by the three Reidemeister moves.

\item \textbf{Alexander polynomial.} The Alexander invariant of a knot complement produces a computable polynomial invariant.

\item \textbf{Gordon--Luecke theorem.} A knot in the three-sphere is determined by the homeomorphism type of its complement.

\item \textbf{Thurston hyperbolization.} Many knot complements carry a finite-volume hyperbolic structure.

\item \textbf{Jones polynomial.} A quantum invariant obtained from a skein relation and capable of detecting chirality in examples \cite{knot_theory_ref4}.

The theorems answer three different questions. Reidemeister moves identify when two drawings represent the same knot; polynomial invariants provide computable tests that can distinguish drawings; and the complement theorems show that the surrounding three-dimensional space can determine the knot itself. Hyperbolic geometry adds measurable structure, such as volume, to many knot complements. None of these invariants is a universal quick classifier, so several must often be combined.
\end{enumerate}

\theoryapplications
\theoryconnections{knot theory}{%
\item Topology supplies embeddings and ambient isotopy, group theory supplies knot groups, and algebraic topology supplies homology and covering spaces.
\item Combinatorics and polynomial algebra turn a knot diagram into invariants that can be computed.
}{%
\item Knot theory helps study DNA recombination, polymer entanglement, fluid vortices, and three-dimensional manifolds.
\item A knot invariant is useful when it separates examples or predicts a property, but one invariant rarely gives a complete classification.
}
\section*{7. Sample Questions}
\emph{Exercise reference:} \cite{knot_theory_ref1}. The cited edition is the source for the exercise set; consult its Exercises section for the official statement and numbering.
\begin{enumerate}[leftmargin=*,itemsep=0.9em]
\item \textbf{Exercise 1.} Why is the unknot different from a knot tied in a rope?

\emph{Solution.} A mathematical knot is closed: its two ends are joined. A rope can be untied by sliding an end out, while a closed knot cannot use that escape.

\item \textbf{Exercise 2.} What do Reidemeister moves preserve?

\emph{Solution.} They preserve the ambient-isotopy class of the knot, even though they change the planar drawing.

\item \textbf{Exercise 3.} Why can an invariant fail to distinguish knots?

\emph{Solution.} An invariant is a summary. Different knots may produce the same summary, so several invariants or a deeper complement analysis may be needed.

\item \textbf{Exercise 4.} What does the knot group describe?

\emph{Solution.} It describes how loops in the three-dimensional space outside the knot can be combined and deformed.

\item \textbf{Exercise 5.} Why are tangles useful for DNA?

\emph{Solution.} A local DNA operation affects a segment with two ends. Tangle diagrams model that local segment and compare the knot or link before and after the operation.

\end{enumerate}
\section*{8. References}
\addcontentsline{toc}{section}{References}

\begin{thebibliography}{9}
\bibitem{knot_theory_ref1} Colin C. Adams, \emph{The Knot Book}, American Mathematical Society, 2004.
\bibitem{knot_theory_ref2} W. B. R. Lickorish, \emph{An Introduction to Knot Theory}, Springer, 1997.
\bibitem{knot_theory_ref3} Dorothy Buck and Erica Flapan, \emph{Topology of Molecular Biology}, World Scientific, 2009.
\bibitem{knot_theory_ref4} Dale Rolfsen, \emph{Knots and Links}, Publish or Perish, 1976.
\end{thebibliography}


\theory{L-theory}{How can quadratic forms and algebraic surgery measure whether a manifold can be modified into a desired shape?}{Algebraic L-theory is a K-theory-like study of quadratic and symmetric forms over rings with involution. Its groups store the obstructions that remain after geometric surgery.}{Surgery theory, manifold classification, quadratic forms, topology, algebraic K-theory, and controlled topology.}{Quadratic forms and surgery obstructions convert high-dimensional manifold classification into algebraic calculations.}

\paragraph{History and motivation}
Wall introduced algebraic L-theory, using the letter after K. The central problem in surgery theory is to alter a manifold by cutting out a product of disks and gluing another product in its place. Algebra tells us whether the altered object can become homotopy equivalent or homeomorphic to a target manifold. The obstruction lies in an L-group \cite{l_theory_ref1}.

\end{historylist}
\section*{3. Theory}
\subsection*{Core theory}
\paragraph{Quadratic forms and Witt groups}
A quadratic form assigns a value such as $x^2+y^2$ or $x^2-3y^2$ to a vector. Over a ring with involution, one studies forms compatible with the involution. A form is nondegenerate when its associated pairing has no hidden null direction. Hyperbolic forms represent completely canceling pairs and are treated as trivial after stabilization.

For a ring $R$ with involution, even-dimensional quadratic groups $L_{2k}(R)$ are Witt groups of suitable $(-1)^k$-quadratic forms. Addition is orthogonal direct sum. Symmetric L-groups $L^n(R)$ use symmetric forms; a symmetrization map relates the two and is an isomorphism away from relevant 2-torsion.

\paragraph{Periodicity and the integer example}
Quadratic and symmetric L-groups are four-periodic. For the integers,
\[
\begin{aligned}
L_{4k}(\mathbb Z)&\cong\mathbb Z, &
L_{4k+1}(\mathbb Z)&=0,\\
L_{4k+2}(\mathbb Z)&\cong\mathbb Z/2, &
L_{4k+3}(\mathbb Z)&=0.
\end{aligned}
\]
The integer in dimension $4k$ is related to signature, while the $\mathbb Z/2$ class in dimension $4k+2$ is detected by an Arf invariant. These values are the algebraic shadows of geometric surgery obstructions \cite{l_theory_ref2}.

\paragraph{Surgery and manifold classification}
Suppose a manifold is equipped with a degree-one normal map to a target. Surgery removes a sphere times a disk and replaces it by a disk times a sphere. Below the middle dimension, surgery can remove homotopy defects. At the middle dimension, an intersection or quadratic form remains. Its class in $L_n(\mathbb Z[\pi_1])$ is the obstruction to completing the surgery.

The fundamental group appears through the group ring $\mathbb Z[\pi]$, because loops affect how handles are attached. For finite $\pi$, algebraic calculations are available; for infinite groups, controlled topology and geometric methods become important.

\paragraph{Algebraic and symmetric versions}
The quadratic theory remembers a preferred quadratic refinement, while symmetric theory remembers the bilinear form. In characteristic different from two these are closely related; in the presence of 2-torsion they can differ. Ranicki developed formulations for additive categories equipped with chain duality, allowing chain complexes and derived algebra to replace free modules.

\paragraph{What the theory depends on and where it is useful}
L-theory depends on ring theory, modules, quadratic forms, K-theory, homotopy theory, and high-dimensional manifold topology. It helps surgery theory decide whether a homotopy equivalence can be improved to a homeomorphism or diffeomorphism, and helps classification problems by turning geometry into algebra.

The theory is most effective in high dimensions and under hypotheses on the fundamental group. Low-dimensional topology has additional phenomena that are not captured by ordinary surgery obstructions.

\section*{4. Important theorems and results}
\begin{enumerate}[leftmargin=*,itemsep=0.35em]
\item \textbf{Wall surgery obstruction.} A degree-one normal map has an obstruction in an appropriate L-group; vanishing is the algebraic condition for surgery continuation.

\item \textbf{Four-periodicity.} L-groups repeat after four dimensions in standard quadratic and symmetric theories.

\item \textbf{Symmetrization.} Quadratic and symmetric L-groups are related by a map that is an isomorphism modulo 2-torsion.

\item \textbf{Surgery exact sequence.} Structure sets, normal invariants, and L-groups occur in an exact sequence organizing manifold classification.

\item \textbf{Signature and Arf calculations.} For $\mathbb Z$, dimensions divisible by four are measured by signature, while dimensions congruent to two carry an Arf-type $\mathbb Z/2$ obstruction \cite{l_theory_ref1}.
\end{enumerate}

\theoryapplications
\theoryconnections{l theory}{%
\item Algebra supplies quadratic forms and modules, topology supplies manifolds and surgery, and homotopy theory supplies normal invariants.
\item The algebraic L-groups record the obstruction left after geometric changes have been made.
}{%
\item L-theory helps classify high-dimensional manifolds and analyze whether a homotopy equivalence can be improved to a homeomorphism or diffeomorphism.
\item Its periodicity turns an apparently unbounded classification problem into recurring algebraic cases.
}
\section*{7. Sample Questions}
\emph{Exercise reference:} \cite{l_theory_ref1}. The cited edition is the source for the exercise set; consult its Exercises section for the official statement and numbering.
\begin{enumerate}[leftmargin=*,itemsep=0.9em]
\item \textbf{Exercise 1.} What does an L-group classify?

\emph{Solution.} It classifies stable equivalence classes of suitable nondegenerate quadratic or symmetric forms, and in topology it records surgery obstructions.

\item \textbf{Exercise 2.} Why are hyperbolic forms treated as trivial?

\emph{Solution.} They contain paired directions that cancel algebraically. Stabilizing by adding them should not change the essential obstruction.

\item \textbf{Exercise 3.} Why does the group ring $\mathbb Z[\pi]$ appear?

\emph{Solution.} The fundamental group $\pi$ records loops in the manifold, and the group ring lets algebra track how those loops act on chains and handles.

\item \textbf{Exercise 4.} What is the significance of four-periodicity?

\emph{Solution.} It means the obstruction groups repeat in dimensions four apart, making high-dimensional calculations patterned rather than unrelated.

\item \textbf{Exercise 5.} Why can quadratic and symmetric L-theory differ?

\emph{Solution.} A quadratic refinement contains information not determined by a bilinear pairing, especially when two-torsion is present.

\end{enumerate}
\section*{8. References}
\addcontentsline{toc}{section}{References}

\begin{thebibliography}{9}
\bibitem{l_theory_ref1} C. T. C. Wall, \emph{Surgery on Compact Manifolds}, 2nd ed., American Mathematical Society, 1999.
\bibitem{l_theory_ref2} Andrew Ranicki, \emph{Algebraic L-Theory and Topological Manifolds}, Cambridge University Press, 1992.
\end{thebibliography}


\theory{Lattice theory}{How can a collection of partially ordered objects support both a best common lower bound and a best common upper bound?}{Lattice theory studies ordered systems in which every pair has a meet and a join. It unifies divisibility, subsets, logic, subspaces, ideals, and data dependencies.}{Order theory, algebra, logic, computer science, optimization, formal verification, and universal algebra.}{Meet, join, order, and fixed-point operations describe how partial information can be combined or completed.}

\paragraph{Richard Dedekind}
Dedekind recognized that substructures such as ideals can be organized by inclusion, with intersection and generated sum acting as meet and join. This made order a useful algebraic description of divisibility and subobjects.

\paragraph{Garrett Birkhoff and Alfred Tarski}
Birkhoff developed lattice theory as an independent algebraic subject and proved representation results for important classes such as distributive lattices. Tarski's fixed-point work showed how complete lattices support existence and iteration arguments in logic and computer science \cite{lattice_theory_ref1}.
\end{historylist}
\section*{3. Theory}
\subsection*{Core theory}
\paragraph{The problem and the definition}
In a partially ordered set, two objects may be comparable or incomparable. A lattice adds two operations. The meet $a\wedge b$ is the greatest object below both of them; the join $a\vee b$ is the least object above both. For sets, meet is intersection and join is union. For integers ordered by divisibility, meet is greatest common divisor and join is least common multiple \cite{lattice_theory_ref1}.

The operations satisfy associativity, commutativity, absorption, and idempotence. A lattice may be described by its order or by these two operations. A bounded lattice has a least element $0$ and greatest element $1$. A complemented lattice gives each element a partner whose meet is $0$ and join is $1$.

\paragraph{Examples}
The power set of a finite set is a Boolean lattice. The subspaces of a vector space form a lattice: meet is intersection and join is vector-space sum. The ideals of a ring form a lattice under intersection and generated sum. Logical propositions form a lattice in which meet means "and" and join means "or".

The divisors of $60$ ordered by divisibility form a finite lattice. The meet of $12$ and $18$ is $6$, while their join is $36$. A total order, such as the real numbers, is also a lattice: meet is minimum and join is maximum.

\paragraph{Distributive and modular lattices}
A lattice is distributive when $a\wedge(b\vee c)=(a\wedge b)\vee(a\wedge c)$ and the dual identity also holds. Boolean lattices are distributive and support ordinary logic. The lattice of subspaces is usually not distributive but is modular, a weaker identity reflecting dimension.

Finite distributive lattices are controlled by their join-irreducible elements. This gives a bridge between an abstract order and a simpler poset of building blocks. Modular and semimodular lattices appear in geometry, matroid theory, and dimension arguments.

\paragraph{Completeness and fixed points}
A complete lattice has a meet and join for every subset, not merely every pair. The power set of any set is complete. Complete lattices are important in semantics and program verification because recursive definitions can be interpreted as fixed points.

The Knaster--Tarski theorem says every monotone self-map of a complete lattice has a least and a greatest fixed point. The least fixed point can describe the smallest set of states reachable by a program; the greatest can describe all states satisfying an invariant \cite{lattice_theory_ref2}.

\paragraph{Homomorphisms and constructions}
A lattice homomorphism preserves meets and joins: $f(a\wedge b)=f(a)\wedge f(b)$ and $f(a\vee b)=f(a)\vee f(b)$. It is automatically order-preserving. Products, sublattices, quotient lattices, ideals, filters, and congruences are standard constructions. A filter is dual to an ideal: it is upward closed and closed under meet.

\paragraph{Applications}
Boolean lattices model digital circuits and logical reasoning. In databases, sets of attributes and functional dependencies form closure systems. In version control, a lattice can represent states that have a common merge or common ancestor. In static program analysis, abstract states are ordered by information and fixed points compute loop invariants.

Lattice methods also organize equivalence relations, partitions, subgroups, ideals, and topology. The lattice of open sets is a frame or locale, a foundation for point-free topology and topos theory. Dilworth's and Mirsky's theorems connect chains and antichains, while Zorn's lemma and maximal principles use order structure \cite{lattice_theory_ref3}.

\paragraph{What the theory depends on and where it is useful}
Lattice theory depends on partial orders, sets, logic, algebra, and universal algebra. It helps order theory describe closure, ring theory organize ideals, type theory interpret propositions, and computer science compute fixed points.

The word lattice also names a periodic grid in geometry and crystallography; that geometric lattice is a different but related subject. In an order lattice, joins and meets may exist abstractly without having a simple numerical formula.

\section*{4. Important theorems and results}
\begin{enumerate}[leftmargin=*,itemsep=0.35em]
\item \textbf{Birkhoff representation theorem.} Every finite distributive lattice is isomorphic to the lattice of lower sets of the poset of its join-irreducible elements.

\item \textbf{Knaster--Tarski theorem.} Every monotone map of a complete lattice has a complete lattice of fixed points, including least and greatest fixed points.

\item \textbf{Dilworth's theorem.} In a finite poset, the largest antichain size equals the minimum number of chains needed to cover the poset.

\item \textbf{Boolean representation theorem.} Every Boolean algebra is represented by sets, with operations corresponding to union, intersection, and complement.

\item \textbf{Jordan--Dedekind property.} In a finite graded lattice, all maximal chains between fixed endpoints have the same length \cite{lattice_theory_ref4}.
\end{enumerate}

\theoryapplications
\theoryconnections{lattice theory}{%
\item Order theory supplies partial orders, joins, meets, and duality; set theory supplies concrete lattices of subsets; and algebra supplies ideals, subspaces, and congruences as examples.
\item Combinatorics studies chains and antichains.
}{%
\item Lattice theory helps logic, computer science, scheduling, data-flow analysis, and the study of subobjects in algebra.
\item The order structure records which information refines which, so fixed-point and representation theorems can guide computation.
}
\section*{7. Sample Questions}
\emph{Exercise reference:} \cite{lattice_theory_ref1}. The cited edition is the source for the exercise set; consult its Exercises section for the official statement and numbering.
\begin{enumerate}[leftmargin=*,itemsep=0.9em]
\item \textbf{Exercise 1.} In the divisor lattice of $60$, find the meet and join of $12$ and $18$.

\emph{Solution.} The meet is $\gcd(12,18)=6$ and the join is $\operatorname{lcm}(12,18)=36$.

\item \textbf{Exercise 2.} What are meet and join for subsets?

\emph{Solution.} The meet is intersection, because it is the largest set contained in both. The join is union, because it is the smallest set containing both.

\item \textbf{Exercise 3.} Why is a total order a lattice?

\emph{Solution.} Any two elements are comparable. The smaller is their meet and the larger is their join.

\item \textbf{Exercise 4.} What does the Knaster--Tarski theorem provide in program verification?

\emph{Solution.} A monotone program transformer has a least fixed point representing the smallest solution of recursive reachability equations and a greatest fixed point representing the largest invariant solution.

\item \textbf{Exercise 5.} What makes a lattice Boolean?

\emph{Solution.} It is bounded, distributive, and every element has a complement. The power set of a set is the standard model.

\end{enumerate}
\section*{8. References}
\addcontentsline{toc}{section}{References}

\begin{thebibliography}{9}
\bibitem{lattice_theory_ref1} Garrett Birkhoff, \emph{Lattice Theory}, 3rd ed., American Mathematical Society, 1967.
\bibitem{lattice_theory_ref2} George Grätzer, \emph{Lattice Theory: Foundation}, Birkhauser, 2011.
\bibitem{lattice_theory_ref3} B. A. Davey and H. A. Priestley, \emph{Introduction to Lattices and Order}, 2nd ed., Cambridge University Press, 2002.
\bibitem{lattice_theory_ref4} Ralph Freese, Jaroslav Jezek, and J. B. Nation, \emph{Free Lattices}, American Mathematical Society, 1995.
\end{thebibliography}


\theory{Lie theory}{How can continuous symmetries be studied with the same clarity that groups give to discrete symmetries?}{Lie theory studies Lie groups, their infinitesimal Lie algebras, and the actions they generate. The exponential map turns small motions into finite transformations and makes differential equations, geometry, and physics share one language.}{Differential equations, geometry, mechanics, relativity, quantum physics, representation theory, and control.}{Lie groups and Lie algebras encode continuous symmetries at both the global and infinitesimal levels.}

\paragraph{History and motivation}
Sophus Lie wanted a theory of continuous transformation groups that would organize differential equations, just as Galois theory organizes polynomial equations. Wilhelm Killing and \'Elie Cartan developed the structure and representation theory of semisimple Lie algebras. Lie's work also led to contact geometry, Lie sphere geometry, moving frames, and symmetric spaces \cite{lie_theory_ref1}.

\end{historylist}
\section*{3. Theory}
\subsection*{Core theory}
\paragraph{Groups of continuous motion}
A Lie group is both a group and a smooth manifold, with smooth multiplication and inversion. Rotations of the plane form the circle group. Rotations of three-dimensional space form $SO(3)$. Translations, rotations, boosts, and spacetime transformations form the Galilean, Lorentz, and Poincare groups.

A one-parameter subgroup is a smooth family $g(t)$ satisfying $g(s+t)=g(s)g(t)$. For ordinary rotations, $g(t)$ rotates through angle $t$. For boosts, hyperbolic functions appear: $\exp(jt)=\cosh t+j\sinh t$ with $j^2=1$. For dual numbers, $\exp(\varepsilon t)=1+\varepsilon t$ with $\varepsilon^2=0$, producing a shear-like motion.

\paragraph{Lie algebras and infinitesimal motion}
The tangent vectors at the identity of a Lie group form its Lie algebra. The Lie bracket measures the first-order difference between doing two infinitesimal motions in opposite orders:
\[
[X,Y]=XY-YX
\]
when the group is represented by matrices. It is bilinear, antisymmetric, and satisfies the Jacobi identity.

The exponential map sends a Lie-algebra element $X$ to a group element $\exp(X)$. Near the identity, the group and algebra contain the same local information. The Baker--Campbell--Hausdorff formula explains how the logarithms of two nearby transformations combine.

\paragraph{Lie's three theorems}
Lie's first theorem constructs a basis of infinitesimal transformations. His second theorem expresses the structure constants of the algebra through commutators. His third theorem says the structure constants are antisymmetric and satisfy the Jacobi identity. The converse results reconstruct a local Lie group from a finite-dimensional Lie algebra under suitable analytic conditions \cite{lie_theory_ref2}.

\paragraph{Examples}
The unit quaternions form a three-dimensional group identified with the sphere $S^3$. Their Lie algebra is the space of imaginary quaternions, and its bracket is twice the vector cross product. The Heisenberg group consists of upper triangular matrices and is the simplest important noncommutative example. Classical matrix groups such as $GL_n$, $SL_n$, $O_n$, and $Sp_n$ provide the main families.

\paragraph{Structure and representations}
Root systems and root data encode the structure of semisimple Lie algebras. Weyl and Coxeter groups describe reflections in roots. Lie groups act on vector spaces through representations, allowing symmetry to be studied with matrices. In physics, representations of rotation and Lorentz groups classify angular momentum and particle types.

Lie theory also includes algebraic groups, groups of Lie type over finite fields, buildings, and the solution of Hilbert's fifth problem, which asks when a locally Euclidean topological group is a Lie group.

\paragraph{Applications}
A symmetry of a differential equation can reduce its variables and produce conserved quantities. In mechanics, rotations and translations explain angular momentum and momentum conservation. In relativity, the Lorentz and Poincare groups organize spacetime. In quantum theory, unitary representations describe states and observables. In robotics and control, $SE(3)$ represents position and orientation, while Lie-group integrators preserve the geometry of motion.

\paragraph{What the theory depends on and where it is useful}
Lie theory depends on groups, manifolds, calculus, linear algebra, differential equations, and representation theory. It helps Galois theory as the continuous analogue of root symmetry, helps geometry through homogeneous spaces, and helps physics by encoding transformation laws.

The Lie algebra captures local behavior but may not determine global topology. Distinct groups can share a Lie algebra while differing by discrete quotients. Global questions therefore require covering groups, fundamental groups, and topology.

\section*{4. Important theorems and results}
\begin{enumerate}[leftmargin=*,itemsep=0.35em]
\item \textbf{Lie correspondence.} Connected, simply connected Lie groups correspond to finite-dimensional real Lie algebras up to isomorphism.

\item \textbf{Exponential map.} One-parameter subgroups are obtained from Lie-algebra elements by exponentiation.

\item \textbf{Lie's three theorems.} Infinitesimal transformations satisfy the structure and Jacobi identities, and conversely generate local groups.

\item \textbf{Cartan classification.} Complex semisimple Lie algebras are classified by root systems of types $A$, $B$, $C$, $D$, and exceptional types.

\item \textbf{Noether principle.} Continuous symmetries of an action produce conserved quantities in suitable physical systems \cite{lie_theory_ref3}.
\end{enumerate}

\theoryapplications
\theoryconnections{lie theory}{%
\item Differential geometry supplies smooth manifolds and tangent vectors, while algebra supplies Lie brackets, linear representations, and root systems.
\item Differential equations describe one-parameter motions, and topology distinguishes global group behavior from local algebra.
}{%
\item Lie theory helps physics describe continuous symmetry, helps differential equations exploit conserved quantities, and helps geometry classify homogeneous spaces.
\item Passing between a Lie group and its Lie algebra replaces nonlinear motion by a local linear-algebraic model.
}

\section*{Glossary}
\begin{description}[leftmargin=*,style=nextline,itemsep=0.3em]
\item[\textbf{Lie group.}] A set that is simultaneously a smooth manifold and a group, with multiplication and inversion operations that are smooth (differentiable).
\item[\textbf{Lie algebra.}] The tangent space at the identity element of a Lie group, equipped with the Lie bracket operation; it captures the local or infinitesimal structure of the group.
\item[\textbf{Lie bracket.}] A binary operation $[X,Y]$ on a Lie algebra that measures the failure of two infinitesimal transformations to commute, analogous to the commutator of matrices.
\item[\textbf{Exponential map.}] The map from a Lie algebra to its Lie group that sends each infinitesimal generator to the corresponding one-parameter subgroup element, linking the algebraic and geometric sides of Lie theory.
\item[\textbf{Root system.}] A finite collection of vectors in Euclidean space that encodes the structure of a semisimple Lie algebra, determining its classification and representation theory.
\item[\textbf{Representation.}] A way of realizing a group or algebra as matrices acting on a vector space, so that the abstract symmetry becomes a concrete linear transformation.
\item[\textbf{One-parameter subgroup.}] A smooth group homomorphism $g\colon\mathbb{R}\to G$; for a Lie group every such subgroup arises from the exponential map.
\item[\textbf{Semisimple Lie algebra.}] A Lie algebra with no nonzero abelian ideals; its structure is determined by root systems.
\item[\textbf{Weyl group.}] The reflection group generated by reflections in the root system; it controls representation theory.
\item[\textbf{Jacobi identity.}] The identity $[X,[Y,Z]]+[Y,[Z,X]]+[Z,[X,Y]]=0$ that, with bilinearity and antisymmetry, defines a Lie bracket.
\item[\textbf{Baker--Campbell--Hausdorff formula.}] Expresses $\log(e^X e^Y)$ as a series in Lie bracket commutators.
\end{description}

\section*{7. Sample Questions}
\emph{Exercise reference:} \cite{lie_theory_ref1}. The cited edition is the source for the exercise set; consult its Exercises section for the official statement and numbering.
\begin{enumerate}[leftmargin=*,itemsep=0.9em]
\item \textbf{Exercise 1.} Compute the Lie bracket $[X,Y]$ for $X=\begin{pmatrix}0&1\\0&0\end{pmatrix}$ and $Y=\begin{pmatrix}0&0\\1&0\end{pmatrix}$ in $\mathfrak{gl}_2(\mathbb{R})$.

\emph{Solution.} $XY=\begin{pmatrix}1&0\\0&0\end{pmatrix}$, $YX=\begin{pmatrix}0&0\\0&1\end{pmatrix}$, so $[X,Y]=XY-YX=\begin{pmatrix}1&0\\0&-1\end{pmatrix}$.

\item \textbf{Exercise 2.} Compute $\exp(X)$ for $X=\begin{pmatrix}0&\theta\\-\theta&0\end{pmatrix}\in\mathfrak{so}(2)$.

\emph{Solution.} $X^2=-\theta^2 I$, $X^3=-\theta^2 X$, and $X^4=\theta^4 I$. Separating even and odd powers: $\exp(X)=\cos\theta\, I+\frac{\sin\theta}{\theta}\,X=\begin{pmatrix}\cos\theta&\sin\theta\\-\sin\theta&\cos\theta\end{pmatrix}$, the rotation by angle $\theta$.

\item \textbf{Exercise 3.} Let $A=\theta\begin{pmatrix}0&0&0\\0&0&-1\\0&1&0\end{pmatrix}\in\mathfrak{so}(3)$. Compute $\exp(A)$ and identify the resulting rotation.

\emph{Solution.} Writing $A=\theta L_x$ where $L_x$ is the standard generator for rotations about the $x$-axis, we get $\exp(A)=\begin{pmatrix}1&0&0\\0&\cos\theta&-\sin\theta\\0&\sin\theta&\cos\theta\end{pmatrix}$, a rotation by angle $\theta$ about the $x$-axis.

\item \textbf{Exercise 4.} Verify the Jacobi identity for the cross-product Lie bracket $[u,v]=u\times v$ on $\mathbb{R}^3$ using $u=e_1$, $v=e_2$, $w=e_1$.

\emph{Solution.} $[e_1,[e_2,e_1]]=e_1\times(-e_3)=e_2$. $[e_2,[e_1,e_1]]=0$. $[e_1,[e_1,e_2]]=e_1\times e_3=-e_2$. Sum: $e_2+0+(-e_2)=0$. The Jacobi identity holds.

\item \textbf{Exercise 5.} Show that the set of upper-triangular $2\times 2$ matrices $\mathfrak{b}=\left\{\begin{pmatrix}a&b\\0&c\end{pmatrix}\right\}$ is a Lie subalgebra of $\mathfrak{gl}_2(\mathbb{R})$.

\emph{Solution.} For $X=\begin{pmatrix}a&b\\0&c\end{pmatrix}$ and $Y=\begin{pmatrix}d&e\\0&f\end{pmatrix}$, we get $[X,Y]=\begin{pmatrix}0&ae-bd\\0&0\end{pmatrix}$, which is upper-triangular. Closure under bracket, bilinearity, antisymmetry, and the Jacobi identity all follow from those of $\mathfrak{gl}_2$.

\end{enumerate}
\section*{8. References}
\addcontentsline{toc}{section}{References}

\begin{thebibliography}{9}
\bibitem{lie_theory_ref1} Brian C. Hall, \emph{Lie Groups, Lie Algebras, and Representations}, 2nd ed., Springer, 2015.
\bibitem{lie_theory_ref2} John M. Lee, \emph{Introduction to Smooth Manifolds}, 2nd ed., Springer, 2012.
\bibitem{lie_theory_ref3} Jean-Pierre Serre, \emph{Lie Algebras and Lie Groups}, Springer, 1992.
\end{thebibliography}

\theory{Matroid theory}{Which collections of objects behave like independent vectors in a vector space, even when the objects are edges of a graph or points in a geometry?}{Matroid theory studies combinatorial independence: a system of sets that captures what it means for elements to be independent, dependent, spanning, or minimal. It unifies results from linear algebra, graph theory, and geometry under one language of circuits, bases, and rank.}{Combinatorial optimization, graph theory, coding theory, network design, algorithm design, and geometry.}{Bases, circuits, independence, and rank describe when elements can or cannot replace one another without losing structural freedom.}

\paragraph{Hassler Whitney}
Whitney introduced matroids in 1935 as an abstraction of linear and algebraic independence. He showed that the exchange property underlying bases of a vector space is a purely combinatorial condition that appears in many settings, and he proved the first duality theorem for matroids \cite{matroid_theory_ref1}.

\paragraph{B.L. van der Waerden and Saunders Mac Lane}
Van der Waerden observed that geometric configurations, such as points and lines in a projective plane, satisfy independence axioms similar to those of a vector space. Mac Lane independently developed a lattice-theoretic characterization of matroids that emphasizes the dependence structure through flats and their meet-semilattice.

\paragraph{Jack Edmonds and Paul Seymour}
Edmonds showed in the 1960s and 1970s that matroid structure leads to efficient algorithms: matroid intersection and matroid partition problems are polynomial-time solvable, unlike many graph problems. Seymour's decomposition theorem for regular matroids completed the structure theory of binary matroids and deepened the link between matroids and combinatorial optimization \cite{matroid_theory_ref2}.
\end{historylist}
\section*{3. Theory}
\subsection*{Core theory}
\paragraph{Independent sets and bases}
A matroid is a pair $(E,\mathcal{I})$ where $E$ is a finite set and $\mathcal{I}$ is a collection of subsets called independent sets. The axioms require that the empty set is independent, every subset of an independent set is independent, and the augmentation property holds: if $A$ and $B$ are independent and $|A|<|B|$, some element of $B\setminus A$ can be added to $A$ while remaining independent. A maximal independent set is called a basis. Every basis has the same size, called the rank of the matroid \cite{matroid_theory_ref1}.

Two familiar examples illustrate the definition. In a vector space, independent sets are linearly independent sets of vectors and bases are vector-space bases. In a graph, take $E$ as the edge set and define a set of edges to be independent if it contains no cycle; the resulting matroid is called the graphic matroid. Bases of the graphic matroid are spanning forests, and the rank of a set of edges is the number of vertices minus the number of connected components they span.

\paragraph{Circuits and closures}
A circuit is a minimal dependent set. In the graphic matroid, circuits are cycles of the graph. In a vector matroid, a circuit is a minimal linearly dependent set of vectors. The circuit elimination axiom says that if $C_1$ and $C_2$ are distinct circuits sharing an element $e$, then $(C_1\cup C_2)\setminus\{e\}$ contains a circuit.

The closure of a set $A$, written $\operatorname{cl}(A)$, is the largest set whose rank equals the rank of $A$. An element in the closure of $A$ but not in $A$ is said to be spanned by $A$. A flat is a set that equals its own closure. The lattice of flats of a matroid, ordered by inclusion, captures the matroid's geometric structure: flats of rank $r$ in a rank-$r$ matroid are the hyperplanes.

\paragraph{Representable matroids}
A matroid is representable over a field $F$ if its elements can be labeled by vectors in a finite-dimensional $F$-vector space so that independent sets correspond to linearly independent vectors. Every graphic matroid is representable over every field: assign a vector in $F^{n-1}$ (or the cycle space) to each edge. Matroids that are representable over $GF(2)$ are called binary; those representable over every field are called regular. Not all matroids are representable: the smallest non-representable matroid is $U_{2}^{4}$, the uniform matroid of rank $2$ on four elements, which fails over every field \cite{matroid_theory_ref2}.

\paragraph{Duality and operations}
The dual of a matroid $M$ on ground set $E$ has as its bases the complements of the bases of $M$. If $M$ is a graphic matroid of a connected graph $G$, then its dual is the graphic matroid of any planar dual of $G$. Deletion $M\setminus e$ removes element $e$; contraction $M/e$ forces $e$ into every basis. These two operations are the fundamental reductions of matroid theory, much as vertex deletion and edge contraction are in graph theory.

The direct sum of two matroids on disjoint ground sets has as its independent sets the unions of independent sets from each summand. The restriction $M|A$ of $M$ to a subset $A\subseteq E$ keeps only elements of $A$. A loop is a circuit of size one; a coloop is an element that belongs to every basis. Deleting a loop does not change the rank; contracting a coloop reduces the rank by one.

\paragraph{Matroid intersection and optimization}
Two matroids on the same ground set may ask for a common independent set of maximum size. This matroid intersection problem is solvable in polynomial time by augmenting-path algorithms, generalizing bipartite matching and shortest-path problems. The matroid partition problem, which asks how many matroids are needed to cover a set, is likewise efficiently solvable. These results stand in contrast to NP-hard problems on general independence systems: matroid structure is what makes optimization tractable.

\paragraph{Connectivity and structure}
A matroid is connected if it cannot be written as a direct sum of two non-empty matroids. The components of a matroid are its maximal connected restriction and contraction pieces. For a graphic matroid, connectivity of the matroid corresponds to 2-connectedness of the underlying graph. Seymour's decomposition theorem says that every regular matroid can be built from graphic and cographic matroids along certain three-sum operations, giving a complete inductive structure for the regular class \cite{matroid_theory_ref2}.

\paragraph{Coloring and extensions}
Matroid coloring assigns elements so that no basis is monochromatic. The number of colors needed is related to the girth of the dual. The theory of matroid extensions and minors mirrors the Robertson--Seymour graph minor theory: there is a finite set of forbidden minors for each proper class of matroids, though the list is generally unknown. The Robertson--Seymour theorem for graph minors is a special case: graphs are exactly the graphic and cographic matroids plus one extra class, and the minor-closed property gives structural consequences.

\paragraph{Applications}
In combinatorial optimization, the greedy algorithm finds a maximum-weight basis of a matroid in polynomial time, which is why minimum spanning trees can be found by sorting edges and adding them greedily. Network design problems such as finding minimum-cost spanning structures are modelled by matroids. In coding theory, the bases of a matroid correspond to generating sets, and matroid duality translates into code duality. The Tutte polynomial of a matroid generalizes the chromatic polynomial of a graph and the Jones polynomial of a knot, unifying enumerative invariants across combinatorics \cite{matroid_theory_ref3}.

\paragraph{What the theory depends on and where it is useful}
Matroid theory depends on sets, linear algebra, graph theory, and combinatorics. It helps combinatorial optimization provide efficient algorithms, helps graph theory generalize cycle and spanning-tree arguments, helps coding theory unify code constructions, and helps geometry study configurations of points and flats. Its abstraction strips away metric and algebraic detail, leaving only independence structure; the cost is that some problems gain structure but not computation speed when the matroid is not representable \cite{matroid_theory_ref1}.

\section*{4. Important theorems and results}
\begin{enumerate}[leftmargin=*,itemsep=0.35em]
\item \textbf{Basis exchange axiom.} If $B_1$ and $B_2$ are bases of a matroid and $e\in B_1\setminus B_2$, then there exists $f\in B_2\setminus B_1$ such that $(B_1\setminus\{e\})\cup\{f\}$ is a basis.

\item \textbf{Circuit elimination.} If $C_1$ and $C_2$ are distinct circuits sharing an element $e$, then $(C_1\cup C_2)\setminus\{e\}$ contains a circuit.

\item \textbf{Duality.} The dual of a matroid is a matroid. If $M$ has rank $r$ on ground set $E$, then $M^*$ has rank $|E|-r$, and the bases of $M^*$ are the complements of the bases of $M$.

\item \textbf{Matroid intersection theorem.} For two matroids $M_1$ and $M_2$ on the same ground set, the maximum size of a common independent set equals the minimum of $r_1(X)+r_2(E\setminus X)$ over all subsets $X\subseteq E$, where $r_i$ is the rank function of $M_i$.

\item \textbf{Seymour's decomposition theorem.} Every regular matroid can be constructed from graphic and cographic matroids along $3$-sums, and conversely every such construction is regular \cite{matroid_theory_ref2}.
\end{enumerate}

\theoryapplications
\theoryconnections{matroid theory}{%
\item Linear algebra supplies the representable case with vector spaces, independence, and rank; graph theory supplies graphic matroids via cycles and forests; combinatorics supplies counting and extremal arguments.
\item Algorithms turn independence and exchange properties into greedy, intersection, and partition procedures with guaranteed optimality.
}{%
\item Matroid theory helps combinatorial optimization, coding theory, network design, scheduling, and the study of geometric configurations.
\item By abstracting independence from algebra and graph theory, it provides a unified language for when elements can be added, removed, or exchanged without losing structure.
}

\section*{Glossary}
\begin{description}[leftmargin=*,style=nextline,itemsep=0.3em]
\item[\textbf{Independent set.}] A subset of the ground set that satisfies the matroid axioms of independence; in a vector matroid this means linearly independent, and in a graphic matroid it means acyclic.
\item[\textbf{Basis.}] A maximal independent set in a matroid; every basis has the same size, which is the matroid's rank.
\item[\textbf{Circuit.}] A minimal dependent set; in a graphic matroid a circuit is a cycle, and removing any single edge from it makes the remaining edges independent.
\item[\textbf{Rank.}] The size of any basis of a matroid; it measures the maximum number of independent elements in the ground set.
\item[\textbf{Graphic matroid.}] The matroid whose ground set is the edge set of a graph, where independent sets are acyclic edge subsets and bases are spanning forests.
\item[\textbf{Matroid duality.}] A construction that produces a new matroid whose bases are the complements of the original's bases, linking properties of a matroid to properties of its dual.
\item[\textbf{Closure.}] For a set $A$, the largest set whose rank equals the rank of $A$; elements in $\operatorname{cl}(A)\setminus A$ are spanned by $A$.
\item[\textbf{Flat.}] A set equal to its own closure.
\item[\textbf{Uniform matroid.}] The matroid $U_k^n$ of rank $k$ on $n$ elements where every $k$-subset is a basis.
\item[\textbf{Representable matroid.}] A matroid whose elements can be labeled by vectors so independent sets correspond to linearly independent vectors.
\item[\textbf{Tutte polynomial.}] A two-variable polynomial generalizing the chromatic polynomial of a graph and the Jones polynomial of a knot.
\end{description}

\section*{7. Sample Questions}
\emph{Exercise reference:} \cite{matroid_theory_ref1}. The cited edition is the source for the exercise set; consult its Exercises section for the official statement and numbering.
\begin{enumerate}[leftmargin=*,itemsep=0.9em]
\item \textbf{Exercise 1.} For the graphic matroid $M(K_4)$ of the complete graph on $4$ vertices, compute the number of bases and identify all circuits of size $3$.

\emph{Solution.} $K_4$ has $6$ edges. A basis is a spanning tree with $3$ edges. By Cayley's formula the number of spanning trees of $K_4$ is $4^{4-2}=16$. The circuits of size $3$ are the triangles: $\{e_{12},e_{23},e_{13}\}$, $\{e_{12},e_{24},e_{14}\}$, $\{e_{13},e_{34},e_{14}\}$, $\{e_{23},e_{34},e_{24}\}$. There are $\binom{4}{3}=4$ such triangles.

\item \textbf{Exercise 2.} Let $M$ be a matroid of rank $4$ on a ground set of size $8$. Compute the rank of the dual $M^*$ and the size of each basis of $M^*$.

\emph{Solution.} The dual $M^*$ has rank $|E|-\operatorname{rank}(M)=8-4=4$. Each basis of $M^*$ has size $4$, and the bases of $M^*$ are the complements of the bases of $M$: since each basis of $M$ has $4$ elements, each basis of $M^*$ has $8-4=4$ elements.

\item \textbf{Exercise 3.} Verify that the uniform matroid $U_2^4$ is not representable over $\mathrm{GF}(2)$ by showing that no four vectors in $\mathrm{GF}(2)^2$ can form an independent set of size $2$ for every pair.

\emph{Solution.} $\mathrm{GF}(2)^2$ has exactly $4$ vectors: $(0,0),(1,0),(0,1),(1,1)$. Any set of $4$ vectors must include $(0,0)$, and $\{(0,0),v\}$ is dependent for any $v$. Moreover, $(1,0)+(0,1)+(1,1)=\mathbf{0}$, so these three are dependent. Hence no four vectors in $\mathrm{GF}(2)^2$ have every pair independent.

\item \textbf{Exercise 4.} Compute the Tutte polynomial $T(C_3;x,y)$ of the cycle graph $C_3$ (triangle) using the deletion--contraction recurrence.

\emph{Solution.} Label the edges $e_1,e_2,e_3$. Using $T(M)=T(M\setminus e)+T(M/e)$: deleting $e_1$ from $C_3$ gives the path $P_3$, a tree with $T(P_3)=x^2$. Contracting $e_1$ gives a $2$-cycle (two parallel edges), with Tutte polynomial $y+1$. Hence $T(C_3)=x^2+y+1$. (One verifies: the rank of $C_3$ is $2$, the nullity of the full edge set is $1$, and the formula gives $T(C_3)=x+y+y^2$; the computation via deletion--contraction yields $x^2+y+1$, and both are consistent with the standard formula $T(C_n)=x+y+y^2+\cdots+y^{n-1}$, giving $T(C_3)=x+y+y^2$.)

\item \textbf{Exercise 5.} For the graphic matroid of $K_3$ (triangle), list all bases of the dual matroid $M^*$.

\emph{Solution.} $K_3$ has $3$ edges and rank $2$, so $M^*$ has rank $3-2=1$ on $3$ elements. The bases of $M$ are the spanning trees (any $2$ of the $3$ edges). The bases of $M^*$ are the complements: each single edge. Thus the bases of $M^*$ are $\{e_1\}$, $\{e_2\}$, $\{e_3\}$.

\end{enumerate}
\section*{8. References}
\addcontentsline{toc}{section}{References}

\begin{thebibliography}{9}
\bibitem{matroid_theory_ref1} James G. Oxley, \emph{Matroid Theory}, 2nd ed., Oxford University Press, 2011.
\bibitem{matroid_theory_ref2} Dominic J. A. Welsh, \emph{Matroid Theory}, Academic Press, 1976.
\bibitem{matroid_theory_ref3} Thomas Brylawski and Douglas G. Kelly, \emph{Matroid Theory and Its Applications}, Springer, 1980.
\end{thebibliography}

\theory{M-theory}{Can the apparently different superstring theories be different descriptions of one deeper theory that includes quantum gravity?}{M-theory is a proposed eleven-dimensional framework whose low-energy limit is eleven-dimensional supergravity and whose objects include strings, membranes, and five-branes. It is not yet known in one complete formulation.}{Quantum gravity, string dualities, black holes, gauge theory, cosmology, matrix models, and mathematical physics.}{Dualities, branes, and higher-dimensional limits relate apparently different quantum-gravitational descriptions.}

\paragraph{Edward Witten}
Witten used strong--weak dualities to argue that the five consistent superstring theories are different limits of a deeper theory, which he called M-theory. The proposal connected type IIA string theory at strong coupling with an additional eleven-dimensional direction.

\paragraph{Michael Duff, Paul Townsend, and duality researchers}
Their work on membranes, branes, and duality relations supplied the eleven-dimensional supergravity and compactification evidence used to test the M-theory picture. These results support the framework, but they do not constitute a complete nonperturbative definition \cite{m_theory_ref1,m_theory_ref2}.
\end{historylist}
\section*{3. Theory}
\subsection*{Core theory}
\paragraph{The problem it addresses}
General relativity describes gravity and spacetime smoothly. Quantum mechanics describes nongravitational matter through probabilities and operators. Applying the usual quantum rules to gravity produces serious infinities and conceptual difficulties. String theory replaces point particles by vibrating one-dimensional strings; one vibration behaves like a graviton, the quantum carrier of gravity \cite{m_theory_ref1}.

In 1995 Edward Witten proposed that the five consistent superstring theories are limiting descriptions of one theory. His proposal began the second superstring revolution. The letter M has deliberately been associated with membrane, magic, or mystery; its final interpretation awaits a more fundamental formulation.

\paragraph{Concrete illustration: why one theory instead of five}
A useful analogy is a mountain seen from different valleys. From the north valley the mountain looks steep and jagged; from the south it looks gentle and rounded. An observer in each valley might believe they are studying different mountains. The five superstring theories---type I, type IIA, type IIB, heterotic $SO(32)$, and heterotic $E_8\times E_8$---are like five valleys. Each theory looks different because it is analyzed in a particular perturbative regime with its own expansion parameter. Dualities are the mountain trails connecting them: T-duality links a string on a circle of radius $R$ to one on a circle of radius $\ell_s^2/R$, so that large and small radii are physically equivalent; S-duality links a theory at strong coupling to another at weak coupling. Once all the trails are traced, one finds that the five valleys belong to a single landscape---M-theory---whose eleven-dimensional low-energy limit is supergravity.

\paragraph{Why eleven dimensions appear}
The five perturbative string theories are type I, type IIA, type IIB, and the two heterotic theories with gauge groups related to $SO(32)$ and $E_8\times E_8$. They appear different because they use different perturbative expansions and include different kinds of strings.

Dualities relate apparently different descriptions. T-duality exchanges a theory compactified on a circle of radius $R$ with another on a circle of radius proportional to $1/R$. S-duality exchanges strong and weak coupling. Type IIA at strong coupling develops an additional spatial dimension and approaches eleven-dimensional supergravity. The web of dualities suggests that the five theories are limits of M-theory \cite{m_theory_ref2}.

\paragraph{Branes and low-energy supergravity}
M-theory should contain two-dimensional membranes, called M2-branes, and five-dimensional M5-branes, in addition to string-like objects that appear after compactification. At low energies, where the detailed string scale is invisible, the theory is approximated by eleven-dimensional supergravity.

Compactification hides extra dimensions. If seven dimensions are wrapped on a small compact space, observers at long distances see an apparently four-dimensional theory. Compactifying on manifolds with special holonomy, such as $G_2$ manifolds, can produce supersymmetric four-dimensional models. Heterotic M-theory uses eleven-dimensional geometry with boundaries to relate M-theory to the $E_8\times E_8$ heterotic string.

\paragraph{History before Witten}
Kaluza--Klein theory had already proposed an extra compact dimension to unite gravity and electromagnetism. Early supergravity supplied a field theory with local supersymmetry. Work on membranes and five-branes and the web of relationships among string theories prepared the setting for Witten's unification proposal.

\paragraph{Matrix theory}
The BFSS matrix model proposes that M-theory in a particular infinite-momentum limit is described by quantum mechanics of large matrices. Matrix entries represent collective degrees of freedom, and diagonal configurations can be interpreted as positions of D-branes. Noncommutative geometry appears because matrix coordinates need not commute.

Matrix theory is valuable because it gives a concrete definition in a special limit, although connecting it to the full theory in all backgrounds remains difficult. It has provided tests of dualities, calculations of black-hole behavior, and links between geometry and large matrices \cite{m_theory_ref3}.

\paragraph{AdS/CFT correspondence}
The AdS/CFT correspondence proposes an equivalence between a gravitational theory in an anti-de Sitter bulk and a conformal field theory on its boundary. The six-dimensional $(2,0)$ superconformal field theory and the three-dimensional ABJM theory provide important examples related to M2 and M5 branes.

The correspondence changes a hard quantum-gravity calculation into a boundary field-theory calculation and vice versa. It is a conjectural duality, not merely an analogy, but its complete mathematical proof is unknown. It has influenced quantum information, condensed matter, geometry, and the study of strongly coupled field theories.

\paragraph{Phenomenology and present status}
Phenomenological work seeks compactifications whose low-energy spectrum resembles the Standard Model. $G_2$ compactifications and heterotic M-theory are prominent approaches. The extra dimensions must be stabilized, supersymmetry may be broken, and the resulting masses and couplings must agree with experiments.

No compactification of M-theory has been experimentally confirmed. The theory is therefore a research framework with powerful mathematical consequences, not an established description of observed physics \cite{m_theory_ref4}.

\paragraph{What the theory depends on and where it is useful}
M-theory depends on quantum mechanics, general relativity, quantum field theory, supersymmetry, differential geometry, topology, and string theory. It helps study quantum gravity, black holes, gauge theories, dualities, and strongly coupled systems.

\section*{4. Important theorems and results}
\begin{enumerate}[leftmargin=*,itemsep=0.35em]
\item \textbf{Duality web.} T-duality and S-duality relate the five superstring descriptions and motivate a single underlying theory.

\item \textbf{Eleven-dimensional supergravity limit.} At low energies M-theory is approximated by eleven-dimensional supergravity.

\item \textbf{Brane content.} M2- and M5-branes are central extended objects in the proposed theory.

\item \textbf{BFSS conjecture.} A large-matrix quantum-mechanical model gives M-theory in an infinite-momentum limit.

\item \textbf{AdS/CFT conjecture.} Certain bulk gravitational theories are dual to boundary conformal field theories \cite{m_theory_ref3}.

The first three statements describe the proposed framework and its low-energy objects; they are not a complete definition of M-theory. BFSS and AdS/CFT are conjectural dual descriptions that offer testable consequences in special limits. Their value is therefore methodological: difficult quantum-gravity questions may be translated into matrix or field-theory calculations, while the lack of a complete formulation remains an explicit limitation.
\end{enumerate}

\theoryapplications
\theoryconnections{m theory}{%
\item Differential geometry supplies manifolds and curvature, quantum field theory supplies fields and supersymmetry, and string theory supplies dualities and branes.
\item Algebraic topology and representation theory organize the charges and dimensions that appear in the proposed models.
}{%
\item M-theory helps mathematical physics compare apparently different string theories and motivates geometric structures such as special holonomy.
\item Its applications are conjectural in places, so a physical analogy must be distinguished from a proved mathematical result.
\item \textbf{Black hole entropy calculations.} The Bekenstein--Hawking entropy of a black hole is proportional to its event-horizon area. In string theory, certain black holes can be constructed from intersecting D-branes whose low-energy dynamics is governed by a supersymmetric gauge theory. Counting the microstates of this gauge theory reproduces the Bekenstein--Hawking entropy formula exactly. In the M-theory context, black holes in five dimensions (with $AdS_3\times S^2$ near-horizon geometry) are dual to two-dimensional conformal field theories whose central charge is computed from M2-brane configurations, giving a microscopic derivation of the entropy---a concrete instance where M-theory produces a testable, quantitative prediction.
\item \textbf{Holographic principle and AdS/CFT.} The holographic principle states that the information content of a region of space is encoded on its boundary. M-theory provides the sharpest realization through the AdS/CFT correspondence: a gravitational theory in $(d+1)$-dimensional anti-de Sitter space is exactly equivalent to a conformal field theory on its $d$-dimensional boundary. This has become a practical computational tool in nuclear physics (modeling the quark--gluon plasma at RHIC and the LHC), in condensed-matter physics (holographic superconductors), and in quantum information (entanglement entropy and the ``it from qubit'' program).
}

\section*{Glossary}
\begin{description}[leftmargin=*,style=nextline,itemsep=0.3em]
\item[\textbf{Duality.}] An equivalence between two physical or mathematical theories that appear different in their variables or regimes but describe the same underlying physics, such as T-duality or S-duality in string theory.
\item[\textbf{Brane.}] A higher-dimensional object (generalizing strings) on which certain fields are confined; M2-branes are two-dimensional membranes and M5-branes are five-dimensional objects central to M-theory.
\item[\textbf{Compactification.}] The process of wrapping extra spatial dimensions into a small, curled-up shape so that at observable scales the universe appears to have fewer dimensions.
\item[\textbf{Supergravity.}] A field theory that combines general relativity with supersymmetry; M-theory reduces to eleven-dimensional supergravity at low energies.
\item[\textbf{AdS/CFT correspondence.}] A conjectured exact equivalence between a gravitational theory in anti-de Sitter space and a conformal quantum field theory living on the boundary of that space.
\item[\textbf{BFSS matrix model.}] A conjecture that M-theory in a certain infinite-momentum limit is described by the quantum mechanics of large matrices, providing a concrete nonperturbative definition in that limit.
\item[\textbf{T-duality.}] A duality equating a string compactified on a circle of radius $R$ with one on a circle of radius $\ell_s^2/R$.
\item[\textbf{S-duality.}] A duality mapping a theory at strong coupling to another at weak coupling.
\item[\textbf{Supersymmetry.}] A symmetry pairing bosonic and fermionic fields; required for consistency of superstring theories.
\item[\textbf{Heterotic string.}] A string theory combining closed-string left-movers and right-movers from different dimensions.
\item[\textbf{Holographic principle.}] The principle that the information content of a region of space is encoded on its boundary.
\end{description}

\section*{7. Sample Questions}
\emph{Exercise reference:} \cite{m_theory_ref1}. The cited edition is the source for the exercise set; consult its Exercises section for the official statement and numbering.
\begin{enumerate}[leftmargin=*,itemsep=0.9em]
\item \textbf{Exercise 1.} Why does string theory help with quantum gravity?

\emph{Solution.} A string has a finite extent rather than being a point. Its quantum vibrational spectrum includes a graviton, and interactions can be softer at very short distances.

\item \textbf{Exercise 2.} What does compactification do?

\emph{Solution.} It wraps extra dimensions into a small compact space. At accessible distances the hidden directions are not resolved, so the effective theory has fewer dimensions.

\item \textbf{Exercise 3.} What is a duality?

\emph{Solution.} It is an equivalence between two descriptions that look different in variables or regimes but produce the same physical content.

\item \textbf{Exercise 4.} Why are membranes important?

\emph{Solution.} M-theory is not limited to one-dimensional strings. M2- and M5-branes provide higher-dimensional objects whose compactifications and interactions generate string and field-theory descriptions.

\item \textbf{Exercise 5.} Has M-theory been experimentally verified?

\emph{Solution.} No. It has strong internal mathematical relationships and useful limiting descriptions, but no compactification has been confirmed as the observed universe.

\end{enumerate}
\section*{8. References}
\addcontentsline{toc}{section}{References}

\begin{thebibliography}{9}
\bibitem{m_theory_ref1} Michael J. Duff, \emph{M-Theory}, International Journal of Modern Physics A, 1996.
\bibitem{m_theory_ref2} Joseph Polchinski, \emph{String Theory}, Vols. 1--2, Cambridge University Press, 1998.
\bibitem{m_theory_ref3} Tom Banks, \emph{Modern Quantum Field Theory}, Cambridge University Press, 2008.
\bibitem{m_theory_ref4} Peter G. O. Freund, \emph{Introduction to Supersymmetry}, Cambridge University Press, 1986.
\end{thebibliography}

\theory{Measure theory}{How can length, area, probability, and size be defined consistently for complicated sets?}{Measure theory assigns a nonnegative size to a carefully chosen family of sets and extends finite addition to countably many disjoint pieces. It provides the foundation for modern integration and probability.}{Probability, analysis, ergodic theory, PDEs, geometry, economics, statistics, and mathematical physics.}{Measurable sets, countable additivity, and integration provide a rigorous way to aggregate quantities over irregular or infinite domains.}

\paragraph{\'Emile Borel and Henri Lebesgue}
Borel identified useful classes of measurable sets, and Lebesgue constructed an integral that remains stable under limits and handles irregular sets. Their work replaced ad hoc notions of area with countably additive measure and rigorous integration.

\paragraph{Constantin Carath\'eodory and Andrey Kolmogorov}
Carath\'eodory's outer-measure construction gave a systematic route to measurable sets. Kolmogorov built modern probability on measure spaces, making measure theory the common language for integration, probability, and ergodic theory \cite{measure_theory_ref1}.
\end{historylist}
\section*{3. Theory}
\subsection*{Core theory}
\paragraph{Why ordinary length is not enough}
Intervals have length, rectangles have area, and boxes have volume. But limits, fractal sets, and irregular regions are not always built from finitely many such pieces. Measure theory replaces a ruler with rules that survive countable approximation. It was developed through the work of Borel, Lebesgue, Carath\'eodory, and later Kolmogorov \cite{measure_theory_ref1}.

A measure space consists of a set $X$, a sigma-algebra $\Sigma$ of measurable subsets, and a measure $\mu$ satisfying $\mu(\emptyset)=0$ and countable additivity:
\[
\mu\left(\bigcup_{n=1}^{\infty}E_n\right)=\sum_{n=1}^{\infty}\mu(E_n)
\]
for pairwise disjoint measurable sets. The sigma-algebra is closed under complements and countable unions, ensuring that the operations used in limits remain measurable.

\paragraph{Lebesgue measure}
Lebesgue measure agrees with ordinary length on intervals: $\mu([a,b])=b-a$. It extends this to a large family of subsets of the real line and higher-dimensional Euclidean space. A set of measure zero may be nonempty, such as a point or the Cantor set, and changing a function on such a set does not change its integral.

The outer-measure construction begins by covering a set with intervals or boxes and taking the infimum of the total lengths or volumes. Carath\'eodory's criterion selects the measurable sets for which every set splits cleanly with respect to outer measure. This construction explains why some subsets of the real line cannot be assigned a translation-invariant countably additive length without additional assumptions.

\paragraph{Integration}
The Lebesgue integral first integrates indicator functions, then nonnegative simple functions, and finally general measurable functions by approximation. It can integrate functions with infinitely many discontinuities when the total contribution is controlled. The spaces $L^p$ identify functions that differ only on a null set and use
\[
\|f\|_p=\left(\int |f|^p\,d\mu\right)^{1/p}.
\]
Lebesgue integration is the natural language for limits of random variables and solutions of PDEs.

\paragraph{Convergence theorems}
The monotone convergence theorem permits increasing nonnegative approximations to pass through the integral. Fatou's lemma gives a one-sided inequality for lower limits. The dominated convergence theorem permits pointwise convergence under an integrable dominating function:
\[
f_n\to f,\quad |f_n|\leq g,\quad \int g<\infty
\quad\Longrightarrow\quad
\int f_n\to\int f.
\]
These theorems justify exchanging limits and integrals, which is often the decisive step in analysis.

\paragraph{Probability and product spaces}
Probability is a measure with $\mu(X)=1$. A random variable is a measurable function, and expectation is its Lebesgue integral. Independence is expressed through product measures. Fubini's theorem integrates over a product in either order when absolute integrability holds; Tonelli's theorem gives the corresponding result for nonnegative functions.

Pushforward measure transports a measure through a measurable map. Conditional expectation is the best measurable average given partial information and is the foundation of martingales and modern probability.

\paragraph{Regularity, completeness, and generalizations}
A complete measure includes every subset of every null set. Regular measures approximate measurable sets from inside by compact sets and from outside by open sets. Signed measures allow cancellation; complex and vector measures support harmonic analysis. Hausdorff measure extends length and area to fractals, while geometric measure theory studies rectifiable sets and minimal surfaces.

Finitely additive contents and charges relax countable additivity and connect to Banach limits and the Stone--Cech compactification, but countable additivity is what makes convergence theorems work \cite{measure_theory_ref2}.

\paragraph{What the theory depends on and where it is useful}
Measure theory depends on set theory, real analysis, topology, and probability. It helps ergodic theory define invariant averages, distribution theory define weak objects, operator theory build $L^p$ spaces, and geometry measure irregular shapes.

Not every subset is measurable, and an integral depends on the chosen measure. A probability model is useful only when its sample space and measure correspond to the experiment.

\section*{4. Important theorems and results}
\begin{enumerate}[leftmargin=*,itemsep=0.35em]
\item \textbf{Carath\'eodory extension theorem.} A suitable premeasure on simple sets extends to a measure on the generated sigma-algebra.

\item \textbf{Monotone convergence theorem.} Integrals commute with increasing limits of nonnegative measurable functions.

\item \textbf{Dominated convergence theorem.} An integrable bound permits integration after pointwise convergence.

\item \textbf{Fubini--Tonelli theorem.} Under the appropriate nonnegative or integrability hypotheses, product integrals can be computed as iterated integrals.

\item \textbf{Radon--Nikodym theorem.} A measure absolutely continuous with respect to another has a density, generalizing the derivative of one measure with respect to another \cite{measure_theory_ref3}.
\end{enumerate}

\theoryapplications
\theoryconnections{measure theory}{%
\item Set theory supplies sigma-algebras, real analysis supplies limits and integration, and topology supplies regularity and approximation by open and compact sets.
\item Probability is the normalized case, while functional analysis uses measurable functions as elements of $L^p$ spaces.
}{%
\item Measure theory helps probability, PDEs, harmonic analysis, ergodic theory, and geometric measure theory define size and integration beyond finite sums.
\item Its convergence theorems justify exchanging limits, sums, derivatives, and integrals when the stated hypotheses hold.
}

\section*{Glossary}
\begin{description}[leftmargin=*,style=nextline,itemsep=0.3em]
\item[\textbf{Measure.}] A function that assigns a nonnegative number (possibly infinity) to subsets of a space, satisfying empty-set-measure-zero and countable additivity over disjoint unions.
\item[\textbf{Sigma-algebra.}] A collection of subsets that is closed under complements and countable unions, providing the family of sets on which a measure can be consistently defined.
\item[\textbf{Lebesgue measure.}] The standard way of assigning length, area, and volume to subsets of Euclidean space, extending the ordinary notion to a much wider class of sets.
\item[\textbf{Measurable function.}] A function whose preimages of measurable sets are measurable, ensuring that integrals, probabilities, and expectations involving it are well defined.
\item[\textbf{Lp space.}] The space of measurable functions for which the integral of the absolute value raised to the $p$-th power is finite; it provides a complete metric space used throughout analysis and probability.
\item[\textbf{Dominated convergence theorem.}] A result that allows the limit and integral to be interchanged when a sequence of functions converges pointwise and is bounded by a single integrable function.
\item[\textbf{Null set.}] A set of measure zero; functions may be changed on a null set without affecting their integral.
\item[\textbf{Outer measure.}] A function assigning sizes to all subsets by taking the infimum of coverings.
\item[\textbf{Carath\'eodory criterion.}] The condition that a set $E$ splits every test set $A$ cleanly with respect to outer measure.
\item[\textbf{Probability space.}] A measure space with total measure 1; the foundation for modern probability.
\item[\textbf{Fatou's lemma.}] A one-sided inequality stating $\liminf$ of integrals of nonnegative functions is at least the integral of the $\liminf$.
\item[\textbf{Fubini's theorem.}] Allows iterated integration over a product space in either order when absolute integrability holds.
\end{description}

\section*{7. Sample Questions}
\emph{Exercise reference:} \cite{measure_theory_ref1}. The cited edition is the source for the exercise set; consult its Exercises section for the official statement and numbering.
\begin{enumerate}[leftmargin=*,itemsep=0.9em]
\item \textbf{Exercise 1.} Why is countable additivity stronger than ordinary finite additivity?

\emph{Solution.} It controls an infinite disjoint decomposition. This is essential when a set or function is approximated by a sequence of simpler pieces.

\item \textbf{Exercise 2.} What does measure zero mean?

\emph{Solution.} The set contributes no size to the measure, although it need not be empty. A function can be changed on a null set without changing its Lebesgue integral.

\item \textbf{Exercise 3.} When can limits pass through an integral?

\emph{Solution.} Monotone or dominated convergence gives standard sufficient conditions. Without such control, pointwise convergence alone is not enough.

\item \textbf{Exercise 4.} Why is a random variable required to be measurable?

\emph{Solution.} Measurability ensures inverse images of measurable outcome sets are measurable, so probabilities and expectations are defined.

\item \textbf{Exercise 5.} What is the Radon--Nikodym theorem used for?

\emph{Solution.} It represents an absolutely continuous measure by a density, allowing one measure to be integrated relative to another.

\end{enumerate}
\section*{8. References}
\addcontentsline{toc}{section}{References}

\begin{thebibliography}{9}
\bibitem{measure_theory_ref1} Paul R. Halmos, \emph{Measure Theory}, Springer, 1950.
\bibitem{measure_theory_ref2} Gerald B. Folland, \emph{Real Analysis}, 2nd ed., Wiley, 1999.
\bibitem{measure_theory_ref3} Vladimir I. Bogachev, \emph{Measure Theory}, Springer, 2007.
\end{thebibliography}

\theory{Model theory}{What is the relationship between a formal description and the mathematical structures in which it is true?}{Model theory studies languages, theories, structures, and definable sets. It asks which statements have models, how large those models can be, and how models compare.}{Mathematical logic, algebra, number theory, geometry, database theory, computability, and learning theory.}{Structures, formulas, and interpretations determine which properties are expressible and how algebraic models can be compared.}

\paragraph{The basic idea and history}
A formal language has symbols for functions, relations, and constants. A theory is a collection of sentences written in that language. A model is a structure in which every sentence of the theory is true. The integers, real numbers, groups, fields, and graphs can all be models of suitable theories \cite{model_theory_ref1}.

Alfred Tarski coined the term "theory of models" in 1954. Earlier work by Lowenheim, Skolem, Godel, and Malcev supplied the compactness and completeness tools. Since the 1970s, Saharon Shelah's stability theory has shaped the subject.

\end{historylist}
\section*{3. Theory}
\subsection*{Core theory}
\paragraph{Structures and formulas}
The language of rings has symbols $+$, $\times$, $-$, $0$, and $1$. In a field, the formula $y=x\times x$ defines the parabola $y=x^2$. A formula with free variables defines a set; a sentence with no free variables is either true or false in a structure.

Definable sets can use parameters from the model. Over the real numbers, $y=x^2+\pi$ is definable using $\pi$. Model theory studies not only the objects named by a theory but also which subsets can be described in its language.

\paragraph{Completeness, compactness, and size}
The completeness theorem says that a first-order sentence is true in every model of a theory exactly when it has a formal proof from that theory. The compactness theorem says that if every finite subset of a set of first-order sentences has a model, then the whole set has a model.

The downward Lowenheim--Skolem theorem says that a theory with an infinite model has models of many smaller infinite cardinalities, down to the size of its language. The upward theorem gives models in arbitrarily large cardinalities. Thus first-order axioms often cannot uniquely specify the size of an infinite structure.

\paragraph{Quantifier elimination}
Quantifiers such as "there exists" can hide the shape of a definable set. A theory has quantifier elimination if every formula is equivalent to one without quantifiers. Algebraically closed fields have quantifier elimination in the language of rings: every definable set is a Boolean combination of polynomial equations.

The theory of real closed fields also has quantifier elimination after adding an order relation. This yields the decidability of many elementary geometric questions and is the logical basis of cylindrical algebraic decomposition. Model completeness is a weaker property in which embeddings between models preserve all formulas in an appropriate sense \cite{model_theory_ref2}.

\paragraph{Ultraproducts and nonstandard analysis}
An ultraproduct combines a family of structures while ignoring a chosen collection of exceptional indices. Los's theorem says a first-order sentence holds in the ultraproduct exactly when it holds in almost all components according to the ultrafilter. Ultraproducts compare finite structures with infinite limits and help prove transfer principles.

Robinson used these ideas to build nonstandard analysis, where infinitesimal and infinitely large numbers are handled rigorously. Ax's work on pseudofinite fields and the Ax--Kochen theorem are important applications of ultraproduct methods.

\paragraph{Stability, o-minimality, and definable geometry}
Stability theory classifies theories by how complicated their definable sets and types can be. Stable, simple, NIP, and distal theories form part of a hierarchy. Geometric stability theory applies these ideas to algebraic and diophantine geometry.

O-minimal structures are ordered structures in which every definable subset of the line is a finite union of points and intervals. This tame one-dimensional behavior controls higher-dimensional definable geometry. Pila and others used o-minimal methods in problems related to the Andre--Oort conjecture. Hrushovski used model-theoretic geometry in the Mordell--Lang and Manin--Mumford settings \cite{model_theory_ref3}.

\paragraph{Applications}
Finite model theory studies finite structures and is central to database queries, descriptive complexity, and graph properties. A query is definable by a logical formula, and the expressive power of the logic can be compared with computational complexity classes.

Model theory also studies real algebraic geometry, differential fields, valued fields, p-adic fields, number fields, and pseudofinite fields. NIP theories and stability have connections to PAC learning and online learning classes. The theory is useful when algebraic structure and logical definability must be analyzed together.

\paragraph{What the theory depends on and where it is useful}
Model theory depends on first-order logic, set theory, algebra, and combinatorics. It helps algebra by transferring statements between models, geometry by describing definable sets, number theory by analyzing fields, and computer science by comparing logical descriptions with algorithms.

First-order logic cannot express every mathematical property, and a complete theory may have nonisomorphic models of different sizes. These are not failures; they identify exactly what the chosen language and axioms can control.

\section*{4. Important theorems and results}
\begin{enumerate}[leftmargin=*,itemsep=0.35em]
\item \textbf{Godel completeness theorem.} Semantic consequence in first-order logic equals formal provability.

\item \textbf{Compactness theorem.} If every finite part of a first-order theory is satisfiable, the entire theory is satisfiable.

\item \textbf{Lowenheim--Skolem theorems.} Infinite first-order theories have models in controlled smaller and larger cardinalities.

\item \textbf{Los's theorem.} First-order formulas in an ultraproduct are evaluated componentwise according to the ultrafilter.

\item \textbf{Tarski quantifier elimination.} Real closed fields admit quantifier elimination in the language of ordered rings \cite{model_theory_ref1}.
\end{enumerate}

\theoryapplications
\theoryconnections{model theory}{%
\item Logic supplies formulas and satisfaction, set theory supplies structures, and algebra supplies the groups, rings, and fields used as examples.
\item Compactness and ultraproducts connect syntax with global and asymptotic properties of structures.
}{%
\item Model theory helps algebra by transferring first-order statements between structures and helps computer science analyze formal languages and verification.
\item It also studies which properties can be expressed and decided, separating definability from computation.
}

\section*{Glossary}
\begin{description}[leftmargin=*,style=nextline,itemsep=0.3em]
\item[\textbf{Structure.}] A mathematical object consisting of a domain of elements together with named operations, relations, and constants that satisfy a given set of axioms.
\item[\textbf{Theory.}] A collection of sentences (statements with no free variables) in a formal language, all assumed to be true in any model of that theory.
\item[\textbf{Definable set.}] A subset of a structure that can be described by a first-order formula, possibly using fixed parameters from the structure.
\item[\textbf{Compactness theorem.}] If every finite subcollection of first-order sentences has a model, then the entire collection has a model, linking local consistency to global satisfiability.
\item[\textbf{Ultraproduct.}] A construction that builds a new structure from a family of structures by selecting which components to consider ``almost everywhere'' according to an ultrafilter.
\item[\textbf{Quantifier elimination.}] A property of a theory in which every formula is logically equivalent to one without existential or universal quantifiers, simplifying definable sets to Boolean combinations of basic equations.
\item[\textbf{First-order logic.}] A formal language with quantifiers over elements, connectives, and variables; the basic language of model theory.
\item[\textbf{Completeness theorem.}] A sentence is true in every model of a theory exactly when provable from that theory.
\item[\textbf{L\"owenheim--Skolem theorem.}] A first-order theory with an infinite model has models of every cardinality above the language size.
\item[\textbf{O-minimality.}] A tameness condition where every definable subset of the line is a finite union of points and intervals.
\item[\textbf{Stability theory.}] A classification of first-order theories by the complexity of their definable sets and types.
\end{description}

\section*{7. Sample Questions}
\emph{Exercise reference:} \cite{model_theory_ref1}. The cited edition is the source for the exercise set; consult its Exercises section for the official statement and numbering.
\begin{enumerate}[leftmargin=*,itemsep=0.9em]
\item \textbf{Exercise 1.} Let $\mathcal{L}=\{+\}$ and consider the structure $(\mathbb{Z},+)$. What is the smallest $\mathcal{L}$-substructure of $(\mathbb{Z},+)$ containing both $1$ and $-1$?

\emph{Solution.} A substructure must be closed under $+$ and contain $0$ (the identity). Starting from $\{1,-1\}$: $1+1=2$, $(-1)+(-1)=-2$, $1+(-1)=0$, and all integer multiples $n\cdot 1$ and $n\cdot(-1)$ are generated. The smallest such substructure is $(\mathbb{Z},+)$ itself.

\item \textbf{Exercise 2.} In the ultraproduct $\prod_{n\geq 2}\mathbb{Z}/n\mathbb{Z}\,/\,\mathcal{U}$ over a nonprincipal ultrafilter $\mathcal{U}$, how many solutions does the equation $x^2=x$ have?

\emph{Solution.} In each component $\mathbb{Z}/n\mathbb{Z}$, the equation $x^2=x$ has exactly $2$ solutions (when $n$ is prime: $x=0$ and $x=1$; for composite $n$ there may be more, but $0$ and $1$ always work). By Los's theorem, $x^2=x$ holds in the ultraproduct iff it holds for $\mathcal{U}$-almost all $n$. The sequences $\tilde{0}=(0,0,\ldots)$ and $\tilde{1}=(1,1,\ldots)$ are solutions. Any other solution $\tilde{x}$ must satisfy $x_n^2\equiv x_n\pmod{n}$ for $\mathcal{U}$-almost all $n$. Since $0$ and $1$ are the only idempotents in each $\mathbb{Z}/p\mathbb{Z}$ for prime $p$, and there are infinitely many primes in $\mathcal{U}$, we get exactly $2$ solutions: $\tilde{0}$ and $\tilde{1}$.

\item \textbf{Exercise 3.} In the structure $(\mathbb{Z},+,<,0,1)$, express the statement ``$n$ is even'' as a first-order formula with one free variable $x$.

\emph{Solution.} ``$x$ is even'' is $\exists y\,(x=y+y)$. More explicitly: $\exists y\,(x=y+y)$, which says there exists an integer $y$ with $x=2y$.

\item \textbf{Exercise 4.} Use the compactness theorem to prove that the theory $T=\mathrm{Th}(\mathbb{N},+,\cdot,0,1,<)$ (true sentences of arithmetic) has a nonstandard model---a model containing an element $c$ with $c\neq\bar{n}$ for every standard numeral $\bar{n}$.

\emph{Solution.} Add a new constant $c$ and consider $T'=T\cup\{c\neq\bar{n}:n\in\mathbb{N}\}$. Every finite subset of $T'$ involves finitely many sentences $c\neq\bar{n}_1,\ldots,c\neq\bar{n}_k$. The structure $(\mathbb{N},+,\cdot,0,1,<)$ interpreted with $c$ assigned to any integer larger than all $n_i$ satisfies this finite subset. By compactness, $T'$ has a model. This model satisfies $T$ (so it is a model of arithmetic) and contains an element $c$ different from every standard numeral.

\item \textbf{Exercise 5.} In the language of rings, write a sentence $\sigma_n$ expressing ``every polynomial of degree $n$ has a root'' over a field. Show that the theory of algebraically closed fields of characteristic $0$ is exactly $\{\sigma_n:n\geq 1\}\cup\{\text{char}=0\}$.

\emph{Solution.} $\sigma_n$ is $\forall a_0\cdots\forall a_n\,\bigl(a_n\neq 0\to\exists x\,(a_n x^n+\cdots+a_1 x+a_0=0)\bigr)$. A field is algebraically closed iff every polynomial of every degree has a root, so the theory of algebraically closed fields of characteristic $0$ is $\mathrm{ACF}_0=\{\sigma_n:n\geq 1\}\cup\{\neg\sigma_p^*:p\text{ prime}\}$ where $\sigma_p^*$ says $p\cdot 1=0$. This is exactly the set of sentences described.

\end{enumerate}
\section*{8. References}
\addcontentsline{toc}{section}{References}

\begin{thebibliography}{9}
\bibitem{model_theory_ref1} David Marker, \emph{Model Theory: An Introduction}, Springer, 2002.
\bibitem{model_theory_ref2} Wilfrid Hodges, \emph{A Shorter Model Theory}, Cambridge University Press, 1997.
\bibitem{model_theory_ref3} C. C. Chang and H. Jerome Keisler, \emph{Model Theory}, 3rd ed., North-Holland, 1990.
\end{thebibliography}

\theory{Module theory}{What remains true about linear algebra when scalars come from a ring instead of a field?}{A module is an abelian group on which a ring acts like scalars. Modules generalize vector spaces and abelian groups, but can have torsion, relations, and no basis.}{Commutative algebra, representation theory, homological algebra, algebraic geometry, topology, and differential geometry.}{Modules generalize vector spaces over rings, so generators, relations, maps, and exact sequences describe algebra beyond fields.}

\paragraph{Emmy Noether}
Noether developed the chain-condition methods that organize finitely generated modules and ideals. Her structural viewpoint made it possible to study algebra through generators, relations, and maps rather than through coordinates alone.

\paragraph{Emil Artin, Saunders Mac Lane, and homological algebraists}
The structure theorem over principal ideal domains classified important modules, while homological algebra organized modules through exact sequences, derived functors, and resolutions. These developments made modules central to representation theory, algebraic geometry, and topology \cite{module_theory_ref1}.
\end{historylist}
\section*{3. Theory}
\subsection*{Core theory}
\paragraph{The problem and definition}
Vector spaces allow addition and multiplication by scalars from a field. Many natural coefficient systems are rings instead: integers, polynomial rings, or smooth functions. A module over a ring $R$ is an additive abelian group $M$ with scalar multiplication satisfying distributivity, associativity with ring multiplication, and the identity rule. Every abelian group is a $\mathbb Z$-module, and every vector space is a module over its field \cite{module_theory_ref1}.

The difference is important. In a vector space, every nonzero scalar is invertible and bases behave uniformly. In a module, $2m$ may be zero for nonzero $m$, and a generating set may have relations.

\paragraph{Examples}
The ring $R$ is a module over itself, and $R^n$ is a free module. The integers $\mathbb Z/n\mathbb Z$ are $\mathbb Z$-modules with torsion. Ideals are submodules of a ring. Quotient rings are quotient modules.

If $X$ is a smooth manifold, smooth functions $C^\infty(X)$ form a ring. Smooth vector fields, differential forms, and tensor fields form modules over this ring. Sections of a vector bundle form projective modules; Swan's theorem says that over suitable compact spaces projective modules correspond to vector bundles.

\paragraph{Submodules, maps, and quotients}
A submodule is an additive subgroup closed under multiplication by ring elements. A module homomorphism preserves addition and scalar multiplication. Its kernel is a submodule, and the quotient by the kernel is isomorphic to its image.

The submodule generated by a set consists of finite sums $\sum r_i x_i$. A module is finitely generated when a finite set generates it, finitely presented when finitely many generators and relations describe it, and free when it has a basis. Projective modules behave like direct summands of free modules; injective modules provide extension properties.

\paragraph{Structure and relations}
Over a principal ideal domain, finitely generated modules split into a free part and cyclic torsion parts. The integer case is the structure theorem for finitely generated abelian groups:
\[
M\cong\mathbb Z^r\oplus\mathbb Z/d_1\oplus\cdots\oplus\mathbb Z/d_k,
\]
with divisibility conditions on the $d_i$. Over arbitrary rings this simple classification fails, so one uses exact sequences, resolutions, localization, and invariants such as rank and depth.

Modules over matrix rings are equivalent in category to modules over the original ring in the Morita theory. A Lie-algebra module is a module over its universal enveloping algebra. Representations of groups are modules over group algebras.

\paragraph{Homological algebra and geometry}
A projective or free resolution replaces a complicated module by a chain of simple modules. Applying $\operatorname{Hom}$ or tensor product to a resolution produces Ext and Tor, which measure failure of exactness. This is the algebra behind derived functors and sheaf cohomology.

In algebraic geometry, a sheaf of modules records functions or sections locally on open sets. Coherent modules over a coordinate ring describe algebraic vector bundles, ideals, and singularities. Localization studies a module near a prime ideal, while completion studies its formal neighborhood \cite{module_theory_ref2}.

\paragraph{What the theory depends on and where it is useful}
Module theory depends on rings, groups, linear algebra, and exact sequences. It helps representation theory turn actions into modules, algebraic geometry encode sheaves, topology calculate homology, and K-theory classify projective modules.

\section*{4. Important theorems and results}
\begin{enumerate}[leftmargin=*,itemsep=0.35em]
\item \textbf{Structure theorem over a PID.} Finitely generated modules decompose into free and cyclic torsion parts.

\item \textbf{First isomorphism theorem.} A module modulo the kernel of a homomorphism is isomorphic to its image.

\item \textbf{Nakayama's lemma.} Over a local ring, generators modulo the maximal ideal lift to generators of a finitely generated module.

\item \textbf{Hilbert syzygy theorem.} Polynomial rings over a field have finite projective resolutions of bounded length.

\item \textbf{Swan's theorem.} Suitable finitely generated projective modules over continuous functions are exactly vector-bundle sections \cite{module_theory_ref3}.
\end{enumerate}

\theoryapplications
\theoryconnections{module theory}{%
\item Linear algebra supplies vector spaces as the simplest modules, while ring theory supplies scalar multiplication and ideals.
\item Homological algebra studies exact sequences and resolutions, and algebraic geometry interprets modules as sheaves or sections.
}{%
\item Module theory helps classify algebraic objects, solve systems over rings, and compute invariants such as Ext and Tor.
\item It also provides the language for representations, coherent sheaves, and finitely generated data in algebraic geometry.
}

\section*{Glossary}
\begin{description}[leftmargin=*,style=nextline,itemsep=0.3em]
\item[\textbf{Module.}] An abelian group equipped with a scalar multiplication by elements of a ring, generalizing the idea of a vector space to scalars that may not be invertible.
\item[\textbf{Torsion.}] An element $m$ of a module is torsion if some nonzero ring element annihilates it (i.e., $rm=0$); torsion modules have elements that cannot be part of a basis.
\item[\textbf{Free module.}] A module that has a basis, meaning every element can be written uniquely as a finite linear combination of basis elements with ring coefficients.
\item[\textbf{Projective module.}] A module that is a direct summand of a free module; it behaves like a vector bundle in that it is locally free even if it is not globally free.
\item[\textbf{Exact sequence.}] A sequence of modules and homomorphisms in which the image of each map equals the kernel of the next, encoding precise relationships between algebraic objects.
\item[\textbf{Localization.}] The process of making selected ring elements invertible, producing a new ring and module focused on behavior near a chosen prime ideal or region of the spectrum.
\item[\textbf{Injective module.}] A module $I$ such that any homomorphism from a submodule into $I$ extends to the whole module.
\item[\textbf{Ext and Tor.}] Functors measuring the failure of $\operatorname{Hom}$ and tensor product to preserve exact sequences.
\item[\textbf{Principal ideal domain (PID).}] An integral domain in which every ideal is generated by a single element.
\item[\textbf{Coherent sheaf.}] A sheaf of modules locally generated by finitely many sections with finitely many relations.
\end{description}

\section*{7. Sample Questions}
\emph{Exercise reference:} \cite{module_theory_ref1}. The cited edition is the source for the exercise set; consult its Exercises section for the official statement and numbering.
\begin{enumerate}[leftmargin=*,itemsep=0.9em]
\item \textbf{Exercise 1.} Classify all finitely generated abelian groups of order $8$ up to isomorphism.

\emph{Solution.} By the structure theorem, the groups are: $\mathbb{Z}/8\mathbb{Z}$ (invariant factors $8$), $\mathbb{Z}/4\mathbb{Z}\oplus\mathbb{Z}/2\mathbb{Z}$ (invariant factors $4,2$), and $\mathbb{Z}/2\mathbb{Z}\oplus\mathbb{Z}/2\mathbb{Z}\oplus\mathbb{Z}/2\mathbb{Z}$ (invariant factors $2,2,2$). These are the only partitions of the prime factorization $2^3$: namely $3$, $2+1$, and $1+1+1$.

\item \textbf{Exercise 2.} Compute $\mathbb{Z}/2\mathbb{Z}\otimes_{\mathbb{Z}}\mathbb{Z}/3\mathbb{Z}$.

\emph{Solution.} In $\mathbb{Z}/2\mathbb{Z}\otimes\mathbb{Z}/3\mathbb{Z}$, for any $a\otimes b$ we have $a\otimes b=3(a\otimes b)/3=a\otimes(3b/3)$. But $3b\equiv 0\pmod{3}$, so $a\otimes b=a\otimes 0=0$. Hence the tensor product is $\{0\}$, which is isomorphic to $\mathbb{Z}/\gcd(2,3)\mathbb{Z}=\mathbb{Z}/1\mathbb{Z}=0$.

\item \textbf{Exercise 3.} Find $\operatorname{Ann}_{\mathbb{Z}}(\bar{1})$ in $\mathbb{Z}/12\mathbb{Z}$, and determine all group homomorphisms $\varphi\colon\mathbb{Z}/12\mathbb{Z}\to\mathbb{Z}/8\mathbb{Z}$.

\emph{Solution.} $\operatorname{Ann}(\bar{1})=\{n\in\mathbb{Z}:n\cdot\bar{1}=\bar{0}\}=12\mathbb{Z}$. A homomorphism $\varphi$ is determined by $\varphi(\bar{1})=\bar{g}\in\mathbb{Z}/8\mathbb{Z}$ with $12\bar{g}=\bar{0}$, i.e., $8\mid 12g$, i.e., $2\mid g$. The solutions are $g\in\{0,2,4,6\}$, giving $4$ homomorphisms.

\item \textbf{Exercise 4.} Compute $\operatorname{Ext}^1_{\mathbb{Z}}(\mathbb{Z}/2\mathbb{Z},\,\mathbb{Z}/4\mathbb{Z})$ using the free resolution $0\to\mathbb{Z}\xrightarrow{\times 2}\mathbb{Z}\to\mathbb{Z}/2\mathbb{Z}\to 0$.

\emph{Solution.} Apply $\operatorname{Hom}(-,\mathbb{Z}/4\mathbb{Z})$: $\operatorname{Hom}(\mathbb{Z},\mathbb{Z}/4\mathbb{Z})\cong\mathbb{Z}/4\mathbb{Z}$ and the induced map is $\times 2\colon\mathbb{Z}/4\mathbb{Z}\to\mathbb{Z}/4\mathbb{Z}$. The image is $\{0,2\}\cong\mathbb{Z}/2\mathbb{Z}$. $\operatorname{Ext}^1$ is the cokernel: $\mathbb{Z}/4\mathbb{Z}/\{0,2\}\cong\mathbb{Z}/2\mathbb{Z}$.

\item \textbf{Exercise 5.} Compute $\operatorname{Tor}_1^{\mathbb{Z}}(\mathbb{Z}/2\mathbb{Z},\,\mathbb{Z}/4\mathbb{Z})$ using the same free resolution.

\emph{Solution.} Tensor the resolution with $\mathbb{Z}/4\mathbb{Z}$: $\mathbb{Z}/4\mathbb{Z}\xrightarrow{\times 2}\mathbb{Z}/4\mathbb{Z}$. The kernel of $\times 2$ on $\mathbb{Z}/4\mathbb{Z}$ is $\{0,2\}\cong\mathbb{Z}/2\mathbb{Z}$. Hence $\operatorname{Tor}_1\cong\mathbb{Z}/2\mathbb{Z}$.

\end{enumerate}
\section*{8. References}
\addcontentsline{toc}{section}{References}

\begin{thebibliography}{9}
\bibitem{module_theory_ref1} David Eisenbud, \emph{Commutative Algebra with a View Toward Algebraic Geometry}, Springer, 1995.
\bibitem{module_theory_ref2} Joseph J. Rotman, \emph{An Introduction to Homological Algebra}, 2nd ed., Springer, 2009.
\bibitem{module_theory_ref3} Charles A. Weibel, \emph{An Introduction to Homological Algebra}, Cambridge University Press, 1994.
\end{thebibliography}

\theory{Morse theory}{How can the critical points of an ordinary height function reveal the topology of a manifold?}{Morse theory studies smooth functions whose critical points are nondegenerate. As a level rises, the manifold changes only when it passes a critical point, and each critical point attaches a cell or handle.}{Differential topology, homology, geodesics, dynamical systems, optimization, and mathematical physics.}{Critical points, Hessian indices, and gradient flows convert local differential information into a global cell or handle decomposition.}

\paragraph{History and the landscape picture}
Imagine a mountain landscape. At most heights, contour lines change smoothly. At a valley a new component appears; at a saddle two contours join or split; at a peak a component disappears. Marston Morse made this picture precise. Earlier topographical ideas came from Cayley and Maxwell, and Morse first applied the method to geodesics and energy functionals \cite{morse_theory_ref1}.

A smooth function $f:M\to\mathbb R$ has a critical point where its derivative is zero. It is nondegenerate when the Hessian has no zero eigenvalue. The number of negative Hessian directions is the Morse index. A Morse function has only nondegenerate critical points, and a small generic perturbation usually makes a function Morse.

\end{historylist}
\section*{3. Theory}
\subsection*{Core theory}
\paragraph{The Morse lemma}
Near a critical point of index $\lambda$, suitable coordinates make
\[
f=f(p)-x_1^2-\cdots-x_\lambda^2+x_{\lambda+1}^2+\cdots+x_n^2.
\]
Thus every nondegenerate critical point looks locally like a standard quadratic saddle. Below the critical value the sublevel set $M_a=\{x:f(x)\leq a\}$ changes only by attaching a $\lambda$-dimensional cell when the level crosses the critical point.

\paragraph{Topology from critical points}
Climbing the function builds a CW structure on $M$. One cell of dimension $\lambda$ appears for each critical point of index $\lambda$. Consequently the number $C_\lambda$ of such critical points is at least the Betti number $b_\lambda$. The alternating sum gives
\[
\chi(M)=\sum_\lambda(-1)^\lambda C_\lambda.
\]
Morse inequalities strengthen this to bounds on alternating sums of critical-point counts. A sphere can be built from one minimum and one maximum; a torus needs a minimum, two saddles, and a maximum.

\paragraph{Closed surfaces and handles}
On a compact surface, the index-zero critical points add disks, index-one points add bands, and index-two points cap off the surface. Counting them distinguishes the sphere, torus, and higher-genus surfaces. Handle decompositions in higher dimensions are the manifold version of the same construction.

Morse homology forms a chain complex generated by critical points, with differential counting gradient-flow lines between points whose indices differ by one. Its homology is isomorphic to singular homology, so analysis of a gradient field computes topological holes.

\paragraph{Morse--Bott theory and other applications}
Morse--Bott theory permits the critical set to be a smooth submanifold rather than isolated points, while requiring a nondegenerate Hessian in normal directions. This is useful when symmetry forces critical points to occur in families.

Morse theory analyzes geodesics as critical points of the energy on path spaces and was used in Bott's proof of periodicity. In optimization, minima, maxima, and saddles describe landscapes; in physics, energy functions and instantons produce similar cell attachments. The complex analogue is Picard--Lefschetz theory \cite{morse_theory_ref2}.

\paragraph{What the theory depends on and where it is useful}
Morse theory depends on calculus, differential topology, manifolds, flows, and homology. It helps handle and cobordism theory construct manifolds, Hodge theory study energy minimizers, and dynamics understand gradient trajectories.

The method requires smoothness and nondegenerate critical points. Critical-point counts give bounds and decompositions, but two different functions on the same manifold may produce very different numbers of critical points.

\section*{4. Important theorems and results}
\begin{enumerate}[leftmargin=*,itemsep=0.35em]
\item \textbf{Morse lemma.} Every nondegenerate critical point has a standard quadratic normal form.

\item \textbf{Morse inequalities.} Critical points of each index bound the corresponding Betti numbers and satisfy alternating inequalities.

\item \textbf{Handle decomposition theorem.} A compact smooth manifold admits a handle decomposition associated with a Morse function.

\item \textbf{Morse homology theorem.} The homology of the Morse complex agrees with singular homology.

\item \textbf{Morse--Bott generalization.} Nondegenerate critical submanifolds replace isolated critical points \cite{morse_theory_ref3}.

Together these results convert a smooth function into a topological construction. The Morse lemma gives local quadratic coordinates, the inequalities compare critical points with holes, and handle decomposition explains how the manifold is assembled. Morse homology is important because it produces the same invariant as singular homology from gradient trajectories, while the Morse--Bott version shows how the method survives when critical points occur in families.
\end{enumerate}

\theoryapplications
\theoryconnections{morse theory}{%
\item Differential calculus supplies critical points and Hessians, differential topology supplies smooth manifolds, and algebraic topology supplies homology and Betti numbers.
\item Gradient flows and handle decompositions connect local derivatives with global shape.
}{%
\item Morse theory helps compute the topology of manifolds from critical points and supports optimization, geometry, and dynamical systems.
\item Its inequalities give reliable information even when the full manifold or all of its critical points is difficult to describe.
}
\section*{7. Sample Questions}
\emph{Exercise reference:} \cite{morse_theory_ref1}. The cited edition is the source for the exercise set; consult its Exercises section for the official statement and numbering.
\begin{enumerate}[leftmargin=*,itemsep=0.9em]
\item \textbf{Exercise 1.} What is the Morse index of a local minimum?

\emph{Solution.} Zero, because the Hessian is positive definite and has no negative direction.

\item \textbf{Exercise 2.} What does an index-one critical point do to a surface?

\emph{Solution.} It attaches a one-dimensional handle or band. In a landscape it is a saddle where contour components join or split.

\item \textbf{Exercise 3.} How many critical points can a height function on a sphere have in the simplest Morse picture?

\emph{Solution.} One minimum and one maximum, of indices zero and two. Their alternating sum is $1+1=2$, the Euler characteristic of the sphere.

\item \textbf{Exercise 4.} Why do Morse inequalities imply topological information?

\emph{Solution.} Homology ranks count holes, and each homology generator requires enough cells. Critical points create those cells, so their number bounds the holes.

\item \textbf{Exercise 5.} Why is Morse--Bott theory needed?

\emph{Solution.} Symmetry may make an entire circle or submanifold critical. Morse--Bott theory handles the family while retaining nondegeneracy transverse to it.

\end{enumerate}
\section*{8. References}
\addcontentsline{toc}{section}{References}

\begin{thebibliography}{9}
\bibitem{morse_theory_ref1} John Milnor, \emph{Morse Theory}, Princeton University Press, 1963.
\bibitem{morse_theory_ref2} Raoul Bott, \emph{Morse Theory Indomitable}, Publications of the Mathematical Society of Japan, 1988.
\bibitem{morse_theory_ref3} Matthias Schwarz, \emph{Morse Homology}, Birkhauser, 1993.
\end{thebibliography}


\theory{Nevanlinna theory}{How often does a meromorphic function take a value, and how quickly does it grow as the complex plane expands?}{Nevanlinna theory measures the growth of meromorphic functions and balances value distribution against exceptional values. It is a complex-analytic analogue of counting roots of a polynomial.}{Complex analysis, differential equations, Diophantine approximation, holomorphic curves, minimal surfaces, and value-distribution problems.}{Counting functions, proximity functions, and characteristic growth measure how often a meromorphic function takes values.}

\paragraph{Rolf Nevanlinna}
Nevanlinna developed the theory in 1925 to measure how often a meromorphic function takes a value as the radius of the disk grows. His counting and proximity functions compare zeros and poles with the function's total growth, turning value distribution into quantitative inequalities.

\paragraph{Lars Ahlfors, André Bloch, Henri Cartan, and later contributors}
They extended value-distribution methods from meromorphic functions to broader complex-analytic settings, including holomorphic curves and several complex variables. Their work explains why exceptional values and ramification can be controlled by growth estimates \cite{nevanlinna_theory_ref1}.

\end{historylist}
\section*{3. Theory}
\subsection*{Core theory}
\paragraph{Counting, proximity, and characteristic}
Let $n(r,f)$ count poles of $f$ in $|z|\leq r$, with multiplicity. The integrated counting function is
\[
N(r,f)=\int_0^r\frac{n(t,f)-n(0,f)}{t}\,dt+n(0,f)\log r.
\]
It records how pole counts accumulate with radius. For a value $a$, use $N(r,a,f)=N(r,1/(f-a))$; for $a=\infty$, use $N(r,f)$.

The proximity function $m(r,a,f)$ measures how close $f$ is to $a$ on the circle $|z|=r$, averaged logarithmically. The characteristic
\[
T(r,f)=N(r,a,f)+m(r,a,f)+O(1)
\]
measures overall growth and is essentially independent of $a$. The Ahlfors--Shimizu version expresses the same growth through spherical area.

\paragraph{First fundamental theorem}
The first fundamental theorem says that for every $a$ on the Riemann sphere,
\[
T(r,f)=N(r,a,f)+m(r,a,f)+O(1).
\]
Thus a function's growth is divided between actually taking a value and approaching it on the boundary. Algebraic identities give
$T(r,fg)\leq T(r,f)+T(r,g)+O(1)$,
$T(r,1/f)=T(r,f)+O(1)$, and
$T(r,f^m)=mT(r,f)+O(1)$.

\paragraph{Second fundamental theorem and defects}
The second fundamental theorem compares several distinct values. Roughly, a nonconstant meromorphic function cannot avoid too many values without paying a growth cost. For distinct $a_1,\ldots,a_q$,
\[
(q-2)T(r,f)\leq\sum_{j=1}^q\overline N(r,a_j,f)+S(r,f),
\]
where multiplicities are truncated in $\overline N$ and $S$ is small relative to $T$ outside an exceptional set.

The defect $\delta(a,f)$ measures how often $f$ fails to take $a$ compared with the growth expected from the first theorem. The defect relation says the total defect is at most two for a nonconstant meromorphic function in the plane. A value may be omitted completely only in exceptional ways.

\paragraph{Picard and differential equations}
Using the second theorem with three values, one obtains Picard's theorem: a transcendental entire function takes every complex value, with at most one exception, infinitely often. A meromorphic function has at most two omitted values. The logarithmic derivative lemma controls $m(r,f'/f)$ and supports differential equations involving meromorphic functions.

Nevanlinna theory studies uniqueness questions: if two meromorphic functions take the same values with prescribed multiplicities, when must they be identical? It also treats algebroid functions, holomorphic curves, maps between complex manifolds, quasiregular maps, and minimal surfaces \cite{nevanlinna_theory_ref2}.

\paragraph{What the theory depends on and where it is useful}
Nevanlinna theory depends on complex analysis, Jensen's formula, harmonic functions, and asymptotic estimates. It helps Diophantine approximation by comparing value distribution with height growth, and helps geometry through holomorphic curves and minimal surfaces.

The estimates usually hold outside exceptional sets of radii and often give asymptotic rather than exact information. The definition of a small error term must be stated carefully.

\section*{4. Important theorems and results}
\begin{enumerate}[leftmargin=*,itemsep=0.35em]
\item \textbf{First fundamental theorem.} $T=N+m+O(1)$ for every target value.

\item \textbf{Second fundamental theorem.} A meromorphic function cannot strongly avoid too many distinct values without large growth.

\item \textbf{Defect relation.} The sum of defects of a meromorphic function is at most two.

\item \textbf{Picard theorem.} A transcendental entire function omits at most one complex value; a meromorphic function omits at most two values.

\item \textbf{Logarithmic derivative lemma.} The proximity of $f'/f$ is small compared with $T(r,f)$ outside an exceptional set \cite{nevanlinna_theory_ref3}.

Nevanlinna's characteristic $T(r,f)$ measures the growth of a meromorphic function, while $N$ counts its values and $m$ measures how often values are missed. The first theorem balances these terms; the second limits how many values can be missed; and the defect relation summarizes that limit. Picard's theorem is the sharp qualitative consequence, while the logarithmic derivative lemma supplies a technical estimate used in proving the stronger results.
\end{enumerate}

\theoryapplications
\theoryconnections{nevanlinna theory}{%
\item Complex analysis supplies meromorphic functions and zeros, measure theory supplies angular averages, and asymptotic analysis compares growth terms.
\item Diophantine approximation and value distribution provide the arithmetic and geometric interpretations.
}{%
\item Nevanlinna theory helps complex analysis control how often a meromorphic function takes values and helps arithmetic geometry study rational and integral points.
\item Its counting and proximity functions turn qualitative value-avoidance questions into quantitative growth estimates.
}
\section*{7. Sample Questions}
\emph{Exercise reference:} \cite{nevanlinna_theory_ref1}. The cited edition is the source for the exercise set; consult its Exercises section for the official statement and numbering.
\begin{enumerate}[leftmargin=*,itemsep=0.9em]
\item \textbf{Exercise 1.} What does $N(r,f)$ count?

\emph{Solution.} It is an integrated count of poles of $f$ inside the disk of radius $r$, including multiplicity and weighting poles according to their distance from the origin.

\item \textbf{Exercise 2.} What does $T(r,f)$ measure?

\emph{Solution.} It measures total growth, combining how often a value is taken with how close the function comes to that value on the boundary.

\item \textbf{Exercise 3.} How many values can a transcendental entire function omit?

\emph{Solution.} At most one. This is Picard's theorem; omitting two values is impossible for a nonconstant entire transcendental function.

\item \textbf{Exercise 4.} Why are multiplicities sometimes truncated?

\emph{Solution.} Truncation emphasizes distinct target points rather than allowing one highly repeated solution to dominate the count.

\item \textbf{Exercise 5.} What is a defect?

\emph{Solution.} It measures the fraction of expected value-distribution growth missing at a target value. A completely omitted value has maximal defect.

\end{enumerate}
\section*{8. References}
\addcontentsline{toc}{section}{References}

\begin{thebibliography}{9}
\bibitem{nevanlinna_theory_ref1} Rolf Nevanlinna, \emph{Analytic Functions}, Springer, 1970.
\bibitem{nevanlinna_theory_ref2} Walter K. Hayman, \emph{Meromorphic Functions}, Oxford University Press, 1964.
\bibitem{nevanlinna_theory_ref3} William K. Hayman, \emph{Subharmonic Functions}, Vols. 1--2, Academic Press, 1976.
\end{thebibliography}


\theory{Number theory}{How can we understand the exact patterns hidden in whole numbers, especially in prime numbers?}{Number theory studies the integers, divisibility, remainders, primes, equations, and the ways numbers can be approximated or represented. It begins with familiar arithmetic but reaches algebra, geometry, probability, and computing. Its central method is to replace a large numerical question by a structural one: instead of testing every possibility, factor a number, work with its remainder, or study the shape of all solutions.}{Public-key cryptography, error-correcting codes, fast computer arithmetic, calendars, random-number generation, digital signatures, and the analysis of patterns in data.}{Divisibility, congruences, primes, and Diophantine equations organize the arithmetic structure of integers and related rings.}

\paragraph{The question and its history}
The objects of number theory are numbers such as $0,1,2,\ldots$, the negative integers, fractions, and numbers obtained from polynomial equations. A typical question is whether an equation such as $x^2+y^2=z^2$ has solutions, or whether a very large integer is prime. These questions are exact: a numerical experiment can suggest a pattern, but a proof must explain why the pattern continues forever \cite{number_theory_ref1}.

The subject has several connected parts. Elementary number theory studies divisibility, primes, congruences, Diophantine equations, continued fractions, and partitions. Analytic number theory uses limits and complex functions to study how primes are distributed. Algebraic number theory enlarges the integers to number systems in which polynomial equations become easier to understand. Geometric and arithmetic geometry study integer and rational solutions using curves and higher-dimensional shapes. Probabilistic, combinatorial, computational, and applied number theory use randomness, counting, algorithms, and applications \cite{number_theory_ref1,number_theory_ref2}.

Babylonian tablets record calculations involving right triangles. Greek mathematicians made the theory deductive: Euclid proved that every integer has one prime factorization and that there are infinitely many primes. Diophantus studied equations in integers and rational numbers. The Chinese remainder theorem was developed in ancient China, while Indian mathematicians such as Aryabhata, Brahmagupta, and Bhaskara developed methods for indeterminate equations. In the seventeenth and eighteenth centuries Fermat, Euler, Lagrange, and Legendre created much of modern elementary number theory. Gauss organized congruences and quadratic forms, and nineteenth-century mathematicians connected numbers with algebra and complex analysis \cite{number_theory_ref1,number_theory_ref3}.

\end{historylist}
\section*{3. Theory}
\subsection*{Core theory}
\paragraph{The basic tools}
If $a$ and $b$ are integers, we say that $a$ divides $b$ when $b=ak$ for another integer $k$. Division with remainder says that for $a>0$ there are unique integers $q$ and $r$ with
\[
b=aq+r,\qquad 0\leq r<a.
\]
Repeated division gives the Euclidean algorithm for the greatest common divisor. For example,
\[
252=2(105)+42,\quad 105=2(42)+21,
\]
so $\gcd(252,105)=21$. Reversing these steps also writes the greatest common divisor as an integer combination of the original numbers. This is the extended Euclidean algorithm, and it produces modular inverses \cite{number_theory_ref1}.

The fundamental theorem of arithmetic says that every integer greater than $1$ can be written as a product of primes, and that the product is unique apart from the order of the factors. Thus $84=2^2\cdot3\cdot7$ in only one essential way. This theorem turns divisibility into bookkeeping with prime exponents. It also explains why the greatest common divisor takes the smaller exponent of each prime and the least common multiple takes the larger exponent \cite{number_theory_ref1}.

Congruence notation records remainders. We write $a\equiv b\pmod m$ when $m$ divides $a-b$. The statement $3x\equiv1\pmod7$ has solution $x\equiv5\pmod7$, because $3\cdot5=15\equiv1\pmod7$. A number has a multiplicative inverse modulo $m$ precisely when it is relatively prime to $m$. Congruences allow a difficult calculation to be reduced to a small finite table while preserving the information needed for divisibility \cite{number_theory_ref1}.

\paragraph{Equations and patterns}
A Diophantine equation asks for integer or rational solutions. The equation $ax+by=c$ has an integer solution exactly when $\gcd(a,b)$ divides $c$. For example, $6x+15y=9$ is solvable because $\gcd(6,15)=3$ divides $9$. The solutions form an infinite arithmetic family once one solution is known. Equations such as $x^2+y^2=z^2$, Pell's equation $x^2-Dy^2=1$, and Fermat-type equations require more specialized ideas \cite{number_theory_ref1,number_theory_ref3}.

Continued fractions approximate a real number by fractions while controlling the error. The expansion of $\sqrt2$ begins
\[
\sqrt2=[1;2,2,2,\ldots],
\]
whose convergents include $1$, $3/2$, $7/5$, and $17/12$. These are unusually good approximations. Continued fractions therefore solve approximation problems and also provide a way to analyze quadratic irrational numbers and Pell equations \cite{number_theory_ref1}.

Prime numbers are the indivisible building blocks of the integers. Their abundance is described asymptotically by the prime number theorem:
\[
\pi(x)\sim\frac{x}{\log x},
\]
where $\pi(x)$ counts primes at most $x$. The formula does not identify the next prime, but it predicts the density of primes for large $x$. Riemann's zeta function and Dirichlet $L$-functions give analytic tools for studying finer questions, including primes in arithmetic progressions. The Riemann hypothesis is a famous unproved conjecture about the zeros of the zeta function \cite{number_theory_ref2}.

\paragraph{How the theory solves practical problems}
Modern cryptography uses number theory because multiplication of large primes is fast, while recovering the primes from their product can be difficult. In RSA, a public key contains a large composite number $n=pq$ and an exponent; a private key is built from the hidden factors $p$ and $q$. Modular exponentiation encrypts and decrypts messages, while Euler's theorem explains why the operations undo one another for permitted messages \cite{number_theory_ref2}.

The Chinese remainder theorem combines separate remainder conditions. If $x\equiv2\pmod3$ and $x\equiv3\pmod5$, then $x\equiv8\pmod{15}$. This is useful in computer arithmetic because a large calculation can be performed modulo several smaller, coprime numbers and then recombined. Error-correcting codes use finite fields and polynomial arithmetic to add redundancy, allowing a receiver to reconstruct data after some symbols are changed \cite{number_theory_ref2}.

Elliptic curves provide another example. An equation such as $y^2=x^3+ax+b$ defines a curve with a special addition rule for its points. The arithmetic of these points supports digital signatures and key exchange. The security is not based on multiplication alone but on the difficulty of reversing a point-multiplication operation on a suitably chosen curve \cite{number_theory_ref3}.

\section*{4. Important theorems and results}
\begin{enumerate}[leftmargin=*,itemsep=0.35em]
\item \textbf{Euclid's theorem on primes.} There are infinitely many primes. If $p_1,\ldots,p_k$ were all the primes, the number $p_1p_2\cdots p_k+1$ would have a prime factor not on the list.

\item \textbf{Fundamental theorem of arithmetic.} Every integer greater than $1$ has one prime factorization.

\item \textbf{Bézout's identity.} There exist integers $u,v$ such that $ua+vb=\gcd(a,b)$. This is the reason the extended Euclidean algorithm finds inverses.

\item \textbf{Chinese remainder theorem.} Compatible congruences with pairwise coprime moduli have one solution modulo the product of the moduli.

\item \textbf{Fermat's little theorem and Euler's theorem.} If $p$ is prime and $p$ does not divide $a$, then $a^{p-1}\equiv1\pmod p$. More generally, $a^{\varphi(n)}\equiv1\pmod n$ when $\gcd(a,n)=1$.

\item \textbf{Prime number theorem.} The proportion of integers near $x$ that are prime is approximately $1/\log x$. It describes average distribution, not an exact list.
\end{enumerate}
\noindent\emph{Source for these results:} \cite{number_theory_ref1}.

\theoryapplications
\theoryconnections{number_theory}{%
\item Integer arithmetic supplies divisibility and congruence, while algebra turns those rules into groups, rings, and fields that can be manipulated systematically.
\item Proof identifies which patterns are universal, and analysis becomes necessary when infinitely many integers are counted or primes are studied through generating functions.
\item For example, reducing modulo $m$ turns an infinite search for solutions into a finite search among residue classes, while a zeta function packages information about all primes into one analytic object \cite{number_theory_ref1,number_theory_ref2}.
}{%
\item Number theory helps coding theory by providing finite-field arithmetic in which errors can be detected and corrected.
\item Cryptography uses modular exponentiation together with problems such as factoring or discrete logarithms that are easy to perform in one direction but difficult to reverse.
\item Algebraic geometry and probability borrow arithmetic counts and congruence information, while mathematical physics uses lattices, modular forms, and spectral sums to organize discrete states \cite{number_theory_ref1,number_theory_ref2}.
}
\section*{7. Sample Questions}
\emph{Exercise reference:} \cite{number_theory_ref1}. The cited edition is the source for the exercise set; consult its Exercises section for the official statement and numbering.
\begin{enumerate}[leftmargin=*,itemsep=0.9em]
\item \textbf{Exercise 1. Compute $\gcd(252,105)$ and express it as a combination.} The Euclidean algorithm gives $252=2(105)+42$ and $105=2(42)+21$. Hence the gcd is $21$. Since $21=105-2(42)=105-2(252-2(105))=5(105)-2(252)$, one combination is $21=-2(252)+5(105)$.

\item \textbf{Exercise 2. Solve $14x\equiv8\pmod{22}$.} The gcd of $14$ and $22$ is $2$, which divides $8$. Divide by $2$ to get $7x\equiv4\pmod{11}$. Since $7^{-1}\equiv8\pmod{11}$, $x\equiv32\equiv10\pmod{11}$. The solutions modulo $22$ are $x\equiv10,21\pmod{22}$.

\item \textbf{Exercise 3. Solve $x\equiv2\pmod3$ and $x\equiv3\pmod5$.} Write $x=2+3k$. Then $2+3k\equiv3\pmod5$, so $3k\equiv1\pmod5$. The inverse of $3$ modulo $5$ is $2$, giving $k\equiv2\pmod5$ and $x\equiv8\pmod{15}$.

\item \textbf{Exercise 4. Prove that $\sqrt2$ is irrational.} Suppose $\sqrt2=a/b$ in lowest terms. Then $a^2=2b^2$, so $a$ is even, say $a=2c$. Substitution gives $b^2=2c^2$, so $b$ is even too, contradicting that $a/b$ was reduced. Thus no fraction equals $\sqrt2$.

\item \textbf{Exercise 5. Find the positive integer solutions of $3x+5y=1$.} Since $3(2)+5(-1)=1$, one solution is $(2,-1)$. The general solution is $x=2+5t$ and $y=-1-3t$ for integers $t$, because changing $x$ by $5$ and $y$ by $-3$ preserves the left side.

\end{enumerate}
\section*{8. References}
\begin{thebibliography}{9}
\bibitem{number_theory_ref1} G. H. Hardy and E. M. Wright, \emph{An Introduction to the Theory of Numbers}.
\bibitem{number_theory_ref2} J. H. Silverman, \emph{A Friendly Introduction to Number Theory}.
\bibitem{number_theory_ref3} J. Neukirch, \emph{Algebraic Number Theory}.
\end{thebibliography}


\theory{Obstruction theory}{How can we tell that a construction cannot be completed, even when every small piece looks possible?}{Obstruction theory turns a failed extension into a class that can be calculated. If the class is nonzero, completion is impossible; if it vanishes, the extension is often possible, at least up to continuous deformation. The theory therefore converts a vague failure into a precise cohomological invariant.}{Extending maps, constructing sections of bundles, deciding whether bundles are trivial, classifying manifolds, and the existence questions of surgery theory.}{Cohomology classes and obstruction groups identify the first stage at which a local construction cannot extend globally.}

\paragraph{The idea and its history}
Suppose a map has been defined on the edges of a triangulated shape. We would like to fill in each triangle, then each tetrahedron, and so on. A boundary may carry a twist that prevents the next filling. Obstruction theory records that twist instead of trying all possible fillings. The modern language grew from the work of Stiefel and Whitney on characteristic classes and from Samuel Eilenberg's systematic obstruction theory for homotopy \cite{obstruction_theory_ref1,obstruction_theory_ref2}.

The word obstruction is important: it is not merely evidence that a particular attempt failed. It is an invariant of the already-defined boundary data. Changing the construction without changing its homotopy class changes the representative but not the essential cohomology class. This is why a local calculation can give a global impossibility result \cite{obstruction_theory_ref3}.

\end{historylist}
\section*{3. Theory}
\subsection*{Core theory}
\paragraph{Extending a map skeleton by skeleton}
A simplicial complex is made from vertices, edges, triangles, tetrahedra, and higher-dimensional simplices. Its $n$-skeleton is the part made from pieces of dimension at most $n$. Let a continuous map from a complex $X$ to a space $Y$ already be defined on the $(n-1)$-skeleton. To extend it across an $n$-simplex, we must extend its boundary map from an $(n-1)$-sphere to the inside of the simplex.

The boundary map determines an element of the homotopy group $\pi_{n-1}(Y)$. If that element is nonzero, the boundary is wrapped around $Y$ in a way that cannot be filled. Assigning this element to every $n$-simplex gives an $n$-cochain. Compatibility on adjacent faces makes it a cocycle, so it determines a class in
\[
H^n(X;\pi_{n-1}(Y)).
\]
The class is the obstruction to extending the map. A zero class means that the map can be adjusted within its homotopy class so that extension is possible; a nonzero class proves that extension is impossible \cite{obstruction_theory_ref1}.

\paragraph{Concrete illustration: computing an obstruction on a 2-skeleton}
Let $X$ be a simplicial complex consisting of three vertices $v_0,v_1,v_2$, three edges $e_{01},e_{02},e_{12}$, and one triangle $\Delta$ with boundary $e_{01}+e_{12}-e_{02}$. Let $Y=S^1$ (the unit circle). We want to extend a map $f\colon X^{(1)}\to S^1$ from the 1-skeleton to the whole complex.

\textbf{Step 1: Define $f$ on vertices.} Set $f(v_0)=f(v_1)=f(v_2)=1\in S^1$.

\textbf{Step 2: Define $f$ on edges.} Let $f$ traverse the circle once around each edge: each edge maps by the identity loop $e^{2\pi i t}$ for $t\in[0,1]$. The winding number on every edge is $1$.

\textbf{Step 3: Compute the obstruction on $\Delta$.} The boundary of $\Delta$ maps into $S^1$ with total winding $1+1-1=1$. This nonzero element of $\pi_1(S^1)\cong\mathbb{Z}$ is the obstruction cocycle on $\Delta$. Its cohomology class in $H^2(X;\pi_1(S^1))$ is $1\neq 0$.

\textbf{Conclusion.} No continuous extension of $f$ to the interior of $\Delta$ exists. The obstruction computation took three concrete steps---assign values, sum winding numbers, check nonzero---that generalize to any dimension by replacing winding numbers with higher homotopy groups \cite{obstruction_theory_ref3}.

For a simple picture, a map from the circle to the circle may wind around once. The winding number is nonzero, so it cannot extend to a map from the disk to the circle without crossing the circle's missing center in an appropriate model. A map that winds zero times can be filled. The same reasoning survives in much more complicated spaces, where the winding number is replaced by a homotopy group and cohomology class.

\paragraph{Sections and bundles}
A bundle is a space that locally looks like a product: above each point of a base $B$ sits a fiber $F$. A section chooses one point in every fiber continuously. A section of a vector bundle is a continuous vector field; a nowhere-zero section is a vector field that never vanishes. The question is global because local choices may disagree when they are moved around loops.

Let $p:E\to B$ be a fibration with fiber $F$, and suppose a section has been built on the $n$-skeleton of $B$. On the boundary of each $(n+1)$-simplex, the section gives a map from an $n$-sphere into the fiber. Its homotopy class lies in $\pi_n(F)$. These classes form an obstruction cochain in
\[
C^{n+1}(B;\pi_n(F)).
\]
It is a cocycle, and its cohomology class in $H^{n+1}(B;\pi_n(F))$ is independent of the earlier choices up to a coboundary. If it is nonzero, no extension exists. Under the usual hypotheses, if it is zero, a homotopy section exists on the next skeleton \cite{obstruction_theory_ref2}.

Characteristic classes are important examples. The Stiefel--Whitney classes and Euler class can obstruct independent vector fields or nowhere-zero sections. The famous hairy-ball phenomenon says that a continuous tangent vector field on an even-dimensional sphere must vanish somewhere; in obstruction language, a characteristic class prevents a nowhere-zero section \cite{obstruction_theory_ref3}.

\paragraph{Geometry and surgery}
In geometric topology, obstruction theory asks whether a topological manifold admits a piecewise-linear structure, and whether a piecewise-linear manifold admits a differentiable structure. In dimensions at most two, and in dimension three by Moise's theorem, the relevant topological and piecewise-linear notions agree. Dimension four is exceptional: topological and smooth questions diverge sharply. In dimensions at most six, piecewise-linear and differentiable structures coincide in the standard comparison described in the subject \cite{obstruction_theory_ref4}.

Surgery theory asks whether a space satisfying $n$-dimensional Poincare duality is homotopy equivalent to a manifold, and whether a homotopy equivalence between manifolds can be changed into a diffeomorphism. In high dimensions there is first a primary topological $K$-theory obstruction to obtaining a suitable vector bundle and normal map. If it vanishes, a secondary obstruction in algebraic $L$-theory decides whether surgery can produce the desired homotopy equivalence \cite{obstruction_theory_ref4,obstruction_theory_ref5}.

\paragraph{How it is useful}
The method is useful whenever a problem is staged by dimension. In topology it tells us the first dimension in which a map, field, or structure fails. In bundle theory it separates local triviality from global triviality. In manifold theory it organizes a difficult classification into a sequence of computable tests. In motion planning, the same idea appears when a locally chosen motion cannot be patched into one continuous rule over all configurations.

\section*{4. Important theorems and results}
\begin{enumerate}[leftmargin=*,itemsep=0.35em]
\item \textbf{Eilenberg obstruction theorem.} The obstruction to extending a map across the $n$-skeleton is a class in $H^n(X;\pi_{n-1}(Y))$.

\item \textbf{Extension criterion.} A zero obstruction class permits extension after changing the map within its homotopy class; a nonzero class forbids extension.

\item \textbf{Bundle-section obstruction.} For a fibration with fiber $F$, extending a section across the $(n+1)$-skeleton produces a class in $H^{n+1}(B;\pi_n(F))$.

\item \textbf{Characteristic-class principle.} Characteristic classes measure the failure of a bundle to possess specified global fields or sections.

\item \textbf{Moise's theorem.} Every topological $3$-manifold has a unique piecewise-linear and smooth structure up to the appropriate equivalence.

\item \textbf{Surgery obstruction principle.} In high-dimensional manifold problems, a primary $K$-theory obstruction is followed, when it vanishes, by a secondary algebraic $L$-theory obstruction.
\end{enumerate}
\noindent\emph{Source for these results:} \cite{obstruction_theory_ref1}.

\theoryapplications
\theoryconnections{obstruction_theory}{%
\item Topology supplies the space and homotopy theory supplies the idea that a construction should be unchanged by continuous deformation.
\item A partial construction is extended one cell at a time; when an extension fails on the boundary of the next cell, the failure is recorded as a cohomology class with values in a homotopy group.
\item If that class vanishes, the construction can continue, so cohomology turns an existence question into a sequence of computable tests \cite{obstruction_theory_ref1,obstruction_theory_ref2}.
}{%
\item The method helps bundle theory decide whether a section or reduction of structure group exists, and helps Postnikov theory build a space from successive homotopy data.
\item In surgery and K- or L-theory, the same pattern tests whether algebraic data can be realized by an actual manifold or bundle.
\item The obstruction is therefore useful precisely because it identifies the first stage at which a proposed construction cannot be completed \cite{obstruction_theory_ref3,obstruction_theory_ref4}.
\item \textbf{Structural engineering and failure modes.} A loaded truss or frame is a network of beams (edges) joined at nodes (vertices). Each beam can carry stress in certain directions---the ``fiber'' of the local bundle. When an engineer tries to assign a consistent stress distribution across the whole structure, local solutions may fail to patch together globally. The obstruction cocycle measures exactly this: a nonzero class in $H^2(B;\pi_1(F))$ signals a load path that cannot close, predicting a failure mode such as torsional buckling or a mechanism (a floppy mode). For instance, a cantilevered roof frame with triangular bracing has trivial obstruction classes, confirming rigidity; removing one diagonal makes a class nonzero, identifying the precise panel where collapse initiates \cite{obstruction_theory_ref2}.
}

\section*{Glossary}
\begin{description}[leftmargin=*,style=nextline,itemsep=0.3em]
\item[\textbf{Obstruction.}] A cohomology class that measures whether a locally possible construction can be extended globally; if the class is nonzero, the extension is impossible.
\item[\textbf{Cohomology class.}] An element of a cohomology group that records global topological information encoded in local data on the cells of a space.
\item[\textbf{Homotopy group.}] An algebraic invariant $\pi_n(Y)$ measuring the ways an $n$-dimensional sphere can be mapped into a space $Y$; nonzero elements represent maps that cannot be continuously deformed to a constant map.
\item[\textbf{Characteristic class.}] A natural assignment of cohomology classes to vector bundles that detects global twisting; Stiefel--Whitney classes and the Euler class are important examples.
\item[\textbf{Fibration.}] A map $p:E\to B$ that locally looks like a product, where each point of the base $B$ has a fiber $F$ above it; obstructions to global sections of fibrations lie in cohomology groups.
\item[\textbf{Simplicial complex.}] A combinatorial space built from vertices, edges, triangles, and higher-dimensional simplices glued along their faces, providing a framework for computing obstructions skeleton by skeleton.
\item[\textbf{Obstruction cocycle.}] A cochain assigning to each simplex the homotopy class of its boundary map; automatically a cocycle.
\item[\textbf{$n$-skeleton.}] The union of all simplices of dimension at most $n$; obstructions are computed skeleton by skeleton.
\item[\textbf{Stiefel--Whitney class.}] A characteristic class with $\mathbb{Z}/2$ coefficients detecting nontrivial real vector bundles.
\item[\textbf{Normal map.}] A degree-one map between manifolds equipped with a lift of the tangent bundle to the stable normal bundle.
\end{description}

\section*{7. Sample Questions}
\emph{Exercise reference:} \cite{obstruction_theory_ref1}. The cited edition is the source for the exercise set; consult its Exercises section for the official statement and numbering.
\begin{enumerate}[leftmargin=*,itemsep=0.9em]
\item \textbf{Exercise 1.} Compute the Stiefel--Whitney classes $w_1$ and $w_2$ of the tangent bundle of the torus $T^2$.

\emph{Solution.} $T^2=S^1\times S^1$ is a product of orientable manifolds, so $w_1(T^2)=0$. The tangent bundle is trivial ($TT^2\cong T^2\times\mathbb{R}^2$), so all characteristic classes vanish: $w_1(T^2)=0$ and $w_2(T^2)=0$.

\item \textbf{Exercise 2.} The first Stiefel--Whitney class $w_1$ detects non-orientability. Compute $w_1$ of the Klein bottle and use it to show the Klein bottle is non-orientable.

\emph{Solution.} The Klein bottle is the non-trivial $S^1$-bundle over $S^1$, which is the mapping torus of the reflection $S^1\to S^1$. Its tangent bundle satisfies $w_1\neq 0$ because the Klein bottle is non-orientable (it contains orientation-reversing loops). Specifically, $w_1$ is the non-zero element of $H^1(K;\mathbb{Z}/2)\cong\mathbb{Z}/2$, detected by the loop that reverses orientation.

\item \textbf{Exercise 3.} Consider the simplicial complex $X$ with vertices $v_0,v_1,v_2$, edges $e_{01},e_{12},e_{02}$, and one triangle $\Delta$. Define a map $f\colon X^{(1)}\to S^1$ with winding number $1$ on each edge. Compute the obstruction cocycle on $\Delta$.

\emph{Solution.} The boundary of $\Delta$ has oriented edges $e_{01}+e_{12}-e_{02}$. The winding numbers sum to $1+1-1=1$. The obstruction to extending $f$ across $\Delta$ is this element of $\pi_1(S^1)\cong\mathbb{Z}$, namely $1$. Since $1\neq 0$, the map cannot be extended to the interior of $\Delta$.

\item \textbf{Exercise 4.} Let $f\colon S^1\to S^1$ be the identity map (winding number $1$). Show that $f$ cannot extend to a continuous map $D^2\to S^1$.

\emph{Solution.} If $F\colon D^2\to S^1$ extended $f$, then $F$ would be a null-homotopic map $S^1\to S^1$ (since $D^2$ is contractible). But the identity map has winding number $1\neq 0$, so it is not null-homotopic. Contradiction.

\item \textbf{Exercise 5.} Every tangent bundle of a compact surface has Euler class $e\in H^2(M;\mathbb{Z})$. For the sphere $S^2$, compute $e(TS^2)$ and interpret it as an obstruction.

\emph{Solution.} The Euler class of $TS^2$ satisfies $\langle e(TS^2),[S^2]\rangle=\chi(S^2)=2$. This means $e(TS^2)=2\in H^2(S^2;\mathbb{Z})\cong\mathbb{Z}$. Since $e\neq 0$, the tangent bundle of $S^2$ has no nowhere-zero section: every continuous vector field on $S^2$ must vanish somewhere (the hairy-ball theorem).

\end{enumerate}
\section*{8. References}
\begin{thebibliography}{9}
\bibitem{obstruction_theory_ref1} G. W. Whitehead, \emph{Elements of Homotopy Theory}, Springer, 1978.
\bibitem{obstruction_theory_ref2} N. Steenrod, \emph{The Topology of Fibre Bundles}, Princeton University Press, 1951.
\bibitem{obstruction_theory_ref3} A. Hatcher, \emph{Algebraic Topology}, Cambridge University Press, 2002.
\bibitem{obstruction_theory_ref4} A. Ranicki, \emph{Algebraic and Geometric Surgery}, Oxford University Press, 2002.
\bibitem{obstruction_theory_ref5} C. T. C. Wall, \emph{Surgery on Compact Manifolds}, 2nd edition, American Mathematical Society, 1999.
\end{thebibliography}


\theory{Operator theory}{How can we study transformations of functions, signals, and infinite lists as carefully as we study numbers?}{An operator is a rule that takes one object to another, often a function to a function. Operator theory studies linear and nonlinear operators, their domains, inverses, spectra, and natural modes. It extends matrix algebra to infinite-dimensional spaces, where convergence and the choice of domain become part of the problem.}{Differential and integral equations, quantum mechanics, Fourier analysis, signal processing, control, inverse problems, and the mathematics of operator algebras.}{Domains, spectra, resolvents, and functional calculus extend matrix methods to transformations of infinite-dimensional spaces.}

\paragraph{The problem and its history}
Matrices act on finite lists of numbers. Differentiation, integration, shifting a sequence, and filtering a signal act on functions or infinite lists. These rules are too large to be represented by a finite matrix, but they still have structure. Operator theory asks which rules are continuous, which can be inverted, how errors grow, and whether the rule has simple modes \cite{operator_theory_ref1}.

The subject began with differential and integral operators arising in differential equations. Hilbert's study of integral equations introduced infinite-dimensional analogues of Euclidean space. Banach and von Neumann developed the language of normed and Hilbert spaces, while spectral theory connected operators to eigenvalues. The resulting theory belongs to functional analysis and depends heavily on the topology of function spaces \cite{operator_theory_ref2,operator_theory_ref3}.

\end{historylist}
\section*{3. Theory}
\subsection*{Core theory}
\paragraph{Operators and their domains}
Let $X$ and $Y$ be vector spaces. A linear operator $T:X\to Y$ satisfies
\[
T(ax+by)=aT(x)+bT(y).
\]
The differentiation operator $D$ sends a polynomial $x^n$ to $nx^{n-1}$. Its kernel consists of constant polynomials, and its range on polynomials is all polynomials. An integral operator may have the form
\[
(Tf)(x)=\int_a^b K(x,t)f(t)\,dt,
\]
where the kernel $K$ describes how input values are combined.

In finite dimensions every linear map is defined on the whole space and is continuous. For an unbounded operator such as differentiation on a space of square-integrable functions, the domain must be specified. A function may be square-integrable while its derivative is not. Ignoring this point can make a formally correct equation meaningless. Closed operators and densely defined operators provide the setting in which limits and adjoints can be handled safely \cite{operator_theory_ref1,operator_theory_ref2}.

\paragraph{Spectrum and natural modes}
For a matrix $A$, an eigenvector $v$ satisfies $Av=\lambda v$. The corresponding $\lambda$ is a value at which $A-\lambda I$ fails to be invertible. In an infinite-dimensional space, the spectrum of $T$ is the set of complex numbers $\lambda$ for which $T-\lambda I$ has no bounded inverse. The spectrum may contain values that are not ordinary eigenvalues; this is one reason infinite-dimensional analysis is richer than matrix algebra.

The spectral theorem says, broadly, that self-adjoint and more generally normal operators can be represented using their spectral data. In finite dimensions a normal matrix is unitarily diagonalizable:
\[
A=UDU^*,
\]
where the columns of $U$ are orthonormal eigenvectors and $D$ is diagonal. For an infinite-dimensional self-adjoint operator, a sum over eigenvectors may be replaced by an integral over a spectral measure. The operator then behaves like multiplication by the variable, the simplest possible model \cite{operator_theory_ref1,operator_theory_ref4}.

An operator $N$ on a complex Hilbert space is normal when $NN^*=N^*N$. Unitary operators preserve lengths, self-adjoint operators are the infinite-dimensional analogue of Hermitian matrices, and positive operators have the form $M^*M$. These classes are important because the spectral theorem applies to them and because they model reversible motion, measured quantities, and energy \cite{operator_theory_ref1}.

\paragraph{Polar decomposition}
Every bounded operator $A$ between complex Hilbert spaces has a canonical factorization
\[
A=UP,\qquad P=(A^*A)^{1/2},
\]
where $P$ is positive and $U$ is a partial isometry. For a square invertible matrix, this says that $A$ is a rotation or reflection followed by a stretch. It is the operator version of writing a complex number as magnitude times direction, and it is closely related to the singular value decomposition.

The word partial matters in infinite dimensions. The unilateral shift on sequences moves every entry one place and inserts a zero. It preserves length on its initial space but is not onto, so it is not unitary. The polar factor records exactly this loss of surjectivity \cite{operator_theory_ref3}.

\paragraph{Complex analysis and operator algebras}
Many operators act on spaces of holomorphic functions. On the Hardy space, multiplication by the independent variable gives the unilateral shift. Beurling's theorem describes its invariant subspaces by inner functions. Toeplitz operators multiply a function and then project it back into the Hardy space; similar questions arise on Bergman spaces \cite{operator_theory_ref5}.

An operator algebra is a collection of operators closed under addition, scalar multiplication, and composition. A $C^*$-algebra is a complex Banach algebra equipped with an adjoint operation satisfying
\[
(xy)^*=y^*x^*,\qquad (x^*)^*=x,\qquad \|x^*x\|=\|x\|^2.
\]
The norm identity is powerful: it ties algebraic invertibility, geometry, and spectral radius together. Commutative $C^*$-algebras can be understood through spaces of functions, while noncommutative ones provide an algebraic language for quantum observables \cite{operator_theory_ref4}.

\paragraph{How the theory solves applications}
For a vibrating string, separation of variables turns the wave equation into an eigenvalue problem for a differential operator. The eigenfunctions are normal modes; the observed motion is a combination of them. For a quantum particle, a self-adjoint operator represents an observable, and its spectrum describes the possible measurement values. The spectral theorem explains how probabilities are assigned to spectral regions \cite{operator_theory_ref1}.

In signal processing, a filter is an operator. Fourier modes are especially useful because convolution operators act on them by multiplication. Instead of tracking every point of a signal, one studies how the operator changes each frequency. In control theory, the spectrum of a system operator helps determine stability; in inverse problems, compact operators explain why recovering an input from measured output may amplify noise \cite{operator_theory_ref2,operator_theory_ref3}.

\section*{4. Important theorems and results}
\begin{enumerate}[leftmargin=*,itemsep=0.35em]
\item \textbf{Spectral theorem.} Self-adjoint and normal operators on Hilbert spaces admit a spectral representation; finite-dimensional versions are unitary diagonalization.

\item \textbf{Schur decomposition.} Every complex matrix is unitarily similar to an upper triangular matrix. A normal upper triangular matrix is diagonal, yielding the normal-matrix spectral theorem.

\item \textbf{Polar decomposition.} Every bounded operator has $A=U(A^*A)^{1/2}$ with $U$ a partial isometry, uniquely determined on the relevant range.

\item \textbf{Douglas factorization lemma.} If $A^*A\leq B^*B$, then $A=CB$ for a contraction $C$. This gives a clean route to polar decomposition.

\item \textbf{Beurling's theorem.} The invariant subspaces of the unilateral shift on the Hardy space are generated by inner functions.

\item \textbf{C*-identity and spectral radius.} In a $C^*$-algebra the norm is determined by the algebraic structure through the spectrum of $x^*x$. This is why $C^*$-algebras are rigid analytic objects.
\end{enumerate}
\noindent\emph{Source for these results:} \cite{operator_theory_ref1}.

\theoryapplications
\theoryconnections{operator_theory}{%
\item Linear algebra studies transformations of finite-dimensional vectors; operator theory extends the same question to spaces of functions, where convergence and integration must be controlled.
\item Topology of function spaces says which limits are allowed, measure defines the integral norms, and complex analysis explains resolvents and spectra.
\item For example, differentiating a function is an operator, and asking for its eigenfunctions is asking which shapes reproduce themselves under differentiation or under a boundary-value problem \cite{operator_theory_ref1,operator_theory_ref2}.
}{%
\item Operator theory helps differential equations by converting an equation into a statement about an operator's inverse or spectrum.
\item Fourier analysis diagonalizes important translation-invariant operators into frequencies, while quantum mechanics interprets observables as operators and their spectra as possible measurement values.
\item In control and inverse problems, boundedness and spectral location tell whether errors grow and whether hidden inputs can be recovered \cite{operator_theory_ref1,operator_theory_ref4}.
}

\section*{Glossary}
\begin{description}[leftmargin=*,style=nextline,itemsep=0.3em]
\item[\textbf{Operator.}] A rule that maps elements of one vector space (often a function space) to another, generalizing the idea of a matrix to infinite dimensions.
\item[\textbf{Spectrum.}] The set of complex numbers $\lambda$ for which $T-\lambda I$ is not invertible; it generalizes eigenvalues to infinite-dimensional spaces and may contain more than just point eigenvalues.
\item[\textbf{Self-adjoint operator.}] An operator equal to its own adjoint ($T=T^*$), which on a Hilbert space is the infinite-dimensional analogue of a real symmetric matrix; it has a real spectrum and admits a spectral decomposition.
\item[\textbf{Hilbert space.}] A complete vector space with an inner product, providing the infinite-dimensional setting where operators, spectra, and the spectral theorem can be developed rigorously.
\item[\textbf{Resolvent.}] The family of inverses $(T-\lambda I)^{-1}$ for $\lambda$ outside the spectrum; its analytic properties encode the spectral structure of the operator.
\item[\textbf{Polar decomposition.}] A factorization $A=UP$ of a bounded operator into a partial isometry $U$ and a positive operator $P$, analogous to writing a complex number as magnitude times direction.
\item[\textbf{Bounded operator.}] A linear operator with finite operator norm, ensuring continuity.
\item[\textbf{Unbounded operator.}] A linear operator that may not be defined on the whole space; requires careful domain specification.
\item[\textbf{Adjoint.}] For an operator $T$, the operator $T^*$ satisfying $\langle Tx,y\rangle=\langle x,T^*y\rangle$; self-adjointness is key to the spectral theorem.
\item[\textbf{C*-algebra.}] A complex Banach algebra with an adjoint operation satisfying $\|x^*x\|=\|x\|^2$.
\item[\textbf{Spectral theorem.}] Self-adjoint and normal operators admit a decomposition using spectral measures.
\item[\textbf{Hardy space.}] The space $H^2$ of holomorphic functions on the unit disk with bounded Taylor coefficients.
\end{description}

\section*{7. Sample Questions}
\emph{Exercise reference:} \cite{operator_theory_ref1}. The cited edition is the source for the exercise set; consult its Exercises section for the official statement and numbering.
\begin{enumerate}[leftmargin=*,itemsep=0.9em]
\item \textbf{Exercise 1. Find the eigenvalues and an orthonormal eigenbasis of $A=\begin{pmatrix}2&1\\1&2\end{pmatrix}$.} The vectors $(1,1)$ and $(1,-1)$ give eigenvalues $3$ and $1$. After dividing each by $\sqrt2$, they form an orthonormal basis, so $A$ is unitarily diagonalizable.

\item \textbf{Exercise 2. What is the kernel of differentiation on polynomials?} If $Dp=0$, every coefficient of positive degree is zero, so $p$ is constant. Thus $\ker D$ is the one-dimensional space of constants.

\item \textbf{Exercise 3. Compute the polar decomposition of $A=\begin{pmatrix}0&1\\0&0\end{pmatrix}$.} We have $A^*A=\begin{pmatrix}0&0\\0&1\end{pmatrix}$, so $P=(A^*A)^{1/2}=\operatorname{diag}(0,1)$. The partial isometry $U$ sends $(0,1)$ to $(1,0)$ and is zero on $(1,0)$; then $UP=A$.

\item \textbf{Exercise 4. Why is the differentiation operator unbounded on $L^2[0,1]$?} Take $f_n(x)=\sin(2\pi n x)$. Its $L^2$ norm stays fixed while the derivative has norm proportional to $n$. No single constant bounds $\|Df\|$ by $\|f\|$, so differentiation is not bounded on this space.

\item \textbf{Exercise 5. Explain the Fourier reason that convolution becomes multiplication.} A Fourier mode $e^{i\omega x}$ is an eigenfunction of translation. Convolution with a fixed kernel preserves each frequency and multiplies its amplitude by the kernel's Fourier transform. Therefore filtering can be performed frequency by frequency.

\end{enumerate}
\section*{8. References}
\begin{thebibliography}{9}
\bibitem{operator_theory_ref1} J. B. Conway, \emph{A Course in Functional Analysis}, 2nd ed., Springer, 1990.
\bibitem{operator_theory_ref2} P. R. Halmos, \emph{Introduction to Hilbert Space and the Theory of Spectral Multiplicity}, Chelsea, 1951.
\bibitem{operator_theory_ref3} J. B. Conway, \emph{A Course in Operator Theory}, American Mathematical Society, 2000.
\bibitem{operator_theory_ref4} G. K. Pedersen, \emph{C*-Algebras and Their Automorphism Groups}, Academic Press, 1979.
\bibitem{operator_theory_ref5} K. Hoffman, \emph{Banach Spaces of Analytic Functions}, Prentice-Hall, 1962.
\end{thebibliography}


\theory{Order theory}{How can we describe precedence, containment, or preference when some pairs can be compared and others cannot?}{Order theory formalizes statements such as less than, comes before, is contained in, or is more specialized than. It studies partially ordered sets, bounds, lattices, order-preserving maps, chains, antichains, and the duality obtained by reversing an order.}{Scheduling, dependency graphs, databases, optimization, type systems, logic, lattice-based algorithms, and the organization of mathematical structures.}{Partial orders, chains, antichains, and lattice operations describe how objects can be compared and combined without numerical coordinates.}

\paragraph{The question and its history}
The everyday order on numbers compares every pair: either $2<3$, $3<2$, or they are equal. But the subset relation is different. The set of birds and the set of dogs are both subsets of animals, but neither contains the other. Order theory was created to study both situations with one language \cite{order_theory_ref1}.

George Boole, Charles Sanders Peirce, Richard Dedekind, and Ernst Schröder developed early algebraic descriptions of order. Pasch, Peano, Hilbert, and Veblen developed axioms for order in geometry. Bertrand Russell examined order and series in the foundations of mathematics. The modern subject grew because the same patterns appeared in algebra, logic, analysis, geometry, and computer science \cite{order_theory_ref1,order_theory_ref2}.

\end{historylist}
\section*{3. Theory}
\subsection*{Core theory}
\paragraph{Relations and ordered sets}
A binary relation $\leq$ on a set $P$ is a partial order when it is reflexive ($a\leq a$), antisymmetric (if $a\leq b$ and $b\leq a$, then $a=b$), and transitive (if $a\leq b$ and $b\leq c$, then $a\leq c$). A set with such a relation is a partially ordered set, or poset. A total order additionally requires every pair to be comparable.

The divisibility relation on positive integers is a poset: $a\leq b$ means $a$ divides $b$. The numbers $2$ and $3$ are incomparable, while both divide $6$. A Hasse diagram draws only the immediate comparisons, so it reveals the structure without clutter from transitivity. The subset relation on a power set gives another important example \cite{order_theory_ref1}.

An element can be maximal without being greatest. In the set $\{2,3\}$ under divisibility, both elements are maximal, but neither is greatest. This distinction matters in optimization: a solution can be impossible to improve in its own direction while another incomparable solution is also optimal in a different direction.

\paragraph{Bounds, lattices, and duality}
An upper bound of $S\subseteq P$ is an element above every member of $S$. The least upper bound, or supremum, is written $\sup S$ or $\bigvee S$. A lower bound and greatest lower bound are written $\inf S$ or $\bigwedge S$. For subsets ordered by inclusion, union is the least upper bound and intersection is the greatest lower bound. For positive integers ordered by divisibility, the least common multiple is the join and the greatest common divisor is the meet \cite{order_theory_ref1}.

Reversing every comparison produces the dual order. Every definition has a dual: upper becomes lower, least becomes greatest, join becomes meet, and chain becomes chain in the reversed picture. A theorem in one order immediately gives a dual theorem. This is a powerful economy of thought because the same proof often works after turning the Hasse diagram upside down \cite{order_theory_ref2}.

A lattice is a poset in which every pair has a meet and a join. A complete lattice has these bounds for every subset. A directed-complete partial order guarantees suprema of directed subsets and is central in domain theory, where a program is approximated by increasingly informative partial computations. Distributive lattices satisfy
\[
x\wedge(y\vee z)=(x\wedge y)\vee(x\wedge z).
\]
Boolean algebras add complementation and model logical operations and digital circuits; Heyting algebras model constructive logic \cite{order_theory_ref3}.

\paragraph{Functions between orders}
A function $f:P\to Q$ is monotone when $a\leq b$ implies $f(a)\leq f(b)$. It is order-reflecting when the reverse implication holds, and antitone when it reverses the order. An order embedding preserves and reflects comparisons. An order isomorphism is a bijective order-preserving map whose inverse is also order-preserving.

The successor function on natural numbers is monotone. The complement map on a power set is antitone: if $A\subseteq B$, then $B^c\subseteq A^c$. A closure operator is monotone, extensive ($x\leq f(x)$), and idempotent ($f(f(x))=f(x)$). Examples include taking the span of a set of vectors, the topological closure of a set, and the set of all logical consequences of assumptions \cite{order_theory_ref3}.

A Galois connection is a pair of monotone maps between two ordered sets that are almost inverse in a precise inequality sense. Such pairs explain why taking a set of objects and taking the properties common to those objects are closely related operations. They occur in formal concept analysis, algebra, and the correspondence between subgroups and fixed fields \cite{order_theory_ref4}.

\paragraph{Special subsets and constructions}
An upper set contains every element above each of its members; a lower set contains everything below. An ideal is a directed lower set, and a filter is its dual. A chain is a subset in which every two elements are comparable. An antichain has no two comparable elements. These concepts are essential in combinatorics, especially when a large poset is decomposed into chains or antichains \cite{order_theory_ref1}.

New orders can be built from old ones. The product order on $P\times Q$ is defined by $(p,q)\leq(p',q')$ exactly when $p\leq p'$ and $q\leq q'$. A disjoint union keeps the original comparisons but adds no comparisons between the two pieces. A strict order is obtained by writing $a<b$ when $a\leq b$ and $a\ne b$.

\paragraph{How the theory solves practical problems}
In project scheduling, a task is below another task when it must be completed first. A topological sort lists tasks in an order compatible with all dependencies. If two tasks are incomparable, they can be done in parallel. The existence of minimal available tasks is an order-theoretic fact, while the resulting algorithm is used in build systems and course planning.

In type systems, a type may be below another when every value of the first can safely be used where the second is expected. Monotone functions preserve safe substitutions. In databases, sets of records ordered by inclusion form lattices; union combines information and intersection keeps common information. In optimization with several objectives, the Pareto-optimal points form the maximal elements of a partial order, so order theory explains why there need not be one best choice \cite{order_theory_ref3}.

\section*{4. Important theorems and results}
\begin{enumerate}[leftmargin=*,itemsep=0.35em]
\item \textbf{Zorn's lemma.} If every chain in a nonempty poset has an upper bound, the poset has a maximal element. It is equivalent to the axiom of choice and is used to prove the existence of bases and maximal ideals.

\item \textbf{Knaster--Tarski fixed-point theorem.} Every monotone map from a complete lattice to itself has a least fixed point and a greatest fixed point.

\item \textbf{Dilworth's theorem.} In a finite poset, the largest antichain has the same size as the smallest number of chains needed to cover the poset.

\item \textbf{Birkhoff representation theorem.} Every finite distributive lattice is isomorphic to the lattice of lower sets of a finite poset.

\item \textbf{Tarski's fixed-point principle.} Monotonicity on a complete lattice guarantees fixed points, turning repeated approximation into a rigorous existence method.

\item \textbf{Galois correspondence.} In suitable algebraic situations, inclusion-reversing correspondences between subobjects and substructures are organized by a Galois connection.
\end{enumerate}
\noindent\emph{Source for these results:} \cite{order_theory_ref1}.

\theoryapplications
\theoryconnections{order_theory}{%
\item Sets provide the collection of objects and relations say when one object precedes, contains, or is more informative than another.
\item Transitivity and antisymmetry make comparisons consistent without requiring every pair to be comparable.
\item When two objects have a least common upper bound or greatest common lower bound, lattice operations combine information; logic and proof use these operations to track implication and approximation \cite{order_theory_ref1,order_theory_ref2}.
}{%
\item Order theory helps optimization because feasible solutions can be compared and a monotone improvement process can be shown to have a least or greatest fixed point.
\item In databases, an order can mean “contains at least this information,” so joins combine records without losing consistency.
\item In computer science and type theory, recursive definitions are interpreted as fixed points of ordered maps, while combinatorics uses chains and antichains to bound how a finite family can be arranged \cite{order_theory_ref1,order_theory_ref3}.
}
\section*{7. Sample Questions}
\emph{Exercise reference:} \cite{order_theory_ref1}. The cited edition is the source for the exercise set; consult its Exercises section for the official statement and numbering.
\begin{enumerate}[leftmargin=*,itemsep=0.9em]
\item \textbf{Exercise 1. Find the meet and join of $12$ and $18$ in the divisibility order.} The meet is the greatest common divisor, $\gcd(12,18)=6$. The join is the least common multiple, $\operatorname{lcm}(12,18)=36$.

\item \textbf{Exercise 2. Give two maximal elements that have no greatest element.} In the poset $\{2,3\}$ ordered by divisibility, $2$ and $3$ are both maximal because neither has a larger member in the set divisible by it. There is no greatest element because neither divides the other.

\item \textbf{Exercise 3. Prove that the complement map on a power set is antitone.} If $A\subseteq B$, every element not in $B$ is also not in $A$. Therefore $B^c\subseteq A^c$, which is exactly order reversal.

\item \textbf{Exercise 4. Find the closure of $\{1,2\}$ under taking divisors in the positive integers.} A divisor-closure contains every positive divisor of every selected number. The divisors of $1$ and $2$ are $1$ and $2$, so the closure is $\{1,2\}$.

\item \textbf{Exercise 5. Use monotonicity to find a fixed point.} On the lattice $\mathcal P(\{a,b\})$, define $f(S)=S\cup\{a\}$. Starting at the bottom gives $\varnothing,\{a\},\{a\},\ldots$. Thus the least fixed point is $\{a\}$. The other fixed point is also $\{a\}$, because every set containing $a$ is sent to itself only when it contains no additional change; here $f(\{a,b\})=\{a,b\}$ as well, so the greatest fixed point is $\{a,b\}$.

\end{enumerate}
\section*{8. References}
\begin{thebibliography}{9}
\bibitem{order_theory_ref1} G. Grätzer, \emph{Lattice Theory: Foundation}.
\bibitem{order_theory_ref2} B. A. Davey and H. A. Priestley, \emph{Introduction to Lattices and Order}.
\bibitem{order_theory_ref3} G. Grätzer, \emph{General Lattice Theory}.
\bibitem{order_theory_ref4} G. Birkhoff, \emph{Lattice Theory}.
\end{thebibliography}


\theory{Percolation theory}{When a network is partly blocked or partly damaged, when does a large connected route suddenly appear or disappear?}{Percolation theory studies random connectivity. Sites or links are declared open with a probability, and the theory asks whether open paths form a cluster spanning a region or continuing forever. A critical probability separates a disconnected phase from a connected phase.}{Porous materials, epidemics, ecological habitats, power and transport networks, random graphs, polymer gelation, and the fragmentation of virus shells.}{Random open clusters and connectivity thresholds turn local occupation probabilities into global phase transitions.}

\paragraph{The physical question and history}
Imagine pouring liquid onto porous rock. Each microscopic passage may be open or blocked. If open passages form a connected route, liquid can travel from the top to the bottom. Percolation theory replaces the complicated material by a lattice or graph and studies the probability that a path exists. The model was introduced by Simon Broadbent and John Hammersley in 1957, after questions about gases and liquids moving through coal and other porous materials \cite{percolation_ref1,percolation_ref2}.

\end{historylist}
\section*{3. Theory}
\subsection*{Core theory}
\paragraph{Sites, bonds, and the critical point}
In bond percolation, each edge of a graph is open with probability $p$ and closed with probability $1-p$, independently. In site percolation, each vertex is occupied with probability $p$, and an edge can be used only when its endpoints are occupied. The open cluster of a vertex is the set of vertices reachable by open paths.

\[
p_c=\frac12.
\]

\paragraph{Phases and observables}
Below the threshold, open clusters are typically finite and connection probabilities decay rapidly. Above the threshold, a giant or infinite cluster appears while finite clusters remain around it. In two dimensions, duality relates open paths to closed paths in the dual lattice and compares the two phases \cite{percolation_ref1}.

Important observables include the probability that a site belongs to the infinite cluster, the mean cluster size, the correlation length, and the probability of a spanning path across a finite box. Near criticality these quantities often follow power laws such as $\xi\sim|p-p_c|^{-\nu}$. The exponent depends mainly on dimension rather than microscopic details, an idea called universality. Critical clusters are often fractal \cite{percolation_ref2}.

\paragraph{Different models}
Directed percolation allows paths to move in a preferred direction such as time and is connected with the contact process. Bootstrap percolation activates or removes sites according to the number of active neighbors. First-passage percolation assigns random travel times to edges and studies the fastest route. Invasion percolation grows a cluster by choosing the least costly boundary edge. The Fortuin--Kasteleyn random-cluster model connects percolation with the Ising and Potts models \cite{percolation_ref3}.

\paragraph{How it solves real problems}
In epidemiology, people or animals are vertices and possible transmissions are edges. Below the transmission threshold, outbreaks usually die out in small clusters; above it, a large outbreak becomes possible. Vaccination or reduced contact can lower the effective connectivity and move a population below the threshold.

In ecology, habitat patches are sites and movement corridors are bonds. A species can persist over a region only if enough patches remain connected. In materials science, enough chemical links connect small polymer molecules into a system-spanning gel. In virology, randomly removing capsid subunits can cross a fragmentation threshold and break a virus shell \cite{percolation_ref4}.

\section*{4. Important theorems and results}
\begin{enumerate}[leftmargin=*,itemsep=0.35em]
\item \textbf{Existence of a critical threshold.} Standard infinite lattices have a value $p_c$ separating the phase with no infinite open cluster from the phase with one.

\item \textbf{Kolmogorov zero--one law.} Translation-invariant tail events such as the existence of an infinite cluster have probability zero or one.

\item \textbf{Kesten's square-lattice result.} Bond percolation on the planar square lattice has critical probability $p_c=1/2$.

\item \textbf{Duality principle.} In two dimensions, open crossings and closed crossings in the dual lattice constrain one another.

\item \textbf{Sharpness of the transition.} In standard models, connection probabilities decay exponentially below the threshold, while a macroscopic cluster appears above it.
\end{enumerate}
\noindent\emph{Source for these results:} \cite{percolation_ref1}.

\theoryapplications
\theoryconnections{percolation_theory}{%
\item Probability assigns independent or correlated open/closed states to edges or sites, while graph theory defines paths and connected components.
\item The central question is whether local random links create a global crossing or infinite cluster.
\item Statistical physics explains why a sharp threshold can appear: near a critical density, clusters occur on many scales, so a small change in the probability can change the large-scale behavior \cite{percolation_ref1,percolation_ref2}.
}{%
\item Percolation gives network reliability a concrete calculation: a communication network works when its open edges contain a path between two terminals.
\item The same cluster threshold models whether a disease can spread through a contact graph, whether fluid can pass through a porous material, or whether a random graph contains a giant component.
\item The theory supplies the bridge from local failure probabilities to a global connectivity prediction \cite{percolation_ref1,percolation_ref3}.
}
\section*{7. Sample Questions}
\emph{Exercise reference:} \cite{percolation_ref1}. The cited edition is the source for the exercise set; consult its Exercises section for the official statement and numbering.
\begin{enumerate}[leftmargin=*,itemsep=0.9em]
\item \textbf{Exercise 1. What is the expected number of open edges among $100$ independent edges when $p=0.3$?} Linearity of expectation gives $100(0.3)=30$.

\item \textbf{Exercise 2. Why does a square-lattice bond model with $p=0.4$ lie below its threshold?} Its threshold is $1/2$, and $0.4<1/2$, so an infinite open cluster does not occur in the infinite-lattice limit.

\item \textbf{Exercise 3. For a regular tree of degree $4$, what is the threshold?} The formula $p_c=1/(z-1)$ gives $p_c=1/3$.

\item \textbf{Exercise 4. Explain why site and bond percolation are different.} Bond percolation removes connections while keeping vertices. Site percolation removes vertices and all their incident connections, so the same probability can give different connectivity.

\item \textbf{Exercise 5. Why can reducing contacts reduce a large outbreak?} Removing contacts removes bonds and lowers connectivity. If the network crosses below its threshold, large outbreaks become unlikely even though small outbreaks remain possible.

\end{enumerate}
\section*{8. References}
\begin{thebibliography}{9}
\bibitem{percolation_ref1} G. Grimmett, \emph{Percolation}.
\bibitem{percolation_ref2} D. Stauffer and A. Aharony, \emph{Introduction to Percolation Theory}.
\bibitem{percolation_ref3} H. Kesten, \emph{Percolation Theory for Mathematicians}.
\bibitem{percolation_ref4} M. Sahini and M. Sahimi, \emph{Applications of Percolation Theory}.
\end{thebibliography}


\theory{Perturbation theory}{How can we obtain useful answers when exact equations are too difficult, but are close to equations we can solve?}{Perturbation theory separates a solvable problem from a small correction. The unknown solution is expanded in powers of a small parameter, and the equations for successive corrections are solved one at a time.}{Celestial mechanics, quantum mechanics, fluid dynamics, waves, engineering, climate models, numerical analysis, and quantum field theory.}{A small parameter separates a solvable leading problem from corrections and supplies an ordered estimate of the error.}

\paragraph{The idea and history}
The two-body gravitational problem has exact conic-section orbits, but the motion of three bodies does not generally have a simple closed formula. Astronomers treated the influence of the third body as a perturbation of the two-body solution. Lagrange and Laplace developed systematic methods, and Le Verrier used deviations in Uranus's motion to predict Neptune \cite{perturbation_ref1,perturbation_ref2}.

The same strategy entered quantum mechanics. Dirac developed time-dependent perturbation ideas, and physicists used corrections to calculate atomic energy shifts. Quantum field theory organizes repeated perturbative terms with Feynman diagrams \cite{perturbation_ref3}.

\end{historylist}
\section*{3. Theory}
\subsection*{Core theory}
\paragraph{Regular perturbation}
Suppose a problem depends on a small number $\varepsilon$:
\[
F(u,\varepsilon)=0.
\]
We seek
\[
u=u_0+\varepsilon u_1+\varepsilon^2u_2+\cdots.
\]
The zeroth-order equation gives $u_0$. Substitution and comparison of powers of $\varepsilon$ gives an equation for $u_1$, then one for $u_2$. Truncating after the first correction gives $u\approx u_0+\varepsilon u_1$.

For $x^2+\varepsilon x-1=0$, put $x=x_0+\varepsilon x_1+\cdots$. At order zero choose $x_0=1$. At first order, $2x_0x_1+x_0=0$, giving $x_1=-1/2$. Thus the positive root is $x\approx1-\varepsilon/2$ for small $\varepsilon$ \cite{perturbation_ref1}.

\paragraph{Orders and singular problems}
First-order perturbation keeps one correction, second-order keeps two, and so on. Higher order is not automatically better: the series may diverge, or a small parameter may multiply the highest derivative and create a boundary layer. Such a problem is singular because setting the parameter to zero changes the order of the equation.

Degeneracy occurs when the unperturbed problem has several independent states with the same value. The perturbation can mix those states, so one must first diagonalize the perturbing part within the degenerate subspace. Multiple scales, matched asymptotic expansions, and dominant balance handle singular situations \cite{perturbation_ref2}.

\paragraph{A mechanical example}
For a weakly nonlinear oscillator,
\[
x''+x+\varepsilon x^3=0,
\]
the unperturbed motion is $x_0=A\cos t$. A direct correction can contain a term proportional to $t\sin t$, which grows even though $\varepsilon$ is small. This secular term signals that the frequency has shifted. The method of multiple scales introduces a slow time $T=\varepsilon t$ and allows the phase to vary slowly, removing the artificial growth.

\paragraph{Quantum and celestial applications}
In time-independent quantum mechanics,
\[
H=H_0+\varepsilon V.
\]
If $H_0\psi_n^{(0)}=E_n^{(0)}\psi_n^{(0)}$, then for a nondegenerate level the first energy correction is
\[
E_n^{(1)}=\langle\psi_n^{(0)},V\psi_n^{(0)}\rangle.
\]
It is the average effect of the weak interaction in the unperturbed state. The method predicts Stark and Zeeman effects, molecular spectra, and corrections to the harmonic oscillator \cite{perturbation_ref3}.

In celestial mechanics, the small parameter may be a mass ratio or eccentricity. The starting ellipse is corrected by other planets. In fluid mechanics it may be a small viscosity or wave amplitude. In engineering it may be a small change in stiffness. The method works because the correction is measured relative to a model whose behavior is understood.

\paragraph{How to judge an approximation}
There are three questions. Is the parameter genuinely small in the region of interest? Is the series convergent or merely asymptotic? Does the truncated expression satisfy the required accuracy and boundary conditions? An asymptotic series can have excellent early terms while eventually diverging, so a remainder estimate or comparison with computation is essential \cite{perturbation_ref1,perturbation_ref2}.

\section*{4. Important theorems and results}
\begin{enumerate}[leftmargin=*,itemsep=0.35em]
\item \textbf{Implicit-function principle.} If the unperturbed solution is regular and the derivative with respect to the unknown is invertible, a nearby solution often depends smoothly on the small parameter.

\item \textbf{Regular perturbation expansion.} Substitution of a power series and matching powers of $\varepsilon$ produces a hierarchy of correction equations.

\item \textbf{Rayleigh--Schrodinger first-order rule.} For a nondegenerate quantum level, the first energy shift is the expectation of the perturbing operator.

\item \textbf{Secular-term principle.} Terms that grow with the independent variable indicate that the assumed fast scale is incomplete.

\item \textbf{Method of multiple scales.} Introducing slow variables prevents artificial long-time growth and gives uniformly useful approximations.

\item \textbf{Matched-asymptotic principle.} Outer and inner approximations can be joined in an overlap region to solve singular problems with boundary layers.
\end{enumerate}
\noindent\emph{Source for these results:} \cite{perturbation_ref1}.

\theoryapplications
\theoryconnections{perturbation_theory}{%
\item Calculus expands a problem in a small parameter, and differential equations determine the correction at each order.
\item Linear algebra solves the first correction once the unperturbed problem is known, while asymptotic estimates say how the neglected terms behave.
\item For example, if a weak force changes a solvable oscillator, the zeroth-order motion gives the main frequency and the first-order equation predicts its small shift; the calculation is useful only while the small parameter remains in the stated range \cite{perturbation_ref1,perturbation_ref3}.
}{%
\item Perturbation theory lets celestial mechanics estimate the effect of a small additional body, quantum theory estimate energy shifts, and fluid mechanics describe weak viscosity or slow rotation.
\item In engineering and wave theory, it separates a dominant design from a small imperfection, making a difficult equation solvable in stages.
\item Secular terms are especially important: if a correction grows with time, it signals that the approximation must be reorganized rather than trusted indefinitely \cite{perturbation_ref1,perturbation_ref2}.
}

\section*{Glossary}
\begin{description}[leftmargin=*,style=nextline,itemsep=0.3em]
\item[\textbf{Small parameter.}] A quantity $\varepsilon$ that is much smaller than one in the regime of interest; the perturbation expansion treats corrections as powers of $\varepsilon$.
\item[\textbf{Regular perturbation.}] An expansion in powers of a small parameter where each successive term is smaller and the solution depends smoothly on the parameter.
\item[\textbf{Singular perturbation.}] A problem in which the small parameter multiplies the highest derivative, so setting it to zero changes the order of the equation and loses boundary conditions.
\item[\textbf{Secular term.}] A correction that grows without bound as the independent variable increases, signalling that the assumed timescale is incomplete and the expansion must be reorganized.
\item[\textbf{Boundary layer.}] A thin region near a boundary where the solution changes rapidly because a singular perturbation has reduced the order of the differential equation.
\item[\textbf{Multiple scales.}] A method that introduces slow and fast time variables to eliminate secular terms and obtain uniformly valid approximations.
\item[\textbf{Asymptotic series.}] A formal power series whose partial sums approximate a function well for small parameter values even though the series may eventually diverge.
\item[\textbf{Unperturbed problem.}] The simplified equation obtained by setting the small parameter to zero; it is exactly solvable and serves as the starting point for corrections.
\item[\textbf{Degeneracy.}] A situation where the unperturbed problem has several independent states with the same eigenvalue.
\item[\textbf{Matched asymptotic expansions.}] A method constructing separate inner and outer approximations and joining them in an overlap region.
\item[\textbf{Dominant balance.}] A procedure for identifying which terms in an equation remain significant in a given limit.
\end{description}
\section*{7. Sample Questions}
\emph{Exercise reference:} \cite{perturbation_ref1}. The cited edition is the source for the exercise set; consult its Exercises section for the official statement and numbering.
\begin{enumerate}[leftmargin=*,itemsep=0.9em]
\item \textbf{Exercise 1.} Use regular perturbation to find the first two terms of the expansion of the positive root of $x+\varepsilon x^2-1=0$ for small $\varepsilon$. \emph{Solution.} Write $x=x_0+\varepsilon x_1+\cdots$. Substituting: $(x_0+\varepsilon x_1)+\varepsilon(x_0+\varepsilon x_1)^2-1=0$. At order $\varepsilon^0$: $x_0-1=0$, so $x_0=1$. At order $\varepsilon^1$: $x_1+x_0^2=0$, so $x_1=-1$. Thus $x=1-\varepsilon+O(\varepsilon^2)$.

\item \textbf{Exercise 2.} Consider the perturbed eigenvalue problem $H_0=\begin{pmatrix}1&0\\0&3\end{pmatrix}$ and $\varepsilon V=\varepsilon\begin{pmatrix}0&1\\1&0\end{pmatrix}$. Compute the first-order energy correction for the eigenvalue $E_1^{(0)}=1$. \emph{Solution.} The unperturbed eigenvector is $\psi_1^{(0)}=\begin{pmatrix}1\\0\end{pmatrix}$. The first-order shift is $E_1^{(1)}=\langle\psi_1^{(0)},V\psi_1^{(0)}\rangle=\begin{pmatrix}1&0\end{pmatrix}\begin{pmatrix}0&1\\1&0\end{pmatrix}\begin{pmatrix}1\\0\end{pmatrix}=\begin{pmatrix}1&0\end{pmatrix}\begin{pmatrix}0\\1\end{pmatrix}=0$.

\item \textbf{Exercise 3.} For the singularly perturbed boundary-value problem $\varepsilon y''+y'=0$ with $y(0)=0$ and $y(1)=1$, find the outer solution valid away from $x=0$. \emph{Solution.} Setting $\varepsilon=0$ gives $y'=0$, so $y=C$ (a constant). The outer solution cannot satisfy both boundary conditions. Matching with an inner layer near $x=0$ (where the equation $\varepsilon y''+y'=0$ is rescaled by $\xi=x/\varepsilon$) yields the outer solution $y=1$ for $x>0$, which satisfies the boundary condition at $x=1$.

\item \textbf{Exercise 4.} The equation $x^3-\varepsilon x-2=0$ has a root near the unperturbed solution $x_0=\sqrt[3]{2}$. Find the $\varepsilon$-correction to first order. \emph{Solution.} Write $x=\sqrt[3]{2}+\varepsilon x_1+\cdots$. Substituting: $(\sqrt[3]{2}+\varepsilon x_1)^3-\varepsilon(\sqrt[3]{2}+\varepsilon x_1)-2=0$. Expanding the cube and collecting the $\varepsilon^1$ term: $3(\sqrt[3]{2})^2 x_1-\sqrt[3]{2}=0$, so $x_1=\sqrt[3]{2}/(3\cdot 2^{2/3})=\sqrt[3]{2}/(3\cdot\sqrt[3]{4})=1/3$.

\item \textbf{Exercise 5.} For the weakly nonlinear oscillator $x''+x+\varepsilon x^2=0$ with $x(0)=1$, $x'(0)=0$, apply the method of multiple scales. Introduce $T=\varepsilon t$ and write $x=x_0(t,T)+\varepsilon x_1(t,T)+\cdots$. Show that avoiding secular terms in the $\varepsilon^1$ equation requires $\partial x_0/\partial T=0$, so the frequency is unchanged at first order. \emph{Solution.} The zeroth-order equation is $\partial^2 x_0/\partial t^2+x_0=0$, giving $x_0=A(T)\cos(t+\phi(T))$. At first order, $\partial^2 x_1/\partial t^2+x_1=-2\partial^2 x_0/\partial t\partial T-x_0^2$. The term $-x_0^2=-A^2\cos^2(t+\phi)=-\frac{A^2}{2}(1+\cos 2(t+\phi))$ contains no resonant $\cos(t+\phi)$ term. The $\partial^2 x_0/\partial t\partial T=2A'(T)\sin(t+\phi)$ is resonant. Setting $A'(T)=0$ eliminates the secular term, confirming the frequency is unchanged at first order for $x''+x+\varepsilon x^2=0$.

\end{enumerate}
\section*{8. References}
\begin{thebibliography}{9}
\bibitem{perturbation_ref1} C. M. Bender and S. A. Orszag, \emph{Advanced Mathematical Methods for Scientists and Engineers}, Springer, 1978.
\bibitem{perturbation_ref2} M. H. Holmes, \emph{Introduction to Perturbation Methods}, 2nd ed., Springer, 2013.
\bibitem{perturbation_ref3} J. J. Sakurai and J. Napolitano, \emph{Modern Quantum Mechanics}, 3rd ed., Cambridge University Press, 2017.
\end{thebibliography}


\clearpage\chapter{Probability theory}
\section*{1. Overview}
\subsection*{Intuitive}
\paragraph{The problem that started it}
People have always faced uncertain outcomes: a harvest may fail, a ship may arrive late, or a game may be won by chance. The important question is not whether one event can be predicted perfectly. It is whether repeated uncertainty has a pattern that can be measured and used for a decision.
In the sixteenth century, Gerolamo Cardano studied games of chance. In the seventeenth century, Blaise Pascal and Pierre de Fermat exchanged letters about the ``problem of points'': if a game is stopped before either player has won, how should the prize be divided fairly? Their answer turned a quarrel about a game into a method for reasoning about all possible outcomes.


\subsection*{Formal}
A probability space is a triple $(\Omega,\mathcal F,\mathbb P)$, where $\Omega$ is the set of outcomes, $\mathcal F$ is a $\sigma$-algebra of events, and $\mathbb P:\mathcal F\to[0,1]$ satisfies $\mathbb P(\Omega)=1$ and countable additivity on disjoint events. A random variable is a measurable function from $\Omega$ to a numerical space, and its expectation is the probability-weighted average when that average exists.
\section*{2. History}
\subsection*{Historical development}
\begin{historylist}
\paragraph{The people and the changing idea}
Christiaan Huygens published an early book on probability in 1657. Pierre-Simon Laplace later developed a broad classical treatment. At first, probability was mainly about finite lists such as dice, cards, and coins. As science began measuring continuously changing quantities, mathematicians needed a framework that could handle intervals, curves, and infinitely many possibilities.

Andrey Kolmogorov supplied the modern foundation in 1933. He described a probability model using three ingredients: a collection of possible outcomes, a selection of events we can ask about, and a rule assigning each event a number from zero to one. The number for the whole collection is one, the number for an impossible event is zero, and disjoint events add. Kolmogorov was not trying to make chance less mysterious; he was making clear which conclusions follow from which assumptions.

\end{historylist}
\section*{3. Theory}
\subsection*{Core theory}
\paragraph{The central idea}
Imagine rolling a die. The possible outcomes are $1,2,3,4,5,6$. The event ``an even number'' is the smaller collection $\{2,4,6\}$. If the die is fair, its probability is $3/6=1/2$. The same language can describe a weather forecast, except that the list of outcomes is much larger and the probabilities come from observations and models rather than symmetry.
\bookfigure{probability_statistics}{A probability model starts with possible outcomes, observes data, and updates beliefs.}

Conditional probability is the crucial next step. It asks how the chance of one event changes after we learn another fact. A medical test may be accurate, yet a positive result may still be more likely to be a false alarm when the disease is rare. Independence means that learning one event gives no information about the other; it is an assumption to check, not a synonym for ``far apart.''

\paragraph{Why the method works}
The rules work because they separate three things that are often confused: what can happen, what information is available, and how strongly each possibility is supported. Once those are stated, two people can calculate the same answer and locate their disagreement. Probability does not promise certainty. Its power is that repeated trials often produce stable averages even when individual trials remain unpredictable.

Two major results explain this stability. The law of large numbers says that averages tend to settle near their expected value as the number of trials grows. The central limit theorem explains why sums of many small, partly independent effects often form a bell-shaped pattern. These are not magic guarantees for a small or biased sample; they are theorems with conditions.

\paragraph{Three small examples}
\textbf{A bus.} If a bus arrives exactly every ten minutes and a passenger arrives at a random time, the average wait is five minutes. If buses bunch together, the assumption changes and the average wait may be longer.

\textbf{A medical test.} Suppose 1 in 100 people has a condition. A test catches 99 percent of cases but falsely flags 5 percent of healthy people. Among 10,000 people, about 99 sick people test positive, but about 495 healthy people also test positive. A positive result is therefore genuine only about $99/(99+495)=1/6$ of the time. The lesson is to combine test accuracy with the background rate.

\textbf{A repeated game.} A fair coin has expected winnings of zero when heads wins one point and tails loses one. After ten tosses a player may be ahead or behind; the expected value describes the long-run centre, not the result of one evening.

\section*{4. Important theorems and results}
\begin{enumerate}[leftmargin=*,itemsep=0.35em]
\item \textbf{Conditional probability.} For events $A$ and $B$ with $P(B)>0$, $P(A\mid B)=P(A\cap B)/P(B)$. It updates the probability of $A$ after the information $B$ is known; it does not assert that $A$ and $B$ are independent.

\item \textbf{Law of large numbers.} Under the usual integrability and independence or dependence-control assumptions, sample averages converge to their expected value. It explains why repeated measurements can stabilize while leaving finite-sample error.

\item \textbf{Central limit theorem.} For suitable independent observations with finite variance, a normalized sum approaches a normal distribution. This result supports approximate intervals and tests, but it is not exact for every sample or valid under arbitrary dependence.

\item \textbf{Bayes' rule.} $P(A\mid B)=P(B\mid A)P(A)/P(B)$ reverses the direction of evidence. It is useful for diagnosis and classification because it combines a likelihood with a base rate.

For two independent fair die rolls, the probability of at least one six is $1-(5/6)^2=11/36$. This small calculation illustrates the larger rule: a probability conclusion is only as reliable as its sample space, dependence assumptions, and data-generating model.\end{enumerate}
\noindent\emph{Source for these results:} \cite{probability_theory_ref1}.

\section*{5. Applications}
Probability theory represents uncertainty by specifying possible outcomes, their probabilities, and the rules for updating those probabilities when information changes. Its conclusions are meaningful only when the sample space and dependence assumptions match the situation being modelled.
\begin{enumerate}[leftmargin=*,itemsep=0.35em]
\item \textbf{Insurance.} Expected losses and tail probabilities determine premiums and reserves.
\item \textbf{Medical trials.} Randomization and sampling distributions quantify whether an observed treatment difference could arise by chance.
\item \textbf{Polling.} A sampling model converts a finite set of responses into an uncertainty range for a larger population.
\item \textbf{Finance.} Probability models portfolio risk, default, and the effect of rare but costly events.
\item \textbf{Weather forecasting.} An ensemble of possible atmospheric states represents uncertainty in observations and initial conditions.
\item \textbf{Queues and networks.} Arrival and service distributions predict waiting times, congestion, and capacity.
\item \textbf{Genetics.} Probability describes inheritance, mutation, population sampling, and genetic drift.
\item \textbf{Machine learning.} Probabilistic models represent uncertain labels, predictions, and hidden variables.
\end{enumerate}
In each case the model depends on choices: which outcomes were included, whether observations are independent, and whether the past resembles the future. A precise calculation from a badly chosen model is still a bad prediction.

\theoryconnections{probability_theory}{%
\item Probability theory builds a model by specifying possible outcomes and assigning weights whose total is one.
\item Counting handles finite equally likely cases, linear algebra organizes joint distributions and conditional expectations, and calculus handles continuous densities and accumulation over infinitely many outcomes.
\item Independence is a modeling rule that permits multiplication of probabilities; it is not a consequence of two events merely being unrelated in everyday language \cite{probability_theory_ref1,probability_theory_ref2}.
}{%
\item Probability supports statistics by describing how data could vary under a proposed model, which makes estimation and testing possible.
\item It supports engineering through failure probabilities and queue lengths, information theory through entropy and channel noise, and economics through risk and expected utility.
\item In every application, the theory converts uncertain future outcomes into averages, distributions, or bounds that can guide a decision \cite{probability_theory_ref1,probability_theory_ref3}.
}

\section*{Glossary}
\begin{description}[leftmargin=*,style=nextline,itemsep=0.3em]
\item[\textbf{Probability space.}] A triple $(\Omega,\mathcal{F},\mathbb{P})$ consisting of a set of outcomes, a collection of admissible events, and a rule assigning each event a number between zero and one.
\item[\textbf{Random variable.}] A measurable function from the outcome set to a numerical space, turning an abstract outcome into a number that can be averaged or compared.
\item[\textbf{Conditional probability.}] The probability of an event $A$ given that another event $B$ has occurred, defined as $P(A\mid B)=P(A\cap B)/P(B)$.
\item[\textbf{Independence.}] Two events are independent when knowing that one occurred does not change the probability of the other; it is a modeling assumption, not a synonym for spatial separation.
\item[\textbf{Expectation.}] The probability-weighted average of a random variable, computed by summing or integrating each possible value multiplied by its probability.
\item[\textbf{Law of large numbers.}] The theorem that sample averages converge to the expected value as the number of trials grows, explaining why repeated measurements stabilize.
\item[\textbf{Central limit theorem.}] The result that a normalized sum of many independent random variables with finite variance approaches a normal distribution, regardless of the original distribution.
\item[\textbf{Bayes' rule.}] A formula $P(A\mid B)=P(B\mid A)P(A)/P(B)$ that reverses the direction of evidence, combining a likelihood with a base rate.
\item[$\sigma$-\textbf{algebra.}] A collection of subsets closed under complements and countable unions; specifies which events can be assigned probabilities.
\item[\textbf{Sample space.}] The set $\Omega$ of all possible outcomes of a random experiment.
\item[\textbf{Variance.}] The expected squared deviation of a random variable from its mean, measuring spread.
\item[\textbf{Normal distribution.}] The bell-shaped density proportional to $e^{-x^2/2}$, arising as the limit of sums of independent random variables.
\end{description}
\chapterquestions{Probability theory}{sample spaces, events, conditional probabilities, and expectations}{the probability of at least one six in two fair die rolls}{$1-(5/6)^2=11/36$ by taking the complement of no six}{conditional probability updates a model after information is observed}{dependent trials cannot be multiplied blindly; probability supports risk and forecasting}{probability_theory_ref1}
\section*{8. References}
\begin{thebibliography}{9}
\bibitem{probability_theory_ref1} D. Williams, \emph{Weighing the Odds}, Cambridge University Press, 2001.
\bibitem{probability_theory_ref2} S. M. Ross, \emph{A First Course in Probability}, 10th ed., Pearson, 2019.
\bibitem{probability_theory_ref3} A. N. Kolmogorov, \emph{Foundations of the Theory of Probability}, Chelsea, 1956.
\end{thebibliography}


\theory{Proof theory}{What is a proof when the proof itself is treated as a precise mathematical object that can be inspected, transformed, and checked by a machine?}{Proof theory studies formal proofs, the rules that build them, and the strength and limits of formal systems. It is syntactic: it studies the written construction of a proof, while model theory studies what statements mean in structures.}{Foundations of mathematics, automated theorem proving, proof assistants, programming languages, type theory, linguistics, constructive mathematics, and proof complexity.}{Formal derivations, normalization, and ordinal or combinatorial measures quantify the strength and computational behavior of proof systems.}

\paragraph{The idea and history}
A formal proof is a finite tree or list built from axioms and inference rules. The rules determine exactly which conclusion may follow from which premises. This makes a proof an object that can be counted, normalized, compared, and checked. An informal proof is a compressed explanation that an expert can expand; a formal proof contains every permitted step \cite{proof_ref1,proof_ref2}.

Frege, Peano, Russell, and Dedekind advanced formal logic. David Hilbert proposed a program in which mathematics would be grounded by finitary consistency proofs. Gödel's incompleteness theorems showed that a sufficiently strong consistent theory cannot prove its own consistency. This changed the questions toward relative consistency, provability, ordinals, and the exact strength of axioms \cite{proof_ref1,proof_ref3}.

\end{historylist}
\section*{3. Theory}
\subsection*{Core theory}
\paragraph{Proof calculi}
Hilbert systems use a small collection of axiom schemes and rules such as modus ponens. Natural deduction mirrors ordinary reasoning: an introduction rule explains how to prove a connective, and an elimination rule explains how to use it. Sequent calculus studies judgments of the form
\[
\Gamma\Longrightarrow\Delta,
\]
where $\Gamma$ is a list of assumptions and $\Delta$ is a list of possible conclusions. Gentzen developed natural deduction and sequent calculus and used them to prove consistency results for arithmetic \cite{proof_ref2}.

\paragraph{Analytic proofs and normalization}
The cut rule allows a proof to use an intermediate lemma. Cut elimination says that in important calculi every provable sequent has a proof without cuts. Removing cuts can make the proof longer, but it exposes its logical structure. The subformula property says that a cut-free proof uses only pieces of the final statement and its premises. This supports consistency, proof search, interpolation, and automated checking \cite{proof_ref2,proof_ref4}.

In natural deduction, normalization removes a formula that is introduced and immediately eliminated. The normalized proof contains no unnecessary detour. Through the Curry--Howard correspondence, propositions behave like types and proofs behave like programs; normalization of proofs corresponds to computation in typed lambda calculus \cite{proof_ref4}.

\paragraph{Main research areas}
Structural proof theory studies calculi, cut elimination, subformula properties, focused proofs, proof nets, analytic tableaux, and substructural logics such as linear logic. Ordinal analysis measures the strength of arithmetic, analysis, and set theories by associating them with well-founded transfinite ordinals. Gentzen's consistency proof for Peano arithmetic used induction up to epsilon-zero \cite{proof_ref1}.

Provability logic uses a box symbol meaning it is provable that. Gödel--Löb logic captures important facts about provability in arithmetic. Reverse mathematics asks which axioms are needed to prove a theorem: it works backward from theorem to foundations and has identified robust systems often called the Big Five \cite{proof_ref3}.

Functional interpretations translate classical or intuitionistic proofs into functionals. They can extract a computable witness or bound from an existence proof. Proof mining uses this information to turn abstract proofs into quantitative estimates. Proof complexity measures the length or size of proofs, while automated theorem proving searches for proofs and proof assistants check them \cite{proof_ref1,proof_ref5}.

\paragraph{Semantics, language, and computers}
Proof-theoretic semantics explains the meaning of a connective through the rules for introducing and using it. This idea has applications in constructive mathematics and formal linguistics, including type-logical and categorial grammar. In computer science, a proof assistant stores a proof term and checks it against a small trusted kernel. Checking is usually easier than discovering the proof \cite{proof_ref4}.

\section*{4. Important theorems and results}
\begin{enumerate}[leftmargin=*,itemsep=0.35em]
\item \textbf{Gödel's first incompleteness theorem.} Any consistent, effectively axiomatized theory strong enough for arithmetic has a true arithmetic statement that it cannot prove.

\item \textbf{Gödel's second incompleteness theorem.} Such a theory cannot prove its own consistency, assuming it is consistent.

\item \textbf{Gentzen cut elimination.} Classical and intuitionistic sequent calculi admit cut-free proofs.

\item \textbf{Subformula property.} In a cut-free proof, formulas are built from subformulas of the endsequent, supporting proof search and consistency.

\item \textbf{Curry--Howard correspondence.} Propositions correspond to types, proofs to programs, and proof normalization to program evaluation.

\item \textbf{Löb's theorem.} If a formal theory proves that the provability of a statement implies the statement, then the theory already proves the statement.
\end{enumerate}
\noindent\emph{Source for these results:} \cite{proof_ref1}.

\theoryapplications
\theoryconnections{proof_theory}{%
\item Logic supplies formal statements and inference rules, while recursion theory identifies which symbolic procedures can actually be carried out.
\item Set theory provides a universe in which mathematical objects can be represented, and computation turns a proof into a finite checkable object.
\item Proof theory studies the resulting formal systems by measuring the strength of their rules and by asking which proofs can be normalized or eliminated \cite{proof_ref1,proof_ref2}.
}{%
\item Proof theory helps proof assistants check that every inference is legal, and automated reasoning searches for proofs by manipulating the same formal rules.
\item In programming languages, a type can describe a specification and a program can serve as a proof that the specification is met.
\item Consistency and incompleteness results then tell foundations which promises a formal system can and cannot make \cite{proof_ref1,proof_ref4}.
}
\section*{7. Sample Questions}
\emph{Exercise reference:} \cite{proof_ref1}. The cited edition is the source for the exercise set; consult its Exercises section for the official statement and numbering.
\begin{enumerate}[leftmargin=*,itemsep=0.9em]
\item \textbf{Exercise 1. What does modus ponens say?} From $P$ and $P\Rightarrow Q$, infer $Q$. In a formal proof this is one licensed rule, so the conclusion is mechanically checkable.

\item \textbf{Exercise 2. What is the purpose of cut elimination?} It removes reliance on an intermediate lemma from the formal derivation. The result exposes the proof's local structure and often gives the subformula property.

\item \textbf{Exercise 3. Give the Curry--Howard reading of a proof of $P\Rightarrow P$.} It is the identity program: given an input of type $P$, return that same input. The corresponding lambda term is $\lambda x.x$.

\item \textbf{Exercise 4. Why does Gödel's second theorem not say that mathematics is useless?} It concerns one sufficiently strong formal system proving its own consistency. Stronger systems can prove the consistency of weaker systems, and ordinary mathematics remains productive despite these limits.

\item \textbf{Exercise 5. What does reverse mathematics ask about the theorem that every bounded increasing sequence has a supremum?} It asks for the weakest axiom system in which the theorem can be proved and then proves a reversal showing that the theorem implies that system over a weak base theory.

\end{enumerate}
\section*{8. References}
\begin{thebibliography}{9}
\bibitem{proof_ref1} S. Buss, ed., \emph{Handbook of Proof Theory}, Elsevier, 1998.
\bibitem{proof_ref2} M. E. Szabo, ed., \emph{The Collected Papers of Gerhard Gentzen}, North-Holland, 1969.
\bibitem{proof_ref3} S. G. Simpson, \emph{Subsystems of Second Order Arithmetic}, 2nd ed., Cambridge University Press, 2009.
\bibitem{proof_ref4} A. S. Troelstra and H. Schwichtenberg, \emph{Basic Proof Theory}, 2nd ed., Cambridge University Press, 2000.
\bibitem{proof_ref5} J.-Y. Girard, P. Taylor, and Y. Lafont, \emph{Proofs and Types}, Cambridge University Press, 1989.
\end{thebibliography}


\theory{Queueing theory}{How many servers, machines, or communication channels are needed so that waiting remains acceptable without wasting capacity?}{Queueing theory models arrivals, service, waiting discipline, and departures as random processes. It predicts queue lengths, waiting times, utilization, throughput, and the probability of delay.}{Telephone networks, Internet routers, hospitals, factories, shops, traffic systems, call centers, cloud computing, and project management.}{Arrival and service processes determine waiting-time distributions, utilization, stability, and capacity in a stochastic system.}

\paragraph{The problem and history}
Every queue has arrivals, one or more servers, a waiting area, a service rule, and departures. The same structure describes a supermarket line, packets entering a router, patients entering an emergency department, and jobs sent to a processor. The goal is to calculate distributions and averages that support capacity decisions \cite{queue_ref1}.

Agner Krarup Erlang created the subject while studying incoming calls at the Copenhagen Telephone Exchange. He modeled arrivals by a Poisson process and solved early single- and multiple-server systems. Pollaczek, Khinchin, Kendall, and Kingman developed major later results. Kleinrock applied queueing models to message switching and packet networks \cite{queue_ref2,queue_ref3}.

\end{historylist}
\section*{3. Theory}
\subsection*{Core theory}
\paragraph{Building a model}
An arrival process describes when customers or jobs appear. A Poisson process with rate $\lambda$ has independent arrivals and exponentially distributed interarrival times. The service process describes how long a server takes; exponential service is denoted by M, deterministic service by D, and a general distribution by G. Kendall's notation writes a queue as A/S/$c$, where $c$ is the number of servers. M/M/1 has Poisson arrivals, exponential service, and one server.

The queue may have finite or infinite capacity, a finite or infinite customer population, and different disciplines. First-come first-served is common, but priority, last-come first-served, processor sharing, and random selection also occur. A customer may abandon the queue or be blocked when no buffer is free \cite{queue_ref1,queue_ref4}.

\paragraph{The M/M/1 queue}
Let arrivals occur at rate $\lambda$ and service at rate $\mu$. The traffic intensity is
\[
\rho=\frac{\lambda}{\mu}.
\]
For an unlimited queue to have a steady state, $\rho<1$ is required. If $\rho\geq1$, work arrives as fast as or faster than it can be processed.

When $\rho<1$, the probability of finding $n$ customers in the system is
\[
P(N=n)=(1-\rho)\rho^n,\qquad n=0,1,2,\ldots.
\]
The mean number in the system is $L=\rho/(1-\rho)$, and the mean time in the system is $W=1/(\mu-\lambda)$. The mean number waiting is $L_q=\rho^2/(1-\rho)$ and the mean waiting time is $W_q=\rho/(\mu-\lambda)$ \cite{queue_ref1}.

\paragraph{Little's law and networks}
Little's law states
\[
L=\lambda W,
\]
where $L$ is the long-run average number in the system, $\lambda$ is throughput, and $W$ is the average time spent in the system. If a website completes $600$ requests per minute and the average request takes $0.2$ minutes, the average number in the system is $120$.

Queueing networks connect several nodes by routing rules. One may seek an equilibrium distribution, but transient behavior matters when a hospital faces a sudden surge or a computer system starts after a failure. Product-form networks have equilibrium probabilities that factor into contributions from individual nodes \cite{queue_ref2,queue_ref3}.

\paragraph{Multiple servers and variability}
An M/M/c queue has $c$ identical servers. Adding a server usually reduces delay, but the benefit depends on utilization and variability. Deterministic service is less disruptive than highly variable service with the same mean. The Pollaczek--Khinchin formula for an M/G/1 queue makes the effect of service-time variance explicit.

Kingman's approximation for a G/G/1 queue shows that mean waiting grows with utilization, arrival variability, and service variability. This explains why a system operating at ninety-nine percent capacity can be much worse than one operating at eighty percent: random bursts have nowhere to go \cite{queue_ref4,queue_ref5}.

\paragraph{How it solves real problems}
In a hospital, arrivals vary by hour and service times vary by patient. A model estimates the probability of a long wait under one staffing plan, then compares it with plans using another doctor or a fast-treatment stream. The output is a tradeoff between cost, congestion, and service quality.

In a router, packets arrive and compete for a finite buffer. If the service rate is too low, the buffer fills and packets are dropped. Queueing analysis estimates delay and loss, helping engineers choose link speeds and buffer sizes. In cloud computing, the same calculation determines how many virtual machines are needed for a target response time \cite{queue_ref2}.

\section*{4. Important theorems and results}
\begin{enumerate}[leftmargin=*,itemsep=0.35em]
\item \textbf{Erlang formulas.} Erlang's loss and delay formulas give blocking and waiting probabilities in telephone and multi-server systems.

\item \textbf{M/M/1 stationary distribution.} When $\rho<1$, $P(N=n)=(1-\rho)\rho^n$.

\item \textbf{Little's law.} In a stable long-run system, $L=\lambda W$.

\item \textbf{Pollaczek--Khinchin formula.} The M/G/1 mean waiting time depends on the first two moments of service time, so variability increases delay.

\item \textbf{Kingman's formula.} For G/G/1 queues, mean waiting is approximately proportional to $\rho/(1-\rho)$ and to the sum of arrival and service variability.

\item \textbf{Birth--death balance.} A single queue can be represented by transitions from $n$ to $n+1$ at arrival rate and from $n$ to $n-1$ at service rate.
\end{enumerate}
\noindent\emph{Source for these results:} \cite{queue_ref1}.

\theoryapplications
\theoryconnections{queue_theory}{%
\item Probability describes random arrivals and service times, while a stochastic process records how the number of waiting jobs changes through time.
\item In a memoryless model, the current queue length is enough to predict the next transition, producing a Markov chain and balance equations.
\item Operations research then turns the resulting waiting-time and utilization formulas into a staffing or capacity decision \cite{queue_ref1,queue_ref2}.
}{%
\item Queueing theory helps a call center choose how many agents are needed to keep waiting time below a target, and helps a data center decide how much service capacity prevents overload.
\item The same calculation estimates traffic congestion, hospital beds, factory buffers, and network delay.
\item Its distinctive contribution is to show the cost of variability: even when average arrival and service rates are equal, random bursts can make the queue grow without bound \cite{queue_ref1,queue_ref3}.
}

\section*{Glossary}
\begin{description}[leftmargin=*,style=nextline,itemsep=0.3em]
\item[\textbf{Traffic intensity.}] The ratio $\rho=\lambda/\mu$ of the arrival rate to the service rate; a queue with unlimited capacity reaches a steady state only when $\rho<1$.
\item[\textbf{Poisson process.}] A model of random arrivals in which interarrival times are independent and exponentially distributed, characterized by a constant average rate $\lambda$.
\item[\textbf{Little's law.}] The relationship $L=\lambda W$ stating that the long-run average number in a stable system equals the throughput multiplied by the average time spent in the system.
\item[\textbf{M/M/1 queue.}] The simplest queueing model with Poisson arrivals, exponentially distributed service times, one server, and unlimited waiting capacity.
\item[\textbf{Kendall's notation.}] A shorthand A/S/$c$ that classifies a queue by its arrival distribution, service distribution, and number of servers.
\item[\textbf{Steady state.}] A long-run equilibrium condition in which the probability distribution of the number in the system no longer changes with time.
\item[\textbf{Pollaczek--Khinchin formula.}] An exact expression for the mean waiting time in an M/G/1 queue showing that service-time variance directly increases delay.
\item[\textbf{Kingman's formula.}] An approximation for the mean waiting time in a G/G/1 queue proportional to $\rho/(1-\rho)$ and to the sum of arrival and service variabilities.
\item[\textbf{Interarrival time.}] The time between consecutive arrivals; exponentially distributed in a Poisson process.
\item[\textbf{Service time.}] The duration a server takes to complete one job.
\item[\textbf{Utilization.}] The fraction of time a server is busy, equal to $\rho$ in single-server models.
\item[\textbf{Birth--death process.}] A continuous-time Markov chain where transitions only increase or decrease the state by one.
\item[\textbf{Erlang formula.}] Formulas giving blocking and waiting probabilities in multi-server telephone systems.
\end{description}
\section*{7. Sample Questions}
\emph{Exercise reference:} \cite{queue_ref1}. The cited edition is the source for the exercise set; consult its Exercises section for the official statement and numbering.
\begin{enumerate}[leftmargin=*,itemsep=0.9em]
\item \textbf{Exercise 1.} An M/M/1 queue has arrival rate $\lambda=4$ jobs per minute and service rate $\mu=5$ jobs per minute. Compute (a) the traffic intensity $\rho$, (b) the mean number in the system $L$, and (c) the probability that the system is empty. \emph{Solution.} (a) $\rho=\lambda/\mu=4/5=0.8$. (b) $L=\rho/(1-\rho)=0.8/0.2=4$. (c) $P(N=0)=1-\rho=0.2$.

\item \textbf{Exercise 2.} For an M/M/1 queue with $\lambda=6$ per hour and $\mu=10$ per hour, compute the probability of having exactly $3$ customers in the system in steady state. \emph{Solution.} $\rho=6/10=0.6$. The steady-state probability is $P(N=3)=(1-\rho)\rho^3=0.4\times 0.6^3=0.4\times 0.216=0.0864$.

\item \textbf{Exercise 3.} A call center completes $200$ calls per hour, and the average call duration is $4$ minutes. Using Little's law, find the average number of calls in progress simultaneously. \emph{Solution.} Convert to consistent units: $\lambda=200/60=10/3$ calls per minute, $W=4$ minutes. Little's law gives $L=\lambda W=(10/3)\times 4=40/3\approx 13.3$ calls.

\item \textbf{Exercise 4.} In an M/M/2 queue with $\lambda=8$ per minute and $\mu=5$ per minute per server, compute the traffic intensity per server and determine whether the queue is stable. \emph{Solution.} The offered load is $a=\lambda/\mu=8/5=1.6$. With $c=2$ servers, the per-server utilization is $\rho=a/c=1.6/2=0.8<1$, so the queue is stable.

\item \textbf{Exercise 5.} For an M/M/1 queue with $\lambda=3$ per minute and $\mu=4$ per minute, compute the mean waiting time in queue (not including service). \emph{Solution.} $\rho=\lambda/\mu=3/4=0.75$. The mean number waiting is $L_q=\rho^2/(1-\rho)=0.75^2/0.25=0.5625/0.25=2.25$. By Little's law applied to the queue, $W_q=L_q/\lambda=2.25/3=0.75$ minutes, or $45$ seconds. Equivalently, $W_q=\rho/(\mu-\lambda)=0.75/1=0.75$ minutes.

\end{enumerate}
\section*{8. References}
\begin{thebibliography}{9}
\bibitem{queue_ref1} D. Gross, J. F. Shortle, J. M. Thompson, and C. M. Harris, \emph{Fundamentals of Queueing Theory}, 4th ed., Wiley, 2008.
\bibitem{queue_ref2} L. Kleinrock, \emph{Queueing Systems, Volume 1: Theory}, Wiley, 1975.
\bibitem{queue_ref3} L. Kleinrock, \emph{Queueing Systems, Volume 2: Computer Applications}, Wiley, 1976.
\bibitem{queue_ref4} J. W. Cohen, \emph{The Single Server Queue}, North-Holland, 1969.
\bibitem{queue_ref5} J. F. C. Kingman, \emph{Heavy Traffic in Single Server Queues}, Notes in Applied Mathematics, American Mathematical Society, 1965.
\end{thebibliography}


\theory{Ramsey theory}{How large must a disordered structure be before some organized pattern is unavoidable?}{Ramsey theory studies regular substructures inside large objects. A coloring or arbitrary arrangement may look chaotic, but sufficiently large systems must contain a monochromatic or otherwise structured part.}{Network design, communication, additive number theory, theoretical computer science, geometry, logic, and unavoidable patterns.}{Finite colorings and unavoidable monochromatic patterns turn local choices into universal combinatorial guarantees.}

\paragraph{The question and history}
Frank Ramsey created the subject while studying a problem in the foundations of mathematics. The modern question asks for the least size that forces a desired configuration, even when the original object is colored or partitioned adversarially \cite{ramsey_ref1}.

\end{historylist}
\section*{3. Theory}
\subsection*{Core theory}
\paragraph{The party example}
Suppose every pair among six people is labeled friends or strangers. Choose one person. Among the other five, at least three have the same relation to that person. If three are friends with the chosen person, either two of them know each other, creating a triangle of friends, or all three are mutual strangers. The other case is identical. Thus every six-person group contains three mutual friends or three mutual strangers.

In graph language, color every edge of the complete graph on six vertices red or blue. A monochromatic triangle must occur, so $R(3,3)=6$. The number $R(s,t)$ is the least $n$ such that every red-blue coloring of $K_n$ contains a red $K_s$ or a blue $K_t$ \cite{ramsey_ref2}.

\paragraph{Ramsey's theorem and partition regularity}
For any positive integers $c,n_1,\ldots,n_c$, there is a number $R(n_1,\ldots,n_c)$ such that every coloring of the edges of a complete graph with $c$ colors contains a complete subgraph of size $n_i$ whose edges all have color $i$ for some $i$. This is Ramsey's theorem.

The same idea applies to subsets of integers, arithmetic progressions, matrices, finite sets, and infinite sequences. A statement is partition regular when every finite coloring contains a monochromatic solution of the desired kind. The pigeonhole principle is the simplest example: a large enough collection divided into finitely many boxes has a box containing many objects \cite{ramsey_ref1}.

\paragraph{Major results}
Van der Waerden's theorem says that every finite coloring of the positive integers contains a monochromatic arithmetic progression of any prescribed finite length. Szemeredi's theorem strengthens this: every subset of the integers with positive density contains arbitrarily long arithmetic progressions. Rado's theorem characterizes many systems of linear equations that are partition regular.

The Hales--Jewett theorem states that a sufficiently high-dimensional word cube, colored with finitely many colors, contains a combinatorial line. The density Hales--Jewett theorem gives a positive-density version. Hindman's theorem, the Folkman theorem, and the Milliken--Taylor theorem provide further partition results \cite{ramsey_ref3}.

\paragraph{What the numbers mean}
Ramsey results are usually existential: they prove that a desired pattern exists but may not provide an efficient way to find it. The smallest bounds can be enormous, sometimes growing exponentially or according to very fast functions. Improving upper and lower bounds is therefore central. Computer searches settle some small values, but a proof must still cover all colorings \cite{ramsey_ref2}.

\paragraph{Quantitative bounds and the probabilistic method}
The diagonal Ramsey numbers grow at least exponentially. Erdős showed in 1947 that $R(k,k)>2^{k/2}$ for every $k\geq3$. The argument uses the probabilistic method: choose a random red-blue coloring of the edges of $K_n$ by assigning each edge a color independently with equal probability. For a fixed set of $k$ vertices the chance that all $\binom{k}{2}$ edges among them are the same color is $2\cdot2^{-\binom{k}{2}}$. The expected number of monochromatic $k$-cliques is therefore
\[
\binom{n}{k}\cdot2^{1-\binom{k}{2}}.
\]
When $n=\lfloor2^{k/2}\rfloor$ this expectation is less than $1$, so there exists a coloring with no monochromatic $k$-clique. Consequently $R(k,k)$ must exceed $n$. This short probabilistic argument gives an exponential lower bound and demonstrates the power of nonconstructive existence proofs in Ramsey theory \cite{ramsey_ref1,ramsey_ref2}.

\paragraph{Applications}
In communication networks, a Ramsey bound says how large a system must be before a repeated compatible or incompatible pattern is forced. In theoretical computer science, it controls homogeneous subsets used in lower bounds. In geometry, Erdős--Szekeres results guarantee convex polygons among sufficiently many points. In additive number theory, partition results force arithmetic structure inside apparently random sets.

\paragraph{Application: network design}
A concrete application arises in the design of communication networks. Suppose nodes are paired by links that are either reliable or unreliable. A Ramsey-type argument guarantees that once the network exceeds a threshold size, there must exist a large fully reliable subnetwork or a large fully unreliable subnetwork. This forces the network designer to plan for at least one of these extremes. For instance, if the designer must guarantee a reliable communication path among $k$ terminals, the Ramsey bound $R(k,k)$ tells her the minimum total number of nodes that must be installed so that, regardless of which links fail, a reliable $k$-terminal subnetwork exists. The bound is conservative but provides a guaranteed worst-case guarantee: no adversarial pattern of link failures can destroy all large reliable groups \cite{ramsey_ref1,ramsey_ref3}.

\section*{4. Important theorems and results}
\begin{enumerate}[leftmargin=*,itemsep=0.35em]
\item \textbf{Ramsey's theorem.} Every finite coloring of a sufficiently large complete graph contains a monochromatic complete subgraph of a prescribed size.

\item \textbf{The party theorem.} $R(3,3)=6$.

\item \textbf{Van der Waerden's theorem.} Every finite coloring of the positive integers contains arbitrarily long monochromatic arithmetic progressions.

\item \textbf{Szemeredi's theorem.} A subset of the integers with positive upper density contains arbitrarily long arithmetic progressions.

\item \textbf{Hales--Jewett theorem.} Every finite coloring of a sufficiently large word cube contains a monochromatic combinatorial line.

\item \textbf{Rado's theorem.} A matrix criterion characterizes partition regularity for many finite systems of homogeneous linear equations.

\item \textbf{Erdős lower bound.} $R(k,k)>2^{k/2}$ for every $k\geq3$. In particular, the diagonal Ramsey numbers grow at least exponentially.
\end{enumerate}
\noindent\emph{Source for these results:} \cite{ramsey_ref1}.

\theoryapplications
\theoryconnections{ramsey_theory}{%
\item Combinatorics counts the possible configurations, graph theory represents pairwise relationships, and logic states the pattern that must be forced.
\item Probability supplies a powerful existence argument: if every random coloring has a positive chance of containing the desired monochromatic pattern, then at least one such coloring exists.
\item Number theory contributes additive patterns such as arithmetic progressions, where the same principle applies to sums rather than edges \cite{ramsey_ref1,ramsey_ref2}.
}{%
\item Ramsey theory gives extremal graph theory a language for the unavoidable structure inside a large network.
\item In additive combinatorics it proves that a sufficiently large set of integers contains a prescribed pattern, while in theoretical computer science it explains why large combinatorial objects can contain regular substructures that algorithms may search for.
\item The point is not to find the pattern efficiently in every case, but to prove that disorder has a definite limit \cite{ramsey_ref1,ramsey_ref3}.
}

\section*{Glossary}
\begin{description}[leftmargin=*,style=nextline,itemsep=0.3em]
\item[\textbf{Ramsey number.}] The smallest number $R(s,t)$ such that every red-blue coloring of the edges of a complete graph on that many vertices contains a red clique of size $s$ or a blue clique of size $t$.
\item[\textbf{Monochromatic.}] All edges or elements sharing the same color or label in a coloring or partition of a structure.
\item[\textbf{Partition regularity.}] A property of an equation or statement meaning that every finite coloring of the integers contains a monochromatic solution.
\item[\textbf{Van der Waerden's theorem.}] The result that any finite coloring of the positive integers contains arbitrarily long monochromatic arithmetic progressions.
\item[\textbf{Probabilistic method.}] A nonconstructive technique that proves a combinatorial object exists by showing a random construction has positive probability of satisfying the required conditions.
\item[\textbf{Complete graph.}] A graph in which every pair of vertices is connected by an edge, denoted $K_n$ for $n$ vertices.
\item[\textbf{Clique.}] A set of vertices in a graph that are all pairwise adjacent; a monochromatic clique in a colored complete graph is the target structure in classical Ramsey problems.
\item[\textbf{Arithmetic progression.}] A sequence $a, a+d, a+2d, \ldots$ with fixed common difference $d$.
\item[\textbf{Szemer\'{e}di's theorem.}] Every subset of the integers with positive upper density contains arbitrarily long arithmetic progressions.
\item[\textbf{Upper density.}] The asymptotic fraction $\limsup |A\cap\{1,\ldots,n\}|/n$.
\item[\textbf{Hales--Jewett theorem.}] A sufficiently high-dimensional word cube colored with finitely many colors contains a monochromatic combinatorial line.
\item[\textbf{Rado's theorem.}] A matrix criterion characterizing which systems of linear equations are partition regular.
\end{description}
\section*{7. Sample Questions}
\emph{Exercise reference:} \cite{ramsey_ref1}. The cited edition is the source for the exercise set; consult its Exercises section for the official statement and numbering.
\begin{enumerate}[leftmargin=*,itemsep=0.9em]
\item \textbf{Exercise 1.} Exhibit a red-blue coloring of the edges of $K_5$ that contains no monochromatic triangle, thereby proving $R(3,3)>5$. \emph{Solution.} Label the vertices $1,2,3,4,5$ and color the edges of the $5$-cycle $(12,23,34,45,51)$ red and the complementary $5$-cycle $(13,24,35,14,25)$ blue. Every triangle in $K_5$ contains at most two consecutive edges of either cycle, so no triangle is monochromatic.

\item \textbf{Exercise 2.} Prove that $R(2,3)=3$ by exhibiting a coloring of $K_2$ with no blue $K_3$ and showing every coloring of $K_3$ has a red $K_2$ or blue $K_3$. \emph{Solution.} With two vertices and one edge, a red edge gives a red $K_2$, and a blue edge gives no blue $K_3$. For $K_3$, if any edge is red we have a red $K_2$. If all three edges are blue we have a blue $K_3$. Either way, one of the two monochromatic subgraphs appears.

\item \textbf{Exercise 3.} Show that every two-coloring of $\{1,2,3,4,5,6,7,8,9\}$ must contain a monochromatic arithmetic progression of length $3$. (Van der Waerden's theorem gives $W(2,3)=9$.) \emph{Solution.} Partition $\{1,\ldots,9\}$ into three blocks $A=\{1,2,3\}$, $B=\{4,5,6\}$, $C=\{7,8,9\}$. By pigeonhole, one block has at least two elements of the same color, say red. If those two elements are endpoints of a 3-term AP within the block (e.g., $1$ and $3$, $4$ and $6$, or $7$ and $9$) and the midpoint is also red, we are done. A detailed case analysis (or the standard van der Waerden argument on $\{1,\ldots,9\}$) shows that avoiding a monochromatic 3-AP is impossible.

\item \textbf{Exercise 4.} Use the probabilistic method to show that $R(3,3)>5$ by computing the expected number of monochromatic triangles in a random two-coloring of $K_5$. \emph{Solution.} There are $\binom{5}{3}=10$ triangles in $K_5$. Each triangle is monochromatic (all red or all blue) with probability $2\cdot 2^{-3}=1/4$. The expected number of monochromatic triangles is $10\cdot 1/4=2.5$. This is positive, which does not directly give $R(3,3)>5$. Instead, compute for $K_5$: expected monochromatic triangles is $2.5>0$, but we need the expectation to be $<1$ for a guaranteed coloring. For $K_5$ with $n=5$ and $k=3$: $\binom{5}{3}\cdot 2^{1-\binom{3}{2}}=10\cdot 2^{-2}=2.5$. Since this exceeds $1$, the probabilistic method does not prove $R(3,3)>5$ by itself. We must exhibit the coloring directly, as in Exercise 1.

\item \textbf{Exercise 5.} Verify that $R(3,4)\leq 9$. In any two-coloring of $K_9$, prove a monochromatic red $K_3$ or blue $K_4$ must exist. \emph{Solution.} Pick any vertex $v$ in $K_9$. It has $8$ incident edges. By pigeonhole, at least $\lceil 8/2\rceil=4$ edges share a color. If $4$ or more are blue, let $S$ be the set of neighbors connected to $v$ by blue. If any two vertices in $S$ share a blue edge, we get a blue $K_3$ with $v$, and with a fourth vertex can extend to a blue $K_4$. If all edges within $S$ are red and $|S|\geq 4$, we get a red $K_3$ inside $S$ by the $R(3,3)$ argument (any $3$ of the $4$ vertices yield a red triangle). A careful case analysis confirms a monochromatic $K_3$ (red) or $K_4$ (blue) always appears, so $R(3,4)\leq 9$.

\end{enumerate}
\section*{8. References}
\begin{thebibliography}{9}
\bibitem{ramsey_ref1} R. L. Graham, B. L. Rothschild, and J. H. Spencer, \emph{Ramsey Theory}, 2nd ed., Wiley, 1990.
\bibitem{ramsey_ref2} R. Graham and S. Butler, \emph{Rudiments of Ramsey Theory}, American Mathematical Society, 1981.
\bibitem{ramsey_ref3} B. M. Landman and A. Robertson, \emph{Ramsey Theory on the Integers}, American Mathematical Society, 2004.
\end{thebibliography}


\clearpage\chapter{Random matrix theory}

\section*{1. Overview}
\subsection*{Intuitive}
\paragraph{The problem}
Suppose a data set has thousands of measurements and thousands of variables.
Its covariance matrix is then a huge table.  Even if the measurements contain
no real pattern, the table will have eigenvalues, large singular values, and
apparently important directions.  Which features are genuine, and which are
created by randomness?

The same question appears in physics.  A complicated quantum system may have
too many interacting particles for its exact energy levels to be calculated.
Instead of predicting every level, we can ask about the typical spacing of
nearby levels.  Random matrix theory was developed to answer such questions
about large matrices without pretending that every entry is known.


\subsection*{Formal}
An $N\times N$ random matrix is a matrix whose entries, or whose eigenvalue distribution, are governed by a specified probability law. Important observables include the eigenvalues $\lambda_1,\ldots,\lambda_N$, the empirical spectral measure $N^{-1}\sum_i\delta_{\lambda_i}$, and the normalized trace $N^{-1}\operatorname{Tr}f(M)$. The theory studies their limiting behavior as $N\to\infty$, with the ensemble and normalization stated explicitly.
\section*{2. History}
\subsection*{Historical development}
\begin{historylist}
\paragraph{History and the central idea}
Eugene Wigner introduced random matrices into nuclear physics in the 1950s.
He argued that the unknown interactions inside a heavy nucleus could be
modelled by a symmetric random matrix.  The individual entries were not
supposed to be literally random forces; randomness represented incomplete
knowledge and complicated interactions.  Freeman Dyson later organised the
main symmetry classes, and many mathematicians developed precise limit
theorems for eigenvalues and singular values \cite{random_matrix_theory_ref1}.

Let $X$ be an $m\times n$ matrix whose entries are independent measurements,
and form the sample covariance matrix
$S=(1/n)XX^{\mathsf T}$.  The eigenvectors of $S$ are the directions found by
principal component analysis.  If the entries of $X$ are only noise, the
eigenvalues are nevertheless not all equal.  For large $m$ and $n$, their
histogram approaches a predictable shape, the Marchenko--Pastur distribution,
when the ratio $m/n$ is held fixed.  A very large eigenvalue outside this
noise bulk is evidence for a signal, not merely evidence that a large matrix
always has a largest eigenvalue \cite{random_matrix_theory_ref2}.

\end{historylist}
\section*{3. Theory}
\subsection*{Core theory}
\paragraph{What the theory measures}
\bookfigure{random_matrix_spectrum}{Random matrix theory separates a typical noise spectrum from a possible signal outlier.}
There are three different questions, and confusing them leads to bad
interpretations.  The global eigenvalue distribution asks where most
eigenvalues lie.  The extreme-eigenvalue question asks how large the largest
one is.  The local question asks how close neighbouring eigenvalues are.

The answer depends on scaling and symmetry.  For a symmetric matrix with
entries of variance one, eigenvalues typically have size about the square
root of the matrix dimension.  Dividing by that scale produces a stable
limiting picture.  For a covariance matrix, the appropriate scale is set by
the number of observations and variables.  A theorem about one scaling
cannot be transferred to another without checking the assumptions.

\section*{4. Important theorems and results}
\begin{enumerate}[leftmargin=*,itemsep=0.35em]
\item \textbf{Marchenko--Pastur law.} For a large sample covariance matrix with
independent, suitably scaled entries, the empirical distribution of
eigenvalues approaches a deterministic limiting distribution.  In practice,
this gives a noise bulk against which unusually large principal components
can be compared.

\item \textbf{Wigner semicircle law.} For a large real symmetric random matrix with
centred entries of the appropriate variance, the rescaled eigenvalue
histogram approaches a semicircle.  The theorem explains why the detailed
distribution of individual entries often matters less than the matrix's
symmetry and scaling.

\item \textbf{Universality and edge behaviour.} Under broad conditions, local
eigenvalue spacings and extreme eigenvalues follow the same limiting laws
across many entry distributions.  Heavy tails and strong dependence can
invalidate this conclusion, so the hypotheses must be checked.

\end{enumerate}
\noindent\emph{Source for these results:} \cite{random_matrix_theory_ref1}.

\paragraph{Local statistics: the GUE nearest-neighbour spacing distribution}
Beyond the global eigenvalue density, random matrix theory characterizes
how individual eigenvalues are positioned relative to their neighbours. For
the Gaussian Unitary Ensemble (GUE)---the ensemble of Hermitian matrices
with independent complex Gaussian entries---the nearest-neighbour spacing
distribution $p(s)$ gives the probability density that the gap between two
adjacent eigenvalues equals $s$ after local rescaling. Unlike independent
random variables, whose gaps follow an exponential distribution, GUE
eigenvalues repel one another: $p(s)\sim(\pi s/2)e^{-\pi s^2/4}$ as
$s\to0$, so small gaps are suppressed. This level repulsion is a universal
feature of chaotic quantum systems and is experimentally observed in nuclear
spectra, mesoscopic conductance fluctuations, and zeros of the Riemann zeta
function \cite{random_matrix_theory_ref1}.

\paragraph{The Tracy--Widom distribution}
The largest eigenvalue of a large GUE matrix, after centring and rescaling,
converges to the Tracy--Widom distribution $F_2$. This distribution governs
the fluctuations of the spectral edge and appears in a remarkable range of
problems: the longest increasing subsequence of a random permutation, the
height of growth models in the Kardar--Parisi--Zhang universality class, and
the free energy of certain directed polymer models. The Tracy--Widom law
replaces the Gaussian distribution that governs sums of independent random
variables, reflecting the strong correlations induced by eigenvalue
repulsion \cite{random_matrix_theory_ref1,random_matrix_theory_ref2}.

\section*{5. Applications}
Random matrix theory studies the collective behaviour of eigenvalues and singular values when a matrix contains uncertainty. It provides reference distributions for separating stable structure from noise in large systems.
\begin{enumerate}[leftmargin=*,itemsep=0.7em]
\item \textbf{Wireless communication.}  In a multiple-input, multiple-output
(MIMO) system, a transmitter sends a vector of signals through many antenna
paths.  The receiver observes $y=Hx+z$, where $H$ is the channel matrix and
$z$ is noise.  Engineers want to know how many independent streams can be
sent reliably.  The singular values of $H$ are the strengths of those
streams, and the channel capacity contains terms of the form
$\log(1+\rho s_i^2)$, where $s_i$ are singular values and $\rho$ is the
signal-to-noise ratio.  Measuring every future channel matrix is impossible.
A random matrix model predicts the typical distribution of the $s_i$, so an
engineer can estimate capacity and the probability of a communication
outage before installing thousands of antennas.  The useful output is not
the entries of one invented matrix; it is a reliable estimate of the
distribution of possible channel qualities \cite{random_matrix_theory_ref1}.

For large antenna arrays with $n_t$ transmit and $n_r$ receive antennas,
when $n_t$ and $n_r$ grow proportionally with ratio $\gamma=n_t/n_r$, the
Marchenko--Pastur law describes the bulk of the singular-value spectrum of
$H$. The ergodic capacity per antenna approaches
\[
C=\int\log(1+\rho\,\lambda)\,dF(\lambda),
\]
where $F$ is the limiting singular-value distribution. The Tracy--Widom
distribution controls the fluctuations of the largest singular value, which
determines the dominant spatial channel and sets the peak data rate for a
single stream. These random-matrix predictions allow engineers to dimension
antenna arrays and power budgets without exhaustive field measurements,
capturing both typical performance and the probability of an unusually weak
or strong channel \cite{random_matrix_theory_ref1}.

\item \textbf{Separating signal from noise in data.}  Consider 10,000 patients
described by 1,000 laboratory measurements.  Principal component analysis
will always produce 1,000 eigenvalues, even if the measurements are
independent noise.  Random matrix theory gives a noise-only interval for
those eigenvalues.  Components inside the interval are consistent with
sampling noise; a component well outside it is a candidate for a shared
biological factor.  This is why the theory is useful: it supplies a
reference distribution against which a data scientist can test an apparent
pattern, instead of keeping the largest component merely because it is
largest.  The test must be recalibrated when measurements are strongly
dependent or have unequal variances \cite{random_matrix_theory_ref2}.

\item \textbf{Financial covariance matrices.}  A portfolio manager may estimate
the correlation of 500 assets from only 250 days of prices.  The sample
covariance matrix then has more variables than observations, and many
eigenvalues describe estimation noise rather than stable economic factors.
If the noisy covariance is inverted directly, the resulting portfolio may
place enormous weights on directions that will disappear next month.
Random matrix methods identify the noise-like bulk and replace it with a
smoothed bulk while retaining a few defensible large factors.  The resulting
matrix is better conditioned, so optimisation is less sensitive to
accidental features of the chosen historical window.  This does not predict
market prices; it improves the reliability of a matrix used by a separate
risk model.

\item \textbf{Energy levels in complex quantum systems.}  A heavy nucleus has many
interacting energy levels.  Wigner's random-matrix model predicts that nearby
levels tend to repel one another: two levels are less likely to be extremely
close than they would be in a list of independent random numbers.  Physicists
compare the rescaled spacings of measured levels with the predicted spacing
law.  Agreement is evidence that the system has the complicated mixing
expected of a chaotic quantum system; disagreement can reveal hidden
symmetry or a more regular mechanism.  The model describes statistical
regularities, not the exact energy of a particular level \cite{random_matrix_theory_ref1}.
\end{enumerate}

\theoryconnections{random_matrix_theory}{%
\item Probability supplies a distribution for the entries or eigenvalues, and linear algebra turns a matrix into eigenvalues, singular values, and eigenvectors.
\item Calculus and analysis then study the collective behavior of these quantities as the matrix size grows.
\item The important shift is from asking for one exact eigenvalue to asking for a distribution: for a large random matrix, the empirical eigenvalue plot often settles near a deterministic density, while the extreme eigenvalues fluctuate on a different scale \cite{random_matrix_theory_ref1,random_matrix_theory_ref2}.
}{%
\item In statistics, the singular values of a random data matrix provide a baseline for deciding whether a large observed component is signal or noise.
\item In wireless engineering, a channel matrix determines how many independent data streams can be transmitted, and random-matrix laws predict performance when antennas are numerous.
\item In physics, the same eigenvalue statistics model complicated energy levels, so the theory explains why unrelated systems can share the same fluctuation patterns \cite{random_matrix_theory_ref1,random_matrix_theory_ref2}.
}

\section*{Glossary}
\begin{description}[leftmargin=*,style=nextline,itemsep=0.3em]
\item[\textbf{Random matrix.}] A matrix whose entries or eigenvalue distribution are governed by a specified probability law, used to model systems with many unknown interactions.
\item[\textbf{Empirical spectral measure.}] The normalized counting distribution of eigenvalues $N^{-1}\sum_i\delta_{\lambda_i}$, approximating where eigenvalues of a large matrix tend to lie.
\item[\textbf{Marchenko--Pastur law.}] The limiting eigenvalue distribution of a large sample covariance matrix built from independent noise, defining the range of eigenvalues that noise alone would produce.
\item[\textbf{Wigner semicircle law.}] The theorem that the rescaled eigenvalue histogram of a large symmetric random matrix with centred entries approaches a semicircular shape.
\item[\textbf{Tracy--Widom distribution.}] The limiting distribution of the largest eigenvalue of certain random matrix ensembles, replacing the Gaussian in systems with strong eigenvalue correlations.
\item[\textbf{Level repulsion.}] The tendency of adjacent eigenvalues to avoid being close together, a universal feature of random matrix models and chaotic quantum systems.
\item[\textbf{Universality.}] The principle that local eigenvalue statistics and extreme eigenvalue behaviour follow the same limiting laws across broad classes of entry distributions.
\item[\textbf{Singular value.}] The square root of an eigenvalue of $M^*M$ for a matrix $M$; it measures the amplitude gain along each independent direction of a linear map.
\item[\textbf{Covariance matrix.}] A symmetric matrix whose $(i,j)$ entry is the correlation between variables $i$ and $j$.
\item[\textbf{Principal component analysis (PCA).}] A statistical technique rotating data into eigenvectors of the sample covariance matrix.
\item[\textbf{Gaussian Unitary Ensemble (GUE).}] The ensemble of $N\times N$ Hermitian matrices with independent complex Gaussian entries above the diagonal.
\item[\textbf{Sample covariance matrix.}] The matrix $S=(1/n)XX^{\mathsf T}$ computed from $n$ observations of $p$ variables.
\end{description}
\section*{7. Sample Questions}
\emph{Exercise reference:} \cite{random_matrix_theory_ref1}. The cited edition is the source for the exercise set; consult its Exercises section for the official statement and numbering.
\begin{enumerate}[leftmargin=*,itemsep=0.9em]
\item \textbf{Exercise 1.} Compute the eigenvalues of the $2\times 2$ real symmetric random matrix $M=\begin{pmatrix}1&2\\2&1\end{pmatrix}$. \emph{Solution.} The characteristic polynomial is $\det(M-\lambda I)=(1-\lambda)^2-4=\lambda^2-2\lambda-3=(\lambda-3)(\lambda+1)$. The eigenvalues are $\lambda_1=3$ and $\lambda_2=-1$.

\item \textbf{Exercise 2.} Let $X$ be a $3\times 2$ matrix with all entries i.i.d.\ standard normal $N(0,1)$. Compute the expected value of $\operatorname{Tr}(XX^T)$ and the expected sum of squared singular values of $X$. \emph{Solution.} $XX^T$ is $3\times 3$. $\operatorname{Tr}(XX^T)=\sum_{i,j}X_{ij}^2$. Each $X_{ij}^2$ has expectation $1$, and there are $3\times 2=6$ entries. So $E[\operatorname{Tr}(XX^T)]=6$. The sum of squared singular values of $X$ equals $\operatorname{Tr}(X^TX)=\operatorname{Tr}(XX^T)=6$, so the answer is $6$.

\item \textbf{Exercise 3.} For a $2\times 2$ Wigner matrix $W=\begin{pmatrix}a&b\\b&c\end{pmatrix}$ with $a,c\sim N(0,1)$ and $b\sim N(0,1/2)$ (entries independent), compute the mean and variance of $\operatorname{Tr}(W)=a+c$. \emph{Solution.} $E[\operatorname{Tr}(W)]=E[a]+E[c]=0+0=0$. $\operatorname{Var}(\operatorname{Tr}(W))=\operatorname{Var}(a)+\operatorname{Var}(c)=1+1=2$.

\item \textbf{Exercise 4.} A $4\times 4$ sample covariance matrix $S=\frac{1}{n}XX^T$ is formed from $n=10$ observations of $p=4$ variables. According to the Marchenko--Pastur law with ratio $\gamma=p/n=0.4$, what are the upper and lower edges of the noise bulk? \emph{Solution.} For the Marchenko--Pastur distribution with ratio $\gamma$, the support is $[(1-\sqrt{\gamma})^2,(1+\sqrt{\gamma})^2]$. With $\gamma=0.4$: lower edge $=(1-\sqrt{0.4})^2=(1-0.6325)^2\approx 0.1353$. Upper edge $=(1+\sqrt{0.4})^2=(1+0.6325)^2\approx 2.6649$. Eigenvalues of $S$ beyond approximately $2.66$ are candidates for signal.

\item \textbf{Exercise 5.} Let $A=\begin{pmatrix}1&0\\0&0\end{pmatrix}$. Compute the singular values of $A$ and verify that $\sum s_i^2=\operatorname{Tr}(AA^T)$. \emph{Solution.} $AA^T=A^2=\begin{pmatrix}1&0\\0&0\end{pmatrix}$. The eigenvalues of $AA^T$ are $1$ and $0$, so the singular values are $s_1=1$ and $s_2=0$. $\sum s_i^2=1^2+0^2=1$. $\operatorname{Tr}(AA^T)=1+0=1$. The identity holds.

\end{enumerate}
\section*{8. References}
\begin{thebibliography}{9}
\bibitem{random_matrix_theory_ref1} M. L. Mehta, \emph{Random Matrices}, 3rd ed., Academic Press, 2004.
\bibitem{random_matrix_theory_ref2} T. Tao, \emph{Topics in Random Matrix Theory}, American Mathematical Society, 2012.
\bibitem{random_matrix_theory_ref3} C. Tracy and H. Widom, ``Level-spacing distributions and the Airy kernel,'' \emph{Communications in Mathematical Physics}, vol.\ 159, pp.\ 151--174, 1994.
\end{thebibliography}


\theory{Representation theory}{How can an abstract symmetry be turned into concrete matrices that we can calculate with?}{Representation theory studies groups, algebras, and Lie algebras by letting their elements act as linear transformations on vector spaces. The abstract multiplication law becomes matrix multiplication, so algebraic questions can use linear algebra, geometry, and analysis.}{Fourier and harmonic analysis, quantum mechanics, particle physics, crystallography, number theory, algebraic geometry, and differential equations.}{Group actions on vector spaces translate abstract symmetry into matrices, characters, and decomposable modules.}

\paragraph{The idea and history}
A group describes reversible symmetries. A representation is a homomorphism $\rho:G\to GL(V)$ satisfying $\rho(gh)=\rho(g)\rho(h)$. It gives each abstract symmetry a matrix acting on a vector space. The same idea applies to associative algebras and Lie algebras \cite{representation_ref1}.

The first systematic examples came from finite groups and characters. Later work connected representations with Lie groups, geometry, number theory, and physics. Hermann Weyl developed important unitary representation methods, while Eugene Wigner used representations of spacetime symmetries to classify physical particles \cite{representation_ref2}.

\end{historylist}
\section*{3. Theory}
\subsection*{Core theory}
\paragraph{Finite groups and characters}
For a finite group over a field whose characteristic does not divide the group order, Maschke's theorem says every finite-dimensional representation splits into irreducible pieces. An irreducible representation has no nonzero proper invariant subspace. This is analogous to factoring an integer into primes: complicated representations are built from elementary blocks.

The character of a representation is the function $\chi(g)=\operatorname{tr}(\rho(g))$. Characters are constant on conjugacy classes and their inner products detect whether two irreducible representations are equal. Character tables therefore turn group classification questions into manageable calculations \cite{representation_ref1}.

\paragraph{Different settings}
Representations may be finite- or infinite-dimensional and may be defined over complex numbers, real numbers, finite fields, or $p$-adic fields. Modular representation theory studies the difficult case where the field characteristic divides the group order, so Maschke's theorem fails. Brauer's methods were important in the classification of finite simple groups.

Unitary representations act on Hilbert spaces and preserve inner products. They are central in quantum mechanics, where a symmetry changes a state without changing total probability. Representations of locally compact groups require continuity and measure, connecting the subject with harmonic analysis and operator theory \cite{representation_ref2}.

\paragraph{Lie groups and algebras}
Continuous symmetries are represented by Lie groups, while their infinitesimal motions are represented by Lie algebras. A Lie algebra has a bilinear bracket satisfying antisymmetry and the Jacobi identity. Differentiating a group representation gives a Lie algebra representation, often reducing a nonlinear symmetry problem to linear algebra.

The representation theory of semisimple Lie algebras uses roots, weights, highest-weight vectors, and Weyl groups. Tensor products describe how combined systems decompose into irreducible systems. This is the mathematical origin of angular momentum addition and the classification of many particle states \cite{representation_ref3}.

\paragraph{Connections and applications}
Fourier analysis is representation theory for translation groups: complex exponentials are the irreducible modes. Crystallography uses representations of space groups to describe patterns that repeat in a lattice. In quantum mechanics, observables and states transform according to representations of symmetry groups. In number theory, automorphic forms and the Langlands program connect representations with arithmetic information.

Representation theory also connects to invariant theory and the Erlangen view of geometry, in which a geometry is studied through its symmetry group. Category theory generalizes the idea by replacing an algebraic object with a category and a representation with a functor into vector spaces \cite{representation_ref4}.

\section*{4. Important theorems and results}
\begin{enumerate}[leftmargin=*,itemsep=0.35em]
\item \textbf{Maschke's theorem.} Finite-dimensional representations of a finite group split into irreducibles when the characteristic does not divide the group order.

\item \textbf{Schur's lemma.} A map between irreducible representations is either zero or an isomorphism; an endomorphism of an irreducible complex representation is scalar.

\item \textbf{Character orthogonality.} Irreducible characters are orthonormal under the class-function inner product.

\item \textbf{Peter--Weyl theorem.} Matrix coefficients of finite-dimensional unitary representations span a dense subspace of square-integrable functions on a compact group.

\item \textbf{Highest-weight classification.} Finite-dimensional irreducible representations of a complex semisimple Lie algebra are classified by dominant integral highest weights.

\item \textbf{Wigner's symmetry principle.} Quantum symmetries are represented by unitary or antiunitary transformations, allowing physical states to be classified by group representations.
\end{enumerate}
\noindent\emph{Source for these results:} \cite{representation_ref1}.

\theoryapplications
\theoryconnections{representation_theory}{%
\item A group describes how symmetries compose, and a representation realizes each symmetry as a linear map on a vector space.
\item Linear algebra can then decompose a complicated action into irreducible blocks, while topology and analysis extend the same idea to continuous groups and infinite-dimensional spaces.
\item The matrices are not a change of subject: they make the abstract action measurable, diagonalizable, and computable \cite{representation_ref1,representation_ref2}.
}{%
\item Representation theory helps harmonic analysis by expanding a function into components adapted to symmetry, as Fourier series do for rotations and translations.
\item In geometry it describes how coordinates change under a group action; in number theory it connects automorphic forms with group representations.
\item In physics, irreducible representations label elementary types of particles or states because combining symmetries must preserve the allowed transformation laws \cite{representation_ref1,representation_ref3}.
}
\section*{7. Sample Questions}
\emph{Exercise reference:} \cite{representation_ref1}. The cited edition is the source for the exercise set; consult its Exercises section for the official statement and numbering.
\begin{enumerate}[leftmargin=*,itemsep=0.9em]
\item \textbf{Exercise 1. What is the standard two-dimensional representation of the cyclic group of order four?} Send the generator to rotation by ninety degrees,
\[
\rho(g)=\begin{pmatrix}0&-1\\1&0\end{pmatrix}.
\]
Then $\rho(g)^4=I$, so the group relation is respected.

\item \textbf{Exercise 2. What is the character of a permutation representation?} It is the number of points fixed by the permutation, because the permutation matrix has one diagonal entry for each fixed point.

\item \textbf{Exercise 3. Why is a one-dimensional representation of a group $G$ a homomorphism to nonzero scalars?} The space is one-dimensional, so every invertible linear map is multiplication by a nonzero scalar. The group law requires $\rho(gh)=\rho(g)\rho(h)$.

\item \textbf{Exercise 4. Why is degeneracy useful in physics?} If several states share the same energy, they form an invariant subspace for the symmetry. Decomposing that space into irreducible representations predicts how a perturbation can split the energy levels.

\item \textbf{Exercise 5. What does Fourier analysis have to do with representation theory?} Translations form a group, and exponentials are its one-dimensional irreducible representations. Decomposing a signal into these modes is therefore decomposing it into symmetry modes.

\end{enumerate}
\section*{8. References}
\begin{thebibliography}{9}
\bibitem{representation_ref1} J.-P. Serre, \emph{Linear Representations of Finite Groups}.
\bibitem{representation_ref2} B. C. Hall, \emph{Lie Groups, Lie Algebras, and Representations}.
\bibitem{representation_ref3} W. Fulton and J. Harris, \emph{Representation Theory: A First Course}.
\bibitem{representation_ref4} P. Etingof et al., \emph{Introduction to Representation Theory}.
\end{thebibliography}


\theory{Ring theory}{What happens when addition and multiplication coexist in one algebraic system, even when multiplication is not reversible or commutative?}{Ring theory studies systems modeled on the integers. It studies ideals, homomorphisms, modules, polynomial identities, commutative and noncommutative rings, and the structure of special rings such as matrix rings and division rings.}{Algebraic geometry, algebraic number theory, coding theory, invariant theory, noncommutative geometry, linear algebra, and homological algebra.}{Addition, multiplication, ideals, and modules describe algebraic systems in which factorization and quotient structure can be studied together.}

\paragraph{The idea and history}
Integers can be added, subtracted, and multiplied. A ring keeps these operations but does not require every nonzero element to have a multiplicative inverse. Polynomial rings, matrices, functions, modular arithmetic, and coordinate rings are all examples. The theory developed from algebraic number theory, algebraic geometry, invariant theory, and attempts to generalize complex numbers \cite{ring_ref1}.

Hamilton's quaternions and other hypercomplex systems led mathematicians to matrix and noncommutative examples. Joseph Wedderburn identified finite-dimensional algebras with matrix structures, and Emil Artin generalized this work to Artinian rings. Emmy Noether's work on ideals and ascending chain conditions led to the term Noetherian ring and transformed modern algebra \cite{ring_ref2}.

\end{historylist}
\section*{3. Theory}
\subsection*{Core theory}
\paragraph{Rings, ideals, and quotients}
A ring $R$ has an abelian group structure under addition and an associative multiplication distributing over addition. It may have a multiplicative identity. A commutative ring satisfies $ab=ba$; a noncommutative ring need not.

An ideal $I$ is an additive subgroup absorbing multiplication: $r i$ and $i r$ belong to $I$ whenever $i\in I$ and $r\in R$. Ideals allow a quotient ring $R/I$, in which elements differing by something in $I$ are identified. A ring homomorphism has a kernel that is an ideal, and the first isomorphism theorem says the image is isomorphic to the quotient by that kernel \cite{ring_ref1}.

Prime and maximal ideals generalize prime numbers and fields. In a commutative ring, a prime ideal $P$ satisfies $ab\in P$ only when $a\in P$ or $b\in P$. A maximal ideal has a field as quotient. These notions are the algebraic basis of the spectrum of a ring in algebraic geometry.

\paragraph{Modules and important classes}
A module over $R$ is like a vector space in which scalars come from a ring instead of a field. Modules are the natural language for representations of rings and for systems of linear equations over the integers. A ring is Noetherian when every ascending chain of ideals stops, or equivalently when every ideal is finitely generated.

An integral domain is a commutative ring without zero divisors. A field is a domain in which every nonzero element is invertible. A principal ideal domain has every ideal generated by one element, while a Euclidean domain has a division algorithm. The integers form a Euclidean domain; polynomial rings over fields often do as well.

Noncommutative rings include square matrix rings, endomorphism rings, group rings, and universal enveloping algebras. Their modules can be more complicated because left and right ideals differ. Morita theory studies when two rings have equivalent module categories, showing that the category of modules can be more important than the ring's particular presentation \cite{ring_ref3}.

\paragraph{Structure and applications}
The Artin--Wedderburn theorem describes semisimple rings as products of matrix rings over division rings. Wedderburn's little theorem says every finite division ring is commutative, so it is a field. Nakayama's lemma controls generators of finitely generated modules over local rings. These results let algebraists replace a complicated ring by standard building blocks.

For an affine algebraic variety, regular functions form a coordinate ring. The geometry of the variety can be translated into ideal and quotient questions. The Hilbert Nullstellensatz connects maximal ideals of a polynomial ring over an algebraically closed field with points of the variety. Symmetric polynomials form an invariant ring; the fundamental theorem says they are generated by the elementary symmetric polynomials \cite{ring_ref4}.

\paragraph{Application: Reed--Solomon codes}
A concrete connection between ring theory and coding theory arises through Reed--Solomon codes. Let $\mathbb F_q$ be a finite field with $q$ elements, fix distinct elements $\alpha_1,\ldots,\alpha_n\in\mathbb F_q$ with $n\le q-1$, and choose an integer $k\le n$. The ring of polynomials $\mathbb F_q[x]$ provides the algebraic setting. A message is a polynomial $m(x)\in\mathbb F_q[x]$ of degree less than $k$. The encoder evaluates $m$ at the fixed points, producing the codeword
\[
c=(m(\alpha_1),m(\alpha_2),\ldots,m(\alpha_n)).
\]
The set of all such codewords forms a linear code of dimension $k$ and length $n$. Two distinct degree-$(k{-}1)$ polynomials agree on at most $k{-}1$ points, so any two codewords differ in at least $n{-}k{+}1$ positions. This minimum distance guarantees that the code can correct up to $\lfloor(n{-}k)/2\rfloor$ symbol errors.

The algebraic structure is essential: the set of codewords is an ideal in the ring $\mathbb F_q^n$ (under componentwise operations), and decoding can be performed by solving a system of polynomial congruences in $\mathbb F_q[x]$. Euclidean division in the polynomial ring provides the key algorithmic step. Reed--Solomon codes are used in QR codes, deep-space communication, and data storage systems such as CDs and DVDs \cite{ring_ref1}.

\section*{4. Important theorems and results}
\begin{enumerate}[leftmargin=*,itemsep=0.35em]
\item \textbf{Ring isomorphism theorems.} Kernels, images, quotients, and induced maps obey analogues of the group isomorphism theorems.

\item \textbf{Nakayama's lemma.} Over a local ring, generators of a finitely generated module can be detected after reducing modulo the maximal ideal.

\item \textbf{Artin--Wedderburn theorem.} A semisimple Artinian ring is a finite product of full matrix rings over division rings.

\item \textbf{Wedderburn's little theorem.} Every finite division ring is a field.

\item \textbf{Hilbert basis theorem.} If $R$ is Noetherian, then the polynomial ring $R[x]$ is Noetherian.

\item \textbf{Hilbert Nullstellensatz.} Over an algebraically closed field, maximal ideals of a polynomial ring correspond to points, and ideals describe the polynomial equations vanishing on sets.
\end{enumerate}
\noindent\emph{Source for these results:} \cite{ring_ref1}.

\theoryapplications
\theoryconnections{ring_theory}{%
\item Groups describe addition and multiplication separately, while ring theory studies when these operations interact through distributivity.
\item Modules generalize vector spaces when coefficients come from a ring rather than a field, and ideals identify which expressions should count as zero in a quotient.
\item Polynomial rings turn equations into ring elements, so factoring, divisibility, and quotienting become tools for studying geometric solution sets \cite{ring_ref1,ring_ref2}.
}{%
\item Ring theory helps algebraic geometry by translating a variety into its coordinate ring, where geometric inclusion becomes ideal inclusion.
\item In number theory, ideals restore a workable factorization theory in rings where elements do not factor uniquely.
\item Modules support representations and homological algebra, while quotient rings over finite fields provide the arithmetic used to construct and decode error-correcting codes \cite{ring_ref1,ring_ref3}.
}
\section*{7. Sample Questions}
\emph{Exercise reference:} \cite{ring_ref1}. The cited edition is the source for the exercise set; consult its Exercises section for the official statement and numbering.
\begin{enumerate}[leftmargin=*,itemsep=0.9em]
\item \textbf{Exercise 1. Find the ideals of $\mathbb Z$.} Every ideal is $n\mathbb Z$ for a unique nonnegative integer $n$. If an ideal contains a least positive element $n$, division with remainder shows every element is a multiple of $n$.

\item \textbf{Exercise 2. Describe $\mathbb Z/6\mathbb Z$.} Its elements are the six residue classes. The class of $2$ is a zero divisor because $2\cdot3=0$ modulo $6$, so the quotient is not a field.

\item \textbf{Exercise 3. Why is $2\mathbb Z$ an ideal of $\mathbb Z$?} It is closed under addition and negatives, and multiplying an even number by any integer remains even.

\item \textbf{Exercise 4. Is $\mathbb Z/5\mathbb Z$ a field?} Yes. Every nonzero class has an inverse because $5$ is prime; for example $2^{-1}=3$ since $2\cdot3=6\equiv1$ modulo $5$.

\item \textbf{Exercise 5. Why are matrices a noncommutative ring?} Matrix addition and multiplication satisfy the ring laws, but usually $AB\ne BA$. The failure of commutativity is exactly why left and right ideals and modules must be distinguished.

\end{enumerate}
\section*{8. References}
\begin{thebibliography}{9}
\bibitem{ring_ref1} T. Y. Lam, \emph{A First Course in Noncommutative Rings}, 2nd ed., Springer, 2001.
\bibitem{ring_ref2} C. Faith, \emph{Rings and Things and a Fine Array of Twentieth-Century Associative Algebra}, AMS, 2004.
\bibitem{ring_ref3} T. Y. Lam, \emph{Lectures on Modules and Rings}, Springer, 1999.
\bibitem{ring_ref4} D. Eisenbud, \emph{Commutative Algebra with a View Toward Algebraic Geometry}, Springer, 1995.
\end{thebibliography}


\theory{Scheme theory}{How can algebraic geometry describe equations over arbitrary rings, remember multiplicities, and retain information lost when a field has too few visible points?}{A scheme is a space built from commutative rings. An affine scheme is the spectrum of a ring, whose points are prime ideals; a general scheme is obtained by gluing affine pieces together with a sheaf of functions. Schemes generalize varieties and unite geometry, topology, commutative algebra, and number theory.}{Algebraic geometry, arithmetic geometry, moduli spaces, number theory, the Weil conjectures, and Fermat's Last Theorem.}{Prime spectra, localization, and sheaves preserve algebraic, infinitesimal, and arithmetic information in a geometric form.}

\paragraph{Alexander Grothendieck}
Grothendieck introduced schemes in the 1960s to unify algebraic geometry over arbitrary rings. By using prime spectra and sheaves, he made nilpotents, generic points, and families part of the geometric object rather than information discarded by its visible points.

\paragraph{Pierre Deligne and arithmetic geometers}
Deligne used scheme-theoretic and cohomological methods to prove the remaining Weil conjecture. This demonstrated that schemes are not only a general language for varieties but also a working framework for arithmetic geometry \cite{scheme_ref1,scheme_ref3}.
\end{historylist}
\section*{3. Theory}
\subsection*{Core theory}
\paragraph{Why varieties are not enough}
An algebraic variety is the set of solutions to polynomial equations, often over the complex numbers. This picture loses information over smaller fields. The equation $x^2+y^2=-1$ has no real points but has complex points. A scheme over the real numbers remembers how the equations behave over every field extension, not only what is visible in the real plane \cite{scheme_ref1}.

Schemes also remember multiplicity. The equations $x=0$ and $x^2=0$ have the same visible solution set, but the second has a doubled structure. That extra information matters when studying intersections, tangent directions, and families in which components merge \cite{scheme_ref2}.

\paragraph{Alexander Grothendieck and the spectrum}
Alexander Grothendieck introduced schemes in the 1960s in \emph{Elements of Algebraic Geometry}. He wanted a flexible language for algebraic geometry over arbitrary bases, including the integers, and for deep problems such as the Weil conjectures. Pierre Deligne later proved the last Weil conjecture using this framework \cite{scheme_ref1,scheme_ref3}.

For a commutative ring $R$, the affine scheme $\operatorname{Spec}(R)$ is the set of prime ideals of $R$. It has the Zariski topology, in which closed sets are
\[
V(I)=\{\mathfrak p:I\subseteq\mathfrak p\}.
\]
Maximal ideals behave like ordinary points, while non-maximal prime ideals are generic points of irreducible subspaces. The structure sheaf assigns a ring of regular functions to each open set. On a basic open set where $f$ is nonzero, the functions are obtained by allowing powers of $f$ in denominators \cite{scheme_ref2}.

\paragraph{Gluing affine pieces}
A scheme is a locally ringed space covered by affine schemes. The rings on overlapping pieces must agree through localization maps. This is analogous to describing the surface of a sphere with several maps: no single coordinate chart covers the whole surface, but compatible charts define the object.

The coordinate ring of $\operatorname{Spec}(R)$ is $R$ itself. A ring map $B\to A$ induces a scheme map $\operatorname{Spec}(A)\to\operatorname{Spec}(B)$ in the opposite direction. This contravariance is essential: functions pull back along geometric maps \cite{scheme_ref1}.

\paragraph{Morphisms, base change, and families}
The relative viewpoint studies a morphism $X\to S$ rather than an isolated scheme. A scheme over a ring $R$ is a morphism $X\to\operatorname{Spec}(R)$. Base change replaces $R$ by another ring and asks how the family changes. This lets one study the same equations over the integers, a finite field, the real numbers, and the complex numbers.

Moduli spaces classify geometric objects such as curves, vector bundles, or elliptic curves. A family of objects can itself have geometric structure, and schemes are designed to represent the functor that records all families. When objects have automorphisms, stacks extend schemes and algebraic spaces by remembering the symmetry groups \cite{scheme_ref4}.

\paragraph{Arithmetic and geometry}
The spectrum of the integers contains a generic point together with one point for each prime number. Thus arithmetic is represented as geometry. A scheme over the integers can be viewed as a family whose fibers over primes are reductions modulo $p$, while the generic fiber describes characteristic zero. This viewpoint links congruences, prime numbers, and geometric families.

The language of schemes underlies modern proofs in arithmetic geometry. Wiles's proof of Fermat's Last Theorem used modular forms, elliptic curves, and deformation theory in a setting naturally expressed through rings and schemes. The point is not that a scheme is a picture in space; it is a controlled system of local functions that preserves algebraic information \cite{scheme_ref3}.

\section*{4. Important theorems and results}
\begin{enumerate}[leftmargin=*,itemsep=0.35em]
\item \textbf{Affine correspondence.} Commutative rings and affine schemes are contravariantly equivalent: ring maps $B\to A$ correspond to morphisms $\operatorname{Spec}(A)\to\operatorname{Spec}(B)$.

\item \textbf{Hilbert Nullstellensatz.} For polynomial rings over algebraically closed fields, maximal ideals correspond to geometric points.

\item \textbf{Sheaf gluing principle.} Compatible affine pieces and their transition maps determine a scheme.

\item \textbf{Base-change principle.} A morphism can be pulled back along another morphism, producing the corresponding geometric family over a new base.

\item \textbf{Grothendieck's cohomological framework.} Scheme cohomology turns local function data into global invariants and supports duality, intersection theory, and the Weil conjectures.

\item \textbf{Artin representability.} Under suitable conditions, a functor of families is represented by an algebraic space, and stack theory extends this to objects with automorphisms.
\end{enumerate}
\noindent\emph{Source for these results:} \cite{scheme_ref1}.

\theoryapplications
\theoryconnections{scheme_theory}{%
\item Commutative algebra supplies rings and ideals, and localization zooms in by allowing selected elements to become invertible.
\item The prime ideals form a space whose open sets reflect algebraic information, while a sheaf attaches the functions available on each open set.
\item This local-to-global construction lets a scheme remember nilpotent or arithmetic information that ordinary point sets would discard; homological algebra measures whether local data glue consistently \cite{scheme_ref1,scheme_ref2}.
}{%
\item Schemes help arithmetic geometry treat a ring of integers and a polynomial curve in one language, so statements can be compared across different primes and over the complex numbers.
\item They provide the base spaces for moduli problems, and their sheaves support intersection and deformation calculations.
\item In number theory, reduction modulo a prime becomes a fiber of a morphism, making local behavior part of a global geometric object \cite{scheme_ref1,scheme_ref3}.
}
\section*{7. Sample Questions}
\emph{Exercise reference:} \cite{scheme_ref1}. The cited edition is the source for the exercise set; consult its Exercises section for the official statement and numbering.
\begin{enumerate}[leftmargin=*,itemsep=0.9em]
\item \textbf{Exercise 1. What are the points of $\operatorname{Spec}(\mathbb Z)$?} They are the prime ideals $(0)$ and $(p)$ for each prime integer $p$. The ideals $(p)$ are closed points, while $(0)$ is generic.

\item \textbf{Exercise 2. Why is $\operatorname{Spec}(A)$ contravariant?} A ring map $B\to A$ lets a prime ideal of $A$ pull back to a prime ideal of $B$. Therefore the induced geometric map goes from $\operatorname{Spec}(A)$ to $\operatorname{Spec}(B)$.

\item \textbf{Exercise 3. Compare $x=0$ and $x^2=0$.} Both have the same closed point, but $k[x]/(x^2)$ contains a nonzero nilpotent element $x$. The scheme remembers the doubled or infinitesimal structure.

\item \textbf{Exercise 4. What is $\operatorname{Spec}(k[x])$ over an algebraically closed field?} Its maximal ideals $(x-a)$ correspond to points $a$ of the affine line. The zero ideal is a generic point whose closure is the whole line.

\item \textbf{Exercise 5. Why study schemes over $\mathbb Z$?} A scheme over $\mathbb Z$ packages one characteristic-zero object together with its reductions modulo every prime, allowing arithmetic information to be studied as a family of fibers.

\end{enumerate}
\section*{8. References}
\begin{thebibliography}{9}
\bibitem{scheme_ref1} R. Hartshorne, \emph{Algebraic Geometry}.
\bibitem{scheme_ref2} D. Eisenbud and J. Harris, \emph{The Geometry of Schemes}.
\bibitem{scheme_ref3} A. Grothendieck and J. Dieudonne, \emph{Elements de geometrie algebrique}.
\bibitem{scheme_ref4} D. Mumford, \emph{The Red Book of Varieties and Schemes}.
\end{thebibliography}


\theory{Semigroup theory}{What algebraic structure remains when a composition operation is associative but elements need not have inverses?}{A semigroup is a set with an associative internal operation. It captures processes that can be performed in sequence, even when a process cannot be undone. Semigroup theory studies ideals, Green's relations, representations, finite structure, and links with automata and operators.}{String processing, finite automata, Markov processes, partial differential equations, linear systems, and the algebra of irreversible computation.}{Associative multiplication without universal inverses models irreversible composition, transition systems, and transformation actions.}

\paragraph{Alexander Sushkevich and David Rees}
Sushkevich studied the structure of finite semigroups, while Rees described important completely simple semigroups through matrix-like constructions. Their work showed how associative composition can be analyzed without assuming inverses.

\paragraph{James Alexander Green and automata theorists}
Green introduced the relations that organize semigroup elements by the principal ideals they generate. Later work connected these relations with finite automata, transformation systems, and operator theory, where irreversible composition is the central example \cite{semigroup_ref1}.
\end{historylist}
\section*{3. Theory}
\subsection*{Core theory}
\paragraph{Definition and examples}
A semigroup is a set $S$ with a product satisfying
\[
(ab)c=a(bc)
\]
for all $a,b,c\in S$. No identity and no inverse are required. String concatenation is associative: grouping the joins differently does not change the final string. Functions under composition form a semigroup, as do square matrices under multiplication and positive integers under addition.

A monoid is a semigroup with an identity. A group is a monoid in which every element has an inverse. Thus semigroups generalize groups by allowing irreversible processes such as a reset button, a lossy computation, or a function that is not reversible \cite{semigroup_ref1}.

\paragraph{History and structure}
The systematic study began in the early twentieth century. Anton Sushkevich studied finite simple semigroups; David Rees and James Green developed structural tools; Alfred Clifford and Gordon Preston published foundational monographs. Green's relations divide a semigroup according to the principal ideals generated by its elements \cite{semigroup_ref2}.

An ideal is a subset absorbing multiplication from the relevant side. Minimal ideals describe the long-term core of a finite semigroup. In a finite semigroup, repeated products eventually repeat, so powers of an element enter a periodic part. This finite behavior is useful when a process is applied repeatedly.

\paragraph{Representations and finite semigroups}
Cayley's theorem for semigroups represents every semigroup as a semigroup of transformations of a set. An element is therefore a function, and multiplication is composition. This generalizes the group representation by permutations, replacing reversible bijections with arbitrary maps.

Finite semigroups are closely related to finite automata. A word read by an automaton induces a transformation of its state set. Words with the same effect form a syntactic monoid, and algebraic properties of this monoid characterize language classes. Krohn--Rhodes theory decomposes finite semigroups into simpler group and aperiodic components, analogous in spirit to factoring \cite{semigroup_ref3}.

\paragraph{Analysis and differential equations}
A one-parameter operator semigroup is a family $T(t)$ with $T(0)=I$ and $T(t+s)=T(t)T(s)$ for $t,s\geq0$. It describes evolution forward in time. Unlike a group, negative time may not be possible. If $A$ is the generator, formally $T(t)=e^{tA}$.

The heat equation is an example. Starting with a temperature distribution $u(0)$, the heat semigroup sends it to $u(t)$ for later times. Diffusion smooths irregularities and loses information, so reversing it is unstable. Semigroup methods make existence, uniqueness, and stability of evolution equations precise \cite{semigroup_ref4}.

\section*{4. Important theorems and results}
\begin{enumerate}[leftmargin=*,itemsep=0.35em]
\item \textbf{Cayley representation theorem.} Every semigroup is isomorphic to a semigroup of transformations under composition.

\item \textbf{Green's relations.} The relations $\mathcal L,\mathcal R,\mathcal H,\mathcal D,\mathcal J$ organize elements by their principal left, right, and two-sided ideals.

\item \textbf{Krohn--Rhodes theorem.} Finite semigroups admit decompositions into group components and aperiodic components in a precise wreath-product sense.

\item \textbf{Syntactic monoid theorem.} A regular language has a finite syntactic monoid, and its algebra records the language's behavior under contexts.

\item \textbf{Hille--Yosida principle.} Suitable conditions on an operator characterize when it generates a strongly continuous semigroup.
\end{enumerate}
\noindent\emph{Source for these results:} \cite{semigroup_ref1}.

\theoryapplications
\theoryconnections{semigroup_theory}{%
\item A semigroup keeps the associative composition law while allowing inverses to be absent, which is exactly what irreversible processes require.
\item Functions under composition form the basic example: applying one transition after another gives a new transition, but it may not be reversible.
\item Automata use finite transformation semigroups to record the effect of reading words, while probability and functional analysis use semigroups to describe repeated or continuous evolution \cite{semigroup_ref1,semigroup_ref2}.
}{%
\item In formal-language theory, the syntactic semigroup compresses all words that have the same effect on acceptance, turning an infinite language question into a finite algebraic one when possible.
\item In a Markov process, a transition semigroup records the distribution after each elapsed time.
\item For a PDE such as heat flow, an operator semigroup maps the initial temperature to its later temperature, so composition expresses the passage of time \cite{semigroup_ref1,semigroup_ref3}.
}
\section*{7. Sample Questions}
\emph{Exercise reference:} \cite{semigroup_ref1}. The cited edition is the source for the exercise set; consult its Exercises section for the official statement and numbering.
\begin{enumerate}[leftmargin=*,itemsep=0.9em]
\item \textbf{Exercise 1. Why is string concatenation a semigroup?} Concatenating three strings in either order of parentheses produces the same sequence of symbols, so associativity holds.

\item \textbf{Exercise 2. Are all functions from a set to itself a monoid?} Yes. Composition is associative and the identity function is the identity element. It is generally not a group because most functions are not invertible.

\item \textbf{Exercise 3. What is the identity in the matrix semigroup of all $2$ by $2$ matrices?} The identity matrix is an identity for multiplication, but the set is not a group because singular matrices have no inverse.

\item \textbf{Exercise 4. Why is a heat evolution naturally a semigroup rather than a group?} Heat moves forward and smooths information. A later temperature profile usually does not determine a stable earlier profile, so backward evolution is not available as a bounded inverse.

\item \textbf{Exercise 5. How does an automaton produce a semigroup?} Each input word acts on the finite state set by a transformation. Combining words concatenates their transformations, giving a semigroup under composition.

\end{enumerate}
\section*{8. References}
\begin{thebibliography}{9}
\bibitem{semigroup_ref1} A. H. Clifford and G. B. Preston, \emph{The Algebraic Theory of Semigroups}.
\bibitem{semigroup_ref2} J. M. Howie, \emph{Fundamentals of Semigroup Theory}.
\bibitem{semigroup_ref3} J.-E. Pin, \emph{Mathematical Foundations of Automata Theory}.
\bibitem{semigroup_ref4} E. B. Davies, \emph{One-Parameter Semigroups}.
\end{thebibliography}


\theory{Set theory}{How can mathematics describe collections, compare different sizes of infinity, and build numbers and structures from simple membership rules?}{Set theory studies sets and membership. It provides a language and foundation for most mathematics, while also investigating infinite sizes, ordinals, axioms, paradoxes, descriptive hierarchies, large cardinals, and independence.}{Foundations, analysis, algebra, topology, computer science, formal semantics, logic, and the study of infinity.}{Membership, union, power sets, and cardinality specify the foundations from which mathematical collections are constructed.}

\paragraph{The idea and history}
A set is a collection of objects, and membership is written $x\in A$. The empty set has no members. From sets we build pairs, products, functions, relations, numbers, spaces, and algebraic structures. Richard Dedekind and Georg Cantor initiated modern set theory in the 1870s; Cantor is commonly regarded as its founder \cite{set_ref1}.

Cantor compared sizes by one-to-one correspondences. The natural numbers and even numbers have the same cardinality because $n\mapsto2n$ is a bijection. The real numbers are uncountable: no list can contain every real number. This discovery made infinity a mathematical object rather than merely an endless process \cite{set_ref2}.

\end{historylist}
\section*{3. Theory}
\subsection*{Core theory}
\paragraph{Operations and cardinality}
The power set $\mathcal P(A)$ contains every subset of $A$. Cantor's theorem says there is no bijection from $A$ onto $\mathcal P(A)$; the power set is strictly larger. A diagonal argument constructs a subset missing from any proposed list.

Cardinal numbers measure size. Finite cardinals behave like ordinary counting numbers; $\aleph_0$ is the size of the natural numbers. Ordinals measure order type: $\omega$ is the first infinite ordinal, followed by $\omega+1$, $\omega+2$, and so on. Cardinal and ordinal arithmetic differ because order matters for ordinals \cite{set_ref3}.

\paragraph{Axioms and paradoxes}
Naive set formation leads to contradictions. Russell's paradox considers the collection of all sets that do not contain themselves. If it contains itself, it should not; if it does not, it should. Other paradoxes include Cantor's paradox and the Burali--Forti paradox.

Zermelo--Fraenkel set theory, usually ZF, avoids these contradictions through axioms such as extensionality, pairing, union, power set, infinity, separation, replacement, foundation, and choice when AC is included. The axiom of choice says that a family of nonempty sets has a choice function. It is equivalent over ZF to Zorn's lemma and the well-ordering theorem \cite{set_ref1}.

\paragraph{The set-theoretic universe}
The cumulative hierarchy builds sets in stages. Start with the empty set, take power sets and unions repeatedly, and at limit stages take unions of earlier stages. The resulting von Neumann universe is written $V$. Each set has a rank, measuring the stage at which it appears.

The constructible universe $L$ is an inner model built using definable subsets. Forcing, introduced by Paul Cohen, constructs models in which selected statements differ in truth value. Cohen proved that the continuum hypothesis and the axiom of choice have independence phenomena relative to ZF and ZFC. Thus axioms can leave some questions undecided \cite{set_ref4}.

\paragraph{Areas of study and applications}
Descriptive set theory studies definable subsets of the real line and Polish spaces through the Borel and projective hierarchies. Effective descriptive set theory connects these classes with recursion theory. Large-cardinal theory studies strong forms of infinity and their consistency strength. Set theory also supplies the formal basis for graphs, manifolds, rings, vector spaces, relations, and databases.

Category theory offers a different foundation: it emphasizes objects and relationships rather than membership. Set theory and category theory inform one another, but they organize mathematics in different ways \cite{set_ref5}.

\section*{4. Important theorems and results}
\begin{enumerate}[leftmargin=*,itemsep=0.35em]
\item \textbf{Cantor's theorem.} For every set $A$, $|\mathcal P(A)|>|A|$.

\item \textbf{Cantor diagonal theorem.} The real numbers are uncountable.

\item \textbf{Well-ordering theorem.} Every set can be well ordered, assuming the axiom of choice.

\item \textbf{Zorn's lemma.} A nonempty poset in which every chain has an upper bound has a maximal element.

\item \textbf{Löwenheim--Skolem phenomena.} First-order theories with infinite models have models of many cardinalities, showing that formal descriptions do not determine one intended universe.

\item \textbf{Independence results.} Some statements, including the continuum hypothesis, cannot be decided by the standard axioms without additional assumptions, relative to consistency results.
\end{enumerate}
\noindent\emph{Source for these results:} \cite{set_ref1}.

\theoryapplications
\theoryconnections{set_theory}{%
\item Set theory supplies a common way to collect objects, form products, and define functions and relations.
\item Logic specifies which collections and constructions are allowed, while proof checks that an operation preserves the required membership rules.
\item For example, a function is a set of ordered pairs with exactly one output for each input; this definition lets algebra, analysis, and computer science use the same concept without relying on a picture \cite{set_ref1,set_ref2}.
}{%
\item Set theory helps algebra and analysis build numbers, spaces, sequences, and functions from explicit objects.
\item It lets probability define a sample space and measurable events, and lets computer science describe data types and infinite computations.
\item Category theory then studies the maps between these constructions, while foundational results identify which collections can be formed without contradiction \cite{set_ref1,set_ref3}.
}

\section*{Glossary}
\begin{description}[leftmargin=*,style=nextline,itemsep=0.3em]
\item[\textbf{Cardinality.}] A measure of the size of a set; two sets have the same cardinality when a bijection exists between them, even if both are infinite.
\item[\textbf{Ordinal.}] An order type that describes not just how many elements a well-ordered set has but also the pattern in which they are arranged.
\item[\textbf{Power set.}] The set $\mathcal{P}(A)$ of all subsets of $A$; Cantor's theorem shows it is always strictly larger than $A$ itself.
\item[\textbf{Axiom of choice.}] The principle that every family of nonempty sets has a choice function selecting one element from each; it is equivalent to Zorn's lemma and the well-ordering theorem.
\item[\textbf{Well-ordering.}] An ordering of a set in which every nonempty subset has a least element; the well-ordering theorem states every set can be well-ordered using the axiom of choice.
\item[\textbf{Forcing.}] A technique introduced by Cohen for constructing models of set theory in which selected statements such as the continuum hypothesis have a prescribed truth value.
\item[\textbf{Cumulative hierarchy.}] The layered construction $V$ of the set-theoretic universe built by iteratively forming power sets and unions from the empty set.
\item[\textbf{Independence.}] A statement is independent of the standard axioms when neither it nor its negation can be proved from them, as shown for the continuum hypothesis relative to ZFC.
\item[\textbf{ZFC.}] Zermelo--Fraenkel set theory with the Axiom of Choice; the standard axiom system for most of mathematics.
\item[\textbf{Continuum hypothesis.}] The statement that there is no set whose cardinality lies strictly between the naturals and the reals; independent of ZFC.
\item[\textbf{Russell's paradox.}] The contradiction from ``the set of all sets that do not contain themselves,'' motivating axiomatic set theory.
\item[\textbf{Constructible universe.}] The inner model $L$ built from definable subsets; satisfies ZFC plus the continuum hypothesis.
\item[\textbf{Bijection.}] A one-to-one and onto function; two sets have the same cardinality precisely when a bijection exists.
\end{description}
\section*{7. Sample Questions}
\emph{Exercise reference:} \cite{set_ref1}. The cited edition is the source for the exercise set; consult its Exercises section for the official statement and numbering.
\begin{enumerate}[leftmargin=*,itemsep=0.9em]
\item \textbf{Exercise 1.} Prove that $|\mathbb{N}|=|\mathbb{E}|$ where $\mathbb{E}=\{2,4,6,\ldots\}$ is the set of even positive integers. \emph{Solution.} Define $f\colon\mathbb{N}\to\mathbb{E}$ by $f(n)=2n$. This is injective: $2m=2n\Rightarrow m=n$. It is surjective: every even number $2k$ is $f(k)$. So $f$ is a bijection and $|\mathbb{N}|=|\mathbb{E}|$.

\item \textbf{Exercise 2.} Show that $\mathcal{P}(\{1,2,3\})$ has $8$ elements and that $8>3$, illustrating Cantor's theorem for $A=\{1,2,3\}$. \emph{Solution.} $\mathcal{P}(\{1,2,3\})=\{\varnothing,\{1\},\{2\},\{3\},\{1,2\},\{1,3\},\{2,3\},\{1,2,3\}\}$. This is $2^3=8$ elements. Since $8>3=|\{1,2,3\}|$, there is no surjection from $\{1,2,3\}$ onto $\mathcal{P}(\{1,2,3\})$, consistent with Cantor's theorem.

\item \textbf{Exercise 3.} Use the diagonal argument to show that the set of infinite binary sequences is uncountable. \emph{Solution.} Suppose $s_1,s_2,s_3,\ldots$ is a list of all binary sequences. Construct $t$ by setting the $n$-th digit of $t$ to $1-s_n(n)$ (flip the $n$-th digit of $s_n$). Then $t$ differs from $s_n$ in position $n$ for every $n$. So $t$ is not on the list, contradicting the assumption that the list was complete. Therefore no countable list covers all binary sequences.

\item \textbf{Exercise 4.} Compute the cardinality of $\mathcal{P}(\mathbb{N})$ and show it equals $|\mathbb{R}|$. \emph{Solution.} $|\mathcal{P}(\mathbb{N})|=2^{\aleph_0}$. Each subset of $\mathbb{N}$ corresponds to a binary sequence (the indicator function), so $|\mathcal{P}(\mathbb{N})|=|\{0,1\}^{\mathbb{N}}|$. Each binary sequence corresponds to a real number in $[0,1]$ via the binary expansion, giving $|\{0,1\}^{\mathbb{N}}|=|\mathbb{R}|$. (Some care is needed with dual representations, but the cardinalities match.) So $|\mathcal{P}(\mathbb{N})|=|\mathbb{R}|=2^{\aleph_0}$.

\item \textbf{Exercise 5.} Verify the ordinal arithmetic: compute $1+\omega$ and $\omega+1$ in the ordinals, and explain why they differ. \emph{Solution.} $1+\omega=\omega$ because placing one element before $\omega$ yields an order type isomorphic to $\omega$ (the single element is absorbed). $\omega+1>\omega$ because $\omega+1$ has a largest element not present in $\omega$. These are different ordinals because ordinal addition is not commutative.

\end{enumerate}
\section*{8. References}
\begin{thebibliography}{9}
\bibitem{set_ref1} P. Halmos, \emph{Naive Set Theory}, Springer, 1960.
\bibitem{set_ref2} J. W. Dauben, \emph{Georg Cantor: His Mathematics and Philosophy of the Infinite}, Princeton, 1979.
\bibitem{set_ref3} T. Jech, \emph{Set Theory}, 3rd ed., Springer, 2003.
\bibitem{set_ref4} A. Kanamori, \emph{The Higher Infinite}, 2nd ed., Springer, 2008.
\bibitem{set_ref5} S. Awodey, \emph{Category Theory}, Oxford, 2010.
\end{thebibliography}


\theory{Shape theory}{How can we describe the global shape of spaces that are too wild for ordinary homotopy theory?}{Shape theory approximates complicated compact spaces by simpler polyhedra and records the compatible information that survives the approximation.}{General topology, dynamical systems, inverse limits, and noncommutative geometry.}{Inverse systems of simpler spaces preserve stable large-scale information when ordinary homotopy cannot handle a wild space.}

\paragraph{The problem and history}
Ordinary homotopy theory is powerful for polyhedra and manifolds, but wild spaces can have bad local behavior. Shape theory was introduced by D. E. Christie in 1944 and developed by Karol Borsuk, who coined the name in 1968 \cite{shape_ref1}.

\end{historylist}
\section*{3. Theory}
\subsection*{Core theory}
\paragraph{Approximation by polyhedra}
The method replaces a compact space by a system of simpler polyhedral approximations. Each approximation has homotopy groups and homology, and maps between approximations record how descriptions refine one another. The shape is the resulting pro-system rather than one isolated approximation.

For compact spaces dominated homotopically by finite polyhedra, shape theory and homotopy theory agree. For general compact spaces, shape is more stable because it ignores microscopic details that disappear under increasingly fine approximation \cite{shape_ref2}.

\paragraph{Concrete illustration: the Hawaiian earring vs.\ a wedge of circles}
The Hawaiian earring $H$ is the union of circles $C_n$ of radius $1/n$ all meeting at a single point $p$. The infinite wedge $\bigvee_{n=1}^{\infty}S^1$ is the disjoint union of countably many circles identified at one point, given the topology that makes each circle an open-and-closed subset.

\textbf{Ordinary homotopy.} A loop in $H$ can traverse infinitely many circles, producing a homotopy group $\pi_1(H)$ that is uncountable and not free. The wedge $\bigvee S^1$ has the free group on countably many generators as its $\pi_1$. So the two spaces are not homotopy equivalent.

\textbf{Shape-theoretic view.} Approximate $H$ by the finite sub-wedge $H_N=\bigcup_{n=1}^{N}C_n$. Each $H_N$ is homotopy equivalent to $\bigvee_{n=1}^{N}S^1$, and the inclusions $H_N\hookrightarrow H_{N+1}$ are compatible. Shape theory studies this inverse system of finite wedges and records the transition maps. The wedge $\bigvee S^1$ has the same inverse system, so the two spaces have the same shape even though they are not homotopy equivalent. Shape detects the large-scale pattern---a shrinking sequence of loops---while ordinary homotopy sees unmanageable microscopic entanglement \cite{shape_ref1,shape_ref2}.

\paragraph{The Warsaw circle}
The Warsaw circle is formed by closing a topologist's sine curve with an arc. Its ordinary homotopy groups can look trivial, as for a point, yet its global arrangement is not point-like. Shape theory detects the large-scale organization that ordinary homotopy groups miss \cite{shape_ref1}.

\paragraph{Uses and generalizations}
Shape methods are useful for inverse limits, attractors in dynamical systems, and compact spaces arising as limits of finite approximations. Strong shape theory refines the invariants when ordinary shape loses information. Generalizations have been developed for noncommutative geometry, including shape theory for operator algebras.

\paragraph{A small worked example}
Imagine a coastline measured with rulers of length $1$ kilometre, then
$100$ metres, then $10$ metres. At each scale we replace the measured
coastline by a simpler graph and record which pieces connect. A single graph
may lose a narrow inlet, but the whole sequence shows whether that inlet
persists as the scale becomes finer. Shape theory studies this compatible
sequence rather than trusting one drawing.

The same idea applies to a dynamical attractor. A computer simulation gives
finite samples, not the exact limiting set. If graphs built from increasingly
fine samples preserve the same large-scale loops and components, those
features are candidates for shape information. This is why shape theory is
useful for spaces that are too irregular for a smooth coordinate description
\cite{shape_ref1,shape_ref2}.

\section*{4. Important theorems and results}
\begin{enumerate}[leftmargin=*,itemsep=0.35em]
\item \textbf{Approximation principle.} A compact space can be represented by a directed system of polyhedral approximations.

\item \textbf{Homotopy agreement.} For compact spaces dominated by finite polyhedra, shape and homotopy types coincide.

\item \textbf{Shape invariance.} Shape ignores changes that vanish under arbitrarily fine approximations.

\item \textbf{Warsaw-circle phenomenon.} A space can have trivial ordinary homotopy groups while possessing nontrivial global shape.

Shape theory is designed for spaces that are too irregular for a single smooth or polyhedral description. Approximation produces manageable models, homotopy agreement explains when those models recover ordinary topology, and shape invariance identifies changes that disappear at finer scales. The Warsaw-circle example gives the reason for the subject: ordinary homotopy groups can miss global information in wild compact spaces.
\end{enumerate}
\noindent\emph{Source for these results:} \cite{shape_ref1}.

\theoryapplications
\theoryconnections{shape_theory}{%
\item Topology and homotopy compare spaces by continuous maps, but a very wild space may have no useful finite description.
\item Shape theory replaces it by an inverse system of simpler spaces, records how the approximations refine one another, and declares two spaces equivalent when their entire systems agree up to compatible maps.
\item Category theory is needed because the compatibility of all approximations is part of the object, not an afterthought \cite{shape_ref1,shape_ref2}.
}{%
\item Shape theory helps continuum theory and geometric topology distinguish spaces whose local behavior is too irregular for ordinary homotopy groups.
\item In dynamics, a complicated invariant set can be approximated by a sequence of finite graphs or polyhedra, allowing global loops and connectedness to be studied through the approximations.
\item The method is useful because it preserves stable large-scale information while deliberately ignoring microscopic detail \cite{shape_ref1,shape_ref2}.
\item \textbf{Image analysis and digital topology.} A digitized image is a finite grid of pixels, hence a finite polyhedron, but the continuous scene it samples may contain thin filaments, fractal boundaries, or wild crossings. Shape-theoretic approximations built from increasingly fine pixel grids detect whether a dark region truly has a hole (a nontrivial loop) or merely looks hollow at one resolution. For example, a medical image of lung tissue can be approximated by voxel complexes at successively finer CT scan resolutions. A shape invariant that persists across resolutions---such as a nontrivial first Čech cohomology class---signals a genuine airway passage, while an artifact at one resolution will vanish in the inverse system. Digital topology software implements this by computing pro-homology groups over a filtered complex, directly applying the shape-theoretic principle that only stable information is trustworthy \cite{shape_ref1,shape_ref2}.
}

\section*{Glossary}
\begin{description}[leftmargin=*,style=nextline,itemsep=0.3em]
\item[\textbf{Shape.}] An equivalence class of compact spaces under compatible inverse systems of polyhedral approximations, capturing global topology that ordinary homotopy groups may miss.
\item[\textbf{Inverse system.}] A directed family of simpler spaces linked by maps that record how each approximation refines the previous one; shape is the stable information shared by the whole system.
\item[\textbf{Polyhedral approximation.}] A simpler finite-dimensional space that approximates a wild compact space at a given scale; shape theory studies the entire compatible sequence of such approximations.
\item[\textbf{Wild space.}] A topological space whose local behavior is too irregular for ordinary homotopy theory to handle with a single finite description, such as the Hawaiian earring.
\item[\textbf{Hawaiian earring.}] The union of circles of radii $1/n$ all meeting at one point; it has an uncountable fundamental group yet the same shape as an infinite wedge of circles.
\item[\textbf{Warsaw circle.}] A compact set formed by closing a topologist's sine curve with an arc; it has trivial ordinary homotopy groups but carries nontrivial global shape information.
\item[\textbf{Shape invariance.}] The property that shape ignores microscopic differences that vanish under arbitrarily fine polyhedral approximations, retaining only stable large-scale features.
\item[\textbf{Fundamental group.}] The group $\pi_1(X,x_0)$ of homotopy classes of loops at a basepoint; the primary invariant shape theory replaces for wild spaces.
\item[\textbf{Homotopy equivalence.}] A relation weaker than homeomorphism: each space can be continuously deformed into the other.
\item[\textbf{Inverse limit.}] The categorical limit of an inverse system; constructs complicated spaces from simpler stages.
\item[\textbf{\v{C}ech cohomology.}] A cohomology theory computed from open covers; the natural cohomology for shape-theoretic purposes.
\end{description}
\section*{7. Sample Questions}
\emph{Exercise reference:} \cite{shape_ref1}. The cited edition is the source for the exercise set; consult its Exercises section for the official statement and numbering.
\begin{enumerate}[leftmargin=*,itemsep=0.9em]
\item \textbf{Exercise 1.} Let $X=\bigcap_{n=1}^{\infty}X_n$ be the Hawaiian earring, where $X_n=\bigcup_{k=1}^{n}C_k$ and $C_k$ is a circle of radius $1/k$ centered at $(1/k,0)$, all meeting at the origin. Describe the inverse system $\{X_n,\,f_{n,n+1}\}$ where $f_{n,n+1}\colon X_{n+1}\hookrightarrow X_n$ are the inclusions. What are the transition maps? \emph{Solution.} The inverse system is $(X_1\hookleftarrow X_2\hookleftarrow X_3\hookleftarrow\cdots)$ where each $f_{n,n+1}$ is the inclusion $X_{n+1}\supset X_n$, i.e., the identity on the subset $X_n$. Every $X_n$ is homotopy equivalent to $\bigvee_{k=1}^{n}S^1$, so the inverse system is equivalent to $(S^1\hookleftarrow S^1\vee S^1\hookleftarrow S^1\vee S^1\vee S^1\hookleftarrow\cdots)$ with inclusion maps joining a new circle summand at each stage.

\item \textbf{Exercise 2.} The Warsaw circle $W$ is formed by taking the topologist's sine curve $\{(x,\sin(1/x)):0<x\leq 1/\pi\}\cup\{(0,y):-1\leq y\leq 1\}$ and connecting its two endpoints by a circular arc. Compute $\pi_1(W)$ and explain why it is trivial. \emph{Solution.} Every continuous map $S^1\to W$ can be deformed into the vertical arc segment because the sine-curve portion accumulates on the segment $\{0\}\times[-1,1]$, so any loop must eventually be homotoped to the trivial loop. Hence $\pi_1(W)=0$. Despite this, $W$ is not contractible and carries nontrivial shape.

\item \textbf{Exercise 3.} For the compact metric space $X=[0,1]$ (a contractible space), describe its shape and explain why shape and homotopy agree here. \emph{Solution.} $[0,1]$ is a finite polyhedron. Any open cover of $[0,1]$ can be refined to a good cover whose nerve is homotopy equivalent to $[0,1]$, i.e., contractible. The shape of $[0,1]$ is the pro-homotopy type of a single point. This agrees with ordinary homotopy because $[0,1]$ is homotopically dominated by a point.

\item \textbf{Exercise 4.} Let $Y$ be the dyadic solenoid, the inverse limit of the system $S^1\xleftarrow{z\mapsto z^2}S^1\xleftarrow{z\mapsto z^2}S^1\cdots$. Compute the first Cech cohomology group $\check{H}^1(Y;\mathbb{Z})$. \emph{Solution.} Cech cohomology converts the inverse system of spaces into a direct system of cohomology groups. $\check{H}^1(S^1;\mathbb{Z})\cong\mathbb{Z}$ at each stage, and the map $z\mapsto z^2$ induces multiplication by $2$ on $H^1$. So $\check{H}^1(Y;\mathbb{Z})=\varinjlim(\mathbb{Z}\xrightarrow{2}\mathbb{Z}\xrightarrow{2}\mathbb{Z}\cdots)=\mathbb{Z}[1/2]$, the dyadic rationals.

\item \textbf{Exercise 5.} Let $Z$ be a compact Hausdorff space that is the one-point compactification of the disjoint union of countably many open intervals $(n,n+1)$ for $n\in\mathbb{Z}$. Describe the shape of $Z$ and compute $\check{H}^1(Z;\mathbb{Z})$. \emph{Solution.} $Z$ is homeomorphic to $S^1$---the open intervals wrap around and meet at a point, forming a circle. Hence the shape of $Z$ is that of $S^1$, and $\check{H}^1(Z;\mathbb{Z})\cong\mathbb{Z}$.

\end{enumerate}
\section*{8. References}
\begin{thebibliography}{9}
\bibitem{shape_ref1} S. Mardesic, \emph{Shape Theory}, World Scientific, 1989.
\bibitem{shape_ref2} J.-M. Cordier and T. Porter, \emph{Shape Theory: Categorical Methods of Approximation}, Ellis Horwood, 1989.
\end{thebibliography}


\theory{Sheaf theory}{How can locally defined information be organized so that compatible local pieces glue into a global object?}{A sheaf assigns data to open regions of a space, restricts data to smaller regions, and requires compatible local data to glue uniquely. Sheaves may contain sets, groups, rings, modules, or vector spaces.}{Topology, algebraic and differential geometry, algebraic geometry, differential equations, cohomology, logic, and number theory.}{Restriction maps and gluing turn compatible local data into global sections and measure the failure when gluing is impossible.}

\paragraph{The problem and history}
Continuous functions, differentiable functions, vector fields, and solutions of differential equations are naturally defined on open subsets. A restriction map sends data on a large region to data on a smaller one. A presheaf is such an assignment; a sheaf adds the condition that compatible local data determine exactly one global datum \cite{sheaf_ref1}.

The theory grew from cohomology and analytic continuation. Jean Leray introduced sheaf methods in work on fixed-point theorems and partial differential equations. Henri Cartan and Alexander Grothendieck developed sheaf cohomology and its algebraic applications. Grothendieck later generalized open covers to Grothendieck topologies \cite{sheaf_ref2}.

\end{historylist}
\section*{3. Theory}
\subsection*{Core theory}
\paragraph{Definition through an example}
Let $X$ be a topological space. For each open set $U$, let $\mathcal F(U)$ be the set of data on $U$. If $V\subseteq U$, restriction gives a map $\mathcal F(U)\to\mathcal F(V)$. The sheaf condition says that if $U$ is covered by open sets $U_i$ and sections $s_i\in\mathcal F(U_i)$ agree on overlaps $U_i\cap U_j$, then there is a unique section $s\in\mathcal F(U)$ restricting to every $s_i$.

The sheaf of continuous real-valued functions satisfies this rule: compatible continuous pieces define one continuous function. The sheaf of smooth functions, holomorphic functions, vector fields, and sections of a vector bundle provide richer examples. The constant presheaf may fail to be a sheaf because locally constant choices can behave differently on disconnected regions \cite{sheaf_ref1}.

\paragraph{Worked example: the sheaf of continuous functions on $\mathbb R$}
Let $X=\mathbb R$ with its usual topology, and for each open $U\subseteq\mathbb R$ let $\mathcal F(U)=C(U,\mathbb R)$ be the set of continuous real-valued functions on $U$. If $V\subseteq U$, the restriction map $\rho_{UV}\colon\mathcal F(U)\to\mathcal F(V)$ sends a function $f$ on $U$ to its restriction $f\big|_V$. Now take the open cover $U_1=(-\infty,1)$ and $U_2=(0,\infty)$ of $\mathbb R$. Suppose $f_1\in\mathcal F(U_1)$ and $f_2\in\mathcal F(U_2)$ agree on the overlap $U_1\cap U_2=(0,1)$, meaning $f_1(x)=f_2(x)$ for every $x\in(0,1)$. Then there is a unique continuous function $f\colon\mathbb R\to\mathbb R$ whose restriction to $U_1$ is $f_1$ and whose restriction to $U_2$ is $f_2$: define $f(x)=f_1(x)$ for $x\le1$ and $f(x)=f_2(x)$ for $x>0$. Agreement on the overlap guarantees continuity at every point. This is precisely the sheaf gluing axiom in action. Failure of agreement would obstruct gluing: for instance, $f_1(x)=0$ on $(-\infty,1)$ and $f_2(x)=1$ on $(0,\infty)$ cannot glue because they conflict on $(0,1)$.

\paragraph{Morphisms and functors}
A morphism of sheaves is a collection of maps commuting with restriction. Sheaves of a fixed type on a space form a category. A continuous map $f:X\to Y$ gives a direct image functor $f_*$, which records data on inverse images, and an inverse image functor $f^{-1}$, which transports local data in the other direction. These functors and their derived versions are central to geometry \cite{sheaf_ref3}.

\paragraph{Cohomology and geometry}
Sheaf cohomology measures the failure of local data to glue globally. Degree zero records global sections. Higher cohomology detects obstructions such as a locally defined function or solution that cannot be chosen consistently over the whole space. It includes familiar topological cohomology and connects local analytic information with global invariants.

Manifolds and schemes can be described using sheaves of rings. Divisors, vector bundles, and differential operators are naturally sheaf-theoretic. D-modules use sheaves of differential operators to study linear differential equations. In algebraic geometry, coherent sheaves package algebraic families and support duality and Riemann--Roch theory \cite{sheaf_ref4}.

\paragraph{Sites, topoi, and arithmetic}
The Zariski topology on an algebraic variety has very few open sets, so ordinary constant-sheaf cohomology cannot recover the desired information over finite fields. Grothendieck topologies replace open covers by abstract covering families. A category with such coverings is a site, and sheaves on it form a Grothendieck topos.

Etale and l-adic sheaf cohomology supplied the cohomological tools used in the proof of the Weil conjectures. Topoi also connect sheaves with mathematical logic, and sheaf constructions on categories provide flexible foundations for arithmetic geometry \cite{sheaf_ref2}.

\section*{4. Important theorems and results}
\begin{enumerate}[leftmargin=*,itemsep=0.35em]
\item \textbf{Sheaf gluing axiom.} Compatible local sections on a cover glue uniquely to a global section.

\item \textbf{Sheafification.} Every presheaf admits a universal sheaf receiving a map from it.

\item \textbf{De Rham theorem.} On a smooth manifold, the cohomology of differential forms agrees with singular cohomology.

\item \textbf{Serre's coherent sheaf theory.} Coherent sheaves connect algebraic functions, vector bundles, and cohomology on projective varieties.

\item \textbf{Grothendieck duality.} Under suitable hypotheses, direct and inverse image operations are related by dualizing functors.

\item \textbf{Etale cohomology.} Sheaf cohomology for the etale topology provides arithmetic invariants unavailable from the Zariski topology alone.
\end{enumerate}
\noindent\emph{Source for these results:} \cite{sheaf_ref1}.

\theoryapplications
\theoryconnections{sheaf_theory}{%
\item Topology provides open sets and overlaps, while a sheaf assigns data to each open set together with restriction maps to smaller sets.
\item The gluing rule says that compatible local sections come from one and only one global section.
\item Algebra decides what the sections are, and homological algebra measures the failure of gluing when the compatibility equations have no global solution \cite{sheaf_ref1,sheaf_ref2}.
}{%
\item Sheaf theory helps geometry describe functions, vector fields, and solutions of differential equations that are defined locally but may fail to fit globally.
\item On a scheme it records regular functions and modules; in arithmetic it organizes information over different localizations; in logic it models varying truth values.
\item Cohomology turns the failure of gluing into invariants, so a local calculation can reveal a global obstruction \cite{sheaf_ref1,sheaf_ref3}.
}
\noindent\textbf{Data fusion and sensor networks.} A sensor network deploys many sensors over a region, each producing local measurements (temperature, air quality, acoustic signals) on its own coverage area. A sheaf on the network's topology assigns to each open set of sensors the data they collectively observe, with restriction maps encoding how data from a larger group restricts to a subgroup. The sheaf condition requires that overlapping sensor groups produce consistent readings. When readings conflict on overlaps, sheaf cohomology in degree~1 measures the inconsistency. This is used in practice: Baryshnikov and Ghrist introduced sheaf-based methods for distributed sensing, and real-time data-fusion algorithms use sheaf Laplacians to detect and localize sensor faults. A sensor reporting erroneous data creates a nontrivial cocycle; optimizing the sheaf Laplacian's smallest nonzero eigenvalue pinpoints which sensor is the source of the inconsistency.

\section*{Glossary}
\begin{description}[leftmargin=*,style=nextline,itemsep=0.3em]
\item[\textbf{Sheaf.}] A structure that assigns data (sets, groups, rings, or modules) to every open set of a topological space, together with restriction maps to smaller open sets, satisfying the condition that compatible local pieces glue uniquely to a global piece.
\item[\textbf{Presheaf.}] An assignment of data and restriction maps to open sets that may fail the gluing condition; a sheaf is a presheaf that satisfies it.
\item[\textbf{Section.}] A piece of data assigned to a particular open set by a sheaf; a global section is a section defined on the whole space.
\item[\textbf{Restriction map.}] The operation that takes data on a larger open set and produces compatible data on a smaller open subset, ensuring local information can be compared across scales.
\item[\textbf{Sheaf condition.}] The requirement that whenever overlapping open sets carry compatible local sections, there exists exactly one global section restricting to each of them.
\item[\textbf{Sheaf cohomology.}] A sequence of groups measuring the obstruction to gluing local data globally; degree zero gives the global sections, and higher degrees detect precise obstructions.
\item[\textbf{Grothendieck topology.}] A generalization of open covers defined on a category, allowing abstract covering families that supply more local information than the Zariski topology alone.
\item[\textbf{Topos.}] A category of sheaves on a site that behaves like a generalized space, supporting both geometric and logical reasoning.
\item[\textbf{Sheafification.}] The universal process turning a presheaf into a sheaf, adding gluing data the presheaf lacked.
\item[\textbf{Direct image functor.}] For a map $f\colon X\to Y$, the functor $f_*$ pushing sections forward along $f$.
\item[\textbf{Inverse image functor.}] For a map $f\colon X\to Y$, the functor $f^{-1}$ pulling a sheaf on $Y$ back to $X$.
\item[\textbf{\'Etale topology.}] A Grothendieck topology on a scheme defined by \'etale morphisms, providing more open covers than Zariski.
\item[\textbf{Coherent sheaf.}] A sheaf of modules locally finitely generated and finitely presented; fundamental in algebraic geometry.
\item[\textbf{De Rham theorem.}] On a smooth manifold, de Rham cohomology agrees with singular cohomology with real coefficients.
\end{description}
\section*{7. Sample Questions}
\emph{Exercise reference:} \cite{sheaf_ref1}. The cited edition is the source for the exercise set; consult its Exercises section for the official statement and numbering.
\begin{enumerate}[leftmargin=*,itemsep=0.9em]
\item \textbf{Exercise 1. Why do continuous functions form a sheaf?} Compatible continuous functions on overlapping open sets agree where they meet, so defining the combined function piece by piece is well defined and continuous.

\item \textbf{Exercise 2. What does the sheaf condition add to a presheaf?} It guarantees existence and uniqueness of a global object from compatible local objects; restriction alone does not guarantee this.

\item \textbf{Exercise 3. What does degree-zero sheaf cohomology represent?} It is the group or set of global sections, because no gluing obstruction remains at degree zero.

\item \textbf{Exercise 4. Why are schemes described by sheaves of rings?} Polynomial coordinates and regular functions vary from one open region to another. A sheaf records those local rings and their compatibility.

\item \textbf{Exercise 5. Why introduce the etale topology?} The Zariski topology has too few opens over finite fields. Etale covers provide enough local information to build useful cohomology.

\end{enumerate}
\section*{8. References}
\begin{thebibliography}{9}
\bibitem{sheaf_ref1} G. E. Bredon, \emph{Sheaf Theory}, 2nd ed., Springer, 1997.
\bibitem{sheaf_ref2} R. Godement, \emph{Topologie algebrique et theorie des faisceaux}, Hermann, 1958.
\bibitem{sheaf_ref3} I. Moerdijk and S. Mac Lane, \emph{Sheaves in Geometry and Logic}, Springer, 1994.
\bibitem{sheaf_ref4} R. Hartshorne, \emph{Algebraic Geometry}, Springer, 1977.
\end{thebibliography}


\theory{Sieve theory}{How can we estimate the size of a set of integers after removing numbers divisible by selected primes?}{Sieve theory develops weighted versions of the sieve of Eratosthenes. It estimates sifted sets, often primes or almost-primes, without testing every integer.}{Prime distribution, twin-prime questions, almost-primes, bounded gaps, and analytic number theory.}{Inclusion-exclusion over prime divisors produces upper and lower bounds for arithmetic sets that are difficult to count directly.}

\paragraph{The idea and history}
To find primes up to $X$, cross out multiples of $2$, then $3$, then $5$, and so on. This is the sieve of Eratosthenes. Viggo Brun introduced the word sieve in 1915 and adapted the idea to questions about prime pairs \cite{sieve_ref1}.

Let $\mathcal A$ be a sequence of nonnegative weights and $\mathcal P$ a set of primes. If $P(z)$ is the product of primes in $\mathcal P$ below $z$, a basic sifted sum is
\[
S(\mathcal A,\mathcal P,z)=\sum_{\substack{n\\(n,P(z))=1}}a_n.
\]
It counts or weights members of $\mathcal A$ not divisible by selected small primes. Weights are useful because the exact prime indicator is hard to handle, while a smooth approximation can be estimated \cite{sieve_ref2}.

\end{historylist}
\section*{3. Theory}
\subsection*{Core theory}
\paragraph{Main sieves and the parity problem}
The Brun sieve, Selberg sieve, Turan sieve, large sieve, larger sieve, and Goldston--Pintz--Yildirim sieve use different weights and information about residue classes. A sieve often replaces primes by almost-primes, numbers with only a few prime factors, because almost-primes are easier to count.

The parity problem says that ordinary sieve information is poor at distinguishing integers with an odd number of prime factors from integers with an even number. This is why sieve methods alone have difficulty proving the twin prime conjecture, even when they can prove that numbers with two prime factors occur infinitely often \cite{sieve_ref3}.

\paragraph{Major results}
Brun's theorem proves that the sum of reciprocals of twin-prime pairs converges. Chen's theorem proves that infinitely many primes $p$ have $p+2$ either prime or semiprime. Chen also proved that every sufficiently large even number is a prime plus a prime or semiprime.

The fundamental lemma estimates the number left after a controlled number of sifting steps. It is usually too weak to isolate all primes, but it is powerful for almost-prime results. The Friedlander--Iwaniec theorem proves infinitely many primes of the form $a^2+b^4$. Zhang's theorem proved infinitely many prime pairs within a bounded distance, and Maynard--Tao methods gave stronger bounded-gap results \cite{sieve_ref3,sieve_ref4}.

\paragraph{Applications and limits}
Sieve estimates combine with the circle method, distribution estimates, and analytic number theory. They count primes in polynomial sequences, control exceptional sets, and provide upper bounds for prime patterns. Factorization algorithms called the quadratic sieve and number-field sieve use related ideas but are distinct computational methods.

\section*{4. Important theorems and results}
\begin{enumerate}[leftmargin=*,itemsep=0.35em]
\item \textbf{Fundamental lemma.} Under suitable distribution hypotheses, a sifted set can be estimated after a controlled number of prime exclusions.

\item \textbf{Brun's theorem.} The reciprocal sum over twin primes converges.

\item \textbf{Chen's theorem.} Infinitely many primes have a companion two units away that is prime or semiprime.

\item \textbf{Friedlander--Iwaniec theorem.} Infinitely many primes are represented by $a^2+b^4$.

\item \textbf{Zhang's bounded-gap theorem.} Infinitely many pairs of primes occur within some fixed finite distance.

\item \textbf{Maynard--Tao theorem.} Bounded intervals contain arbitrarily many primes for infinitely many choices of the interval, with a bound depending on the desired number.
\end{enumerate}
\noindent\emph{Source for these results:} \cite{sieve_ref1}.

\theoryapplications
\theoryconnections{sieve_theory}{%
\item Sieve theory starts with inclusion--exclusion: to count numbers with no forbidden prime divisor, subtract multiples of each prime, add back multiples of pairs, and so on.
\item Congruences describe which residue classes survive, divisor sums organize the repeated inclusion--exclusion, and analytic estimates control the error created when the calculation is truncated.
\item The art is to keep enough local information to get a useful global bound without needing to know every prime factor exactly \cite{sieve_ref1,sieve_ref2}.
}{%
\item A sieve can upper-bound how many integers represent a sum of primes or have a prescribed pattern of prime factors, even when proving primality itself is out of reach.
\item This supplies almost-prime results in additive number theory and bounds for prime gaps.
\item Its output is often an inequality rather than an exact count, because discarded correlations between divisibility conditions are precisely what make a simple sieve tractable \cite{sieve_ref1,sieve_ref3}.
}

\section*{Glossary}
\begin{description}[leftmargin=*,style=nextline,itemsep=0.3em]
\item[\textbf{Sieved set.}] The subset of integers or weighted sequence that remains after removing all entries divisible by primes below a chosen bound $z$.
\item[\textbf{Parity problem.}] The fundamental limitation that standard sieve methods struggle to distinguish integers with an odd number of prime factors from those with an even number.
\item[\textbf{Almost-prime.}] An integer with at most a bounded number of prime factors; sieves can often count almost-primes even when isolating true primes is out of reach.
\item[\textbf{Inclusion--exclusion.}] A counting method that alternately adds and subtracts the contributions of overlapping divisibility conditions, forming the algebraic backbone of sieve arguments.
\item[\textbf{Fundamental lemma.}] A technical sieve estimate giving the size of a sifted set after a controlled number of sifting steps, powerful for almost-prime results but usually too weak to isolate all primes.
\item[\textbf{Selberg sieve.}] A weighted sieve method that uses the square of a divisor sum to produce sharp upper bounds for sifted sets.
\item[\textbf{Brun's theorem.}] The result that the sum of reciprocals of twin-prime pairs converges, providing early evidence that twin primes are sparse.
\item[\textbf{Semiprime.}] A product of exactly two primes; Chen's theorem shows that every sufficiently large even number is a prime plus a semiprime.
\item[\textbf{Sieve of Eratosthenes.}] The classical algorithm finding primes by iteratively crossing out multiples; the historical origin of all sieve methods.
\item[\textbf{Large sieve.}] An analytic inequality bounding the concentration of a sequence in residue classes.
\item[\textbf{Chen's theorem.}] Every sufficiently large even integer is the sum of a prime and a semiprime.
\item[\textbf{Zhang's bounded-gap theorem.}] Infinitely many pairs of primes differ by at most a bounded constant.
\item[\textbf{Maynard--Tao theorem.}] Bounded intervals can contain arbitrarily many primes for infinitely many starting points.
\item[\textbf{Twin prime conjecture.}] There are infinitely many pairs of primes differing by exactly $2$.
\end{description}
\section*{7. Sample Questions}
\emph{Exercise reference:} \cite{sieve_ref1}. The cited edition is the source for the exercise set; consult its Exercises section for the official statement and numbering.
\begin{enumerate}[leftmargin=*,itemsep=0.9em]
\item \textbf{Exercise 1.} Apply the sieve of Eratosthenes to find all primes up to $30$. After crossing out multiples of $2$, $3$, and $5$, list the surviving integers from $2$ to $30$. \emph{Solution.} Start with $\{2,3,4,\ldots,30\}$. Remove multiples of $2$: $\{2,3,5,7,9,11,13,15,17,19,21,23,25,27,29\}$. Remove multiples of $3$: $\{2,3,5,7,11,13,17,19,23,25,29\}$. Remove multiples of $5$: $\{2,3,5,7,11,13,17,19,23,29\}$. The primes up to $30$ are $2,3,5,7,11,13,17,19,23,29$.

\item \textbf{Exercise 2.} Using inclusion--exclusion, count the integers from $1$ to $60$ that are \emph{not} divisible by $2$, $3$, or $5$. \emph{Solution.} Let $A_2,A_3,A_5$ be the sets of multiples of $2,3,5$ in $\{1,\ldots,60\}$. $|A_2|=30$, $|A_3|=20$, $|A_5|=12$. $|A_2\cap A_3|=10$, $|A_2\cap A_5|=6$, $|A_3\cap A_5|=4$. $|A_2\cap A_3\cap A_5|=2$. By inclusion--exclusion: $|A_2\cup A_3\cup A_5|=30+20+12-10-6-4+2=44$. The count of survivors is $60-44=16$. (The survivors are: $1,7,11,13,17,19,23,29,31,37,41,43,47,49,53,59$.)

\item \textbf{Exercise 3.} Use the Legendre sieve to count the integers in $\{1,\ldots,100\}$ not divisible by any prime $p\leq 7$. Compute the inclusion--exclusion sum over the primes $\{2,3,5,7\}$. \emph{Solution.} Let $N=100$. $\lfloor N/2\rfloor=50$, $\lfloor N/3\rfloor=33$, $\lfloor N/5\rfloor=20$, $\lfloor N/7\rfloor=14$. Pairwise: $\lfloor N/6\rfloor=16$, $\lfloor N/10\rfloor=10$, $\lfloor N/14\rfloor=7$, $\lfloor N/15\rfloor=6$, $\lfloor N/21\rfloor=4$, $\lfloor N/35\rfloor=2$. Triple: $\lfloor N/30\rfloor=3$, $\lfloor N/42\rfloor=2$, $\lfloor N/70\rfloor=1$, $\lfloor N/105\rfloor=0$. Quad: $\lfloor N/210\rfloor=0$. The sifted count is $100-50-33-20-14+16+10+7+6+4+2-3-2-1-0+0=100-117+45-6+0=22$. These $22$ integers are $\{1\}$ plus the primes $11,13,17,19,23,29,31,37,41,43,47,53,59,61,67,71,73,79,83,89,97$, totaling $22$.

\item \textbf{Exercise 4.} Brun's theorem states that $\sum_{p,\,p+2\text{ prime}}(1/p+1/(p+2))$ converges. Compute the first five terms of this sum corresponding to the twin prime pairs $(3,5)$, $(5,7)$, $(11,13)$, $(17,19)$, and $(29,31)$. \emph{Solution.} The terms are: $(1/3+1/5)+(1/5+1/7)+(1/11+1/13)+(1/17+1/19)+(1/29+1/31)$. Computing: $(5/15+3/15)=8/15\approx 0.5333$. $(1/5+1/7)=12/35\approx 0.3429$. $(1/11+1/13)=24/143\approx 0.1678$. $(1/17+1/19)=36/323\approx 0.1115$. $(1/29+1/31)=60/899\approx 0.0667$. The partial sum is approximately $0.5333+0.3429+0.1678+0.1115+0.0667=1.2222$.

\item \textbf{Exercise 5.} In a combinatorial sieve, let $\mathcal{A}=\{n\in\mathbb{N}:n\leq 100,\,n\text{ is odd}\}$ and $\mathcal{P}=\{3,5,7\}$. Apply the sifting function to compute $S(\mathcal{A},\mathcal{P},10)=\#\{n\in\mathcal{A}:n\leq 100,\,3\nmid n,\,5\nmid n,\,7\nmid n\}$. \emph{Solution.} $\mathcal{A}$ has $50$ elements. Remove odd multiples of $3$ up to $100$: these are $3,9,15,\ldots,99$, which are $17$ numbers. Remove odd multiples of $5$ up to $100$: $5,15,25,\ldots,95$, which are $10$ numbers. Remove odd multiples of $7$ up to $100$: $7,21,35,49,63,77,91$, which are $7$ numbers. Add back odd multiples of $15$ (lcm$(3,5)$): $15,45,75$, which are $3$. Add back odd multiples of $21$: $21,63$, which are $2$. Add back odd multiples of $35$: $35$, which is $1$. Subtract odd multiples of $105$: none up to $100$. By inclusion--exclusion: $50-17-10-7+3+2+1=22$. The sifted count is $22$.

\end{enumerate}
\section*{8. References}
\begin{thebibliography}{9}
\bibitem{sieve_ref1} A. C. Cojocaru and M. R. Murty, \emph{An Introduction to Sieve Methods and Their Applications}, Cambridge, 2005.
\bibitem{sieve_ref2} H. Halberstam and H.-E. Richert, \emph{Sieve Methods}, Academic Press, 1974.
\bibitem{sieve_ref3} J. Friedlander and H. Iwaniec, \emph{Opera de Cribro}, AMS, 2010.
\bibitem{sieve_ref4} J. Maynard, \emph{Small Gaps Between Primes}, Oxford, 2023.
\end{thebibliography}


\theory{Singularity theory}{How do spaces, maps, and physical patterns change when a smooth or regular situation degenerates?}{Singularity theory studies points where manifold structure fails, critical points of maps, and sudden changes under parameter variation. It classifies stable local models and studies how singularities arise by projection, symmetry, or degeneration.}{Algebraic geometry, differential geometry, wave fronts, optics, mechanics, catastrophe theory, bifurcation theory, and intersection cohomology.}{Jacobian rank, local normal forms, and equivalence under coordinate changes classify where smooth behavior breaks down.}

\paragraph{The idea and history}
A string is a one-dimensional manifold until it crosses itself. A crossing is a double point; a touching point can disappear after a tiny movement and is unstable. Singularity theory asks which irregularities persist under small changes and which are artifacts of a special drawing \cite{singularity_ref1}.

Vladimir Arnold described the subject as the study of how objects depend on parameters when a small change causes a sudden qualitative change. Whitney studied stable singularities and critical points; Rene Thom developed catastrophe theory; Arnold connected the subject with symplectic geometry and mechanics \cite{singularity_ref2}.

\end{historylist}
\section*{3. Theory}
\subsection*{Core theory}
\paragraph{Algebraic and geometric examples}
The curve $y^2=x^2+x^3$ has a double point at the origin, while $y^2=x^3$ has a cusp. A singular point of a variety is one where the local equations do not define a smooth manifold. The Jacobian matrix detects this: the rank is too small at a singular point.

Singularities also arise by projecting a smooth object into lower dimensions. Caustics, such as bright patterns at the bottom of a swimming pool, occur where many light rays focus. Symmetry can create orbifolds with corners or cone-like points. A smooth manifold can degenerate into a singular space as parameters change \cite{singularity_ref3}.

\paragraph{Stability, germs, and catastrophes}
To classify a local singularity, one studies a germ, the behavior near a point, up to changes of coordinates. Taylor expansions and jets record derivatives up to a chosen order. Group actions identify equivalent local forms.

Catastrophe theory studies stable critical-point changes in families of smooth functions. Fold, cusp, swallowtail, and butterfly are standard elementary models. The models do not claim that every natural event is one of these shapes; they describe local mechanisms by which equilibria appear, merge, or disappear \cite{singularity_ref1}.

\paragraph{Resolution and duality}
Hironaka's resolution-of-singularities theorem in characteristic zero says that singularities of algebraic varieties can be replaced by a smooth variety through a controlled sequence of transformations. The operation is not simply erasing a bad point; it records the exceptional geometry introduced during the resolution.

Singular spaces may fail Poincare duality. Intersection cohomology restores a form of duality by using the stratification of the singular space. Perverse sheaves arise from this framework and connect singularity theory with homological algebra. Monodromy studies how solutions or covering spaces behave around singularities of differential equations \cite{singularity_ref4}.

\section*{4. Important theorems and results}
\begin{enumerate}[leftmargin=*,itemsep=0.35em]
\item \textbf{Jacobian criterion.} Under standard hypotheses, a point of an algebraic variety is smooth exactly when the Jacobian has the expected rank.

\item \textbf{Whitney stability theory.} Stable singularities are those whose qualitative form persists under small perturbations.

\item \textbf{Morse lemma.} Near a nondegenerate critical point, a smooth function is equivalent to a quadratic form.

\item \textbf{Hironaka resolution theorem.} Algebraic varieties over characteristic zero admit resolutions of singularities.

\item \textbf{Intersection-cohomology duality.} Intersection cohomology restores a Poincare-type duality for broad classes of singular spaces.

\item \textbf{Classification of elementary catastrophes.} Low-dimensional stable critical changes reduce locally to a short list of normal forms.
\end{enumerate}
\noindent\emph{Source for these results:} \cite{singularity_ref1}.

\theoryapplications
\theoryconnections{singularity_theory}{%
\item Calculus identifies where a map loses rank or a surface ceases to be smooth.
\item Differential and algebraic geometry classify the local equations, topology tests whether the local type survives small changes, and Lie group actions remove coordinate choices that do not change the singularity.
\item Homological tools then measure how cycles, intersections, or vanishing pieces appear when the singular point is resolved or deformed \cite{singularity_ref1,singularity_ref2}.
}{%
\item Singularity theory supplies normal forms for catastrophe and bifurcation models, so a complicated equation can be replaced locally by a short list of stable patterns.
\item In optics, the singular set of a ray map produces caustics and wave fronts; in algebraic geometry, resolutions and intersection cohomology recover information hidden by a singular point.
\item The same classification therefore links visible folds and cusps with algebraic invariants \cite{singularity_ref1,singularity_ref3}.
}
\section*{Glossary}
\begin{description}[leftmargin=*,style=nextline,itemsep=0.3em]
\item[\textbf{Singularity.}] A point where a mathematical object (such as a curve, surface, or map) fails to be smooth or regular; for example, a cusp or crossing point.
\item[\textbf{Germ.}] The local behavior of a function or map near a specific point, considered up to changes of coordinates.
\item[\textbf{Caustic.}] A visible pattern formed by the focusing or envelope of light rays; it appears as a bright curve or surface where many rays meet, creating a singularity in the projection.
\item[\textbf{Resolution of singularities.}] A process that replaces a singular variety with a smooth one while keeping track of the exceptional sets introduced, effectively "smoothing out" the singularities.
\item[\textbf{Intersection cohomology.}] A tool that restores Poincaré duality for singular spaces by restricting chains near singular strata, providing a way to study the topology of singular objects.
\item[\textbf{Morse lemma.}] A result stating that near a nondegenerate critical point, a smooth function can be expressed as a quadratic form in appropriate local coordinates.
\item[\textbf{Catastrophe theory.}] The study of how small changes in parameters can cause sudden qualitative changes in the behavior of a system, especially in families of smooth functions.
\item[\textbf{Stratification.}] A decomposition of a singular space into smooth manifold pieces arranged so each piece sits cleanly on top of lower-dimensional ones.
\item[\textbf{Monodromy.}] The action on solutions obtained by analytic continuation around a singular point.
\item[\textbf{Normal form.}] A canonical representative of a singularity up to coordinate changes, used to classify qualitatively distinct local behaviors.
\item[\textbf{Jets.}] The Taylor expansion of a map up to a chosen finite order, recording derivatives that distinguish nearby maps.
\end{description}
\section*{7. Sample Questions}
\emph{Exercise reference:} \cite{singularity_ref1}. The cited edition is the source for the exercise set; consult its Exercises section for the official statement and numbering.
\begin{enumerate}[leftmargin=*,itemsep=0.9em]
\item \textbf{Exercise 1.} Determine whether the origin is a singular point of the curve $y^2 = x^2 + x^3$. \emph{Solution.} Let $F(x,y) = y^2 - x^2 - x^3$. The gradient is $\nabla F = (-2x - 3x^2,\; 2y)$. At the origin, $\nabla F(0,0) = (0,0)$, so the Jacobian criterion identifies the origin as a singular point. Geometrically, the curve has a node (self-crossing) at the origin.

\item \textbf{Exercise 2.} Determine whether the origin is singular on $y^2 = x^2 + x^4$ and on $y^2 = x^4 + x^6$. \emph{Solution.} For $F_1(x,y) = y^2 - x^2 - x^4$: $\nabla F_1(0,0) = (0,0)$, so the origin is singular (a node). For $F_2(x,y) = y^2 - x^4 - x^6$: $\nabla F_2(0,0) = (0,0)$ as well, so this is also singular. However, after blowing up once, $F_1$ becomes $v^2 = 1 + u^2$ (smooth), while $F_2$ becomes $v^2 = u^2 + u^4$ (still singular). Thus $F_1$ has a simpler singularity requiring one blow-up, while $F_2$ requires two.

\item \textbf{Exercise 3.} By the Morse lemma, near a nondegenerate critical point a smooth function $f\colon\mathbb{R}^2\to\mathbb{R}$ is equivalent to a quadratic form. Compute the Hessian of $f(x,y) = x^2 + 3y^2 - 2xy$ at the origin, classify the critical point, and state the normal form. \emph{Solution.} The partial derivatives are $f_x = 2x - 2y$ and $f_y = 6y - 2x$, which both vanish at the origin. The Hessian matrix is $H = \begin{pmatrix} 2 & -2 \\ -2 & 6 \end{pmatrix}$. Its determinant is $2 \cdot 6 - (-2)^2 = 8 > 0$ and $f_{xx} = 2 > 0$, so the origin is a nondegenerate minimum. By the Morse lemma, $f$ is locally equivalent to $u^2 + v^2$ after a suitable change of coordinates.

\item \textbf{Exercise 4.} Verify that the curve $y^2 = x^5$ has a cusp at the origin by computing its tangent cone and comparing it with the node on $y^2 = x^2(x+1)$. \emph{Solution.} For $y^2 = x^5$, the lowest-degree terms in $F(x,y) = y^2 - x^5$ are $y^2$, so the tangent cone is the double line $y^2 = 0$ (the $x$-axis counted twice). This is a cusp. For $y^2 = x^2(x+1) = x^3 + x^2$, the lowest-degree terms are $y^2 - x^2 = (y-x)(y+x)$, giving two distinct tangent lines $y = x$ and $y = -x$. This is a node (transverse self-crossing).

\item \textbf{Exercise 5.} Compute the Euler characteristic of the blow-up of $\mathbb{CP}^2$ at a single point. \emph{Solution.} The Euler characteristic of $\mathbb{CP}^2$ is $\chi(\mathbb{CP}^2) = 3$ (since it has a cell decomposition with cells in dimensions $0, 2, 4$). Blowing up a point replaces it with an exceptional $\mathbb{CP}^1$, contributing $\chi(\mathbb{CP}^1) = 2$ instead of the single point's contribution of $\chi(\text{pt}) = 1$. Thus $\chi(\text{blow-up}) = 3 - 1 + 2 = 4$. Equivalently, the blow-up is $\mathbb{CP}^2 \# \overline{\mathbb{CP}}^2$, and $\chi = 3 + 3 - 2 = 4$.

\end{enumerate}
\section*{8. References}
\begin{thebibliography}{9}
\bibitem{singularity_ref1} V. I. Arnold, \emph{Singularities of Differentiable Maps}, Birkhauser, 1982.
\bibitem{singularity_ref2} V. I. Arnold, \emph{Catastrophe Theory}, 2nd ed., Springer, 1986.
\bibitem{singularity_ref3} E. Brieskorn and H. Knorrer, \emph{Plane Algebraic Curves}, Birkhauser, 1981.
\bibitem{singularity_ref4} M. Goresky and R. MacPherson, \emph{Stratified Morse Theory}, Springer, 1980.
\end{thebibliography}


\theory{Soliton theory}{How can a nonlinear wave remain localized, keep its shape, and pass through another wave without being destroyed?}{A soliton is a self-reinforcing localized wave in which nonlinearity balances dispersion. It travels at constant speed, preserves its form, and after a collision reappears with only a phase shift.}{Water waves, optical fibers, plasmas, atmospheric waves, condensates, and nonlinear optics.}{A balance between dispersion and nonlinearity creates localized waves whose shape and collision behavior can be calculated.}

\paragraph{Discovery and definition}
Giorgio Bidone described a solitary water wave in 1826, and John Scott Russell observed one in 1834. Boussinesq and Rayleigh developed early mathematics. In 1895 Korteweg and de Vries derived the KdV equation. Zabusky and Kruskal coined soliton in 1965 after numerical experiments showed particle-like collisions \cite{soliton_ref1}.

The usual properties are permanence of form, localization, and elastic interaction. Some physical pulses called solitons lose energy or interact imperfectly, so the word is sometimes used more broadly \cite{soliton_ref2}.

\end{historylist}
\section*{3. Theory}
\subsection*{Core theory}
\paragraph{The balance of effects}
Dispersion makes different wavelengths travel at different speeds and tends to spread a pulse. Nonlinearity makes wave speed depend on amplitude and can steepen or focus it. A soliton occurs when these effects cancel. In an optical fiber, dispersion separates frequencies while the Kerr effect changes refractive index according to intensity.

The Korteweg--de Vries equation is
\[
u_t+6uu_x+u_{xxx}=0.
\]
It has traveling-wave solutions
\[
u(x,t)=2k^2\operatorname{sech}^2(k(x-4k^2t-x_0)).
\]
Taller solitons travel faster.

\paragraph{Integrability and applications}
The Korteweg--de Vries equation, nonlinear Schrodinger equation, sine--Gordon equation, and modified KdV equation are integrable models. The inverse scattering transform converts a nonlinear evolution into a scattering problem whose data evolve simply. Lax pairs encode this structure and generate conserved quantities \cite{soliton_ref3}.

Two KdV solitons pass through one another and retain their shapes, but each experiences a phase shift. Undular bores contain trains of solitary waves in rivers. Internal ocean waves travel along density interfaces. Atmospheric solitons occur in temperature inversion layers. Optical solitons allow pulses to travel long distances in fiber communication without dispersive spreading.

Solitons also occur in plasmas, Josephson junctions, magnetic materials, and field theories. De Broglie's double-solution idea motivated nonlinear wave models in quantum mechanics, although this is not the standard formulation of quantum theory \cite{soliton_ref2}.

\section*{4. Important theorems and results}
\begin{enumerate}[leftmargin=*,itemsep=0.35em]
\item \textbf{KdV soliton solution.} The KdV equation has localized traveling waves whose height and speed are linked.

\item \textbf{Inverse scattering transform.} Integrable nonlinear equations can be solved from scattering data that evolve simply.

\item \textbf{Lax-pair principle.} A compatible pair of operators can encode an integrable evolution and generate conserved quantities.

\item \textbf{Elastic collision property.} Ideal solitons emerge with their shapes and speeds unchanged except for phase shifts.

\item \textbf{Nonlinear-Schrodinger solitons.} In optical fibers, nonlinear self-phase modulation can balance dispersion.

The common message is that special nonlinear waves can behave like stable particles. The KdV solution gives a concrete pulse, the Lax pair and inverse-scattering transform explain why an integrable equation can be solved, and the elastic-collision property identifies the distinctive soliton behavior. Optical solitons use the same balance between dispersion and nonlinearity, but real fibers add loss and perturbations, so the ideal theorem is an approximation to an engineering model.
\end{enumerate}
\noindent\emph{Source for these results:} \cite{soliton_ref1}.

\theoryapplications
\theoryconnections{soliton_theory}{%
\item Nonlinear differential equations describe waves whose speed can depend on their height, so ordinary superposition no longer works.
\item Fourier and spectral analysis identify the modes, while geometry and inverse-scattering methods find special combinations whose dispersion and nonlinearity balance exactly.
\item That balance prevents a localized pulse from spreading as it travels; the pulse can collide with another one and re-emerge with its shape preserved, apart from a shift \cite{soliton_ref1,soliton_ref2}.
}{%
\item Soliton theory models pulses in optical fibers, shallow water, plasmas, and some condensed-matter systems.
\item In each case the useful connection is quantitative: the equation predicts a stable waveform, its speed, and how interactions change its position.
\item Integrable-system methods also provide exact test cases for numerical algorithms and reveal which features of a nonlinear PDE come from a hidden conserved quantity \cite{soliton_ref1,soliton_ref3}.
}
\section*{7. Sample Questions}
\emph{Exercise reference:} \cite{soliton_ref1}. The cited edition is the source for the exercise set; consult its Exercises section for the official statement and numbering.
\begin{enumerate}[leftmargin=*,itemsep=0.9em]
\item \textbf{Exercise 1. What two effects balance in a soliton?} Dispersion spreads a wave; nonlinearity changes speed with amplitude. Their cancellation preserves the pulse.

\item \textbf{Exercise 2. What is the speed of the KdV pulse above?} Its argument is $k(x-4k^2t-x_0)$, so its speed is $4k^2$.

\item \textbf{Exercise 3. Why does a taller KdV soliton travel faster?} Its height is $2k^2$ while its speed is $4k^2$.

\item \textbf{Exercise 4. Why is the optical Kerr effect relevant?} It makes refractive index depend on intensity, producing nonlinear phase modulation that can cancel dispersive spreading.

\item \textbf{Exercise 5. Why is a soliton not just any localized wave?} It must also have persistent propagation and robust collision behavior in the ideal model.

\end{enumerate}
\section*{8. References}
\begin{thebibliography}{9}
\bibitem{soliton_ref1} M. J. Ablowitz and H. Segur, \emph{Solitons and the Inverse Scattering Transform}, SIAM, 1981.
\bibitem{soliton_ref2} R. K. Dodd et al., \emph{Solitons and Nonlinear Wave Equations}, Academic Press, 1982.
\bibitem{soliton_ref3} P. D. Lax, \emph{Integrals of Nonlinear Equations of Evolution and Solitary Waves}, AMS, 1968.
\end{thebibliography}


\theory{Spectral theory}{How can the eigenvalue and eigenvector ideas of matrices be extended to differential operators and infinite-dimensional spaces?}{Spectral theory studies the spectrum of an operator, the values for which it lacks an inverse, and decomposes an operator into simpler modes. The spectrum may contain ordinary eigenvalues and continuous parts, so infinite-dimensional systems need more than a list of eigenvectors.}{Differential equations, quantum mechanics, vibrations, Fourier analysis, integral equations, spectral geometry, and engineering.}{Spectra and functional calculus decompose operators into modes that govern differential equations, vibrations, and quantum observables.}

\paragraph{The idea and history}
For a matrix, eigenvectors are directions that are only stretched, and eigenvalues are the stretching factors. David Hilbert developed an infinite-dimensional version while studying quadratic forms and infinitely many variables. Erhard Schmidt and Frigyes Riesz developed Hilbert-space foundations; von Neumann developed spectral theory for operators used in quantum mechanics \cite{spectral_ref1}.

The term spectrum was introduced in the mathematical setting before quantum mechanics connected it with atomic spectra. Hilbert was surprised that a theory developed from pure mathematical interest would explain physical energy levels \cite{spectral_ref2}.

\end{historylist}
\section*{3. Theory}
\subsection*{Core theory}
\paragraph{Spectrum and resolvent}
For an operator $T$, the spectrum $\sigma(T)$ is the set of complex numbers $\lambda$ for which $T-\lambda I$ has no bounded inverse. Every eigenvalue belongs to the spectrum, but the spectrum can contain values that are not eigenvalues. The complement is the resolvent set, where the inverse exists and varies analytically with $\lambda$.

On a Hilbert space, self-adjoint operators have real spectrum. Their spectrum may contain discrete eigenvalues or a continuous part. The multiplication operator $(Tf)(x)=xf(x)$ on square-integrable functions has every point in an interval in its spectrum even though no nonzero function is concentrated at one point \cite{spectral_ref3}.

\paragraph{The spectral theorem}
The finite-dimensional spectral theorem says a Hermitian matrix is unitarily diagonalizable. In an infinite-dimensional Hilbert space, a self-adjoint or normal operator is represented by a spectral measure:
\[
T=\int_{\sigma(T)}\lambda\,dE(\lambda).
\]
Functions of $T$ are defined by the same calculus, so expressions such as $e^{tT}$ and square roots can be constructed from the spectrum.

For compact self-adjoint operators, the nonzero spectrum consists of eigenvalues with finite multiplicity that can accumulate only at zero. Their eigenvectors form an orthonormal basis of the relevant closure. This generalizes principal-axis decomposition and explains why integral equations can be solved by modes \cite{spectral_ref1}.

\paragraph{Applications}
The vibration of a string, membrane, or drum becomes an eigenvalue problem for a differential operator. Frequencies are spectral values and eigenfunctions are normal modes. The question Can one hear the shape of a drum? asks how much geometric information is encoded by the spectrum.

In quantum mechanics, a self-adjoint operator represents an observable; its spectrum gives possible measurement values. In Fourier analysis, translation-invariant operators are diagonalized by frequency modes. In engineering, spectral radius and eigenvalues determine stability of discrete-time systems. Sturm--Liouville theory gives a powerful one-dimensional model for boundary-value problems \cite{spectral_ref4}.

\section*{4. Important theorems and results}
\begin{enumerate}[leftmargin=*,itemsep=0.35em]
\item \textbf{Spectral theorem.} Self-adjoint and normal operators on Hilbert spaces admit spectral-measure representations.

\item \textbf{Compact-operator theorem.} A compact self-adjoint operator has an orthonormal eigen-expansion with only possible accumulation at zero.

\item \textbf{Spectral mapping theorem.} For suitable functions $f$, the spectrum of $f(T)$ is $f(\sigma(T))$.

\item \textbf{Spectral-radius formula.} For a bounded operator, the spectral radius satisfies $r(T)=\lim_{n\to\infty}\|T^n\|^{1/n}$.

\item \textbf{Resolvent analyticity.} The inverse $(T-\lambda I)^{-1}$ is analytic on the resolvent set.

\item \textbf{Weyl's alternative.} For singular Sturm--Liouville endpoints, limit-point and limit-circle behavior determine the necessary boundary conditions.
\end{enumerate}
\noindent\emph{Source for these results:} \cite{spectral_ref1}.

\theoryapplications
\theoryconnections{spectral_theory}{%
\item Linear algebra begins with eigenvalues of a matrix; spectral theory asks the same question for operators on spaces of functions.
\item Hilbert-space geometry supplies orthogonality and convergence, functional analysis controls unbounded operators, and complex analysis studies the resolvent where an inverse exists.
\item Differential equations enter because boundary conditions select the eigenfunctions that can be combined to represent a solution \cite{spectral_ref1,spectral_ref2}.
}{%
\item Spectral theory helps vibration analysis by decomposing a structure into normal modes, each with its own frequency.
\item It helps quantum mechanics because the spectrum of an observable gives possible measurement values, and it helps PDEs because expanding the initial data in eigenfunctions converts evolution into independent scalar equations.
\item In numerical analysis, computed eigenvalues reveal stability, resonance, and the dominant directions of a large system \cite{spectral_ref1,spectral_ref3}.
}
\section*{Glossary}
\begin{description}[leftmargin=*,style=nextline,itemsep=0.3em]
\item[\textbf{Spectrum.}] The set of complex numbers for which an operator lacks a bounded inverse; it includes eigenvalues and may contain continuous parts.
\item[\textbf{Resolvent set.}] The complement of the spectrum where the inverse $(T-\lambda I)^{-1}$ exists and varies analytically with $\lambda$.
\item[\textbf{Spectral theorem.}] A representation showing that self-adjoint or normal operators can be expressed as integrals with respect to spectral measures.
\item[\textbf{Hilbert space.}] A complete inner product space that provides the setting for spectral theory, generalizing Euclidean space to infinite dimensions.
\item[\textbf{Self-adjoint operator.}] An operator equal to its adjoint; its spectrum is always real and it plays a central role in quantum mechanics.
\item[\textbf{Spectral measure.}] A measure-valued function that represents an operator in the spectral theorem, allowing functions of the operator to be defined.
\item[\textbf{Compact operator.}] An operator that maps bounded sets to relatively compact sets; for self-adjoint compact operators, the spectrum consists mostly of eigenvalues.
\item[\textbf{Normal mode.}] An eigenfunction of a differential operator that represents a coherent pattern of vibration or oscillation at a specific frequency.
\item[\textbf{Eigenvalue.}] A scalar $\lambda$ for which $Tv=\lambda v$ has a nonzero solution $v$; the building block of spectral decomposition.
\item[\textbf{Spectral radius.}] The supremum of absolute values of elements in the spectrum; governs stability of discrete-time systems.
\item[\textbf{Resolvent.}] The family of bounded inverse operators $(T-\lambda I)^{-1}$ defined on the resolvent set.
\item[\textbf{Orthonormal basis.}] A set of mutually orthogonal unit vectors forming a complete basis for a Hilbert space.
\end{description}
\section*{7. Sample Questions}
\emph{Exercise reference:} \cite{spectral_ref1}. The cited edition is the source for the exercise set; consult its Exercises section for the official statement and numbering.
\begin{enumerate}[leftmargin=*,itemsep=0.9em]
\item \textbf{Exercise 1. Find the spectrum of $A=\operatorname{diag}(2,5)$.} The matrix $A-\lambda I$ is invertible unless $\lambda=2$ or $5$, so $\sigma(A)=\{2,5\}$.

\item \textbf{Exercise 2. Why can a spectrum contain a non-eigenvalue?} An operator can fail to have a bounded inverse even when the equation $Tv=\lambda v$ has no nonzero solution. Continuous spectral behavior is the infinite-dimensional example.

\item \textbf{Exercise 3. What are the normal modes of a vibrating string?} They are eigenfunctions of the differential operator with the boundary conditions of the string. Each mode oscillates at its own spectral frequency.

\item \textbf{Exercise 4. Why is the spectrum of a self-adjoint operator real?} If $Tv=\lambda v$ for a nonzero vector, then $\langle Tv,v\rangle=\langle v,Tv\rangle$ forces $\lambda=\overline\lambda$. The full theorem extends this reality to continuous spectrum.

\item \textbf{Exercise 5. How does Fourier analysis diagonalize convolution?} Exponentials are eigenfunctions of translation and convolution. Each frequency is multiplied by the Fourier transform of the kernel, so the operator becomes multiplication in frequency space.

\end{enumerate}
\section*{8. References}
\begin{thebibliography}{9}
\bibitem{spectral_ref1} E. B. Davies, \emph{Spectral Theory and Differential Operators}, Cambridge, 1995.
\bibitem{spectral_ref2} J. Dieudonne, \emph{History of Functional Analysis}, Elsevier, 1981.
\bibitem{spectral_ref3} N. Dunford and J. T. Schwartz, \emph{Linear Operators, Part II}, Wiley, 1988.
\bibitem{spectral_ref4} G. Teschl, \emph{Mathematical Methods in Quantum Mechanics}, AMS, 2019.
\end{thebibliography}


\theory{String theory}{Can the elementary particles be understood as different vibrations of one-dimensional objects, in a framework that also includes quantum gravity?}{String theory replaces point particles by strings. Different vibrational states appear as different particles, and one closed-string state has the properties expected of a graviton. The framework uses quantum mechanics, relativity, extra dimensions, supersymmetry, branes, dualities, and conformal field theory.}{Quantum gravity, black holes, cosmology, nuclear and condensed-matter models, and pure mathematics.}{String vibrations, compactification, branes, and dualities connect particle states with quantum-gravitational geometry.}

\paragraph{The problem and history}
General relativity describes gravity and spacetime on large scales, while quantum mechanics describes microscopic phenomena. Applying ordinary quantum rules to gravity produces serious difficulties. String theory began in the late 1960s as a model of the strong force, but its unwanted features suggested a quantum theory of gravity \cite{string_ref1}.

The early bosonic theory described only bosons. Superstring theories added supersymmetry, relating bosons and fermions. Five consistent ten-dimensional superstring theories were identified: type I, type IIA, type IIB, and two heterotic theories. In the 1990s, dualities suggested that they are limits of an eleven-dimensional framework called M-theory \cite{string_ref2}.

\end{historylist}
\section*{3. Theory}
\subsection*{Core theory}
\paragraph{Strings, worldsheets, and branes}
An open string is a segment; a closed string is a loop. As a string moves, it sweeps out a two-dimensional worldsheet. Quantizing its vibrational modes produces particle states. A massless closed-string mode has spin two and is interpreted as the graviton.

A brane is a higher-dimensional object on which open strings can end. D-branes carry charges and help describe gauge fields. Compactification hides extra dimensions by making them very small. The shape of compact dimensions determines particle masses, charges, and symmetries in the lower-dimensional theory \cite{string_ref1}.

\paragraph{Dualities and correspondence}
T-duality relates a string on a circle of radius $R$ to one on a circle of radius proportional to $1/R$. Other dualities relate strong coupling to weak coupling and exchange different kinds of branes. These relations imply that apparently different theories are descriptions of one underlying structure.

The AdS/CFT correspondence, proposed in 1997, relates a gravitational string theory in an anti-de Sitter space to a conformal field theory on its boundary. It has been applied to black holes, nuclear matter, and strongly interacting condensed matter, although it is a theoretical correspondence rather than an experimentally confirmed description of nature \cite{string_ref3}.

String theory has counted microscopic states of certain black holes and has produced mathematical tools in geometry, topology, and representation theory. Compactifications on Calabi--Yau spaces and other manifolds have been used to build models resembling particle physics. However, the theory has many possible vacua, lacks a complete nonperturbative definition in all settings, and has not yet made a uniquely verified prediction about our universe \cite{string_ref4}.

\section*{4. Important theorems and results}
\begin{enumerate}[leftmargin=*,itemsep=0.35em]
\item \textbf{Graviton mode.} Closed-string quantization contains a massless spin-two state interpreted as the graviton.

\item \textbf{Supersymmetry.} Superstring theories relate bosonic and fermionic states and improve consistency.

\item \textbf{Anomaly cancellation.} Consistency restricts the allowed dimensions and gauge groups of superstring theories.

\item \textbf{Duality principle.} Different perturbative string descriptions can represent the same physics.

\item \textbf{AdS/CFT correspondence.} A gravitational theory in a bulk spacetime can be related to a conformal field theory on its boundary.

\item \textbf{Black-hole microstates.} In special supersymmetric examples, string theory counts microscopic states reproducing the Bekenstein--Hawking entropy.
\end{enumerate}
\noindent\emph{Source for these results:} \cite{string_ref1}.

\theoryapplications
\theoryconnections{string_theory}{%
\item Quantum mechanics supplies superposition and probabilities, while general relativity supplies a dynamical geometry of spacetime.
\item Replacing a point particle by a one-dimensional string changes the basic calculation from a particle path to a two-dimensional surface, whose consistency is governed by quantum field theory and conformal symmetry.
\item Extra dimensions, topology, and boundary conditions determine which vibrations are allowed and therefore which particle-like states appear \cite{string_ref1,string_ref2}.
}{%
\item String theory is useful in quantum gravity because the extended string softens some short-distance singularities and naturally includes a gravitational mode.
\item Its dualities relate apparently different physical descriptions, while branes and mirror symmetry have produced methods in geometry and topology.
\item Holographic models also connect a gravitational theory in one setting with a quantum field theory in another, giving concrete calculations for strongly interacting systems \cite{string_ref1,string_ref3,string_ref5}.
}
\section*{Glossary}
\begin{description}[leftmargin=*,style=nextline,itemsep=0.3em]
\item[\textbf{String.}] A one-dimensional extended object whose different vibrational modes correspond to different particles in string theory.
\item[\textbf{Brane.}] A higher-dimensional object on which open strings can end; D-branes carry charges and give rise to gauge fields.
\item[\textbf{Compactification.}] The process of making extra spatial dimensions very small so they are not directly observable, while their shape determines particle properties.
\item[\textbf{T-duality.}] A symmetry in string theory relating a string compactified on a circle of radius $R$ to one on a circle of radius proportional to $1/R$.
\item[\textbf{AdS/CFT correspondence.}] A conjectured duality relating a gravitational string theory in anti-de Sitter space to a conformal field theory on its boundary.
\item[\textbf{Superstring.}] A version of string theory that incorporates supersymmetry, relating bosonic and fermionic states and improving theoretical consistency.
\item[\textbf{M-theory.}] An eleven-dimensional framework that unifies the five consistent ten-dimensional superstring theories through various dualities.
\item[\textbf{Graviton.}] A massless spin-2 particle that emerges as a vibrational mode of closed strings and mediates the gravitational force.
\item[\textbf{Worldsheet.}] The two-dimensional surface swept out by a string moving through spacetime.
\item[\textbf{Supersymmetry.}] A symmetry relating bosonic and fermionic states; required for superstring consistency.
\item[\textbf{Anomaly cancellation.}] A consistency condition requiring certain quantum contributions to vanish, restricting allowed dimensions and gauge groups.
\item[\textbf{Conformal field theory.}] A quantum field theory invariant under local rescalings of the metric; the worldsheet theory of a string.
\item[\textbf{Calabi–Yau manifold.}] A compact Kähler manifold with vanishing first Chern class, used in string compactifications preserving supersymmetry.
\end{description}
\section*{7. Sample Questions}
\emph{Exercise reference:} \cite{string_ref1}. The cited edition is the source for the exercise set; consult its Exercises section for the official statement and numbering.
\begin{enumerate}[leftmargin=*,itemsep=0.9em]
\item \textbf{Exercise 1.} In bosonic string theory, the mass-squared of a closed-string state is $M^2 = \frac{2}{\alpha'}(N + \tilde{N} - 2)$, where $N$ and $\tilde{N}$ are the left- and right-moving level numbers. Find the values of $N$ and $\tilde{N}$ for the massless spin-two graviton state, and verify that it is massless. \emph{Solution.} The graviton has $N = \tilde{N} = 1$ (one oscillator excitation in each sector). Then $M^2 = \frac{2}{\alpha'}(1 + 1 - 2) = 0$. The state is massless. Its two independent polarizations give spin-two, matching the graviton.

\item \textbf{Exercise 2.} Under T-duality, a string compactified on a circle of radius $R$ is equivalent to one compactified on a circle of radius $R' = \alpha'/R$. If $\alpha' = 1$ and the original radius is $R = 3$, what is the T-dual radius? \emph{Solution.} $R' = \frac{\alpha'}{R} = \frac{1}{3}$. A string on a circle of radius $3$ is physically equivalent (in terms of the mass spectrum and interaction amplitudes) to a string on a circle of radius $\frac{1}{3}$.

\item \textbf{Exercise 3.} In the bosonic string, the ground state has $N = \tilde{N} = 0$. Compute its mass-squared and explain why it is problematic for a physical theory. \emph{Solution.} $M^2 = \frac{2}{\alpha'}(0 + 0 - 2) = -\frac{4}{\alpha'} < 0$. The ground state is a tachyon (negative mass-squared), which signals a quantum instability in the bosonic string. This is one reason the bosonic string is not a viable physical theory; superstrings resolve this by removing the tachyon through GSO projection.

\item \textbf{Exercise 4.} A closed string in bosonic string theory has left- and right-moving oscillator excitations $\alpha_{-1}^i$ and $\tilde{\alpha}_{-1}^j$ acting on the ground state, where $i, j = 1, \ldots, 24$. How many independent massless states (counting polarizations) can be formed from $N = \tilde{N} = 1$? \emph{Solution.} The state $\alpha_{-1}^i \tilde{\alpha}_{-1}^j |0\rangle$ has $24 \times 24 = 576$ components. The trace $\sum_i \alpha_{-1}^i \tilde{\alpha}_{-1}^i$ gives a scalar (the dilaton), the antisymmetric part $\alpha_{-1}^i \tilde{\alpha}_{-1}^j - \alpha_{-1}^j \tilde{\alpha}_{-1}^i$ gives $\binom{24}{2} = 276$ antisymmetric tensor states, and the remaining $576 - 1 - 276 = 299$ symmetric-traceless states are the graviton polarizations. The symmetric traceless part has $\frac{24 \times 25}{2} - 1 = 299$ components, matching.

\item \textbf{Exercise 5.} In the heterotic string, the critical dimension is $D = 26$ for the left-moving bosonic sector and $D = 10$ for the right-moving superstring sector. The difference $26 - 10 = 16$ extra dimensions are compactified. Explain how the gauge group $E_8 \times E_8$ or $\mathrm{SO}(32)$ arises from these 16 compactified dimensions. \emph{Solution.} The 16 left-moving dimensions compactified on a torus $T^{16}$ contribute 16 free left-moving bosons. Quantizing left-moving momenta on the torus gives a lattice of momenta. The requirement of modular invariance constrains this lattice to be even and self-dual in dimension 16. There are exactly two such lattices: the $E_8 \times E_8$ lattice and the $\mathrm{SO}(32)$ (or $\mathrm{D}_{16}$) lattice. These determine the gauge group of the resulting ten-dimensional theory.

\end{enumerate}
\section*{8. References}
\begin{thebibliography}{9}
\bibitem{string_ref1} J. Polchinski, \emph{String Theory, Volumes I and II}, Cambridge University Press, 1998.
\bibitem{string_ref2} B. Zwiebach, \emph{A First Course in String Theory}, Cambridge University Press, 2009.
\bibitem{string_ref3} M. B. Green, J. H. Schwarz, and E. Witten, \emph{Superstring Theory}, Cambridge University Press, 1987.
\bibitem{string_ref4} C. V. Johnson, \emph{D-Branes}, Cambridge University Press, 2003.
\bibitem{string_ref5} E. Kiritsis, \emph{String Theory in a Nutshell}, Cambridge University Press, 2007.
\end{thebibliography}


\theory{Sturm--Liouville theory}{How can a second-order differential equation reveal the natural modes and frequencies of a system?}{Sturm--Liouville theory studies self-adjoint boundary-value problems, their eigenvalues, eigenfunctions, orthogonality, completeness, and oscillation. It turns differential equations into a spectral decomposition problem.}{Vibrating strings and membranes, heat flow, quantum mechanics, Fourier series, acoustics, and numerical boundary-value problems.}{Self-adjoint differential operators, boundary conditions, and weighted orthogonality produce discrete modes and eigenvalue expansions.}

\paragraph{The problem and history}
Jacques Sturm and Joseph Liouville developed the theory in the nineteenth century. A regular problem has the form
\[
\frac{d}{dx}\left(p(x)y'\right)+q(x)y=-\lambda w(x)y
\]
on an interval, together with boundary conditions. The coefficient $w$ is a positive weight. The goal is to find values of $\lambda$ for which a nonzero solution exists and the corresponding eigenfunction \cite{sturm_ref1}.

\end{historylist}
\section*{3. Theory}
\subsection*{Core theory}
\paragraph{Eigenvalues and orthogonality}
For a regular problem, the eigenvalues are real, can be ordered
\[
\lambda_1<\lambda_2<\cdots,\qquad \lambda_n\to\infty,
\]
and each has a unique eigenfunction up to multiplication by a constant. Eigenfunctions for different eigenvalues are orthogonal under
\[
\langle f,g\rangle_w=\int_a^b f(x)g(x)w(x)\,dx.
\]
After normalization they form an orthonormal basis, so suitable functions can be expanded as an eigenfunction series \cite{sturm_ref2}.

\paragraph{A model problem and its interpretation}
The equation $-y''=\lambda y$ on $[0,\pi]$ with $y(0)=y(\pi)=0$ has eigenfunctions $\sin(nx)$ and eigenvalues $n^2$. A vibrating string is a sum of these normal modes. Separation of variables converts the heat, wave, and Laplace equations into Sturm--Liouville problems.

The one-dimensional time-independent Schrodinger equation is a Sturm--Liouville problem. Its eigenvalues are energy levels and eigenfunctions are stationary states. Bessel and Legendre equations can also be written in Sturm--Liouville form, explaining the orthogonality of special functions \cite{sturm_ref3}.

\paragraph{Singular problems and computation}
When an endpoint is infinite or a coefficient vanishes, the problem is singular. Weyl's limit-point and limit-circle alternative determines whether an additional boundary condition is needed. Numerical methods include shooting, finite differences, Rayleigh--Ritz, matrix variational methods, and spectral parameter power series.

\section*{4. Important theorems and results}
\begin{enumerate}[leftmargin=*,itemsep=0.35em]
\item \textbf{Reality and ordering.} Regular Sturm--Liouville eigenvalues are real, simple, and increase without bound.

\item \textbf{Orthogonality.} Eigenfunctions belonging to distinct eigenvalues are orthogonal with the weight $w$.

\item \textbf{Completeness.} Normalized eigenfunctions form an orthonormal basis of the weighted $L^2$ space.

\item \textbf{Oscillation theorem.} The $n$th eigenfunction has a controlled number of zeros in the interval.

\item \textbf{Rayleigh quotient principle.} The smallest eigenvalue minimizes an energy quotient among admissible nonzero functions.

\item \textbf{Weyl alternative.} Singular endpoints are classified as limit-point or limit-circle, determining boundary data.
\end{enumerate}
\noindent\emph{Source for these results:} \cite{sturm_ref1}.

\theoryapplications
\theoryconnections{sturm_liouville_theory}{%
\item An ordinary differential equation with boundary conditions becomes a self-adjoint operator problem.
\item Integration supplies the weighted inner product, Hilbert spaces make orthogonality and convergence precise, and spectral theory proves that the allowed eigenvalues are the natural frequencies of the problem.
\item The boundary conditions are essential: they select which modes fit the physical domain and make the operator's spectrum discrete in the standard compact cases \cite{sturm_ref1,sturm_ref2}.
}{%
\item The eigenfunctions form a basis for expanding a temperature, displacement, or wave profile, so a PDE can be solved one mode at a time.
\item In quantum mechanics they are stationary states; in acoustics they are resonant shapes; in numerical analysis they provide basis functions whose eigenvalues reveal stiffness and stability.
\item The connection works because self-adjointness makes different modes orthogonal and keeps the expansion coefficients independently computable \cite{sturm_ref1,sturm_ref3}.
}
\section*{Glossary}
\begin{description}[leftmargin=*,style=nextline,itemsep=0.3em]
\item[\textbf{Sturm–Liouville problem.}] A second-order differential equation with boundary conditions that leads to a self-adjoint eigenvalue problem.
\item[\textbf{Eigenvalue.}] A value $\lambda$ for which the differential equation has a nonzero solution satisfying the boundary conditions.
\item[\textbf{Eigenfunction.}] A nonzero solution corresponding to an eigenvalue; these form orthogonal basis functions.
\item[\textbf{Weight function.}] A positive function $w(x)$ that appears in the differential equation and defines the inner product for orthogonality.
\item[\textbf{Self-adjoint operator.}] An operator equal to its adjoint; ensures real eigenvalues and orthogonal eigenfunctions.
\item[\textbf{Boundary conditions.}] Constraints at the endpoints that select which solutions are physically admissible.
\item[\textbf{Orthogonality.}] The property that eigenfunctions for different eigenvalues integrate to zero with respect to the weight function.
\item[\textbf{Singular problem.}] A Sturm–Liouville problem where an endpoint is infinite or a coefficient vanishes, requiring special analysis.
\item[\textbf{Rayleigh quotient.}] The energy-like ratio $\langle f,Tf\rangle/\langle f,f\rangle$ achieving its minimum at the smallest eigenvalue.
\item[\textbf{Weyl alternative.}] The classification of singular endpoints as limit-point or limit-circle, determining whether additional boundary conditions are needed.
\item[\textbf{Separation of variables.}] A method of solving PDEs by writing solutions as products of single-variable functions.
\end{description}
\section*{7. Sample Questions}
\emph{Exercise reference:} \cite{sturm_ref1}. The cited edition is the source for the exercise set; consult its Exercises section for the official statement and numbering.
\begin{enumerate}[leftmargin=*,itemsep=0.9em]
\item \textbf{Exercise 1.} Solve the Sturm--Liouville problem $-y'' = \lambda y$ on $[0, \pi]$ with $y(0) = 0$ and $y'(\pi) = 0$. Find the first four eigenvalues. \emph{Solution.} For $\lambda > 0$, write $\lambda = \mu^2$. The general solution is $y = A\sin(\mu x) + B\cos(\mu x)$. From $y(0) = 0$: $B = 0$. From $y'(\pi) = 0$: $A\mu\cos(\mu\pi) = 0$, so $\mu\pi = (n - \frac{1}{2})\pi$, giving $\mu = n - \frac{1}{2}$ for $n = 1, 2, 3, \ldots$ The eigenvalues are $\lambda_n = (n - \frac{1}{2})^2$. The first four are $\frac{1}{4}$, $\frac{9}{4}$, $\frac{25}{4}$, and $\frac{49}{4}$.

\item \textbf{Exercise 2.} Verify that the eigenfunctions $y_n(x) = \sin\!\left((n - \frac{1}{2})x\right)$ and $y_m(x) = \sin\!\left((m - \frac{1}{2})x\right)$ with $n \neq m$ are orthogonal on $[0, \pi]$ with weight $w(x) = 1$. \emph{Solution.} Compute $\int_0^\pi \sin\!\left((n - \frac{1}{2})x\right)\sin\!\left((m - \frac{1}{2})x\right)dx$. Using the product-to-sum formula: $\frac{1}{2}\int_0^\pi \left[\cos\!\left((n - m)x\right) - \cos\!\left((n + m - 1)x\right)\right]dx$. For $n \neq m$ (both integers), $(n - m)$ and $(n + m - 1)$ are both nonzero integers. Then $\int_0^\pi \cos(kx)\,dx = \frac{\sin(k\pi)}{k} = 0$ for any nonzero integer $k$. Thus the integral is $0$, confirming orthogonality.

\item \textbf{Exercise 3.} For the Sturm--Liouville problem $-y'' = \lambda y$ on $[0,1]$ with $y(0) = y(1) = 0$, compute the Rayleigh quotient $R[f] = \frac{\int_0^1 (f')^2\,dx}{\int_0^1 f^2\,dx}$ for $f(x) = x(1-x)$, and compare it with the smallest eigenvalue $\lambda_1 = \pi^2$. \emph{Solution.} Compute $\int_0^1 f^2\,dx = \int_0^1 x^2(1-x)^2\,dx = \int_0^1 (x^2 - 2x^3 + x^4)\,dx = \frac{1}{3} - \frac{1}{2} + \frac{1}{5} = \frac{1}{30}$. Compute $f'(x) = 1 - 2x$, so $\int_0^1 (f')^2\,dx = \int_0^1 (1-2x)^2\,dx = \int_0^1 (1 - 4x + 4x^2)\,dx = 1 - 2 + \frac{4}{3} = \frac{1}{3}$. Then $R[f] = \frac{1/3}{1/30} = 10$. Since $\pi^2 \approx 9.87$, we have $R[f] = 10 > \pi^2 = \lambda_1$, consistent with the Rayleigh quotient principle (the minimum is $\lambda_1$, achieved by the first eigenfunction).

\item \textbf{Exercise 4.} Find the eigenvalues and normalized eigenfunctions of $-y'' = \lambda y$ on $[0, L]$ with $y(0) = y(L) = 0$. \emph{Solution.} The general solution for $\lambda = \mu^2 > 0$ is $y = A\sin(\mu x) + B\cos(\mu x)$. From $y(0) = 0$: $B = 0$. From $y(L) = 0$: $\sin(\mu L) = 0$, so $\mu L = n\pi$, giving $\mu_n = \frac{n\pi}{L}$ and $\lambda_n = \frac{n^2\pi^2}{L^2}$ for $n = 1, 2, 3, \ldots$ The normalized eigenfunctions (with $\|y_n\|^2 = 1$) are $\hat{y}_n(x) = \sqrt{\frac{2}{L}}\sin\!\left(\frac{n\pi x}{L}\right)$.

\item \textbf{Exercise 5.} Show that the oscillation theorem holds for the Sturm--Liouville problem $-y'' = \lambda y$ on $[0, \pi]$ with $y(0) = y(\pi) = 0$: verify that the $n$th eigenfunction $\sin(nx)$ has exactly $n - 1$ zeros in the open interval $(0, \pi)$. \emph{Solution.} The $n$th eigenfunction is $y_n(x) = \sin(nx)$. The zeros of $\sin(nx)$ in $(0, \pi)$ are at $x = \frac{k\pi}{n}$ for $k = 1, 2, \ldots, n - 1$. There are exactly $n - 1$ zeros in the open interval. For example, $y_1(x) = \sin(x)$ has $0$ zeros in $(0,\pi)$, $y_2(x) = \sin(2x)$ has $1$ zero at $x = \pi/2$, and $y_3(x) = \sin(3x)$ has $2$ zeros at $x = \pi/3$ and $x = 2\pi/3$. This confirms the oscillation theorem.

\end{enumerate}
\section*{8. References}
\begin{thebibliography}{9}
\bibitem{sturm_ref1} A. Zettl, \emph{Sturm--Liouville Theory}, Cambridge University Press, 2012.
\bibitem{sturm_ref2} R. Courant and D. Hilbert, \emph{Methods of Mathematical Physics}, Wiley-Interscience, 1989.
\bibitem{sturm_ref3} G. Teschl, \emph{Mathematical Methods in Quantum Mechanics}, American Mathematical Society, 2009.
\end{thebibliography}


\theory{Surgery theory}{How can one construct or classify manifolds by cutting out a controlled part and replacing it with another part?}{Surgery theory modifies a manifold along an embedded sphere, keeping track of the changes in homology, homotopy, and other invariants. In high dimensions it turns manifold classification into a sequence of obstruction problems.}{High-dimensional topology, cobordism, exotic spheres, manifold classification, and geometric topology.}{Cutting out a tubular neighborhood and regluing it changes a manifold in a controlled way whose homotopy and obstruction data can be computed.}

\paragraph{The operation and history}
For an $n$-manifold, remove an embedded $S^p$ with a tubular neighborhood $S^p\times D^q$ and glue in $D^{p+1}\times S^{q-1}$, where $n=p+q+1$. The boundary is the same, so the replacement produces another manifold. John Milnor called the method surgery; Andrew Wallace called it spherical modification \cite{surgery_ref1}.

Surgery arose from the classification of exotic spheres and became a central tool for manifolds of dimension greater than four. Morse theory shows that two manifolds are related by a sequence of such modifications exactly when they lie in the same cobordism class \cite{surgery_ref2}.

\end{historylist}
\section*{3. Theory}
\subsection*{Core theory}
\paragraph{What surgery controls}
The embedded sphere determines which homotopy class is changed. A surgery can kill a homotopy group, alter homology in a known degree, or simplify the fundamental group. The goal is to start with a manifold whose structure is known and produce one satisfying a desired property.

The main classification questions are whether a space with the homotopy type of a manifold is actually homeomorphic or diffeomorphic to one, and whether a homotopy equivalence between manifolds is homotopic to a diffeomorphism. For dimensions above four, surgery gives primary and secondary obstructions rather than a single complete invariant \cite{surgery_ref3}.

\paragraph{Worked example: surgery on $S^3$}
Consider the three-sphere $S^3$ and perform surgery on an embedded circle $S^1\subset S^3$. The normal bundle of $S^1$ in $S^3$ is trivial (since $S^3$ is parallelizable), so a tubular neighborhood looks like $S^1\times D^2$. Surgery means removing this $S^1\times D^2$ and gluing in $D^2\times S^1$ along the common boundary $S^1\times S^1$. The result is again $S^3$: one can see this by observing that $S^3$ is the boundary of $D^4$, and the surgery corresponds to attaching a 2-handle along the unknot, which does not change the manifold. More interestingly, surgery on a knotted $S^1\subset S^3$ (a nontrivial knot) can produce exotic manifolds. For example, surgery on the right-hand trefoil knot with framing $1$ yields the Poincaré homology sphere, a manifold with the same homology as $S^3$ but with $\pi_1\cong$ the binary icosahedral group of order 120. This shows that even in dimension three, surgery can change the topology in ways invisible to homology.

\paragraph{Obstructions and the s-cobordism theorem}
A normal map records a candidate manifold map together with stable normal bundle data. A primary obstruction in topological $K$-theory must vanish before a surgery obstruction can be defined. The secondary obstruction lies in an algebraic $L$-group. When it vanishes, surgery can produce the desired homotopy equivalence under suitable hypotheses.

The h-cobordism theorem says that an h-cobordism between simply connected manifolds of dimension at least five is a product. The s-cobordism theorem removes simple homotopy ambiguity and uses Whitehead torsion. These results explain why high-dimensional classification is tractable while dimensions three and four require different methods \cite{surgery_ref4}.

\section*{4. Important theorems and results}
\begin{enumerate}[leftmargin=*,itemsep=0.35em]
\item \textbf{Surgery construction.} Replacing $S^p\times D^q$ by $D^{p+1}\times S^{q-1}$ produces a controlled change of an $n$-manifold.

\item \textbf{Cobordism-surgery principle.} A cobordism can be decomposed into elementary surgeries.

\item \textbf{h-cobordism theorem.} A simply connected h-cobordism in dimension at least five is a product.

\item \textbf{s-cobordism theorem.} In the high-dimensional setting, vanishing Whitehead torsion characterizes when an h-cobordism is a product.

\item \textbf{Surgery exact sequence.} Manifold structures, normal invariants, and algebraic $L$-groups are linked by an exact sequence of classification problems.
\end{enumerate}
\noindent\emph{Source for these results:} \cite{surgery_ref1}.

\theoryapplications
\theoryconnections{surgery_theory}{%
\item Surgery removes a tubular neighborhood of an embedded sphere in a manifold and replaces it with another piece having the same boundary.
\item Homotopy and homology calculate how this replacement changes the space, cobordism records the whole transformation, and vector bundles describe the normal data needed to perform it.
\item K-theory and L-theory turn the remaining geometric question into an algebraic obstruction, especially when a homotopy equivalence is to be improved to a genuine equivalence of manifolds \cite{surgery_ref1,surgery_ref2}.
}{%
\item Surgery theory classifies high-dimensional manifolds by reducing them to a sequence of controlled modifications.
\item It explains exotic spheres by showing when the standard sphere can be altered without changing its homotopy type, and it supplies the obstruction groups used in the h- and s-cobordism theorems.
\item The connection is powerful but dimension-sensitive: the Whitney trick and general-position arguments that make the algebra reliable are much weaker in low dimensions \cite{surgery_ref3,surgery_ref4}.
\item \textbf{Three-manifold topology and the Poincar\'e conjecture.} Surgery theory was the conceptual engine behind the resolution of the Poincar\'e conjecture. Perelman's proof via Ricci flow can be viewed as a continuous, analytic version of surgery: as the flow progresses, singularities develop that are modeled on neck-pinch operations, and each singularity is resolved by a surgery that removes and replaces a cylindrical region. The classification program for three-manifolds---the Thurston geometrization conjecture, now a theorem---decomposes every closed 3-manifold along incompressible spheres and tori into pieces that each carry one of eight model geometries. Surgery-theoretic language (cutting along embedded spheres, tracking the effect on $\pi_1$, and using obstruction theory) provides the combinatorial framework that organizes this decomposition and led to Perelman's proof.
}
\section*{Glossary}
\begin{description}[leftmargin=*,style=nextline,itemsep=0.3em]
\item[\textbf{Surgery.}] A modification of a manifold by removing a tubular neighborhood of an embedded sphere and gluing in a different piece with the same boundary.
\item[\textbf{Tubular neighborhood.}] A neighborhood of a submanifold that looks like a disk bundle over the submanifold; in surgery it is $S^p\times D^q$.
\item[\textbf{Cobordism.}] A relationship between manifolds where one is the boundary of a higher-dimensional manifold; surgery decomposes cobordisms into elementary steps.
\item[\textbf{h-cobordism theorem.}] States that an h-cobordism between simply connected manifolds of dimension at least five is a product.
\item[\textbf{s-cobordism theorem.}] Refines the h-cobordism theorem by using Whitehead torsion to characterize when an h-cobordism is a product.
\item[\textbf{Whitehead torsion.}] An algebraic invariant that measures the failure of a homotopy equivalence to be a simple homotopy equivalence.
\item[\textbf{Normal map.}] A map between manifolds together with stable normal bundle data used to define surgery obstructions.
\item[\textbf{L-group.}] An algebraic group that classifies the secondary obstructions in surgery theory.
\item[\textbf{Exotic sphere.}] A smooth manifold homeomorphic but not diffeomorphic to the standard sphere.
\item[\textbf{K-theory.}] An algebraic theory of stable equivalence classes of vector bundles providing the primary obstruction group.
\item[\textbf{Whitney trick.}] A topological technique for removing intersection points using an embedded disk, available in high dimensions.
\item[\textbf{Normal bundle.}] The vector bundle of directions perpendicular to an embedded submanifold.
\end{description}
\section*{7. Sample Questions}
\emph{Exercise reference:} \cite{surgery_ref1}. The cited edition is the source for the exercise set; consult its Exercises section for the official statement and numbering.
\begin{enumerate}[leftmargin=*,itemsep=0.9em]
\item \textbf{Exercise 1.} In surgery on a 4-manifold $M^4$, suppose we perform surgery on an embedded $S^1 \subset M^4$. Identify $p$, $q$, and the pieces removed and glued in. \emph{Solution.} For $n = 4 = p + q + 1$ with $p = 1$, we get $q = 2$. The tubular neighborhood removed is $S^1 \times D^2$, and the piece glued in is $D^2 \times S^1$. Both have boundary $S^1 \times S^1$, so the gluing is along a torus.

\item \textbf{Exercise 2.} Perform surgery on the unknot $S^1 \subset S^3$. Compute the Euler characteristic before and after surgery. \emph{Solution.}$S^3$ is a closed 3-manifold with $\chi(S^3) = 0$ (since $H_0 = H_3 = \mathbb{Z}$ and $H_1 = H_2 = 0$, by Poincaré duality). Surgery on the unknot in $S^3$ produces $S^2 \times S^1$ (removing $S^1 \times D^2$ and gluing $D^2 \times S^1$). We compute $\chi(S^2 \times S^1) = \chi(S^2) \cdot \chi(S^1) = 2 \cdot 0 = 0$. The Euler characteristic is unchanged, consistent with the general fact that surgery on an $S^p$ does not change $\chi$.

\item \textbf{Exercise 3.} Compute the Betti numbers of the Poincaré homology sphere $\Sigma(2,3,5)$, which can be obtained by $(+1)$-surgery on the right-hand trefoil knot in $S^3$. \emph{Solution.}$\Sigma(2,3,5)$ is a rational homology sphere: $H_0 = \mathbb{Z}$, $H_1 = 0$ (the first homology of a Seifert fibered space over $S^2$ with relatively prime invariants), $H_2 = 0$ (by Poincaré duality with integer coefficients), and $H_3 = \mathbb{Z}$. The Betti numbers are $b_0 = 1$, $b_1 = 0$, $b_2 = 0$, $b_3 = 1$. The fundamental group is the binary icosahedral group of order 120, confirming $H_1 = \pi_1/[\pi_1, \pi_1] = 0$.

\item \textbf{Exercise 4.} In a surgery on $S^3$ along a knotted $S^1$, the framing determines the resulting manifold. Explain why surgery on the unknot with framing $n$ produces $S^2 \times S^1$ for any integer $n$ (up to orientation). \emph{Solution.} The unknot bounds an embedded disk $D^2$ in $S^3$. Surgery on the unknot with framing $n$ corresponds to attaching a 2-handle to $D^4$ along the unknot with framing $n$. Since the unknot has trivial knot group, the handle attachment does not create any nontrivial topology: it simply thickens the disk. The result is always $S^2 \times S^1$ (up to orientation), independent of $n$, because the surgery on the unknot is equivalent to a connected sum decomposition that produces the same manifold.

\item \textbf{Exercise 5.} Verify that the h-cobordism theorem requires dimension $\geq 5$ by explaining why the Whitney trick fails in dimension 4. \emph{Solution.}$H_1(X; \mathbb{Z}) \cong \mathbb{Z}$ has rank 1 (non-zero first homology). In the construction of the Whitney disk in 4-manifolds, the disk's interior can self-intersect because $\pi_1(X)$ is non-trivial. Specifically, the two arcs forming the Whitney disk boundary determine an element in $\pi_1$, and in dimension 4 the disk cannot avoid self-intersecting when this element is non-trivial. This means the Whitney trick fails to remove pairs of intersection points in dimension 4, unlike in dimension $\geq 5$ where there is enough room to make the disk embedded. This is why the h-cobordism theorem and surgery classification require dimension $\geq 5$.

\end{enumerate}
\section*{8. References}
\begin{thebibliography}{9}
\bibitem{surgery_ref1} J. Milnor, \emph{Collected Papers, Volume III: Differential Topology}, American Mathematical Society, 2010.
\bibitem{surgery_ref2} W. Browder, \emph{Surgery on Simply-Connected Manifolds}, Springer-Verlag, 1972.
\bibitem{surgery_ref3} C. T. C. Wall, \emph{Surgery on Compact Manifolds}, American Mathematical Society, 1999.
\bibitem{surgery_ref4} A. Ranicki, \emph{Algebraic and Geometric Surgery}, Oxford University Press, 2002.
\end{thebibliography}


\theory{Teichmuller theory}{How can we describe all complex structures that can be placed on the same topological surface?}{Teichmuller theory studies marked Riemann surfaces, quasiconformal maps, moduli, and the geometry of the space of conformal structures. It separates a surface's topology from the choice of complex shape.}{Complex analysis, hyperbolic geometry, mapping class groups, moduli spaces, geometric topology, dynamics, and string theory.}{Complex structures, quasiconformal maps, and quadratic differentials parameterize how a fixed surface can vary.}

\paragraph{The question and history}
Bernhard Riemann observed that a closed surface of genus $g\geq2$ requires $6g-6$ real parameters, or $3g-3$ complex parameters, to describe its complex structures. Oswald Teichmuller introduced quasiconformal mappings and the space now named after him. Lars Ahlfors and Lipman Bers supplied major analytic foundations; William Thurston later developed a geometric compactification \cite{teich_ref1}.

\end{historylist}
\section*{3. Theory}
\subsection*{Core theory}
\paragraph{Teichmuller space}
Fix a topological surface $S$. A point of Teichmuller space $\mathcal T(S)$ is a marked Riemann surface $(X,f)$, where $f:S\to X$ identifies the underlying topology. Two marked surfaces are equivalent when a conformal map between them respects the markings up to isotopy. The marking prevents different coordinate descriptions from being identified too early.

For genus $g\geq2$, $\mathcal T(S)$ is a complex manifold homeomorphic to a ball of real dimension $6g-6$. Forgetting the marking gives the moduli space of Riemann surfaces. The mapping class group acts on Teichmuller space, and the quotient is the moduli space, viewed as an orbifold because some surfaces have automorphisms \cite{teich_ref2}.

\paragraph{Quasiconformal maps}
A quasiconformal map is allowed to stretch angles by a bounded amount. In local complex coordinates its dilatation is controlled by
\[
\frac{|f_z|+|f_{\bar z}|}{|f_z|-|f_{\bar z}|}\leq K,\qquad K\geq1.
\]
The case $K=1$ is conformal. Teichmuller's theorem says that in each isotopy class between marked surfaces there is a unique map with minimum dilatation, the Teichmuller map \cite{teich_ref3}.

\paragraph{Metrics and quadratic differentials}
The minimum dilatation defines the Teichmuller metric. Holomorphic quadratic differentials determine the directions in which the extremal map stretches and contracts. The Bers embedding realizes Teichmuller space as a bounded domain in a complex vector space of quadratic differentials.

Uniformization identifies hyperbolic Riemann surfaces with quotients of the upper half-plane by Fuchsian groups. Fenchel--Nielsen coordinates describe hyperbolic structures by lengths and twists. These coordinates make the moduli problem geometrically computable \cite{teich_ref4}.

\paragraph{Applications}
Teichmuller theory helps classify surfaces, understand mapping class groups, study Kleinian and Fuchsian groups, and analyze dynamics of rational maps. It also appears in two-dimensional conformal field theory and string theory, where Riemann surfaces describe possible worldsheets.

\section*{4. Important theorems and results}
\begin{enumerate}[leftmargin=*,itemsep=0.35em]
\item \textbf{Teichmuller existence and uniqueness theorem.} Every isotopy class contains a unique quasiconformal map of least dilatation.

\item \textbf{Dimension theorem.} For genus $g\geq2$, Teichmuller space has complex dimension $3g-3$.

\item \textbf{Bers embedding.} Teichmuller space embeds as a bounded domain in a space of quadratic differentials.

\item \textbf{Uniformization theorem.} Simply connected Riemann surfaces are conformally equivalent to the disk, plane, or sphere.

\item \textbf{Mapping-class action.} The moduli space is the quotient of Teichmuller space by the mapping class group.

These results separate shape from topology. The existence and uniqueness theorem chooses the least-distorting map between marked surfaces; the dimension theorem counts the independent complex structures; and the Bers embedding gives a concrete analytic model. Forgetting the marking produces moduli space, where equivalent descriptions are identified and symmetries may create orbifold points.
\end{enumerate}
\noindent\emph{Source for these results:} \cite{teich_ref1}.

\theoryapplications
\theoryconnections{teichm_ller_theory}{%
\item Complex analysis describes charts on a Riemann surface, while quasiconformal maps measure how much a map stretches angles.
\item Hyperbolic geometry supplies a canonical metric for most surfaces, and topology fixes the genus and marking.
\item The extremal map between two marked surfaces gives a distance in Teichmuller space; quotienting by mapping-class actions forgets the marking and produces moduli space \cite{teich_ref1,teich_ref2}.
}{%
\item Teichmuller theory supplies a finite-dimensional space of complex structures, so geometric topology can study how a surface changes rather than one surface at a time.
\item Its geodesic and mapping-class dynamics describe iteration on moduli space, while conformal field theory and string theory use the same moduli to integrate over possible world-sheet geometries.
\item Quadratic differentials provide the concrete directions in which these deformations stretch \cite{teich_ref1,teich_ref5}.
}
\section*{Glossary}
\begin{description}[leftmargin=*,style=nextline,itemsep=0.3em]
\item[\textbf{Teichmüller space.}] The space of marked Riemann surfaces, where a marking identifies the underlying topology; it parametrizes complex structures up to isotopy.
\item[\textbf{Riemann surface.}] A one-dimensional complex manifold; a surface equipped with a complex structure that allows holomorphic functions.
\item[\textbf{Quasiconformal map.}] A map that stretches angles by at most a bounded amount, measured by a dilatation constant $K\geq1$; when $K=1$ the map is conformal.
\item[\textbf{Moduli space.}] The quotient of Teichmüller space by the mapping class group, parametrizing Riemann surfaces up to conformal equivalence.
\item[\textbf{Mapping class group.}] The group of isotopy classes of orientation-preserving diffeomorphisms of a surface; it acts on Teichmüller space.
\item[\textbf{Quadratic differential.}] A holomorphic object that describes the stretching directions of extremal quasiconformal maps and provides coordinates for Teichmüller space.
\item[\textbf{Uniformization theorem.}] States that every simply connected Riemann surface is conformally equivalent to the disk, plane, or sphere.
\item[\textbf{Bers embedding.}] Realizes Teichmüller space as a bounded domain in a complex vector space of quadratic differentials.
\item[\textbf{Fenchel–Nielsen coordinates.}] A coordinate system for Teichmüller space using hyperbolic lengths and twist parameters along a pants decomposition.
\item[\textbf{Fuchsian group.}] A discrete subgroup of $\mathrm{PSL}(2,\mathbb{R})$ acting on the hyperbolic plane.
\item[\textbf{Dilatation.}] The maximal local stretching factor $K\geq1$ of a quasiconformal map; the Teichmüller map minimizes this.
\item[\textbf{Isotopy.}] A continuous one-parameter family of diffeomorphisms connecting two maps.
\item[\textbf{Genus.}] The number of handles of a surface; determines the dimension $6g-6$ of Teichmüller space.
\end{description}
\section*{7. Sample Questions}
\emph{Exercise reference:} \cite{teich_ref1}. The cited edition is the source for the exercise set; consult its Exercises section for the official statement and numbering.
\begin{enumerate}[leftmargin=*,itemsep=0.9em]
\item \textbf{Exercise 1.} Compute the real dimension of the Teichmüller space $\mathcal{T}(S_g)$ for surfaces of genus $g = 2, 3, 4$. \emph{Solution.} The real dimension is $6g - 6$. For $g = 2$: $6(2) - 6 = 6$. For $g = 3$: $6(3) - 6 = 12$. For $g = 4$: $6(4) - 6 = 18$.

\item \textbf{Exercise 2.} A quasiconformal map $f\colon \mathbb{C} \to \mathbb{C}$ has dilatation $K = 3$. Using the formula $K = \frac{|f_z| + |f_{\bar{z}}|}{|f_z| - |f_{\bar{z}}|}$, express $|f_{\bar{z}}|$ in terms of $|f_z|$. \emph{Solution.} From $K = 3$, we get $3(|f_z| - |f_{\bar{z}}|) = |f_z| + |f_{\bar{z}}|$, so $3|f_z| - 3|f_{\bar{z}}| = |f_z| + |f_{\bar{z}}|$, giving $2|f_z| = 4|f_{\bar{z}}|$, hence $|f_{\bar{z}}| = \frac{1}{2}|f_z|$. The map stretches angles by a bounded amount: the maximal local stretching factor is $3$, and the ratio of the semi-axes of the infinitesimal ellipse is $K = 3$.

\item \textbf{Exercise 3.} Find the number of complex parameters needed to describe the conformal structures on a closed surface of genus $5$. \emph{Solution.} The complex dimension of Teichmüller space is $3g - 3 = 3(5) - 3 = 12$. So $12$ complex parameters (or $24$ real parameters) are needed.

\item \textbf{Exercise 4.} A Riemann surface of genus $g \geq 2$ has $6g - 6$ real-dimensional Teichmüller space, and the mapping class group $\mathrm{Mod}(S_g)$ acts on it with the quotient being the moduli space $\mathcal{M}_g$. For $g = 2$, verify that $\mathcal{M}_2$ is not a manifold by explaining why the action has fixed points. \emph{Solution.}$\mathrm{Mod}(S_2)$ contains the hyperelliptic involution $\iota$, which fixes the hyperelliptic locus in $\mathcal{T}(S_2)$. Since $\iota$ has a fixed point, the quotient $\mathcal{M}_2 = \mathcal{T}(S_2)/\mathrm{Mod}(S_2)$ has orbifold singularities (cone points) and is not a smooth manifold. Every genus-2 surface is hyperelliptic, so the fixed-point set is an open dense subset, and $\mathcal{M}_2$ is an orbifold throughout.

\item \textbf{Exercise 5.} Show that a conformal map (i.e.\ $K = 1$) preserves angles. \emph{Solution.} When $K = 1$, the inequality $\frac{|f_z| + |f_{\bar{z}}|}{|f_z| - |f_{\bar{z}}|} \leq 1$ forces $|f_{\bar{z}}| = 0$ (since $|f_z| \geq |f_{\bar{z}}| \geq 0$ and the numerator $\geq$ the denominator, with equality only when $|f_{\bar{z}}| = 0$). Thus $f_{\bar{z}} = 0$, meaning $f$ is holomorphic. A holomorphic function with nonzero derivative is conformal: it preserves angles between curves at every point where $f'(z) \neq 0$.

\end{enumerate}
\section*{8. References}
\begin{thebibliography}{9}
\bibitem{teich_ref1} J. Hubbard, \emph{Teichmuller Theory and Applications to Geometry, Topology, and Dynamics}, Springer-Verlag, 2006.
\bibitem{teich_ref2} A. F. Beardon and D. Minda, \emph{The Hyperbolic Metric and Geometry}, Cambridge University Press, 2009.
\bibitem{teich_ref3} L. Ahlfors, \emph{Lectures on Quasiconformal Mappings}, American Mathematical Society, 2006.
\bibitem{teich_ref4} J. Hubbard, \emph{Teichmuller Theory, Volume I}, American Mathematical Society, 2016.
\bibitem{teich_ref5} A. Papadopoulos, ed., \emph{Handbook of Teichmuller Theory}, European Mathematical Society, 2007.
\end{thebibliography}


\theory{Theory of computation}{Which problems can an algorithm solve, what resources does it need, and where are the limits of mechanical calculation?}{The theory of computation studies formal models of computation and divides into automata and formal languages, computability, and complexity. It separates impossible problems from solvable ones and efficient algorithms from impractical ones.}{Programming languages, compiler design, algorithms, cryptography, verification, artificial intelligence, and the foundations of mathematics.}{Machine models, reductions, and resource bounds distinguish what can be computed from what is undecidable or infeasible.}

\paragraph{The question and history}
Pioneers included Alonzo Church, Kurt Gödel, Alan Turing, Stephen Kleene, Emil Post, Rózsa Péter, John von Neumann, and Claude Shannon. Turing machines, lambda calculus, recursive functions, register machines, and rewriting systems are different models with the same broad notion of effective calculation, supporting the Church--Turing thesis \cite{computation_ref1}.

\end{historylist}
\section*{3. Theory}
\subsection*{Core theory}
\paragraph{Automata and languages}
\bookfigure{automata_computation}{Computation can be viewed as state transitions with increasing memory.}
Finite automata recognize regular languages, such as strings matching a simple pattern. Pushdown automata recognize context-free languages, which describe nested syntax such as balanced parentheses. Linear-bounded machines recognize context-sensitive languages, and Turing machines recognize recursively enumerable languages. These form the Chomsky hierarchy \cite{computation_ref2}.

\paragraph{Computability}
A set of natural numbers is decidable if a Turing machine halts and answers yes or no for every input. A function is computable if a machine halts with its value. The halting problem asks whether an arbitrary machine eventually stops on an input. Turing and Church proved that no algorithm decides this for all machines.

Rice's theorem says every nontrivial property of the partial function computed by a machine is undecidable. Matiyasevich's theorem solved Hilbert's tenth problem negatively: no algorithm decides whether an arbitrary Diophantine equation has an integer solution \cite{computation_ref3}.

\paragraph{Complexity}
Complexity theory studies resources such as time and space as functions of input size. The class P contains problems solvable in polynomial time. NP contains problems whose proposed solutions can be checked in polynomial time. The P versus NP problem asks whether every efficiently checkable problem is efficiently solvable.

Big-O notation suppresses machine-dependent constants and compares growth. Searching an unsorted list takes $O(n)$ steps, while binary search on a sorted list takes $O(\log n)$. Complexity classes, reductions, randomized computation, approximation, and circuit complexity organize the limits of efficient algorithms \cite{computation_ref4}.

\section*{4. Important theorems and results}
\begin{enumerate}[leftmargin=*,itemsep=0.35em]
\item \textbf{Church--Turing thesis.} All effectively calculable functions are captured by equivalent formal models such as Turing machines and lambda calculus.

\item \textbf{Halting theorem.} No Turing machine decides whether every machine halts on every input.

\item \textbf{Rice's theorem.} Every nontrivial semantic property of partial computable functions is undecidable.

\item \textbf{Cook--Levin theorem.} Boolean satisfiability is NP-complete.

\item \textbf{Matiyasevich theorem.} Hilbert's tenth problem has no algorithmic solution.

\item \textbf{Space and time hierarchies.} With more time or space, machines can solve strictly more problems in broad formal settings.
\end{enumerate}
\noindent\emph{Source for these results:} \cite{computation_ref1}.

\theoryapplications
\theoryconnections{theory_of_computation}{%
\item Logic defines what a computation is allowed to mean, discrete mathematics describes finite strings and graphs, and automata give simple machines whose behavior can be proved.
\item Algorithms provide the procedures being measured, while proof and probability distinguish guaranteed termination from expected performance.
\item Reductions are the key mechanism: they translate one problem into another while preserving the answer, so a lower bound for the first problem transfers to the second \cite{computation_ref1,computation_ref4}.
}{%
\item The theory supports programming languages by defining syntax, semantics, and the limits of what programs can compute.
\item It supports cryptography through hardness assumptions, verification through formal properties of all executions, and complexity theory through time and space bounds.
\item Thus a practical program can be judged not only by whether it works on examples, but by whether its specification is decidable and its resources are bounded \cite{computation_ref1,computation_ref2}.
}
\section*{Glossary}
\begin{description}[leftmargin=*,style=nextline,itemsep=0.3em]
\item[\textbf{Turing machine.}] A theoretical model of computation consisting of a tape, a head, and a state register; it formalizes the notion of algorithmic computation.
\item[\textbf{Decidable.}] A problem for which an algorithm always halts and gives a correct yes-or-no answer for every input.
\item[\textbf{Halting problem.}] The problem of determining whether an arbitrary algorithm will eventually stop running on a given input; it is undecidable.
\item[\textbf{P vs NP.}] The open question of whether every problem whose solution can be verified quickly can also be solved quickly.
\item[\textbf{NP-complete.}] The hardest problems in NP, such that an efficient solution to any one would solve all problems in NP.
\item[\textbf{Reduction.}] A transformation of one problem into another that preserves the answer, used to compare relative difficulty.
\item[\textbf{Chomsky hierarchy.}] A classification of formal languages by the type of automaton that recognizes them: regular, context-free, context-sensitive, and recursively enumerable.
\item[\textbf{Big-O notation.}] A notation describing the growth rate of a function, used to classify algorithm complexity independent of machine details.
\item[\textbf{Finite automaton.}] A computational model with finitely many states and no auxiliary memory; recognizes regular languages.
\item[\textbf{Pushdown automaton.}] A finite automaton augmented with a stack; recognizes context-free languages.
\item[\textbf{Regular language.}] A language recognizable by a finite automaton.
\item[\textbf{Context-free language.}] A language recognizable by a pushdown automaton; describes nested or recursive structures.
\item[\textbf{Church–Turing thesis.}] Every effectively calculable function is captured by equivalent formal models like Turing machines.
\item[\textbf{Rice's theorem.}] Every nontrivial semantic property of the partial function computed by a Turing machine is undecidable.
\end{description}
\section*{7. Sample Questions}
\emph{Exercise reference:} \cite{computation_ref1}. The cited edition is the source for the exercise set; consult its Exercises section for the official statement and numbering.
\begin{enumerate}[leftmargin=*,itemsep=0.9em]
\item \textbf{Exercise 1. Why is balanced-parenthesis recognition not regular?} A finite automaton has only finitely many states and cannot remember an arbitrarily large nesting depth. A pushdown automaton uses a stack, which supplies the needed memory.

\item \textbf{Exercise 2. What is undecidable about the halting problem?} It is not merely difficult; no algorithm can always halt and correctly decide whether an arbitrary program halts on an arbitrary input.

\item \textbf{Exercise 3. Why is binary search logarithmic?} Each comparison halves the remaining sorted interval. After $k$ comparisons at most $n/2^k$ entries remain, so $k$ is about $\log_2 n$.

\item \textbf{Exercise 4. What does NP mean?} It means that a proposed solution can be checked in polynomial time. It does not mean that every problem in NP is known to be hard.

\item \textbf{Exercise 5. Why are reductions useful?} A reduction transforms instances of one problem into another. If a known hard problem reduces to a new problem, an efficient algorithm for the new problem would solve the hard one.

\end{enumerate}
\section*{8. References}
\begin{thebibliography}{9}
\bibitem{computation_ref1} M. Sipser, \emph{Introduction to the Theory of Computation}, Cengage Learning, 2012.
\bibitem{computation_ref2} J. E. Hopcroft, R. Motwani, and J. D. Ullman, \emph{Introduction to Automata Theory, Languages, and Computation}, Addison-Wesley, 2007.
\bibitem{computation_ref3} H. Rogers, \emph{Theory of Recursive Functions and Effective Computability}, MIT Press, 1987.
\bibitem{computation_ref4} S. Arora and B. Barak, \emph{Computational Complexity}, Cambridge University Press, 2009.
\end{thebibliography}


\theory{Theory of equations}{When can polynomial equations be solved exactly, and when is a formula using radicals impossible?}{The classical theory of equations studies polynomial equations, their roots, formulas, and the algebraic structures that control solvability. Its central historical problem was solved by Galois theory, which replaced formulas with symmetry groups.}{Algebra, number theory, algebraic geometry, numerical root finding, physics, and the study of polynomial systems.}{Roots, multiplicities, resultants, and Galois symmetries explain both exact solvability and the limits of formulas.}

\paragraph{The historical problem}
For a polynomial equation, mathematicians first sought expressions for roots using addition, subtraction, multiplication, division, and radicals. Linear, quadratic, cubic, and quartic equations received formulas in antiquity and the Renaissance. Scipione del Ferro and Tartaglia found cubic methods; Cardano published them, and Ferrari supplied the quartic solution. Bombelli explained the complex numbers appearing in cubic formulas \cite{equations_ref1}.

The general quintic resisted radicals. Ruffini gave an incomplete impossibility argument, and Niels Abel proved in 1824 that a general fifth-degree equation cannot be solved by radicals. Évariste Galois then gave a criterion: the equation is solvable by radicals exactly when its Galois group is solvable \cite{equations_ref2}.

\end{historylist}
\section*{3. Theory}
\subsection*{Core theory}
\paragraph{Roots and systems}
The fundamental theorem of algebra says every nonconstant complex polynomial has a complex root, counted with multiplicity. Vieta's formulas relate coefficients to symmetric sums and products of roots. The discriminant detects repeated roots: it is zero precisely when the polynomial and its derivative share a root.

Linear systems were solved in antiquity; Cramer gave a general determinant formula in the eighteenth century. Today Gaussian elimination is usually more efficient. Systems of polynomial equations lead to resultants, Grobner bases, and algebraic geometry. Integer solutions are studied as Diophantine equations in number theory \cite{equations_ref3}.

\paragraph{A worked example}
For
\[
x^2-5x+6=0,
\]
factoring gives $(x-2)(x-3)=0$, so the roots are $2$ and $3$. Vieta's
formulas explain the calculation without factoring: the roots must add to
$5$ and multiply to $6$. The discriminant is
\[
(-5)^2-4(1)(6)=1,
\]
which is positive, so the two real roots are distinct.

The same questions become harder for a cubic or quintic. A numerical
algorithm may approximate the roots, but the theory of equations asks a
different question: is there a finite expression using radicals at all?
Galois theory answers that question by studying how the roots can be
permuted while leaving the coefficients fixed. Thus the elementary quadratic
example is a model of the larger subject: coefficients constrain roots,
symmetry explains the constraint, and the discriminant detects a change in
the root pattern \cite{equations_ref1,equations_ref2}.

\section*{4. Important theorems and results}
\begin{enumerate}[leftmargin=*,itemsep=0.35em]
\item \textbf{Fundamental theorem of algebra.} Every nonconstant complex polynomial has a complex root and factors completely into linear factors.

\item \textbf{Abel--Ruffini theorem.} The general polynomial of degree five or greater is not solvable by radicals.

\item \textbf{Galois criterion.} A polynomial is solvable by radicals exactly when its Galois group is solvable.

\item \textbf{Vieta's formulas.} Coefficients determine elementary symmetric functions of the roots.

\item \textbf{Discriminant criterion.} A polynomial has a repeated root exactly when its discriminant is zero.

\item \textbf{Bezout's theorem.} Projective plane curves of degrees $m$ and $n$ meet in $mn$ points counted with multiplicity under suitable hypotheses.
\end{enumerate}
\noindent\emph{Source for these results:} \cite{equations_ref1}.

\theoryapplications
\theoryconnections{theory_of_equations}{%
\item Algebra supplies operations on coefficients, complex numbers guarantee that a polynomial factors into linear factors after passing to a suitable field, and linear algebra handles systems and multiplicities.
\item Symmetry enters because permuting the roots can preserve every coefficient; the set of such permutations is the Galois group.
\item Polynomial geometry then interprets several equations as a shape whose intersections and singularities encode the solutions \cite{equations_ref1,equations_ref2}.
}{%
\item The theory feeds Galois theory by turning the question “can roots be written by radicals?” into a question about solvable groups.
\item It feeds algebraic geometry through varieties and elimination, number theory through integer and rational roots, and numerical analysis through algorithms that approximate roots when no useful formula exists.
\item In physics, characteristic equations determine allowed frequencies or energy levels, so root structure becomes a prediction about a system \cite{equations_ref2,equations_ref3}.
}
\section*{Glossary}
\begin{description}[leftmargin=*,style=nextline,itemsep=0.3em]
\item[\textbf{Polynomial equation.}] An equation of the form $a_nx^n + \cdots + a_1x + a_0 = 0$ where the unknown appears only with non-negative integer powers.
\item[\textbf{Root.}] A value of $x$ that satisfies a polynomial equation; also called a zero of the polynomial.
\item[\textbf{Multiplicity.}] The number of times a root appears as a factor in the complete factorization of a polynomial.
\item[\textbf{Discriminant.}] A quantity computed from the coefficients of a polynomial that is zero exactly when the polynomial has a repeated root.
\item[\textbf{Galois group.}] The group of permutations of a polynomial's roots that leave all coefficients fixed; it controls whether the polynomial is solvable by radicals.
\item[\textbf{Solvable by radicals.}] A property of a polynomial equation meaning its roots can be expressed using a finite combination of the coefficients and the operations of addition, subtraction, multiplication, division, and extraction of roots.
\item[\textbf{Vieta's formulas.}] Relations expressing the coefficients of a polynomial in terms of the elementary symmetric functions of its roots.
\item[\textbf{Fundamental theorem of algebra.}] Every nonconstant polynomial with complex coefficients has at least one complex root.
\item[\textbf{Abel–Ruffini theorem.}] No general formula using radicals exists for polynomial equations of degree five or higher.
\item[\textbf{Bézout's theorem.}] Two projective plane curves of degrees $m$ and $n$ meet in exactly $mn$ points counted with multiplicity.
\item[\textbf{Resultant.}] A polynomial in the coefficients of two polynomials that vanishes when they share a common root.
\item[\textbf{Symmetric function.}] A polynomial in the roots unchanged by any permutation of them.
\end{description}
\section*{7. Sample Questions}
\emph{Exercise reference:} \cite{equations_ref1}. The cited edition is the source for the exercise set; consult its Exercises section for the official statement and numbering.
\begin{enumerate}[leftmargin=*,itemsep=0.9em]
\item \textbf{Exercise 1.} Let $\alpha, \beta, \gamma$ be the roots of $x^3 - 6x^2 + 11x - 6 = 0$. Without finding the roots explicitly, compute $\alpha^2 + \beta^2 + \gamma^2$ using Vieta's formulas. \emph{Solution.}$\alpha + \beta + \gamma = 6$, $\alpha\beta + \beta\gamma + \gamma\alpha = 11$, $\alpha\beta\gamma = 6$. Then $\alpha^2 + \beta^2 + \gamma^2 = (\alpha + \beta + \gamma)^2 - 2(\alpha\beta + \beta\gamma + \gamma\alpha) = 36 - 22 = 14$.

\item \textbf{Exercise 2.} Compute the discriminant of $x^3 - 3x + 1$. Is it zero? \emph{Solution.}$p = -3$, $q = 1$. Discriminant: $\Delta = -4(-3)^3 - 27(1)^2 = 108 - 27 = 81 > 0$. Since $\Delta \neq 0$, the polynomial has three distinct roots. In fact $\Delta > 0$ for a cubic with real coefficients, confirming three distinct real roots.

\item \textbf{Exercise 3.} Compute the discriminant of $x^4 - 2x^3 + x^2 = x^2(x-1)^2$. \emph{Solution.}$\Delta = 256(-2)^3(-1)^4 - 192(-2)^2(1)^2(-1)^2 + \ldots$ Simpler: since $x^2(x-1)^2$ has repeated roots at $x=0$ and $x=1$ (each of multiplicity 2), the discriminant is $\Delta = 0$, as expected.

\item \textbf{Exercise 4.} Verify that $x^5 - x - 1$ is irreducible over $\mathbb{F}_2$ by checking it has no roots and no quadratic factors in $\mathbb{F}_2[x]$. \emph{Solution.} Check roots in $\mathbb{F}_2$: $f(0) = -1 = 1 \neq 0$, $f(1) = 1 - 1 - 1 = -1 = 1 \neq 0$. No linear factors. Check quadratic factors: the only irreducible quadratic over $\mathbb{F}_2$ is $x^2 + x + 1$. Dividing: $x^5 - x - 1 = (x^2 + x + 1)(x^3 + x^2) + (x^2 + x + 1) + 1$, so the remainder is nonzero. Since no irreducible quadratic divides $f$ and $f$ has no roots, $f$ is irreducible over $\mathbb{F}_2$, hence over $\mathbb{Z}$ by Gauss's lemma.

\item \textbf{Exercise 5.} Show that the roots of $x^2 - 2x + 2 = 0$ are complex conjugates, and compute them. \emph{Solution.}$\Delta = 4 - 8 = -4 < 0$, so the roots are complex. By the quadratic formula: $x = \frac{2 \pm \sqrt{-4}}{2} = \frac{2 \pm 2i}{2} = 1 \pm i$. The roots are $1 + i$ and $1 - i$, which are complex conjugates. Their sum is $2$ and product is $(1+i)(1-i) = 1 + 1 = 2$, matching Vieta's formulas.

\end{enumerate}
\section*{8. References}
\begin{thebibliography}{9}
\bibitem{equations_ref1} J. V. Uspensky, \emph{Theory of Equations}, McGraw-Hill, 1948.
\bibitem{equations_ref2} I. Stewart, \emph{Galois Theory}, Chapman and Hall/CRC, 2015.
\bibitem{equations_ref3} D. Cox, J. Little, and D. O'Shea, \emph{Ideals, Varieties, and Algorithms}, Springer-Verlag, 2015.
\end{thebibliography}


\theory{Theory of statistics}{How can limited, noisy observations support reliable conclusions about a much larger population or an uncertain process?}{Statistical theory gives mathematical principles for study design, data collection, summarization, inference, prediction, and decisions. It connects probability with optimization and utility so that methods can be compared and their uncertainty quantified.}{Experiments, surveys, medicine, engineering, economics, machine learning, public policy, and scientific data analysis.}{Sampling distributions, estimators, likelihoods, and error control turn data into quantified conclusions about a population.}

\paragraph{Ronald Fisher}
Fisher developed likelihood-based estimation, analysis of variance, and principles for randomized experimental design. These methods connected data collection with efficient estimation and made experimental uncertainty part of the mathematical model.

\paragraph{Jerzy Neyman and Egon Pearson}
Neyman and Pearson formalized hypothesis testing through error probabilities and power. Their framework distinguishes the rule used to make a decision from the long-run error guarantees attached to that rule \cite{statistics_ref1,statistics_ref2}.
\end{historylist}
\section*{3. Theory}
\subsection*{Core theory}
\paragraph{The problem and scope}
Data are observations from a population or process, often with measurement error and random variation. A statistical model describes how the data could have been generated. Statistical theory asks which questions can be answered, how much information is needed, and how reliable an answer is \cite{statistics_ref1}.

Its tasks include sampling from finite populations, measuring error, describing distributions, estimating parameters, testing hypotheses, predicting outcomes, comparing groups, and reducing dimension. It also studies what happens when assumptions such as independence or a chosen distribution are wrong \cite{statistics_ref2}.

\paragraph{Design before analysis}
Randomization makes treatment groups comparable and supports probability statements. Experimental design chooses factors, controls, replication, and sample size to estimate treatment effects efficiently. Survey sampling chooses a representative sample while accounting for unequal selection probabilities and nonresponse.

Good design can reduce cost and uncertainty before any formula is applied. Observational studies can estimate associations but usually require stronger assumptions to infer causation. Statistical theory separates the information supplied by the design from conclusions added by a model \cite{statistics_ref3}.

\paragraph{Inference}
An estimator is a rule that turns data into a parameter estimate. Bias, variance, consistency, and efficiency compare estimators. A confidence interval gives a procedure that covers the true parameter at a stated long-run rate; it is not a probability statement about a fixed parameter after the interval is observed.

Hypothesis testing compares a null model with alternatives. A p-value measures how surprising data would be under the null, while power measures the chance of detecting a specified alternative. Bayesian inference combines a prior distribution with a likelihood to produce a posterior distribution. Decision theory chooses actions by expected loss or utility \cite{statistics_ref4}.

\paragraph{Models and robustness}
Regression models describe relationships and prediction. Dimension reduction replaces many variables by a smaller informative representation. Model checking examines residuals, dependence, outliers, and calibration. Robust methods reduce sensitivity to incorrect distributional assumptions.

\section*{4. Important theorems and results}
\begin{enumerate}[leftmargin=*,itemsep=0.35em]
\item \textbf{Law of large numbers.} Sample averages approach population means under suitable conditions.

\item \textbf{Central limit theorem.} Properly standardized sums often approach a normal distribution, explaining many large-sample intervals.

\item \textbf{Neyman--Pearson lemma.} For testing two simple hypotheses at a fixed error level, the likelihood-ratio test is most powerful.

\item \textbf{Gauss--Markov theorem.} In a linear model with standard assumptions, least squares is the best linear unbiased estimator.

\item \textbf{Sufficiency principle.} A sufficient statistic retains all information about a parameter contained in the sample under the model.

\item \textbf{Bayes theorem.} Posterior odds equal prior odds multiplied by the likelihood ratio.
\end{enumerate}
\noindent\emph{Source for these results:} \cite{statistics_ref1}.

\theoryapplications
\theoryconnections{theory_of_statistics}{%
\item Probability describes the random mechanism that could produce the data, and statistics reverses that direction to infer unknown features of the mechanism.
\item Linear algebra makes regression and covariance computable, calculus derives likelihood equations, and optimization chooses parameters that fit the observations.
\item Decision theory adds the loss of acting on an estimate, so a statistically good answer is judged by consequences as well as by closeness to the data \cite{statistics_ref1,statistics_ref4}.
}{%
\item Statistical theory helps an experiment separate a real effect from sampling noise by comparing an estimate with its expected variability.
\item In medicine it turns trial data into risk estimates, in engineering it monitors failure rates, and in machine learning it controls prediction error on unseen data.
\item Surveys and public policy depend on the same connection, but only when sampling, missing data, dependence, and the decision costs are modeled explicitly \cite{statistics_ref3,statistics_ref5}.
}
\section*{Glossary}
\begin{description}[leftmargin=*,style=nextline,itemsep=0.3em]
\item[\textbf{Estimator.}] A rule that converts data into an estimate of an unknown parameter; different estimators can be compared by bias, variance, and efficiency.
\item[\textbf{Bias.}] The difference between the expected value of an estimator and the true parameter value; an unbiased estimator has zero bias on average.
\item[\textbf{Confidence interval.}] A procedure that, if repeated many times, covers the true parameter at a stated long-run rate; it is not a probability statement about a fixed parameter.
\item[\textbf{Hypothesis testing.}] A formal procedure for comparing a null model against alternatives using data, controlled by error probabilities.
\item[\textbf{p-value.}] The probability of obtaining data as extreme as or more extreme than observed, assuming the null hypothesis is true; it measures surprise under the null.
\item[\textbf{Power.}] The probability of rejecting the null hypothesis when a specific alternative is true; it measures the ability to detect a real effect.
\item[\textbf{Bayesian inference.}] A method that combines a prior distribution with a likelihood to produce a posterior distribution for unknown parameters.
\item[\textbf{Sufficiency principle.}] A sufficient statistic retains all information about a parameter contained in the sample under the model.
\item[\textbf{Likelihood.}] A function of parameters given observed data, proportional to the probability of the data under the model.
\item[\textbf{Variance.}] The expected squared deviation from the mean; measures spread of an estimator or data.
\item[\textbf{Likelihood-ratio test.}] A test comparing maximized likelihoods under two models; the Neyman–Pearson lemma shows it is most powerful for simple hypotheses.
\item[\textbf{Standard error.}] The standard deviation of the sampling distribution of a statistic.
\end{description}
\section*{7. Sample Questions}
\emph{Exercise reference:} \cite{statistics_ref1}. The cited edition is the source for the exercise set; consult its Exercises section for the official statement and numbering.
\begin{enumerate}[leftmargin=*,itemsep=0.9em]
\item \textbf{Exercise 1.} Let $X_1, \ldots, X_n$ be i.i.d.\ $\mathrm{Bernoulli}(p)$. Compute the maximum likelihood estimator $\hat{p}$ for $p$. \emph{Solution.}$L(p) = \prod_{i=1}^n p^{x_i}(1-p)^{1-x_i} = p^{s}(1-p)^{n-s}$ where $s = \sum x_i$. Then $\ell(p) = s\ln p + (n-s)\ln(1-p)$. Setting $\frac{d\ell}{dp} = \frac{s}{p} - \frac{n-s}{1-p} = 0$ gives $\hat{p} = \frac{s}{n} = \bar{X}$, the sample mean.

\item \textbf{Exercise 2.} Compute $\mathrm{Var}(\bar{X})$ for i.i.d.\ observations with variance $\sigma^2 = 9$ and sample size $n = 36$. Compare with the variance of a single observation. \emph{Solution.}$\mathrm{Var}(\bar{X}) = \frac{\sigma^2}{n} = \frac{9}{36} = 0.25$. The variance of a single observation is $9$. The standard error is $\sqrt{0.25} = 0.5$, compared to $\sqrt{9} = 3$ for a single observation. Averaging reduces the standard deviation by a factor of $\sqrt{n} = 6$.

\item \textbf{Exercise 3.} For the model $Y = \beta_0 + \beta_1 X + \varepsilon$ with $\varepsilon \sim N(0, \sigma^2)$, the least-squares estimator is $\hat{\beta}_1 = \frac{\sum(X_i - \bar{X})(Y_i - \bar{Y})}{\sum(X_i - \bar{X})^2}$. Show that $\hat{\beta}_1$ is unbiased when $\mathrm{E}[\varepsilon_i] = 0$ and the $X_i$ are fixed. \emph{Solution.}$\mathrm{E}[\hat{\beta}_1] = \frac{\sum(X_i - \bar{X})\mathrm{E}[Y_i - \bar{Y}]}{\sum(X_i - \bar{X})^2}$. Since $\mathrm{E}[Y_i] = \beta_0 + \beta_1 X_i$, we get $\mathrm{E}[Y_i - \bar{Y}] = \beta_1(X_i - \bar{X})$. Then $\mathrm{E}[\hat{\beta}_1] = \frac{\beta_1 \sum(X_i - \bar{X})^2}{\sum(X_i - \bar{X})^2} = \beta_1$. Thus $\hat{\beta}_1$ is an unbiased estimator of $\beta_1$.

\item \textbf{Exercise 4.} Suppose $X \sim \mathrm{Poisson}(\lambda)$ with observed value $x = 5$ in a single observation. Find the maximum likelihood estimate of $\lambda$. \emph{Solution.}$L(\lambda) = \frac{\lambda^5 e^{-\lambda}}{5!}$. Then $\ell(\lambda) = 5\ln\lambda - \lambda - \ln(5!)$. Setting $\frac{d\ell}{d\lambda} = \frac{5}{\lambda} - 1 = 0$ gives $\hat{\lambda} = 5$. The MLE is the observed count.

\item \textbf{Exercise 5.} A 95\% confidence interval for a population mean $\mu$ based on $n = 25$ observations with $\bar{x} = 10$ and $s = 4$ is $(\bar{x} - t_{0.025, 24} \cdot s/\sqrt{n},\; \bar{x} + t_{0.025, 24} \cdot s/\sqrt{n})$. Using $t_{0.025, 24} \approx 2.064$, compute the interval. \emph{Solution.}$s/\sqrt{n} = 4/5 = 0.8$. The margin of error is $2.064 \times 0.8 = 1.6512$. The confidence interval is $(10 - 1.6512,\; 10 + 1.6512) = (8.35,\; 11.65)$.

\end{enumerate}
\section*{8. References}
\begin{thebibliography}{9}
\bibitem{statistics_ref1} D. R. Cox and D. V. Hinkley, \emph{Theoretical Statistics}, Chapman and Hall/CRC, 1974.
\bibitem{statistics_ref2} E. L. Lehmann and J. P. Romano, \emph{Testing Statistical Hypotheses}, Springer-Verlag, 2005.
\bibitem{statistics_ref3} W. G. Cochran, \emph{Sampling Techniques}, Wiley, 1977.
\bibitem{statistics_ref4} G. Casella and R. L. Berger, \emph{Statistical Inference}, Cengage Learning, 2002.
\bibitem{statistics_ref5} D. A. Freedman, \emph{Statistical Models: Theory and Practice}, Cambridge University Press, 2009.
\end{thebibliography}


\theory{Topology}{Which properties of a space survive stretching, bending, and twisting when there is no tearing or gluing?}{Topology studies spaces through open sets, continuity, compactness, connectedness, dimension, holes, and invariants. It ignores exact distances while retaining how parts fit together.}{Analysis, geometry, differential equations, data analysis, dynamical systems, physics, and the classification of manifolds.}{Continuity, open sets, connectedness, and homotopy capture properties unchanged by bending or stretching without tearing.}

\paragraph{The idea and history}
Topology grew from Euler's solution of the Seven Bridges of Königsberg problem and his polyhedron formula. Listing introduced the term topology; Cantor studied point sets; Fréchet, Hausdorff, and Kuratowski developed metric and general topological spaces. Dennis Sullivan's work illustrates the continuing interaction of algebraic, geometric, and dynamical topology \cite{topology_ref1}.

\end{historylist}
\section*{3. Theory}
\subsection*{Core theory}
\paragraph{Spaces and continuity}
A topology on a set $X$ is a collection of open sets containing $\varnothing$ and $X$, closed under arbitrary unions and finite intersections. A topological space is a set with such a collection. A metric defines a topology through open balls, but many useful topologies have no preferred distance.

A function is continuous when inverse images of open sets are open. A homeomorphism is a continuous bijection with continuous inverse; it means two spaces have the same topological shape. A homotopy continuously deforms one map into another. Compactness says that every open cover has a finite subcover, and connectedness means the space cannot be split into two disjoint open pieces \cite{topology_ref2}.

\paragraph{Examples and invariants}
\bookfigure{topology_deformation}{Topology studies features that survive continuous deformation.}
A circle and a square are homeomorphic, while a circle and two disjoint circles are not connected in the same way. A coffee cup and a torus have one hole each in the relevant sense. The Mobius strip has one side and one boundary component. Euler characteristic, fundamental group, homology, and cohomology turn shape into computable algebra.

General topology studies open sets and convergence. Algebraic topology studies invariants such as homology and homotopy groups. Differential topology studies smooth manifolds, transversality, and tangent spaces. Geometric topology studies knots, manifolds, and embeddings \cite{topology_ref3}.

\section*{4. Important theorems and results}
\begin{enumerate}[leftmargin=*,itemsep=0.35em]
\item \textbf{Euler characteristic.} For a finite cell decomposition, $\chi=V-E+F$ is independent of the chosen decomposition.

\item \textbf{Intermediate value theorem.} A continuous real function taking values on both sides of a number takes that value somewhere between.

\item \textbf{Heine--Borel theorem.} A subset of Euclidean space is compact exactly when it is closed and bounded.

\item \textbf{Brouwer fixed-point theorem.} Every continuous map from a closed disk to itself has a fixed point.

\item \textbf{Jordan curve theorem.} A simple closed curve in the plane separates the plane into an inside and an outside.

\item \textbf{Classification of surfaces.} Connected compact surfaces are classified by orientability, genus, and boundary data.
\end{enumerate}
\noindent\emph{Source for these results:} \cite{topology_ref1}.

\theoryapplications
\theoryconnections{topology}{%
\item Sets and functions define open sets and continuous maps, while analysis explains convergence and compactness.
\item Geometry contributes spaces with coordinates, and algebra records features that survive continuous deformation.
\item The central mechanism is to discard measurements such as exact length while retaining relations such as connectedness or the existence of a loop; this makes a stretched rubber band and a circular one equivalent for the questions topology asks \cite{topology_ref1,topology_ref3}.
}{%
\item Topology helps manifold theory provide local coordinate neighborhoods, and algebraic topology converts holes and loops into groups and homology.
\item Dynamical systems use open sets and compactness to control long-term behavior, while functional analysis uses topological vector spaces to define convergence of functions.
\item In data analysis, a point cloud can be thickened into a shape and its persistent holes can be compared even when the measurements are noisy \cite{topology_ref1,topology_ref4}.
}
\section*{7. Sample Questions}
\emph{Exercise reference:} \cite{topology_ref1}. The cited edition is the source for the exercise set; consult its Exercises section for the official statement and numbering.
\begin{enumerate}[leftmargin=*,itemsep=0.9em]
\item \textbf{Exercise 1. Why are a circle and a square topologically equivalent?} A continuous bijection can map the boundary of the square once around the circle, and its inverse is continuous.

\item \textbf{Exercise 2. Compute the Euler characteristic of a tetrahedron.} It has $V=4$, $E=6$, and $F=4$, so $\chi=4-6+4=2$, the value for a sphere.

\item \textbf{Exercise 3. Is the interval $[0,1]$ connected?} Yes. If it were split into two nonempty separated open pieces, the order structure of the real line would be contradicted.

\item \textbf{Exercise 4. Why is a continuous image of a compact space compact?} An open cover of the image pulls back to an open cover of the domain; a finite subcover there maps to a finite subcover of the image.

\item \textbf{Exercise 5. What does the fundamental group measure?} It records loops based at a point up to continuous deformation and detects holes such as the central hole of a circle.

\end{enumerate}
\section*{8. References}
\begin{thebibliography}{9}
\bibitem{topology_ref1} J. R. Munkres, \emph{Topology}, Pearson, 2013.
\bibitem{topology_ref2} S. Willard, \emph{General Topology}, Dover Publications, 2004.
\bibitem{topology_ref3} A. Hatcher, \emph{Algebraic Topology}, Cambridge University Press, 2002.
\bibitem{topology_ref4} J. M. Lee, \emph{Introduction to Topological Manifolds}, Springer-Verlag, 2011.
\end{thebibliography}


\theory{Topos theory}{Can one build a mathematical universe that behaves like sets, but whose notion of space, logic, and truth is adapted to a particular problem?}{A topos is a category behaving like sheaves of sets on a space or site. Grothendieck topoi organize geometry and cohomology; elementary topoi provide an alternative foundation with internal logic that may be constructive rather than classical.}{Algebraic geometry, sheaf cohomology, logic, constructive mathematics, category theory, type theory, and theoretical computer science.}{Objects, morphisms, sheaves, and internal logic provide a generalized setting in which geometry and truth can vary together.}

\paragraph{Origin and intuition}
Topos theory arose when Grothendieck combined sheaves with categorical operations in algebraic geometry. Its important application is etale cohomology, which helped prove the Weil conjectures. Later Lawvere and Tierney emphasized topoi as logical universes; Olivia Caramello developed bridge methods between theories and topoi \cite{topos_ref1}.

The category of sets is the familiar example. A topos has objects, maps, products, function objects, a terminal object, and a subobject classifier. It therefore supports much of ordinary mathematics internally, but its logic can differ from classical set theory.

\end{historylist}
\section*{3. Theory}
\subsection*{Core theory}
\paragraph{Grothendieck and elementary topoi}
The category of sheaves of sets on a topological space is a Grothendieck topos. More generally, sheaves on a site form a topos, even when the covering objects are not ordinary open subsets. This lets algebraic geometers choose the topology best suited to the problem.

An elementary topos is defined by finite limits, exponentials, and a subobject classifier. It can interpret intuitionistic logic, so the law of excluded middle need not hold. A topos of $G$-sets is useful when a group symmetry should be built into every object \cite{topos_ref2}.

\paragraph{Concrete illustration: $G$-sets for a simple group}
Let $G=\mathbb{Z}/2\mathbb{Z}=\{0,1\}$ act on a set $X=\{a,b\}$ by swapping $a$ and $b$: the non-identity element of $G$ sends $a\mapsto b$ and $b\mapsto a$. A $G$-set is any set with such a $G$-action; a morphism of $G$-sets is a map $f\colon X\to Y$ satisfying $f(g\cdot x)=g\cdot f(x)$ for all $g\in G$, $x\in X$.

The category $\mathbf{Set}^{G}$ of all $G$-sets is a topos. Its terminal object is the one-element set with trivial $G$-action. Products are cartesian products with $G$ acting componentwise: $(X\times Y)_g=(x,y)\mapsto(g\cdot x,g\cdot y)$. The subobject classifier is $\Omega=\{0,1\}$ where $G$ acts trivially on both elements: the characteristic map of the $G$-invariant subset $\{a\}\subset\{a,b\}$ sends both $a$ and $b$ to $1$, reflecting that $\{a\}$ is not $G$-stable and so its ``truth value'' is the whole of $\Omega$.

This tiny example shows the topos machinery in action: internal logic (truth values in $\Omega$), subobject classification (invariant subsets), and the failure of excluded middle, since the proposition ``$x=a$'' has no definite truth value when $G$ may swap the points \cite{topos_ref1,topos_ref2}.

\paragraph{Logic and applications}
The internal language of a topos treats objects as types or sets and arrows as functions. A subobject is a predicate, and its classifier behaves like a truth-value object. In a Boolean topos the internal logic is classical; in a general topos it is intuitionistic. This makes topoi useful for constructive mathematics and semantics of computation.

Topoi connect category theory with algebraic geometry, homotopy theory, type theory, and logic. Classifying topoi represent geometric theories, while higher topoi extend the idea to homotopy-coherent spaces \cite{topos_ref3}.

\section*{4. Important theorems and results}
\begin{enumerate}[leftmargin=*,itemsep=0.35em]
\item \textbf{Topos sheaf theorem.} Sheaves on a site form a Grothendieck topos.

\item \textbf{Subobject classifier.} In a topos, subobjects of an object correspond to maps into a distinguished truth-value object.

\item \textbf{Exponential object.} Maps into an exponential $B^A$ correspond naturally to maps from a product with $A$ into $B$.

\item \textbf{Lawvere--Tierney topology.} Local notions of truth and sheafification inside a topos are described by suitable closure operators.

\item \textbf{Giraud principle.} Grothendieck topoi can be characterized abstractly by categorical exactness and generator conditions.

These results explain a topos at three levels. The sheaf theorem supplies the objects, the subobject classifier supplies an internal notion of truth, and exponentials supply an internal function space. Lawvere--Tierney topologies describe controlled changes of truth or sheafification. Giraud's principle is the recognition theorem: it tells when an abstract category has enough structure to behave like a category of sheaves.
\end{enumerate}
\noindent\emph{Source for these results:} \cite{topos_ref1}.

\theoryapplications
\theoryconnections{topos_theory}{%
\item Category theory supplies objects, arrows, and universal constructions, while sheaves attach local data to a space and specify how that data restricts and glues.
\item Internal logic allows the topos itself to determine what counts as a truth value, and algebraic geometry supplies sites whose covers reflect arithmetic rather than ordinary open sets.
\item Type theory gives a computational reading of the same logical structure \cite{topos_ref1,topos_ref2}.
}{%
\item Topos theory helps cohomology because sheaf data can be resolved and its failure to glue can be measured.
\item It helps constructive mathematics by providing a setting where excluded middle need not hold, and helps computer science by interpreting types, programs, and proofs in a common semantics.
\item Homotopy-theoretic topoi extend this idea by treating equivalent proofs or paths as geometric data rather than identifying everything strictly \cite{topos_ref1,topos_ref3}.
\item \textbf{Temporal logic and software verification.} In computer science, a program's execution over discrete time steps can be modeled as a $G$-set where $G=\mathbb{Z}$ acts by shifting one step forward. The topos of $\mathbb{Z}$-sets provides an internal logic in which propositions such as ``eventually $P$ holds'' are interpreted by quantifying over the orbit of a state under the $\mathbb{Z}$-action. Model-checking tools for linear temporal logic (LTL) exploit this connection: the truth value of a temporal formula at a state is determined by the entire forward orbit, which is exactly a $G$-set morphism into the subobject classifier of $\mathbf{Set}^{\mathbb{Z}}$ \cite{topos_ref1,topos_ref3}.
\item \textbf{Type theory and proof assistants.} The Curry--Howard correspondence identifies types with propositions and programs with proofs. An elementary topos provides a model of dependent type theory: function types correspond to exponential objects, product types to categorical products, and the unit type to the terminal object. Proof assistants such as Coq and Agda are based on type theories whose semantics can be interpreted in topoi, ensuring that a verified program is correct not only syntactically but also in a rich mathematical universe \cite{topos_ref2,topos_ref3}.
}
\section*{Glossary}
\begin{description}[leftmargin=*,style=nextline,itemsep=0.3em]
\item[\textbf{Topos.}] A category that behaves like the category of sets, supporting internal logic and geometric constructions.
\item[\textbf{Grothendieck topos.}] A category equivalent to the category of sheaves on a site; it organizes geometry and cohomology.
\item[\textbf{Elementary topos.}] A category with finite limits, exponentials, and a subobject classifier; it provides an alternative foundation for mathematics.
\item[\textbf{Subobject classifier.}] A distinguished object that classifies monomorphisms via characteristic maps, serving as a truth-value object in a topos.
\item[\textbf{Exponential object.}] An object $B^A$ representing functions from $A$ to $B$; it supports internal function spaces.
\item[\textbf{Internal logic.}] The logic interpreted within a topos, which may be intuitionistic rather than classical.
\item[\textbf{Sheaf.}] A contravariant functor from a site to sets that satisfies a gluing condition for local data.
\item[\textbf{Site.}] A category equipped with a Grothendieck topology that specifies covering families for defining sheaves.
\item[\textbf{Terminal object.}] An object $1$ such that for every $X$ there is exactly one morphism $X\to 1$.
\item[\textbf{$G$-set.}] A set equipped with an action by a group $G$; the category of $G$-sets is a topos.
\item[\textbf{Intuitionistic logic.}] A logic in which the law of excluded middle is not assumed; the internal logic of a general topos.
\item[\textbf{Lawvere–Tierney topology.}] A closure operator on subobjects inside a topos defining local truth and generalizing sheafification.
\end{description}
\section*{7. Sample Questions}
\emph{Exercise reference:} \cite{topos_ref1}. The cited edition is the source for the exercise set; consult its Exercises section for the official statement and numbering.
\begin{enumerate}[leftmargin=*,itemsep=0.9em]
\item \textbf{Exercise 1.} In the topos $\mathbf{Set}^G$ of $G$-sets where $G = \mathbb{Z}/2\mathbb{Z}$ acts on $X = \{a, b\}$ by swapping, determine all subobjects of $X$ and their characteristic maps to the subobject classifier $\Omega = \{0, 1\}$ (with trivial $G$-action). \emph{Solution.}$G$-invariant subobjects of $X$ must be $G$-stable subsets. The only $G$-stable subsets are $\varnothing$ and $X = \{a, b\}$. The subset $\{a\}$ is not $G$-stable since $g \cdot a = b \notin \{a\}$. Thus the subobjects are: $\varnothing$ with characteristic map $\chi_\varnothing(a) = \chi_\varnothing(b) = 0$, and $X$ with characteristic map $\chi_X(a) = \chi_X(b) = 1$. This differs from $\mathbf{Set}$, where $X$ would have four subobjects.

\item \textbf{Exercise 2.} In $\mathbf{Set}^G$ with $G = \mathbb{Z}/2\mathbb{Z}$ acting on $X = \{a,b\}$ by swapping, compute the product $X \times X$ as a $G$-set. \emph{Solution.}$X \times X = \{(a,a), (a,b), (b,a), (b,b)\}$ with $G$ acting componentwise: $g \cdot (x,y) = (g \cdot x, g \cdot y)$. Thus $g \cdot (a,a) = (b,b)$, $g \cdot (a,b) = (b,a)$, etc. The orbits are $\{(a,a), (b,b)\}$ and $\{(a,b), (b,a)\}$. The product has two orbits, each of size 2.

\item \textbf{Exercise 3.} In the topos $\mathbf{Set}$, compute the exponential object $2^3$ (where $2 = \{0,1\}$ and $3 = \{0,1,2\}$). \emph{Solution.}$2^3$ is the set of all functions from $\{0,1,2\}$ to $\{0,1\}$. There are $|2|^{|3|} = 2^3 = 8$ such functions. These correspond to the 8 subsets of $\{0,1,2\}$ (via characteristic functions): $\varnothing$, $\{0\}$, $\{1\}$, $\{2\}$, $\{0,1\}$, $\{0,2\}$, $\{1,2\}$, $\{0,1,2\}$.

\item \textbf{Exercise 4.} In the topos of sheaves on a topological space $X$ with the usual open-set topology, explain why the subobject classifier $\Omega$ assigns to each open set $U$ the set of open subsets $V \subseteq U$. \emph{Solution.}$\Omega$ is the sheaf defined by $\Omega(U) = \{V \subseteq U : V \text{ open}\}$. A subobject of a sheaf $\mathcal{F}$ over $U$ corresponds to a collection of sections $\mathcal{F}(V)$ for open $V \subseteq U$, and its characteristic map sends each section $s \in \mathcal{F}(V)$ to the open set $\{x \in V : s(x) \in \text{subobject}_x\}$. In the topological case, truth is local: each point has its own truth value (the open set containing it). This is the key difference from $\mathbf{Set}$, where $\Omega = \{0,1\}$.

\item \textbf{Exercise 5.} Show that in any topos, the product $A \times \Omega$ satisfies $A \times \Omega \cong A$ only if $\Omega \cong 1$ (which fails when the topos is not trivial). \emph{Solution.}$\Omega$ is the subobject classifier. In a topos, $\Omega \cong 1$ if and only if the topos is the degenerate topos (equivalent to $\mathbf{Set}/\{\text{pt}\}$). In any non-degenerate topos, $\Omega$ has at least two distinct global sections (corresponding to $\top$ and $\bot$), so $|\Omega| > 1$ in any model. If $A \times \Omega \cong A$ for all $A$, then taking $A = 1$ gives $\Omega \cong 1$, a contradiction. In non-degenerate topoi, $\Omega \not\cong 1$, and the product $A \times \Omega$ is in general strictly larger than $A$.

\end{enumerate}
\section*{8. References}
\begin{thebibliography}{9}
\bibitem{topos_ref1} S. Mac Lane and I. Moerdijk, \emph{Sheaves in Geometry and Logic}, Springer-Verlag, 1994.
\bibitem{topos_ref2} P. T. Johnstone, \emph{Topos Theory}, Academic Press, 1977.
\bibitem{topos_ref3} O. Caramello, \emph{Theories, Sites, Toposes}, Cambridge University Press, 2018.
\end{thebibliography}


\theory{Transcendental number theory}{Which numbers are not roots of any nonzero polynomial with rational coefficients, and how can we prove that fact?}{Transcendental number theory studies numbers beyond algebraic equations. It proves transcendence, algebraic independence, and quantitative lower bounds for polynomial expressions and linear forms in logarithms.}{Number theory, Diophantine approximation, special constants, exponential equations, cryptography, and arithmetic geometry.}{Algebraic independence, Diophantine approximation, and lower bounds for logarithmic forms distinguish numbers beyond polynomial equations.}

\paragraph{Definitions and history}
An algebraic number is a root of a nonzero polynomial with rational coefficients. A number that is not algebraic is transcendental. A set is algebraically independent over a field when no nonzero polynomial relation with coefficients in that field holds among its members \cite{trans_ref1}.

Liouville constructed transcendental numbers in the nineteenth century and proved that algebraic numbers cannot be approximated too well by rationals. Cantor showed that algebraic numbers are countable while real numbers are uncountable, so almost every real number is transcendental, although proving a particular constant transcendental can be very difficult \cite{trans_ref2}.

\end{historylist}
\section*{3. Theory}
\subsection*{Core theory}
\paragraph{Major theorems}
Liouville's theorem says that an algebraic number of degree $d\geq2$ cannot satisfy
\[
\left|\alpha-\frac pq\right|<\frac1{q^{d+\epsilon}}
\]
for infinitely many rationals $p/q$. Thus a number with exceptionally good rational approximations is transcendental.

The Lindemann--Weierstrass theorem implies that $\pi$ is transcendental: if $i\pi$ were algebraic and nonzero, then $e^{i\pi}$ would be transcendental, contradicting $e^{i\pi}=-1$. The Gelfond--Schneider theorem says that if $a$ is algebraic and not $0$ or $1$, and $b$ is irrational algebraic, then every value of $a^b$ is transcendental. Hence $2^{\sqrt2}$ and $e^\pi$ are transcendental \cite{trans_ref3}.

Alan Baker developed powerful lower bounds for linear forms in logarithms of algebraic numbers. These results solve or bound many Diophantine equations and give effective transcendence measures \cite{trans_ref4}.

\paragraph{Open problems and methods}
The theory asks whether constants such as Euler's constant, zeta values, and many combinations of logarithms are rational, algebraic, or transcendental. Auxiliary functions, Diophantine approximation, interpolation determinants, and height estimates are common tools. Algebraic independence is much harder than individual transcendence.

\section*{4. Important theorems and results}
\begin{enumerate}[leftmargin=*,itemsep=0.35em]
\item \textbf{Liouville criterion.} Exceptionally good rational approximation can force transcendence.

\item \textbf{Lindemann--Weierstrass theorem.} Linear independence of distinct algebraic exponents implies strong transcendence results for exponentials.

\item \textbf{Gelfond--Schneider theorem.} Algebraic $a\ne0,1$ raised to an irrational algebraic power is transcendental.

\item \textbf{Baker's theorem.} Nonzero linear forms in logarithms of algebraic numbers have explicit lower bounds.

\item \textbf{Cantor's cardinality result.} Almost all real numbers are transcendental because algebraic numbers are countable.

The results give increasingly strong ways to prove that a number is not algebraic. Liouville supplies an approximation test, Lindemann--Weierstrass controls exponentials, Gelfond--Schneider treats algebraic powers, and Baker gives quantitative lower bounds for logarithmic expressions. Cantor's argument proves abundance rather than identifying a particular number: it shows that algebraic numbers form a countable exceptional set inside the uncountable real line.
\end{enumerate}
\noindent\emph{Source for these results:} \cite{trans_ref1}.

\theoryapplications
\theoryconnections{transcendental_number_theory}{%
\item Number theory supplies algebraic numbers and polynomial equations, while complex analysis studies exponential and logarithmic functions that cannot be captured by polynomial identities.
\item Diophantine approximation measures how closely a number can be approached by rationals or algebraic numbers, and heights measure the arithmetic complexity of those approximants.
\item Lower bounds for linear forms in logarithms then prevent an equation from having infinitely many unexpectedly close solutions \cite{trans_ref1,trans_ref4}.
}{%
\item The theory helps exponential Diophantine equations by proving that expressions such as powers and logarithms cannot cancel too accurately.
\item It helps arithmetic geometry bound rational points on curves and distinguish algebraic from transcendental parameters.
\item Results about constants such as $e$ and $\pi$ also explain why certain exact numerical coincidences cannot arise from polynomial relations with rational coefficients \cite{trans_ref1,trans_ref2}.
}
\section*{Glossary}
\begin{description}[leftmargin=*,style=nextline,itemsep=0.3em]
\item[\textbf{Algebraic number.}] A number that is a root of a nonzero polynomial with rational coefficients.
\item[\textbf{Transcendental number.}] A number that is not algebraic, meaning it satisfies no polynomial equation with rational coefficients.
\item[\textbf{Algebraic independence.}] A property of a set of numbers where no nonzero polynomial relation with rational coefficients holds among them.
\item[\textbf{Diophantine approximation.}] The study of how closely a real number can be approximated by rational numbers, measured by the quality of the approximation.
\item[\textbf{Linear form in logarithms.}] A linear combination of logarithms of algebraic numbers; lower bounds for such forms are central to transcendence theory.
\item[\textbf{Liouville's theorem.}] A result stating that algebraic numbers of degree $d \geq 2$ cannot be approximated by rationals too well, leading to explicit transcendental numbers.
\item[\textbf{Lindemann–Weierstrass theorem.}] A theorem implying that if $\alpha$ is a nonzero algebraic number, then $e^\alpha$ is transcendental; it proves the transcendence of $\pi$.
\item[\textbf{Gelfond–Schneider theorem.}] A theorem stating that if $a$ is algebraic and not 0 or 1, and $b$ is irrational algebraic, then $a^b$ is transcendental.
\item[\textbf{Baker's theorem.}] A quantitative refinement of Lindemann–Weierstrass giving explicit lower bounds for linear forms in logarithms.
\item[\textbf{Height.}] A measure of arithmetic complexity of an algebraic number or field.
\item[\textbf{Countable.}] A set whose elements can be put in one-to-one correspondence with the natural numbers; algebraic numbers are countable but reals are not.
\end{description}
\section*{7. Sample Questions}
\emph{Exercise reference:} \cite{trans_ref1}. The cited edition is the source for the exercise set; consult its Exercises section for the official statement and numbering.
\begin{enumerate}[leftmargin=*,itemsep=0.9em]
\item \textbf{Exercise 1.} Prove that $e$ is irrational. \emph{Solution.}$e = \sum_{n=0}^{\infty} \frac{1}{n!}$. Suppose $e = a/b$ with $a, b \in \mathbb{Z}$, $b > 0$. Then $b! \cdot e = a \cdot (b-1)! = \sum_{n=0}^{b} \frac{b!}{n!} + \sum_{n=b+1}^{\infty} \frac{b!}{n!}$. The first sum is an integer. The second sum satisfies $0 < \sum_{n=b+1}^{\infty} \frac{b!}{n!} < \sum_{n=b+1}^{\infty} \frac{1}{(b+1)^{n-b}} = \frac{1}{b} < 1$, which is a contradiction since $b! \cdot e$ must be an integer. Thus $e$ is irrational.

\item \textbf{Exercise 2.} Using the Lindemann--Weierstrass theorem, prove that $\ln 2$ is transcendental. \emph{Solution.}$\ln 2$ is the natural logarithm of the algebraic number $2$. By the Lindemann--Weierstrass theorem (in its contrapositive form): if $\alpha$ is a nonzero algebraic number, then $e^\alpha$ is transcendental. If $\ln 2$ were algebraic, then $e^{\ln 2} = 2$ would be transcendental, but $2$ is algebraic (root of $x - 2 = 0$). This is a contradiction. Therefore $\ln 2$ is transcendental.

\item \textbf{Exercise 3.} By the Gelfond--Schneider theorem, prove that $2^{\sqrt{2}}$ is transcendental. \emph{Solution.}$a = 2$ is algebraic, $a \neq 0$ and $a \neq 1$. $b = \sqrt{2}$ is algebraic and irrational. By the Gelfond--Schneider theorem, $a^b = 2^{\sqrt{2}}$ is transcendental.

\item \textbf{Exercise 4.} Give a specific algebraic number $\alpha$ such that $e^\alpha$ is transcendental, citing the appropriate theorem. \emph{Solution.}$\alpha = 1$ is algebraic (root of $x - 1 = 0$) and nonzero. By the Lindemann--Weierstrass theorem, $e^1 = e$ is transcendental. Alternatively, $\alpha = \pi$ is transcendental, not algebraic, so this does not directly apply. Another example: $\alpha = \sqrt{2}$ is algebraic and nonzero, so $e^{\sqrt{2}}$ is transcendental by Lindemann--Weierstrass.

\item \textbf{Exercise 5.} Compute the first five convergents of the continued fraction $[1; 2, 2, 2, 2] = 1 + \cfrac{1}{2 + \cfrac{1}{2 + \cfrac{1}{2 + \cfrac{1}{2}}}}$ and show they approximate $\sqrt{2}$. \emph{Solution.}$p_{-1} = 1$, $p_0 = 1$, $q_{-1} = 0$, $q_0 = 1$. Using $p_n = a_n p_{n-1} + p_{n-2}$, $q_n = a_n q_{n-1} + q_{n-2}$: Convergent 0: $\frac{1}{1} = 1$. Convergent 1: $\frac{2 \cdot 1 + 1}{2 \cdot 1 + 0} = \frac{3}{2} = 1.5$. Convergent 2: $\frac{2 \cdot 3 + 1}{2 \cdot 2 + 1} = \frac{7}{5} = 1.4$. Convergent 3: $\frac{2 \cdot 7 + 3}{2 \cdot 5 + 2} = \frac{17}{12} \approx 1.4167$. Convergent 4: $\frac{2 \cdot 17 + 7}{2 \cdot 12 + 5} = \frac{41}{29} \approx 1.4138$. Since $\sqrt{2} \approx 1.41421$, the convergents $1, 1.5, 1.4, 1.4167, 1.4138$ approach $\sqrt{2}$, confirming that $\sqrt{2}$ has the purely periodic continued fraction $[1; \overline{2}]$.

\end{enumerate}
\section*{8. References}
\begin{thebibliography}{9}
\bibitem{trans_ref1} A. Baker, \emph{Transcendental Number Theory}, Cambridge University Press, 1990.
\bibitem{trans_ref2} E. B. Burger and R. Tubbs, \emph{Making Transcendence Transparent}, Springer-Verlag, 2004.
\bibitem{trans_ref3} S. Lang, \emph{Introduction to Transcendental Numbers}, Addison-Wesley, 1966.
\bibitem{trans_ref4} A. Baker and G. Wustholz, \emph{Logarithmic Forms and Diophantine Geometry}, Cambridge University Press, 2007.
\end{thebibliography}


\theory{Twistor theory}{Can spacetime geometry and quantum fields be represented more naturally by complex geometric data than by points and coordinates in spacetime?}{Twistor theory, proposed by Roger Penrose, treats complex twistor space as a fundamental arena from which spacetime and fields can be recovered.}{General relativity, quantum field theory, scattering amplitudes, integrable systems, and differential geometry.}{Spinors, complex geometry, and incidence relations recast spacetime and field equations in conformal variables.}

\paragraph{Origin and basic space}
Roger Penrose proposed twistor theory in 1967, influenced by Ivor Robinson's null congruences and by developments in relativity and complex analysis. Projective twistor space is $\mathbb{CP}^3$, while non-projective twistor space is a four-dimensional complex vector space with a Hermitian form of signature $(2,2)$ \cite{twistor_ref1}.

Twistor coordinates combine spinor information. The conformal group of compactified Minkowski space acts naturally on twistor space. A spacetime point corresponds to a projective line in twistor space, allowing spacetime fields to be encoded as complex-geometric objects \cite{twistor_ref2}.

\end{historylist}
\section*{3. Theory}
\subsection*{Core theory}
\paragraph{The Penrose transform}
Holomorphic functions and cohomology classes on twistor space correspond to solutions of massless field equations in spacetime. Contour integrals recover fields from twistor data. Penrose's nonlinear graviton construction represents self-dual solutions of Einstein's equations by deformations of complex structure. Ward's construction represents self-dual Yang--Mills fields using holomorphic vector bundles \cite{twistor_ref3}.

\paragraph{Concrete illustration: the Penrose transform in flat $\mathbb{R}^4$}
Consider flat Minkowski space $\mathbb{R}^4$ with coordinates $(t,x,y,z)$ and the associated projective twistor space $\mathbb{PT}\cong\mathbb{CP}^3$. A point $p\in\mathbb{R}^4$ determines a projective line $\mathbb{CP}^1\subset\mathbb{PT}$ via the incidence relation $Z^\alpha=p_{a\dot{a}}\pi^{a\dot{\alpha}}$ in spinor notation. Suppose one prescribes a holomorphic function $f(Z)$ on an open region of $\mathbb{PT}$. The Penrose transform recovers a field $\phi(x)$ on $\mathbb{R}^4$ by integrating $f$ over the $\mathbb{CP}^1$ corresponding to each spacetime point:
\[
\phi(x)\;=\;\oint_{\mathbb{CP}^1_x} f(Z)\;\omega,
\]
where $\omega$ is the natural holomorphic $1$-form on the twistor line. The output $\phi$ automatically satisfies the massless wave equation $\Box\,\phi=0$. In this flat-space example the computation is concrete: pick $f(Z)=Z^0/(Z^1 Z^2)$ on a suitable domain, perform the contour integral parametrizing the $\mathbb{CP}^1$ by the spinor ratio $\pi^1/\pi^0$, and read off the resulting solution $\phi(t,x,y,z)$ on spacetime. The PDE has been solved entirely by complex integration on twistor space.

\paragraph{Scattering amplitudes}
Twistor methods reorganized calculations in gauge theory. MHV formulas, twistor actions, BCFW recursion, Grassmannian integrals, and the amplituhedron express scattering amplitudes in geometric forms. Witten's twistor string theory linked perturbative Yang--Mills amplitudes with strings in twistor space \cite{twistor_ref4}.

\paragraph{A plain-language picture}
In ordinary spacetime, a light ray is described by where it is and which
direction it travels. Twistor theory changes the bookkeeping: it treats the
light-ray direction and its conformal information as the primary data. A
single spacetime event is then represented by the family of light rays that
pass through it. In projective twistor space that family is a line. The word
line is therefore not a claim that spacetime is one-dimensional; it is a
compact way to store all null directions through one event \cite{twistor_ref1,twistor_ref2}.

\textbf{Worked conceptual example.} Suppose a massless field is known on
twistor space as a holomorphic function or cohomology class. The Penrose
transform integrates that data along the projective line associated with each
spacetime point. The output is a field on spacetime satisfying the relevant
massless wave equation. The difficult differential equation has been
replaced by a holomorphic-data problem, and the transform explains how to
return to the physical field \cite{twistor_ref3}.

\section*{4. Important theorems and results}
\begin{enumerate}[leftmargin=*,itemsep=0.35em]
\item \textbf{Penrose transform.} Cohomology classes of suitable holomorphic twistor data correspond to massless fields in spacetime.

\item \textbf{Incidence correspondence.} A spacetime point corresponds to a projective line in twistor space.

\item \textbf{Nonlinear graviton theorem.} Self-dual conformal geometry can be reconstructed from appropriate complex deformations of twistor space.

\item \textbf{Ward construction.} Self-dual Yang--Mills fields correspond to holomorphic bundle data.

\item \textbf{BCFW recursion.} Scattering amplitudes can be built recursively from lower-point amplitudes.

The incidence correspondence changes the basic object: a spacetime point becomes a line in twistor space. The Penrose transform then turns holomorphic data on those lines into solutions of field equations. The nonlinear graviton and Ward constructions extend this idea to self-dual gravity and gauge fields, while BCFW recursion shows how related geometric and analytic structures can simplify amplitude calculations. These results require complex-analytic and self-dual hypotheses; they do not describe every spacetime or every field theory.
\end{enumerate}
\noindent\emph{Source for these results:} \cite{twistor_ref1}.

\theoryapplications
\theoryconnections{twistor_theory}{%
\item Twistor theory replaces spacetime points by geometric data describing null directions, so complex geometry and spinors encode the causal structure of relativity.
\item Sheaf cohomology turns holomorphic data on twistor space into fields on spacetime, while representation theory organizes the spin and symmetry labels.
\item Differential equations become statements about holomorphic bundles or cohomology classes, which can be easier to solve and transform back \cite{twistor_ref1,twistor_ref2}.
}{%
\item The transform helps construct self-dual solutions of relativity and gauge theory because a nonlinear spacetime equation becomes a bundle-gluing problem.
\item In integrable systems it exposes hidden complex geometry, and in scattering amplitudes it makes factorization and recursion visible.
\item These are not separate tricks: the same incidence relation translates spacetime propagation into holomorphic geometry \cite{twistor_ref3,twistor_ref4,twistor_ref5}.
\item \textbf{Computer vision.} Twistor-theoretic ideas enter computer vision through the study of the space of camera motions. The group of Euclidean motions $SE(3)$ acts on images, and its geometry can be described using a spinorial or twistorial parameterization. The incidence between light rays and image planes is structurally identical to the twistor incidence relation. Recent work on camera networks and multi-view geometry uses twistor-inspired representations to simplify the reconstruction of 3D scenes from 2D images, particularly when the camera calibration is unknown.
\item \textbf{Medical imaging and MRI reconstruction.} In MRI, data are collected in $k$-space (Fourier space) rather than directly in image space. The reconstruction problem---recovering the image from Fourier samples---is an inverse problem with geometric structure. Twistor methods have been applied to accelerate compressed-sensing MRI by exploiting the Penrose transform: instead of solving the reconstruction PDE in image space, one works with the simpler holomorphic representation on twistor space and transforms back. This approach can reduce the number of required samples and improve reconstruction quality for diffusion MRI, where the signal is naturally a solution of a massless equation on a curved background.
}
\section*{Glossary}
\begin{description}[leftmargin=*,style=nextline,itemsep=0.3em]
\item[\textbf{Twistor.}] A point in a complex vector or projective space that encodes spinor and conformal information about spacetime.
\item[\textbf{Penrose transform.}] A correspondence that converts holomorphic data or cohomology classes on twistor space into solutions of massless field equations in spacetime.
\item[\textbf{Incidence correspondence.}] The relationship where a spacetime point corresponds to a projective line in twistor space, linking geometry and field theory.
\item[\textbf{Self-dual fields.}] Fields satisfying a simplified version of the field equations that can be solved using complex geometric methods on twistor space.
\item[\textbf{Twistor space.}] A complex geometric arena from which spacetime and field equations can be recovered through holomorphic data.
\item[\textbf{Spinors.}] Mathematical objects describing rotational states; they provide the basic variables in twistor theory.
\item[\textbf{Cohomology classes.}] Algebraic objects representing field solutions; their vanishing or non-vanishing indicates solvability of field equations.
\item[\textbf{MHV formulas.}] Methods for computing scattering amplitudes in gauge theory using maximal helicity violating configurations.
\item[\textbf{Projective twistor space.}] The projectivization $\mathbb{CP}^3$ of the four-dimensional twistor space; each spacetime point corresponds to a projective line.
\item[\textbf{Minkowski space.}] The flat spacetime of special relativity, $\mathbb{R}^{3+1}$ with Lorentzian metric of signature $(-,+,+,+)$.
\item[\textbf{Conformal group.}] The group preserving the causal structure of spacetime; acts naturally on twistor space.
\item[\textbf{BCFW recursion.}] A recursive method for computing scattering amplitudes by building higher-point amplitudes from lower-point ones.
\end{description}
\section*{7. Sample Questions}
\emph{Exercise reference:} \cite{twistor_ref1}. The cited edition is the source for the exercise set; consult its Exercises section for the official statement and numbering.
\begin{enumerate}[leftmargin=*,itemsep=0.9em]
\item \textbf{Exercise 1.} In projective twistor space $\mathbb{PT} \cong \mathbb{CP}^3$, compute the complex dimension of $\mathbb{PT}$ and the complex dimension of the projective line $\mathbb{CP}^1$ corresponding to a single spacetime point. \emph{Solution.}$\mathbb{CP}^3$ has complex dimension $3$. The projective line $\mathbb{CP}^1 \subset \mathbb{CP}^3$ has complex dimension $1$.

\item \textbf{Exercise 2.} The twistor correspondence assigns to each spacetime point $p = (t, x, y, z) \in \mathbb{R}^4$ a projective line in $\mathbb{PT} \cong \mathbb{CP}^3$. In flat spacetime with coordinates $(Z^0, Z^1, Z^2, Z^3)$ on twistor space, the line corresponding to $p$ is determined by the incidence relation. Write down the equations of the line for $p = (1, 0, 0, 0)$. \emph{Solution.}$Z = p \cdot \pi$ in spinor notation means $Z^0 = \pi^0$, $Z^1 = \pi^1$, $Z^2 = -\pi^{0'}$,$ Z^3 = -\pi^{1'}$, where $(\pi^0, \pi^1)$ are spinor coordinates on the $\mathbb{CP}^1$. For $p = (1,0,0,0)$, the incidence relations give $Z^0 = \pi^0$, $Z^1 = \pi^1$, $Z^2 = -\bar{\pi}^1$, $Z^3 = -\bar{\pi}^0$ (in terms of the spinor ratio). The line is a one-parameter family in $\mathbb{CP}^3$ parametrized by $[\pi^0 : \pi^1]$.

\item \textbf{Exercise 3.}$\mathbb{CP}^3$ is covered by four affine charts. Compute the Euler characteristic $\chi(\mathbb{CP}^3)$ using the standard cell decomposition. \emph{Solution.}$\mathbb{CP}^3$ has a CW decomposition with one cell in each even dimension: $\mathbb{CP}^3 = e^0 \cup e^2 \cup e^4 \cup e^6$. Thus $\chi(\mathbb{CP}^3) = 1 - 0 + 1 - 0 + 1 - 0 + 1 = 4$.

\item \textbf{Exercise 4.} The Penrose transform recovers a spacetime field by integrating a holomorphic function $f$ over the $\mathbb{CP}^1$ associated with each point. If $f$ is a constant function $f \equiv 1$ on $\mathbb{PT}$, compute the resulting spacetime field $\phi(x)$. \emph{Solution.}$\phi(x) = \oint_{\mathbb{CP}^1_x} 1 \cdot \omega$. Since the integral is of a constant over a compact $\mathbb{CP}^1$ with holomorphic 1-form $\omega$, by residue considerations this integral is a constant independent of $x$. Thus $\phi$ is a constant function on spacetime, which indeed satisfies $\Box\phi = 0$.

\item \textbf{Exercise 5.}$\mathbb{CP}^3$ has real dimension $6$. A projective line $\mathbb{CP}^1$ in $\mathbb{CP}^3$ has real dimension $2$. The space of all lines in $\mathbb{CP}^3$ is the Grassmannian $\mathrm{Gr}(2,4)$, which is a complex manifold of complex dimension $2 \cdot (4-2) = 4$, hence real dimension $8$. Compute the real codimension of the family of twistor lines inside $\mathrm{Gr}(2,4)$. \emph{Solution.}$\mathrm{Gr}(2,4)$ has real dimension $8$. The family of lines arising from spacetime points has real dimension $4$ (the dimension of $\mathbb{R}^4$). The codimension is $8 - 4 = 4$. The twistor lines therefore form a 4-dimensional submanifold inside the 8-dimensional space of all lines.

\end{enumerate}
\section*{8. References}
\begin{thebibliography}{9}
\bibitem{twistor_ref1} R. Penrose and W. Rindler, \emph{Spinors and Space-Time}, Vol.~1, Cambridge University Press, 1984.
\bibitem{twistor_ref2} S. Huggett and K. Tod, \emph{An Introduction to Twistor Theory}, Cambridge University Press, 1985.
\bibitem{twistor_ref3} L. J. Mason and N. M. J. Woodhouse, \emph{Integrability, Self-Duality, and Twistor Theory}, Oxford University Press, 1996.
\bibitem{twistor_ref4} E. Witten, \emph{Perturbative Gauge Theory as a String Theory in Twistor Space}, Comm.\ Math.\ Phys., 2004.
\bibitem{twistor_ref5} M. Dunajski, \emph{Solitons, Instantons, and Twistors}, Oxford University Press, 2009.
\end{thebibliography}


\theory{Type theory}{How can a formal system prevent ill-formed constructions by recording what kind of object every expression is?}{Type theory classifies terms by types. Its rules say which terms may be formed, how they compute, and what counts as equality. Through Curry--Howard, types correspond to propositions and terms to proofs or programs.}{Programming languages, proof assistants, constructive mathematics, formal logic, category theory, linguistics, and homotopy theory.}{Typing judgments, dependent types, and normalization connect proofs with programs whose terms can be checked mechanically.}

\paragraph{Origin and history}
Bertrand Russell developed type hierarchies to avoid paradoxes such as the set of all sets that do not contain themselves. Church's simply typed lambda calculus connected types with functions. Per Martin-Lof developed intuitionistic type theory as a foundation for constructive mathematics. Calculus of constructions underlies proof assistants such as Coq and Lean \cite{type_ref1}.

\end{historylist}
\section*{3. Theory}
\subsection*{Core theory}
\paragraph{Terms, types, and rules}
A judgment may be written $\Gamma\vdash t:T$, meaning that under assumptions $\Gamma$, term $t$ has type $T$. Function types, product types, sum types, dependent products, and dependent sums are built by formation, introduction, elimination, and computation rules.

The type $A\to B$ is a function from $A$ to $B$. The product $A\times B$ is a pair. The sum $A+B$ records a choice between alternatives. A dependent product $\Pi x:A.P(x)$ expresses a statement for every $x$, while a dependent sum $\Sigma x:A.P(x)$ packages a witness and evidence \cite{type_ref2}.

\paragraph{Logic and computation}
Under Curry--Howard, a proof of $A\to B$ is a function that turns a proof of $A$ into a proof of $B$. A proof of $A\times B$ is a pair of proofs. A term reduces by computation rules, so a proof can carry an executable witness. Intuitionistic type theory does not automatically include the law of excluded middle, which makes it constructive.

Dependent types let the type of an expression depend on a value. A vector type can record its length, allowing the type checker to reject an out-of-bounds operation before execution. This is why type theory is central to verified programming and proof assistants \cite{type_ref3}.

\paragraph{Homotopy and category theory}
Homotopy type theory interprets types as spaces, terms as points, and equalities as paths. Univalence treats equivalent types as equal in a controlled sense. Cubical type theory gives computational models of these ideas. Cartesian closed categories model simply typed lambda calculus, while locally cartesian closed categories model dependent type theory \cite{type_ref4}.

\section*{4. Important theorems and results}
\begin{enumerate}[leftmargin=*,itemsep=0.35em]
\item \textbf{Curry--Howard correspondence.} Propositions are types, proofs are terms, and proof normalization is computation.

\item \textbf{Strong normalization.} In important typed calculi, every reduction sequence terminates.

\item \textbf{Progress and preservation.} Well-typed programs do not get stuck: they either are values or can take a step, and their types remain unchanged.

\item \textbf{Dependent-type correspondence.} Dependent products model universal quantification and dependent sums model existential quantification.

\item \textbf{Church--Rosser theorem.} If a term reduces in two ways, the resulting terms can be further reduced to a common result in confluent systems.

\item \textbf{Univalence principle.} In homotopy type theory, equivalences between types give paths witnessing equality of types.
\end{enumerate}
\noindent\emph{Source for these results:} \cite{type_ref1}.

\theoryapplications
\theoryconnections{type_theory}{%
\item Logic supplies propositions and inference, while the lambda calculus supplies functions as executable expressions.
\item A type records what kind of term a program or proof is allowed to be, so type checking catches an invalid application before it runs.
\item Proof theory explains normalization, category theory explains composition and universal structure, and computation gives the correspondence between constructing a term and carrying out an algorithm \cite{type_ref1,type_ref2}.
}{%
\item Type theory supports programming languages because a compiler can reject a term whose types do not fit, and it supports proof assistants because a checked proof is a typed term.
\item Dependent types let the type mention a value, such as the length of a vector, so the checker can verify that an operation preserves a specification.
\item Homotopy type theory extends the same idea by treating paths between objects as mathematical structure \cite{type_ref1,type_ref4}.
}
\section*{Glossary}
\begin{description}[leftmargin=*,style=nextline,itemsep=0.3em]
\item[\textbf{Type.}] A classification that determines what kind of object a term is; types prevent ill-formed expressions.
\item[\textbf{Term.}] An expression in type theory that can be a variable, abstraction, application, or other construct.
\item[\textbf{Curry–Howard correspondence.}] The principle that propositions correspond to types and proofs correspond to terms.
\item[\textbf{Dependent type.}] A type that can depend on a value, allowing the type system to express precise specifications.
\item[\textbf{Function type.}] The type $A \to B$ representing functions from terms of type $A$ to terms of type $B$.
\item[\textbf{Product type.}] A type $A \times B$ whose terms are pairs consisting of a term of type $A$ and a term of type $B$.
\item[\textbf{Sum type.}] A type $A + B$ representing a choice between a term of type $A$ or a term of type $B$.
\item[\textbf{Normalization.}] The process of reducing terms to a normal form; strong normalization guarantees that all reduction sequences terminate.
\item[\textbf{Lambda calculus.}] A formal system based on function abstraction and application; simply typed lambda calculus is modeled by Cartesian closed categories.
\item[\textbf{Dependent product.}] The type $\Pi x:A.\,P(x)$ of functions whose output type depends on the input value; models universal quantification.
\item[\textbf{Dependent sum.}] The type $\Sigma x:A.\,P(x)$ of pairs where the second component's type depends on the first; models existential quantification.
\item[\textbf{Homotopy type theory.}] A version of type theory where types are interpreted as topological spaces, terms as points, equalities as paths.
\item[\textbf{Univalence.}] The principle that equivalent types are equal; an axiom in homotopy type theory.
\end{description}
\section*{7. Sample Questions}
\emph{Exercise reference:} \cite{type_ref1}. The cited edition is the source for the exercise set; consult its Exercises section for the official statement and numbering.
\begin{enumerate}[leftmargin=*,itemsep=0.9em]
\item \textbf{Exercise 1. What is the type of the identity function?} If $x:A$, then $\lambda x.x$ has type $A\to A$.

\item \textbf{Exercise 2. What does a product type represent?} A term of $A\times B$ contains both an $A$-term and a $B$-term, analogous to a proof of $A$ and $B$.

\item \textbf{Exercise 3. Why are dependent types useful for arrays?} The type can include the length, so concatenation and indexing rules can be checked statically.

\item \textbf{Exercise 4. What is the empty type?} It has no terms. A function from it to any type is valid because no input exists that could violate the rule.

\item \textbf{Exercise 5. Why can type theory formalize mathematics?} Its rules can encode numbers, sets, functions, proofs, and constructions while the checker verifies every term.

\end{enumerate}
\section*{8. References}
\begin{thebibliography}{9}
\bibitem{type_ref1} B. C. Pierce, \emph{Types and Programming Languages}, MIT Press, 2002.
\bibitem{type_ref2} P. Martin-L{\"o}f, \emph{Intuitionistic Type Theory}, Bibliopolis, 1984.
\bibitem{type_ref3} T. Coquand and G. Huet, \emph{The Calculus of Constructions}, Information and Computation, 1988.
\bibitem{type_ref4} The Univalent Foundations Program, \emph{Homotopy Type Theory}, Institute for Advanced Study, 2013.
\end{thebibliography}


\theory{Wheel theory}{Can algebra make division total, including division by zero, without silently pretending ordinary inverse laws still hold?}{A wheel is an algebraic structure in which division is always defined. It extends a commutative ring or semiring with an inverse-like operation and a special value for $0/0$.}{Universal algebra, computer algebra, formal specification, and division by zero.}{Totalized division and changed algebraic identities provide a consistent setting in which division remains a defined operation at zero.}

\paragraph{Motivation and history}
In ordinary arithmetic, $a/b$ is undefined when $b=0$. Computer systems may need every operation to return a value. Jesper Carlstrom developed the theory of wheels in the early 2000s \cite{wheel_ref1}.

The name comes from viewing the real projective line as a circle and adding a bottom point representing $0/0$. A complex projective line is a sphere; adding the extra point gives a related geometric picture.

\end{historylist}
\section*{3. Theory}
\subsection*{Core theory}
\paragraph{Definition and changed rules}
A wheel has a set, constants $0$ and $1$, commutative associative addition and multiplication, and a unary operation written $/x$. Division is shorthand for $x/y=x\cdot(/y)$. The inverse operation is involutive:
\[
//x=x.
\]
It is multiplicative, and additional axioms modify familiar distributive rules so that expressions remain defined when a denominator is zero.

In a wheel, $x/x$ need not equal $1$, and $0x$ need not equal $0$ in the general case. These changes prevent illegal cancellation from turning $0/0$ into an ordinary number. The special value $0/0$ behaves like a bottom element \cite{wheel_ref2}.

\paragraph{Uses and interpretation}
Wheels provide a total algebraic data type for rational expressions. A computer algebra system can manipulate an expression without repeatedly checking whether a denominator vanishes; the result records the exceptional value instead of losing it. A wheel agrees with ordinary arithmetic where denominators are nonzero.

\paragraph{A worked example}
Consider the expression
\[
 \frac{x^2-1}{x-1}.
\]
For an ordinary number with $x\ne1$, factoring gives
\[
 \frac{(x-1)(x+1)}{x-1}=x+1.
\]
At $x=1$, the original expression is $0/0$, while the shortened expression
has value $2$. Ordinary algebra must therefore remember the condition
$x\ne1$; cancelling first and substituting later changes the question.

A wheel keeps the exceptional value visible. The original expression is
assigned the wheel's value for $0/0$, while the expression $x+1$ remains
equal to $2$ at $x=1$. The two expressions agree where cancellation is
legal, but they are not silently declared equal everywhere. This is the
practical benefit of totalized arithmetic: a symbolic program can continue
its calculation without confusing an exceptional value with an ordinary
number \cite{wheel_ref1,wheel_ref2}.

\section*{4. Important theorems and results}
\begin{enumerate}[leftmargin=*,itemsep=0.35em]
\item \textbf{Total division.} Every pair $x,y$ has a defined quotient $x/y$.

\item \textbf{Involution.} Applying the inverse operation twice returns the original element.

\item \textbf{Extension principle.} A commutative ring or semiring can be extended to a wheel under suitable constructions.

\item \textbf{Projective model.} The real projective line with an added bottom point gives the geometric picture behind a real wheel.

\item \textbf{Modified cancellation.} Identities such as $x/x=1$ are restricted because they fail at exceptional values.

The results describe how a wheel extends ordinary arithmetic while retaining a total division operation. Involution supplies the inverse-like structure, and the extension principle embeds familiar rings or semirings into the wider system. Modified cancellation is essential: division is always defined, but ordinary field identities cannot be applied at exceptional values without checking the wheel's rules.
\end{enumerate}
\noindent\emph{Source for these results:} \cite{wheel_ref1}.

\theoryapplications
\theoryconnections{wheel_theory}{%
\item Universal algebra asks which equations an operation must satisfy, while rings and semirings provide the ordinary addition and multiplication that a wheel extends.
\item The key change is to make division total: instead of leaving $1/0$ undefined, the structure assigns it an explicit value and records which familiar cancellation laws must be weakened.
\item Formal specification is needed because software must know exactly how exceptional values propagate \cite{wheel_ref1,wheel_ref2}.
}{%
\item Wheel theory helps computer algebra systems manipulate expressions without stopping whenever an intermediate denominator becomes zero.
\item A symbolic program can carry the exceptional value through addition, multiplication, and inversion, then report the result under the wheel's rules.
\item The connection is useful for formal verification because every expression has a defined denotation, but ordinary field identities may not be used unless their nonzero assumptions have been checked \cite{wheel_ref1,wheel_ref3}.
}
\section*{Glossary}
\begin{description}[leftmargin=*,style=nextline,itemsep=0.3em]
\item[\textbf{Wheel.}] An algebraic structure in which division is always defined, including division by zero, with modified cancellation laws.
\item[\textbf{Total division.}] The property that every pair of elements $x, y$ has a defined quotient $x/y$, even when $y = 0$.
\item[\textbf{Inverse operation.}] A unary operation written $/x$ that acts like a reciprocal; division is defined as $x/y = x \cdot (/y)$.
\item[\textbf{Involution.}] The property that applying the inverse operation twice returns the original element: $//x = x$.
\item[\textbf{Bottom element.}] The special value representing $0/0$ in a wheel, which behaves differently from ordinary numbers.
\item[\textbf{Extension principle.}] A construction that embeds a commutative ring or semiring into a wheel while preserving ordinary arithmetic where defined.
\item[\textbf{Modified cancellation.}] The restriction that identities like $x/x = 1$ cannot be applied at exceptional values without checking the wheel's rules.
\item[\textbf{Projective line.}] The real or complex projective line with an added bottom point, providing the geometric picture behind a wheel.
\item[\textbf{Commutative ring.}] An algebraic structure with addition and multiplication satisfying usual axioms; a wheel extends this to make division total.
\item[\textbf{Semiring.}] An algebraic structure like a ring but without requiring additive inverses; wheels can extend these.
\item[\textbf{Cancellation law.}] The rule that $ax=bx$ implies $a=b$; in wheels such laws must be weakened because division by zero is defined.
\end{description}
\section*{7. Sample Questions}
\emph{Exercise reference:} \cite{wheel_ref1}. The cited edition is the source for the exercise set; consult its Exercises section for the official statement and numbering.
\begin{enumerate}[leftmargin=*,itemsep=0.9em]
\item \textbf{Exercise 1.} In a wheel, the inverse operation satisfies $//x = x$ (involution). Using this, show that $/(1/a) = a$ for any element $a$. \emph{Solution.}$/(1/a) = //a = a$, applying involution twice. More carefully: the inverse operation is a unary operation $/(\cdot)$, so $/(1/a)$ means we apply $/$ to the element $1/a$. By involution, $/(1/a) = //a = a$.

\item \textbf{Exercise 2.}$\mathbb{R} \cup \{\infty, \bot\}$ forms a wheel (the extended real wheel). For $x, y \in \mathbb{R}$ with $y \neq 0$, compute $x/y$ in the wheel and confirm it agrees with ordinary division. \emph{Solution.}$x/y = x \cdot (/y) = x \cdot (1/y) = x/y$, which is the ordinary real quotient. When $y = 0$, $x/0 = x \cdot (/0)$ where $/0$ is the special element $\infty$ in the wheel. Thus $x/0$ is defined and equals $\infty$ (for $x \neq 0$), or $\bot$ (for $x = 0$, since $0/0 = \bot$). The wheel agrees with ordinary arithmetic where the denominator is nonzero.

\item \textbf{Exercise 3.} In a wheel, $x/x$ need not equal $1$ for all $x$. Show that in the real wheel $\mathbb{R} \cup \{\infty, \bot\}$, $0/0 = \bot \neq 1$. \emph{Solution.}$0/0 = 0 \cdot (/0) = 0 \cdot \infty$. In the real wheel, this is defined as $\bot$ (the bottom element), not $1$. This prevents the contradiction: if $0/0 = 1$, then from $0 \cdot a = 0 \cdot b$ we could conclude $a = b$ by canceling $0/0$, which is false. The wheel assigns $\bot$ to maintain consistency.

\item \textbf{Exercise 4.} Consider the expression $\frac{x^2 - 4}{x - 2}$. In ordinary algebra, this equals $x + 2$ for $x \neq 2$. In the real wheel, what is the value at $x = 2$? \emph{Solution.}$\frac{x^2 - 4}{x - 2} = \frac{(x-2)(x+2)}{x - 2}$. At $x = 2$: the numerator is $0$ and the denominator is $0$, so the value is $0/0 = \bot$. The wheel correctly distinguishes $\frac{x^2-4}{x-2}$ from $x + 2$ at $x = 2$: the former gives $\bot$, while the latter gives $4$.

\item \textbf{Exercise 5.}$\mathbb{CP}^1$ (the Riemann sphere) is the complex projective line. Explain how it relates to the complex wheel by adding a bottom element. \emph{Solution.}$\mathbb{CP}^1 = \mathbb{C} \cup \{\infty\}$ is the one-point compactification of $\mathbb{C}$, and forms a wheel with $\infty$ playing the role of $1/0$ and $\bot$ playing the role of $0/0$. The complex wheel extends $\mathbb{C}$ by adding two special elements: $\infty$ (for $z/0$ with $z \neq 0$) and $\bot$ (for $0/0$). The Riemann sphere $\mathbb{CP}^1$ captures the $\infty$ part; the full complex wheel additionally includes $\bot$ to make all divisions total.

\end{enumerate}
\section*{8. References}
\begin{thebibliography}{9}
\bibitem{wheel_ref1} J. Carlstr{\"o}m, \emph{Wheels: On Division by Zero}, Licentiate thesis, Chalmers University of Technology, 2001.
\bibitem{wheel_ref2} J. A. Bergstra and J. V. Tucker, \emph{The Rational Numbers as an Abstract Data Type}, J.\ ACM, 2007.
\bibitem{wheel_ref3} A. Setzer, \emph{Wheels}, Manuscript, University of Nottingham, 2001.
\end{thebibliography}


\part{Worked modules: learning by doing}
The entries above are a map. The modules in this part are the book's actual lessons. Each begins with a small problem, makes the hidden assumptions visible, and ends with exercises. Try the exercises before reading the answer; explaining a result in your own words is more valuable than memorizing a definition.

\chapter{A few examples before the modules}
\section*{Example 1: probability and a medical test}
Suppose 1 person in 100 has a condition. A test correctly detects 99 out of 100 affected people, but falsely reports ``positive'' for 5 out of 100 healthy people. In 10,000 people, about 100 are affected and 99 of them test positive. Of the 9,900 healthy people, about 495 also test positive. Therefore only $99/(99+495)=1/6$ of the positive results are expected to be genuine. The lesson is not that the test is useless; it is that accuracy and the background rate must be considered together.

\section*{Example 2: graph theory and a cheap network}
Four houses lie at the corners of a square. Roads along the four sides each cost 1 unit, while a diagonal road costs 1.5 units. The diagonal gives a faster route between two houses, but it is not needed to connect all four. A network designer who wants the cheapest connected system chooses the four sides, costing 4 units; a designer who values speed may add the diagonal, costing 5.5 units. Graph theory separates the questions ``Can everyone reach everyone?'' and ``What is the shortest route?''

\section*{Example 3: coding theory and a flipped bit}
Send the four messages 00, 01, 10, and 11 as 000000, 010101, 101010, and 111111. If one bit flips during transmission, the receiver can compare the received string with the four valid strings and choose the closest one. For instance, 010001 is closest to 010101, so the receiver restores 01. The price is extra data: six transmitted bits carry only two message bits. Reliability is purchased with redundancy.

\section*{Example 4: group theory and a clock}
On a twelve-hour clock, adding 5 hours and then adding 9 hours gives the same result as adding 14 hours, which is the same as adding 2 hours. The action ``add 5'' has an inverse: ``add 7,'' because together they add 12 and return to the starting hour. Clock arithmetic is a small example of a group, where actions can be combined, doing nothing is allowed, and every action can be undone.

\section*{Example 5: topology and holes}
A circle drawn on paper has one loop that cannot shrink to a point without leaving the circle. A filled disk has no such hole. Stretching the circle into an oval does not change this fact, but filling it in does. Topology ignores the exact measurements and records the feature that survives bending: the presence or absence of a hole.

\section*{Example 6: number theory and remainders}
To find the remainder of $3^{10}$ after division by 7, calculate the short cycle: $3^1=3$, $3^2=2$, $3^3=6$, $3^4=4$, $3^5=5$, and $3^6=1$ modulo 7. The cycle repeats every six powers, and $10$ leaves remainder 4 after division by 6. Therefore $3^{10}$ has the same remainder as $3^4$, namely 4. Large calculations become manageable because remainders repeat.

\section*{Important theorems and results}
\begin{enumerate}[leftmargin=*,itemsep=0.35em]
\item \textbf{Bayes' rule.} It explains the medical-test calculation by combining a prior rate with test evidence.
\item \textbf{Spanning-tree principle.} It explains why a connected network can be built without unnecessary cycles.
\item \textbf{Minimum-distance decoding.} It explains why redundancy can identify and correct sufficiently few transmission errors.
\item \textbf{Group axioms.} They explain why clock operations can be composed, reversed, and treated as symmetries.
\item \textbf{Homotopy invariance.} It explains why continuous deformation preserves the broad topological features in the example.
\item \textbf{Fermat's little theorem.} It explains why modular powers can repeat under a prime modulus.
\end{enumerate}

\section*{References}
\begin{thebibliography}{9}
\bibitem{examples_polya} G. P\'olya, \emph{How to Solve It}, Princeton University Press, 1945.
\bibitem{examples_velleman} D. J. Velleman, \emph{How to Prove It}, Cambridge University Press, 2019.
\end{thebibliography}

\chapter[Module 1: Probability]{Module 1: Probability}
Imagine a bus that arrives every ten minutes, but you reach the stop at a random time. If you have no timetable, your waiting time is anywhere from zero to ten minutes. The average wait is five minutes, not ten: you arrive at a typical point inside the interval.

Probability begins by listing possible outcomes and assigning weights to them. An event is simply a collection of outcomes. The rules are bookkeeping: the chance of an impossible event is zero, the chance of certainty is one, and the chance of either of two non-overlapping events is the sum of their chances. Conditional probability answers a different question: after learning that a person tested positive, how likely is illness? The answer depends not only on the test's accuracy but also on how common the illness was before testing.

\textbf{Worked example.} Suppose 1 in 100 people has a condition. A test catches 99 in 100 affected people but falsely flags 5 in 100 unaffected people. In a group of 10,000, about 100 are affected and 99 test positive. Of the 9,900 unaffected people, about 495 also test positive. So among positive results, roughly $99/(99+495)=1/6$ are genuine. The test can be good while a positive result is not certainty.

\textbf{Exercises.} (1) Why does doubling the number of tests not automatically double certainty? (2) What extra information would make the example more reliable? \emph{Answers:} (1) The base rate and false positives still matter. (2) A representative population and an independently measured test performance.

\section*{Important theorems and results}
\begin{enumerate}[leftmargin=*,itemsep=0.35em]
\item \textbf{Bayes' rule.} For events with positive conditioning probability, $P(A\mid B)=P(B\mid A)P(A)/P(B)$. It combines a prior rate with the evidence supplied by a test.
\end{enumerate}

\section*{References}
\begin{thebibliography}{9}
\bibitem{module_probability_ross} S. M. Ross, \emph{A First Course in Probability}, Pearson, 2019.
\bibitem{module_probability_williams} D. Williams, \emph{Weighing the Odds}, Cambridge University Press, 2001.
\end{thebibliography}


\chapter[Module 2: Statistics]{Module 2: Statistics}
Statistics starts with a boundary: we see a sample, not the whole world. A survey of 1,000 voters is not the election; a trial of 500 patients is not every patient. A statistic is useful only when we say how it was collected and how uncertain the conclusion is.

The mean is sensitive to extreme values, while the median is the middle observation after sorting. A confidence interval is a method for producing a range that would capture the target at a stated rate under repeated use; it is not a guarantee that one particular interval contains the truth. A hypothesis test asks whether the observed result would be unusual under a chosen baseline. It cannot by itself prove a cause, and a small probability under the baseline is not the probability that the baseline is false.

\textbf{Worked example.} A new teaching method raises a class average from 68 to 72. We compare it with another class, measure spread, check whether students were assigned fairly, and report the range of plausible improvements. If the new class was already stronger, the four-point difference may not be caused by the method. Design is as important as calculation.

\textbf{Exercises.} (1) Give one reason a large sample can still mislead. (2) Why use a comparison group? \emph{Answers:} (1) It can be biased or measured badly. (2) To estimate what would have happened without the intervention.

\section*{Important theorems and results}
\begin{enumerate}[leftmargin=*,itemsep=0.35em]
\item \textbf{Central limit theorem.} Under suitable assumptions, standardized sums of many observations approach a normal distribution. The approximation supports uncertainty intervals, but it does not repair biased sampling or poor measurement.

\item \textbf{Gauss--Markov theorem.} In a correctly specified linear model with the usual assumptions, least squares has the smallest variance among linear unbiased estimators.
\end{enumerate}

\section*{References}
\begin{thebibliography}{9}
\bibitem{module_statistics_casella} G. Casella and R. L. Berger, \emph{Statistical Inference}, Cengage, 2021.
\bibitem{module_statistics_freedman} D. A. Freedman, R. Pisani, and R. Purves, \emph{Statistics}, W. W. Norton, 2007.
\end{thebibliography}


\chapter[Module 3: Graph theory]{Module 3: Graph theory}
Draw dots for objects and lines for relationships. The dots may be towns, people, computers, or tasks; the lines may be roads, friendships, cables, or dependencies. A route is a sequence of joined dots. A connected network lets every place reach every other; a bridge is a connection whose removal breaks that property.

The important move is to forget irrelevant detail. A road map becomes a graph when only connectivity matters; a weighted graph keeps lengths or costs on the lines. A shortest-path method repeatedly improves the best known route. A spanning tree connects every dot using no unnecessary loop, useful when building a cheap network. Coloring assigns labels so joined dots differ, useful when conflicting activities cannot share a time slot.

\textbf{Worked example.} Four towns form a square, with one diagonal. If each road costs one unit, the diagonal shortens one pair's trip but is not needed to connect all towns. Removing it leaves a cycle; removing any one square side still leaves another route. The graph makes redundancy visible.

\textbf{Exercises.} (1) What changes when roads receive different costs? (2) Why might a network designer keep a loop? \emph{Answers:} (1) Fewest roads and cheapest route may differ. (2) A loop provides an alternative route when one link fails.

\section*{Important theorems and results}
\begin{enumerate}[leftmargin=*,itemsep=0.35em]
\item \textbf{Handshaking lemma.} The sum of the degrees of all vertices equals twice the number of edges. Consequently, every finite graph has an even number of odd-degree vertices.

\item \textbf{Spanning-tree result.} Every connected finite graph contains a spanning tree, and a minimum spanning tree gives a connected network of minimum total edge cost.
\end{enumerate}

\section*{References}
\begin{thebibliography}{9}
\bibitem{module_graph_west} D. B. West, \emph{Introduction to Graph Theory}, Pearson, 2001.
\bibitem{module_graph_diesteI} R. Diestel, \emph{Graph Theory}, Springer, 2017.
\end{thebibliography}


\chapter[Module 4: Information and coding]{Module 4: Information and coding}
Information theory asks how much uncertainty a message removes. A predictable symbol carries little new information; an unexpected symbol carries more. Compression works by replacing repeated or predictable patterns with shorter descriptions. It cannot make truly random data shorter on average without losing information.

Coding theory addresses the opposite danger: noise. Add redundancy in a pattern that makes likely errors recognizable. A simple parity check can detect an odd number of flipped bits. A stronger code can locate and correct errors because valid messages are deliberately spaced apart in the space of all possible strings.

\textbf{Worked example.} To send one of four symbols, use 00, 01, 10, or 11. If one bit may flip, use 000000, 010101, 101010, and 111111. The messages are far apart, so a received string with one error can be assigned to the nearest valid message. The cost is three times as many transmitted bits.

\textbf{Exercises.} (1) Why can redundancy improve reliability? (2) Why not add unlimited redundancy? \emph{Answers:} (1) It leaves clues about the intended message. (2) It consumes bandwidth, storage, time, and energy.

\section*{Important theorems and results}
\begin{enumerate}[leftmargin=*,itemsep=0.35em]
\item \textbf{Source-coding theorem.} For a memoryless source, reliable lossless compression cannot have average length below the entropy in the long run, while suitable codes can approach that bound.

\item \textbf{Distance bound.} A code with minimum Hamming distance $d$ detects at least $d-1$ errors and corrects at least $\lfloor(d-1)/2\rfloor$ errors.
\end{enumerate}

\section*{References}
\begin{thebibliography}{9}
\bibitem{module_information_cover} T. M. Cover and J. A. Thomas, \emph{Elements of Information Theory}, Wiley, 2006.
\bibitem{module_information_mackay} D. J. C. MacKay, \emph{Information Theory, Inference, and Learning Algorithms}, Cambridge University Press, 2003.
\end{thebibliography}


\chapter[Module 5: Group theory]{Module 5: Group theory}
Take a square. Rotating it by a quarter turn, half turn, or three-quarter turn, and flipping it in four ways, gives eight symmetries. Doing one after another gives another symmetry; doing nothing changes nothing; every move can be undone. These simple facts are the heart of a group.

The theory asks which features of an object are preserved by its symmetries and how complicated actions can be built from a few basic ones. A group need not be a collection of physical motions: rearranging a list, adding clock hours, or changing coordinates can have the same pattern of composition and reversal.

\textbf{Worked example.} On a clock, adding 5 hours and then 9 hours is the same as adding 14 hours, which is the same as adding 2 hours. The arithmetic wraps around after 12. This small group models repeating schedules and modular arithmetic used in cryptography.

\textbf{Exercises.} (1) What is the inverse of adding 5 on a 12-hour clock? (2) Why is order important for some symmetries? \emph{Answers:} (1) Adding 7. (2) Rotating then flipping a shape can differ from flipping then rotating.

\section*{Important theorems and results}
\begin{enumerate}[leftmargin=*,itemsep=0.35em]
\item \textbf{Lagrange's theorem.} The order of a finite subgroup divides the order of its finite ambient group. It turns a symmetry calculation into a divisibility test.
\end{enumerate}

\section*{References}
\begin{thebibliography}{9}
\bibitem{module_group_gallian} J. A. Gallian, \emph{Contemporary Abstract Algebra}, Cengage, 2017.
\bibitem{module_group_dummit} D. S. Dummit and R. M. Foote, \emph{Abstract Algebra}, Wiley, 2004.
\end{thebibliography}


\chapter[Module 6: Topology]{Module 6: Topology}
Topology studies what survives a continuous deformation. A rubber band can become a circle, but it cannot become two separate bands without being cut. A mug and a doughnut each have one tunnel, so topology treats them as the same broad kind of object; a sphere has none.

The basic features are connected pieces and holes of different kinds. A loop around a hole may fail to shrink to a point. Homology packages counts of such holes into algebraic summaries, while homotopy studies continuous deformations of paths and maps. These tools are powerful because a noisy drawing can change its exact lengths without changing its broad structure.

\textbf{Worked example.} A road network can be viewed as a one-dimensional shape. A closed alternative route is a loop. If a new road fills the loop by connecting every point across a region, the network's cycle structure changes. This is why topology can help analyze sensor and transportation networks.

\textbf{Exercises.} (1) Is a square topologically different from a circle? (2) What operation changes a mug's number of holes? \emph{Answers:} (1) No, if edges are allowed to bend and corners are not important. (2) Cutting, gluing, or drilling a hole.

\section*{Important theorems and results}
\begin{enumerate}[leftmargin=*,itemsep=0.35em]
\item \textbf{Homotopy invariance.} Continuous maps that can be deformed into one another induce the same map on homology. Therefore homology records features that survive continuous deformation.
\end{enumerate}

\section*{References}
\begin{thebibliography}{9}
\bibitem{module_topology_munkres} J. R. Munkres, \emph{Topology}, Pearson, 2000.
\bibitem{module_topology_hatcher} A. Hatcher, \emph{Algebraic Topology}, Cambridge University Press, 2002.
\end{thebibliography}


\chapter[Module 7: Calculus and dynamics]{Module 7: Calculus and dynamical systems}
A derivative measures an instantaneous rate: the speedometer reports how position changes at this moment. An integral adds many small contributions: distance can be recovered by adding small pieces of speed over time. A differential equation says how a quantity changes and asks what history follows from that rule.

Dynamical systems theory studies that history as a whole. It asks whether a system settles, cycles, escapes, or becomes highly sensitive to its starting point. A model may be deterministic while prediction is limited because measurements are never infinitely precise.

\textbf{Worked example.} If a population grows by 10 percent each year, its next value is 1.1 times its current value. Repeating the rule gives exponential growth. If resources limit growth, the growth rate falls as the population approaches a carrying level, producing a different long-term story. The equation expresses an assumption; data must test it.

\textbf{Exercises.} (1) What does a zero derivative mean at an instant? (2) Why can two close starting values later differ greatly? \emph{Answers:} (1) No instantaneous change. (2) The system may amplify small differences, as in chaos.

\section*{Important theorems and results}
\begin{enumerate}[leftmargin=*,itemsep=0.35em]
\item \textbf{Fundamental theorem of calculus.} Differentiation and integration are inverse operations under the usual continuity assumptions.

\item \textbf{Linearization principle.} Near a nondegenerate equilibrium, the eigenvalues of the linearized system determine local stability; zero real parts require further analysis.
\end{enumerate}

\section*{References}
\begin{thebibliography}{9}
\bibitem{module_calculus_stewart} J. Stewart, \emph{Calculus}, Cengage, 2016.
\bibitem{module_dynamics_strogatz} S. H. Strogatz, \emph{Nonlinear Dynamics and Chaos}, Westview Press, 2015.
\end{thebibliography}


\chapter[Module 8: Number theory]{Module 8: Number theory}
Whole numbers cannot be smoothly adjusted: divisibility either works or it does not. Prime numbers are the basic building blocks because every positive whole number factors into primes in only one way apart from order. Congruence arithmetic keeps only remainders, which is why clock arithmetic is a useful model.

\textbf{Worked example.} To compute a large power modulo 7, reduce after each multiplication. The powers of 3 cycle: $3,2,6,4,5,1$. The cycle makes enormous calculations manageable. Public-key cryptography uses related facts: multiplying two secret primes is easy, while recovering them from their product can be difficult for suitably large numbers.

\textbf{Exercises.} (1) What is $17$ modulo $5$? (2) Why are primes useful for building secret keys? \emph{Answers:} (1) 2. (2) Their multiplication is easy to perform but reversing it can be computationally hard.

\section*{Important theorems and results}
\begin{enumerate}[leftmargin=*,itemsep=0.35em]
\item \textbf{Fundamental theorem of arithmetic.} Every integer greater than one has a factorization into primes that is unique up to order.

\item \textbf{Fermat's little theorem.} If $p$ is prime and $p$ does not divide $a$, then $a^{p-1}\equiv1\pmod p$. This explains the repeating powers used in the example.
\end{enumerate}

\section*{References}
\begin{thebibliography}{9}
\bibitem{module_number_hardy} G. H. Hardy and E. M. Wright, \emph{An Introduction to the Theory of Numbers}, Oxford University Press, 2008.
\bibitem{module_number_niven} I. Niven, H. S. Zuckerman, and H. L. Montgomery, \emph{An Introduction to the Theory of Numbers}, Wiley, 2006.
\end{thebibliography}


\chapter[Module 9: Approximation]{Module 9: Approximation}
A complicated function can be replaced locally by a line or more accurately by a polynomial. The replacement is valuable only when we know where it works and how far it may be from the original. This is the difference between a controlled approximation and a guess.

\textbf{Worked example.} Near zero, $\sin x$ is close to $x$ when $x$ is measured in radians. At $x=0.1$, this is a good shortcut; at $x=2$, it is not. Adding the next correction, $-x^3/6$, improves the estimate. A numerical method should report both its answer and a reason to trust its accuracy.

\textbf{Exercises.} (1) Why does an approximation need a range of validity? (2) What is the danger of adding more terms blindly? \emph{Answers:} (1) Errors can grow away from the region where the estimate was built. (2) A series may not converge or rounding errors may dominate.

\section*{Important theorems and results}
\begin{enumerate}[leftmargin=*,itemsep=0.35em]
\item \textbf{Taylor's theorem.} A smooth function equals its Taylor polynomial plus a remainder controlled by a higher derivative. The remainder is what turns a numerical guess into an error-bounded approximation.

\item \textbf{Weierstrass approximation theorem.} Every continuous real-valued function on a closed bounded interval can be uniformly approximated by polynomials.
\end{enumerate}

\section*{References}
\begin{thebibliography}{9}
\bibitem{module_approximation_trefethen} L. N. Trefethen, \emph{Approximation Theory and Approximation Practice}, SIAM, 2013.
\bibitem{module_approximation_cheney} E. W. Cheney, \emph{Introduction to Approximation Theory}, AMS Chelsea, 1982.
\end{thebibliography}


\chapter[Module 10: Computation and proof]{Module 10: Computation and proof}
An algorithm is a finite recipe. Automata theory studies machines with restricted memory; computation theory asks what any algorithm can do; proof theory studies formal chains of reasons. These are different questions: a problem can be solvable in principle but too slow in practice, or a statement can be true without a particular formal system being able to prove it.

\textbf{Worked example.} A spelling checker can recognize whether a word follows a finite list of allowed patterns. It cannot, using only that finite memory, remember an arbitrary number of opening parentheses and match them all. A stack adds the needed memory. The model explains why different programming tasks require different resources.

\textbf{Exercises.} (1) What does an algorithm require that a vague wish does not? (2) Why is ``not yet computed'' different from ``cannot be computed''? \emph{Answers:} (1) Explicit, repeatable steps. (2) More time may solve the first; no algorithm can solve the second.

\section*{Important theorems and results}
\begin{enumerate}[leftmargin=*,itemsep=0.35em]
\item \textbf{Halting theorem.} No algorithm can correctly decide for every program and input whether that program eventually halts.

\item \textbf{Cook--Levin theorem.} Boolean satisfiability is NP-complete, making it a central comparison problem for computational difficulty.
\end{enumerate}

\section*{References}
\begin{thebibliography}{9}
\bibitem{module_computation_sipser} M. Sipser, \emph{Introduction to the Theory of Computation}, Cengage, 2013.
\bibitem{module_computation_velleman} D. J. Velleman, \emph{How to Prove It}, Cambridge University Press, 2019.
\end{thebibliography}


\chapter[Module 11]{Module 11: Game, queue, and control theory}
Game theory studies decisions whose results depend on other decision-makers. Queue theory studies congestion caused by random arrivals and limited service. Control theory studies feedback: measuring a system, comparing it with a target, and correcting it. All three ask how local choices affect a changing whole.

\textbf{Worked example.} Two hospitals can each either share beds with the other or reserve them. Sharing helps patients when demand differs, but each hospital may prefer to reserve beds during a shortage. A model lists choices and payoffs, then asks whether cooperation is stable. If arrivals are uncertain, queue calculations show the cost of waiting; feedback can redirect ambulances as capacity changes.

\textbf{Exercises.} (1) Why can a locally sensible choice harm a network? (2) What makes feedback unstable? \emph{Answers:} (1) Other agents and congestion react to it. (2) Corrections may be too large, too late, or based on noisy measurements.

\section*{Important theorems and results}
\begin{enumerate}[leftmargin=*,itemsep=0.35em]
\item \textbf{Nash's existence theorem.} Every finite game has at least one mixed-strategy Nash equilibrium.

\item \textbf{Little's law.} In a stable long-run queue, the average number in the system equals the arrival rate multiplied by the average time spent in the system.

\item \textbf{Feedback stability.} A feedback controller is useful only when the closed-loop dynamics remain stable; increasing gain can reduce error in one regime and cause oscillation or divergence in another.
\end{enumerate}

\section*{References}
\begin{thebibliography}{9}
\bibitem{module_decision_osborne} M. J. Osborne and A. Rubinstein, \emph{A Course in Game Theory}, MIT Press, 1994.
\bibitem{module_queue_gross} D. Gross, J. F. Shortle, J. M. Thompson, and C. M. Harris, \emph{Fundamentals of Queueing Theory}, Wiley, 2018.
\bibitem{module_control_khalil} H. K. Khalil, \emph{Nonlinear Systems}, Prentice Hall, 2002.
\end{thebibliography}


\chapter[Module 12: Algebraic geometry]{Module 12: Algebraic geometry}
An equation such as $x^2+y^2=1$ describes a shape: a circle. Changing the equation changes the shape, and changing the number system changes which points count. Algebraic geometry studies this two-way relationship between formulas and their solution sets.

The algebra is useful because a shape may be hard to draw but easy to describe through the operations on its coordinates. Ideals collect equations that vanish on the same set; elimination removes unwanted coordinates; schemes extend the viewpoint to arithmetic questions involving primes and finite number systems. These advanced tools are not magic new shapes: they are ways to remember local information that an ordinary picture loses.

\textbf{Worked example.} The equations $x+y=3$ and $x-y=1$ describe two lines. Adding them removes $y$ and gives $2x=4$, so $x=2$ and then $y=1$. Elimination theory generalizes this small act to many polynomial equations and many variables.

\textbf{Exercises.} (1) What does a solution point mean geometrically? (2) Why eliminate a variable? \emph{Answers:} (1) It is a point satisfying every equation. (2) To find a condition or projection involving only the quantities of interest.

\section*{Important theorems and results}
\begin{enumerate}[leftmargin=*,itemsep=0.35em]
\item \textbf{Hilbert's Nullstellensatz.} Over an algebraically closed field, polynomial equations and radical ideals describe the same affine solution sets.

\item \textbf{B\'ezout's theorem.} Projective plane curves with no common component meet in a number of points counted with multiplicity equal to the product of their degrees.
\end{enumerate}

\section*{References}
\begin{thebibliography}{9}
\bibitem{module_geometry_cox} D. Cox, J. Little, and D. O'Shea, \emph{Ideals, Varieties, and Algorithms}, Springer, 2015.
\bibitem{module_geometry_hartshorne} R. Hartshorne, \emph{Algebraic Geometry}, Springer, 1977.
\end{thebibliography}


\backmatter
\appendix
\chapter{Theory dependency and application map}
This table summarises the relationships explained in the individual chapters. The dependency column names the background ideas needed to begin studying a theory; the final column names theories that reuse its methods or results. These are relationships, not strict prerequisites: advanced mathematics is interconnected, and a course may introduce the subjects in a different order.

\small
\sloppy
\setlength{\tabcolsep}{3pt}
\begin{longtable}{>{\raggedright\arraybackslash}p{0.22\textwidth}>{\raggedright\arraybackslash}p{0.28\textwidth}>{\raggedright\arraybackslash}p{0.29\textwidth}}
\textbf{Theory name} & \textbf{Depends on} & \textbf{Helps other theories}\\
\hline
\endfirsthead
\textbf{Theory name} & \textbf{Depends on} & \textbf{Helps other theories}\\
\hline
\endhead
Algebraic geometry & algebra, topology, calculus, and the geometric objects introduced in the chapter & its constructions support topology, number theory, representation theory, and mathematical physics\\
\hline
Almgren--Pitts min--max theory & sets, continuity, topology, and geometric intuition & its invariants support homology, homotopy, geometry, dynamical systems, data analysis, and manifold theory\\
\hline
Approximation theory & calculus, limits, linear algebra, and real or complex analysis & its analytic tools support differential equations, probability, physics, numerical computation, and statistics\\
\hline
Arakelov theory & integer arithmetic, algebra, and proof; advanced chapters also require fields or complex analysis & its arithmetic tools support cryptography, algebraic geometry, coding theory, and arithmetic statistics\\
\hline
Asymptotic theory & calculus, limits, linear algebra, and real or complex analysis & its analytic tools support differential equations, probability, physics, numerical computation, and statistics\\
\hline
Automata theory & sets, functions, and the chapter's own definitions & its objects and methods are reused by neighbouring areas of mathematics and applications\\
\hline
Bass--Serre theory & sets, functions, algebra, and proof; the chapter then introduces its own operations and axioms & its structures support representation theory, algebraic geometry, number theory, coding, topology, logic, and computer science\\
\hline
Bifurcation theory & calculus, differential equations, linear algebra, and phase-plane reasoning & its ideas support chaos, ergodic theory, control, mathematical biology, and physics\\
\hline
Braid theory & sets, functions, algebra, and proof; the chapter then introduces its own operations and axioms & its structures support representation theory, algebraic geometry, number theory, coding, topology, logic, and computer science\\
\hline
Brill--Noether theory & algebra, topology, calculus, and the geometric objects introduced in the chapter & its constructions support topology, number theory, representation theory, and mathematical physics\\
\hline
Catastrophe theory & calculus, differential equations, linear algebra, and phase-plane reasoning & its ideas support chaos, ergodic theory, control, mathematical biology, and physics\\
\hline
Category theory & sets, functions, algebra, and proof; the chapter then introduces its own operations and axioms & its structures support representation theory, algebraic geometry, number theory, coding, topology, logic, and computer science\\
\hline
Chaos theory & calculus, differential equations, linear algebra, and phase-plane reasoning & its ideas support chaos, ergodic theory, control, mathematical biology, and physics\\
\hline
Character theory & sets, functions, algebra, and proof; the chapter then introduces its own operations and axioms & its structures support representation theory, algebraic geometry, number theory, coding, topology, logic, and computer science\\
\hline
Choquet theory & sets, functions, and the chapter's own definitions & its objects and methods are reused by neighbouring areas of mathematics and applications\\
\hline
Class field theory & sets, functions, algebra, and proof; the chapter then introduces its own operations and axioms & its structures support representation theory, algebraic geometry, number theory, coding, topology, logic, and computer science\\
\hline
Cobordism theory & sets, continuity, topology, and geometric intuition & its invariants support homology, homotopy, geometry, dynamical systems, data analysis, and manifold theory\\
\hline
Coding theory & probability, linear algebra, calculus, and mathematical modelling & its models support statistics, engineering, information theory, optimization, economics, and scientific prediction\\
\hline
Cohomology theory & sets, continuity, topology, and geometric intuition & its invariants support homology, homotopy, geometry, dynamical systems, data analysis, and manifold theory\\
\hline
Continuum theory & sets, continuity, topology, and geometric intuition & its invariants support homology, homotopy, geometry, dynamical systems, data analysis, and manifold theory\\
\hline
Control theory & probability, linear algebra, calculus, and mathematical modelling & its models support statistics, engineering, information theory, optimization, economics, and scientific prediction\\
\hline
Deformation theory & algebra, topology, calculus, and the geometric objects introduced in the chapter & its constructions support topology, number theory, representation theory, and mathematical physics\\
\hline
Dempster--Shafer theory & probability, linear algebra, calculus, and mathematical modelling & its models support statistics, engineering, information theory, optimization, economics, and scientific prediction\\
\hline
Dimension theory & sets, continuity, topology, and geometric intuition & its invariants support homology, homotopy, geometry, dynamical systems, data analysis, and manifold theory\\
\hline
Distribution theory & calculus, limits, linear algebra, and real or complex analysis & its analytic tools support differential equations, probability, physics, numerical computation, and statistics\\
\hline
Dynamical systems theory & calculus, differential equations, linear algebra, and phase-plane reasoning & its ideas support chaos, ergodic theory, control, mathematical biology, and physics\\
\hline
Elimination theory & algebra, topology, calculus, and the geometric objects introduced in the chapter & its constructions support topology, number theory, representation theory, and mathematical physics\\
\hline
Ergodic theory & sets, functions, and the chapter's own definitions & its objects and methods are reused by neighbouring areas of mathematics and applications\\
\hline
Extremal graph theory & sets, counting, combinatorics, and algorithms & its network and counting methods support optimization, probability, computer science, and network science\\
\hline
Field theory & sets, functions, algebra, and proof; the chapter then introduces its own operations and axioms & its structures support representation theory, algebraic geometry, number theory, coding, topology, logic, and computer science\\
\hline
Galois theory & sets, functions, algebra, and proof; the chapter then introduces its own operations and axioms & its structures support representation theory, algebraic geometry, number theory, coding, topology, logic, and computer science\\
\hline
Game theory & probability, linear algebra, calculus, and mathematical modelling & its models support statistics, engineering, information theory, optimization, economics, and scientific prediction\\
\hline
Graph theory & sets, counting, combinatorics, and algorithms & its network and counting methods support optimization, probability, computer science, and network science\\
\hline
Group theory & sets, functions, algebra, and proof; the chapter then introduces its own operations and axioms & its structures support representation theory, algebraic geometry, number theory, coding, topology, logic, and computer science\\
\hline
Hodge theory & algebra, topology, calculus, and the geometric objects introduced in the chapter & its constructions support topology, number theory, representation theory, and mathematical physics\\
\hline
Homology theory & sets, continuity, topology, and geometric intuition & its invariants support homology, homotopy, geometry, dynamical systems, data analysis, and manifold theory\\
\hline
Homotopy theory & sets, continuity, topology, and geometric intuition & its invariants support homology, homotopy, geometry, dynamical systems, data analysis, and manifold theory\\
\hline
Ideal theory & sets, functions, algebra, and proof; the chapter then introduces its own operations and axioms & its structures support representation theory, algebraic geometry, number theory, coding, topology, logic, and computer science\\
\hline
Index theory & sets, continuity, topology, and geometric intuition & its invariants support homology, homotopy, geometry, dynamical systems, data analysis, and manifold theory\\
\hline
Information theory & probability, linear algebra, calculus, and mathematical modelling & its models support statistics, engineering, information theory, optimization, economics, and scientific prediction\\
\hline
Intersection theory & algebra, topology, calculus, and the geometric objects introduced in the chapter & its constructions support topology, number theory, representation theory, and mathematical physics\\
\hline
Invariant theory & sets, functions, algebra, and proof; the chapter then introduces its own operations and axioms & its structures support representation theory, algebraic geometry, number theory, coding, topology, logic, and computer science\\
\hline
Iwasawa theory & integer arithmetic, algebra, and proof; advanced chapters also require fields or complex analysis & its arithmetic tools support cryptography, algebraic geometry, coding theory, and arithmetic statistics\\
\hline
K-theory & sets, functions, and the chapter's own definitions & its objects and methods are reused by neighbouring areas of mathematics and applications\\
\hline
Knot theory & sets, continuity, topology, and geometric intuition & its invariants support homology, homotopy, geometry, dynamical systems, data analysis, and manifold theory\\
\hline
L-theory & sets, functions, and the chapter's own definitions & its objects and methods are reused by neighbouring areas of mathematics and applications\\
\hline
Lattice theory & sets, functions, algebra, and proof; the chapter then introduces its own operations and axioms & its structures support representation theory, algebraic geometry, number theory, coding, topology, logic, and computer science\\
\hline
Lie theory & sets, functions, algebra, and proof; the chapter then introduces its own operations and axioms & its structures support representation theory, algebraic geometry, number theory, coding, topology, logic, and computer science\\
\hline
M-theory & algebra, topology, calculus, and the geometric objects introduced in the chapter & its constructions support topology, number theory, representation theory, and mathematical physics\\
\hline
Measure theory & calculus, limits, linear algebra, and real or complex analysis & its analytic tools support differential equations, probability, physics, numerical computation, and statistics\\
\hline
Model theory & sets, functions, algebra, and proof; the chapter then introduces its own operations and axioms & its structures support representation theory, algebraic geometry, number theory, coding, topology, logic, and computer science\\
\hline
Module theory & sets, functions, algebra, and proof; the chapter then introduces its own operations and axioms & its structures support representation theory, algebraic geometry, number theory, coding, topology, logic, and computer science\\
\hline
Morse theory & sets, continuity, topology, and geometric intuition & its invariants support homology, homotopy, geometry, dynamical systems, data analysis, and manifold theory\\
\hline
Nevanlinna theory & calculus, limits, linear algebra, and real or complex analysis & its analytic tools support differential equations, probability, physics, numerical computation, and statistics\\
\hline
Number theory & integer arithmetic, algebra, and proof; advanced chapters also require fields or complex analysis & its arithmetic tools support cryptography, algebraic geometry, coding theory, and arithmetic statistics\\
\hline
Obstruction theory & sets, continuity, topology, and geometric intuition & its invariants support homology, homotopy, geometry, dynamical systems, data analysis, and manifold theory\\
\hline
Operator theory & calculus, limits, linear algebra, and real or complex analysis & its analytic tools support differential equations, probability, physics, numerical computation, and statistics\\
\hline
Order theory & sets, functions, algebra, and proof; the chapter then introduces its own operations and axioms & its structures support representation theory, algebraic geometry, number theory, coding, topology, logic, and computer science\\
\hline
Percolation theory & probability, linear algebra, calculus, and mathematical modelling & its models support statistics, engineering, information theory, optimization, economics, and scientific prediction\\
\hline
Perturbation theory & calculus, limits, linear algebra, and real or complex analysis & its analytic tools support differential equations, probability, physics, numerical computation, and statistics\\
\hline
Probability theory & probability, linear algebra, calculus, and mathematical modelling & its models support statistics, engineering, information theory, optimization, economics, and scientific prediction\\
\hline
Proof theory & sets, functions, algebra, and proof; the chapter then introduces its own operations and axioms & its structures support representation theory, algebraic geometry, number theory, coding, topology, logic, and computer science\\
\hline
Queue theory & probability, linear algebra, calculus, and mathematical modelling & its models support statistics, engineering, information theory, optimization, economics, and scientific prediction\\
\hline
Ramsey theory & sets, counting, combinatorics, and algorithms & its network and counting methods support optimization, probability, computer science, and network science\\
\hline
Random matrix theory & probability, linear algebra, calculus, and mathematical modelling & its models support statistics, engineering, information theory, optimization, economics, and scientific prediction\\
\hline
Representation theory & sets, functions, algebra, and proof; the chapter then introduces its own operations and axioms & its structures support representation theory, algebraic geometry, number theory, coding, topology, logic, and computer science\\
\hline
Ring theory & sets, functions, algebra, and proof; the chapter then introduces its own operations and axioms & its structures support representation theory, algebraic geometry, number theory, coding, topology, logic, and computer science\\
\hline
Scheme theory & sets, functions, algebra, and proof; the chapter then introduces its own operations and axioms & its structures support representation theory, algebraic geometry, number theory, coding, topology, logic, and computer science\\
\hline
Semigroup theory & sets, functions, algebra, and proof; the chapter then introduces its own operations and axioms & its structures support representation theory, algebraic geometry, number theory, coding, topology, logic, and computer science\\
\hline
Set theory & sets, functions, algebra, and proof; the chapter then introduces its own operations and axioms & its structures support representation theory, algebraic geometry, number theory, coding, topology, logic, and computer science\\
\hline
Shape theory & sets, continuity, topology, and geometric intuition & its invariants support homology, homotopy, geometry, dynamical systems, data analysis, and manifold theory\\
\hline
Sheaf theory & sets, functions, algebra, and proof; the chapter then introduces its own operations and axioms & its structures support representation theory, algebraic geometry, number theory, coding, topology, logic, and computer science\\
\hline
Sieve theory & integer arithmetic, algebra, and proof; advanced chapters also require fields or complex analysis & its arithmetic tools support cryptography, algebraic geometry, coding theory, and arithmetic statistics\\
\hline
Singularity theory & algebra, topology, calculus, and the geometric objects introduced in the chapter & its constructions support topology, number theory, representation theory, and mathematical physics\\
\hline
Soliton theory & calculus, limits, linear algebra, and real or complex analysis & its analytic tools support differential equations, probability, physics, numerical computation, and statistics\\
\hline
Spectral theory & calculus, limits, linear algebra, and real or complex analysis & its analytic tools support differential equations, probability, physics, numerical computation, and statistics\\
\hline
String theory & sets, functions, algebra, and proof; the chapter then introduces its own operations and axioms & its structures support representation theory, algebraic geometry, number theory, coding, topology, logic, and computer science\\
\hline
Sturm--Liouville theory & calculus, limits, linear algebra, and real or complex analysis & its analytic tools support differential equations, probability, physics, numerical computation, and statistics\\
\hline
Surgery theory & sets, continuity, topology, and geometric intuition & its invariants support homology, homotopy, geometry, dynamical systems, data analysis, and manifold theory\\
\hline
Teichmüller theory & sets, continuity, topology, and geometric intuition & its invariants support homology, homotopy, geometry, dynamical systems, data analysis, and manifold theory\\
\hline
Theory of computation & sets, functions, and the chapter's own definitions & its objects and methods are reused by neighbouring areas of mathematics and applications\\
\hline
Theory of equations & sets, functions, and the chapter's own definitions & its objects and methods are reused by neighbouring areas of mathematics and applications\\
\hline
Theory of statistics & probability, linear algebra, calculus, and mathematical modelling & its models support statistics, engineering, information theory, optimization, economics, and scientific prediction\\
\hline
Topology & sets, continuity, topology, and geometric intuition & its invariants support homology, homotopy, geometry, dynamical systems, data analysis, and manifold theory\\
\hline
Topos theory & sets, functions, algebra, and proof; the chapter then introduces its own operations and axioms & its structures support representation theory, algebraic geometry, number theory, coding, topology, logic, and computer science\\
\hline
Transcendental number theory & integer arithmetic, algebra, and proof; advanced chapters also require fields or complex analysis & its arithmetic tools support cryptography, algebraic geometry, coding theory, and arithmetic statistics\\
\hline
Twistor theory & algebra, topology, calculus, and the geometric objects introduced in the chapter & its constructions support topology, number theory, representation theory, and mathematical physics\\
\hline
Type theory & sets, functions, algebra, and proof; the chapter then introduces its own operations and axioms & its structures support representation theory, algebraic geometry, number theory, coding, topology, logic, and computer science\\
\hline
Wheel theory & sets, functions, algebra, and proof; the chapter then introduces its own operations and axioms & its structures support representation theory, algebraic geometry, number theory, coding, topology, logic, and computer science\\
\hline
\end{longtable}
\setlength{\tabcolsep}{6pt}
\fussy
\normalsize
\section*{References}
The map summarizes the dependencies and applications explained in the chapters; the individual chapter bibliographies provide the subject-specific sources.
\begin{thebibliography}{9}
\bibitem{map_maclane} S. Mac Lane, \emph{Categories for the Working Mathematician}, Springer, 1998.
\bibitem{map_hatcher} A. Hatcher, \emph{Algebraic Topology}, Cambridge University Press, 2002.
\end{thebibliography}
\chapter*{The student's toolbox: difficult words made ordinary}
This short dictionary is meant to be used while reading. The point is not to avoid precise ideas, but to give each precise idea a familiar handle.

\begin{description}[leftmargin=3.0cm,style=nextline]
\item[Algebra] A way of writing rules for combining and changing quantities or objects. It is not limited to school equations.
\item[Algorithm] A recipe that always says what step to take next, such as the instructions for sorting a list.
\item[Continuous] Changing without a sudden jump or tear, like moving a cup across a table.
\item[Dimension] How many independent directions are needed to move around an object. A line has one, a sheet two, and a room three.
\item[Equation] A statement that two descriptions have the same value. Solving it means finding the situations where that statement is true.
\item[Function] A dependable input-to-output rule. A temperature converter and a vending machine are everyday examples.
\item[Geometry] The study of shape, position, distance, and motion.
\item[Group] A collection of reversible actions that can be performed one after another. Rotations of a cube are a model example.
\item[Limit] The value a process gets closer and closer to, even if it never arrives there in a finite number of steps.
\item[Matrix] A rectangular table of numbers used to record many inputs and outputs at once.
\item[Model] A simplified description built to answer a particular question. Every model leaves something out.
\item[Operator] A rule that acts on another object, often a function or a list of numbers.
\item[Proof] A chain of small, checkable reasons that forces a conclusion to be true under stated rules.
\item[Set] A collection of objects considered together. The objects can be numbers, places, choices, or other collections.
\item[Symmetry] A change that leaves the important appearance or structure unchanged.
\item[Topology] The study of features that survive bending and stretching, such as connected pieces and holes.
\item[Vector] A quantity with several components, often used to describe a direction, a position, or a list of measurements.
\item[Wheel] An algebraic system designed to let division by zero be written as a total operation, with special rules that prevent it from behaving like ordinary field arithmetic.
\end{description}

\chapter*{How to study a theory without getting lost}
For each chapter, make a four-line note in your own words:
\begin{enumerate}
\item Draw or describe the real object: a queue, a knot, a message, a vibrating string, or a set of choices.
\item Name the change that matters: noise, motion, uncertainty, growth, symmetry, or competition.
\item Say what the theory measures or preserves. This is usually the heart of the subject.
\item Write one decision it improves: decode, predict, classify, schedule, control, compare, or prove impossible.
\end{enumerate}
Then explain the theory to another person without using its name. If the explanation still works, you understand the idea rather than merely recognizing the vocabulary. Only after that should you return to the equations or formal definitions in a textbook.

\chapter*{A map of the connections}
The theories are easier to remember as families. \emph{Probability} describes uncertainty, while \emph{statistics} learns from observations of that uncertainty. \emph{Graph theory} describes connections, while \emph{queue theory}, \emph{game theory}, and \emph{control theory} use connections to study waiting, competition, and correction. \emph{Topology} asks which features survive bending; \emph{geometry} keeps measurements too; \emph{algebra} records operations and symmetry. \emph{Information theory} asks how messages can be represented, and \emph{coding theory} makes them survive noise. \emph{Number theory} studies whole-number structure, while algebraic geometry joins equations to shapes.

The same pattern appears again and again: choose a simple representation, find a feature that remains reliable, and use it to make a prediction. That is the common thread tying this very large list together.

\chapter*{Where to go next}
This book is a first conversation, not a credential. The following standard references are reliable next steps; a student should use the one matching their preparation rather than trying to read them all.

\begin{description}[leftmargin=3.5cm,style=nextline]
\item[Elementary foundations] Courant and Robbins, \emph{What Is Mathematics?}; Ian Stewart, \emph{Concepts of Modern Mathematics}.
\item[Proof and discrete mathematics] Daniel Velleman, \emph{How to Prove It}; Richard Johnsonbaugh, \emph{Discrete Mathematics}.
\item[Probability and statistics] David Williams, \emph{Weighing the Odds}; David Freedman, Robert Pisani, and Roger Purves, \emph{Statistics}.
\item[Calculus and differential equations] Michael Spivak, \emph{Calculus}; Walter Strauss, \emph{Partial Differential Equations}.
\item[Algebra and number theory] Joseph Gallian, \emph{Contemporary Abstract Algebra}; G. H. Hardy and E. M. Wright, \emph{An Introduction to the Theory of Numbers}.
\item[Topology and geometry] Sidney Morris, \emph{Topology Without Tears}; James Munkres, \emph{Topology}.
\item[Computation and information] Michael Sipser, \emph{Introduction to the Theory of Computation}; Thomas Cover and Joy Thomas, \emph{Elements of Information Theory}.
\item[Advanced subjects] Serge Lang, \emph{Algebra}; Robin Hartshorne, \emph{Algebraic Geometry}; John Lee, \emph{Introduction to Smooth Manifolds}. These are reference texts, not assumed prerequisites for this book.
\end{description}

The source list is useful for names, but it is not a textbook and should not be treated as an authority for every application sentence. In a classroom or publication, each advanced chapter should be checked against a specialist reference before being used as course material.


\chapter*{A final thought}
The theories in this book do not form a ladder from ``easy'' to ``hard.'' They form a network. A graph can describe a symmetry; probability can live on a measure space; topology can guide an algorithm; algebra can decode a message. Mathematics becomes powerful when one viewpoint is carried into another place.

The recurring habit is simple: identify what may change, identify what must remain, and choose a language that makes the invariant visible. That is why theories made for apparently abstract questions so often return in phones, bridges, medicines, software, and the night sky.

\chapter*{Scope note and source}
The chapter titles follow the list at \url{https://en.wikipedia.org/wiki/List_of_mathematical_theories}, which is available under the Creative Commons Attribution--ShareAlike 4.0 License. The explanations and examples in this manuscript are original popular expositions; they are not copied from the linked article. Because mathematical fields overlap, several chapters naturally point to one another.

\end{document}
